% Generated by `lake env lean --run scripts/BlueprintGraph.lean`; do not edit.

\chapter{The statement}

\section{The Hodge conjecture}\label{sec:HodgeConjecture.Statement}
\noindent\emph{File:} \href{https://github.com/Paul-Lez/HodgeConjecture/blob/main/HodgeConjecture/Statement.lean}{\texttt{HodgeConjecture/Statement.lean}}.

This file states (but does not prove!) the Hodge conjecture.

The Hodge conjecture says that, for every nonsingular complex projective variety \texttt{X} and natural
number \texttt{p}, every rational Hodge class of degree \texttt{2p} on \texttt{X} is a rational linear combination of
classes of algebraic subvarieties of \texttt{X} of codimension \texttt{p}.

\paragraph{References}

\href{https://www.claymath.org/wp-content/uploads/2022/02/MPPc.pdf}{P. Deligne, *The Hodge Conjecture*}

\begin{conjecture}[Hodge conjecture]\label{conj:HodgeConjecture}
\lean{HodgeConjecture}\leanok
\uses{def:AlgebraicGeometry.ComplexPoint.algebraicCycleClassSpan, def:AlgebraicGeometry.ComplexPoint.hodgeClasses}
Statement of the \textbf{Hodge conjecture}.

For every nonsingular complex projective variety \texttt{X} and natural number \texttt{p}, every rational Hodge
class of degree \texttt{2p} on \texttt{X} is a rational linear combination of classes of algebraic subvarieties of
\texttt{X} of codimension \texttt{p}. The algebraic subspace is \texttt{algebraicCycleClassSpan}, the span of the
component classes of \texttt{CycleComponentSheafClass} indexed by the points of codimension \texttt{p}.

\noindent\emph{Lean:} \texttt{def HodgeConjecture}, \href{https://github.com/Paul-Lez/HodgeConjecture/blob/main/HodgeConjecture/Statement.lean#L37}{\texttt{HodgeConjecture/Statement.lean:37}}.
\end{conjecture}

\chapter{Definitions: Algebra}

\section{A coefficient field inside the complex numbers}\label{sec:HodgeConjecture.Definitions.Algebra.FieldToComplex}
\noindent\emph{File:} \href{https://github.com/Paul-Lez/HodgeConjecture/blob/main/HodgeConjecture/Definitions/Algebra/FieldToComplex.lean}{\texttt{HodgeConjecture/Definitions/Algebra/FieldToComplex.lean}}.

A field \texttt{K} with an algebra map to \texttt{ℂ} embeds \texttt{K}-linearly via \texttt{Algebra.linearMap}, and that
embedding splits: an
injective linear map of vector spaces over a field has a left inverse. Choosing one gives a
\texttt{K}-linear retraction \texttt{ℂ → K}, which is what lets a complex cohomology class be pushed back to
\texttt{K}-coefficients.

Multiplication by a scalar, as an additive endomorphism, is recorded here alongside it. It is
kept free of the geometry so the constant-sheaf and cohomology developments can use it without
carrying the algebra hypothesis.

\begin{definition}[fieldScalarAddHom]\label{def:fieldScalarAddHom}
\lean{fieldScalarAddHom}\leanok
Multiplication by a rational scalar as an additive endomorphism of \texttt{K}.

\noindent\emph{Lean:} \texttt{def fieldScalarAddHom}, \href{https://github.com/Paul-Lez/HodgeConjecture/blob/main/HodgeConjecture/Definitions/Algebra/FieldToComplex.lean#L45}{\texttt{HodgeConjecture/Definitions/Algebra/FieldToComplex.lean:45}}.
\end{definition}

\section{The linear dual of a complex of vector spaces}\label{sec:HodgeConjecture.Definitions.Algebra.Homology.LinearDual}
\noindent\emph{File:} \href{https://github.com/Paul-Lez/HodgeConjecture/blob/main/HodgeConjecture/Definitions/Algebra/Homology/LinearDual.lean}{\texttt{HodgeConjecture/Definitions/Algebra/Homology/LinearDual.lean}}.

This file dualises complexes of vector spaces. The reversed algebraic-dual short complex of a
short complex has homology canonically linearly equivalent to the dual of the original homology:
this is the universal-coefficient identification, and no finite-dimensionality hypothesis is
needed.

Applying the dual degreewise turns a nonnegatively graded chain complex into a cochain complex,
whose degree-\texttt{n} short complex is the reversed dual of the original one. A chain-homotopy
equivalence therefore induces, contravariantly, a linear equivalence on dual homology.

Duality is contravariantly functorial for maps, isomorphisms, chain homotopies and
chain-homotopy equivalences, so all of these transport to the dual cochain complex.

\begin{definition}[CategoryTheory.ShortComplex.linearDual]\label{def:CategoryTheory.ShortComplex.linearDual}
\lean{CategoryTheory.ShortComplex.linearDual}\leanok
The reversed algebraic-dual short complex.

\noindent\emph{Lean:} \texttt{abbrev CategoryTheory.ShortComplex.linearDual}, \href{https://github.com/Paul-Lez/HodgeConjecture/blob/main/HodgeConjecture/Definitions/Algebra/Homology/LinearDual.lean#L53}{\texttt{HodgeConjecture/Definitions/Algebra/Homology/LinearDual.lean:53}}.
\end{definition}

\begin{definition}[CategoryTheory.ShortComplex.dualCycleToHomologyFunctional]\label{def:CategoryTheory.ShortComplex.dualCycleToHomologyFunctional}
\lean{CategoryTheory.ShortComplex.dualCycleToHomologyFunctional}\leanok
\noindent\emph{Lean:} \texttt{def CategoryTheory.ShortComplex.dualCycleToHomologyFunctional}, \href{https://github.com/Paul-Lez/HodgeConjecture/blob/main/HodgeConjecture/Definitions/Algebra/Homology/LinearDual.lean#L58}{\texttt{HodgeConjecture/Definitions/Algebra/Homology/LinearDual.lean:58}}.
\end{definition}

\begin{definition}[CategoryTheory.ShortComplex.dualHomologyComparisonExplicit]\label{def:CategoryTheory.ShortComplex.dualHomologyComparisonExplicit}
\lean{CategoryTheory.ShortComplex.dualHomologyComparisonExplicit}\leanok
\uses{def:CategoryTheory.ShortComplex.dualCycleToHomologyFunctional, def:CategoryTheory.ShortComplex.linearDual}
The canonical pairing from the explicit homology of the dual complex to the dual of the
explicit homology of the original complex.

\noindent\emph{Lean:} \texttt{def CategoryTheory.ShortComplex.dualHomologyComparisonExplicit}, \href{https://github.com/Paul-Lez/HodgeConjecture/blob/main/HodgeConjecture/Definitions/Algebra/Homology/LinearDual.lean#L85}{\texttt{HodgeConjecture/Definitions/Algebra/Homology/LinearDual.lean:85}}.
\end{definition}

\begin{definition}[CategoryTheory.ShortComplex.linearDualHomologyEquiv]\label{def:CategoryTheory.ShortComplex.linearDualHomologyEquiv}
\lean{CategoryTheory.ShortComplex.linearDualHomologyEquiv}\leanok
\uses{def:CategoryTheory.ShortComplex.dualHomologyComparisonExplicit}
Universal coefficients for a short complex of vector spaces: the homology of its reversed
dual is canonically linearly equivalent to the dual of its homology.

\noindent\emph{Lean:} \texttt{def CategoryTheory.ShortComplex.linearDualHomologyEquiv}, \href{https://github.com/Paul-Lez/HodgeConjecture/blob/main/HodgeConjecture/Definitions/Algebra/Homology/LinearDual.lean#L135}{\texttt{HodgeConjecture/Definitions/Algebra/Homology/LinearDual.lean:135}}.
\end{definition}

\begin{definition}[HomologicalComplex.linearDualCochainComplex]\label{def:HomologicalComplex.linearDualCochainComplex}
\lean{HomologicalComplex.linearDualCochainComplex}\leanok
The algebraic-dual cochain complex of a nonnegatively graded chain complex of vector
spaces.

\noindent\emph{Lean:} \texttt{def HomologicalComplex.linearDualCochainComplex}, \href{https://github.com/Paul-Lez/HodgeConjecture/blob/main/HodgeConjecture/Definitions/Algebra/Homology/LinearDual.lean#L151}{\texttt{HodgeConjecture/Definitions/Algebra/Homology/LinearDual.lean:151}}.
\end{definition}

\begin{definition}[HomologicalComplex.linearDualCochainComplexScIso]\label{def:HomologicalComplex.linearDualCochainComplexScIso}
\lean{HomologicalComplex.linearDualCochainComplexScIso}\leanok
\uses{def:CategoryTheory.ShortComplex.linearDual, def:HomologicalComplex.linearDualCochainComplex}
The degree-\texttt{n} short complex of a linear-dual cochain complex is the reversed dual of the
degree-\texttt{n} short complex of the original chain complex.

\noindent\emph{Lean:} \texttt{def HomologicalComplex.linearDualCochainComplexScIso}, \href{https://github.com/Paul-Lez/HodgeConjecture/blob/main/HodgeConjecture/Definitions/Algebra/Homology/LinearDual.lean#L170}{\texttt{HodgeConjecture/Definitions/Algebra/Homology/LinearDual.lean:170}}.
\end{definition}

\begin{definition}[HomologicalComplex.linearDualHomologyEquiv]\label{def:HomologicalComplex.linearDualHomologyEquiv}
\lean{HomologicalComplex.linearDualHomologyEquiv}\leanok
\uses{def:CategoryTheory.ShortComplex.linearDualHomologyEquiv, def:HomologicalComplex.linearDualCochainComplexScIso}
Universal coefficients for a complex of vector spaces: the degree-\texttt{n} cohomology of the
linear-dual cochain complex is canonically the linear dual of the degree-\texttt{n} homology.

\noindent\emph{Lean:} \texttt{def HomologicalComplex.linearDualHomologyEquiv}, \href{https://github.com/Paul-Lez/HodgeConjecture/blob/main/HodgeConjecture/Definitions/Algebra/Homology/LinearDual.lean#L179}{\texttt{HodgeConjecture/Definitions/Algebra/Homology/LinearDual.lean:179}}.
\end{definition}

\begin{definition}[HomologicalComplex.linearDualMap]\label{def:HomologicalComplex.linearDualMap}
\lean{HomologicalComplex.linearDualMap}\leanok
\uses{def:HomologicalComplex.linearDualCochainComplex}
Algebraic duality sends a map of nonnegative chain complexes contravariantly to a map of
cochain complexes.

\noindent\emph{Lean:} \texttt{def HomologicalComplex.linearDualMap}, \href{https://github.com/Paul-Lez/HodgeConjecture/blob/main/HodgeConjecture/Definitions/Algebra/Homology/LinearDual.lean#L188}{\texttt{HodgeConjecture/Definitions/Algebra/Homology/LinearDual.lean:188}}.
\end{definition}

\begin{definition}[HomologicalComplex.linearDualIso]\label{def:HomologicalComplex.linearDualIso}
\lean{HomologicalComplex.linearDualIso}\leanok
\uses{def:HomologicalComplex.linearDualMap}
Algebraic duality sends an isomorphism of chain complexes to an isomorphism of cochain
complexes, reversing its direction.

\noindent\emph{Lean:} \texttt{def HomologicalComplex.linearDualIso}, \href{https://github.com/Paul-Lez/HodgeConjecture/blob/main/HodgeConjecture/Definitions/Algebra/Homology/LinearDual.lean#L209}{\texttt{HodgeConjecture/Definitions/Algebra/Homology/LinearDual.lean:209}}.
\end{definition}

\begin{definition}[HomologicalComplex.linearDualHomotopy]\label{def:HomologicalComplex.linearDualHomotopy}
\lean{HomologicalComplex.linearDualHomotopy}\leanok
\uses{def:HomologicalComplex.linearDualMap}
Algebraic duality sends a chain homotopy contravariantly to a cochain homotopy.

\noindent\emph{Lean:} \texttt{def HomologicalComplex.linearDualHomotopy}, \href{https://github.com/Paul-Lez/HodgeConjecture/blob/main/HodgeConjecture/Definitions/Algebra/Homology/LinearDual.lean#L224}{\texttt{HodgeConjecture/Definitions/Algebra/Homology/LinearDual.lean:224}}.
\end{definition}

\begin{definition}[HomologicalComplex.linearDualHomotopyEquiv]\label{def:HomologicalComplex.linearDualHomotopyEquiv}
\lean{HomologicalComplex.linearDualHomotopyEquiv}\leanok
\uses{def:HomologicalComplex.linearDualHomotopy}
Algebraic duality sends a chain-homotopy equivalence contravariantly to a cochain-homotopy
equivalence.

\noindent\emph{Lean:} \texttt{def HomologicalComplex.linearDualHomotopyEquiv}, \href{https://github.com/Paul-Lez/HodgeConjecture/blob/main/HodgeConjecture/Definitions/Algebra/Homology/LinearDual.lean#L246}{\texttt{HodgeConjecture/Definitions/Algebra/Homology/LinearDual.lean:246}}.
\end{definition}

\section{The canonical homotopy fiber comparison for a short exact sequence}\label{sec:HodgeConjecture.Definitions.Algebra.Homology.DerivedCategory.MappingCoconeShortExact}
\noindent\emph{File:} \href{https://github.com/Paul-Lez/HodgeConjecture/blob/main/HodgeConjecture/Definitions/Algebra/Homology/DerivedCategory/MappingCoconeShortExact.lean}{\texttt{HodgeConjecture/Definitions/Algebra/Homology/DerivedCategory/MappingCoconeShortExact.lean}}.

For \texttt{0 → A → B → C → 0} this file constructs the canonical quasi-isomorphism
\texttt{A → mappingCocone (B → C)}. Its normalization is the actual inclusion \texttt{A → B}.
The proof uses the explicit mapping-cone rotation homotopy equivalence and the
canonical quasi-isomorphism from the cone of \texttt{A → B} to \texttt{C}.

\begin{definition}[CochainComplex.mappingCocone.shiftedLiftShortComplex]\label{def:CochainComplex.mappingCocone.shiftedLiftShortComplex}
\lean{CochainComplex.mappingCocone.shiftedLiftShortComplex}\leanok
The explicit rotated-cone comparison, before shifting back to the homotopy
fiber. It is built from canonical chain maps, not from a choice of a completion
of a morphism of distinguished triangles.

\noindent\emph{Lean:} \texttt{def CochainComplex.mappingCocone.shiftedLiftShortComplex}, \href{https://github.com/Paul-Lez/HodgeConjecture/blob/main/HodgeConjecture/Definitions/Algebra/Homology/DerivedCategory/MappingCoconeShortExact.lean#L35}{\texttt{HodgeConjecture/Definitions/Algebra/Homology/DerivedCategory/MappingCoconeShortExact.lean:35}}.
\end{definition}

\section{Naturality of the canonical short-exact-sequence cone comparison}\label{sec:HodgeConjecture.Definitions.Algebra.Homology.DerivedCategory.MappingCoconeShortExactNaturality}
\noindent\emph{File:} \href{https://github.com/Paul-Lez/HodgeConjecture/blob/main/HodgeConjecture/Definitions/Algebra/Homology/DerivedCategory/MappingCoconeShortExactNaturality.lean}{\texttt{HodgeConjecture/Definitions/Algebra/Homology/DerivedCategory/MappingCoconeShortExactNaturality.lean}}.

\begin{definition}[CochainComplex.mappingCocone.shortExactHomologyIsoCone]\label{def:CochainComplex.mappingCocone.shortExactHomologyIsoCone}
\lean{CochainComplex.mappingCocone.shortExactHomologyIsoCone}\leanok
\uses{def:CochainComplex.mappingCocone.shiftedLiftShortComplex}
Homology comparison induced by the canonical short-exact-sequence
lift. Its source degree is one larger than the cone degree.

\noindent\emph{Lean:} \texttt{def CochainComplex.mappingCocone.shortExactHomologyIsoCone}, \href{https://github.com/Paul-Lez/HodgeConjecture/blob/main/HodgeConjecture/Definitions/Algebra/Homology/DerivedCategory/MappingCoconeShortExactNaturality.lean#L22}{\texttt{HodgeConjecture/Definitions/Algebra/Homology/DerivedCategory/MappingCoconeShortExactNaturality.lean:22}}.
\end{definition}

\section{Algebraic de Rham forms}\label{sec:HodgeConjecture.Definitions.Algebra.DeRham.Basic}
\noindent\emph{File:} \href{https://github.com/Paul-Lez/HodgeConjecture/blob/main/HodgeConjecture/Definitions/Algebra/DeRham/Basic.lean}{\texttt{HodgeConjecture/Definitions/Algebra/DeRham/Basic.lean}}.

Relative differential forms of degree \texttt{p} are the exterior powers \texttt{⋀[A]\^{}p Ω[A⁄R]} of the module
of relative Kähler differentials, built from Mathlib's \texttt{exteriorPower} and \texttt{KaehlerDifferential}.
Since \texttt{⋀[A]\^{}p M} is a submodule of \texttt{ExteriorAlgebra A M}, all degrees live inside the single
algebra \texttt{ExtAlg R A = ⋀ Ω[A⁄R]}, and the exterior derivative is constructed once and for all on
that algebra.

\paragraph{The exterior derivative}

The exterior derivative is the odd degree-one \texttt{R}-derivation \texttt{d} of \texttt{⋀ Ω[A⁄R]} extending the
universal derivation \texttt{A → Ω[A⁄R]}. It is obtained from a universal property rather than from a
choice of generators: an odd derivation is the same thing as an algebra map \texttt{x ↦ x + (d x) * ε}
into the square-zero extension \texttt{Sq R A = ⋀ Ω[A⁄R] ⊕ ⋀ Ω[A⁄R] * ε} by an odd square-zero element,
whose multiplication \texttt{(x, u) * (y, v) = (x * y, involute x * v + u * y)} encodes the graded
Leibniz rule. Here \texttt{involute} is the grade involution, so \texttt{involute x = (-1) \^{} p * x} on a form of
degree \texttt{p}, and the rule read off from that product is the usual
\texttt{d (x * y) = d x * y + (-1) \^{} p * x * d y}.

Giving \texttt{Sq R A} the twisted \texttt{A}-algebra structure \texttt{a ↦ a + (d a) * ε} turns that algebra map into
an \texttt{A}-algebra map, so it is produced in two steps:

\begin{itemize}
\item the universal property of \texttt{Ω[A⁄R]} applied to the derivation \texttt{a ↦ (d a) * ε} gives an
\texttt{A}-linear \texttt{phi : Ω[A⁄R] → Sq R A}, necessarily of the form \texttt{ω ↦ ω + (d ω) * ε};
\item the universal property of the exterior algebra applied to \texttt{phi} gives \texttt{deRhamHom}, and \texttt{d} is
its \texttt{ε}-component \texttt{extDeriv}.
\end{itemize}

The identities \texttt{d ∘ d = 0}, the graded Leibniz rule, and the fact that \texttt{d} raises the degree by
one then follow by induction on the exterior algebra.

\paragraph{Main definitions}

\begin{itemize}
\item \texttt{Algebra.DeRham.Form R A p}: differential forms of degree \texttt{p}, that is \texttt{⋀[A]\^{}p Ω[A⁄R]};
\item \texttt{Algebra.DeRham.differential}: the exterior derivative \texttt{Form R A p →ₗ[R] Form R A (p + 1)};
\item \texttt{Algebra.DeRham.mk R A p a₀ v}: the form \texttt{a₀ * d v₀ ∧ ... ∧ d vₚ₋₁}; these span \texttt{Form R A p}
over \texttt{R}, by \texttt{Algebra.DeRham.span\_mk};
\item \texttt{Algebra.DeRham.map}: functoriality along an \texttt{R}-algebra homomorphism.
\end{itemize}

\begin{definition}[Algebra.DeRham.ExtAlg]\label{def:Algebra.DeRham.ExtAlg}
\lean{Algebra.DeRham.ExtAlg}\leanok
The exterior algebra of the module of relative Kähler differentials. Differential forms of
each degree live inside it as the exterior powers of \texttt{Ω[A⁄R]}.

\noindent\emph{Lean:} \texttt{abbrev Algebra.DeRham.ExtAlg}, \href{https://github.com/Paul-Lez/HodgeConjecture/blob/main/HodgeConjecture/Definitions/Algebra/DeRham/Basic.lean#L72}{\texttt{HodgeConjecture/Definitions/Algebra/DeRham/Basic.lean:72}}.
\end{definition}

\begin{definition}[Algebra.DeRham.Sq]\label{def:Algebra.DeRham.Sq}
\lean{Algebra.DeRham.Sq}\leanok
\uses{def:Algebra.DeRham.ExtAlg}
The square-zero extension of \texttt{⋀ Ω[A⁄R]} by an odd square-zero element \texttt{ε}, with elements
written \texttt{x + u * ε} and multiplication \texttt{(x, u) * (y, v) = (x * y, involute x * v + u * y)}.
An \texttt{R}-algebra map \texttt{x ↦ (x, d x)} into it is precisely an odd degree-one \texttt{R}-derivation \texttt{d}.

\noindent\emph{Lean:} \texttt{def Algebra.DeRham.Sq}, \href{https://github.com/Paul-Lez/HodgeConjecture/blob/main/HodgeConjecture/Definitions/Algebra/DeRham/Basic.lean#L94}{\texttt{HodgeConjecture/Definitions/Algebra/DeRham/Basic.lean:94}}.
\end{definition}

\begin{definition}[Algebra.DeRham.Sq.mk]\label{def:Algebra.DeRham.Sq.mk}
\lean{Algebra.DeRham.Sq.mk}\leanok
\uses{def:Algebra.DeRham.Sq}
The element \texttt{x + u * ε} of the square-zero extension.

\noindent\emph{Lean:} \texttt{def Algebra.DeRham.Sq.mk}, \href{https://github.com/Paul-Lez/HodgeConjecture/blob/main/HodgeConjecture/Definitions/Algebra/DeRham/Basic.lean#L103}{\texttt{HodgeConjecture/Definitions/Algebra/DeRham/Basic.lean:103}}.
\end{definition}

\begin{definition}[Algebra.DeRham.Sq.fst]\label{def:Algebra.DeRham.Sq.fst}
\lean{Algebra.DeRham.Sq.fst}\leanok
\uses{def:Algebra.DeRham.Sq}
The part of an element of the square-zero extension not involving \texttt{ε}.

\noindent\emph{Lean:} \texttt{def Algebra.DeRham.Sq.fst}, \href{https://github.com/Paul-Lez/HodgeConjecture/blob/main/HodgeConjecture/Definitions/Algebra/DeRham/Basic.lean#L106}{\texttt{HodgeConjecture/Definitions/Algebra/DeRham/Basic.lean:106}}.
\end{definition}

\begin{definition}[Algebra.DeRham.Sq.snd]\label{def:Algebra.DeRham.Sq.snd}
\lean{Algebra.DeRham.Sq.snd}\leanok
\uses{def:Algebra.DeRham.Sq}
The coefficient of \texttt{ε} in an element of the square-zero extension.

\noindent\emph{Lean:} \texttt{def Algebra.DeRham.Sq.snd}, \href{https://github.com/Paul-Lez/HodgeConjecture/blob/main/HodgeConjecture/Definitions/Algebra/DeRham/Basic.lean#L109}{\texttt{HodgeConjecture/Definitions/Algebra/DeRham/Basic.lean:109}}.
\end{definition}

\begin{definition}[Algebra.DeRham.twist]\label{def:Algebra.DeRham.twist}
\lean{Algebra.DeRham.twist}\leanok
\uses{def:Algebra.DeRham.Sq.fst, def:Algebra.DeRham.Sq.mk, def:Algebra.DeRham.Sq.snd}
The twisted embedding \texttt{a ↦ a + (d a) * ε} of \texttt{A} into the square-zero extension.

\noindent\emph{Lean:} \texttt{def Algebra.DeRham.twist}, \href{https://github.com/Paul-Lez/HodgeConjecture/blob/main/HodgeConjecture/Definitions/Algebra/DeRham/Basic.lean#L171}{\texttt{HodgeConjecture/Definitions/Algebra/DeRham/Basic.lean:171}}.
\end{definition}

\begin{definition}[Algebra.DeRham.Sq.fstAlgHom]\label{def:Algebra.DeRham.Sq.fstAlgHom}
\lean{Algebra.DeRham.Sq.fstAlgHom}\leanok
\uses{def:Algebra.DeRham.twist}
The projection to the \texttt{ε}-free part, as an \texttt{A}-algebra map.

\noindent\emph{Lean:} \texttt{def Algebra.DeRham.Sq.fstAlgHom}, \href{https://github.com/Paul-Lez/HodgeConjecture/blob/main/HodgeConjecture/Definitions/Algebra/DeRham/Basic.lean#L234}{\texttt{HodgeConjecture/Definitions/Algebra/DeRham/Basic.lean:234}}.
\end{definition}

\begin{definition}[Algebra.DeRham.Sq.sndLinear]\label{def:Algebra.DeRham.Sq.sndLinear}
\lean{Algebra.DeRham.Sq.sndLinear}\leanok
\uses{def:Algebra.DeRham.twist}
The coefficient of \texttt{ε}, as an \texttt{R}-linear map.

\noindent\emph{Lean:} \texttt{def Algebra.DeRham.Sq.sndLinear}, \href{https://github.com/Paul-Lez/HodgeConjecture/blob/main/HodgeConjecture/Definitions/Algebra/DeRham/Basic.lean#L245}{\texttt{HodgeConjecture/Definitions/Algebra/DeRham/Basic.lean:245}}.
\end{definition}

\begin{definition}[Algebra.DeRham.epsDerivation]\label{def:Algebra.DeRham.epsDerivation}
\lean{Algebra.DeRham.epsDerivation}\leanok
\uses{def:Algebra.DeRham.twist}
The tautological \texttt{R}-derivation \texttt{a ↦ (d a) * ε} of \texttt{A} into the square-zero extension.

\noindent\emph{Lean:} \texttt{def Algebra.DeRham.epsDerivation}, \href{https://github.com/Paul-Lez/HodgeConjecture/blob/main/HodgeConjecture/Definitions/Algebra/DeRham/Basic.lean#L255}{\texttt{HodgeConjecture/Definitions/Algebra/DeRham/Basic.lean:255}}.
\end{definition}

\begin{definition}[Algebra.DeRham.phi]\label{def:Algebra.DeRham.phi}
\lean{Algebra.DeRham.phi}\leanok
\uses{def:Algebra.DeRham.epsDerivation}
Every differential \texttt{ω} determines the element \texttt{ω + (d ω) * ε} of the square-zero extension.
This is the lift of \texttt{epsDerivation} along the universal property of Kähler differentials.

\noindent\emph{Lean:} \texttt{def Algebra.DeRham.phi}, \href{https://github.com/Paul-Lez/HodgeConjecture/blob/main/HodgeConjecture/Definitions/Algebra/DeRham/Basic.lean#L266}{\texttt{HodgeConjecture/Definitions/Algebra/DeRham/Basic.lean:266}}.
\end{definition}

\begin{definition}[Algebra.DeRham.deRhamHom]\label{def:Algebra.DeRham.deRhamHom}
\lean{Algebra.DeRham.deRhamHom}\leanok
\uses{def:Algebra.DeRham.Sq.fstAlgHom, def:Algebra.DeRham.phi}
The exterior derivative on \texttt{⋀ Ω[A⁄R]}, packaged as the algebra map \texttt{x ↦ x + (d x) * ε} into
the square-zero extension.

\noindent\emph{Lean:} \texttt{def Algebra.DeRham.deRhamHom}, \href{https://github.com/Paul-Lez/HodgeConjecture/blob/main/HodgeConjecture/Definitions/Algebra/DeRham/Basic.lean#L320}{\texttt{HodgeConjecture/Definitions/Algebra/DeRham/Basic.lean:320}}.
\end{definition}

\begin{definition}[Algebra.DeRham.extDeriv]\label{def:Algebra.DeRham.extDeriv}
\lean{Algebra.DeRham.extDeriv}\leanok
\uses{def:Algebra.DeRham.Sq.sndLinear, def:Algebra.DeRham.deRhamHom}
The exterior derivative on the exterior algebra of Kähler differentials.

\noindent\emph{Lean:} \texttt{def Algebra.DeRham.extDeriv}, \href{https://github.com/Paul-Lez/HodgeConjecture/blob/main/HodgeConjecture/Definitions/Algebra/DeRham/Basic.lean#L336}{\texttt{HodgeConjecture/Definitions/Algebra/DeRham/Basic.lean:336}}.
\end{definition}

\begin{definition}[Algebra.DeRham.Form]\label{def:Algebra.DeRham.Form}
\lean{Algebra.DeRham.Form}\leanok
Relative Kähler differential forms of degree \texttt{p}: the \texttt{p}-th exterior power over \texttt{A} of the
module of relative Kähler differentials \texttt{Ω[A⁄R]}.

\noindent\emph{Lean:} \texttt{abbrev Algebra.DeRham.Form}, \href{https://github.com/Paul-Lez/HodgeConjecture/blob/main/HodgeConjecture/Definitions/Algebra/DeRham/Basic.lean#L410}{\texttt{HodgeConjecture/Definitions/Algebra/DeRham/Basic.lean:410}}.
\end{definition}

\begin{definition}[Algebra.DeRham.differential]\label{def:Algebra.DeRham.differential}
\lean{Algebra.DeRham.differential}\leanok
\uses{def:Algebra.DeRham.Form, def:Algebra.DeRham.extDeriv}
The exterior derivative on differential forms of degree \texttt{p}.

\noindent\emph{Lean:} \texttt{def Algebra.DeRham.differential}, \href{https://github.com/Paul-Lez/HodgeConjecture/blob/main/HodgeConjecture/Definitions/Algebra/DeRham/Basic.lean#L414}{\texttt{HodgeConjecture/Definitions/Algebra/DeRham/Basic.lean:414}}.
\end{definition}

\begin{definition}[Algebra.DeRham.exact]\label{def:Algebra.DeRham.exact}
\lean{Algebra.DeRham.exact}\leanok
\uses{def:Algebra.DeRham.Form}
The exact form \texttt{d a₁ ∧ ⋯ ∧ d aₚ}.

\noindent\emph{Lean:} \texttt{def Algebra.DeRham.exact}, \href{https://github.com/Paul-Lez/HodgeConjecture/blob/main/HodgeConjecture/Definitions/Algebra/DeRham/Basic.lean#L430}{\texttt{HodgeConjecture/Definitions/Algebra/DeRham/Basic.lean:430}}.
\end{definition}

\begin{definition}[Algebra.DeRham.mk]\label{def:Algebra.DeRham.mk}
\lean{Algebra.DeRham.mk}\leanok
\uses{def:Algebra.DeRham.exact}
The differential form \texttt{a₀ * d a₁ ∧ ⋯ ∧ d aₚ}. Such forms span the degree \texttt{p} forms over the
base ring.

\noindent\emph{Lean:} \texttt{def Algebra.DeRham.mk}, \href{https://github.com/Paul-Lez/HodgeConjecture/blob/main/HodgeConjecture/Definitions/Algebra/DeRham/Basic.lean#L453}{\texttt{HodgeConjecture/Definitions/Algebra/DeRham/Basic.lean:453}}.
\end{definition}

\begin{definition}[Algebra.DeRham.ofFunction]\label{def:Algebra.DeRham.ofFunction}
\lean{Algebra.DeRham.ofFunction}\leanok
\uses{def:Algebra.DeRham.mk}
A function regarded as a differential form of degree zero.

\noindent\emph{Lean:} \texttt{def Algebra.DeRham.ofFunction}, \href{https://github.com/Paul-Lez/HodgeConjecture/blob/main/HodgeConjecture/Definitions/Algebra/DeRham/Basic.lean#L468}{\texttt{HodgeConjecture/Definitions/Algebra/DeRham/Basic.lean:468}}.
\end{definition}

\begin{definition}[Algebra.DeRham.ofConstant]\label{def:Algebra.DeRham.ofConstant}
\lean{Algebra.DeRham.ofConstant}\leanok
\uses{def:Algebra.DeRham.ofFunction}
A scalar from the base ring regarded as a differential form of degree zero.

\noindent\emph{Lean:} \texttt{def Algebra.DeRham.ofConstant}, \href{https://github.com/Paul-Lez/HodgeConjecture/blob/main/HodgeConjecture/Definitions/Algebra/DeRham/Basic.lean#L476}{\texttt{HodgeConjecture/Definitions/Algebra/DeRham/Basic.lean:476}}.
\end{definition}

\begin{definition}[Algebra.DeRham.extAlgMap]\label{def:Algebra.DeRham.extAlgMap}
\lean{Algebra.DeRham.extAlgMap}\leanok
\uses{def:Algebra.DeRham.ExtAlg}
The algebra map on exterior algebras of Kähler differentials attached to a tower
\texttt{R → A → B}.

\noindent\emph{Lean:} \texttt{def Algebra.DeRham.extAlgMap}, \href{https://github.com/Paul-Lez/HodgeConjecture/blob/main/HodgeConjecture/Definitions/Algebra/DeRham/Basic.lean#L554}{\texttt{HodgeConjecture/Definitions/Algebra/DeRham/Basic.lean:554}}.
\end{definition}

\begin{definition}[Algebra.DeRham.mapTower]\label{def:Algebra.DeRham.mapTower}
\lean{Algebra.DeRham.mapTower}\leanok
\uses{def:Algebra.DeRham.Form, def:Algebra.DeRham.extAlgMap}
Pull differential forms forward along a tower \texttt{R → A → B}.

\noindent\emph{Lean:} \texttt{def Algebra.DeRham.mapTower}, \href{https://github.com/Paul-Lez/HodgeConjecture/blob/main/HodgeConjecture/Definitions/Algebra/DeRham/Basic.lean#L608}{\texttt{HodgeConjecture/Definitions/Algebra/DeRham/Basic.lean:608}}.
\end{definition}

\begin{definition}[Algebra.DeRham.map]\label{def:Algebra.DeRham.map}
\lean{Algebra.DeRham.map}\leanok
\uses{def:Algebra.DeRham.mapTower}
Pull differential forms forward along an algebra homomorphism.

\noindent\emph{Lean:} \texttt{def Algebra.DeRham.map}, \href{https://github.com/Paul-Lez/HodgeConjecture/blob/main/HodgeConjecture/Definitions/Algebra/DeRham/Basic.lean#L626}{\texttt{HodgeConjecture/Definitions/Algebra/DeRham/Basic.lean:626}}.
\end{definition}

\chapter{Definitions: AlgebraicGeometry}

\section{Projective space and explicit projective presentations}\label{sec:HodgeConjecture.Definitions.AlgebraicGeometry.ProjectiveSpace}
\noindent\emph{File:} \href{https://github.com/Paul-Lez/HodgeConjecture/blob/main/HodgeConjecture/Definitions/AlgebraicGeometry/ProjectiveSpace.lean}{\texttt{HodgeConjecture/Definitions/AlgebraicGeometry/ProjectiveSpace.lean}}.

This file constructs projective space over a scheme as a pullback from \texttt{Proj}, proves its
structure morphism proper for finitely many homogeneous coordinates, and packages a closed
embedding into finite-dimensional projective space as an explicit projective presentation.

The original projective-space construction was contributed to Mathlib in
\href{https://github.com/leanprover-community/mathlib4/pull/26061}{\texttt{mathlib4} pull request \#26061}.

\begin{definition}[ProjectiveSpace]\label{def:ProjectiveSpace}
\lean{ProjectiveSpace}\leanok
The projective space over a scheme \texttt{S}, with homogeneous coordinates indexed by \texttt{n}

\noindent\emph{Lean:} \texttt{def ProjectiveSpace}, \href{https://github.com/Paul-Lez/HodgeConjecture/blob/main/HodgeConjecture/Definitions/AlgebraicGeometry/ProjectiveSpace.lean#L50}{\texttt{HodgeConjecture/Definitions/AlgebraicGeometry/ProjectiveSpace.lean:50}}.
\end{definition}

\begin{definition}[ProjectiveSpace.degreeZeroEquiv]\label{def:ProjectiveSpace.degreeZeroEquiv}
\lean{ProjectiveSpace.degreeZeroEquiv}\leanok
The degree-zero part of the standard grading on a polynomial ring consists exactly of
the constant polynomials.

\noindent\emph{Lean:} \texttt{def ProjectiveSpace.degreeZeroEquiv}, \href{https://github.com/Paul-Lez/HodgeConjecture/blob/main/HodgeConjecture/Definitions/AlgebraicGeometry/ProjectiveSpace.lean#L61}{\texttt{HodgeConjecture/Definitions/AlgebraicGeometry/ProjectiveSpace.lean:61}}.
\end{definition}

\begin{definition}[ProjectiveSpace.toBase]\label{def:ProjectiveSpace.toBase}
\lean{ProjectiveSpace.toBase}\leanok
\uses{def:ProjectiveSpace}
The structure morphism from projective space to its base scheme.

\noindent\emph{Lean:} \texttt{def ProjectiveSpace.toBase}, \href{https://github.com/Paul-Lez/HodgeConjecture/blob/main/HodgeConjecture/Definitions/AlgebraicGeometry/ProjectiveSpace.lean#L108}{\texttt{HodgeConjecture/Definitions/AlgebraicGeometry/ProjectiveSpace.lean:108}}.
\end{definition}

\begin{definition}[ProjectiveSpace.Presentation]\label{def:ProjectiveSpace.Presentation}
\lean{ProjectiveSpace.Presentation}\leanok
\uses{def:ProjectiveSpace.toBase}
An explicit closed embedding over \texttt{S} into a finite-dimensional projective space.

\noindent\emph{Lean:} \texttt{structure ProjectiveSpace.Presentation}, \href{https://github.com/Paul-Lez/HodgeConjecture/blob/main/HodgeConjecture/Definitions/AlgebraicGeometry/ProjectiveSpace.lean#L121}{\texttt{HodgeConjecture/Definitions/AlgebraicGeometry/ProjectiveSpace.lean:121}}.
\end{definition}

\begin{definition}[AlgebraicGeometry.IsProjective]\label{def:AlgebraicGeometry.IsProjective}
\lean{AlgebraicGeometry.IsProjective}\leanok
\uses{def:ProjectiveSpace.Presentation}
A scheme morphism is projective if it admits a finite-dimensional projective presentation.

\noindent\emph{Lean:} \texttt{class AlgebraicGeometry.IsProjective}, \href{https://github.com/Paul-Lez/HodgeConjecture/blob/main/HodgeConjecture/Definitions/AlgebraicGeometry/ProjectiveSpace.lean#L142}{\texttt{HodgeConjecture/Definitions/AlgebraicGeometry/ProjectiveSpace.lean:142}}.
\end{definition}

\section{Points over a commutative ring, and analytification}\label{sec:HodgeConjecture.Definitions.AlgebraicGeometry.Points}
\noindent\emph{File:} \href{https://github.com/Paul-Lez/HodgeConjecture/blob/main/HodgeConjecture/Definitions/AlgebraicGeometry/Points.lean}{\texttt{HodgeConjecture/Definitions/AlgebraicGeometry/Points.lean}}.

For a scheme \texttt{X : Over (Spec R)} over a commutative ring \texttt{R}, an \texttt{R}-point is a morphism
\texttt{Over.mk (𝟙 (Spec R)) ⟶ X}. The structure map is bundled with \texttt{X}, and \texttt{map} is composition
in this over category. These constructions require only that \texttt{R} is a commutative ring.

Everything that speaks of *the* point underlying an \texttt{R}-point needs \texttt{Spec R} to have a unique
closed point, so from \texttt{stalkData} onwards \texttt{R} is local. An \texttt{R}-point then determines a scheme point
together with a local homomorphism out of the stalk there, and evaluating the sections on an open
\texttt{U} containing it gives a ring homomorphism \texttt{Γ(X.left, U) ⟶ ↧R}. Over an affine open this
identifies the \texttt{R}-points with a subspace of that hom set, which carries Mathlib's topology of pointwise
convergence; the analytic topology is the topology glued from these affine charts. It exists as
soon as multiplication on \texttt{R} is continuous and the units of \texttt{R} are open, so it applies to \texttt{ℝ},
to a \texttt{p}-adic field, and to \texttt{ℤ\_p} alike.

\texttt{analyticTopology\_eq\_generateFrom} identifies it with the more elementary description: the
topology generated by the sets on which a local regular function takes values in a prescribed open
subset of \texttt{R}.

Over a field the stalk homomorphism factors through the residue field, which is the classical
description; \texttt{residueData} and the lemmas beside it record it, and there nonvanishing replaces
invertibility.

Complex points are the case \texttt{K = ℂ}. Analytification is then packaged as a functor from schemes
over \texttt{Spec ℂ} to topological spaces, so that Betti (co)homology of a complex variety is Mathlib's
singular (co)homology of the resulting space.

\texttt{isNoetherian\_of\_isProjective} records that a projective complex scheme is Noetherian.

\begin{definition}[AlgebraicGeometry.Point]\label{def:AlgebraicGeometry.Point}
\lean{AlgebraicGeometry.Point}\leanok
An \texttt{R}-point of a scheme over \texttt{Spec R}, as a morphism in the over category.

\noindent\emph{Lean:} \texttt{abbrev AlgebraicGeometry.Point}, \href{https://github.com/Paul-Lez/HodgeConjecture/blob/main/HodgeConjecture/Definitions/AlgebraicGeometry/Points.lean#L68}{\texttt{HodgeConjecture/Definitions/AlgebraicGeometry/Points.lean:68}}.
\end{definition}

\begin{definition}[AlgebraicGeometry.Point.map]\label{def:AlgebraicGeometry.Point.map}
\lean{AlgebraicGeometry.Point.map}\leanok
\uses{def:AlgebraicGeometry.Point}
The map on \texttt{R}-points induced by a morphism over \texttt{Spec R}.

\noindent\emph{Lean:} \texttt{def AlgebraicGeometry.Point.map}, \href{https://github.com/Paul-Lez/HodgeConjecture/blob/main/HodgeConjecture/Definitions/AlgebraicGeometry/Points.lean#L80}{\texttt{HodgeConjecture/Definitions/AlgebraicGeometry/Points.lean:80}}.
\end{definition}

\begin{definition}[AlgebraicGeometry.Point.stalkData]\label{def:AlgebraicGeometry.Point.stalkData}
\lean{AlgebraicGeometry.Point.stalkData}\leanok
\uses{def:AlgebraicGeometry.Point}
The scheme point and local stalk homomorphism corresponding to an \texttt{R}-point.

\texttt{Spec R} has a unique closed point when \texttt{R} is local, and a morphism out of it is exactly a point
of \texttt{X.left} together with a local homomorphism from the stalk there. Everything below that speaks of
*the* point underlying an \texttt{R}-point rests on this.

\noindent\emph{Lean:} \texttt{def AlgebraicGeometry.Point.stalkData}, \href{https://github.com/Paul-Lez/HodgeConjecture/blob/main/HodgeConjecture/Definitions/AlgebraicGeometry/Points.lean#L101}{\texttt{HodgeConjecture/Definitions/AlgebraicGeometry/Points.lean:101}}.
\end{definition}

\begin{definition}[AlgebraicGeometry.Point.underlying]\label{def:AlgebraicGeometry.Point.underlying}
\lean{AlgebraicGeometry.Point.underlying}\leanok
\uses{def:AlgebraicGeometry.Point.stalkData}
The underlying point of the scheme.

\noindent\emph{Lean:} \texttt{def AlgebraicGeometry.Point.underlying}, \href{https://github.com/Paul-Lez/HodgeConjecture/blob/main/HodgeConjecture/Definitions/AlgebraicGeometry/Points.lean#L110}{\texttt{HodgeConjecture/Definitions/AlgebraicGeometry/Points.lean:110}}.
\end{definition}

\begin{definition}[AlgebraicGeometry.Point.stalkHom]\label{def:AlgebraicGeometry.Point.stalkHom}
\lean{AlgebraicGeometry.Point.stalkHom}\leanok
\uses{def:AlgebraicGeometry.Point.underlying}
The local homomorphism from the stalk at the underlying point.

\noindent\emph{Lean:} \texttt{def AlgebraicGeometry.Point.stalkHom}, \href{https://github.com/Paul-Lez/HodgeConjecture/blob/main/HodgeConjecture/Definitions/AlgebraicGeometry/Points.lean#L113}{\texttt{HodgeConjecture/Definitions/AlgebraicGeometry/Points.lean:113}}.
\end{definition}

\begin{definition}[AlgebraicGeometry.Point.overOpen]\label{def:AlgebraicGeometry.Point.overOpen}
\lean{AlgebraicGeometry.Point.overOpen}\leanok
\uses{def:AlgebraicGeometry.Point.underlying}
The \texttt{R}-points whose underlying scheme point lies in \texttt{U}.

\noindent\emph{Lean:} \texttt{def AlgebraicGeometry.Point.overOpen}, \href{https://github.com/Paul-Lez/HodgeConjecture/blob/main/HodgeConjecture/Definitions/AlgebraicGeometry/Points.lean#L119}{\texttt{HodgeConjecture/Definitions/AlgebraicGeometry/Points.lean:119}}.
\end{definition}

\begin{definition}[AlgebraicGeometry.Point.OverOpen]\label{def:AlgebraicGeometry.Point.OverOpen}
\lean{AlgebraicGeometry.Point.OverOpen}\leanok
\uses{def:AlgebraicGeometry.Point.underlying}
An \texttt{R}-point whose underlying scheme point lies in \texttt{U}.

Using this subtype is preferable whenever a local regular function is evaluated: it records the
domain condition in the type instead of assigning an arbitrary value outside \texttt{U}.

\noindent\emph{Lean:} \texttt{abbrev AlgebraicGeometry.Point.OverOpen}, \href{https://github.com/Paul-Lez/HodgeConjecture/blob/main/HodgeConjecture/Definitions/AlgebraicGeometry/Points.lean#L123}{\texttt{HodgeConjecture/Definitions/AlgebraicGeometry/Points.lean:123}}.
\end{definition}

\begin{definition}[AlgebraicGeometry.Point.evaluate]\label{def:AlgebraicGeometry.Point.evaluate}
\lean{AlgebraicGeometry.Point.evaluate}\leanok
\uses{def:AlgebraicGeometry.Point.stalkHom}
Evaluation of a local regular function at an \texttt{R}-point. Outside the function's domain this is
defined to be zero; all uses in the analytic topology are restricted to \texttt{overOpen U}.

\noindent\emph{Lean:} \texttt{def AlgebraicGeometry.Point.evaluate}, \href{https://github.com/Paul-Lez/HodgeConjecture/blob/main/HodgeConjecture/Definitions/AlgebraicGeometry/Points.lean#L134}{\texttt{HodgeConjecture/Definitions/AlgebraicGeometry/Points.lean:134}}.
\end{definition}

\begin{definition}[AlgebraicGeometry.Point.evaluationHom]\label{def:AlgebraicGeometry.Point.evaluationHom}
\lean{AlgebraicGeometry.Point.evaluationHom}\leanok
\uses{def:AlgebraicGeometry.Point.OverOpen, def:AlgebraicGeometry.Point.stalkHom}
Evaluation of all regular functions on \texttt{U} at an \texttt{R}-point lying over \texttt{U}, bundled as a
ring homomorphism \texttt{Γ(X.left, U) ⟶ ↧R}.

This is the chart map of the analytic topology: over an affine open it identifies the \texttt{R}-points
with Mathlib's space of ring homomorphisms into \texttt{R}.

\noindent\emph{Lean:} \texttt{def AlgebraicGeometry.Point.evaluationHom}, \href{https://github.com/Paul-Lez/HodgeConjecture/blob/main/HodgeConjecture/Definitions/AlgebraicGeometry/Points.lean#L165}{\texttt{HodgeConjecture/Definitions/AlgebraicGeometry/Points.lean:165}}.
\end{definition}

\begin{definition}[AlgebraicGeometry.Point.chartTopology]\label{def:AlgebraicGeometry.Point.chartTopology}
\lean{AlgebraicGeometry.Point.chartTopology}\leanok
\uses{def:AlgebraicGeometry.Point.evaluationHom}
The chart topology on the \texttt{R}-points lying over an open subset: the topology of pointwise
convergence, pulled back along \texttt{evaluationHom} from Mathlib's topology on \texttt{Γ(X.left, U) ⟶ ↧R}.

\noindent\emph{Lean:} \texttt{def AlgebraicGeometry.Point.chartTopology}, \href{https://github.com/Paul-Lez/HodgeConjecture/blob/main/HodgeConjecture/Definitions/AlgebraicGeometry/Points.lean#L242}{\texttt{HodgeConjecture/Definitions/AlgebraicGeometry/Points.lean:242}}.
\end{definition}

\begin{definition}[AlgebraicGeometry.Point.analyticSubbasis]\label{def:AlgebraicGeometry.Point.analyticSubbasis}
\lean{AlgebraicGeometry.Point.analyticSubbasis}\leanok
\uses{def:AlgebraicGeometry.Point.evaluate, def:AlgebraicGeometry.Point.overOpen}
Sets obtained by restricting an inverse image of an open subset of \texttt{R} to the domain of a local
regular function.

\noindent\emph{Lean:} \texttt{def AlgebraicGeometry.Point.analyticSubbasis}, \href{https://github.com/Paul-Lez/HodgeConjecture/blob/main/HodgeConjecture/Definitions/AlgebraicGeometry/Points.lean#L249}{\texttt{HodgeConjecture/Definitions/AlgebraicGeometry/Points.lean:249}}.
\end{definition}

\begin{definition}[AlgebraicGeometry.Point.continuousMap]\label{def:AlgebraicGeometry.Point.continuousMap}
\lean{AlgebraicGeometry.Point.continuousMap}\leanok
\uses{def:AlgebraicGeometry.Point.analyticSubbasis, def:AlgebraicGeometry.Point.chartTopology, def:AlgebraicGeometry.Point.map}
A morphism over \texttt{Spec R}, bundled as a continuous map on \texttt{R}-points.

\noindent\emph{Lean:} \texttt{def AlgebraicGeometry.Point.continuousMap}, \href{https://github.com/Paul-Lez/HodgeConjecture/blob/main/HodgeConjecture/Definitions/AlgebraicGeometry/Points.lean#L387}{\texttt{HodgeConjecture/Definitions/AlgebraicGeometry/Points.lean:387}}.
\end{definition}

\begin{definition}[AlgebraicGeometry.Point.isoMapHomeomorph]\label{def:AlgebraicGeometry.Point.isoMapHomeomorph}
\lean{AlgebraicGeometry.Point.isoMapHomeomorph}\leanok
\uses{def:AlgebraicGeometry.Point.analyticSubbasis, def:AlgebraicGeometry.Point.chartTopology, def:AlgebraicGeometry.Point.map}
An isomorphism of schemes over \texttt{Spec R} induces a homeomorphism on \texttt{R}-points.

\noindent\emph{Lean:} \texttt{def AlgebraicGeometry.Point.isoMapHomeomorph}, \href{https://github.com/Paul-Lez/HodgeConjecture/blob/main/HodgeConjecture/Definitions/AlgebraicGeometry/Points.lean#L393}{\texttt{HodgeConjecture/Definitions/AlgebraicGeometry/Points.lean:393}}.
\end{definition}

\begin{definition}[AlgebraicGeometry.Point.residueData]\label{def:AlgebraicGeometry.Point.residueData}
\lean{AlgebraicGeometry.Point.residueData}\leanok
\uses{def:AlgebraicGeometry.Point}
Over a field, a \texttt{K}-point is a scheme point together with an embedding of its residue field.

This refines \texttt{stalkData}: the stalk homomorphism out of a local ring into a field factors through
the residue field.

\noindent\emph{Lean:} \texttt{def AlgebraicGeometry.Point.residueData}, \href{https://github.com/Paul-Lez/HodgeConjecture/blob/main/HodgeConjecture/Definitions/AlgebraicGeometry/Points.lean#L421}{\texttt{HodgeConjecture/Definitions/AlgebraicGeometry/Points.lean:421}}.
\end{definition}

\begin{definition}[AlgebraicGeometry.ComplexPoint]\label{def:AlgebraicGeometry.ComplexPoint}
\lean{AlgebraicGeometry.ComplexPoint}\leanok
\uses{def:AlgebraicGeometry.Point}
A complex point of a scheme over \texttt{Spec ℂ}.

\noindent\emph{Lean:} \texttt{abbrev AlgebraicGeometry.ComplexPoint}, \href{https://github.com/Paul-Lez/HodgeConjecture/blob/main/HodgeConjecture/Definitions/AlgebraicGeometry/Points.lean#L433}{\texttt{HodgeConjecture/Definitions/AlgebraicGeometry/Points.lean:433}}.
\end{definition}

\section{Geometric support of algebraic cycles}\label{sec:HodgeConjecture.Definitions.AlgebraicGeometry.AlgebraicCycleSupport}
\noindent\emph{File:} \href{https://github.com/Paul-Lez/HodgeConjecture/blob/main/HodgeConjecture/Definitions/AlgebraicGeometry/AlgebraicCycleSupport.lean}{\texttt{HodgeConjecture/Definitions/AlgebraicGeometry/AlgebraicCycleSupport.lean}}.

An algebraic cycle in Mathlib is indexed by the generic points of its irreducible components.
This file constructs the reduced integral closed subscheme attached to such a point and the
corresponding closed subset of complex points. It also constructs the geometric support of a
whole cycle as the union of the closures of the generic points with nonzero coefficient.

\begin{definition}[AlgebraicGeometry.cycleComponent]\label{def:AlgebraicGeometry.cycleComponent}
\lean{AlgebraicGeometry.cycleComponent}\leanok
The reduced closed subscheme whose underlying space is the closure of \texttt{x}.

\noindent\emph{Lean:} \texttt{def AlgebraicGeometry.cycleComponent}, \href{https://github.com/Paul-Lez/HodgeConjecture/blob/main/HodgeConjecture/Definitions/AlgebraicGeometry/AlgebraicCycleSupport.lean#L44}{\texttt{HodgeConjecture/Definitions/AlgebraicGeometry/AlgebraicCycleSupport.lean:44}}.
\end{definition}

\begin{definition}[AlgebraicGeometry.cycleComponentι]\label{def:AlgebraicGeometry.cycleComponentι}
\lean{AlgebraicGeometry.cycleComponentι}\leanok
\uses{def:AlgebraicGeometry.cycleComponent}
The canonical closed immersion of the reduced closure of \texttt{x}.

\noindent\emph{Lean:} \texttt{def AlgebraicGeometry.cycleComponentι}, \href{https://github.com/Paul-Lez/HodgeConjecture/blob/main/HodgeConjecture/Definitions/AlgebraicGeometry/AlgebraicCycleSupport.lean#L49}{\texttt{HodgeConjecture/Definitions/AlgebraicGeometry/AlgebraicCycleSupport.lean:49}}.
\end{definition}

\begin{definition}[AlgebraicGeometry.cycleComponentSupport]\label{def:AlgebraicGeometry.cycleComponentSupport}
\lean{AlgebraicGeometry.cycleComponentSupport}\leanok
\uses{def:AlgebraicGeometry.ComplexPoint, def:AlgebraicGeometry.IsProjective, def:AlgebraicGeometry.Point.underlying}
The complex points supported on the irreducible closed subset with generic point \texttt{x}.

\noindent\emph{Lean:} \texttt{def AlgebraicGeometry.cycleComponentSupport}, \href{https://github.com/Paul-Lez/HodgeConjecture/blob/main/HodgeConjecture/Definitions/AlgebraicGeometry/AlgebraicCycleSupport.lean#L143}{\texttt{HodgeConjecture/Definitions/AlgebraicGeometry/AlgebraicCycleSupport.lean:143}}.
\end{definition}

\section{Analytic local left inverses from actual section lifting}\label{sec:HodgeConjecture.Definitions.AlgebraicGeometry.ClosedImmersionAnalyticLeftInverse}
\noindent\emph{File:} \href{https://github.com/Paul-Lez/HodgeConjecture/blob/main/HodgeConjecture/Definitions/AlgebraicGeometry/ClosedImmersionAnalyticLeftInverse.lean}{\texttt{HodgeConjecture/Definitions/AlgebraicGeometry/ClosedImmersionAnalyticLeftInverse.lean}}.

Intrinsic regular coordinate functions lift through a closed immersion on a common affine
ambient neighborhood. Their analytic evaluations give a local left inverse to the inclusion
written in complex charts. In particular derivative injectivity is proved, not supplied as
an immersion or purity field.

\begin{definition}[AlgebraicGeometry.ComplexPoint.inclusionInComplexCharts]\label{def:AlgebraicGeometry.ComplexPoint.inclusionInComplexCharts}
\lean{AlgebraicGeometry.ComplexPoint.inclusionInComplexCharts}\leanok
\uses{def:AlgebraicGeometry.ComplexPoint.localChart}
The actual inclusion written in the canonical intrinsic and ambient complex charts.

\noindent\emph{Lean:} \texttt{def AlgebraicGeometry.ComplexPoint.inclusionInComplexCharts}, \href{https://github.com/Paul-Lez/HodgeConjecture/blob/main/HodgeConjecture/Definitions/AlgebraicGeometry/ClosedImmersionAnalyticLeftInverse.lean#L37}{\texttt{HodgeConjecture/Definitions/AlgebraicGeometry/ClosedImmersionAnalyticLeftInverse.lean:37}}.
\end{definition}

\section{Local dual classes on cycle components}\label{sec:HodgeConjecture.Definitions.AlgebraicGeometry.CycleComponentPurity}
\noindent\emph{File:} \href{https://github.com/Paul-Lez/HodgeConjecture/blob/main/HodgeConjecture/Definitions/AlgebraicGeometry/CycleComponentPurity.lean}{\texttt{HodgeConjecture/Definitions/AlgebraicGeometry/CycleComponentPurity.lean}}.

The local fundamental class obtained from exact coordinates on the smooth locus of a cycle
component is nonzero.  Since it generates local homology, there is a unique local cohomology
class evaluating to one on it, and that normalized dual class generates cohomology supported at
the chosen point.

This is intrinsic purity on the smooth component neighborhood.  It does not yet give ambient
purity in codimension \texttt{p}: that requires a Thom or Gysin comparison between the intrinsic local
homology in degree \texttt{2 * (d - p)} and ambient cohomology supported on the component in degree
\texttt{2 * p}.

\begin{definition}[AlgebraicTopology.Singular.normalizedDual]\label{def:AlgebraicTopology.Singular.normalizedDual}
\lean{AlgebraicTopology.Singular.normalizedDual}\leanok
The normalized dual of a nonzero vector.

\noindent\emph{Lean:} \texttt{def AlgebraicTopology.Singular.normalizedDual}, \href{https://github.com/Paul-Lez/HodgeConjecture/blob/main/HodgeConjecture/Definitions/AlgebraicGeometry/CycleComponentPurity.lean#L48}{\texttt{HodgeConjecture/Definitions/AlgebraicGeometry/CycleComponentPurity.lean:48}}.
\end{definition}

\begin{definition}[AlgebraicTopology.Singular.normalizedRelativeCoclass]\label{def:AlgebraicTopology.Singular.normalizedRelativeCoclass}
\lean{AlgebraicTopology.Singular.normalizedRelativeCoclass}\leanok
\uses{def:AlgebraicTopology.Singular.normalizedDual, def:AlgebraicTopology.Singular.relativeCohomologyEquivDualHomology}
The relative cohomology class normalized to pair to one with a given nonzero relative
homology class. Cohomology is the homology of the dual cochain complex rather than the dual of
homology, so the normalized functional is transported along the universal-coefficient
equivalence.

\noindent\emph{Lean:} \texttt{def AlgebraicTopology.Singular.normalizedRelativeCoclass}, \href{https://github.com/Paul-Lez/HodgeConjecture/blob/main/HodgeConjecture/Definitions/AlgebraicGeometry/CycleComponentPurity.lean#L57}{\texttt{HodgeConjecture/Definitions/AlgebraicGeometry/CycleComponentPurity.lean:57}}.
\end{definition}

\section{A finite smooth decomposition of reduced closed subschemes}\label{sec:HodgeConjecture.Definitions.AlgebraicGeometry.ReducedSmoothStratification}
\noindent\emph{File:} \href{https://github.com/Paul-Lez/HodgeConjecture/blob/main/HodgeConjecture/Definitions/AlgebraicGeometry/ReducedSmoothStratification.lean}{\texttt{HodgeConjecture/Definitions/AlgebraicGeometry/ReducedSmoothStratification.lean}}.

Each step takes the actual reduced closed subscheme on a closed subset, removes its smooth
locus, and continues on the closed remainder. Noetherian induction makes this construction
finite. The pieces are actual smooth locally closed subschemes, not supplied stratification
data. This is an algebraic prerequisite for dimension induction; it does not assert a
Whitney or frontier condition, triangulation, analytic homology vanishing, or extension of an
orientation across singularities.

\begin{definition}[AlgebraicGeometry.reducedClosedSubscheme]\label{def:AlgebraicGeometry.reducedClosedSubscheme}
\lean{AlgebraicGeometry.reducedClosedSubscheme}\leanok
The actual reduced closed subscheme on a closed subset.

\noindent\emph{Lean:} \texttt{def AlgebraicGeometry.reducedClosedSubscheme}, \href{https://github.com/Paul-Lez/HodgeConjecture/blob/main/HodgeConjecture/Definitions/AlgebraicGeometry/ReducedSmoothStratification.lean#L33}{\texttt{HodgeConjecture/Definitions/AlgebraicGeometry/ReducedSmoothStratification.lean:33}}.
\end{definition}

\begin{definition}[AlgebraicGeometry.reducedClosedSubschemeι]\label{def:AlgebraicGeometry.reducedClosedSubschemeι}
\lean{AlgebraicGeometry.reducedClosedSubschemeι}\leanok
\uses{def:AlgebraicGeometry.reducedClosedSubscheme}
The canonical inclusion of the reduced closed subscheme.

\noindent\emph{Lean:} \texttt{def AlgebraicGeometry.reducedClosedSubschemeι}, \href{https://github.com/Paul-Lez/HodgeConjecture/blob/main/HodgeConjecture/Definitions/AlgebraicGeometry/ReducedSmoothStratification.lean#L37}{\texttt{HodgeConjecture/Definitions/AlgebraicGeometry/ReducedSmoothStratification.lean:37}}.
\end{definition}

\begin{definition}[AlgebraicGeometry.reducedClosedStructureMap]\label{def:AlgebraicGeometry.reducedClosedStructureMap}
\lean{AlgebraicGeometry.reducedClosedStructureMap}\leanok
\uses{def:AlgebraicGeometry.reducedClosedSubschemeι}
The structure morphism of the actual reduced closed subscheme.

\noindent\emph{Lean:} \texttt{def AlgebraicGeometry.reducedClosedStructureMap}, \href{https://github.com/Paul-Lez/HodgeConjecture/blob/main/HodgeConjecture/Definitions/AlgebraicGeometry/ReducedSmoothStratification.lean#L64}{\texttt{HodgeConjecture/Definitions/AlgebraicGeometry/ReducedSmoothStratification.lean:64}}.
\end{definition}

\begin{definition}[AlgebraicGeometry.reducedClosedSmoothPiece]\label{def:AlgebraicGeometry.reducedClosedSmoothPiece}
\lean{AlgebraicGeometry.reducedClosedSmoothPiece}\leanok
\uses{def:AlgebraicGeometry.reducedClosedStructureMap}
The smooth piece removed from a closed subset at one step.

\noindent\emph{Lean:} \texttt{def AlgebraicGeometry.reducedClosedSmoothPiece}, \href{https://github.com/Paul-Lez/HodgeConjecture/blob/main/HodgeConjecture/Definitions/AlgebraicGeometry/ReducedSmoothStratification.lean#L74}{\texttt{HodgeConjecture/Definitions/AlgebraicGeometry/ReducedSmoothStratification.lean:74}}.
\end{definition}

\begin{definition}[AlgebraicGeometry.reducedClosedSmoothPieceι]\label{def:AlgebraicGeometry.reducedClosedSmoothPieceι}
\lean{AlgebraicGeometry.reducedClosedSmoothPieceι}\leanok
\uses{def:AlgebraicGeometry.reducedClosedSmoothPiece}
This piece is locally closed in the original ambient scheme.

\noindent\emph{Lean:} \texttt{def AlgebraicGeometry.reducedClosedSmoothPieceι}, \href{https://github.com/Paul-Lez/HodgeConjecture/blob/main/HodgeConjecture/Definitions/AlgebraicGeometry/ReducedSmoothStratification.lean#L78}{\texttt{HodgeConjecture/Definitions/AlgebraicGeometry/ReducedSmoothStratification.lean:78}}.
\end{definition}

\begin{definition}[AlgebraicGeometry.reducedClosedSingularRemainder]\label{def:AlgebraicGeometry.reducedClosedSingularRemainder}
\lean{AlgebraicGeometry.reducedClosedSingularRemainder}\leanok
\uses{def:AlgebraicGeometry.reducedClosedStructureMap}
The singular remainder, regarded as a closed subset of the original scheme.

\noindent\emph{Lean:} \texttt{def AlgebraicGeometry.reducedClosedSingularRemainder}, \href{https://github.com/Paul-Lez/HodgeConjecture/blob/main/HodgeConjecture/Definitions/AlgebraicGeometry/ReducedSmoothStratification.lean#L95}{\texttt{HodgeConjecture/Definitions/AlgebraicGeometry/ReducedSmoothStratification.lean:95}}.
\end{definition}

\begin{definition}[AlgebraicGeometry.reducedSmoothStratificationWellFoundedRelation]\label{def:AlgebraicGeometry.reducedSmoothStratificationWellFoundedRelation}
\lean{AlgebraicGeometry.reducedSmoothStratificationWellFoundedRelation}\leanok
\noindent\emph{Lean:} \texttt{def AlgebraicGeometry.reducedSmoothStratificationWellFoundedRelation}, \href{https://github.com/Paul-Lez/HodgeConjecture/blob/main/HodgeConjecture/Definitions/AlgebraicGeometry/ReducedSmoothStratification.lean#L123}{\texttt{HodgeConjecture/Definitions/AlgebraicGeometry/ReducedSmoothStratification.lean:123}}.
\end{definition}

\begin{definition}[AlgebraicGeometry.reducedSmoothStratification]\label{def:AlgebraicGeometry.reducedSmoothStratification}
\lean{AlgebraicGeometry.reducedSmoothStratification}\leanok
\uses{def:AlgebraicGeometry.reducedClosedSingularRemainder, def:AlgebraicGeometry.reducedSmoothStratificationWellFoundedRelation}
The finite, explicitly recursive sequence of nonempty closed remainders. The actual
smooth strata are \texttt{reducedClosedSmoothPiece f S} for the members of this list.

\noindent\emph{Lean:} \texttt{def AlgebraicGeometry.reducedSmoothStratification}, \href{https://github.com/Paul-Lez/HodgeConjecture/blob/main/HodgeConjecture/Definitions/AlgebraicGeometry/ReducedSmoothStratification.lean#L127}{\texttt{HodgeConjecture/Definitions/AlgebraicGeometry/ReducedSmoothStratification.lean:127}}.
\end{definition}

\section{The canonical smooth decomposition as a closed filtration}\label{sec:HodgeConjecture.Definitions.AlgebraicGeometry.ReducedSmoothClosedFiltration}
\noindent\emph{File:} \href{https://github.com/Paul-Lez/HodgeConjecture/blob/main/HodgeConjecture/Definitions/AlgebraicGeometry/ReducedSmoothClosedFiltration.lean}{\texttt{HodgeConjecture/Definitions/AlgebraicGeometry/ReducedSmoothClosedFiltration.lean}}.

This module indexes the already constructed recursive smooth decomposition by natural
numbers. Consecutive closed supports differ by the actual smooth piece. The filtration
is empty at the length of the existing finite list and stays empty thereafter. This
format exposes precisely the nested closed supports required for localization induction;
it introduces no stratification choices or assumed cohomology vanishing.

\begin{definition}[AlgebraicGeometry.reducedSmoothClosedFiltration]\label{def:AlgebraicGeometry.reducedSmoothClosedFiltration}
\lean{AlgebraicGeometry.reducedSmoothClosedFiltration}\leanok
\uses{def:AlgebraicGeometry.reducedClosedSingularRemainder}
Iteration of the actual reduced singular remainder, padded only by empty supports
after the already constructed finite decomposition terminates.

\noindent\emph{Lean:} \texttt{def AlgebraicGeometry.reducedSmoothClosedFiltration}, \href{https://github.com/Paul-Lez/HodgeConjecture/blob/main/HodgeConjecture/Definitions/AlgebraicGeometry/ReducedSmoothClosedFiltration.lean#L29}{\texttt{HodgeConjecture/Definitions/AlgebraicGeometry/ReducedSmoothClosedFiltration.lean:29}}.
\end{definition}

\section{The singular locus has smaller algebraic dimension}\label{sec:HodgeConjecture.Definitions.AlgebraicGeometry.SingularLocusDimension}
\noindent\emph{File:} \href{https://github.com/Paul-Lez/HodgeConjecture/blob/main/HodgeConjecture/Definitions/AlgebraicGeometry/SingularLocusDimension.lean}{\texttt{HodgeConjecture/Definitions/AlgebraicGeometry/SingularLocusDimension.lean}}.

Over a perfect field the complement of the smooth locus of a reduced irreducible scheme is
a proper closed subset. This file constructs that reduced closed subscheme and proves its
strict Krull-dimension bound, including the \texttt{d - p} bound for cycle components. The bounds
are on algebraic dimension throughout.

\begin{definition}[AlgebraicGeometry.singularLocusClosed]\label{def:AlgebraicGeometry.singularLocusClosed}
\lean{AlgebraicGeometry.singularLocusClosed}\leanok
The closed complement of the actual smooth locus.

\noindent\emph{Lean:} \texttt{def AlgebraicGeometry.singularLocusClosed}, \href{https://github.com/Paul-Lez/HodgeConjecture/blob/main/HodgeConjecture/Definitions/AlgebraicGeometry/SingularLocusDimension.lean#L35}{\texttt{HodgeConjecture/Definitions/AlgebraicGeometry/SingularLocusDimension.lean:35}}.
\end{definition}

\begin{definition}[AlgebraicGeometry.reducedSingularLocus]\label{def:AlgebraicGeometry.reducedSingularLocus}
\lean{AlgebraicGeometry.reducedSingularLocus}\leanok
\uses{def:AlgebraicGeometry.reducedClosedSubscheme, def:AlgebraicGeometry.singularLocusClosed}
The singular locus equipped with its reduced closed-subscheme structure.

\noindent\emph{Lean:} \texttt{def AlgebraicGeometry.reducedSingularLocus}, \href{https://github.com/Paul-Lez/HodgeConjecture/blob/main/HodgeConjecture/Definitions/AlgebraicGeometry/SingularLocusDimension.lean#L38}{\texttt{HodgeConjecture/Definitions/AlgebraicGeometry/SingularLocusDimension.lean:38}}.
\end{definition}

\begin{definition}[AlgebraicGeometry.cycleComponentSingularStratification]\label{def:AlgebraicGeometry.cycleComponentSingularStratification}
\lean{AlgebraicGeometry.cycleComponentSingularStratification}\leanok
\uses{def:AlgebraicGeometry.IsProjective, def:AlgebraicGeometry.cycleComponentι, def:AlgebraicGeometry.reducedSmoothStratification, def:AlgebraicGeometry.singularLocusClosed, def:ProjectiveSpace.degreeZeroEquiv}
The singular locus of every cycle component admits the actual finite smooth
decomposition constructed by Noetherian recursion.

\noindent\emph{Lean:} \texttt{def AlgebraicGeometry.cycleComponentSingularStratification}, \href{https://github.com/Paul-Lez/HodgeConjecture/blob/main/HodgeConjecture/Definitions/AlgebraicGeometry/SingularLocusDimension.lean#L55}{\texttt{HodgeConjecture/Definitions/AlgebraicGeometry/SingularLocusDimension.lean:55}}.
\end{definition}

\section{Complex points of the constructed smooth decomposition}\label{sec:HodgeConjecture.Definitions.AlgebraicGeometry.SmoothStratificationAnalytification}
\noindent\emph{File:} \href{https://github.com/Paul-Lez/HodgeConjecture/blob/main/HodgeConjecture/Definitions/AlgebraicGeometry/SmoothStratificationAnalytification.lean}{\texttt{HodgeConjecture/Definitions/AlgebraicGeometry/SmoothStratificationAnalytification.lean}}.

The finite algebraic decomposition induces an actual partition of complex-point sets into
the images of smooth complex schemes. The images are analytically locally closed. The
lifting assertion follows from the existing equivalence between complex points and closed
scheme points; it is not a supplied parametrization of a stratum.

No analytic triangulation, homology-dimension theorem, or frontier condition is asserted.

\begin{definition}[AlgebraicGeometry.ComplexPoint.smoothStratificationAnalyticTopology]\label{def:AlgebraicGeometry.ComplexPoint.smoothStratificationAnalyticTopology}
\lean{AlgebraicGeometry.ComplexPoint.smoothStratificationAnalyticTopology}\leanok
\uses{def:AlgebraicGeometry.ComplexPoint, def:AlgebraicGeometry.Point.chartTopology}
\noindent\emph{Lean:} \texttt{def AlgebraicGeometry.ComplexPoint.smoothStratificationAnalyticTopology}, \href{https://github.com/Paul-Lez/HodgeConjecture/blob/main/HodgeConjecture/Definitions/AlgebraicGeometry/SmoothStratificationAnalytification.lean#L30}{\texttt{HodgeConjecture/Definitions/AlgebraicGeometry/SmoothStratificationAnalytification.lean:30}}.
\end{definition}

\section{Ambient closed supports for singular-component localization induction}\label{sec:HodgeConjecture.Definitions.AlgebraicGeometry.CycleComponentSingularClosedFiltration}
\noindent\emph{File:} \href{https://github.com/Paul-Lez/HodgeConjecture/blob/main/HodgeConjecture/Definitions/AlgebraicGeometry/CycleComponentSingularClosedFiltration.lean}{\texttt{HodgeConjecture/Definitions/AlgebraicGeometry/CycleComponentSingularClosedFiltration.lean}}.

The singular boundary of an integral cycle component carries its canonical finite reduced
smooth filtration, whose images are closed both algebraically and analytically in the ambient
variety, and whose successive differences are the complex-point images of smooth locally
closed strata. Their local relative dimensions are strictly below \texttt{d - p}, so their ambient
normal codimensions are at least \texttt{p + 1}.

\begin{definition}[AlgebraicGeometry.cycleComponentSingularClosedFiltration]\label{def:AlgebraicGeometry.cycleComponentSingularClosedFiltration}
\lean{AlgebraicGeometry.cycleComponentSingularClosedFiltration}\leanok
\uses{def:AlgebraicGeometry.IsProjective, def:AlgebraicGeometry.cycleComponentι, def:AlgebraicGeometry.reducedSmoothClosedFiltration, def:AlgebraicGeometry.singularLocusClosed, def:ProjectiveSpace.degreeZeroEquiv}
The canonical closed remainders inside the integral component's singular boundary.

\noindent\emph{Lean:} \texttt{abbrev AlgebraicGeometry.cycleComponentSingularClosedFiltration}, \href{https://github.com/Paul-Lez/HodgeConjecture/blob/main/HodgeConjecture/Definitions/AlgebraicGeometry/CycleComponentSingularClosedFiltration.lean#L31}{\texttt{HodgeConjecture/Definitions/AlgebraicGeometry/CycleComponentSingularClosedFiltration.lean:31}}.
\end{definition}

\begin{definition}[AlgebraicGeometry.cycleComponentSingularFiltrationLength]\label{def:AlgebraicGeometry.cycleComponentSingularFiltrationLength}
\lean{AlgebraicGeometry.cycleComponentSingularFiltrationLength}\leanok
\uses{def:AlgebraicGeometry.cycleComponentSingularStratification}
The exact terminal index is read from the already constructed finite decomposition.

\noindent\emph{Lean:} \texttt{abbrev AlgebraicGeometry.cycleComponentSingularFiltrationLength}, \href{https://github.com/Paul-Lez/HodgeConjecture/blob/main/HodgeConjecture/Definitions/AlgebraicGeometry/CycleComponentSingularClosedFiltration.lean#L36}{\texttt{HodgeConjecture/Definitions/AlgebraicGeometry/CycleComponentSingularClosedFiltration.lean:36}}.
\end{definition}

\begin{definition}[AlgebraicGeometry.cycleComponentSingularFiltrationStratum]\label{def:AlgebraicGeometry.cycleComponentSingularFiltrationStratum}
\lean{AlgebraicGeometry.cycleComponentSingularFiltrationStratum}\leanok
\uses{def:AlgebraicGeometry.cycleComponentSingularClosedFiltration, def:AlgebraicGeometry.reducedClosedSmoothPiece}
The actual smooth scheme occurring between two consecutive closed supports.

\noindent\emph{Lean:} \texttt{abbrev AlgebraicGeometry.cycleComponentSingularFiltrationStratum}, \href{https://github.com/Paul-Lez/HodgeConjecture/blob/main/HodgeConjecture/Definitions/AlgebraicGeometry/CycleComponentSingularClosedFiltration.lean#L40}{\texttt{HodgeConjecture/Definitions/AlgebraicGeometry/CycleComponentSingularClosedFiltration.lean:40}}.
\end{definition}

\begin{definition}[AlgebraicGeometry.cycleComponentSingularFiltrationStratumι]\label{def:AlgebraicGeometry.cycleComponentSingularFiltrationStratumι}
\lean{AlgebraicGeometry.cycleComponentSingularFiltrationStratumι}\leanok
\uses{def:AlgebraicGeometry.cycleComponentSingularFiltrationStratum, def:AlgebraicGeometry.reducedClosedSmoothPieceι}
Its actual locally closed immersion into the original smooth ambient scheme.

\noindent\emph{Lean:} \texttt{def AlgebraicGeometry.cycleComponentSingularFiltrationStratumι}, \href{https://github.com/Paul-Lez/HodgeConjecture/blob/main/HodgeConjecture/Definitions/AlgebraicGeometry/CycleComponentSingularClosedFiltration.lean#L45}{\texttt{HodgeConjecture/Definitions/AlgebraicGeometry/CycleComponentSingularClosedFiltration.lean:45}}.
\end{definition}

\begin{definition}[AlgebraicGeometry.cycleComponentSingularFiltrationStratumOver]\label{def:AlgebraicGeometry.cycleComponentSingularFiltrationStratumOver}
\lean{AlgebraicGeometry.cycleComponentSingularFiltrationStratumOver}\leanok
\uses{def:AlgebraicGeometry.cycleComponentSingularFiltrationStratumι}
A singular-filtration stratum with its induced structure map to \texttt{Spec ℂ}.

\noindent\emph{Lean:} \texttt{abbrev AlgebraicGeometry.cycleComponentSingularFiltrationStratumOver}, \href{https://github.com/Paul-Lez/HodgeConjecture/blob/main/HodgeConjecture/Definitions/AlgebraicGeometry/CycleComponentSingularClosedFiltration.lean#L51}{\texttt{HodgeConjecture/Definitions/AlgebraicGeometry/CycleComponentSingularClosedFiltration.lean:51}}.
\end{definition}

\begin{definition}[AlgebraicGeometry.cycleComponentSingularFiltrationStratumOverι]\label{def:AlgebraicGeometry.cycleComponentSingularFiltrationStratumOverι}
\lean{AlgebraicGeometry.cycleComponentSingularFiltrationStratumOverι}\leanok
\uses{def:AlgebraicGeometry.cycleComponentSingularFiltrationStratumOver}
The stratum immersion bundled over \texttt{Spec ℂ}.

\noindent\emph{Lean:} \texttt{def AlgebraicGeometry.cycleComponentSingularFiltrationStratumOverι}, \href{https://github.com/Paul-Lez/HodgeConjecture/blob/main/HodgeConjecture/Definitions/AlgebraicGeometry/CycleComponentSingularClosedFiltration.lean#L55}{\texttt{HodgeConjecture/Definitions/AlgebraicGeometry/CycleComponentSingularClosedFiltration.lean:55}}.
\end{definition}

\begin{definition}[AlgebraicGeometry.cycleComponentSingularAmbientClosedFiltration]\label{def:AlgebraicGeometry.cycleComponentSingularAmbientClosedFiltration}
\lean{AlgebraicGeometry.cycleComponentSingularAmbientClosedFiltration}\leanok
\uses{def:AlgebraicGeometry.cycleComponentSingularClosedFiltration}
Every stage is a closed support in the original algebraic ambient scheme.

\noindent\emph{Lean:} \texttt{def AlgebraicGeometry.cycleComponentSingularAmbientClosedFiltration}, \href{https://github.com/Paul-Lez/HodgeConjecture/blob/main/HodgeConjecture/Definitions/AlgebraicGeometry/CycleComponentSingularClosedFiltration.lean#L84}{\texttt{HodgeConjecture/Definitions/AlgebraicGeometry/CycleComponentSingularClosedFiltration.lean:84}}.
\end{definition}

\begin{definition}[AlgebraicGeometry.cycleComponentSingularStratumAmbientOpen]\label{def:AlgebraicGeometry.cycleComponentSingularStratumAmbientOpen}
\lean{AlgebraicGeometry.cycleComponentSingularStratumAmbientOpen}\leanok
\uses{def:AlgebraicGeometry.cycleComponentSingularAmbientClosedFiltration}
The exact ambient open used by the consecutive-support localization triangle.

\noindent\emph{Lean:} \texttt{def AlgebraicGeometry.cycleComponentSingularStratumAmbientOpen}, \href{https://github.com/Paul-Lez/HodgeConjecture/blob/main/HodgeConjecture/Definitions/AlgebraicGeometry/CycleComponentSingularClosedFiltration.lean#L100}{\texttt{HodgeConjecture/Definitions/AlgebraicGeometry/CycleComponentSingularClosedFiltration.lean:100}}.
\end{definition}

\begin{definition}[AlgebraicGeometry.cycleComponentSingularStratumClosedLift]\label{def:AlgebraicGeometry.cycleComponentSingularStratumClosedLift}
\lean{AlgebraicGeometry.cycleComponentSingularStratumClosedLift}\leanok
\uses{def:AlgebraicGeometry.cycleComponentSingularFiltrationStratumι, def:AlgebraicGeometry.cycleComponentSingularStratumAmbientOpen}
The actual smooth stratum factors into the complement of the next closed support.

\noindent\emph{Lean:} \texttt{def AlgebraicGeometry.cycleComponentSingularStratumClosedLift}, \href{https://github.com/Paul-Lez/HodgeConjecture/blob/main/HodgeConjecture/Definitions/AlgebraicGeometry/CycleComponentSingularClosedFiltration.lean#L104}{\texttt{HodgeConjecture/Definitions/AlgebraicGeometry/CycleComponentSingularClosedFiltration.lean:104}}.
\end{definition}

\begin{definition}[AlgebraicGeometry.cycleComponentSingularStratumAmbientOpenOver]\label{def:AlgebraicGeometry.cycleComponentSingularStratumAmbientOpenOver}
\lean{AlgebraicGeometry.cycleComponentSingularStratumAmbientOpenOver}\leanok
\uses{def:AlgebraicGeometry.ComplexPoint.openScheme, def:AlgebraicGeometry.cycleComponentSingularStratumAmbientOpen}
The localization open with its induced structure map to \texttt{Spec ℂ}.

\noindent\emph{Lean:} \texttt{abbrev AlgebraicGeometry.cycleComponentSingularStratumAmbientOpenOver}, \href{https://github.com/Paul-Lez/HodgeConjecture/blob/main/HodgeConjecture/Definitions/AlgebraicGeometry/CycleComponentSingularClosedFiltration.lean#L114}{\texttt{HodgeConjecture/Definitions/AlgebraicGeometry/CycleComponentSingularClosedFiltration.lean:114}}.
\end{definition}

\begin{definition}[AlgebraicGeometry.cycleComponentSingularStratumClosedLiftOver]\label{def:AlgebraicGeometry.cycleComponentSingularStratumClosedLiftOver}
\lean{AlgebraicGeometry.cycleComponentSingularStratumClosedLiftOver}\leanok
\uses{def:AlgebraicGeometry.cycleComponentSingularFiltrationStratumOver, def:AlgebraicGeometry.cycleComponentSingularStratumAmbientOpenOver, def:AlgebraicGeometry.cycleComponentSingularStratumClosedLift}
The closed stratum lift bundled over \texttt{Spec ℂ}.

\noindent\emph{Lean:} \texttt{def AlgebraicGeometry.cycleComponentSingularStratumClosedLiftOver}, \href{https://github.com/Paul-Lez/HodgeConjecture/blob/main/HodgeConjecture/Definitions/AlgebraicGeometry/CycleComponentSingularClosedFiltration.lean#L132}{\texttt{HodgeConjecture/Definitions/AlgebraicGeometry/CycleComponentSingularClosedFiltration.lean:132}}.
\end{definition}

\begin{definition}[AlgebraicGeometry.ComplexPoint.cycleComponentSingularClosedFiltrationAnalyticTopology]\label{def:AlgebraicGeometry.ComplexPoint.cycleComponentSingularClosedFiltrationAnalyticTopology}
\lean{AlgebraicGeometry.ComplexPoint.cycleComponentSingularClosedFiltrationAnalyticTopology}\leanok
\uses{def:AlgebraicGeometry.ComplexPoint, def:AlgebraicGeometry.Point.chartTopology}
\noindent\emph{Lean:} \texttt{def AlgebraicGeometry.ComplexPoint.cycleComponentSingularClosedFiltrationAnalyticTopology}, \href{https://github.com/Paul-Lez/HodgeConjecture/blob/main/HodgeConjecture/Definitions/AlgebraicGeometry/CycleComponentSingularClosedFiltration.lean#L198}{\texttt{HodgeConjecture/Definitions/AlgebraicGeometry/CycleComponentSingularClosedFiltration.lean:198}}.
\end{definition}

\begin{definition}[AlgebraicGeometry.ComplexPoint.cycleComponentSingularAnalyticClosedFiltration]\label{def:AlgebraicGeometry.ComplexPoint.cycleComponentSingularAnalyticClosedFiltration}
\lean{AlgebraicGeometry.ComplexPoint.cycleComponentSingularAnalyticClosedFiltration}\leanok
\uses{def:AlgebraicGeometry.ComplexPoint.cycleComponentSingularClosedFiltrationAnalyticTopology, def:AlgebraicGeometry.ComplexPoint.smoothStratificationAnalyticTopology, def:AlgebraicGeometry.Point.evaluate, def:AlgebraicGeometry.Point.overOpen, def:AlgebraicGeometry.cycleComponentSingularAmbientClosedFiltration}
The actual analytically closed ambient supports for nested-support localization.

\noindent\emph{Lean:} \texttt{def AlgebraicGeometry.ComplexPoint.cycleComponentSingularAnalyticClosedFiltration}, \href{https://github.com/Paul-Lez/HodgeConjecture/blob/main/HodgeConjecture/Definitions/AlgebraicGeometry/CycleComponentSingularClosedFiltration.lean#L201}{\texttt{HodgeConjecture/Definitions/AlgebraicGeometry/CycleComponentSingularClosedFiltration.lean:201}}.
\end{definition}

\begin{definition}[AlgebraicGeometry.ComplexPoint.cycleComponentSmoothLocusOver]\label{def:AlgebraicGeometry.ComplexPoint.cycleComponentSmoothLocusOver}
\lean{AlgebraicGeometry.ComplexPoint.cycleComponentSmoothLocusOver}\leanok
\uses{def:AlgebraicGeometry.IsProjective, def:AlgebraicGeometry.cycleComponentι}
The smooth locus of the component, bundled over \texttt{Spec ℂ}.

\noindent\emph{Lean:} \texttt{abbrev AlgebraicGeometry.ComplexPoint.cycleComponentSmoothLocusOver}, \href{https://github.com/Paul-Lez/HodgeConjecture/blob/main/HodgeConjecture/Definitions/AlgebraicGeometry/CycleComponentSingularClosedFiltration.lean#L207}{\texttt{HodgeConjecture/Definitions/AlgebraicGeometry/CycleComponentSingularClosedFiltration.lean:207}}.
\end{definition}

\section{Holomorphic functions on complex points}\label{sec:HodgeConjecture.Definitions.AlgebraicGeometry.ComplexAnalyticSheaf}
\noindent\emph{File:} \href{https://github.com/Paul-Lez/HodgeConjecture/blob/main/HodgeConjecture/Definitions/AlgebraicGeometry/ComplexAnalyticSheaf.lean}{\texttt{HodgeConjecture/Definitions/AlgebraicGeometry/ComplexAnalyticSheaf.lean}}.

For a complex scheme smooth of relative dimension \texttt{d}, the algebraic étale charts constructed in
\texttt{SmoothComplexCoordinates} give a canonical charted-space structure on its complex points. This
file defines its sheaf of holomorphic functions: a section is a complex-valued function which is
complex analytic in those charts. Mathlib expresses complex analyticity as manifold
differentiability of order \texttt{ω}.

The sheaf and its ring operations are constructed by the local-predicate sheaf machinery. No
analytic atlas or sheaf is supplied as data.

\begin{definition}[AlgebraicGeometry.ComplexPoint.holomorphicFunctionSheafToTypes]\label{def:AlgebraicGeometry.ComplexPoint.holomorphicFunctionSheafToTypes}
\lean{AlgebraicGeometry.ComplexPoint.holomorphicFunctionSheafToTypes}\leanok
\uses{def:AlgebraicGeometry.ComplexPoint.localChart}
The sheaf of complex-valued functions which are analytic in the algebraically constructed
étale charts, initially regarded as a sheaf of types.

\noindent\emph{Lean:} \texttt{def AlgebraicGeometry.ComplexPoint.holomorphicFunctionSheafToTypes}, \href{https://github.com/Paul-Lez/HodgeConjecture/blob/main/HodgeConjecture/Definitions/AlgebraicGeometry/ComplexAnalyticSheaf.lean#L46}{\texttt{HodgeConjecture/Definitions/AlgebraicGeometry/ComplexAnalyticSheaf.lean:46}}.
\end{definition}

\begin{definition}[AlgebraicGeometry.ComplexPoint.holomorphicFunctionPresheaf]\label{def:AlgebraicGeometry.ComplexPoint.holomorphicFunctionPresheaf}
\lean{AlgebraicGeometry.ComplexPoint.holomorphicFunctionPresheaf}\leanok
\uses{def:AlgebraicGeometry.ComplexPoint.holomorphicFunctionSheafToTypes}
The presheaf of rings underlying the holomorphic-function sheaf.

\noindent\emph{Lean:} \texttt{def AlgebraicGeometry.ComplexPoint.holomorphicFunctionPresheaf}, \href{https://github.com/Paul-Lez/HodgeConjecture/blob/main/HodgeConjecture/Definitions/AlgebraicGeometry/ComplexAnalyticSheaf.lean#L59}{\texttt{HodgeConjecture/Definitions/AlgebraicGeometry/ComplexAnalyticSheaf.lean:59}}.
\end{definition}

\begin{definition}[AlgebraicGeometry.ComplexPoint.holomorphicFunctionSheaf]\label{def:AlgebraicGeometry.ComplexPoint.holomorphicFunctionSheaf}
\lean{AlgebraicGeometry.ComplexPoint.holomorphicFunctionSheaf}\leanok
\uses{def:AlgebraicGeometry.ComplexPoint.holomorphicFunctionPresheaf}
The sheaf of complex-valued functions which are analytic in the algebraically constructed
étale charts.

\noindent\emph{Lean:} \texttt{def AlgebraicGeometry.ComplexPoint.holomorphicFunctionSheaf}, \href{https://github.com/Paul-Lez/HodgeConjecture/blob/main/HodgeConjecture/Definitions/AlgebraicGeometry/ComplexAnalyticSheaf.lean#L70}{\texttt{HodgeConjecture/Definitions/AlgebraicGeometry/ComplexAnalyticSheaf.lean:70}}.
\end{definition}

\section{Analytic differential forms}\label{sec:HodgeConjecture.Definitions.AlgebraicGeometry.AnalyticDifferentialForms}
\noindent\emph{File:} \href{https://github.com/Paul-Lez/HodgeConjecture/blob/main/HodgeConjecture/Definitions/AlgebraicGeometry/AnalyticDifferentialForms.lean}{\texttt{HodgeConjecture/Definitions/AlgebraicGeometry/AnalyticDifferentialForms.lean}}.

This file constructs analytic differential forms from holomorphic functions and their actual
manifold derivatives. The pointwise algebra of wedges lives in
\texttt{HodgeConjecture.Mathlib.Analysis.NormedSpace.WedgeCovectors} and
\texttt{HodgeConjecture.Lemmas.Analysis.Calculus.DifferentialForm.ExactWedge}.

The starting point is \texttt{Algebra.DeRham.Form ℂ 𝒪(U) p}, the Kähler differential forms of the
holomorphic functions on \texttt{U}, that is \texttt{⋀\^{}p Ω[𝒪(U)⁄ℂ]}. In a fixed chart at \texttt{z} such a form is
evaluated pointwise as an alternating continuous multilinear map: taking coordinate derivatives
is a \texttt{ℂ}-derivation of \texttt{𝒪(U)} into covector fields along the coordinate domain, so the universal
property of \texttt{Ω[𝒪(U)⁄ℂ]} and then that of the exterior power turn the pointwise wedge of covectors
into \texttt{chartEvaluation}. By construction it takes the algebraic exterior derivative to the
analytic one.

Analytic differential forms are then \texttt{Algebra.DeRham.Form ℂ 𝒪(U) p} modulo the forms whose
evaluation vanishes in every chart. These relations include analytic identities such as the
chain rule.

In degrees above the complex dimension the evaluation target is zero, so the resulting forms
vanish.

\begin{definition}[AlgebraicGeometry.ComplexPoint.holomorphicRestrictionAlgHom]\label{def:AlgebraicGeometry.ComplexPoint.holomorphicRestrictionAlgHom}
\lean{AlgebraicGeometry.ComplexPoint.holomorphicRestrictionAlgHom}\leanok
\uses{def:AlgebraicGeometry.ComplexPoint.holomorphicFunctionSheaf}
Restriction of holomorphic functions as an algebra homomorphism over the constants.

\noindent\emph{Lean:} \texttt{def AlgebraicGeometry.ComplexPoint.holomorphicRestrictionAlgHom}, \href{https://github.com/Paul-Lez/HodgeConjecture/blob/main/HodgeConjecture/Definitions/AlgebraicGeometry/AnalyticDifferentialForms.lean#L65}{\texttt{HodgeConjecture/Definitions/AlgebraicGeometry/AnalyticDifferentialForms.lean:65}}.
\end{definition}

\begin{definition}[AlgebraicGeometry.ComplexPoint.OpenHolomorphicFunctions]\label{def:AlgebraicGeometry.ComplexPoint.OpenHolomorphicFunctions}
\lean{AlgebraicGeometry.ComplexPoint.OpenHolomorphicFunctions}\leanok
\uses{def:AlgebraicGeometry.ComplexPoint.holomorphicFunctionPresheaf}
\noindent\emph{Lean:} \texttt{abbrev AlgebraicGeometry.ComplexPoint.OpenHolomorphicFunctions}, \href{https://github.com/Paul-Lez/HodgeConjecture/blob/main/HodgeConjecture/Definitions/AlgebraicGeometry/AnalyticDifferentialForms.lean#L88}{\texttt{HodgeConjecture/Definitions/AlgebraicGeometry/AnalyticDifferentialForms.lean:88}}.
\end{definition}

\begin{definition}[AlgebraicGeometry.ComplexPoint.chartSectionDomain]\label{def:AlgebraicGeometry.ComplexPoint.chartSectionDomain}
\lean{AlgebraicGeometry.ComplexPoint.chartSectionDomain}\leanok
\uses{def:AlgebraicGeometry.ComplexPoint.localChart}
The part of a fixed chart target whose inverse image belongs to \texttt{U}.

\noindent\emph{Lean:} \texttt{def AlgebraicGeometry.ComplexPoint.chartSectionDomain}, \href{https://github.com/Paul-Lez/HodgeConjecture/blob/main/HodgeConjecture/Definitions/AlgebraicGeometry/AnalyticDifferentialForms.lean#L92}{\texttt{HodgeConjecture/Definitions/AlgebraicGeometry/AnalyticDifferentialForms.lean:92}}.
\end{definition}

\begin{definition}[AlgebraicGeometry.ComplexPoint.extendedSection]\label{def:AlgebraicGeometry.ComplexPoint.extendedSection}
\lean{AlgebraicGeometry.ComplexPoint.extendedSection}\leanok
\uses{def:AlgebraicGeometry.ComplexPoint.OpenHolomorphicFunctions}
An arbitrary total extension of a section from its open domain. Its values outside the
domain play no role in derivatives taken within that domain.

\noindent\emph{Lean:} \texttt{def AlgebraicGeometry.ComplexPoint.extendedSection}, \href{https://github.com/Paul-Lez/HodgeConjecture/blob/main/HodgeConjecture/Definitions/AlgebraicGeometry/AnalyticDifferentialForms.lean#L108}{\texttt{HodgeConjecture/Definitions/AlgebraicGeometry/AnalyticDifferentialForms.lean:108}}.
\end{definition}

\begin{definition}[AlgebraicGeometry.ComplexPoint.chartSection]\label{def:AlgebraicGeometry.ComplexPoint.chartSection}
\lean{AlgebraicGeometry.ComplexPoint.chartSection}\leanok
\uses{def:AlgebraicGeometry.ComplexPoint.extendedSection}
The expression of a holomorphic section in one fixed chart.

\noindent\emph{Lean:} \texttt{def AlgebraicGeometry.ComplexPoint.chartSection}, \href{https://github.com/Paul-Lez/HodgeConjecture/blob/main/HodgeConjecture/Definitions/AlgebraicGeometry/AnalyticDifferentialForms.lean#L133}{\texttt{HodgeConjecture/Definitions/AlgebraicGeometry/AnalyticDifferentialForms.lean:133}}.
\end{definition}

\begin{definition}[AlgebraicGeometry.ComplexPoint.chartSectionDifferential]\label{def:AlgebraicGeometry.ComplexPoint.chartSectionDifferential}
\lean{AlgebraicGeometry.ComplexPoint.chartSectionDifferential}\leanok
\uses{def:AlgebraicGeometry.ComplexPoint.chartSection, def:AlgebraicGeometry.ComplexPoint.chartSectionDomain}
The derivative of a section in one fixed chart.

\noindent\emph{Lean:} \texttt{def AlgebraicGeometry.ComplexPoint.chartSectionDifferential}, \href{https://github.com/Paul-Lez/HodgeConjecture/blob/main/HodgeConjecture/Definitions/AlgebraicGeometry/AnalyticDifferentialForms.lean#L183}{\texttt{HodgeConjecture/Definitions/AlgebraicGeometry/AnalyticDifferentialForms.lean:183}}.
\end{definition}

\begin{definition}[AlgebraicGeometry.ComplexPoint.chartSectionAlgHom]\label{def:AlgebraicGeometry.ComplexPoint.chartSectionAlgHom}
\lean{AlgebraicGeometry.ComplexPoint.chartSectionAlgHom}\leanok
\uses{def:AlgebraicGeometry.ComplexPoint.chartSection, def:AlgebraicGeometry.ComplexPoint.chartSectionDomain}
Expressing holomorphic functions on \texttt{U} in a fixed chart, as a homomorphism of \texttt{ℂ}-algebras
into functions on the coordinate domain of that chart.

\noindent\emph{Lean:} \texttt{def AlgebraicGeometry.ComplexPoint.chartSectionAlgHom}, \href{https://github.com/Paul-Lez/HodgeConjecture/blob/main/HodgeConjecture/Definitions/AlgebraicGeometry/AnalyticDifferentialForms.lean#L268}{\texttt{HodgeConjecture/Definitions/AlgebraicGeometry/AnalyticDifferentialForms.lean:268}}.
\end{definition}

\begin{definition}[AlgebraicGeometry.ComplexPoint.chartFieldModule]\label{def:AlgebraicGeometry.ComplexPoint.chartFieldModule}
\lean{AlgebraicGeometry.ComplexPoint.chartFieldModule}\leanok
\uses{def:AlgebraicGeometry.ComplexPoint.chartSectionAlgHom}
Functions on the coordinate domain of a fixed chart form a module over the holomorphic
functions on \texttt{U}, acting through their expression in that chart.

\noindent\emph{Lean:} \texttt{abbrev AlgebraicGeometry.ComplexPoint.chartFieldModule}, \href{https://github.com/Paul-Lez/HodgeConjecture/blob/main/HodgeConjecture/Definitions/AlgebraicGeometry/AnalyticDifferentialForms.lean#L291}{\texttt{HodgeConjecture/Definitions/AlgebraicGeometry/AnalyticDifferentialForms.lean:291}}.
\end{definition}

\begin{definition}[AlgebraicGeometry.ComplexPoint.chartDerivation]\label{def:AlgebraicGeometry.ComplexPoint.chartDerivation}
\lean{AlgebraicGeometry.ComplexPoint.chartDerivation}\leanok
\uses{def:AlgebraicGeometry.ComplexAlgHom.isStandardEtaleBaseAlgHom, def:AlgebraicGeometry.ComplexAlgHom.standardEtaleBivariateEvaluation, def:AlgebraicGeometry.ComplexPoint.affineSpecComplexPointMap, def:AlgebraicGeometry.ComplexPoint.chartFieldModule, def:AlgebraicGeometry.ComplexPoint.chartSectionDifferential, def:AlgebraicGeometry.ComplexPoint.holomorphicFunctionSheaf, def:AlgebraicGeometry.ComplexPoint.polynomialSpecToAffineSpacePointMap, def:AlgebraicGeometry.LocalEtaleCoordinates.algebraicCoordinates, def:AlgebraicGeometry.LocalEtaleCoordinates.ambientAnalyticCoordinates, def:AlgebraicGeometry.LocalEtaleCoordinates.ambientCoordinateSection}
The coordinate differentials of holomorphic sections form a \texttt{ℂ}-derivation into covector
fields along the coordinate domain of a fixed chart.

\noindent\emph{Lean:} \texttt{def AlgebraicGeometry.ComplexPoint.chartDerivation}, \href{https://github.com/Paul-Lez/HodgeConjecture/blob/main/HodgeConjecture/Definitions/AlgebraicGeometry/AnalyticDifferentialForms.lean#L322}{\texttt{HodgeConjecture/Definitions/AlgebraicGeometry/AnalyticDifferentialForms.lean:322}}.
\end{definition}

\begin{definition}[AlgebraicGeometry.ComplexPoint.chartWedge]\label{def:AlgebraicGeometry.ComplexPoint.chartWedge}
\lean{AlgebraicGeometry.ComplexPoint.chartWedge}\leanok
\uses{def:AlgebraicGeometry.ComplexPoint.chartFieldModule, def:ContinuousAlternatingMap.wedgeCovectors}
The pointwise wedge of covector fields along a fixed chart.

\noindent\emph{Lean:} \texttt{def AlgebraicGeometry.ComplexPoint.chartWedge}, \href{https://github.com/Paul-Lez/HodgeConjecture/blob/main/HodgeConjecture/Definitions/AlgebraicGeometry/AnalyticDifferentialForms.lean#L356}{\texttt{HodgeConjecture/Definitions/AlgebraicGeometry/AnalyticDifferentialForms.lean:356}}.
\end{definition}

\begin{definition}[AlgebraicGeometry.ComplexPoint.chartEvaluationHolo]\label{def:AlgebraicGeometry.ComplexPoint.chartEvaluationHolo}
\lean{AlgebraicGeometry.ComplexPoint.chartEvaluationHolo}\leanok
\uses{def:Algebra.DeRham.Form, def:AlgebraicGeometry.ComplexPoint.chartDerivation, def:AlgebraicGeometry.ComplexPoint.chartWedge}
Fixed-chart evaluation of differential forms, as a \texttt{ℂ}-linear map into alternating-form
fields on the coordinate domain.

\noindent\emph{Lean:} \texttt{def AlgebraicGeometry.ComplexPoint.chartEvaluationHolo}, \href{https://github.com/Paul-Lez/HodgeConjecture/blob/main/HodgeConjecture/Definitions/AlgebraicGeometry/AnalyticDifferentialForms.lean#L390}{\texttt{HodgeConjecture/Definitions/AlgebraicGeometry/AnalyticDifferentialForms.lean:390}}.
\end{definition}

\begin{definition}[AlgebraicGeometry.ComplexPoint.extendByZero]\label{def:AlgebraicGeometry.ComplexPoint.extendByZero}
\lean{AlgebraicGeometry.ComplexPoint.extendByZero}\leanok
\uses{def:AlgebraicGeometry.ComplexPoint.chartSectionDomain}
Extend a function on the coordinate domain of a fixed chart by zero.

\noindent\emph{Lean:} \texttt{def AlgebraicGeometry.ComplexPoint.extendByZero}, \href{https://github.com/Paul-Lez/HodgeConjecture/blob/main/HodgeConjecture/Definitions/AlgebraicGeometry/AnalyticDifferentialForms.lean#L415}{\texttt{HodgeConjecture/Definitions/AlgebraicGeometry/AnalyticDifferentialForms.lean:415}}.
\end{definition}

\begin{definition}[AlgebraicGeometry.ComplexPoint.extendByZeroLinear]\label{def:AlgebraicGeometry.ComplexPoint.extendByZeroLinear}
\lean{AlgebraicGeometry.ComplexPoint.extendByZeroLinear}\leanok
\uses{def:AlgebraicGeometry.ComplexPoint.extendByZero}
Extension by zero, as a \texttt{ℂ}-linear map on chart fields.

\noindent\emph{Lean:} \texttt{def AlgebraicGeometry.ComplexPoint.extendByZeroLinear}, \href{https://github.com/Paul-Lez/HodgeConjecture/blob/main/HodgeConjecture/Definitions/AlgebraicGeometry/AnalyticDifferentialForms.lean#L438}{\texttt{HodgeConjecture/Definitions/AlgebraicGeometry/AnalyticDifferentialForms.lean:438}}.
\end{definition}

\begin{definition}[AlgebraicGeometry.ComplexPoint.chartEvaluation]\label{def:AlgebraicGeometry.ComplexPoint.chartEvaluation}
\lean{AlgebraicGeometry.ComplexPoint.chartEvaluation}\leanok
\uses{def:AlgebraicGeometry.ComplexPoint.chartEvaluationHolo, def:AlgebraicGeometry.ComplexPoint.extendByZeroLinear}
Fixed-chart evaluation of a differential form, as a \texttt{ℂ}-linear map into alternating-form
fields on the chart target, extended by zero outside the coordinate domain.

\noindent\emph{Lean:} \texttt{def AlgebraicGeometry.ComplexPoint.chartEvaluation}, \href{https://github.com/Paul-Lez/HodgeConjecture/blob/main/HodgeConjecture/Definitions/AlgebraicGeometry/AnalyticDifferentialForms.lean#L455}{\texttt{HodgeConjecture/Definitions/AlgebraicGeometry/AnalyticDifferentialForms.lean:455}}.
\end{definition}

\begin{definition}[AlgebraicGeometry.ComplexPoint.chartGeneratorEvaluation]\label{def:AlgebraicGeometry.ComplexPoint.chartGeneratorEvaluation}
\lean{AlgebraicGeometry.ComplexPoint.chartGeneratorEvaluation}\leanok
\uses{def:AlgebraicGeometry.ComplexPoint.chartSectionDifferential, def:ContinuousAlternatingMap.wedgeCovectors}
Fixed-chart evaluation of the generator \texttt{a₀ * d v₀ ∧ ⋯ ∧ d vₚ₋₁}.

\noindent\emph{Lean:} \texttt{def AlgebraicGeometry.ComplexPoint.chartGeneratorEvaluation}, \href{https://github.com/Paul-Lez/HodgeConjecture/blob/main/HodgeConjecture/Definitions/AlgebraicGeometry/AnalyticDifferentialForms.lean#L502}{\texttt{HodgeConjecture/Definitions/AlgebraicGeometry/AnalyticDifferentialForms.lean:502}}.
\end{definition}

\begin{definition}[AlgebraicGeometry.ComplexPoint.chartEvaluationAt]\label{def:AlgebraicGeometry.ComplexPoint.chartEvaluationAt}
\lean{AlgebraicGeometry.ComplexPoint.chartEvaluationAt}\leanok
\uses{def:AlgebraicGeometry.ComplexPoint.chartEvaluation}
Evaluation at one point of one fixed coordinate chart.

\noindent\emph{Lean:} \texttt{def AlgebraicGeometry.ComplexPoint.chartEvaluationAt}, \href{https://github.com/Paul-Lez/HodgeConjecture/blob/main/HodgeConjecture/Definitions/AlgebraicGeometry/AnalyticDifferentialForms.lean#L633}{\texttt{HodgeConjecture/Definitions/AlgebraicGeometry/AnalyticDifferentialForms.lean:633}}.
\end{definition}

\begin{definition}[AlgebraicGeometry.ComplexPoint.chartEvaluationKernel]\label{def:AlgebraicGeometry.ComplexPoint.chartEvaluationKernel}
\lean{AlgebraicGeometry.ComplexPoint.chartEvaluationKernel}\leanok
\uses{def:AlgebraicGeometry.ComplexPoint.chartEvaluationAt}
Differential forms that vanish in every fixed chart at every point of its coordinate
domain.

\noindent\emph{Lean:} \texttt{def AlgebraicGeometry.ComplexPoint.chartEvaluationKernel}, \href{https://github.com/Paul-Lez/HodgeConjecture/blob/main/HodgeConjecture/Definitions/AlgebraicGeometry/AnalyticDifferentialForms.lean#L641}{\texttt{HodgeConjecture/Definitions/AlgebraicGeometry/AnalyticDifferentialForms.lean:641}}.
\end{definition}

\begin{definition}[AlgebraicGeometry.ComplexPoint.formRestriction]\label{def:AlgebraicGeometry.ComplexPoint.formRestriction}
\lean{AlgebraicGeometry.ComplexPoint.formRestriction}\leanok
\uses{def:Algebra.DeRham.map, def:AlgebraicGeometry.ComplexPoint.OpenHolomorphicFunctions, def:AlgebraicGeometry.ComplexPoint.holomorphicRestrictionAlgHom}
Restriction of differential forms along an inclusion of open sets.

\noindent\emph{Lean:} \texttt{def AlgebraicGeometry.ComplexPoint.formRestriction}, \href{https://github.com/Paul-Lez/HodgeConjecture/blob/main/HodgeConjecture/Definitions/AlgebraicGeometry/AnalyticDifferentialForms.lean#L675}{\texttt{HodgeConjecture/Definitions/AlgebraicGeometry/AnalyticDifferentialForms.lean:675}}.
\end{definition}

\begin{definition}[AlgebraicGeometry.ComplexPoint.holomorphicFormRelations]\label{def:AlgebraicGeometry.ComplexPoint.holomorphicFormRelations}
\lean{AlgebraicGeometry.ComplexPoint.holomorphicFormRelations}\leanok
\uses{def:AlgebraicGeometry.ComplexPoint.chartEvaluationKernel}
Relations for analytic de Rham forms detected by actual complex derivatives.

\noindent\emph{Lean:} \texttt{def AlgebraicGeometry.ComplexPoint.holomorphicFormRelations}, \href{https://github.com/Paul-Lez/HodgeConjecture/blob/main/HodgeConjecture/Definitions/AlgebraicGeometry/AnalyticDifferentialForms.lean#L790}{\texttt{HodgeConjecture/Definitions/AlgebraicGeometry/AnalyticDifferentialForms.lean:790}}.
\end{definition}

\begin{definition}[AlgebraicGeometry.ComplexPoint.HolomorphicForm]\label{def:AlgebraicGeometry.ComplexPoint.HolomorphicForm}
\lean{AlgebraicGeometry.ComplexPoint.HolomorphicForm}\leanok
\uses{def:AlgebraicGeometry.ComplexPoint.holomorphicFormRelations}
Analytic differential forms of degree \texttt{p} on \texttt{U}: Kähler differential forms of the
holomorphic functions, modulo the identities that actual complex derivatives detect.

\noindent\emph{Lean:} \texttt{abbrev AlgebraicGeometry.ComplexPoint.HolomorphicForm}, \href{https://github.com/Paul-Lez/HodgeConjecture/blob/main/HodgeConjecture/Definitions/AlgebraicGeometry/AnalyticDifferentialForms.lean#L813}{\texttt{HodgeConjecture/Definitions/AlgebraicGeometry/AnalyticDifferentialForms.lean:813}}.
\end{definition}

\begin{definition}[AlgebraicGeometry.ComplexPoint.holomorphicFormOfConstant]\label{def:AlgebraicGeometry.ComplexPoint.holomorphicFormOfConstant}
\lean{AlgebraicGeometry.ComplexPoint.holomorphicFormOfConstant}\leanok
\uses{def:Algebra.DeRham.ofConstant, def:AlgebraicGeometry.ComplexPoint.HolomorphicForm}
A complex constant regarded as an analytic differential zero-form.

\noindent\emph{Lean:} \texttt{def AlgebraicGeometry.ComplexPoint.holomorphicFormOfConstant}, \href{https://github.com/Paul-Lez/HodgeConjecture/blob/main/HodgeConjecture/Definitions/AlgebraicGeometry/AnalyticDifferentialForms.lean#L820}{\texttt{HodgeConjecture/Definitions/AlgebraicGeometry/AnalyticDifferentialForms.lean:820}}.
\end{definition}

\begin{definition}[AlgebraicGeometry.ComplexPoint.holomorphicFormDifferential]\label{def:AlgebraicGeometry.ComplexPoint.holomorphicFormDifferential}
\lean{AlgebraicGeometry.ComplexPoint.holomorphicFormDifferential}\leanok
\uses{def:Algebra.DeRham.differential, def:Algebra.DeRham.mk, def:AlgebraicGeometry.ComplexPoint.HolomorphicForm, def:AlgebraicGeometry.ComplexPoint.chartGeneratorEvaluation, def:ContinuousAlternatingMap.standardVolumeForm, def:DifferentialForm.exactWedgeWithin}
\noindent\emph{Lean:} \texttt{def AlgebraicGeometry.ComplexPoint.holomorphicFormDifferential}, \href{https://github.com/Paul-Lez/HodgeConjecture/blob/main/HodgeConjecture/Definitions/AlgebraicGeometry/AnalyticDifferentialForms.lean#L827}{\texttt{HodgeConjecture/Definitions/AlgebraicGeometry/AnalyticDifferentialForms.lean:827}}.
\end{definition}

\begin{definition}[AlgebraicGeometry.ComplexPoint.holomorphicFormRestriction]\label{def:AlgebraicGeometry.ComplexPoint.holomorphicFormRestriction}
\lean{AlgebraicGeometry.ComplexPoint.holomorphicFormRestriction}\leanok
\uses{def:Algebra.DeRham.mk, def:AlgebraicGeometry.ComplexPoint.HolomorphicForm, def:AlgebraicGeometry.ComplexPoint.chartGeneratorEvaluation, def:AlgebraicGeometry.ComplexPoint.formRestriction}
\noindent\emph{Lean:} \texttt{def AlgebraicGeometry.ComplexPoint.holomorphicFormRestriction}, \href{https://github.com/Paul-Lez/HodgeConjecture/blob/main/HodgeConjecture/Definitions/AlgebraicGeometry/AnalyticDifferentialForms.lean#L838}{\texttt{HodgeConjecture/Definitions/AlgebraicGeometry/AnalyticDifferentialForms.lean:838}}.
\end{definition}

\section{The holomorphic de Rham complex}\label{sec:HodgeConjecture.Definitions.AlgebraicGeometry.HolomorphicDeRham}
\noindent\emph{File:} \href{https://github.com/Paul-Lez/HodgeConjecture/blob/main/HodgeConjecture/Definitions/AlgebraicGeometry/HolomorphicDeRham.lean}{\texttt{HodgeConjecture/Definitions/AlgebraicGeometry/HolomorphicDeRham.lean}}.

This file assembles the analytic differential forms constructed from actual manifold derivatives
into a presheaf complex and then sheafifies it degree by degree. Restriction of functions induces
restriction of forms and commutes with the exterior derivative.

Constant complex-valued functions define a canonical morphism from the constant sheaf complex to
the holomorphic de Rham complex. Both complexes are also extended by zero to integer degrees for
use in the derived category. No differential forms, differentials, or comparison maps are supplied
as data.

\begin{definition}[AlgebraicGeometry.ComplexPoint.holomorphicDeRhamPresheaf]\label{def:AlgebraicGeometry.ComplexPoint.holomorphicDeRhamPresheaf}
\lean{AlgebraicGeometry.ComplexPoint.holomorphicDeRhamPresheaf}\leanok
\uses{def:AlgebraicGeometry.ComplexPoint.holomorphicFormRestriction}
Holomorphic de Rham forms in a fixed degree, as a presheaf of additive groups.

\noindent\emph{Lean:} \texttt{def AlgebraicGeometry.ComplexPoint.holomorphicDeRhamPresheaf}, \href{https://github.com/Paul-Lez/HodgeConjecture/blob/main/HodgeConjecture/Definitions/AlgebraicGeometry/HolomorphicDeRham.lean#L55}{\texttt{HodgeConjecture/Definitions/AlgebraicGeometry/HolomorphicDeRham.lean:55}}.
\end{definition}

\begin{definition}[AlgebraicGeometry.ComplexPoint.holomorphicDeRhamDifferential]\label{def:AlgebraicGeometry.ComplexPoint.holomorphicDeRhamDifferential}
\lean{AlgebraicGeometry.ComplexPoint.holomorphicDeRhamDifferential}\leanok
\uses{def:AlgebraicGeometry.ComplexPoint.holomorphicDeRhamPresheaf, def:AlgebraicGeometry.ComplexPoint.holomorphicFormDifferential}
The exterior derivative as a morphism of presheaves.

\noindent\emph{Lean:} \texttt{def AlgebraicGeometry.ComplexPoint.holomorphicDeRhamDifferential}, \href{https://github.com/Paul-Lez/HodgeConjecture/blob/main/HodgeConjecture/Definitions/AlgebraicGeometry/HolomorphicDeRham.lean#L86}{\texttt{HodgeConjecture/Definitions/AlgebraicGeometry/HolomorphicDeRham.lean:86}}.
\end{definition}

\begin{definition}[AlgebraicGeometry.ComplexPoint.holomorphicDeRhamPresheafComplex]\label{def:AlgebraicGeometry.ComplexPoint.holomorphicDeRhamPresheafComplex}
\lean{AlgebraicGeometry.ComplexPoint.holomorphicDeRhamPresheafComplex}\leanok
\uses{def:AlgebraicGeometry.ComplexPoint.holomorphicDeRhamDifferential}
The holomorphic de Rham complex before sheafification.

\noindent\emph{Lean:} \texttt{def AlgebraicGeometry.ComplexPoint.holomorphicDeRhamPresheafComplex}, \href{https://github.com/Paul-Lez/HodgeConjecture/blob/main/HodgeConjecture/Definitions/AlgebraicGeometry/HolomorphicDeRham.lean#L103}{\texttt{HodgeConjecture/Definitions/AlgebraicGeometry/HolomorphicDeRham.lean:103}}.
\end{definition}

\begin{definition}[AlgebraicGeometry.ComplexPoint.constantComplexAddCommGrpPresheaf]\label{def:AlgebraicGeometry.ComplexPoint.constantComplexAddCommGrpPresheaf}
\lean{AlgebraicGeometry.ComplexPoint.constantComplexAddCommGrpPresheaf}\leanok
\uses{def:AlgebraicGeometry.ComplexPoint, def:AlgebraicGeometry.Point.chartTopology}
The constant presheaf of additive groups with value \texttt{ℂ}.

\noindent\emph{Lean:} \texttt{def AlgebraicGeometry.ComplexPoint.constantComplexAddCommGrpPresheaf}, \href{https://github.com/Paul-Lez/HodgeConjecture/blob/main/HodgeConjecture/Definitions/AlgebraicGeometry/HolomorphicDeRham.lean#L118}{\texttt{HodgeConjecture/Definitions/AlgebraicGeometry/HolomorphicDeRham.lean:118}}.
\end{definition}

\begin{definition}[AlgebraicGeometry.ComplexPoint.constantsToHolomorphicDeRhamZero]\label{def:AlgebraicGeometry.ComplexPoint.constantsToHolomorphicDeRhamZero}
\lean{AlgebraicGeometry.ComplexPoint.constantsToHolomorphicDeRhamZero}\leanok
\uses{def:AlgebraicGeometry.ComplexPoint.constantComplexAddCommGrpPresheaf, def:AlgebraicGeometry.ComplexPoint.holomorphicDeRhamPresheaf, def:AlgebraicGeometry.ComplexPoint.holomorphicFormOfConstant}
Constants as holomorphic de Rham forms of degree zero.

\noindent\emph{Lean:} \texttt{def AlgebraicGeometry.ComplexPoint.constantsToHolomorphicDeRhamZero}, \href{https://github.com/Paul-Lez/HodgeConjecture/blob/main/HodgeConjecture/Definitions/AlgebraicGeometry/HolomorphicDeRham.lean#L124}{\texttt{HodgeConjecture/Definitions/AlgebraicGeometry/HolomorphicDeRham.lean:124}}.
\end{definition}

\begin{definition}[AlgebraicGeometry.ComplexPoint.holomorphicDeRhamSheaf]\label{def:AlgebraicGeometry.ComplexPoint.holomorphicDeRhamSheaf}
\lean{AlgebraicGeometry.ComplexPoint.holomorphicDeRhamSheaf}\leanok
\uses{def:AlgebraicGeometry.ComplexPoint.holomorphicDeRhamPresheaf}
Holomorphic de Rham forms in a fixed degree, after additive sheafification.

\noindent\emph{Lean:} \texttt{def AlgebraicGeometry.ComplexPoint.holomorphicDeRhamSheaf}, \href{https://github.com/Paul-Lez/HodgeConjecture/blob/main/HodgeConjecture/Definitions/AlgebraicGeometry/HolomorphicDeRham.lean#L226}{\texttt{HodgeConjecture/Definitions/AlgebraicGeometry/HolomorphicDeRham.lean:226}}.
\end{definition}

\begin{definition}[AlgebraicGeometry.ComplexPoint.holomorphicDeRhamSheafDifferential]\label{def:AlgebraicGeometry.ComplexPoint.holomorphicDeRhamSheafDifferential}
\lean{AlgebraicGeometry.ComplexPoint.holomorphicDeRhamSheafDifferential}\leanok
\uses{def:AlgebraicGeometry.ComplexPoint.holomorphicDeRhamDifferential, def:AlgebraicGeometry.ComplexPoint.holomorphicDeRhamSheaf}
The sheafified exterior derivative.

\noindent\emph{Lean:} \texttt{def AlgebraicGeometry.ComplexPoint.holomorphicDeRhamSheafDifferential}, \href{https://github.com/Paul-Lez/HodgeConjecture/blob/main/HodgeConjecture/Definitions/AlgebraicGeometry/HolomorphicDeRham.lean#L241}{\texttt{HodgeConjecture/Definitions/AlgebraicGeometry/HolomorphicDeRham.lean:241}}.
\end{definition}

\begin{definition}[AlgebraicGeometry.ComplexPoint.holomorphicDeRhamComplex]\label{def:AlgebraicGeometry.ComplexPoint.holomorphicDeRhamComplex}
\lean{AlgebraicGeometry.ComplexPoint.holomorphicDeRhamComplex}\leanok
\uses{def:AlgebraicGeometry.ComplexPoint.holomorphicDeRhamSheafDifferential}
The sheafified holomorphic de Rham complex.

\noindent\emph{Lean:} \texttt{def AlgebraicGeometry.ComplexPoint.holomorphicDeRhamComplex}, \href{https://github.com/Paul-Lez/HodgeConjecture/blob/main/HodgeConjecture/Definitions/AlgebraicGeometry/HolomorphicDeRham.lean#L249}{\texttt{HodgeConjecture/Definitions/AlgebraicGeometry/HolomorphicDeRham.lean:249}}.
\end{definition}

\begin{definition}[AlgebraicGeometry.ComplexPoint.holomorphicDeRhamSheafificationUnit]\label{def:AlgebraicGeometry.ComplexPoint.holomorphicDeRhamSheafificationUnit}
\lean{AlgebraicGeometry.ComplexPoint.holomorphicDeRhamSheafificationUnit}\leanok
\uses{def:AlgebraicGeometry.ComplexPoint.holomorphicDeRhamComplex, def:AlgebraicGeometry.ComplexPoint.holomorphicDeRhamPresheafComplex}
The degreewise sheafification unit from holomorphic forms to the underlying presheaf of the
holomorphic de Rham sheaf complex.

\noindent\emph{Lean:} \texttt{def AlgebraicGeometry.ComplexPoint.holomorphicDeRhamSheafificationUnit}, \href{https://github.com/Paul-Lez/HodgeConjecture/blob/main/HodgeConjecture/Definitions/AlgebraicGeometry/HolomorphicDeRham.lean#L270}{\texttt{HodgeConjecture/Definitions/AlgebraicGeometry/HolomorphicDeRham.lean:270}}.
\end{definition}

\begin{definition}[AlgebraicGeometry.ComplexPoint.scalarHolomorphicDeRhamPresheaf]\label{def:AlgebraicGeometry.ComplexPoint.scalarHolomorphicDeRhamPresheaf}
\lean{AlgebraicGeometry.ComplexPoint.scalarHolomorphicDeRhamPresheaf}\leanok
\uses{def:AlgebraicGeometry.ComplexPoint.holomorphicDeRhamPresheaf}
Multiplication by a complex scalar on the presheaf of holomorphic de Rham forms.

\noindent\emph{Lean:} \texttt{def AlgebraicGeometry.ComplexPoint.scalarHolomorphicDeRhamPresheaf}, \href{https://github.com/Paul-Lez/HodgeConjecture/blob/main/HodgeConjecture/Definitions/AlgebraicGeometry/HolomorphicDeRham.lean#L329}{\texttt{HodgeConjecture/Definitions/AlgebraicGeometry/HolomorphicDeRham.lean:329}}.
\end{definition}

\begin{definition}[AlgebraicGeometry.ComplexPoint.scalarHolomorphicDeRhamComplex]\label{def:AlgebraicGeometry.ComplexPoint.scalarHolomorphicDeRhamComplex}
\lean{AlgebraicGeometry.ComplexPoint.scalarHolomorphicDeRhamComplex}\leanok
\uses{def:AlgebraicGeometry.ComplexPoint.holomorphicDeRhamComplex, def:AlgebraicGeometry.ComplexPoint.scalarHolomorphicDeRhamPresheaf}
Multiplication by a complex scalar as an endomorphism of the sheafified de Rham complex.

\noindent\emph{Lean:} \texttt{def AlgebraicGeometry.ComplexPoint.scalarHolomorphicDeRhamComplex}, \href{https://github.com/Paul-Lez/HodgeConjecture/blob/main/HodgeConjecture/Definitions/AlgebraicGeometry/HolomorphicDeRham.lean#L365}{\texttt{HodgeConjecture/Definitions/AlgebraicGeometry/HolomorphicDeRham.lean:365}}.
\end{definition}

\begin{definition}[AlgebraicGeometry.ComplexPoint.constantComplexSheaf]\label{def:AlgebraicGeometry.ComplexPoint.constantComplexSheaf}
\lean{AlgebraicGeometry.ComplexPoint.constantComplexSheaf}\leanok
\uses{def:AlgebraicGeometry.ComplexPoint, def:AlgebraicGeometry.Point.chartTopology}
The constant sheaf with value the additive group of complex numbers.

\noindent\emph{Lean:} \texttt{def AlgebraicGeometry.ComplexPoint.constantComplexSheaf}, \href{https://github.com/Paul-Lez/HodgeConjecture/blob/main/HodgeConjecture/Definitions/AlgebraicGeometry/HolomorphicDeRham.lean#L392}{\texttt{HodgeConjecture/Definitions/AlgebraicGeometry/HolomorphicDeRham.lean:392}}.
\end{definition}

\begin{definition}[AlgebraicGeometry.ComplexPoint.complexScalarPresheaf]\label{def:AlgebraicGeometry.ComplexPoint.complexScalarPresheaf}
\lean{AlgebraicGeometry.ComplexPoint.complexScalarPresheaf}\leanok
\uses{def:AlgebraicGeometry.ComplexPoint.constantComplexAddCommGrpPresheaf}
Scalar multiplication on the constant complex presheaf.

\noindent\emph{Lean:} \texttt{def AlgebraicGeometry.ComplexPoint.complexScalarPresheaf}, \href{https://github.com/Paul-Lez/HodgeConjecture/blob/main/HodgeConjecture/Definitions/AlgebraicGeometry/HolomorphicDeRham.lean#L399}{\texttt{HodgeConjecture/Definitions/AlgebraicGeometry/HolomorphicDeRham.lean:399}}.
\end{definition}

\begin{definition}[AlgebraicGeometry.ComplexPoint.complexScalarSheaf]\label{def:AlgebraicGeometry.ComplexPoint.complexScalarSheaf}
\lean{AlgebraicGeometry.ComplexPoint.complexScalarSheaf}\leanok
\uses{def:AlgebraicGeometry.ComplexPoint.complexScalarPresheaf, def:AlgebraicGeometry.ComplexPoint.constantComplexSheaf}
Scalar multiplication on the constant complex sheaf.

\noindent\emph{Lean:} \texttt{def AlgebraicGeometry.ComplexPoint.complexScalarSheaf}, \href{https://github.com/Paul-Lez/HodgeConjecture/blob/main/HodgeConjecture/Definitions/AlgebraicGeometry/HolomorphicDeRham.lean#L408}{\texttt{HodgeConjecture/Definitions/AlgebraicGeometry/HolomorphicDeRham.lean:408}}.
\end{definition}

\begin{definition}[AlgebraicGeometry.ComplexPoint.conjConstantComplexPresheaf]\label{def:AlgebraicGeometry.ComplexPoint.conjConstantComplexPresheaf}
\lean{AlgebraicGeometry.ComplexPoint.conjConstantComplexPresheaf}\leanok
\uses{def:AlgebraicGeometry.ComplexPoint.constantComplexAddCommGrpPresheaf}
Complex conjugation on the constant complex presheaf.

Conjugation is a ring automorphism of \texttt{ℂ}, so it acts on the constant complex sheaf exactly the
way a scalar does; unlike a scalar it is only additive over \texttt{ℂ}, which is what makes the induced
map on cohomology conjugate-linear rather than linear.

\noindent\emph{Lean:} \texttt{def AlgebraicGeometry.ComplexPoint.conjConstantComplexPresheaf}, \href{https://github.com/Paul-Lez/HodgeConjecture/blob/main/HodgeConjecture/Definitions/AlgebraicGeometry/HolomorphicDeRham.lean#L416}{\texttt{HodgeConjecture/Definitions/AlgebraicGeometry/HolomorphicDeRham.lean:416}}.
\end{definition}

\begin{definition}[AlgebraicGeometry.ComplexPoint.conjConstantComplexSheaf]\label{def:AlgebraicGeometry.ComplexPoint.conjConstantComplexSheaf}
\lean{AlgebraicGeometry.ComplexPoint.conjConstantComplexSheaf}\leanok
\uses{def:AlgebraicGeometry.ComplexPoint.conjConstantComplexPresheaf, def:AlgebraicGeometry.ComplexPoint.constantComplexSheaf}
Complex conjugation on the constant complex sheaf.

\noindent\emph{Lean:} \texttt{def AlgebraicGeometry.ComplexPoint.conjConstantComplexSheaf}, \href{https://github.com/Paul-Lez/HodgeConjecture/blob/main/HodgeConjecture/Definitions/AlgebraicGeometry/HolomorphicDeRham.lean#L429}{\texttt{HodgeConjecture/Definitions/AlgebraicGeometry/HolomorphicDeRham.lean:429}}.
\end{definition}

\begin{definition}[AlgebraicGeometry.ComplexPoint.constantsToHolomorphicDeRhamZeroSheaf]\label{def:AlgebraicGeometry.ComplexPoint.constantsToHolomorphicDeRhamZeroSheaf}
\lean{AlgebraicGeometry.ComplexPoint.constantsToHolomorphicDeRhamZeroSheaf}\leanok
\uses{def:AlgebraicGeometry.ComplexPoint.constantComplexSheaf, def:AlgebraicGeometry.ComplexPoint.constantsToHolomorphicDeRhamZero, def:AlgebraicGeometry.ComplexPoint.holomorphicDeRhamSheaf}
The sheafified inclusion of constants as de Rham zero-forms.

\noindent\emph{Lean:} \texttt{def AlgebraicGeometry.ComplexPoint.constantsToHolomorphicDeRhamZeroSheaf}, \href{https://github.com/Paul-Lez/HodgeConjecture/blob/main/HodgeConjecture/Definitions/AlgebraicGeometry/HolomorphicDeRham.lean#L437}{\texttt{HodgeConjecture/Definitions/AlgebraicGeometry/HolomorphicDeRham.lean:437}}.
\end{definition}

\begin{definition}[AlgebraicGeometry.ComplexPoint.constantsToHolomorphicDeRhamPresheafShortComplex]\label{def:AlgebraicGeometry.ComplexPoint.constantsToHolomorphicDeRhamPresheafShortComplex}
\lean{AlgebraicGeometry.ComplexPoint.constantsToHolomorphicDeRhamPresheafShortComplex}\leanok
\uses{def:AlgebraicGeometry.ComplexPoint.constantsToHolomorphicDeRhamZero, def:AlgebraicGeometry.ComplexPoint.holomorphicDeRhamDifferential}
The augmented holomorphic de Rham short complex before sheafification.

\noindent\emph{Lean:} \texttt{def AlgebraicGeometry.ComplexPoint.constantsToHolomorphicDeRhamPresheafShortComplex}, \href{https://github.com/Paul-Lez/HodgeConjecture/blob/main/HodgeConjecture/Definitions/AlgebraicGeometry/HolomorphicDeRham.lean#L456}{\texttt{HodgeConjecture/Definitions/AlgebraicGeometry/HolomorphicDeRham.lean:456}}.
\end{definition}

\begin{definition}[AlgebraicGeometry.ComplexPoint.constantsToHolomorphicDeRhamSheafShortComplex]\label{def:AlgebraicGeometry.ComplexPoint.constantsToHolomorphicDeRhamSheafShortComplex}
\lean{AlgebraicGeometry.ComplexPoint.constantsToHolomorphicDeRhamSheafShortComplex}\leanok
\uses{def:AlgebraicGeometry.ComplexPoint.constantsToHolomorphicDeRhamZeroSheaf, def:AlgebraicGeometry.ComplexPoint.holomorphicDeRhamSheafDifferential}
The augmented holomorphic de Rham short complex after sheafification.

\noindent\emph{Lean:} \texttt{def AlgebraicGeometry.ComplexPoint.constantsToHolomorphicDeRhamSheafShortComplex}, \href{https://github.com/Paul-Lez/HodgeConjecture/blob/main/HodgeConjecture/Definitions/AlgebraicGeometry/HolomorphicDeRham.lean#L465}{\texttt{HodgeConjecture/Definitions/AlgebraicGeometry/HolomorphicDeRham.lean:465}}.
\end{definition}

\begin{definition}[AlgebraicGeometry.ComplexPoint.constantsToHolomorphicDeRhamShortComplexSheafificationUnit]\label{def:AlgebraicGeometry.ComplexPoint.constantsToHolomorphicDeRhamShortComplexSheafificationUnit}
\lean{AlgebraicGeometry.ComplexPoint.constantsToHolomorphicDeRhamShortComplexSheafificationUnit}\leanok
\uses{def:AlgebraicGeometry.ComplexPoint.constantsToHolomorphicDeRhamPresheafShortComplex, def:AlgebraicGeometry.ComplexPoint.constantsToHolomorphicDeRhamSheafShortComplex}
The sheafification unit between the augmented presheaf and sheaf short complexes.

\noindent\emph{Lean:} \texttt{def AlgebraicGeometry.ComplexPoint.constantsToHolomorphicDeRhamShortComplexSheafificationUnit}, \href{https://github.com/Paul-Lez/HodgeConjecture/blob/main/HodgeConjecture/Definitions/AlgebraicGeometry/HolomorphicDeRham.lean#L475}{\texttt{HodgeConjecture/Definitions/AlgebraicGeometry/HolomorphicDeRham.lean:475}}.
\end{definition}

\begin{definition}[AlgebraicGeometry.ComplexPoint.constantsToHolomorphicDeRhamComplex]\label{def:AlgebraicGeometry.ComplexPoint.constantsToHolomorphicDeRhamComplex}
\lean{AlgebraicGeometry.ComplexPoint.constantsToHolomorphicDeRhamComplex}\leanok
\uses{def:AlgebraicGeometry.ComplexPoint.constantsToHolomorphicDeRhamZeroSheaf, def:AlgebraicGeometry.ComplexPoint.holomorphicDeRhamComplex}
The comparison from the constant sheaf complex to the holomorphic de Rham complex.

\noindent\emph{Lean:} \texttt{def AlgebraicGeometry.ComplexPoint.constantsToHolomorphicDeRhamComplex}, \href{https://github.com/Paul-Lez/HodgeConjecture/blob/main/HodgeConjecture/Definitions/AlgebraicGeometry/HolomorphicDeRham.lean#L559}{\texttt{HodgeConjecture/Definitions/AlgebraicGeometry/HolomorphicDeRham.lean:559}}.
\end{definition}

\begin{definition}[AlgebraicGeometry.ComplexPoint.complexScalarComplex]\label{def:AlgebraicGeometry.ComplexPoint.complexScalarComplex}
\lean{AlgebraicGeometry.ComplexPoint.complexScalarComplex}\leanok
\uses{def:AlgebraicGeometry.ComplexPoint.complexScalarSheaf}
Scalar multiplication on the constant complex-valued complex concentrated in degree zero.

\noindent\emph{Lean:} \texttt{def AlgebraicGeometry.ComplexPoint.complexScalarComplex}, \href{https://github.com/Paul-Lez/HodgeConjecture/blob/main/HodgeConjecture/Definitions/AlgebraicGeometry/HolomorphicDeRham.lean#L601}{\texttt{HodgeConjecture/Definitions/AlgebraicGeometry/HolomorphicDeRham.lean:601}}.
\end{definition}

\begin{definition}[AlgebraicGeometry.ComplexPoint.constantComplexSheafComplexInt]\label{def:AlgebraicGeometry.ComplexPoint.constantComplexSheafComplexInt}
\lean{AlgebraicGeometry.ComplexPoint.constantComplexSheafComplexInt}\leanok
\uses{def:AlgebraicGeometry.ComplexPoint.constantComplexSheaf}
The constant sheaf complex, extended by zero from natural to integer degrees.

\noindent\emph{Lean:} \texttt{def AlgebraicGeometry.ComplexPoint.constantComplexSheafComplexInt}, \href{https://github.com/Paul-Lez/HodgeConjecture/blob/main/HodgeConjecture/Definitions/AlgebraicGeometry/HolomorphicDeRham.lean#L612}{\texttt{HodgeConjecture/Definitions/AlgebraicGeometry/HolomorphicDeRham.lean:612}}.
\end{definition}

\begin{definition}[AlgebraicGeometry.ComplexPoint.complexScalarComplexInt]\label{def:AlgebraicGeometry.ComplexPoint.complexScalarComplexInt}
\lean{AlgebraicGeometry.ComplexPoint.complexScalarComplexInt}\leanok
\uses{def:AlgebraicGeometry.ComplexPoint.complexScalarComplex, def:AlgebraicGeometry.ComplexPoint.constantComplexSheafComplexInt}
Scalar multiplication on the integer-indexed constant complex-valued complex.

\noindent\emph{Lean:} \texttt{def AlgebraicGeometry.ComplexPoint.complexScalarComplexInt}, \href{https://github.com/Paul-Lez/HodgeConjecture/blob/main/HodgeConjecture/Definitions/AlgebraicGeometry/HolomorphicDeRham.lean#L621}{\texttt{HodgeConjecture/Definitions/AlgebraicGeometry/HolomorphicDeRham.lean:621}}.
\end{definition}

\begin{definition}[AlgebraicGeometry.ComplexPoint.conjConstantComplexComplex]\label{def:AlgebraicGeometry.ComplexPoint.conjConstantComplexComplex}
\lean{AlgebraicGeometry.ComplexPoint.conjConstantComplexComplex}\leanok
\uses{def:AlgebraicGeometry.ComplexPoint.conjConstantComplexSheaf}
Complex conjugation on the constant complex-valued complex concentrated in degree zero.

\noindent\emph{Lean:} \texttt{def AlgebraicGeometry.ComplexPoint.conjConstantComplexComplex}, \href{https://github.com/Paul-Lez/HodgeConjecture/blob/main/HodgeConjecture/Definitions/AlgebraicGeometry/HolomorphicDeRham.lean#L628}{\texttt{HodgeConjecture/Definitions/AlgebraicGeometry/HolomorphicDeRham.lean:628}}.
\end{definition}

\begin{definition}[AlgebraicGeometry.ComplexPoint.conjConstantComplexSheafComplexInt]\label{def:AlgebraicGeometry.ComplexPoint.conjConstantComplexSheafComplexInt}
\lean{AlgebraicGeometry.ComplexPoint.conjConstantComplexSheafComplexInt}\leanok
\uses{def:AlgebraicGeometry.ComplexPoint.conjConstantComplexComplex, def:AlgebraicGeometry.ComplexPoint.constantComplexSheafComplexInt}
Complex conjugation on the integer-indexed constant complex-valued complex.

The holomorphic de Rham complex carries no such map: conjugation is not \texttt{ℂ}-linear, so it exists
only on the constant-sheaf side of the comparison.

\noindent\emph{Lean:} \texttt{def AlgebraicGeometry.ComplexPoint.conjConstantComplexSheafComplexInt}, \href{https://github.com/Paul-Lez/HodgeConjecture/blob/main/HodgeConjecture/Definitions/AlgebraicGeometry/HolomorphicDeRham.lean#L638}{\texttt{HodgeConjecture/Definitions/AlgebraicGeometry/HolomorphicDeRham.lean:638}}.
\end{definition}

\begin{definition}[AlgebraicGeometry.ComplexPoint.holomorphicDeRhamComplexInt]\label{def:AlgebraicGeometry.ComplexPoint.holomorphicDeRhamComplexInt}
\lean{AlgebraicGeometry.ComplexPoint.holomorphicDeRhamComplexInt}\leanok
\uses{def:AlgebraicGeometry.ComplexPoint.holomorphicDeRhamComplex, def:TopologicalSpace.dim}
The holomorphic de Rham complex, extended by zero to negative degrees.

\noindent\emph{Lean:} \texttt{def AlgebraicGeometry.ComplexPoint.holomorphicDeRhamComplexInt}, \href{https://github.com/Paul-Lez/HodgeConjecture/blob/main/HodgeConjecture/Definitions/AlgebraicGeometry/HolomorphicDeRham.lean#L647}{\texttt{HodgeConjecture/Definitions/AlgebraicGeometry/HolomorphicDeRham.lean:647}}.
\end{definition}

\begin{definition}[AlgebraicGeometry.ComplexPoint.scalarHolomorphicDeRhamComplexInt]\label{def:AlgebraicGeometry.ComplexPoint.scalarHolomorphicDeRhamComplexInt}
\lean{AlgebraicGeometry.ComplexPoint.scalarHolomorphicDeRhamComplexInt}\leanok
\uses{def:AlgebraicGeometry.ComplexPoint.holomorphicDeRhamComplexInt, def:AlgebraicGeometry.ComplexPoint.scalarHolomorphicDeRhamComplex}
Multiplication by a complex scalar on the integer-indexed holomorphic de Rham complex.

\noindent\emph{Lean:} \texttt{def AlgebraicGeometry.ComplexPoint.scalarHolomorphicDeRhamComplexInt}, \href{https://github.com/Paul-Lez/HodgeConjecture/blob/main/HodgeConjecture/Definitions/AlgebraicGeometry/HolomorphicDeRham.lean#L681}{\texttt{HodgeConjecture/Definitions/AlgebraicGeometry/HolomorphicDeRham.lean:681}}.
\end{definition}

\begin{definition}[AlgebraicGeometry.ComplexPoint.constantsToHolomorphicDeRhamComplexInt]\label{def:AlgebraicGeometry.ComplexPoint.constantsToHolomorphicDeRhamComplexInt}
\lean{AlgebraicGeometry.ComplexPoint.constantsToHolomorphicDeRhamComplexInt}\leanok
\uses{def:AlgebraicGeometry.ComplexPoint.constantComplexSheafComplexInt, def:AlgebraicGeometry.ComplexPoint.constantsToHolomorphicDeRhamComplex, def:AlgebraicGeometry.ComplexPoint.holomorphicDeRhamComplexInt}
The constant-to-de Rham comparison on integer-indexed complexes.

\noindent\emph{Lean:} \texttt{def AlgebraicGeometry.ComplexPoint.constantsToHolomorphicDeRhamComplexInt}, \href{https://github.com/Paul-Lez/HodgeConjecture/blob/main/HodgeConjecture/Definitions/AlgebraicGeometry/HolomorphicDeRham.lean#L689}{\texttt{HodgeConjecture/Definitions/AlgebraicGeometry/HolomorphicDeRham.lean:689}}.
\end{definition}

\section{The Hodge filtration}\label{sec:HodgeConjecture.Definitions.AlgebraicGeometry.HodgeFiltration}
\noindent\emph{File:} \href{https://github.com/Paul-Lez/HodgeConjecture/blob/main/HodgeConjecture/Definitions/AlgebraicGeometry/HodgeFiltration.lean}{\texttt{HodgeConjecture/Definitions/AlgebraicGeometry/HodgeFiltration.lean}}.

This file defines rational sheaf cohomology and holomorphic de Rham hypercohomology on the
analytic complex-point space of a smooth complex scheme. Hypercohomology is expressed with
Mathlib's small shifted morphisms in the localization at quasi-isomorphisms. This avoids exposing
a noncanonical choice of derived category in the public types.

The stupid truncation of the holomorphic de Rham complex in form degrees at least \texttt{p} maps into
the full complex. Its image on hypercohomology is the Hodge filtration \texttt{F\^{}p}. A rational Hodge
class is then a rational class whose de Rham image belongs to \texttt{F\^{}p}.

\begin{definition}[AlgebraicGeometry.ComplexPoint.hodgeFiltrationTopology]\label{def:AlgebraicGeometry.ComplexPoint.hodgeFiltrationTopology}
\lean{AlgebraicGeometry.ComplexPoint.hodgeFiltrationTopology}\leanok
\uses{def:AlgebraicGeometry.ComplexPoint, def:AlgebraicGeometry.Point.chartTopology}
\noindent\emph{Lean:} \texttt{def AlgebraicGeometry.ComplexPoint.hodgeFiltrationTopology}, \href{https://github.com/Paul-Lez/HodgeConjecture/blob/main/HodgeConjecture/Definitions/AlgebraicGeometry/HodgeFiltration.lean#L53}{\texttt{HodgeConjecture/Definitions/AlgebraicGeometry/HodgeFiltration.lean:53}}.
\end{definition}

\begin{definition}[AlgebraicGeometry.ComplexPoint.AnalyticAdditiveSheaf]\label{def:AlgebraicGeometry.ComplexPoint.AnalyticAdditiveSheaf}
\lean{AlgebraicGeometry.ComplexPoint.AnalyticAdditiveSheaf}\leanok
\uses{def:AlgebraicGeometry.ComplexPoint.hodgeFiltrationTopology}
Sheaves of additive groups on the analytic complex-point space.

\noindent\emph{Lean:} \texttt{abbrev AlgebraicGeometry.ComplexPoint.AnalyticAdditiveSheaf}, \href{https://github.com/Paul-Lez/HodgeConjecture/blob/main/HodgeConjecture/Definitions/AlgebraicGeometry/HodgeFiltration.lean#L56}{\texttt{HodgeConjecture/Definitions/AlgebraicGeometry/HodgeFiltration.lean:56}}.
\end{definition}

\begin{definition}[AlgebraicGeometry.ComplexPoint.analyticHasDerivedCategory]\label{def:AlgebraicGeometry.ComplexPoint.analyticHasDerivedCategory}
\lean{AlgebraicGeometry.ComplexPoint.analyticHasDerivedCategory}\leanok
\uses{def:AlgebraicGeometry.ComplexPoint.AnalyticAdditiveSheaf}
\noindent\emph{Lean:} \texttt{def AlgebraicGeometry.ComplexPoint.analyticHasDerivedCategory}, \href{https://github.com/Paul-Lez/HodgeConjecture/blob/main/HodgeConjecture/Definitions/AlgebraicGeometry/HodgeFiltration.lean#L60}{\texttt{HodgeConjecture/Definitions/AlgebraicGeometry/HodgeFiltration.lean:60}}.
\end{definition}

\begin{definition}[AlgebraicGeometry.ComplexPoint.constantFieldSheaf]\label{def:AlgebraicGeometry.ComplexPoint.constantFieldSheaf}
\lean{AlgebraicGeometry.ComplexPoint.constantFieldSheaf}\leanok
\uses{def:AlgebraicGeometry.ComplexPoint.AnalyticAdditiveSheaf}
The constant rational sheaf on the analytic complex-point space.

\noindent\emph{Lean:} \texttt{abbrev AlgebraicGeometry.ComplexPoint.constantFieldSheaf}, \href{https://github.com/Paul-Lez/HodgeConjecture/blob/main/HodgeConjecture/Definitions/AlgebraicGeometry/HodgeFiltration.lean#L64}{\texttt{HodgeConjecture/Definitions/AlgebraicGeometry/HodgeFiltration.lean:64}}.
\end{definition}

\begin{definition}[AlgebraicGeometry.ComplexPoint.fieldToComplexConstantSheaf]\label{def:AlgebraicGeometry.ComplexPoint.fieldToComplexConstantSheaf}
\lean{AlgebraicGeometry.ComplexPoint.fieldToComplexConstantSheaf}\leanok
\uses{def:AlgebraicGeometry.ComplexPoint.constantComplexSheaf, def:AlgebraicGeometry.ComplexPoint.constantFieldSheaf}
The inclusion of the rational constant sheaf into the complex constant sheaf.

\noindent\emph{Lean:} \texttt{abbrev AlgebraicGeometry.ComplexPoint.fieldToComplexConstantSheaf}, \href{https://github.com/Paul-Lez/HodgeConjecture/blob/main/HodgeConjecture/Definitions/AlgebraicGeometry/HodgeFiltration.lean#L69}{\texttt{HodgeConjecture/Definitions/AlgebraicGeometry/HodgeFiltration.lean:69}}.
\end{definition}

\begin{definition}[AlgebraicGeometry.ComplexPoint.constantFieldSheafComplexInt]\label{def:AlgebraicGeometry.ComplexPoint.constantFieldSheafComplexInt}
\lean{AlgebraicGeometry.ComplexPoint.constantFieldSheafComplexInt}\leanok
\uses{def:AlgebraicGeometry.ComplexPoint.constantFieldSheaf}
The constant rational sheaf complex, extended by zero to integer degrees.

\noindent\emph{Lean:} \texttt{def AlgebraicGeometry.ComplexPoint.constantFieldSheafComplexInt}, \href{https://github.com/Paul-Lez/HodgeConjecture/blob/main/HodgeConjecture/Definitions/AlgebraicGeometry/HodgeFiltration.lean#L83}{\texttt{HodgeConjecture/Definitions/AlgebraicGeometry/HodgeFiltration.lean:83}}.
\end{definition}

\begin{definition}[AlgebraicGeometry.ComplexPoint.fieldToComplexConstantSheafComplexInt]\label{def:AlgebraicGeometry.ComplexPoint.fieldToComplexConstantSheafComplexInt}
\lean{AlgebraicGeometry.ComplexPoint.fieldToComplexConstantSheafComplexInt}\leanok
\uses{def:AlgebraicGeometry.ComplexPoint.constantComplexSheafComplexInt, def:AlgebraicGeometry.ComplexPoint.constantFieldSheafComplexInt, def:AlgebraicGeometry.ComplexPoint.fieldToComplexConstantSheaf}
Extension of rational constants to complex constants as a map of integer complexes.

\noindent\emph{Lean:} \texttt{def AlgebraicGeometry.ComplexPoint.fieldToComplexConstantSheafComplexInt}, \href{https://github.com/Paul-Lez/HodgeConjecture/blob/main/HodgeConjecture/Definitions/AlgebraicGeometry/HodgeFiltration.lean#L94}{\texttt{HodgeConjecture/Definitions/AlgebraicGeometry/HodgeFiltration.lean:94}}.
\end{definition}

\begin{definition}[AlgebraicGeometry.ComplexPoint.fieldToHolomorphicDeRhamComplexInt]\label{def:AlgebraicGeometry.ComplexPoint.fieldToHolomorphicDeRhamComplexInt}
\lean{AlgebraicGeometry.ComplexPoint.fieldToHolomorphicDeRhamComplexInt}\leanok
\uses{def:AlgebraicGeometry.ComplexPoint.constantsToHolomorphicDeRhamComplexInt, def:AlgebraicGeometry.ComplexPoint.fieldToComplexConstantSheafComplexInt}
Rational constants mapped canonically into the holomorphic de Rham complex.

\noindent\emph{Lean:} \texttt{def AlgebraicGeometry.ComplexPoint.fieldToHolomorphicDeRhamComplexInt}, \href{https://github.com/Paul-Lez/HodgeConjecture/blob/main/HodgeConjecture/Definitions/AlgebraicGeometry/HodgeFiltration.lean#L110}{\texttt{HodgeConjecture/Definitions/AlgebraicGeometry/HodgeFiltration.lean:110}}.
\end{definition}

\begin{definition}[AlgebraicGeometry.ComplexPoint.constantIntegerSheaf]\label{def:AlgebraicGeometry.ComplexPoint.constantIntegerSheaf}
\lean{AlgebraicGeometry.ComplexPoint.constantIntegerSheaf}\leanok
\uses{def:AlgebraicGeometry.ComplexPoint.AnalyticAdditiveSheaf}
The constant integer sheaf on the analytic complex-point space.

\noindent\emph{Lean:} \texttt{def AlgebraicGeometry.ComplexPoint.constantIntegerSheaf}, \href{https://github.com/Paul-Lez/HodgeConjecture/blob/main/HodgeConjecture/Definitions/AlgebraicGeometry/HodgeFiltration.lean#L117}{\texttt{HodgeConjecture/Definitions/AlgebraicGeometry/HodgeFiltration.lean:117}}.
\end{definition}

\begin{definition}[AlgebraicGeometry.ComplexPoint.constantIntegerSheafComplexInt]\label{def:AlgebraicGeometry.ComplexPoint.constantIntegerSheafComplexInt}
\lean{AlgebraicGeometry.ComplexPoint.constantIntegerSheafComplexInt}\leanok
\uses{def:AlgebraicGeometry.ComplexPoint.constantIntegerSheaf}
The constant integer sheaf complex, extended by zero to integer degrees.

\noindent\emph{Lean:} \texttt{def AlgebraicGeometry.ComplexPoint.constantIntegerSheafComplexInt}, \href{https://github.com/Paul-Lez/HodgeConjecture/blob/main/HodgeConjecture/Definitions/AlgebraicGeometry/HodgeFiltration.lean#L122}{\texttt{HodgeConjecture/Definitions/AlgebraicGeometry/HodgeFiltration.lean:122}}.
\end{definition}

\begin{definition}[AlgebraicGeometry.ComplexPoint.fieldScalarSheaf]\label{def:AlgebraicGeometry.ComplexPoint.fieldScalarSheaf}
\lean{AlgebraicGeometry.ComplexPoint.fieldScalarSheaf}\leanok
\uses{def:AlgebraicGeometry.ComplexPoint.constantFieldSheaf, def:fieldScalarAddHom}
Scalar multiplication on the rational constant sheaf.

\noindent\emph{Lean:} \texttt{abbrev AlgebraicGeometry.ComplexPoint.fieldScalarSheaf}, \href{https://github.com/Paul-Lez/HodgeConjecture/blob/main/HodgeConjecture/Definitions/AlgebraicGeometry/HodgeFiltration.lean#L183}{\texttt{HodgeConjecture/Definitions/AlgebraicGeometry/HodgeFiltration.lean:183}}.
\end{definition}

\begin{definition}[AlgebraicGeometry.ComplexPoint.fieldScalarComplex]\label{def:AlgebraicGeometry.ComplexPoint.fieldScalarComplex}
\lean{AlgebraicGeometry.ComplexPoint.fieldScalarComplex}\leanok
\uses{def:AlgebraicGeometry.ComplexPoint.constantFieldSheafComplexInt, def:AlgebraicGeometry.ComplexPoint.fieldScalarSheaf}
Scalar multiplication on the rational constant sheaf complex.

\noindent\emph{Lean:} \texttt{def AlgebraicGeometry.ComplexPoint.fieldScalarComplex}, \href{https://github.com/Paul-Lez/HodgeConjecture/blob/main/HodgeConjecture/Definitions/AlgebraicGeometry/HodgeFiltration.lean#L292}{\texttt{HodgeConjecture/Definitions/AlgebraicGeometry/HodgeFiltration.lean:292}}.
\end{definition}

\begin{definition}[AlgebraicGeometry.ComplexPoint.analyticQuasiIsomorphisms]\label{def:AlgebraicGeometry.ComplexPoint.analyticQuasiIsomorphisms}
\lean{AlgebraicGeometry.ComplexPoint.analyticQuasiIsomorphisms}\leanok
\uses{def:AlgebraicGeometry.ComplexPoint.AnalyticAdditiveSheaf}
Quasi-isomorphisms of analytic sheaf complexes.

\noindent\emph{Lean:} \texttt{abbrev AlgebraicGeometry.ComplexPoint.analyticQuasiIsomorphisms}, \href{https://github.com/Paul-Lez/HodgeConjecture/blob/main/HodgeConjecture/Definitions/AlgebraicGeometry/HodgeFiltration.lean#L354}{\texttt{HodgeConjecture/Definitions/AlgebraicGeometry/HodgeFiltration.lean:354}}.
\end{definition}

\begin{definition}[AlgebraicGeometry.ComplexPoint.Hypercohomology]\label{def:AlgebraicGeometry.ComplexPoint.Hypercohomology}
\lean{AlgebraicGeometry.ComplexPoint.Hypercohomology}\leanok
\uses{def:AlgebraicGeometry.ComplexPoint.analyticHasDerivedCategory, def:AlgebraicGeometry.ComplexPoint.analyticQuasiIsomorphisms, def:AlgebraicGeometry.ComplexPoint.constantIntegerSheafComplexInt}
Hypercohomology of an analytic sheaf complex in integer degree \texttt{n}.

\noindent\emph{Lean:} \texttt{def AlgebraicGeometry.ComplexPoint.Hypercohomology}, \href{https://github.com/Paul-Lez/HodgeConjecture/blob/main/HodgeConjecture/Definitions/AlgebraicGeometry/HodgeFiltration.lean#L366}{\texttt{HodgeConjecture/Definitions/AlgebraicGeometry/HodgeFiltration.lean:366}}.
\end{definition}

\begin{definition}[AlgebraicGeometry.ComplexPoint.ComplexConstantCohomology]\label{def:AlgebraicGeometry.ComplexPoint.ComplexConstantCohomology}
\lean{AlgebraicGeometry.ComplexPoint.ComplexConstantCohomology}\leanok
\uses{def:AlgebraicGeometry.ComplexPoint.Hypercohomology, def:AlgebraicGeometry.ComplexPoint.constantComplexSheafComplexInt}
Complex constant-sheaf cohomology in integer degree \texttt{n}.

\noindent\emph{Lean:} \texttt{abbrev AlgebraicGeometry.ComplexPoint.ComplexConstantCohomology}, \href{https://github.com/Paul-Lez/HodgeConjecture/blob/main/HodgeConjecture/Definitions/AlgebraicGeometry/HodgeFiltration.lean#L381}{\texttt{HodgeConjecture/Definitions/AlgebraicGeometry/HodgeFiltration.lean:381}}.
\end{definition}

\begin{definition}[AlgebraicGeometry.ComplexPoint.DeRhamHypercohomology]\label{def:AlgebraicGeometry.ComplexPoint.DeRhamHypercohomology}
\lean{AlgebraicGeometry.ComplexPoint.DeRhamHypercohomology}\leanok
\uses{def:AlgebraicGeometry.ComplexPoint.Hypercohomology, def:AlgebraicGeometry.ComplexPoint.holomorphicDeRhamComplexInt}
Hypercohomology of the holomorphic de Rham complex in integer degree \texttt{n}.

\noindent\emph{Lean:} \texttt{abbrev AlgebraicGeometry.ComplexPoint.DeRhamHypercohomology}, \href{https://github.com/Paul-Lez/HodgeConjecture/blob/main/HodgeConjecture/Definitions/AlgebraicGeometry/HodgeFiltration.lean#L441}{\texttt{HodgeConjecture/Definitions/AlgebraicGeometry/HodgeFiltration.lean:441}}.
\end{definition}

\begin{definition}[AlgebraicGeometry.ComplexPoint.complexConstantCohomologyDeRhamEquiv]\label{def:AlgebraicGeometry.ComplexPoint.complexConstantCohomologyDeRhamEquiv}
\lean{AlgebraicGeometry.ComplexPoint.complexConstantCohomologyDeRhamEquiv}\leanok
\uses{def:AlgebraicGeometry.ComplexPoint.ComplexConstantCohomology, def:AlgebraicGeometry.ComplexPoint.DeRhamHypercohomology, def:AlgebraicGeometry.ComplexPoint.constantsToHolomorphicDeRhamComplexInt}
A proved constant-to-holomorphic-de Rham quasi-isomorphism induces the corresponding
equivalence on hypercohomology.

\noindent\emph{Lean:} \texttt{def AlgebraicGeometry.ComplexPoint.complexConstantCohomologyDeRhamEquiv}, \href{https://github.com/Paul-Lez/HodgeConjecture/blob/main/HodgeConjecture/Definitions/AlgebraicGeometry/HodgeFiltration.lean#L445}{\texttt{HodgeConjecture/Definitions/AlgebraicGeometry/HodgeFiltration.lean:445}}.
\end{definition}

\begin{definition}[AlgebraicGeometry.ComplexPoint.hypercohomologyMap]\label{def:AlgebraicGeometry.ComplexPoint.hypercohomologyMap}
\lean{AlgebraicGeometry.ComplexPoint.hypercohomologyMap}\leanok
\uses{def:AlgebraicGeometry.ComplexPoint.Hypercohomology}
Postcomposition on hypercohomology by a map of complexes.

\noindent\emph{Lean:} \texttt{def AlgebraicGeometry.ComplexPoint.hypercohomologyMap}, \href{https://github.com/Paul-Lez/HodgeConjecture/blob/main/HodgeConjecture/Definitions/AlgebraicGeometry/HodgeFiltration.lean#L455}{\texttt{HodgeConjecture/Definitions/AlgebraicGeometry/HodgeFiltration.lean:455}}.
\end{definition}

\begin{definition}[AlgebraicGeometry.ComplexPoint.hypercohomologyModule]\label{def:AlgebraicGeometry.ComplexPoint.hypercohomologyModule}
\lean{AlgebraicGeometry.ComplexPoint.hypercohomologyModule}\leanok
\uses{def:AlgebraicGeometry.ComplexPoint.hypercohomologyMap}
Scalars acting on a coefficient complex by an additive, unital and antimultiplicative family
of endomorphisms act on its hypercohomology.

\noindent\emph{Lean:} \texttt{abbrev AlgebraicGeometry.ComplexPoint.hypercohomologyModule}, \href{https://github.com/Paul-Lez/HodgeConjecture/blob/main/HodgeConjecture/Definitions/AlgebraicGeometry/HodgeFiltration.lean#L543}{\texttt{HodgeConjecture/Definitions/AlgebraicGeometry/HodgeFiltration.lean:543}}.
\end{definition}

\begin{definition}[AlgebraicGeometry.ComplexPoint.fieldCohomologySMul]\label{def:AlgebraicGeometry.ComplexPoint.fieldCohomologySMul}
\lean{AlgebraicGeometry.ComplexPoint.fieldCohomologySMul}\leanok
\uses{def:AlgebraicGeometry.ComplexPoint.fieldScalarComplex, def:AlgebraicGeometry.ComplexPoint.hypercohomologyMap}
The rational action on constant-sheaf cohomology, induced by scalar multiplication on the
coefficient sheaf.

\noindent\emph{Lean:} \texttt{def AlgebraicGeometry.ComplexPoint.fieldCohomologySMul}, \href{https://github.com/Paul-Lez/HodgeConjecture/blob/main/HodgeConjecture/Definitions/AlgebraicGeometry/HodgeFiltration.lean#L575}{\texttt{HodgeConjecture/Definitions/AlgebraicGeometry/HodgeFiltration.lean:575}}.
\end{definition}

\begin{definition}[AlgebraicGeometry.ComplexPoint.deRhamComplexSMul]\label{def:AlgebraicGeometry.ComplexPoint.deRhamComplexSMul}
\lean{AlgebraicGeometry.ComplexPoint.deRhamComplexSMul}\leanok
\uses{def:AlgebraicGeometry.ComplexPoint.DeRhamHypercohomology, def:AlgebraicGeometry.ComplexPoint.hypercohomologyMap, def:AlgebraicGeometry.ComplexPoint.scalarHolomorphicDeRhamComplexInt}
The complex action on holomorphic de Rham hypercohomology, induced by scalar multiplication
on the holomorphic de Rham complex.

\noindent\emph{Lean:} \texttt{def AlgebraicGeometry.ComplexPoint.deRhamComplexSMul}, \href{https://github.com/Paul-Lez/HodgeConjecture/blob/main/HodgeConjecture/Definitions/AlgebraicGeometry/HodgeFiltration.lean#L597}{\texttt{HodgeConjecture/Definitions/AlgebraicGeometry/HodgeFiltration.lean:597}}.
\end{definition}

\begin{definition}[AlgebraicGeometry.ComplexPoint.fieldToDeRhamCohomology]\label{def:AlgebraicGeometry.ComplexPoint.fieldToDeRhamCohomology}
\lean{AlgebraicGeometry.ComplexPoint.fieldToDeRhamCohomology}\leanok
\uses{def:AlgebraicGeometry.ComplexPoint.DeRhamHypercohomology, def:AlgebraicGeometry.ComplexPoint.fieldToHolomorphicDeRhamComplexInt, def:AlgebraicGeometry.ComplexPoint.hypercohomologyMap}
The derived comparison from rational cohomology to holomorphic de Rham hypercohomology.

\noindent\emph{Lean:} \texttt{def AlgebraicGeometry.ComplexPoint.fieldToDeRhamCohomology}, \href{https://github.com/Paul-Lez/HodgeConjecture/blob/main/HodgeConjecture/Definitions/AlgebraicGeometry/HodgeFiltration.lean#L667}{\texttt{HodgeConjecture/Definitions/AlgebraicGeometry/HodgeFiltration.lean:667}}.
\end{definition}

\begin{definition}[AlgebraicGeometry.ComplexPoint.fieldToDeRhamCohomologyLinear]\label{def:AlgebraicGeometry.ComplexPoint.fieldToDeRhamCohomologyLinear}
\lean{AlgebraicGeometry.ComplexPoint.fieldToDeRhamCohomologyLinear}\leanok
\uses{def:AlgebraicGeometry.ComplexPoint.complexScalarComplexInt, def:AlgebraicGeometry.ComplexPoint.deRhamComplexSMul, def:AlgebraicGeometry.ComplexPoint.fieldCohomologySMul, def:AlgebraicGeometry.ComplexPoint.fieldToDeRhamCohomology, def:AlgebraicGeometry.ComplexPoint.hypercohomologyModule}
The rational-to-de Rham comparison as a rational-linear map.

\noindent\emph{Lean:} \texttt{def AlgebraicGeometry.ComplexPoint.fieldToDeRhamCohomologyLinear}, \href{https://github.com/Paul-Lez/HodgeConjecture/blob/main/HodgeConjecture/Definitions/AlgebraicGeometry/HodgeFiltration.lean#L684}{\texttt{HodgeConjecture/Definitions/AlgebraicGeometry/HodgeFiltration.lean:684}}.
\end{definition}

\begin{definition}[AlgebraicGeometry.ComplexPoint.hodgeFilteredDeRhamComplex]\label{def:AlgebraicGeometry.ComplexPoint.hodgeFilteredDeRhamComplex}
\lean{AlgebraicGeometry.ComplexPoint.hodgeFilteredDeRhamComplex}\leanok
\uses{def:AlgebraicGeometry.ComplexPoint.AnalyticAdditiveSheaf, def:AlgebraicGeometry.ComplexPoint.holomorphicDeRhamComplexInt}
The de Rham complex with only form degrees at least \texttt{p} retained.

\noindent\emph{Lean:} \texttt{def AlgebraicGeometry.ComplexPoint.hodgeFilteredDeRhamComplex}, \href{https://github.com/Paul-Lez/HodgeConjecture/blob/main/HodgeConjecture/Definitions/AlgebraicGeometry/HodgeFiltration.lean#L693}{\texttt{HodgeConjecture/Definitions/AlgebraicGeometry/HodgeFiltration.lean:693}}.
\end{definition}

\begin{definition}[AlgebraicGeometry.ComplexPoint.hodgeFilteredDeRhamInclusion]\label{def:AlgebraicGeometry.ComplexPoint.hodgeFilteredDeRhamInclusion}
\lean{AlgebraicGeometry.ComplexPoint.hodgeFilteredDeRhamInclusion}\leanok
\uses{def:AlgebraicGeometry.ComplexPoint.hodgeFilteredDeRhamComplex, def:HomologicalComplex.stupidTruncInclusion}
Inclusion of the degree-at-least-\texttt{p} de Rham complex into the full complex.

\noindent\emph{Lean:} \texttt{def AlgebraicGeometry.ComplexPoint.hodgeFilteredDeRhamInclusion}, \href{https://github.com/Paul-Lez/HodgeConjecture/blob/main/HodgeConjecture/Definitions/AlgebraicGeometry/HodgeFiltration.lean#L699}{\texttt{HodgeConjecture/Definitions/AlgebraicGeometry/HodgeFiltration.lean:699}}.
\end{definition}

\begin{definition}[AlgebraicGeometry.ComplexPoint.hodgeFilteredDeRhamComplexScalar]\label{def:AlgebraicGeometry.ComplexPoint.hodgeFilteredDeRhamComplexScalar}
\lean{AlgebraicGeometry.ComplexPoint.hodgeFilteredDeRhamComplexScalar}\leanok
\uses{def:AlgebraicGeometry.ComplexPoint.hodgeFilteredDeRhamComplex, def:AlgebraicGeometry.ComplexPoint.scalarHolomorphicDeRhamComplexInt}
Complex scalar multiplication on the filtered de Rham complex.

\noindent\emph{Lean:} \texttt{def AlgebraicGeometry.ComplexPoint.hodgeFilteredDeRhamComplexScalar}, \href{https://github.com/Paul-Lez/HodgeConjecture/blob/main/HodgeConjecture/Definitions/AlgebraicGeometry/HodgeFiltration.lean#L706}{\texttt{HodgeConjecture/Definitions/AlgebraicGeometry/HodgeFiltration.lean:706}}.
\end{definition}

\begin{definition}[AlgebraicGeometry.ComplexPoint.FilteredDeRhamHypercohomology]\label{def:AlgebraicGeometry.ComplexPoint.FilteredDeRhamHypercohomology}
\lean{AlgebraicGeometry.ComplexPoint.FilteredDeRhamHypercohomology}\leanok
\uses{def:AlgebraicGeometry.ComplexPoint.Hypercohomology, def:AlgebraicGeometry.ComplexPoint.hodgeFilteredDeRhamComplex}
Hypercohomology of the degree-at-least-\texttt{p} part of the de Rham complex.

\noindent\emph{Lean:} \texttt{abbrev AlgebraicGeometry.ComplexPoint.FilteredDeRhamHypercohomology}, \href{https://github.com/Paul-Lez/HodgeConjecture/blob/main/HodgeConjecture/Definitions/AlgebraicGeometry/HodgeFiltration.lean#L727}{\texttt{HodgeConjecture/Definitions/AlgebraicGeometry/HodgeFiltration.lean:727}}.
\end{definition}

\begin{definition}[AlgebraicGeometry.ComplexPoint.filteredToDeRhamCohomology]\label{def:AlgebraicGeometry.ComplexPoint.filteredToDeRhamCohomology}
\lean{AlgebraicGeometry.ComplexPoint.filteredToDeRhamCohomology}\leanok
\uses{def:AlgebraicGeometry.ComplexPoint.DeRhamHypercohomology, def:AlgebraicGeometry.ComplexPoint.FilteredDeRhamHypercohomology, def:AlgebraicGeometry.ComplexPoint.hodgeFilteredDeRhamInclusion, def:AlgebraicGeometry.ComplexPoint.hypercohomologyMap}
The map from filtered to full de Rham hypercohomology.

\noindent\emph{Lean:} \texttt{def AlgebraicGeometry.ComplexPoint.filteredToDeRhamCohomology}, \href{https://github.com/Paul-Lez/HodgeConjecture/blob/main/HodgeConjecture/Definitions/AlgebraicGeometry/HodgeFiltration.lean#L732}{\texttt{HodgeConjecture/Definitions/AlgebraicGeometry/HodgeFiltration.lean:732}}.
\end{definition}

\begin{definition}[AlgebraicGeometry.ComplexPoint.hodgeFiltration]\label{def:AlgebraicGeometry.ComplexPoint.hodgeFiltration}
\lean{AlgebraicGeometry.ComplexPoint.hodgeFiltration}\leanok
\uses{def:AlgebraicGeometry.ComplexPoint.filteredToDeRhamCohomology}
The Hodge filtration \texttt{F\^{}p} on de Rham hypercohomology.

\noindent\emph{Lean:} \texttt{def AlgebraicGeometry.ComplexPoint.hodgeFiltration}, \href{https://github.com/Paul-Lez/HodgeConjecture/blob/main/HodgeConjecture/Definitions/AlgebraicGeometry/HodgeFiltration.lean#L738}{\texttt{HodgeConjecture/Definitions/AlgebraicGeometry/HodgeFiltration.lean:738}}.
\end{definition}

\begin{definition}[AlgebraicGeometry.ComplexPoint.hodgeFiltrationComplexSubmodule]\label{def:AlgebraicGeometry.ComplexPoint.hodgeFiltrationComplexSubmodule}
\lean{AlgebraicGeometry.ComplexPoint.hodgeFiltrationComplexSubmodule}\leanok
\uses{def:AlgebraicGeometry.ComplexPoint.deRhamComplexSMul, def:AlgebraicGeometry.ComplexPoint.hodgeFilteredDeRhamComplexScalar, def:AlgebraicGeometry.ComplexPoint.hodgeFiltration, def:AlgebraicGeometry.ComplexPoint.hypercohomologyModule}
The Hodge filtration bundled as a complex subspace of de Rham hypercohomology.

\noindent\emph{Lean:} \texttt{def AlgebraicGeometry.ComplexPoint.hodgeFiltrationComplexSubmodule}, \href{https://github.com/Paul-Lez/HodgeConjecture/blob/main/HodgeConjecture/Definitions/AlgebraicGeometry/HodgeFiltration.lean#L756}{\texttt{HodgeConjecture/Definitions/AlgebraicGeometry/HodgeFiltration.lean:756}}.
\end{definition}

\begin{definition}[AlgebraicGeometry.ComplexPoint.complexConstantCohomologyDeRhamAddEquiv]\label{def:AlgebraicGeometry.ComplexPoint.complexConstantCohomologyDeRhamAddEquiv}
\lean{AlgebraicGeometry.ComplexPoint.complexConstantCohomologyDeRhamAddEquiv}\leanok
\uses{def:AlgebraicGeometry.ComplexPoint.complexConstantCohomologyDeRhamEquiv, def:AlgebraicGeometry.ComplexPoint.hypercohomologyMap}
The constant-to-de Rham comparison equivalence, upgraded to an additive equivalence.

\noindent\emph{Lean:} \texttt{def AlgebraicGeometry.ComplexPoint.complexConstantCohomologyDeRhamAddEquiv}, \href{https://github.com/Paul-Lez/HodgeConjecture/blob/main/HodgeConjecture/Definitions/AlgebraicGeometry/HodgeFiltration.lean#L770}{\texttt{HodgeConjecture/Definitions/AlgebraicGeometry/HodgeFiltration.lean:770}}.
\end{definition}

\begin{definition}[AlgebraicGeometry.ComplexPoint.deRhamConj]\label{def:AlgebraicGeometry.ComplexPoint.deRhamConj}
\lean{AlgebraicGeometry.ComplexPoint.deRhamConj}\leanok
\uses{def:AlgebraicGeometry.ComplexPoint.complexConstantCohomologyDeRhamAddEquiv, def:AlgebraicGeometry.ComplexPoint.conjConstantComplexSheafComplexInt, def:AlgebraicGeometry.ComplexPoint.constantsToHolomorphicDeRhamShortComplexSheafificationUnit, def:AlgebraicGeometry.ComplexPoint.fixedChartTransition, def:AlgebraicGeometry.ComplexPoint.formOfAnalyticField, def:AlgebraicGeometry.ComplexPoint.holomorphicDeRhamSheafificationUnit, def:ContinuousAlternatingMap.covectorProduct, def:DifferentialForm.curryDerivativeAt, def:DifferentialForm.radialHomotopy, def:DifferentialForm.radialIntegrandFDeriv, def:DifferentialForm.radialPrimitiveSeries, def:DifferentialForm.radialPullback, def:DifferentialForm.radialPullbackDeriv, def:DifferentialForm.translateForm, def:DifferentialForm.untranslateForm, def:DifferentialForm.zeroFormCoeff}
Complex conjugation on de Rham hypercohomology, transported from the constant sheaf \texttt{ℂ}.

\noindent\emph{Lean:} \texttt{def AlgebraicGeometry.ComplexPoint.deRhamConj}, \href{https://github.com/Paul-Lez/HodgeConjecture/blob/main/HodgeConjecture/Definitions/AlgebraicGeometry/HodgeFiltration.lean#L811}{\texttt{HodgeConjecture/Definitions/AlgebraicGeometry/HodgeFiltration.lean:811}}.
\end{definition}

\begin{definition}[AlgebraicGeometry.ComplexPoint.deRhamConjSemilinear]\label{def:AlgebraicGeometry.ComplexPoint.deRhamConjSemilinear}
\lean{AlgebraicGeometry.ComplexPoint.deRhamConjSemilinear}\leanok
\uses{def:AlgebraicGeometry.ComplexPoint.complexScalarComplexInt, def:AlgebraicGeometry.ComplexPoint.deRhamComplexSMul, def:AlgebraicGeometry.ComplexPoint.deRhamConj, def:AlgebraicGeometry.ComplexPoint.hypercohomologyModule}
Conjugation on de Rham hypercohomology, bundled as a conjugate-linear map.

\noindent\emph{Lean:} \texttt{def AlgebraicGeometry.ComplexPoint.deRhamConjSemilinear}, \href{https://github.com/Paul-Lez/HodgeConjecture/blob/main/HodgeConjecture/Definitions/AlgebraicGeometry/HodgeFiltration.lean#L835}{\texttt{HodgeConjecture/Definitions/AlgebraicGeometry/HodgeFiltration.lean:835}}.
\end{definition}

\begin{definition}[AlgebraicGeometry.ComplexPoint.conjHodgeFiltrationComplexSubmodule]\label{def:AlgebraicGeometry.ComplexPoint.conjHodgeFiltrationComplexSubmodule}
\lean{AlgebraicGeometry.ComplexPoint.conjHodgeFiltrationComplexSubmodule}\leanok
\uses{def:AlgebraicGeometry.ComplexPoint.deRhamConjSemilinear, def:AlgebraicGeometry.ComplexPoint.hodgeFiltrationComplexSubmodule}
The conjugate Hodge filtration \texttt{conj F\^{}p}. Conjugation is an involution, so the preimage of
\texttt{F\^{}p} is also its image.

\noindent\emph{Lean:} \texttt{def AlgebraicGeometry.ComplexPoint.conjHodgeFiltrationComplexSubmodule}, \href{https://github.com/Paul-Lez/HodgeConjecture/blob/main/HodgeConjecture/Definitions/AlgebraicGeometry/HodgeFiltration.lean#L842}{\texttt{HodgeConjecture/Definitions/AlgebraicGeometry/HodgeFiltration.lean:842}}.
\end{definition}

\begin{definition}[AlgebraicGeometry.ComplexPoint.hodgePiece]\label{def:AlgebraicGeometry.ComplexPoint.hodgePiece}
\lean{AlgebraicGeometry.ComplexPoint.hodgePiece}\leanok
\uses{def:AlgebraicGeometry.ComplexPoint.conjHodgeFiltrationComplexSubmodule}
The intersection \texttt{F\^{}p ⊓ conj F\^{}q} in degree \texttt{n}. For smooth projective varieties and
\texttt{p + q = n}, this is the usual \texttt{(p,q)} Hodge piece.

\noindent\emph{Lean:} \texttt{def AlgebraicGeometry.ComplexPoint.hodgePiece}, \href{https://github.com/Paul-Lez/HodgeConjecture/blob/main/HodgeConjecture/Definitions/AlgebraicGeometry/HodgeFiltration.lean#L848}{\texttt{HodgeConjecture/Definitions/AlgebraicGeometry/HodgeFiltration.lean:848}}.
\end{definition}

\begin{definition}[AlgebraicGeometry.ComplexPoint.hodgeClasses]\label{def:AlgebraicGeometry.ComplexPoint.hodgeClasses}
\lean{AlgebraicGeometry.ComplexPoint.hodgeClasses}\leanok
\uses{def:AlgebraicGeometry.ComplexPoint.fieldToDeRhamCohomologyLinear, def:AlgebraicGeometry.ComplexPoint.hodgePiece}
Cohomology classes with coefficients in \texttt{K} whose de Rham images lie in \texttt{F\^{}p ⊓ conj F\^{}p}
in degree \texttt{2p}. When conjugation fixes the image of \texttt{K} in \texttt{ℂ}, see
\texttt{hodgeClasses\_eq\_comap\_hodgeFiltrationComplexSubmodule} for the equivalent \texttt{F\^{}p} condition.

\noindent\emph{Lean:} \texttt{def AlgebraicGeometry.ComplexPoint.hodgeClasses}, \href{https://github.com/Paul-Lez/HodgeConjecture/blob/main/HodgeConjecture/Definitions/AlgebraicGeometry/HodgeFiltration.lean#L868}{\texttt{HodgeConjecture/Definitions/AlgebraicGeometry/HodgeFiltration.lean:868}}.
\end{definition}

\section{Betti comparison with support}\label{sec:HodgeConjecture.Definitions.AlgebraicGeometry.BettiCohomologyWithSupportComparison}
\noindent\emph{File:} \href{https://github.com/Paul-Lez/HodgeConjecture/blob/main/HodgeConjecture/Definitions/AlgebraicGeometry/BettiCohomologyWithSupportComparison.lean}{\texttt{HodgeConjecture/Definitions/AlgebraicGeometry/BettiCohomologyWithSupportComparison.lean}}.

This file proves the categorical input needed to replace the constant-sheaf term in the
mapping-cone model for supported cohomology by a singular-cochain resolution. In particular,
direct image along an open embedding preserves injective additive sheaves.

\begin{definition}[CochainComplex.liftToInjective]\label{def:CochainComplex.liftToInjective}
\lean{CochainComplex.liftToInjective}\leanok
A map into a bounded-below degreewise-injective complex extends strictly across a monic
quasi-isomorphism. This is the lifting property in the injective model structure.

\noindent\emph{Lean:} \texttt{def CochainComplex.liftToInjective}, \href{https://github.com/Paul-Lez/HodgeConjecture/blob/main/HodgeConjecture/Definitions/AlgebraicGeometry/BettiCohomologyWithSupportComparison.lean#L91}{\texttt{HodgeConjecture/Definitions/AlgebraicGeometry/BettiCohomologyWithSupportComparison.lean:91}}.
\end{definition}

\section{Betti cohomology and global sections}\label{sec:HodgeConjecture.Definitions.AlgebraicGeometry.BettiGlobalSectionsComparison}
\noindent\emph{File:} \href{https://github.com/Paul-Lez/HodgeConjecture/blob/main/HodgeConjecture/Definitions/AlgebraicGeometry/BettiGlobalSectionsComparison.lean}{\texttt{HodgeConjecture/Definitions/AlgebraicGeometry/BettiGlobalSectionsComparison.lean}}.

This file identifies rational singular cohomology with the cohomology of the global-section
complex of the sheafified singular-cochain resolution on a paracompact Hausdorff space. It also
computes the hypercohomology of that resolution on a hereditarily paracompact Hausdorff space.

The derived comparison uses a bounded-below termwise-injective replacement. The mapping-cone
argument in \texttt{FlasqueQuasiIsoGlobalSections} proves that its quasi-isomorphism remains a
quasi-isomorphism after taking global sections. Thus no spectral sequence or acyclic-resolution
theorem is assumed.

\begin{definition}[CochainComplex.HomComplex.fromSingleZeroIsoPreadditiveCoyoneda]\label{def:CochainComplex.HomComplex.fromSingleZeroIsoPreadditiveCoyoneda}
\lean{CochainComplex.HomComplex.fromSingleZeroIsoPreadditiveCoyoneda}\leanok
The Hom complex from an object placed in degree zero is the degreewise preadditive
coyoneda functor applied to the target complex.

\noindent\emph{Lean:} \texttt{def CochainComplex.HomComplex.fromSingleZeroIsoPreadditiveCoyoneda}, \href{https://github.com/Paul-Lez/HodgeConjecture/blob/main/HodgeConjecture/Definitions/AlgebraicGeometry/BettiGlobalSectionsComparison.lean#L52}{\texttt{HodgeConjecture/Definitions/AlgebraicGeometry/BettiGlobalSectionsComparison.lean:52}}.
\end{definition}

\begin{definition}[TopCat.Sheaf.integerConstantHomEquivGlobalSections]\label{def:TopCat.Sheaf.integerConstantHomEquivGlobalSections}
\lean{TopCat.Sheaf.integerConstantHomEquivGlobalSections}\leanok
Morphisms from the constant integer sheaf are the same as global sections. This is the
degree-zero adjunction underlying the global-sections comparison below.

\noindent\emph{Lean:} \texttt{def TopCat.Sheaf.integerConstantHomEquivGlobalSections}, \href{https://github.com/Paul-Lez/HodgeConjecture/blob/main/HodgeConjecture/Definitions/AlgebraicGeometry/BettiGlobalSectionsComparison.lean#L89}{\texttt{HodgeConjecture/Definitions/AlgebraicGeometry/BettiGlobalSectionsComparison.lean:89}}.
\end{definition}

\begin{definition}[TopCat.Sheaf.integerConstantHomAddEquivGlobalSections]\label{def:TopCat.Sheaf.integerConstantHomAddEquivGlobalSections}
\lean{TopCat.Sheaf.integerConstantHomAddEquivGlobalSections}\leanok
Additive form of \texttt{integerConstantHomEquivGlobalSections}. The explicit local Hom-group
instance avoids depending on reducibility-sensitive typeclass search through the sheaf
subcategory.

\noindent\emph{Lean:} \texttt{def TopCat.Sheaf.integerConstantHomAddEquivGlobalSections}, \href{https://github.com/Paul-Lez/HodgeConjecture/blob/main/HodgeConjecture/Definitions/AlgebraicGeometry/BettiGlobalSectionsComparison.lean#L100}{\texttt{HodgeConjecture/Definitions/AlgebraicGeometry/BettiGlobalSectionsComparison.lean:100}}.
\end{definition}

\begin{definition}[TopCat.Sheaf.integerConstantSingleComplex]\label{def:TopCat.Sheaf.integerConstantSingleComplex}
\lean{TopCat.Sheaf.integerConstantSingleComplex}\leanok
The integer sheaf placed in cohomological degree zero.

\noindent\emph{Lean:} \texttt{def TopCat.Sheaf.integerConstantSingleComplex}, \href{https://github.com/Paul-Lez/HodgeConjecture/blob/main/HodgeConjecture/Definitions/AlgebraicGeometry/BettiGlobalSectionsComparison.lean#L134}{\texttt{HodgeConjecture/Definitions/AlgebraicGeometry/BettiGlobalSectionsComparison.lean:134}}.
\end{definition}

\begin{definition}[TopCat.Sheaf.globalSectionsComplexInt]\label{def:TopCat.Sheaf.globalSectionsComplexInt}
\lean{TopCat.Sheaf.globalSectionsComplexInt}\leanok
\uses{def:TopCat.Sheaf.IsFlasque.BoundedBelowComplex.globalSectionsComplex}
Evaluation of an integer-indexed sheaf complex on the top open subset.

\noindent\emph{Lean:} \texttt{def TopCat.Sheaf.globalSectionsComplexInt}, \href{https://github.com/Paul-Lez/HodgeConjecture/blob/main/HodgeConjecture/Definitions/AlgebraicGeometry/BettiGlobalSectionsComparison.lean#L141}{\texttt{HodgeConjecture/Definitions/AlgebraicGeometry/BettiGlobalSectionsComparison.lean:141}}.
\end{definition}

\begin{definition}[TopCat.Sheaf.integerConstantHomIsoGlobalSectionsFunctor]\label{def:TopCat.Sheaf.integerConstantHomIsoGlobalSectionsFunctor}
\lean{TopCat.Sheaf.integerConstantHomIsoGlobalSectionsFunctor}\leanok
\uses{def:TopCat.Sheaf.IsFlasque.BoundedBelowComplex.globalSectionsFunctor, def:TopCat.Sheaf.integerConstantHomAddEquivGlobalSections, def:TopCat.Sheaf.integerConstantHomEquivGlobalSections}
The additive constant-sheaf adjunction identifies the coyoneda functor represented by the
constant integer sheaf with the global-sections functor.

\noindent\emph{Lean:} \texttt{def TopCat.Sheaf.integerConstantHomIsoGlobalSectionsFunctor}, \href{https://github.com/Paul-Lez/HodgeConjecture/blob/main/HodgeConjecture/Definitions/AlgebraicGeometry/BettiGlobalSectionsComparison.lean#L149}{\texttt{HodgeConjecture/Definitions/AlgebraicGeometry/BettiGlobalSectionsComparison.lean:149}}.
\end{definition}

\begin{definition}[TopCat.Sheaf.homComplexSingleIntegerIsoGlobalSections]\label{def:TopCat.Sheaf.homComplexSingleIntegerIsoGlobalSections}
\lean{TopCat.Sheaf.homComplexSingleIntegerIsoGlobalSections}\leanok
\uses{def:CochainComplex.HomComplex.fromSingleZeroIsoPreadditiveCoyoneda, def:TopCat.Sheaf.globalSectionsComplexInt, def:TopCat.Sheaf.integerConstantHomIsoGlobalSectionsFunctor, def:TopCat.Sheaf.integerConstantSingleComplex}
The Hom complex from the degree-zero integer sheaf is canonically the complex of global
sections.

\noindent\emph{Lean:} \texttt{def TopCat.Sheaf.homComplexSingleIntegerIsoGlobalSections}, \href{https://github.com/Paul-Lez/HodgeConjecture/blob/main/HodgeConjecture/Definitions/AlgebraicGeometry/BettiGlobalSectionsComparison.lean#L168}{\texttt{HodgeConjecture/Definitions/AlgebraicGeometry/BettiGlobalSectionsComparison.lean:168}}.
\end{definition}

\begin{definition}[AlgebraicGeometry.ComplexPoint.constantIntegerSheafComplexIntIsoSingle]\label{def:AlgebraicGeometry.ComplexPoint.constantIntegerSheafComplexIntIsoSingle}
\lean{AlgebraicGeometry.ComplexPoint.constantIntegerSheafComplexIntIsoSingle}\leanok
\uses{def:AlgebraicGeometry.ComplexPoint.constantIntegerSheafComplexInt, def:TopCat.Sheaf.integerConstantSingleComplex}
The integer constant-sheaf complex used to define hypercohomology is the degree-zero
integer constant sheaf, after extending its natural-number grading to integer degrees.

\noindent\emph{Lean:} \texttt{def AlgebraicGeometry.ComplexPoint.constantIntegerSheafComplexIntIsoSingle}, \href{https://github.com/Paul-Lez/HodgeConjecture/blob/main/HodgeConjecture/Definitions/AlgebraicGeometry/BettiGlobalSectionsComparison.lean#L192}{\texttt{HodgeConjecture/Definitions/AlgebraicGeometry/BettiGlobalSectionsComparison.lean:192}}.
\end{definition}

\section{Hypercohomology and singular cohomology with support}\label{sec:HodgeConjecture.Definitions.AlgebraicGeometry.BettiSupportSingularHypercohomologyComparison}
\noindent\emph{File:} \href{https://github.com/Paul-Lez/HodgeConjecture/blob/main/HodgeConjecture/Definitions/AlgebraicGeometry/BettiSupportSingularHypercohomologyComparison.lean}{\texttt{HodgeConjecture/Definitions/AlgebraicGeometry/BettiSupportSingularHypercohomologyComparison.lean}}.

This file develops the bounded-below flasque comparison needed for cohomology with support.
It applies the same explicit injective-replacement argument used for ordinary Betti cohomology
to the mapping cone of singular restriction. No hypercohomology spectral sequence is assumed.

\begin{definition}[AlgebraicGeometry.ComplexPoint.bettiSupportHypercohomologyComparisonHasDerivedCategory]\label{def:AlgebraicGeometry.ComplexPoint.bettiSupportHypercohomologyComparisonHasDerivedCategory}
\lean{AlgebraicGeometry.ComplexPoint.bettiSupportHypercohomologyComparisonHasDerivedCategory}\leanok
\uses{def:AlgebraicGeometry.ComplexPoint.AnalyticAdditiveSheaf}
\noindent\emph{Lean:} \texttt{def AlgebraicGeometry.ComplexPoint.bettiSupportHypercohomologyComparisonHasDerivedCategory}, \href{https://github.com/Paul-Lez/HodgeConjecture/blob/main/HodgeConjecture/Definitions/AlgebraicGeometry/BettiSupportSingularHypercohomologyComparison.lean#L44}{\texttt{HodgeConjecture/Definitions/AlgebraicGeometry/BettiSupportSingularHypercohomologyComparison.lean:44}}.
\end{definition}

\begin{definition}[AlgebraicGeometry.ComplexPoint.isoHomCongrAddEquiv]\label{def:AlgebraicGeometry.ComplexPoint.isoHomCongrAddEquiv}
\lean{AlgebraicGeometry.ComplexPoint.isoHomCongrAddEquiv}\leanok
Conjugating a morphism by two isomorphisms is an additive equivalence of Hom groups.

\noindent\emph{Lean:} \texttt{def AlgebraicGeometry.ComplexPoint.isoHomCongrAddEquiv}, \href{https://github.com/Paul-Lez/HodgeConjecture/blob/main/HodgeConjecture/Definitions/AlgebraicGeometry/BettiSupportSingularHypercohomologyComparison.lean#L128}{\texttt{HodgeConjecture/Definitions/AlgebraicGeometry/BettiSupportSingularHypercohomologyComparison.lean:128}}.
\end{definition}

\begin{definition}[AlgebraicGeometry.ComplexPoint.hypercohomologyAddEquivDerived]\label{def:AlgebraicGeometry.ComplexPoint.hypercohomologyAddEquivDerived}
\lean{AlgebraicGeometry.ComplexPoint.hypercohomologyAddEquivDerived}\leanok
\uses{def:AlgebraicGeometry.ComplexPoint.Hypercohomology, def:AlgebraicGeometry.ComplexPoint.bettiSupportHypercohomologyComparisonHasDerivedCategory}
The chosen additive structure on hypercohomology is transported from shifted morphisms in
the derived category.

\noindent\emph{Lean:} \texttt{def AlgebraicGeometry.ComplexPoint.hypercohomologyAddEquivDerived}, \href{https://github.com/Paul-Lez/HodgeConjecture/blob/main/HodgeConjecture/Definitions/AlgebraicGeometry/BettiSupportSingularHypercohomologyComparison.lean#L136}{\texttt{HodgeConjecture/Definitions/AlgebraicGeometry/BettiSupportSingularHypercohomologyComparison.lean:136}}.
\end{definition}

\begin{definition}[AlgebraicGeometry.ComplexPoint.kInjectiveDerivedHomAddEquivCohomologyClass]\label{def:AlgebraicGeometry.ComplexPoint.kInjectiveDerivedHomAddEquivCohomologyClass}
\lean{AlgebraicGeometry.ComplexPoint.kInjectiveDerivedHomAddEquivCohomologyClass}\leanok
\uses{def:AlgebraicGeometry.ComplexPoint.isoHomCongrAddEquiv}
For a K-injective target, shifted derived morphisms are additively identified with
cohomology classes in the Hom complex.

\noindent\emph{Lean:} \texttt{def AlgebraicGeometry.ComplexPoint.kInjectiveDerivedHomAddEquivCohomologyClass}, \href{https://github.com/Paul-Lez/HodgeConjecture/blob/main/HodgeConjecture/Definitions/AlgebraicGeometry/BettiSupportSingularHypercohomologyComparison.lean#L148}{\texttt{HodgeConjecture/Definitions/AlgebraicGeometry/BettiSupportSingularHypercohomologyComparison.lean:148}}.
\end{definition}

\section{Naturality of the hypercohomology/global-sections comparison}\label{sec:HodgeConjecture.Definitions.AlgebraicGeometry.HypercohomologyGlobalSectionsNaturality}
\noindent\emph{File:} \href{https://github.com/Paul-Lez/HodgeConjecture/blob/main/HodgeConjecture/Definitions/AlgebraicGeometry/HypercohomologyGlobalSectionsNaturality.lean}{\texttt{HodgeConjecture/Definitions/AlgebraicGeometry/HypercohomologyGlobalSectionsNaturality.lean}}.

\begin{definition}[AlgebraicGeometry.ComplexPoint.hypercohomologyNaturalitySheafDerivedCategory]\label{def:AlgebraicGeometry.ComplexPoint.hypercohomologyNaturalitySheafDerivedCategory}
\lean{AlgebraicGeometry.ComplexPoint.hypercohomologyNaturalitySheafDerivedCategory}\leanok
\uses{def:AlgebraicGeometry.ComplexPoint.AnalyticAdditiveSheaf}
\noindent\emph{Lean:} \texttt{def AlgebraicGeometry.ComplexPoint.hypercohomologyNaturalitySheafDerivedCategory}, \href{https://github.com/Paul-Lez/HodgeConjecture/blob/main/HodgeConjecture/Definitions/AlgebraicGeometry/HypercohomologyGlobalSectionsNaturality.lean#L31}{\texttt{HodgeConjecture/Definitions/AlgebraicGeometry/HypercohomologyGlobalSectionsNaturality.lean:31}}.
\end{definition}

\begin{definition}[AlgebraicGeometry.ComplexPoint.derivedHomAddEquivGlobalSectionsKInjective]\label{def:AlgebraicGeometry.ComplexPoint.derivedHomAddEquivGlobalSectionsKInjective}
\lean{AlgebraicGeometry.ComplexPoint.derivedHomAddEquivGlobalSectionsKInjective}\leanok
\uses{def:AlgebraicGeometry.ComplexPoint.hypercohomologyNaturalitySheafDerivedCategory, def:AlgebraicGeometry.ComplexPoint.kInjectiveDerivedHomAddEquivCohomologyClass, def:TopCat.Sheaf.homComplexSingleIntegerIsoGlobalSections}
On a K-injective sheaf complex, derived morphisms from the integer
constant sheaf are computed by actual global sections, with no further
replacement complex.

\noindent\emph{Lean:} \texttt{def AlgebraicGeometry.ComplexPoint.derivedHomAddEquivGlobalSectionsKInjective}, \href{https://github.com/Paul-Lez/HodgeConjecture/blob/main/HodgeConjecture/Definitions/AlgebraicGeometry/HypercohomologyGlobalSectionsNaturality.lean#L35}{\texttt{HodgeConjecture/Definitions/AlgebraicGeometry/HypercohomologyGlobalSectionsNaturality.lean:35}}.
\end{definition}

\begin{definition}[AlgebraicGeometry.ComplexPoint.hypercohomologyAddEquivGlobalSectionsKInjective]\label{def:AlgebraicGeometry.ComplexPoint.hypercohomologyAddEquivGlobalSectionsKInjective}
\lean{AlgebraicGeometry.ComplexPoint.hypercohomologyAddEquivGlobalSectionsKInjective}\leanok
\uses{def:AlgebraicGeometry.ComplexPoint.constantIntegerSheafComplexIntIsoSingle, def:AlgebraicGeometry.ComplexPoint.derivedHomAddEquivGlobalSectionsKInjective, def:AlgebraicGeometry.ComplexPoint.hypercohomologyAddEquivDerived}
Hypercohomology of an actual K-injective complex is its global-section
cohomology. This direct form exposes naturality without choosing another
injective resolution.

\noindent\emph{Lean:} \texttt{def AlgebraicGeometry.ComplexPoint.hypercohomologyAddEquivGlobalSectionsKInjective}, \href{https://github.com/Paul-Lez/HodgeConjecture/blob/main/HodgeConjecture/Definitions/AlgebraicGeometry/HypercohomologyGlobalSectionsNaturality.lean#L50}{\texttt{HodgeConjecture/Definitions/AlgebraicGeometry/HypercohomologyGlobalSectionsNaturality.lean:50}}.
\end{definition}

\section{Normalized injective models for the rational support cone}\label{sec:HodgeConjecture.Definitions.AlgebraicGeometry.DerivedSupportRationalConeComparison}
\noindent\emph{File:} \href{https://github.com/Paul-Lez/HodgeConjecture/blob/main/HodgeConjecture/Definitions/AlgebraicGeometry/DerivedSupportRationalConeComparison.lean}{\texttt{HodgeConjecture/Definitions/AlgebraicGeometry/DerivedSupportRationalConeComparison.lean}}.

The existing rational support object resolves constants on the complement
independently of the ambient space. This file connects that model with the
actual restriction of an ambient injective resolution. All comparison maps
extend the given constant restriction, rather than choosing an abstract
equivalence between cohomology groups.

\begin{definition}[AlgebraicGeometry.ComplexPoint.ambientRationalInjectiveComplex]\label{def:AlgebraicGeometry.ComplexPoint.ambientRationalInjectiveComplex}
\lean{AlgebraicGeometry.ComplexPoint.ambientRationalInjectiveComplex}\leanok
\uses{def:AlgebraicGeometry.ComplexPoint.AnalyticAdditiveSheaf, def:TopCat.Sheaf.ambientConstantInjectiveResolution}
The standard ambient rational injective resolution in integer degrees.

\noindent\emph{Lean:} \texttt{def AlgebraicGeometry.ComplexPoint.ambientRationalInjectiveComplex}, \href{https://github.com/Paul-Lez/HodgeConjecture/blob/main/HodgeConjecture/Definitions/AlgebraicGeometry/DerivedSupportRationalConeComparison.lean#L29}{\texttt{HodgeConjecture/Definitions/AlgebraicGeometry/DerivedSupportRationalConeComparison.lean:29}}.
\end{definition}

\begin{definition}[AlgebraicGeometry.ComplexPoint.ambientRationalInjectiveAugmentation]\label{def:AlgebraicGeometry.ComplexPoint.ambientRationalInjectiveAugmentation}
\lean{AlgebraicGeometry.ComplexPoint.ambientRationalInjectiveAugmentation}\leanok
\uses{def:AlgebraicGeometry.ComplexPoint.ambientRationalInjectiveComplex, def:AlgebraicGeometry.ComplexPoint.constantFieldSheafComplexInt}
Its actual constant augmentation.

\noindent\emph{Lean:} \texttt{def AlgebraicGeometry.ComplexPoint.ambientRationalInjectiveAugmentation}, \href{https://github.com/Paul-Lez/HodgeConjecture/blob/main/HodgeConjecture/Definitions/AlgebraicGeometry/DerivedSupportRationalConeComparison.lean#L36}{\texttt{HodgeConjecture/Definitions/AlgebraicGeometry/DerivedSupportRationalConeComparison.lean:36}}.
\end{definition}

\section{Constructed normal coordinates for smooth closed immersions}\label{sec:HodgeConjecture.Definitions.AlgebraicGeometry.ClosedImmersionNormalCoordinates}
\noindent\emph{File:} \href{https://github.com/Paul-Lez/HodgeConjecture/blob/main/HodgeConjecture/Definitions/AlgebraicGeometry/ClosedImmersionNormalCoordinates.lean}{\texttt{HodgeConjecture/Definitions/AlgebraicGeometry/ClosedImmersionNormalCoordinates.lean}}.

The derivative projection here comes from lifted algebraic coordinate sections. The normal
chart is then constructed by the inverse function theorem. Finally the actual topological
embedding theorem excludes remote branches and identifies the whole local analytic support
with zero normal coordinate. No analytic immersion or flattening equivalence is supplied.

\begin{definition}[AlgebraicGeometry.ComplexPoint.closedImmersionDerivativeProjection]\label{def:AlgebraicGeometry.ComplexPoint.closedImmersionDerivativeProjection}
\lean{AlgebraicGeometry.ComplexPoint.closedImmersionDerivativeProjection}\leanok
\uses{def:AlgebraicGeometry.ComplexAlgHom.isStandardEtaleBaseAlgHom, def:AlgebraicGeometry.ComplexAlgHom.standardEtaleBivariateEvaluation, def:AlgebraicGeometry.ComplexPoint.affineSpecComplexPointMap, def:AlgebraicGeometry.ComplexPoint.inclusionInComplexCharts, def:AlgebraicGeometry.ComplexPoint.polynomialSpecToAffineSpacePointMap, def:AlgebraicGeometry.LocalEtaleCoordinates.algebraicCoordinates, def:AlgebraicGeometry.LocalEtaleCoordinates.ambientAnalyticCoordinates, def:AlgebraicGeometry.LocalEtaleCoordinates.ambientCoordinateSection}
The actual derivative left inverse obtained from lifted intrinsic coordinates.

\noindent\emph{Lean:} \texttt{def AlgebraicGeometry.ComplexPoint.closedImmersionDerivativeProjection}, \href{https://github.com/Paul-Lez/HodgeConjecture/blob/main/HodgeConjecture/Definitions/AlgebraicGeometry/ClosedImmersionNormalCoordinates.lean#L31}{\texttt{HodgeConjecture/Definitions/AlgebraicGeometry/ClosedImmersionNormalCoordinates.lean:31}}.
\end{definition}

\begin{definition}[AlgebraicGeometry.ComplexPoint.closedImmersionNormalChart]\label{def:AlgebraicGeometry.ComplexPoint.closedImmersionNormalChart}
\lean{AlgebraicGeometry.ComplexPoint.closedImmersionNormalChart}\leanok
\uses{def:AlgebraicGeometry.ComplexPoint.closedImmersionDerivativeProjection, def:HasStrictFDerivAt.normalChart}
Actual normal-coordinate parametrization, with the kernel of the constructed
derivative projection as its complex normal space.

\noindent\emph{Lean:} \texttt{def AlgebraicGeometry.ComplexPoint.closedImmersionNormalChart}, \href{https://github.com/Paul-Lez/HodgeConjecture/blob/main/HodgeConjecture/Definitions/AlgebraicGeometry/ClosedImmersionNormalCoordinates.lean#L43}{\texttt{HodgeConjecture/Definitions/AlgebraicGeometry/ClosedImmersionNormalCoordinates.lean:43}}.
\end{definition}

\begin{definition}[AlgebraicGeometry.ComplexPoint.closedImmersionFlatteningChart]\label{def:AlgebraicGeometry.ComplexPoint.closedImmersionFlatteningChart}
\lean{AlgebraicGeometry.ComplexPoint.closedImmersionFlatteningChart}\leanok
\uses{def:AlgebraicGeometry.ComplexPoint.closedImmersionNormalChart}
The ambient flattening chart, restricted so that its zero-normal locus is exactly
the actual embedded support on its entire source. Every ingredient has been constructed
from the given smooth closed immersion.

\noindent\emph{Lean:} \texttt{def AlgebraicGeometry.ComplexPoint.closedImmersionFlatteningChart}, \href{https://github.com/Paul-Lez/HodgeConjecture/blob/main/HodgeConjecture/Definitions/AlgebraicGeometry/ClosedImmersionNormalCoordinates.lean#L167}{\texttt{HodgeConjecture/Definitions/AlgebraicGeometry/ClosedImmersionNormalCoordinates.lean:167}}.
\end{definition}

\begin{definition}[AlgebraicGeometry.ComplexPoint.closedImmersionNormalKernelEquiv]\label{def:AlgebraicGeometry.ComplexPoint.closedImmersionNormalKernelEquiv}
\lean{AlgebraicGeometry.ComplexPoint.closedImmersionNormalKernelEquiv}\leanok
\uses{def:AlgebraicGeometry.ComplexPoint.closedImmersionDerivativeProjection, def:ContinuousLinearMap.splitKernelEquiv}
Identifying the actual normal space with standard complex coordinates uses only its
proved complex dimension. This is a complex-linear coordinate choice, not a choice of
homology generator or an orientation class.

\noindent\emph{Lean:} \texttt{def AlgebraicGeometry.ComplexPoint.closedImmersionNormalKernelEquiv}, \href{https://github.com/Paul-Lez/HodgeConjecture/blob/main/HodgeConjecture/Definitions/AlgebraicGeometry/ClosedImmersionNormalCoordinates.lean#L189}{\texttt{HodgeConjecture/Definitions/AlgebraicGeometry/ClosedImmersionNormalCoordinates.lean:189}}.
\end{definition}

\begin{definition}[AlgebraicGeometry.ComplexPoint.closedImmersionStandardFlatteningChart]\label{def:AlgebraicGeometry.ComplexPoint.closedImmersionStandardFlatteningChart}
\lean{AlgebraicGeometry.ComplexPoint.closedImmersionStandardFlatteningChart}\leanok
\uses{def:AlgebraicGeometry.ComplexPoint.closedImmersionFlatteningChart, def:AlgebraicGeometry.ComplexPoint.closedImmersionNormalKernelEquiv}
The actual ambient support-flattening chart in standard tangent and normal complex
spaces, ready for the normal-slice pair calculation.

\noindent\emph{Lean:} \texttt{def AlgebraicGeometry.ComplexPoint.closedImmersionStandardFlatteningChart}, \href{https://github.com/Paul-Lez/HodgeConjecture/blob/main/HodgeConjecture/Definitions/AlgebraicGeometry/ClosedImmersionNormalCoordinates.lean#L199}{\texttt{HodgeConjecture/Definitions/AlgebraicGeometry/ClosedImmersionNormalCoordinates.lean:199}}.
\end{definition}

\section{Holomorphic support-flattening charts}\label{sec:HodgeConjecture.Definitions.AlgebraicGeometry.HolomorphicClosedImmersionCharts}
\noindent\emph{File:} \href{https://github.com/Paul-Lez/HodgeConjecture/blob/main/HodgeConjecture/Definitions/AlgebraicGeometry/HolomorphicClosedImmersionCharts.lean}{\texttt{HodgeConjecture/Definitions/AlgebraicGeometry/HolomorphicClosedImmersionCharts.lean}}.

Restricting to the open loci where a normal coordinate change and its inverse are analytic
upgrades centerwise analyticity to analyticity throughout each selected chart. All coordinate
changes come from the smooth closed immersion constructed earlier.

\begin{definition}[OpenPartialHomeomorph.biAnalyticLocus]\label{def:OpenPartialHomeomorph.biAnalyticLocus}
\lean{OpenPartialHomeomorph.biAnalyticLocus}\leanok
The actual open locus on which both directions of a coordinate change are analytic.

\noindent\emph{Lean:} \texttt{def OpenPartialHomeomorph.biAnalyticLocus}, \href{https://github.com/Paul-Lez/HodgeConjecture/blob/main/HodgeConjecture/Definitions/AlgebraicGeometry/HolomorphicClosedImmersionCharts.lean#L28}{\texttt{HodgeConjecture/Definitions/AlgebraicGeometry/HolomorphicClosedImmersionCharts.lean:28}}.
\end{definition}

\begin{definition}[OpenPartialHomeomorph.biAnalyticRestrict]\label{def:OpenPartialHomeomorph.biAnalyticRestrict}
\lean{OpenPartialHomeomorph.biAnalyticRestrict}\leanok
\uses{def:OpenPartialHomeomorph.biAnalyticLocus}
Restricting to this proved open set introduces no analytic-equivalence input.

\noindent\emph{Lean:} \texttt{def OpenPartialHomeomorph.biAnalyticRestrict}, \href{https://github.com/Paul-Lez/HodgeConjecture/blob/main/HodgeConjecture/Definitions/AlgebraicGeometry/HolomorphicClosedImmersionCharts.lean#L35}{\texttt{HodgeConjecture/Definitions/AlgebraicGeometry/HolomorphicClosedImmersionCharts.lean:35}}.
\end{definition}

\begin{definition}[AlgebraicGeometry.ComplexPoint.closedImmersionNormalCoordinatesLinearEquiv]\label{def:AlgebraicGeometry.ComplexPoint.closedImmersionNormalCoordinatesLinearEquiv}
\lean{AlgebraicGeometry.ComplexPoint.closedImmersionNormalCoordinatesLinearEquiv}\leanok
\uses{def:AlgebraicGeometry.ComplexPoint.closedImmersionNormalKernelEquiv}
The actual complex-linear identification of tangent and normal product coordinates.

\noindent\emph{Lean:} \texttt{def AlgebraicGeometry.ComplexPoint.closedImmersionNormalCoordinatesLinearEquiv}, \href{https://github.com/Paul-Lez/HodgeConjecture/blob/main/HodgeConjecture/Definitions/AlgebraicGeometry/HolomorphicClosedImmersionCharts.lean#L56}{\texttt{HodgeConjecture/Definitions/AlgebraicGeometry/HolomorphicClosedImmersionCharts.lean:56}}.
\end{definition}

\begin{definition}[AlgebraicGeometry.ComplexPoint.closedImmersionNormalCoordinateChange]\label{def:AlgebraicGeometry.ComplexPoint.closedImmersionNormalCoordinateChange}
\lean{AlgebraicGeometry.ComplexPoint.closedImmersionNormalCoordinateChange}\leanok
\uses{def:AlgebraicGeometry.ComplexPoint.closedImmersionNormalChart, def:AlgebraicGeometry.ComplexPoint.closedImmersionNormalCoordinatesLinearEquiv}
The normal coordinate change inside the canonical ambient complex chart.

\noindent\emph{Lean:} \texttt{def AlgebraicGeometry.ComplexPoint.closedImmersionNormalCoordinateChange}, \href{https://github.com/Paul-Lez/HodgeConjecture/blob/main/HodgeConjecture/Definitions/AlgebraicGeometry/HolomorphicClosedImmersionCharts.lean#L64}{\texttt{HodgeConjecture/Definitions/AlgebraicGeometry/HolomorphicClosedImmersionCharts.lean:64}}.
\end{definition}

\begin{definition}[AlgebraicGeometry.ComplexPoint.closedImmersionHolomorphicFlatteningChart]\label{def:AlgebraicGeometry.ComplexPoint.closedImmersionHolomorphicFlatteningChart}
\lean{AlgebraicGeometry.ComplexPoint.closedImmersionHolomorphicFlatteningChart}\leanok
\uses{def:AlgebraicGeometry.ComplexPoint.closedImmersionNormalCoordinateChange, def:AlgebraicGeometry.ComplexPoint.closedImmersionStandardFlatteningChart, def:OpenPartialHomeomorph.biAnalyticRestrict}
An actual support-flattening chart whose normal coordinate change is holomorphic in
both directions throughout its source. It contains the distinguished support point.

\noindent\emph{Lean:} \texttt{def AlgebraicGeometry.ComplexPoint.closedImmersionHolomorphicFlatteningChart}, \href{https://github.com/Paul-Lez/HodgeConjecture/blob/main/HodgeConjecture/Definitions/AlgebraicGeometry/HolomorphicClosedImmersionCharts.lean#L70}{\texttt{HodgeConjecture/Definitions/AlgebraicGeometry/HolomorphicClosedImmersionCharts.lean:70}}.
\end{definition}

\begin{definition}[AlgebraicGeometry.ComplexPoint.closedImmersionNormalTransition]\label{def:AlgebraicGeometry.ComplexPoint.closedImmersionNormalTransition}
\lean{AlgebraicGeometry.ComplexPoint.closedImmersionNormalTransition}\leanok
\uses{def:AlgebraicGeometry.ComplexPoint.closedImmersionHolomorphicFlatteningChart}
The genuine transition between two constructed holomorphic support-flattening charts.

\noindent\emph{Lean:} \texttt{def AlgebraicGeometry.ComplexPoint.closedImmersionNormalTransition}, \href{https://github.com/Paul-Lez/HodgeConjecture/blob/main/HodgeConjecture/Definitions/AlgebraicGeometry/HolomorphicClosedImmersionCharts.lean#L105}{\texttt{HodgeConjecture/Definitions/AlgebraicGeometry/HolomorphicClosedImmersionCharts.lean:105}}.
\end{definition}

\section{Exactly normalized smooth-support coclasses on actual overlaps}\label{sec:HodgeConjecture.Definitions.AlgebraicGeometry.SmoothClosedSupportCoclassOverlap}
\noindent\emph{File:} \href{https://github.com/Paul-Lez/HodgeConjecture/blob/main/HodgeConjecture/Definitions/AlgebraicGeometry/SmoothClosedSupportCoclassOverlap.lean}{\texttt{HodgeConjecture/Definitions/AlgebraicGeometry/SmoothClosedSupportCoclassOverlap.lean}}.

The previously constructed local normal coclass is the pullback of the fixed complex
normal coclass along the actual normal coordinate projection. The actual holomorphic
closed-immersion charts prove that these coclasses agree on sufficiently small common
ambient neighborhoods. No transition compatibility or purity equivalence is supplied.

\begin{definition}[AlgebraicGeometry.ComplexPoint.smoothClosedSupportChartCoclass]\label{def:AlgebraicGeometry.ComplexPoint.smoothClosedSupportChartCoclass}
\lean{AlgebraicGeometry.ComplexPoint.smoothClosedSupportChartCoclass}\leanok
\uses{def:AlgebraicGeometry.ComplexPoint.closedImmersionHolomorphicFlatteningChart, def:AlgebraicTopology.Singular.chartNormalProjectionCoclass}
Actual normal-projection coclass on any subset of a holomorphic flattening chart.

\noindent\emph{Lean:} \texttt{def AlgebraicGeometry.ComplexPoint.smoothClosedSupportChartCoclass}, \href{https://github.com/Paul-Lez/HodgeConjecture/blob/main/HodgeConjecture/Definitions/AlgebraicGeometry/SmoothClosedSupportCoclassOverlap.lean#L32}{\texttt{HodgeConjecture/Definitions/AlgebraicGeometry/SmoothClosedSupportCoclassOverlap.lean:32}}.
\end{definition}

\section{Restricting the source of a closed immersion without losing closedness}\label{sec:HodgeConjecture.Definitions.AlgebraicGeometry.ClosedImmersionSourceOpen}
\noindent\emph{File:} \href{https://github.com/Paul-Lez/HodgeConjecture/blob/main/HodgeConjecture/Definitions/AlgebraicGeometry/ClosedImmersionSourceOpen.lean}{\texttt{HodgeConjecture/Definitions/AlgebraicGeometry/ClosedImmersionSourceOpen.lean}}.

For a closed immersion \texttt{i : Y ⟶ X} and an open \texttt{A ⊆ Y}, delete \texttt{i(Y \textbackslash{} A)}
from the target. The actual map from \texttt{A} into this target open is closed, and its
image is exactly the restriction of the original image. This allows a smooth
source of varying local dimensions to be treated on fixed-dimensional opens.

\begin{definition}[AlgebraicGeometry.closedImmersionSourceOpenTarget]\label{def:AlgebraicGeometry.closedImmersionSourceOpenTarget}
\lean{AlgebraicGeometry.closedImmersionSourceOpenTarget}\leanok
Delete the actual closed image of the discarded source complement.

\noindent\emph{Lean:} \texttt{def AlgebraicGeometry.closedImmersionSourceOpenTarget}, \href{https://github.com/Paul-Lez/HodgeConjecture/blob/main/HodgeConjecture/Definitions/AlgebraicGeometry/ClosedImmersionSourceOpen.lean#L26}{\texttt{HodgeConjecture/Definitions/AlgebraicGeometry/ClosedImmersionSourceOpen.lean:26}}.
\end{definition}

\begin{definition}[AlgebraicGeometry.closedImmersionSourceOpenLift]\label{def:AlgebraicGeometry.closedImmersionSourceOpenLift}
\lean{AlgebraicGeometry.closedImmersionSourceOpenLift}\leanok
\uses{def:AlgebraicGeometry.closedImmersionSourceOpenTarget}
The actual factor of the source open into the corresponding target open.

\noindent\emph{Lean:} \texttt{def AlgebraicGeometry.closedImmersionSourceOpenLift}, \href{https://github.com/Paul-Lez/HodgeConjecture/blob/main/HodgeConjecture/Definitions/AlgebraicGeometry/ClosedImmersionSourceOpen.lean#L36}{\texttt{HodgeConjecture/Definitions/AlgebraicGeometry/ClosedImmersionSourceOpen.lean:36}}.
\end{definition}

\section{The smooth-locus closed lift of an integral cycle component}\label{sec:HodgeConjecture.Definitions.AlgebraicGeometry.CycleComponentSmoothClosedLift}
\noindent\emph{File:} \href{https://github.com/Paul-Lez/HodgeConjecture/blob/main/HodgeConjecture/Definitions/AlgebraicGeometry/CycleComponentSmoothClosedLift.lean}{\texttt{HodgeConjecture/Definitions/AlgebraicGeometry/CycleComponentSmoothClosedLift.lean}}.

The ambient open is the complement of the first canonical singular-boundary support, and
the smooth locus embeds closed in it. Its constant relative dimension is \texttt{d-p}, by smooth
equidimensionality together with the coheight formula for a closed point of the integral
component.

\begin{definition}[AlgebraicGeometry.cycleComponentSmoothLocusAmbientOpen]\label{def:AlgebraicGeometry.cycleComponentSmoothLocusAmbientOpen}
\lean{AlgebraicGeometry.cycleComponentSmoothLocusAmbientOpen}\leanok
\uses{def:AlgebraicGeometry.cycleComponentSingularAmbientClosedFiltration}
The precise algebraic open complementary to the canonical singular boundary.

\noindent\emph{Lean:} \texttt{def AlgebraicGeometry.cycleComponentSmoothLocusAmbientOpen}, \href{https://github.com/Paul-Lez/HodgeConjecture/blob/main/HodgeConjecture/Definitions/AlgebraicGeometry/CycleComponentSmoothClosedLift.lean#L29}{\texttt{HodgeConjecture/Definitions/AlgebraicGeometry/CycleComponentSmoothClosedLift.lean:29}}.
\end{definition}

\begin{definition}[AlgebraicGeometry.cycleComponentSmoothLocusAmbientOpenOver]\label{def:AlgebraicGeometry.cycleComponentSmoothLocusAmbientOpenOver}
\lean{AlgebraicGeometry.cycleComponentSmoothLocusAmbientOpenOver}\leanok
\uses{def:AlgebraicGeometry.ComplexPoint.openScheme, def:AlgebraicGeometry.cycleComponentSmoothLocusAmbientOpen}
The ambient boundary complement with its induced structure map to \texttt{Spec ℂ}.

\noindent\emph{Lean:} \texttt{abbrev AlgebraicGeometry.cycleComponentSmoothLocusAmbientOpenOver}, \href{https://github.com/Paul-Lez/HodgeConjecture/blob/main/HodgeConjecture/Definitions/AlgebraicGeometry/CycleComponentSmoothClosedLift.lean#L33}{\texttt{HodgeConjecture/Definitions/AlgebraicGeometry/CycleComponentSmoothClosedLift.lean:33}}.
\end{definition}

\begin{definition}[AlgebraicGeometry.cycleComponentSmoothLocusClosedLift]\label{def:AlgebraicGeometry.cycleComponentSmoothLocusClosedLift}
\lean{AlgebraicGeometry.cycleComponentSmoothLocusClosedLift}\leanok
\uses{def:AlgebraicGeometry.closedImmersionSourceOpenLift, def:AlgebraicGeometry.cycleComponentSmoothLocusAmbientOpen}
The actual smooth locus, closed in the complement of its singular boundary.

\noindent\emph{Lean:} \texttt{def AlgebraicGeometry.cycleComponentSmoothLocusClosedLift}, \href{https://github.com/Paul-Lez/HodgeConjecture/blob/main/HodgeConjecture/Definitions/AlgebraicGeometry/CycleComponentSmoothClosedLift.lean#L48}{\texttt{HodgeConjecture/Definitions/AlgebraicGeometry/CycleComponentSmoothClosedLift.lean:48}}.
\end{definition}

\begin{definition}[AlgebraicGeometry.cycleComponentSmoothLocusClosedLiftOver]\label{def:AlgebraicGeometry.cycleComponentSmoothLocusClosedLiftOver}
\lean{AlgebraicGeometry.cycleComponentSmoothLocusClosedLiftOver}\leanok
\uses{def:AlgebraicGeometry.ComplexPoint.cycleComponentSmoothLocusOver, def:AlgebraicGeometry.cycleComponentSmoothLocusAmbientOpenOver, def:AlgebraicGeometry.cycleComponentSmoothLocusClosedLift}
The smooth-locus closed lift bundled over \texttt{Spec ℂ}.

\noindent\emph{Lean:} \texttt{def AlgebraicGeometry.cycleComponentSmoothLocusClosedLiftOver}, \href{https://github.com/Paul-Lez/HodgeConjecture/blob/main/HodgeConjecture/Definitions/AlgebraicGeometry/CycleComponentSmoothClosedLift.lean#L67}{\texttt{HodgeConjecture/Definitions/AlgebraicGeometry/CycleComponentSmoothClosedLift.lean:67}}.
\end{definition}

\begin{definition}[AlgebraicGeometry.ComplexPoint.cycleComponentSmoothClosedLiftAnalyticTopology]\label{def:AlgebraicGeometry.ComplexPoint.cycleComponentSmoothClosedLiftAnalyticTopology}
\lean{AlgebraicGeometry.ComplexPoint.cycleComponentSmoothClosedLiftAnalyticTopology}\leanok
\uses{def:AlgebraicGeometry.ComplexPoint, def:AlgebraicGeometry.Point.chartTopology}
\noindent\emph{Lean:} \texttt{def AlgebraicGeometry.ComplexPoint.cycleComponentSmoothClosedLiftAnalyticTopology}, \href{https://github.com/Paul-Lez/HodgeConjecture/blob/main/HodgeConjecture/Definitions/AlgebraicGeometry/CycleComponentSmoothClosedLift.lean#L99}{\texttt{HodgeConjecture/Definitions/AlgebraicGeometry/CycleComponentSmoothClosedLift.lean:99}}.
\end{definition}

\begin{definition}[AlgebraicGeometry.ComplexPoint.cycleComponentAnalyticClosedSupport]\label{def:AlgebraicGeometry.ComplexPoint.cycleComponentAnalyticClosedSupport}
\lean{AlgebraicGeometry.ComplexPoint.cycleComponentAnalyticClosedSupport}\leanok
\uses{def:AlgebraicGeometry.ComplexPoint.cycleComponentSmoothClosedLiftAnalyticTopology, def:AlgebraicGeometry.Point.evaluate, def:AlgebraicGeometry.Point.overOpen, def:AlgebraicGeometry.cycleComponentSupport}
The full cycle support as an actual closed analytic subset.

\noindent\emph{Lean:} \texttt{def AlgebraicGeometry.ComplexPoint.cycleComponentAnalyticClosedSupport}, \href{https://github.com/Paul-Lez/HodgeConjecture/blob/main/HodgeConjecture/Definitions/AlgebraicGeometry/CycleComponentSmoothClosedLift.lean#L102}{\texttt{HodgeConjecture/Definitions/AlgebraicGeometry/CycleComponentSmoothClosedLift.lean:102}}.
\end{definition}

\section{Supported singular models on smooth projective complex varieties}\label{sec:HodgeConjecture.Definitions.AlgebraicGeometry.ComplexSupportedSingularModel}
\noindent\emph{File:} \href{https://github.com/Paul-Lez/HodgeConjecture/blob/main/HodgeConjecture/Definitions/AlgebraicGeometry/ComplexSupportedSingularModel.lean}{\texttt{HodgeConjecture/Definitions/AlgebraicGeometry/ComplexSupportedSingularModel.lean}}.

The generic supported singular/injective comparison is specialized using the
constructed analytic contractible neighborhoods and hereditary paracompactness
of smooth projective analytifications. Only the geometric scheme hypotheses
remain: no purity, fundamental class, or comparison equivalence is an input.
The target is literally the ambient injective complex used by the derived
rational support comparison.

\begin{definition}[AlgebraicGeometry.ComplexPoint.complexSupportedSingularToAmbientInjective]\label{def:AlgebraicGeometry.ComplexPoint.complexSupportedSingularToAmbientInjective}
\lean{AlgebraicGeometry.ComplexPoint.complexSupportedSingularToAmbientInjective}\leanok
\uses{def:AlgebraicGeometry.ComplexPoint.ambientRationalInjectiveComplex, def:AlgebraicGeometry.ComplexPoint.localChart, def:AlgebraicTopology.Singular.supportedSingularToInjectiveComplex, def:TopologicalSpace.dim}
Its actual supported version for any open complement, not just a smooth support.

\noindent\emph{Lean:} \texttt{def AlgebraicGeometry.ComplexPoint.complexSupportedSingularToAmbientInjective}, \href{https://github.com/Paul-Lez/HodgeConjecture/blob/main/HodgeConjecture/Definitions/AlgebraicGeometry/ComplexSupportedSingularModel.lean#L45}{\texttt{HodgeConjecture/Definitions/AlgebraicGeometry/ComplexSupportedSingularModel.lean:45}}.
\end{definition}

\begin{definition}[AlgebraicGeometry.ComplexPoint.complexSupportedSingularInjectiveHomologyIso]\label{def:AlgebraicGeometry.ComplexPoint.complexSupportedSingularInjectiveHomologyIso}
\lean{AlgebraicGeometry.ComplexPoint.complexSupportedSingularInjectiveHomologyIso}\leanok
\uses{def:AlgebraicGeometry.ComplexPoint.affineSpaceHomeomorph, def:AlgebraicGeometry.ComplexPoint.ambientRationalInjectiveComplex, def:AlgebraicGeometry.ComplexPoint.localChart, def:AlgebraicGeometry.ComplexProjectiveSpace.awayHomogeneousEvaluation, def:AlgebraicGeometry.ComplexProjectiveSpace.chartAffineComplexPointMap, def:AlgebraicGeometry.ComplexProjectiveSpace.normalize, def:AlgebraicGeometry.ComplexProjectiveSpace.projectiveSpaceBasicOpen, def:AlgebraicGeometry.ComplexProjectiveSpace.projectivizationToComplexPoint, def:AlgebraicGeometry.ComplexProjectiveSpace.sphereToProjectivization, def:AlgebraicGeometry.ComplexProjectiveSpace.toVectorChart, def:AlgebraicGeometry.ComplexProjectiveSpace.vectorChartToAffinePoint, def:AlgebraicGeometry.ComplexProjectiveSpace.vectorOfAwayRingHom, def:AlgebraicGeometry.IsClosedImmersion.residueFieldIso, def:AlgebraicGeometry.IsProjective, def:AlgebraicGeometry.Point.residueData, def:AlgebraicGeometry.ProjectiveSpace.Presentation.analyticImmersion, def:AlgebraicTopology.Singular.supportedSingularInjectiveHomologyIso, def:ProjectiveSpace.degreeZeroEquiv, def:TopologicalSpace.dim}
Local supported singular cohomology and the literal supported injective
model are canonically isomorphic in every integer degree.

\noindent\emph{Lean:} \texttt{def AlgebraicGeometry.ComplexPoint.complexSupportedSingularInjectiveHomologyIso}, \href{https://github.com/Paul-Lez/HodgeConjecture/blob/main/HodgeConjecture/Definitions/AlgebraicGeometry/ComplexSupportedSingularModel.lean#L67}{\texttt{HodgeConjecture/Definitions/AlgebraicGeometry/ComplexSupportedSingularModel.lean:67}}.
\end{definition}

\section{Cohomology-sheaf concentration for smooth closed supports}\label{sec:HodgeConjecture.Definitions.AlgebraicGeometry.SmoothClosedSupportCohomologySheaf}
\noindent\emph{File:} \href{https://github.com/Paul-Lez/HodgeConjecture/blob/main/HodgeConjecture/Definitions/AlgebraicGeometry/SmoothClosedSupportCohomologySheaf.lean}{\texttt{HodgeConjecture/Definitions/AlgebraicGeometry/SmoothClosedSupportCohomologySheaf.lean}}.

The complex is the supported-sections kernel applied to the fixed ambient rational injective
resolution. Its open-section homology is compared to relative singular cohomology by the
singular resolution and restriction-cone maps, and cofinal normal neighborhoods then give
stalkwise and sheafwise concentration in degree twice the complex codimension. Identifying
the surviving cohomology sheaf with rational constants on the support is left to a later file.

\begin{definition}[AlgebraicGeometry.ComplexPoint.smoothClosedSupportCohomologySheafAnalyticTopology]\label{def:AlgebraicGeometry.ComplexPoint.smoothClosedSupportCohomologySheafAnalyticTopology}
\lean{AlgebraicGeometry.ComplexPoint.smoothClosedSupportCohomologySheafAnalyticTopology}\leanok
\uses{def:AlgebraicGeometry.ComplexPoint, def:AlgebraicGeometry.Point.chartTopology}
\noindent\emph{Lean:} \texttt{def AlgebraicGeometry.ComplexPoint.smoothClosedSupportCohomologySheafAnalyticTopology}, \href{https://github.com/Paul-Lez/HodgeConjecture/blob/main/HodgeConjecture/Definitions/AlgebraicGeometry/SmoothClosedSupportCohomologySheaf.lean#L32}{\texttt{HodgeConjecture/Definitions/AlgebraicGeometry/SmoothClosedSupportCohomologySheaf.lean:32}}.
\end{definition}

\begin{definition}[AlgebraicGeometry.ComplexPoint.complexSupportInjectiveComplex]\label{def:AlgebraicGeometry.ComplexPoint.complexSupportInjectiveComplex}
\lean{AlgebraicGeometry.ComplexPoint.complexSupportInjectiveComplex}\leanok
\uses{def:AlgebraicGeometry.ComplexPoint.ambientRationalInjectiveComplex, def:AlgebraicGeometry.ComplexPoint.smoothClosedSupportCohomologySheafAnalyticTopology, def:TopCat.Sheaf.sheafSectionsSupportedOutside}
The literal supported ambient rational injective complex for a closed support.

\noindent\emph{Lean:} \texttt{def AlgebraicGeometry.ComplexPoint.complexSupportInjectiveComplex}, \href{https://github.com/Paul-Lez/HodgeConjecture/blob/main/HodgeConjecture/Definitions/AlgebraicGeometry/SmoothClosedSupportCohomologySheaf.lean#L35}{\texttt{HodgeConjecture/Definitions/AlgebraicGeometry/SmoothClosedSupportCohomologySheaf.lean:35}}.
\end{definition}

\begin{definition}[AlgebraicGeometry.ComplexPoint.complexSupportInjectiveSectionCohomologyEquiv]\label{def:AlgebraicGeometry.ComplexPoint.complexSupportInjectiveSectionCohomologyEquiv}
\lean{AlgebraicGeometry.ComplexPoint.complexSupportInjectiveSectionCohomologyEquiv}\leanok
\uses{def:AlgebraicGeometry.ComplexPoint.complexSupportInjectiveComplex, def:AlgebraicGeometry.ComplexPoint.complexSupportedSingularInjectiveHomologyIso, def:AlgebraicTopology.Singular.supportedRationalSingularSectionCohomologyEquivSupportComplement}
Actual open-section cohomology of the supported injective model is relative
singular cohomology of the same literal local support pair.

\noindent\emph{Lean:} \texttt{def AlgebraicGeometry.ComplexPoint.complexSupportInjectiveSectionCohomologyEquiv}, \href{https://github.com/Paul-Lez/HodgeConjecture/blob/main/HodgeConjecture/Definitions/AlgebraicGeometry/SmoothClosedSupportCohomologySheaf.lean#L47}{\texttt{HodgeConjecture/Definitions/AlgebraicGeometry/SmoothClosedSupportCohomologySheaf.lean:47}}.
\end{definition}

\section{Purity along the smooth locus of an integral cycle component}\label{sec:HodgeConjecture.Definitions.AlgebraicGeometry.CycleComponentSmoothSupportPurity}
\noindent\emph{File:} \href{https://github.com/Paul-Lez/HodgeConjecture/blob/main/HodgeConjecture/Definitions/AlgebraicGeometry/CycleComponentSmoothSupportPurity.lean}{\texttt{HodgeConjecture/Definitions/AlgebraicGeometry/CycleComponentSmoothSupportPurity.lean}}.

The coefficient complex is the ambient injective resolution, with support in the full cycle
component, restricted to the complement of the canonical singular boundary. Normal
neighborhoods of the smooth-locus closed lift show its cohomology sheaves are concentrated
in degree \texttt{2p}.

\begin{definition}[AlgebraicGeometry.ComplexPoint.cycleComponentSmoothSupportPurityAnalyticTopology]\label{def:AlgebraicGeometry.ComplexPoint.cycleComponentSmoothSupportPurityAnalyticTopology}
\lean{AlgebraicGeometry.ComplexPoint.cycleComponentSmoothSupportPurityAnalyticTopology}\leanok
\uses{def:AlgebraicGeometry.ComplexPoint, def:AlgebraicGeometry.Point.chartTopology}
\noindent\emph{Lean:} \texttt{def AlgebraicGeometry.ComplexPoint.cycleComponentSmoothSupportPurityAnalyticTopology}, \href{https://github.com/Paul-Lez/HodgeConjecture/blob/main/HodgeConjecture/Definitions/AlgebraicGeometry/CycleComponentSmoothSupportPurity.lean#L34}{\texttt{HodgeConjecture/Definitions/AlgebraicGeometry/CycleComponentSmoothSupportPurity.lean:34}}.
\end{definition}

\begin{definition}[AlgebraicGeometry.ComplexPoint.cycleComponentSmoothSupportAmbientOpen]\label{def:AlgebraicGeometry.ComplexPoint.cycleComponentSmoothSupportAmbientOpen}
\lean{AlgebraicGeometry.ComplexPoint.cycleComponentSmoothSupportAmbientOpen}\leanok
\uses{def:AlgebraicGeometry.ComplexPoint.cycleComponentSingularAnalyticClosedFiltration, def:AlgebraicGeometry.ComplexPoint.cycleComponentSmoothSupportPurityAnalyticTopology}
The exact analytic complement of the canonical singular boundary.

\noindent\emph{Lean:} \texttt{abbrev AlgebraicGeometry.ComplexPoint.cycleComponentSmoothSupportAmbientOpen}, \href{https://github.com/Paul-Lez/HodgeConjecture/blob/main/HodgeConjecture/Definitions/AlgebraicGeometry/CycleComponentSmoothSupportPurity.lean#L37}{\texttt{HodgeConjecture/Definitions/AlgebraicGeometry/CycleComponentSmoothSupportPurity.lean:37}}.
\end{definition}

\begin{definition}[AlgebraicGeometry.ComplexPoint.cycleComponentSmoothRestrictedInjectiveComplex]\label{def:AlgebraicGeometry.ComplexPoint.cycleComponentSmoothRestrictedInjectiveComplex}
\lean{AlgebraicGeometry.ComplexPoint.cycleComponentSmoothRestrictedInjectiveComplex}\leanok
\uses{def:AlgebraicGeometry.ComplexPoint.complexSupportInjectiveComplex, def:AlgebraicGeometry.ComplexPoint.cycleComponentAnalyticClosedSupport, def:AlgebraicGeometry.ComplexPoint.cycleComponentSmoothSupportAmbientOpen}
Restriction of the original full-support injective model to the boundary complement.

\noindent\emph{Lean:} \texttt{def AlgebraicGeometry.ComplexPoint.cycleComponentSmoothRestrictedInjectiveComplex}, \href{https://github.com/Paul-Lez/HodgeConjecture/blob/main/HodgeConjecture/Definitions/AlgebraicGeometry/CycleComponentSmoothSupportPurity.lean#L117}{\texttt{HodgeConjecture/Definitions/AlgebraicGeometry/CycleComponentSmoothSupportPurity.lean:117}}.
\end{definition}

\begin{definition}[AlgebraicGeometry.ComplexPoint.cycleComponentSmoothSupportLowestSectionCohomologyIso]\label{def:AlgebraicGeometry.ComplexPoint.cycleComponentSmoothSupportLowestSectionCohomologyIso}
\lean{AlgebraicGeometry.ComplexPoint.cycleComponentSmoothSupportLowestSectionCohomologyIso}\leanok
\uses{def:AlgebraicGeometry.ComplexPoint.complexSupportInjectiveSectionCohomologyEquiv, def:AlgebraicGeometry.ComplexPoint.cycleComponentSmoothRestrictedInjectiveComplex, def:AlgebraicGeometry.ComplexPoint.smoothClosedSupportRelativeHomologyIso, def:AlgebraicGeometry.cycleComponentOrderIsoIic, def:AlgebraicGeometry.cycleComponentSmoothLocusClosedLiftOver, def:AlgebraicTopology.Singular.neighborhoodSupportPairImageCohomologyEquiv, def:AlgebraicTopology.Singular.rationalSingularHomologyIsoOfHomotopyEquiv, def:AlgebraicTopology.Singular.relativeCohomologyEquivDualHomology, def:AlgebraicTopology.Singular.relativeHomologyProjection, def:AlgebraicTopology.Singular.relativeSingularBoundary, def:AlgebraicTopology.Singular.standardAffineBoundaryChainHomotopyEquiv, def:AlgebraicTopology.Singular.standardComplexRealPairIso, def:AlgebraicTopology.Singular.standardSphereSuccNormalizedChainsXIsoStandard, def:PolynomialCatenary.Triangular.T, def:SSet.chainComplexXZeroIsoCyclesZero, def:TopCat.Sheaf.openRestrictedLowestSectionCohomologyIso, def:TopCat.Sheaf.sectionCohomologyPresheafStalkIso, def:TopCat.Sheaf.sheafSectionsSupportedOutsideOnOpenIso, def:TopCat.Sheaf.supportedOutsideGlueZero}
The canonical lowest-degree isomorphism using the original ambient resolution and
the original ambient cohomology sheaf, both evaluated on the boundary complement.

\noindent\emph{Lean:} \texttt{def AlgebraicGeometry.ComplexPoint.cycleComponentSmoothSupportLowestSectionCohomologyIso}, \href{https://github.com/Paul-Lez/HodgeConjecture/blob/main/HodgeConjecture/Definitions/AlgebraicGeometry/CycleComponentSmoothSupportPurity.lean#L142}{\texttt{HodgeConjecture/Definitions/AlgebraicGeometry/CycleComponentSmoothSupportPurity.lean:142}}.
\end{definition}

\section{The global normalized smooth-support coclass section}\label{sec:HodgeConjecture.Definitions.AlgebraicGeometry.SmoothClosedSupportCoclassSection}
\noindent\emph{File:} \href{https://github.com/Paul-Lez/HodgeConjecture/blob/main/HodgeConjecture/Definitions/AlgebraicGeometry/SmoothClosedSupportCoclassSection.lean}{\texttt{HodgeConjecture/Definitions/AlgebraicGeometry/SmoothClosedSupportCoclassSection.lean}}.

Holomorphic normal charts supply local relative coclasses whose ambient overlap agreement
identifies their sheaf germs, and those germs vanish off the closed image because the
relative complexes do. Unique sheaf gluing then produces the global section, with its
complex normalization.

\begin{definition}[AlgebraicGeometry.ComplexPoint.smoothClosedSupportChartOpen]\label{def:AlgebraicGeometry.ComplexPoint.smoothClosedSupportChartOpen}
\lean{AlgebraicGeometry.ComplexPoint.smoothClosedSupportChartOpen}\leanok
\uses{def:AlgebraicGeometry.ComplexPoint.closedImmersionHolomorphicFlatteningChart}
The actual source open of a constructed holomorphic normal chart.

\noindent\emph{Lean:} \texttt{def AlgebraicGeometry.ComplexPoint.smoothClosedSupportChartOpen}, \href{https://github.com/Paul-Lez/HodgeConjecture/blob/main/HodgeConjecture/Definitions/AlgebraicGeometry/SmoothClosedSupportCoclassSection.lean#L31}{\texttt{HodgeConjecture/Definitions/AlgebraicGeometry/SmoothClosedSupportCoclassSection.lean:31}}.
\end{definition}

\begin{definition}[AlgebraicGeometry.ComplexPoint.smoothClosedSupportCoclassSheaf]\label{def:AlgebraicGeometry.ComplexPoint.smoothClosedSupportCoclassSheaf}
\lean{AlgebraicGeometry.ComplexPoint.smoothClosedSupportCoclassSheaf}\leanok
\uses{def:AlgebraicGeometry.ComplexPoint, def:AlgebraicGeometry.Point.chartTopology, def:AlgebraicGeometry.Point.map, def:AlgebraicTopology.Singular.supportRelativeCohomologySheaf}
The target is the sheafification of the literal relative-cohomology presheaf.

\noindent\emph{Lean:} \texttt{abbrev AlgebraicGeometry.ComplexPoint.smoothClosedSupportCoclassSheaf}, \href{https://github.com/Paul-Lez/HodgeConjecture/blob/main/HodgeConjecture/Definitions/AlgebraicGeometry/SmoothClosedSupportCoclassSection.lean#L41}{\texttt{HodgeConjecture/Definitions/AlgebraicGeometry/SmoothClosedSupportCoclassSection.lean:41}}.
\end{definition}

\begin{definition}[AlgebraicGeometry.ComplexPoint.smoothClosedSupportChartSheafSection]\label{def:AlgebraicGeometry.ComplexPoint.smoothClosedSupportChartSheafSection}
\lean{AlgebraicGeometry.ComplexPoint.smoothClosedSupportChartSheafSection}\leanok
\uses{def:AlgebraicGeometry.ComplexPoint.smoothClosedSupportChartCoclass, def:AlgebraicGeometry.ComplexPoint.smoothClosedSupportChartOpen, def:AlgebraicGeometry.ComplexPoint.smoothClosedSupportCoclassSheaf, def:AlgebraicTopology.Singular.supportRelativeCohomologyToSheaf}
The exact normal coclass determines an actual section on its full chart source.

\noindent\emph{Lean:} \texttt{def AlgebraicGeometry.ComplexPoint.smoothClosedSupportChartSheafSection}, \href{https://github.com/Paul-Lez/HodgeConjecture/blob/main/HodgeConjecture/Definitions/AlgebraicGeometry/SmoothClosedSupportCoclassSection.lean#L47}{\texttt{HodgeConjecture/Definitions/AlgebraicGeometry/SmoothClosedSupportCoclassSection.lean:47}}.
\end{definition}

\begin{definition}[AlgebraicGeometry.ComplexPoint.smoothClosedSupportChartCoclassGerm]\label{def:AlgebraicGeometry.ComplexPoint.smoothClosedSupportChartCoclassGerm}
\lean{AlgebraicGeometry.ComplexPoint.smoothClosedSupportChartCoclassGerm}\leanok
\uses{def:AlgebraicGeometry.ComplexPoint.smoothClosedSupportChartCoclass, def:AlgebraicGeometry.ComplexPoint.smoothClosedSupportChartOpen, def:AlgebraicGeometry.ComplexPoint.smoothClosedSupportCoclassSheaf, def:AlgebraicTopology.Singular.supportRelativeCohomologyGerm}
Its germ is, by definition, the germ of the fixed normal-projection coclass.

\noindent\emph{Lean:} \texttt{def AlgebraicGeometry.ComplexPoint.smoothClosedSupportChartCoclassGerm}, \href{https://github.com/Paul-Lez/HodgeConjecture/blob/main/HodgeConjecture/Definitions/AlgebraicGeometry/SmoothClosedSupportCoclassSection.lean#L56}{\texttt{HodgeConjecture/Definitions/AlgebraicGeometry/SmoothClosedSupportCoclassSection.lean:56}}.
\end{definition}

\begin{definition}[AlgebraicGeometry.ComplexPoint.smoothClosedSupportCoclassStalk]\label{def:AlgebraicGeometry.ComplexPoint.smoothClosedSupportCoclassStalk}
\lean{AlgebraicGeometry.ComplexPoint.smoothClosedSupportCoclassStalk}\leanok
\uses{def:AlgebraicGeometry.ComplexPoint.smoothClosedSupportChartCoclassGerm}
A pointwise normalized germ family. Choice selects a preimage point only;
the proved chart-overlap theorem below proves independence of that selection.

\noindent\emph{Lean:} \texttt{def AlgebraicGeometry.ComplexPoint.smoothClosedSupportCoclassStalk}, \href{https://github.com/Paul-Lez/HodgeConjecture/blob/main/HodgeConjecture/Definitions/AlgebraicGeometry/SmoothClosedSupportCoclassSection.lean#L99}{\texttt{HodgeConjecture/Definitions/AlgebraicGeometry/SmoothClosedSupportCoclassSection.lean:99}}.
\end{definition}

\begin{definition}[AlgebraicGeometry.ComplexPoint.smoothClosedSupportCoclassSection]\label{def:AlgebraicGeometry.ComplexPoint.smoothClosedSupportCoclassSection}
\lean{AlgebraicGeometry.ComplexPoint.smoothClosedSupportCoclassSection}\leanok
\uses{def:AlgebraicGeometry.ComplexPoint.closedImmersionNormalTransition, def:AlgebraicGeometry.ComplexPoint.smoothClosedSupportChartSheafSection, def:AlgebraicGeometry.ComplexPoint.smoothClosedSupportCoclassStalk, def:AlgebraicGeometry.IsClosedImmersion.residueFieldIso, def:AlgebraicGeometry.Point.residueData, def:AlgebraicTopology.Singular.ambientPairLeftToPointComplementChains, def:AlgebraicTopology.Singular.centeredComplexEmbeddingPairHomotopy, def:AlgebraicTopology.Singular.chartNormalFiberPair, def:AlgebraicTopology.Singular.complexMatrixOfContinuousLinearMap, def:AlgebraicTopology.Singular.complexMatrixPuncturedPairHomotopy, def:AlgebraicTopology.Singular.complexStraightLineLocalPairHomotopy, def:AlgebraicTopology.Singular.complexStraightLineNeighborhoodPairHomotopy, def:AlgebraicTopology.Singular.flattenedSupportNormalClass, def:AlgebraicTopology.Singular.neighborhoodAmbientChainMap, def:AlgebraicTopology.Singular.neighborhoodRelativePointExcisionSmallIso, def:AlgebraicTopology.Singular.normalTransitionDerivativeEquiv, def:AlgebraicTopology.Singular.normalTransitionDomain, def:AlgebraicTopology.Singular.normalTransitionMap, def:AlgebraicTopology.Singular.pointComplementAmbientChainMap, def:AlgebraicTopology.Singular.pointExcisionSmallToAmbientCokernelSequenceHom, def:AlgebraicTopology.Singular.relativeHomologyMap, def:AlgebraicTopology.Singular.standardAffineBoundaryChainHomotopyEquiv, def:AlgebraicTopology.Singular.standardLocalRelativeHomologyZeroIso, def:AlgebraicTopology.Singular.standardSphereSuccSimplicialHomologyTopIsoRat, def:CategoryTheory.ShortComplex.linearDualMap, def:CategoryTheory.ShortComplex.moduleCatCycleMap, def:CategoryTheory.ShortComplex.moduleCatHomologyClass, def:Matrix.ComplexIsotopyToOne.diagonal, def:Matrix.ComplexIsotopyToOne.mul, def:Matrix.ComplexIsotopyToOne.transvection, def:TopCat.Sheaf.sectionOfLocallyRepresentable}
The actual unique global gluing of exactly normalized smooth normal coclasses.

\noindent\emph{Lean:} \texttt{def AlgebraicGeometry.ComplexPoint.smoothClosedSupportCoclassSection}, \href{https://github.com/Paul-Lez/HodgeConjecture/blob/main/HodgeConjecture/Definitions/AlgebraicGeometry/SmoothClosedSupportCoclassSection.lean#L151}{\texttt{HodgeConjecture/Definitions/AlgebraicGeometry/SmoothClosedSupportCoclassSection.lean:151}}.
\end{definition}

\section{The normalized component coclass on the original ambient smooth-support open}\label{sec:HodgeConjecture.Definitions.AlgebraicGeometry.CycleComponentSmoothSupportCoclassSection}
\noindent\emph{File:} \href{https://github.com/Paul-Lez/HodgeConjecture/blob/main/HodgeConjecture/Definitions/AlgebraicGeometry/CycleComponentSmoothSupportCoclassSection.lean}{\texttt{HodgeConjecture/Definitions/AlgebraicGeometry/CycleComponentSmoothSupportCoclassSection.lean}}.

The component's actual smooth locus is closed in the complement of its singular
boundary. We construct its exactly normalized smooth-support section there and
transport it through the actual analytic open embedding. The result is a section
of the ORIGINAL ambient relative-cohomology sheaf on the singular-boundary
complement. No section, purity comparison, or orientation coherence is an input.

\begin{definition}[AlgebraicGeometry.ComplexPoint.cycleComponentSmoothClosedLiftStructureMap]\label{def:AlgebraicGeometry.ComplexPoint.cycleComponentSmoothClosedLiftStructureMap}
\lean{AlgebraicGeometry.ComplexPoint.cycleComponentSmoothClosedLiftStructureMap}\leanok
\uses{def:AlgebraicGeometry.ComplexPoint.cycleComponentSmoothLocusOver}
The actual structure morphism of the closed-lift source.

\noindent\emph{Lean:} \texttt{abbrev AlgebraicGeometry.ComplexPoint.cycleComponentSmoothClosedLiftStructureMap}, \href{https://github.com/Paul-Lez/HodgeConjecture/blob/main/HodgeConjecture/Definitions/AlgebraicGeometry/CycleComponentSmoothSupportCoclassSection.lean#L52}{\texttt{HodgeConjecture/Definitions/AlgebraicGeometry/CycleComponentSmoothSupportCoclassSection.lean:52}}.
\end{definition}

\begin{definition}[AlgebraicGeometry.ComplexPoint.cycleComponentSmoothClosedLiftCoclassSection]\label{def:AlgebraicGeometry.ComplexPoint.cycleComponentSmoothClosedLiftCoclassSection}
\lean{AlgebraicGeometry.ComplexPoint.cycleComponentSmoothClosedLiftCoclassSection}\leanok
\uses{def:AlgebraicGeometry.ComplexPoint.cycleComponentSmoothClosedLiftStructureMap, def:AlgebraicGeometry.ComplexPoint.smoothClosedSupportCoclassSection, def:AlgebraicGeometry.cycleComponentOrderIsoIic, def:AlgebraicGeometry.cycleComponentSmoothLocusClosedLiftOver, def:PolynomialCatenary.Triangular.T}
The normalized section in the auxiliary algebraic ambient open, in the proved
degree 2p. The class is the general normal-chart gluing, not supplied data.

\noindent\emph{Lean:} \texttt{def AlgebraicGeometry.ComplexPoint.cycleComponentSmoothClosedLiftCoclassSection}, \href{https://github.com/Paul-Lez/HodgeConjecture/blob/main/HodgeConjecture/Definitions/AlgebraicGeometry/CycleComponentSmoothSupportCoclassSection.lean#L72}{\texttt{HodgeConjecture/Definitions/AlgebraicGeometry/CycleComponentSmoothSupportCoclassSection.lean:72}}.
\end{definition}

\begin{definition}[AlgebraicGeometry.ComplexPoint.cycleComponentSmoothClosedLiftAmbientMap]\label{def:AlgebraicGeometry.ComplexPoint.cycleComponentSmoothClosedLiftAmbientMap}
\lean{AlgebraicGeometry.ComplexPoint.cycleComponentSmoothClosedLiftAmbientMap}\leanok
\uses{def:AlgebraicGeometry.ComplexPoint, def:AlgebraicGeometry.ComplexPoint.openInclusion, def:AlgebraicGeometry.Point.continuousMap, def:AlgebraicGeometry.cycleComponentSmoothLocusAmbientOpenOver}
The actual analytic open-embedding map back to the original ambient space.

\noindent\emph{Lean:} \texttt{def AlgebraicGeometry.ComplexPoint.cycleComponentSmoothClosedLiftAmbientMap}, \href{https://github.com/Paul-Lez/HodgeConjecture/blob/main/HodgeConjecture/Definitions/AlgebraicGeometry/CycleComponentSmoothSupportCoclassSection.lean#L86}{\texttt{HodgeConjecture/Definitions/AlgebraicGeometry/CycleComponentSmoothSupportCoclassSection.lean:86}}.
\end{definition}

\begin{definition}[AlgebraicGeometry.ComplexPoint.cycleComponentSmoothSupportCoclassSection]\label{def:AlgebraicGeometry.ComplexPoint.cycleComponentSmoothSupportCoclassSection}
\lean{AlgebraicGeometry.ComplexPoint.cycleComponentSmoothSupportCoclassSection}\leanok
\uses{def:AlgebraicGeometry.ComplexPoint.cycleComponentSmoothClosedLiftAmbientMap, def:AlgebraicGeometry.ComplexPoint.cycleComponentSmoothClosedLiftCoclassSection, def:AlgebraicGeometry.ComplexPoint.cycleComponentSmoothSupportAmbientOpen, def:AlgebraicGeometry.cycleComponentSupport, def:AlgebraicTopology.Singular.supportRelativeCohomologySectionOnOpen}
The actual normalized component coclass section, living on the singular-boundary
complement in the ORIGINAL ambient relative-cohomology sheaf.

\noindent\emph{Lean:} \texttt{def AlgebraicGeometry.ComplexPoint.cycleComponentSmoothSupportCoclassSection}, \href{https://github.com/Paul-Lez/HodgeConjecture/blob/main/HodgeConjecture/Definitions/AlgebraicGeometry/CycleComponentSmoothSupportCoclassSection.lean#L114}{\texttt{HodgeConjecture/Definitions/AlgebraicGeometry/CycleComponentSmoothSupportCoclassSection.lean:114}}.
\end{definition}

\section{Support-forgetting in the actual rational injective model}\label{sec:HodgeConjecture.Definitions.AlgebraicGeometry.DerivedSupportRationalConeForget}
\noindent\emph{File:} \href{https://github.com/Paul-Lez/HodgeConjecture/blob/main/HodgeConjecture/Definitions/AlgebraicGeometry/DerivedSupportRationalConeForget.lean}{\texttt{HodgeConjecture/Definitions/AlgebraicGeometry/DerivedSupportRationalConeForget.lean}}.

\begin{definition}[AlgebraicGeometry.ComplexPoint.rationalCohomologyAddEquivAmbientInjectiveHomology]\label{def:AlgebraicGeometry.ComplexPoint.rationalCohomologyAddEquivAmbientInjectiveHomology}
\lean{AlgebraicGeometry.ComplexPoint.rationalCohomologyAddEquivAmbientInjectiveHomology}\leanok
\uses{def:AlgebraicGeometry.ComplexPoint.ambientRationalInjectiveAugmentation, def:AlgebraicGeometry.ComplexPoint.hypercohomologyAddEquivGlobalSectionsKInjective, def:AlgebraicGeometry.ComplexPoint.hypercohomologyMap}
Ordinary rational cohomology computed by the actual ambient rational
injective resolution. This has the ordinary augmentation normalization.

\noindent\emph{Lean:} \texttt{def AlgebraicGeometry.ComplexPoint.rationalCohomologyAddEquivAmbientInjectiveHomology}, \href{https://github.com/Paul-Lez/HodgeConjecture/blob/main/HodgeConjecture/Definitions/AlgebraicGeometry/DerivedSupportRationalConeForget.lean#L29}{\texttt{HodgeConjecture/Definitions/AlgebraicGeometry/DerivedSupportRationalConeForget.lean:29}}.
\end{definition}

\section{Unique extension across a cycle component's singular boundary}\label{sec:HodgeConjecture.Definitions.AlgebraicGeometry.CycleComponentSupportExtension}
\noindent\emph{File:} \href{https://github.com/Paul-Lez/HodgeConjecture/blob/main/HodgeConjecture/Definitions/AlgebraicGeometry/CycleComponentSupportExtension.lean}{\texttt{HodgeConjecture/Definitions/AlgebraicGeometry/CycleComponentSupportExtension.lean}}.

The canonical finite smooth filtration and its layerwise vanishing give the singular
boundary zero supported cohomology below \texttt{2(p+1)}, in particular in degrees \texttt{2p} and
\texttt{2p+1}. The nested-support localization sequence then makes restriction from the full
component support to its smooth-locus ambient open an isomorphism in degree \texttt{2p}, whose
inverse is the unique extension operation.

\begin{definition}[AlgebraicGeometry.ComplexPoint.cycleComponentSupportExtensionAnalyticTopology]\label{def:AlgebraicGeometry.ComplexPoint.cycleComponentSupportExtensionAnalyticTopology}
\lean{AlgebraicGeometry.ComplexPoint.cycleComponentSupportExtensionAnalyticTopology}\leanok
\uses{def:AlgebraicGeometry.ComplexPoint, def:AlgebraicGeometry.Point.chartTopology}
\noindent\emph{Lean:} \texttt{def AlgebraicGeometry.ComplexPoint.cycleComponentSupportExtensionAnalyticTopology}, \href{https://github.com/Paul-Lez/HodgeConjecture/blob/main/HodgeConjecture/Definitions/AlgebraicGeometry/CycleComponentSupportExtension.lean#L31}{\texttt{HodgeConjecture/Definitions/AlgebraicGeometry/CycleComponentSupportExtension.lean:31}}.
\end{definition}

\begin{definition}[AlgebraicGeometry.ComplexPoint.cycleComponentSupportSectionRestriction]\label{def:AlgebraicGeometry.ComplexPoint.cycleComponentSupportSectionRestriction}
\lean{AlgebraicGeometry.ComplexPoint.cycleComponentSupportSectionRestriction}\leanok
\uses{def:AlgebraicGeometry.ComplexPoint.complexSupportInjectiveComplex, def:AlgebraicGeometry.ComplexPoint.cycleComponentAnalyticClosedSupport, def:AlgebraicGeometry.ComplexPoint.cycleComponentSmoothSupportAmbientOpen, def:AlgebraicGeometry.ComplexPoint.cycleComponentSupportExtensionAnalyticTopology, def:TopCat.Sheaf.sectionComplexRestriction}
Restriction of the original supported injective section complex to the
actual smooth-locus ambient open.

\noindent\emph{Lean:} \texttt{def AlgebraicGeometry.ComplexPoint.cycleComponentSupportSectionRestriction}, \href{https://github.com/Paul-Lez/HodgeConjecture/blob/main/HodgeConjecture/Definitions/AlgebraicGeometry/CycleComponentSupportExtension.lean#L98}{\texttt{HodgeConjecture/Definitions/AlgebraicGeometry/CycleComponentSupportExtension.lean:98}}.
\end{definition}

\begin{definition}[AlgebraicGeometry.ComplexPoint.cycleComponentSupportExtensionIso]\label{def:AlgebraicGeometry.ComplexPoint.cycleComponentSupportExtensionIso}
\lean{AlgebraicGeometry.ComplexPoint.cycleComponentSupportExtensionIso}\leanok
\uses{def:AlgebraicGeometry.ComplexPoint.complexSupportInjectiveSectionCohomologyEquiv, def:AlgebraicGeometry.ComplexPoint.cycleComponentSupportSectionRestriction, def:AlgebraicGeometry.ComplexPoint.singularFiltrationLocalSupportVanishingAnalyticTopology, def:AlgebraicGeometry.ComplexPoint.smoothClosedSupportRelativeHomologyIso, def:AlgebraicGeometry.closedImmersionSourceOpenLift, def:AlgebraicGeometry.cycleComponentOrderIsoIic, def:AlgebraicGeometry.cycleComponentSingularFiltrationLength, def:AlgebraicGeometry.cycleComponentSingularFiltrationStratumOverι, def:AlgebraicGeometry.cycleComponentSingularStratumClosedLiftOver, def:AlgebraicGeometry.reducedSingularLocus, def:AlgebraicTopology.Singular.neighborhoodSupportPairImageCohomologyEquiv, def:AlgebraicTopology.Singular.rationalSingularHomologyIsoOfHomotopyEquiv, def:AlgebraicTopology.Singular.relativeCohomologyEquivDualHomology, def:AlgebraicTopology.Singular.relativeHomologyProjection, def:AlgebraicTopology.Singular.relativeSingularBoundary, def:AlgebraicTopology.Singular.standardAffineBoundaryChainHomotopyEquiv, def:AlgebraicTopology.Singular.standardComplexRealPairIso, def:AlgebraicTopology.Singular.standardSphereSuccNormalizedChainsXIsoStandard, def:CategoryTheory.kernelCompositionFactorizationIso, def:CategoryTheory.kernelFactorizationShortComplexMapIso, def:SSet.chainComplexXZeroIsoCyclesZero, def:TopCat.Sheaf.lowestSectionCohomologyIso, def:TopCat.Sheaf.nestedSupportRestrictionLastComplexIso, def:TopCat.Sheaf.nestedSupportRestrictionShortComplex, def:TopCat.Sheaf.sectionCohomologyPresheafStalkIso, def:TopCat.Sheaf.supportedOutsideGlueZero}
The extension equivalence is the actual restriction map with its proved
inverse. It has no boundary-vanishing or fundamental-class input.

\noindent\emph{Lean:} \texttt{def AlgebraicGeometry.ComplexPoint.cycleComponentSupportExtensionIso}, \href{https://github.com/Paul-Lez/HodgeConjecture/blob/main/HodgeConjecture/Definitions/AlgebraicGeometry/CycleComponentSupportExtension.lean#L134}{\texttt{HodgeConjecture/Definitions/AlgebraicGeometry/CycleComponentSupportExtension.lean:134}}.
\end{definition}

\section{Constructed sheaf cycle classes in arbitrary codimension}\label{sec:HodgeConjecture.Definitions.AlgebraicGeometry.CycleComponentSheafClass}
\noindent\emph{File:} \href{https://github.com/Paul-Lez/HodgeConjecture/blob/main/HodgeConjecture/Definitions/AlgebraicGeometry/CycleComponentSheafClass.lean}{\texttt{HodgeConjecture/Definitions/AlgebraicGeometry/CycleComponentSheafClass.lean}}.

The exactly normalized normal-chart coclass on a component's smooth locus is
transported to the actual supported cohomology sheaf. Lowest-degree purity
and the proved unique extension across the singular boundary then give an
actual supported class on the original ambient variety. Forgetting support
lands in the repository's ordinary rational cohomology.

All comparison maps, purity statements, and extension isomorphisms are
constructed. No fundamental-class, orientation, duality, or vanishing datum
is an argument. The corresponding Borel–Moore fundamental class is obtained
through the previously constructed complex-orientation duality for the actual
ambient chain sheaf. This does not assert intrinsic compactification
independence or rational-equivalence invariance.

\begin{definition}[AlgebraicGeometry.ComplexPoint.cycleComponentSheafClassAnalyticTopology]\label{def:AlgebraicGeometry.ComplexPoint.cycleComponentSheafClassAnalyticTopology}
\lean{AlgebraicGeometry.ComplexPoint.cycleComponentSheafClassAnalyticTopology}\leanok
\uses{def:AlgebraicGeometry.ComplexPoint, def:AlgebraicGeometry.Point.chartTopology}
\noindent\emph{Lean:} \texttt{def AlgebraicGeometry.ComplexPoint.cycleComponentSheafClassAnalyticTopology}, \href{https://github.com/Paul-Lez/HodgeConjecture/blob/main/HodgeConjecture/Definitions/AlgebraicGeometry/CycleComponentSheafClass.lean#L38}{\texttt{HodgeConjecture/Definitions/AlgebraicGeometry/CycleComponentSheafClass.lean:38}}.
\end{definition}

\begin{definition}[AlgebraicGeometry.ComplexPoint.complexSupportInjectiveCohomologySheafIsoRelative]\label{def:AlgebraicGeometry.ComplexPoint.complexSupportInjectiveCohomologySheafIsoRelative}
\lean{AlgebraicGeometry.ComplexPoint.complexSupportInjectiveCohomologySheafIsoRelative}\leanok
\uses{def:AlgebraicGeometry.ComplexPoint.affineSpaceHomeomorph, def:AlgebraicGeometry.ComplexPoint.complexSupportInjectiveComplex, def:AlgebraicGeometry.ComplexPoint.complexSupportedSingularToAmbientInjective, def:AlgebraicGeometry.ComplexPoint.cycleComponentSheafClassAnalyticTopology, def:AlgebraicGeometry.ComplexProjectiveSpace.awayHomogeneousEvaluation, def:AlgebraicGeometry.ComplexProjectiveSpace.chartAffineComplexPointMap, def:AlgebraicGeometry.ComplexProjectiveSpace.normalize, def:AlgebraicGeometry.ComplexProjectiveSpace.projectiveSpaceBasicOpen, def:AlgebraicGeometry.ComplexProjectiveSpace.projectivizationToComplexPoint, def:AlgebraicGeometry.ComplexProjectiveSpace.sphereToProjectivization, def:AlgebraicGeometry.ComplexProjectiveSpace.toVectorChart, def:AlgebraicGeometry.ComplexProjectiveSpace.vectorChartToAffinePoint, def:AlgebraicGeometry.ComplexProjectiveSpace.vectorOfAwayRingHom, def:AlgebraicGeometry.IsClosedImmersion.residueFieldIso, def:AlgebraicGeometry.IsProjective, def:AlgebraicGeometry.Point.residueData, def:AlgebraicGeometry.ProjectiveSpace.Presentation.analyticImmersion, def:AlgebraicTopology.Singular.supportedSingularCohomologySheafIsoRelative, def:ProjectiveSpace.degreeZeroEquiv, def:TopCat.Sheaf.IsFlasque.BoundedBelowComplex.cyclesShortComplex, def:TopCat.Sheaf.IsFlasque.BoundedBelowComplex.globalSectionsComplex, def:TopCat.Sheaf.freeAbelianYonedaSheafHomEquiv, def:TopCat.Sheaf.openRestrictionGlobalSectionsIso, def:TopCat.Sheaf.sectionCohomologyPresheafStalkIso, def:TopCat.Sheaf.supportRestrictionComplexShortComplexMap}
The actual supported injective cohomology sheaf is the sheaf of local
relative cohomology, by the constructed singular resolution and its literal
restriction-natural comparison.

\noindent\emph{Lean:} \texttt{def AlgebraicGeometry.ComplexPoint.complexSupportInjectiveCohomologySheafIsoRelative}, \href{https://github.com/Paul-Lez/HodgeConjecture/blob/main/HodgeConjecture/Definitions/AlgebraicGeometry/CycleComponentSheafClass.lean#L41}{\texttt{HodgeConjecture/Definitions/AlgebraicGeometry/CycleComponentSheafClass.lean:41}}.
\end{definition}

\begin{definition}[AlgebraicGeometry.ComplexPoint.CycleComponentSupportedCohomology]\label{def:AlgebraicGeometry.ComplexPoint.CycleComponentSupportedCohomology}
\lean{AlgebraicGeometry.ComplexPoint.CycleComponentSupportedCohomology}\leanok
\uses{def:AlgebraicGeometry.ComplexPoint.complexSupportInjectiveComplex, def:AlgebraicGeometry.ComplexPoint.cycleComponentAnalyticClosedSupport, def:AlgebraicGeometry.ComplexPoint.cycleComponentSheafClassAnalyticTopology, def:TopCat.Sheaf.supportEvaluation}
Degree-\texttt{2p} cohomology of global sections supported on the component,
computed in the fixed ambient injective resolution.

\noindent\emph{Lean:} \texttt{abbrev AlgebraicGeometry.ComplexPoint.CycleComponentSupportedCohomology}, \href{https://github.com/Paul-Lez/HodgeConjecture/blob/main/HodgeConjecture/Definitions/AlgebraicGeometry/CycleComponentSheafClass.lean#L57}{\texttt{HodgeConjecture/Definitions/AlgebraicGeometry/CycleComponentSheafClass.lean:57}}.
\end{definition}

\begin{definition}[AlgebraicGeometry.ComplexPoint.CycleComponentSmoothCoclassSections]\label{def:AlgebraicGeometry.ComplexPoint.CycleComponentSmoothCoclassSections}
\lean{AlgebraicGeometry.ComplexPoint.CycleComponentSmoothCoclassSections}\leanok
\uses{def:AlgebraicGeometry.ComplexPoint.cycleComponentSheafClassAnalyticTopology, def:AlgebraicGeometry.ComplexPoint.cycleComponentSmoothSupportAmbientOpen, def:AlgebraicGeometry.cycleComponentSupport, def:AlgebraicTopology.Singular.supportRelativeCohomologySheaf}
Sections of the local relative-cohomology sheaf on the smooth-locus ambient open.

\noindent\emph{Lean:} \texttt{abbrev AlgebraicGeometry.ComplexPoint.CycleComponentSmoothCoclassSections}, \href{https://github.com/Paul-Lez/HodgeConjecture/blob/main/HodgeConjecture/Definitions/AlgebraicGeometry/CycleComponentSheafClass.lean#L64}{\texttt{HodgeConjecture/Definitions/AlgebraicGeometry/CycleComponentSheafClass.lean:64}}.
\end{definition}

\begin{definition}[AlgebraicGeometry.ComplexPoint.cycleComponentSupportedClassNormalizationIso]\label{def:AlgebraicGeometry.ComplexPoint.cycleComponentSupportedClassNormalizationIso}
\lean{AlgebraicGeometry.ComplexPoint.cycleComponentSupportedClassNormalizationIso}\leanok
\uses{def:AlgebraicGeometry.ComplexPoint.CycleComponentSmoothCoclassSections, def:AlgebraicGeometry.ComplexPoint.CycleComponentSupportedCohomology, def:AlgebraicGeometry.ComplexPoint.complexSupportInjectiveCohomologySheafIsoRelative, def:AlgebraicGeometry.ComplexPoint.cycleComponentSmoothSupportLowestSectionCohomologyIso, def:AlgebraicGeometry.ComplexPoint.cycleComponentSupportExtensionIso}
Supported cohomology on the full component is identified with sections
of the local relative-cohomology sheaf on its smooth-locus ambient open.
Each of the three arrows is an actual proved isomorphism.

\noindent\emph{Lean:} \texttt{def AlgebraicGeometry.ComplexPoint.cycleComponentSupportedClassNormalizationIso}, \href{https://github.com/Paul-Lez/HodgeConjecture/blob/main/HodgeConjecture/Definitions/AlgebraicGeometry/CycleComponentSheafClass.lean#L70}{\texttt{HodgeConjecture/Definitions/AlgebraicGeometry/CycleComponentSheafClass.lean:70}}.
\end{definition}

\begin{definition}[AlgebraicGeometry.ComplexPoint.cycleComponentExtendSmoothCoclass]\label{def:AlgebraicGeometry.ComplexPoint.cycleComponentExtendSmoothCoclass}
\lean{AlgebraicGeometry.ComplexPoint.cycleComponentExtendSmoothCoclass}\leanok
\uses{def:AlgebraicGeometry.ComplexPoint.cycleComponentSupportedClassNormalizationIso}
Extend a smooth-locus coclass uniquely across the singular boundary.
The inverse comes from proved purity and boundary vanishing; no extension datum is supplied.

\noindent\emph{Lean:} \texttt{def AlgebraicGeometry.ComplexPoint.cycleComponentExtendSmoothCoclass}, \href{https://github.com/Paul-Lez/HodgeConjecture/blob/main/HodgeConjecture/Definitions/AlgebraicGeometry/CycleComponentSheafClass.lean#L87}{\texttt{HodgeConjecture/Definitions/AlgebraicGeometry/CycleComponentSheafClass.lean:87}}.
\end{definition}

\begin{definition}[AlgebraicGeometry.ComplexPoint.cycleComponentSupportedInjectiveClass]\label{def:AlgebraicGeometry.ComplexPoint.cycleComponentSupportedInjectiveClass}
\lean{AlgebraicGeometry.ComplexPoint.cycleComponentSupportedInjectiveClass}\leanok
\uses{def:AlgebraicGeometry.ComplexPoint.cycleComponentExtendSmoothCoclass, def:AlgebraicGeometry.ComplexPoint.cycleComponentSmoothSupportCoclassSection}
The actual globally supported class extending the exact complex-normal
coclass. The inverse is that of the proved normalization isomorphism.

\noindent\emph{Lean:} \texttt{def AlgebraicGeometry.ComplexPoint.cycleComponentSupportedInjectiveClass}, \href{https://github.com/Paul-Lez/HodgeConjecture/blob/main/HodgeConjecture/Definitions/AlgebraicGeometry/CycleComponentSheafClass.lean#L94}{\texttt{HodgeConjecture/Definitions/AlgebraicGeometry/CycleComponentSheafClass.lean:94}}.
\end{definition}

\begin{definition}[AlgebraicGeometry.ComplexPoint.cycleComponentSheafClass]\label{def:AlgebraicGeometry.ComplexPoint.cycleComponentSheafClass}
\lean{AlgebraicGeometry.ComplexPoint.cycleComponentSheafClass}\leanok
\uses{def:AlgebraicGeometry.ComplexPoint.cycleComponentSupportedInjectiveClass, def:AlgebraicGeometry.ComplexPoint.rationalCohomologyAddEquivAmbientInjectiveHomology}
The unconditional ordinary class of an arbitrary integral component.
This uses the literal inclusion of supported injective sections.

\noindent\emph{Lean:} \texttt{def AlgebraicGeometry.ComplexPoint.cycleComponentSheafClass}, \href{https://github.com/Paul-Lez/HodgeConjecture/blob/main/HodgeConjecture/Definitions/AlgebraicGeometry/CycleComponentSheafClass.lean#L108}{\texttt{HodgeConjecture/Definitions/AlgebraicGeometry/CycleComponentSheafClass.lean:108}}.
\end{definition}

\section{The span of normalized algebraic component classes}\label{sec:HodgeConjecture.Definitions.AlgebraicGeometry.AlgebraicCycleClassSpan}
\noindent\emph{File:} \href{https://github.com/Paul-Lez/HodgeConjecture/blob/main/HodgeConjecture/Definitions/AlgebraicGeometry/AlgebraicCycleClassSpan.lean}{\texttt{HodgeConjecture/Definitions/AlgebraicGeometry/AlgebraicCycleClassSpan.lean}}.

\begin{definition}[AlgebraicGeometry.ComplexPoint.algebraicCycleClassSpan]\label{def:AlgebraicGeometry.ComplexPoint.algebraicCycleClassSpan}
\lean{AlgebraicGeometry.ComplexPoint.algebraicCycleClassSpan}\leanok
\uses{def:AlgebraicGeometry.ComplexPoint.cycleComponentSheafClass, def:AlgebraicGeometry.ComplexPoint.fieldCohomologySMul, def:AlgebraicGeometry.ComplexPoint.hypercohomologyModule}
The rational span of the actually constructed codimension-\texttt{p} component classes.

The relative dimension is the canonical \texttt{dim X}, whose certificate is proved from smoothness and
integrality. This definition spans explicit normalized component classes.

\noindent\emph{Lean:} \texttt{def AlgebraicGeometry.ComplexPoint.algebraicCycleClassSpan}, \href{https://github.com/Paul-Lez/HodgeConjecture/blob/main/HodgeConjecture/Definitions/AlgebraicGeometry/AlgebraicCycleClassSpan.lean#L21}{\texttt{HodgeConjecture/Definitions/AlgebraicGeometry/AlgebraicCycleClassSpan.lean:21}}.
\end{definition}

\chapter{Definitions: AlgebraicTopology}

\section{Singular cohomology over a field}\label{sec:HodgeConjecture.Definitions.AlgebraicTopology.SingularCohomology}
\noindent\emph{File:} \href{https://github.com/Paul-Lez/HodgeConjecture/blob/main/HodgeConjecture/Definitions/AlgebraicTopology/SingularCohomology.lean}{\texttt{HodgeConjecture/Definitions/AlgebraicTopology/SingularCohomology.lean}}.

This file constructs field-valued singular cohomology from Mathlib's singular chain complex.

\textbf{Cohomology is defined by dualising the chain complex, not by dualising homology.} The singular
cochain complex \texttt{C\^{}*(X; R)} is the degreewise \texttt{R}-linear dual of the singular chain complex
\texttt{C\_*(X; R)}, and \texttt{Cohomology R X n} is the degree-\texttt{n} homology of that cochain complex, exactly as
Mathlib defines the homology of any \texttt{HomologicalComplex}. Pullback in cohomology is the map on
homology induced by the dualised chain map; the definition never mentions singular homology.

Dualising homology instead would be wrong in general: over a ring it silently discards the \texttt{Ext}
term of the universal coefficient theorem, and even over a field it produces an object with no
cochain-level representative, so cochain-level constructions — cup products, mapping cones of
restriction, comparisons with sheaf cohomology — cannot be expressed against it. The coefficients
are a field throughout this development, and there the universal coefficient theorem does identify
the two; that identification is recorded as \texttt{cohomologyEquivDualHomology} (and
\texttt{relativeCohomologyEquivDualHomology}) rather than being taken as the definition, so every
statement that pairs a cohomology class with a homology class goes through it explicitly.

For a topological pair \texttt{A ⊆ X}, the relative chain complex is the cokernel of the actual chain map
\texttt{C\_*(A) ⟶ C\_*(X)}; the relative cochain complex is its dual, and relative cohomology is again the
homology of that complex. Cohomology with support and the map that forgets support are derived
from this construction.

\begin{definition}[AlgebraicTopology.Singular.ChainCategory]\label{def:AlgebraicTopology.Singular.ChainCategory}
\lean{AlgebraicTopology.Singular.ChainCategory}\leanok
The category containing field-valued singular chain complexes.

\noindent\emph{Lean:} \texttt{abbrev AlgebraicTopology.Singular.ChainCategory}, \href{https://github.com/Paul-Lez/HodgeConjecture/blob/main/HodgeConjecture/Definitions/AlgebraicTopology/SingularCohomology.lean#L62}{\texttt{HodgeConjecture/Definitions/AlgebraicTopology/SingularCohomology.lean:62}}.
\end{definition}

\begin{definition}[AlgebraicTopology.Singular.SingularChainComplex]\label{def:AlgebraicTopology.Singular.SingularChainComplex}
\lean{AlgebraicTopology.Singular.SingularChainComplex}\leanok
\uses{def:AlgebraicTopology.Singular.ChainCategory}
The ordinary singular chain complex of \texttt{X} with coefficients in \texttt{R}.

\noindent\emph{Lean:} \texttt{abbrev AlgebraicTopology.Singular.SingularChainComplex}, \href{https://github.com/Paul-Lez/HodgeConjecture/blob/main/HodgeConjecture/Definitions/AlgebraicTopology/SingularCohomology.lean#L66}{\texttt{HodgeConjecture/Definitions/AlgebraicTopology/SingularCohomology.lean:66}}.
\end{definition}

\begin{definition}[AlgebraicTopology.Singular.Homology]\label{def:AlgebraicTopology.Singular.Homology}
\lean{AlgebraicTopology.Singular.Homology}\leanok
Singular homology of a topological space with coefficients in a field.

\noindent\emph{Lean:} \texttt{abbrev AlgebraicTopology.Singular.Homology}, \href{https://github.com/Paul-Lez/HodgeConjecture/blob/main/HodgeConjecture/Definitions/AlgebraicTopology/SingularCohomology.lean#L75}{\texttt{HodgeConjecture/Definitions/AlgebraicTopology/SingularCohomology.lean:75}}.
\end{definition}

\begin{definition}[AlgebraicTopology.Singular.chainPairFunctor]\label{def:AlgebraicTopology.Singular.chainPairFunctor}
\lean{AlgebraicTopology.Singular.chainPairFunctor}\leanok
\uses{def:AlgebraicTopology.Singular.ChainCategory}
A topological pair \texttt{A ⊆ X}, sent to the induced arrow \texttt{C\_*(A) ⟶ C\_*(X)} of singular
chain complexes.

\noindent\emph{Lean:} \texttt{def AlgebraicTopology.Singular.chainPairFunctor}, \href{https://github.com/Paul-Lez/HodgeConjecture/blob/main/HodgeConjecture/Definitions/AlgebraicTopology/SingularCohomology.lean#L112}{\texttt{HodgeConjecture/Definitions/AlgebraicTopology/SingularCohomology.lean:112}}.
\end{definition}

\begin{definition}[AlgebraicTopology.Singular.relativeChainFunctor]\label{def:AlgebraicTopology.Singular.relativeChainFunctor}
\lean{AlgebraicTopology.Singular.relativeChainFunctor}\leanok
\uses{def:AlgebraicTopology.Singular.chainPairFunctor}
The relative singular chain complex \texttt{C\_*(X, A)}, defined as the cokernel of
\texttt{C\_*(A) ⟶ C\_*(X)}.

\noindent\emph{Lean:} \texttt{def AlgebraicTopology.Singular.relativeChainFunctor}, \href{https://github.com/Paul-Lez/HodgeConjecture/blob/main/HodgeConjecture/Definitions/AlgebraicTopology/SingularCohomology.lean#L119}{\texttt{HodgeConjecture/Definitions/AlgebraicTopology/SingularCohomology.lean:119}}.
\end{definition}

\begin{definition}[AlgebraicTopology.Singular.relativeHomologyFunctor]\label{def:AlgebraicTopology.Singular.relativeHomologyFunctor}
\lean{AlgebraicTopology.Singular.relativeHomologyFunctor}\leanok
\uses{def:AlgebraicTopology.Singular.relativeChainFunctor}
Relative singular homology.

\noindent\emph{Lean:} \texttt{def AlgebraicTopology.Singular.relativeHomologyFunctor}, \href{https://github.com/Paul-Lez/HodgeConjecture/blob/main/HodgeConjecture/Definitions/AlgebraicTopology/SingularCohomology.lean#L125}{\texttt{HodgeConjecture/Definitions/AlgebraicTopology/SingularCohomology.lean:125}}.
\end{definition}

\begin{definition}[AlgebraicTopology.Singular.RelativeHomology]\label{def:AlgebraicTopology.Singular.RelativeHomology}
\lean{AlgebraicTopology.Singular.RelativeHomology}\leanok
\uses{def:AlgebraicTopology.Singular.relativeHomologyFunctor}
Relative singular homology of the pair \texttt{A ⊆ X}.

\noindent\emph{Lean:} \texttt{abbrev AlgebraicTopology.Singular.RelativeHomology}, \href{https://github.com/Paul-Lez/HodgeConjecture/blob/main/HodgeConjecture/Definitions/AlgebraicTopology/SingularCohomology.lean#L130}{\texttt{HodgeConjecture/Definitions/AlgebraicTopology/SingularCohomology.lean:130}}.
\end{definition}

\begin{definition}[AlgebraicTopology.Singular.relativeHomologyMap]\label{def:AlgebraicTopology.Singular.relativeHomologyMap}
\lean{AlgebraicTopology.Singular.relativeHomologyMap}\leanok
\uses{def:AlgebraicTopology.Singular.RelativeHomology}
The map on relative homology induced by a map of pairs.

\noindent\emph{Lean:} \texttt{def AlgebraicTopology.Singular.relativeHomologyMap}, \href{https://github.com/Paul-Lez/HodgeConjecture/blob/main/HodgeConjecture/Definitions/AlgebraicTopology/SingularCohomology.lean#L135}{\texttt{HodgeConjecture/Definitions/AlgebraicTopology/SingularCohomology.lean:135}}.
\end{definition}

\begin{definition}[AlgebraicTopology.Singular.RelativeCochainComplex]\label{def:AlgebraicTopology.Singular.RelativeCochainComplex}
\lean{AlgebraicTopology.Singular.RelativeCochainComplex}\leanok
\uses{def:AlgebraicTopology.Singular.relativeChainFunctor, def:HomologicalComplex.linearDualCochainComplex}
The relative singular cochain complex \texttt{C\^{}*(X, A)}, the degreewise \texttt{R}-linear dual of the
relative singular chain complex.

\noindent\emph{Lean:} \texttt{abbrev AlgebraicTopology.Singular.RelativeCochainComplex}, \href{https://github.com/Paul-Lez/HodgeConjecture/blob/main/HodgeConjecture/Definitions/AlgebraicTopology/SingularCohomology.lean#L140}{\texttt{HodgeConjecture/Definitions/AlgebraicTopology/SingularCohomology.lean:140}}.
\end{definition}

\begin{definition}[AlgebraicTopology.Singular.relativeCochainComplexMap]\label{def:AlgebraicTopology.Singular.relativeCochainComplexMap}
\lean{AlgebraicTopology.Singular.relativeCochainComplexMap}\leanok
\uses{def:AlgebraicTopology.Singular.RelativeCochainComplex, def:HomologicalComplex.linearDualMap}
The relative singular cochain map induced by a map of pairs, obtained by dualising the
relative chain map.

\noindent\emph{Lean:} \texttt{abbrev AlgebraicTopology.Singular.relativeCochainComplexMap}, \href{https://github.com/Paul-Lez/HodgeConjecture/blob/main/HodgeConjecture/Definitions/AlgebraicTopology/SingularCohomology.lean#L146}{\texttt{HodgeConjecture/Definitions/AlgebraicTopology/SingularCohomology.lean:146}}.
\end{definition}

\begin{definition}[AlgebraicTopology.Singular.RelativeCohomology]\label{def:AlgebraicTopology.Singular.RelativeCohomology}
\lean{AlgebraicTopology.Singular.RelativeCohomology}\leanok
\uses{def:AlgebraicTopology.Singular.RelativeCochainComplex}
Relative singular cohomology over a field, defined as the homology of the relative singular
cochain complex — again by dualising the chain complex, not by dualising homology.

\noindent\emph{Lean:} \texttt{abbrev AlgebraicTopology.Singular.RelativeCohomology}, \href{https://github.com/Paul-Lez/HodgeConjecture/blob/main/HodgeConjecture/Definitions/AlgebraicTopology/SingularCohomology.lean#L152}{\texttt{HodgeConjecture/Definitions/AlgebraicTopology/SingularCohomology.lean:152}}.
\end{definition}

\begin{definition}[AlgebraicTopology.Singular.relativeCohomologyMap]\label{def:AlgebraicTopology.Singular.relativeCohomologyMap}
\lean{AlgebraicTopology.Singular.relativeCohomologyMap}\leanok
\uses{def:AlgebraicTopology.Singular.RelativeCohomology, def:AlgebraicTopology.Singular.relativeCochainComplexMap}
Pullback in relative singular cohomology, induced by the dualised relative chain map.

\noindent\emph{Lean:} \texttt{def AlgebraicTopology.Singular.relativeCohomologyMap}, \href{https://github.com/Paul-Lez/HodgeConjecture/blob/main/HodgeConjecture/Definitions/AlgebraicTopology/SingularCohomology.lean#L157}{\texttt{HodgeConjecture/Definitions/AlgebraicTopology/SingularCohomology.lean:157}}.
\end{definition}

\begin{definition}[AlgebraicTopology.Singular.relativeCohomologyEquivDualHomology]\label{def:AlgebraicTopology.Singular.relativeCohomologyEquivDualHomology}
\lean{AlgebraicTopology.Singular.relativeCohomologyEquivDualHomology}\leanok
\uses{def:AlgebraicTopology.Singular.RelativeCohomology, def:AlgebraicTopology.Singular.RelativeHomology, def:HomologicalComplex.linearDualHomologyEquiv}
Universal coefficients over a field for a pair: relative singular cohomology is canonically
the linear dual of relative singular homology.

\noindent\emph{Lean:} \texttt{def AlgebraicTopology.Singular.relativeCohomologyEquivDualHomology}, \href{https://github.com/Paul-Lez/HodgeConjecture/blob/main/HodgeConjecture/Definitions/AlgebraicTopology/SingularCohomology.lean#L162}{\texttt{HodgeConjecture/Definitions/AlgebraicTopology/SingularCohomology.lean:162}}.
\end{definition}

\begin{definition}[AlgebraicTopology.Singular.relativeChainProjection]\label{def:AlgebraicTopology.Singular.relativeChainProjection}
\lean{AlgebraicTopology.Singular.relativeChainProjection}\leanok
\uses{def:AlgebraicTopology.Singular.SingularChainComplex, def:AlgebraicTopology.Singular.relativeChainFunctor}
The quotient map from absolute singular chains of \texttt{X} to relative chains of \texttt{(X, A)}.

\noindent\emph{Lean:} \texttt{def AlgebraicTopology.Singular.relativeChainProjection}, \href{https://github.com/Paul-Lez/HodgeConjecture/blob/main/HodgeConjecture/Definitions/AlgebraicTopology/SingularCohomology.lean#L168}{\texttt{HodgeConjecture/Definitions/AlgebraicTopology/SingularCohomology.lean:168}}.
\end{definition}

\begin{definition}[AlgebraicTopology.Singular.relativeHomologyProjection]\label{def:AlgebraicTopology.Singular.relativeHomologyProjection}
\lean{AlgebraicTopology.Singular.relativeHomologyProjection}\leanok
\uses{def:AlgebraicTopology.Singular.Homology, def:AlgebraicTopology.Singular.RelativeHomology, def:AlgebraicTopology.Singular.relativeChainProjection}
The quotient map from absolute homology to relative homology.

\noindent\emph{Lean:} \texttt{def AlgebraicTopology.Singular.relativeHomologyProjection}, \href{https://github.com/Paul-Lez/HodgeConjecture/blob/main/HodgeConjecture/Definitions/AlgebraicTopology/SingularCohomology.lean#L173}{\texttt{HodgeConjecture/Definitions/AlgebraicTopology/SingularCohomology.lean:173}}.
\end{definition}

\section{A standard local fundamental cycle}\label{sec:HodgeConjecture.Definitions.AlgebraicTopology.LocalFundamentalClass}
\noindent\emph{File:} \href{https://github.com/Paul-Lez/HodgeConjecture/blob/main/HodgeConjecture/Definitions/AlgebraicTopology/LocalFundamentalClass.lean}{\texttt{HodgeConjecture/Definitions/AlgebraicTopology/LocalFundamentalClass.lean}}.

This file constructs a canonical relative singular cycle in
\texttt{H\_d(ℝ\^{}d, ℝ\^{}d ∖ \{0\}; ℚ)}. The affine simplex has vertices the standard basis and
\texttt{(-1, ..., -1)}. Its barycenter is the unique point that maps to the origin, while each face
misses the origin. Its relative boundary therefore vanishes.

This explicit cycle fixes the ordering and sign convention needed to normalize local Thom and
fundamental classes without adding an orientation as arbitrary data.

\begin{definition}[AlgebraicTopology.Singular.StandardRealModel]\label{def:AlgebraicTopology.Singular.StandardRealModel}
\lean{AlgebraicTopology.Singular.StandardRealModel}\leanok
The ordered real coordinate space used for the standard local class.

\noindent\emph{Lean:} \texttt{abbrev AlgebraicTopology.Singular.StandardRealModel}, \href{https://github.com/Paul-Lez/HodgeConjecture/blob/main/HodgeConjecture/Definitions/AlgebraicTopology/LocalFundamentalClass.lean#L39}{\texttt{HodgeConjecture/Definitions/AlgebraicTopology/LocalFundamentalClass.lean:39}}.
\end{definition}

\begin{definition}[AlgebraicTopology.Singular.standardPuncturedPair]\label{def:AlgebraicTopology.Singular.standardPuncturedPair}
\lean{AlgebraicTopology.Singular.standardPuncturedPair}\leanok
\uses{def:AlgebraicTopology.Singular.StandardRealModel}
The pair consisting of real coordinate space and the complement of its origin.

\noindent\emph{Lean:} \texttt{abbrev AlgebraicTopology.Singular.standardPuncturedPair}, \href{https://github.com/Paul-Lez/HodgeConjecture/blob/main/HodgeConjecture/Definitions/AlgebraicTopology/LocalFundamentalClass.lean#L42}{\texttt{HodgeConjecture/Definitions/AlgebraicTopology/LocalFundamentalClass.lean:42}}.
\end{definition}

\begin{definition}[AlgebraicTopology.Singular.standardAffineSimplex]\label{def:AlgebraicTopology.Singular.standardAffineSimplex}
\lean{AlgebraicTopology.Singular.standardAffineSimplex}\leanok
\uses{def:AlgebraicTopology.Singular.StandardRealModel}
The affine \texttt{d}-simplex with vertices the standard basis and the vector \texttt{(-1, ..., -1)}.

\noindent\emph{Lean:} \texttt{def AlgebraicTopology.Singular.standardAffineSimplex}, \href{https://github.com/Paul-Lez/HodgeConjecture/blob/main/HodgeConjecture/Definitions/AlgebraicTopology/LocalFundamentalClass.lean#L47}{\texttt{HodgeConjecture/Definitions/AlgebraicTopology/LocalFundamentalClass.lean:47}}.
\end{definition}

\begin{definition}[AlgebraicTopology.Singular.standardAffineSimplexMap]\label{def:AlgebraicTopology.Singular.standardAffineSimplexMap}
\lean{AlgebraicTopology.Singular.standardAffineSimplexMap}\leanok
\uses{def:AlgebraicTopology.Singular.standardAffineSimplex}
The continuous affine simplex underlying the standard local cycle.

\noindent\emph{Lean:} \texttt{def AlgebraicTopology.Singular.standardAffineSimplexMap}, \href{https://github.com/Paul-Lez/HodgeConjecture/blob/main/HodgeConjecture/Definitions/AlgebraicTopology/LocalFundamentalClass.lean#L87}{\texttt{HodgeConjecture/Definitions/AlgebraicTopology/LocalFundamentalClass.lean:87}}.
\end{definition}

\begin{definition}[AlgebraicTopology.Singular.standardSingularSimplex]\label{def:AlgebraicTopology.Singular.standardSingularSimplex}
\lean{AlgebraicTopology.Singular.standardSingularSimplex}\leanok
\uses{def:AlgebraicTopology.Singular.standardAffineSimplexMap, def:AlgebraicTopology.Singular.standardPuncturedPair}
The affine simplex as a singular simplex of real coordinate space.

\noindent\emph{Lean:} \texttt{def AlgebraicTopology.Singular.standardSingularSimplex}, \href{https://github.com/Paul-Lez/HodgeConjecture/blob/main/HodgeConjecture/Definitions/AlgebraicTopology/LocalFundamentalClass.lean#L93}{\texttt{HodgeConjecture/Definitions/AlgebraicTopology/LocalFundamentalClass.lean:93}}.
\end{definition}

\begin{definition}[AlgebraicTopology.Singular.standardFaceMap]\label{def:AlgebraicTopology.Singular.standardFaceMap}
\lean{AlgebraicTopology.Singular.standardFaceMap}\leanok
\uses{def:AlgebraicTopology.Singular.standardAffineSimplex}
A face of the positive-dimensional standard simplex, lifted to the punctured space.

\noindent\emph{Lean:} \texttt{def AlgebraicTopology.Singular.standardFaceMap}, \href{https://github.com/Paul-Lez/HodgeConjecture/blob/main/HodgeConjecture/Definitions/AlgebraicTopology/LocalFundamentalClass.lean#L99}{\texttt{HodgeConjecture/Definitions/AlgebraicTopology/LocalFundamentalClass.lean:99}}.
\end{definition}

\begin{definition}[AlgebraicTopology.Singular.standardFaceSimplex]\label{def:AlgebraicTopology.Singular.standardFaceSimplex}
\lean{AlgebraicTopology.Singular.standardFaceSimplex}\leanok
\uses{def:AlgebraicTopology.Singular.standardFaceMap, def:AlgebraicTopology.Singular.standardPuncturedPair}
A face of the positive-dimensional standard simplex as a singular simplex of the
punctured space.

\noindent\emph{Lean:} \texttt{def AlgebraicTopology.Singular.standardFaceSimplex}, \href{https://github.com/Paul-Lez/HodgeConjecture/blob/main/HodgeConjecture/Definitions/AlgebraicTopology/LocalFundamentalClass.lean#L108}{\texttt{HodgeConjecture/Definitions/AlgebraicTopology/LocalFundamentalClass.lean:108}}.
\end{definition}

\begin{definition}[AlgebraicTopology.Singular.standardLocalRelativeChainComplex]\label{def:AlgebraicTopology.Singular.standardLocalRelativeChainComplex}
\lean{AlgebraicTopology.Singular.standardLocalRelativeChainComplex}\leanok
\uses{def:AlgebraicTopology.Singular.relativeChainFunctor, def:AlgebraicTopology.Singular.standardPuncturedPair}
The relative singular chain complex of the standard punctured real coordinate space.

\noindent\emph{Lean:} \texttt{abbrev AlgebraicTopology.Singular.standardLocalRelativeChainComplex}, \href{https://github.com/Paul-Lez/HodgeConjecture/blob/main/HodgeConjecture/Definitions/AlgebraicTopology/LocalFundamentalClass.lean#L124}{\texttt{HodgeConjecture/Definitions/AlgebraicTopology/LocalFundamentalClass.lean:124}}.
\end{definition}

\begin{definition}[AlgebraicTopology.Singular.standardLocalProjectionComponent]\label{def:AlgebraicTopology.Singular.standardLocalProjectionComponent}
\lean{AlgebraicTopology.Singular.standardLocalProjectionComponent}\leanok
\uses{def:AlgebraicTopology.Singular.relativeChainProjection, def:AlgebraicTopology.Singular.standardLocalRelativeChainComplex}
A component of the projection to the standard relative chain complex.

\noindent\emph{Lean:} \texttt{def AlgebraicTopology.Singular.standardLocalProjectionComponent}, \href{https://github.com/Paul-Lez/HodgeConjecture/blob/main/HodgeConjecture/Definitions/AlgebraicTopology/LocalFundamentalClass.lean#L128}{\texttt{HodgeConjecture/Definitions/AlgebraicTopology/LocalFundamentalClass.lean:128}}.
\end{definition}

\begin{definition}[AlgebraicTopology.Singular.standardAmbientSimplexChain]\label{def:AlgebraicTopology.Singular.standardAmbientSimplexChain}
\lean{AlgebraicTopology.Singular.standardAmbientSimplexChain}\leanok
\uses{def:AlgebraicTopology.Singular.chainPairFunctor, def:AlgebraicTopology.Singular.standardSingularSimplex}
The standard affine simplex as an absolute singular chain.

\noindent\emph{Lean:} \texttt{def AlgebraicTopology.Singular.standardAmbientSimplexChain}, \href{https://github.com/Paul-Lez/HodgeConjecture/blob/main/HodgeConjecture/Definitions/AlgebraicTopology/LocalFundamentalClass.lean#L134}{\texttt{HodgeConjecture/Definitions/AlgebraicTopology/LocalFundamentalClass.lean:134}}.
\end{definition}

\begin{definition}[AlgebraicTopology.Singular.standardLocalChain]\label{def:AlgebraicTopology.Singular.standardLocalChain}
\lean{AlgebraicTopology.Singular.standardLocalChain}\leanok
\uses{def:AlgebraicTopology.Singular.standardAmbientSimplexChain, def:AlgebraicTopology.Singular.standardLocalProjectionComponent}
The standard affine simplex, projected to the relative singular chain complex of
\texttt{(ℝ\^{}d, ℝ\^{}d ∖ \{0\})}.

\noindent\emph{Lean:} \texttt{def AlgebraicTopology.Singular.standardLocalChain}, \href{https://github.com/Paul-Lez/HodgeConjecture/blob/main/HodgeConjecture/Definitions/AlgebraicTopology/LocalFundamentalClass.lean#L141}{\texttt{HodgeConjecture/Definitions/AlgebraicTopology/LocalFundamentalClass.lean:141}}.
\end{definition}

\begin{definition}[AlgebraicTopology.Singular.standardAmbientFaceChain]\label{def:AlgebraicTopology.Singular.standardAmbientFaceChain}
\lean{AlgebraicTopology.Singular.standardAmbientFaceChain}\leanok
\uses{def:AlgebraicTopology.Singular.chainPairFunctor, def:AlgebraicTopology.Singular.standardSingularSimplex}
A face of the standard simplex as an absolute singular chain.

\noindent\emph{Lean:} \texttt{def AlgebraicTopology.Singular.standardAmbientFaceChain}, \href{https://github.com/Paul-Lez/HodgeConjecture/blob/main/HodgeConjecture/Definitions/AlgebraicTopology/LocalFundamentalClass.lean#L147}{\texttt{HodgeConjecture/Definitions/AlgebraicTopology/LocalFundamentalClass.lean:147}}.
\end{definition}

\begin{definition}[AlgebraicTopology.Singular.standardSubspaceFaceChain]\label{def:AlgebraicTopology.Singular.standardSubspaceFaceChain}
\lean{AlgebraicTopology.Singular.standardSubspaceFaceChain}\leanok
\uses{def:AlgebraicTopology.Singular.chainPairFunctor, def:AlgebraicTopology.Singular.standardFaceSimplex}
A face of the standard simplex as a chain in the punctured subspace.

\noindent\emph{Lean:} \texttt{def AlgebraicTopology.Singular.standardSubspaceFaceChain}, \href{https://github.com/Paul-Lez/HodgeConjecture/blob/main/HodgeConjecture/Definitions/AlgebraicTopology/LocalFundamentalClass.lean#L155}{\texttt{HodgeConjecture/Definitions/AlgebraicTopology/LocalFundamentalClass.lean:155}}.
\end{definition}

\begin{definition}[AlgebraicTopology.Singular.standardLocalCycle]\label{def:AlgebraicTopology.Singular.standardLocalCycle}
\lean{AlgebraicTopology.Singular.standardLocalCycle}\leanok
\uses{def:AlgebraicTopology.Singular.standardAmbientFaceChain, def:AlgebraicTopology.Singular.standardLocalChain, def:AlgebraicTopology.Singular.standardSubspaceFaceChain}
The standard relative cycle represented by an affine simplex meeting the origin once.

\noindent\emph{Lean:} \texttt{def AlgebraicTopology.Singular.standardLocalCycle}, \href{https://github.com/Paul-Lez/HodgeConjecture/blob/main/HodgeConjecture/Definitions/AlgebraicTopology/LocalFundamentalClass.lean#L232}{\texttt{HodgeConjecture/Definitions/AlgebraicTopology/LocalFundamentalClass.lean:232}}.
\end{definition}

\begin{definition}[AlgebraicTopology.Singular.standardLocalClass]\label{def:AlgebraicTopology.Singular.standardLocalClass}
\lean{AlgebraicTopology.Singular.standardLocalClass}\leanok
\uses{def:AlgebraicTopology.Singular.RelativeHomology, def:AlgebraicTopology.Singular.standardLocalCycle}
The canonical class represented by the standard affine local cycle in
\texttt{H\_d(ℝ\^{}d, ℝ\^{}d ∖ \{0\}; ℚ)}.

\noindent\emph{Lean:} \texttt{def AlgebraicTopology.Singular.standardLocalClass}, \href{https://github.com/Paul-Lez/HodgeConjecture/blob/main/HodgeConjecture/Definitions/AlgebraicTopology/LocalFundamentalClass.lean#L238}{\texttt{HodgeConjecture/Definitions/AlgebraicTopology/LocalFundamentalClass.lean:238}}.
\end{definition}

\section{The standard complex local class}\label{sec:HodgeConjecture.Definitions.AlgebraicTopology.ComplexOrientation}
\noindent\emph{File:} \href{https://github.com/Paul-Lez/HodgeConjecture/blob/main/HodgeConjecture/Definitions/AlgebraicTopology/ComplexOrientation.lean}{\texttt{HodgeConjecture/Definitions/AlgebraicTopology/ComplexOrientation.lean}}.

This file identifies \texttt{ℂ\^{}p} with an ordered real coordinate space by listing the real and imaginary
part of each complex coordinate consecutively. It transports the explicit standard local cycle to
\texttt{H\_\{2p\}(ℂ\^{}p, ℂ\^{}p ∖ \{0\}; ℚ)}. Thus the complex orientation is constructed from complex coordinates
rather than supplied as data.

\begin{definition}[AlgebraicTopology.Singular.complexRealHomeomorph]\label{def:AlgebraicTopology.Singular.complexRealHomeomorph}
\lean{AlgebraicTopology.Singular.complexRealHomeomorph}\leanok
\uses{def:AlgebraicTopology.Singular.StandardRealModel, def:Complex.piCoordCLE}
The orientation-ordered homeomorphism from complex coordinate space to real coordinate space.

This is the coordinate map of \texttt{Complex.piBasisOneI}, so it lists the real and imaginary part of
each complex coordinate consecutively; continuity in both directions comes from
\texttt{Basis.equivFunL}.

\noindent\emph{Lean:} \texttt{def AlgebraicTopology.Singular.complexRealHomeomorph}, \href{https://github.com/Paul-Lez/HodgeConjecture/blob/main/HodgeConjecture/Definitions/AlgebraicTopology/ComplexOrientation.lean#L37}{\texttt{HodgeConjecture/Definitions/AlgebraicTopology/ComplexOrientation.lean:37}}.
\end{definition}

\begin{definition}[AlgebraicTopology.Singular.standardComplexPuncturedPair]\label{def:AlgebraicTopology.Singular.standardComplexPuncturedPair}
\lean{AlgebraicTopology.Singular.standardComplexPuncturedPair}\leanok
Complex coordinate space paired with the complement of its origin.

\noindent\emph{Lean:} \texttt{abbrev AlgebraicTopology.Singular.standardComplexPuncturedPair}, \href{https://github.com/Paul-Lez/HodgeConjecture/blob/main/HodgeConjecture/Definitions/AlgebraicTopology/ComplexOrientation.lean#L45}{\texttt{HodgeConjecture/Definitions/AlgebraicTopology/ComplexOrientation.lean:45}}.
\end{definition}

\begin{definition}[AlgebraicTopology.Singular.puncturedRealToComplex]\label{def:AlgebraicTopology.Singular.puncturedRealToComplex}
\lean{AlgebraicTopology.Singular.puncturedRealToComplex}\leanok
\uses{def:AlgebraicTopology.Singular.StandardRealModel, def:Complex.piCoordCLE}
The inverse coordinate homeomorphism restricted to the punctured spaces.

\noindent\emph{Lean:} \texttt{def AlgebraicTopology.Singular.puncturedRealToComplex}, \href{https://github.com/Paul-Lez/HodgeConjecture/blob/main/HodgeConjecture/Definitions/AlgebraicTopology/ComplexOrientation.lean#L49}{\texttt{HodgeConjecture/Definitions/AlgebraicTopology/ComplexOrientation.lean:49}}.
\end{definition}

\begin{definition}[AlgebraicTopology.Singular.standardRealToComplexPair]\label{def:AlgebraicTopology.Singular.standardRealToComplexPair}
\lean{AlgebraicTopology.Singular.standardRealToComplexPair}\leanok
\uses{def:AlgebraicTopology.Singular.complexRealHomeomorph, def:AlgebraicTopology.Singular.puncturedRealToComplex, def:AlgebraicTopology.Singular.standardComplexPuncturedPair, def:AlgebraicTopology.Singular.standardPuncturedPair}
The inverse coordinate homeomorphism as a morphism of punctured pairs.

\noindent\emph{Lean:} \texttt{def AlgebraicTopology.Singular.standardRealToComplexPair}, \href{https://github.com/Paul-Lez/HodgeConjecture/blob/main/HodgeConjecture/Definitions/AlgebraicTopology/ComplexOrientation.lean#L59}{\texttt{HodgeConjecture/Definitions/AlgebraicTopology/ComplexOrientation.lean:59}}.
\end{definition}

\begin{definition}[AlgebraicTopology.Singular.puncturedComplexToReal]\label{def:AlgebraicTopology.Singular.puncturedComplexToReal}
\lean{AlgebraicTopology.Singular.puncturedComplexToReal}\leanok
\uses{def:AlgebraicTopology.Singular.StandardRealModel, def:Complex.piCoordCLE}
The complex-to-real coordinate homeomorphism restricted to the punctured spaces.

\noindent\emph{Lean:} \texttt{def AlgebraicTopology.Singular.puncturedComplexToReal}, \href{https://github.com/Paul-Lez/HodgeConjecture/blob/main/HodgeConjecture/Definitions/AlgebraicTopology/ComplexOrientation.lean#L68}{\texttt{HodgeConjecture/Definitions/AlgebraicTopology/ComplexOrientation.lean:68}}.
\end{definition}

\begin{definition}[AlgebraicTopology.Singular.standardComplexToRealPair]\label{def:AlgebraicTopology.Singular.standardComplexToRealPair}
\lean{AlgebraicTopology.Singular.standardComplexToRealPair}\leanok
\uses{def:AlgebraicTopology.Singular.complexRealHomeomorph, def:AlgebraicTopology.Singular.puncturedComplexToReal, def:AlgebraicTopology.Singular.standardComplexPuncturedPair, def:AlgebraicTopology.Singular.standardPuncturedPair}
The coordinate homeomorphism as a morphism from the complex pair to the real pair.

\noindent\emph{Lean:} \texttt{def AlgebraicTopology.Singular.standardComplexToRealPair}, \href{https://github.com/Paul-Lez/HodgeConjecture/blob/main/HodgeConjecture/Definitions/AlgebraicTopology/ComplexOrientation.lean#L78}{\texttt{HodgeConjecture/Definitions/AlgebraicTopology/ComplexOrientation.lean:78}}.
\end{definition}

\begin{definition}[AlgebraicTopology.Singular.standardComplexRealPairIso]\label{def:AlgebraicTopology.Singular.standardComplexRealPairIso}
\lean{AlgebraicTopology.Singular.standardComplexRealPairIso}\leanok
\uses{def:AlgebraicTopology.Singular.standardComplexToRealPair, def:AlgebraicTopology.Singular.standardRealToComplexPair}
The isomorphism of punctured pairs induced by ordered real and imaginary coordinates.

\noindent\emph{Lean:} \texttt{def AlgebraicTopology.Singular.standardComplexRealPairIso}, \href{https://github.com/Paul-Lez/HodgeConjecture/blob/main/HodgeConjecture/Definitions/AlgebraicTopology/ComplexOrientation.lean#L86}{\texttt{HodgeConjecture/Definitions/AlgebraicTopology/ComplexOrientation.lean:86}}.
\end{definition}

\begin{definition}[AlgebraicTopology.Singular.standardComplexRealRelativeHomologyIso]\label{def:AlgebraicTopology.Singular.standardComplexRealRelativeHomologyIso}
\lean{AlgebraicTopology.Singular.standardComplexRealRelativeHomologyIso}\leanok
\uses{def:AlgebraicTopology.Singular.RelativeHomology, def:AlgebraicTopology.Singular.standardComplexRealPairIso}
The induced isomorphism between complex and real local homology.

\noindent\emph{Lean:} \texttt{def AlgebraicTopology.Singular.standardComplexRealRelativeHomologyIso}, \href{https://github.com/Paul-Lez/HodgeConjecture/blob/main/HodgeConjecture/Definitions/AlgebraicTopology/ComplexOrientation.lean#L104}{\texttt{HodgeConjecture/Definitions/AlgebraicTopology/ComplexOrientation.lean:104}}.
\end{definition}

\begin{definition}[AlgebraicTopology.Singular.standardComplexLocalClass]\label{def:AlgebraicTopology.Singular.standardComplexLocalClass}
\lean{AlgebraicTopology.Singular.standardComplexLocalClass}\leanok
\uses{def:AlgebraicTopology.Singular.standardComplexRealRelativeHomologyIso, def:AlgebraicTopology.Singular.standardLocalClass}
The standard complex local class obtained from the explicit real local cycle.

\noindent\emph{Lean:} \texttt{def AlgebraicTopology.Singular.standardComplexLocalClass}, \href{https://github.com/Paul-Lez/HodgeConjecture/blob/main/HodgeConjecture/Definitions/AlgebraicTopology/ComplexOrientation.lean#L110}{\texttt{HodgeConjecture/Definitions/AlgebraicTopology/ComplexOrientation.lean:110}}.
\end{definition}

\section{Concrete sheaf sections with support and their right derived functor}\label{sec:HodgeConjecture.Definitions.AlgebraicTopology.DerivedSheafSupport}
\noindent\emph{File:} \href{https://github.com/Paul-Lez/HodgeConjecture/blob/main/HodgeConjecture/Definitions/AlgebraicTopology/DerivedSheafSupport.lean}{\texttt{HodgeConjecture/Definitions/AlgebraicTopology/DerivedSheafSupport.lean}}.

For an open subset \texttt{U} of \texttt{X}, restriction gives a natural map \texttt{F → j\_* (F\textbar{}\_U)}.
Its kernel is the sheaf of sections supported on the closed complement of \texttt{U}.
This file constructs that additive functor and its right derived functor on the
bounded-below derived category using Mathlib's injective-resolution machinery.

Both the sheaf-valued local cohomology operation, conventionally \texttt{RΓ\_Z}, and
the abelian-group-valued derived global sections \texttt{RΓ\_Z} are constructed, with
different source and target categories displayed explicitly.
No dualizing complex or orientation is assumed or constructed. In particular, this
file does not identify Borel--Moore homology with supported cohomology. Comparison
with the project's restriction mapping-cone model, and coherent commutation with
shifts, remain separate theorems.

\begin{definition}[TopCat.Sheaf.openRestrictionPushforward]\label{def:TopCat.Sheaf.openRestrictionPushforward}
\lean{TopCat.Sheaf.openRestrictionPushforward}\leanok
Restrict a sheaf to an open subspace and push it forward again. The pullback
here is the concrete open-embedding pullback, obtained by evaluating on the
corresponding ambient open sets.

\noindent\emph{Lean:} \texttt{def TopCat.Sheaf.openRestrictionPushforward}, \href{https://github.com/Paul-Lez/HodgeConjecture/blob/main/HodgeConjecture/Definitions/AlgebraicTopology/DerivedSheafSupport.lean#L40}{\texttt{HodgeConjecture/Definitions/AlgebraicTopology/DerivedSheafSupport.lean:40}}.
\end{definition}

\begin{definition}[TopCat.Sheaf.toOpenRestrictionPushforward]\label{def:TopCat.Sheaf.toOpenRestrictionPushforward}
\lean{TopCat.Sheaf.toOpenRestrictionPushforward}\leanok
\uses{def:TopCat.Sheaf.openRestrictionPushforward}
The actual restriction morphism, functorial in the coefficient sheaf.

\noindent\emph{Lean:} \texttt{def TopCat.Sheaf.toOpenRestrictionPushforward}, \href{https://github.com/Paul-Lez/HodgeConjecture/blob/main/HodgeConjecture/Definitions/AlgebraicTopology/DerivedSheafSupport.lean#L48}{\texttt{HodgeConjecture/Definitions/AlgebraicTopology/DerivedSheafSupport.lean:48}}.
\end{definition}

\begin{definition}[TopCat.Sheaf.sheafSectionsSupportedOutside]\label{def:TopCat.Sheaf.sheafSectionsSupportedOutside}
\lean{TopCat.Sheaf.sheafSectionsSupportedOutside}\leanok
\uses{def:TopCat.Sheaf.toOpenRestrictionPushforward}
The sheaf of sections vanishing on \texttt{U}, defined as the kernel of the
coefficient-wise restriction map.

\noindent\emph{Lean:} \texttt{def TopCat.Sheaf.sheafSectionsSupportedOutside}, \href{https://github.com/Paul-Lez/HodgeConjecture/blob/main/HodgeConjecture/Definitions/AlgebraicTopology/DerivedSheafSupport.lean#L62}{\texttt{HodgeConjecture/Definitions/AlgebraicTopology/DerivedSheafSupport.lean:62}}.
\end{definition}

\begin{definition}[TopCat.Sheaf.sheafSectionsSupportedOutsideInclusion]\label{def:TopCat.Sheaf.sheafSectionsSupportedOutsideInclusion}
\lean{TopCat.Sheaf.sheafSectionsSupportedOutsideInclusion}\leanok
\uses{def:TopCat.Sheaf.sheafSectionsSupportedOutside}
Inclusion of supported sections into the original coefficient sheaf.

\noindent\emph{Lean:} \texttt{def TopCat.Sheaf.sheafSectionsSupportedOutsideInclusion}, \href{https://github.com/Paul-Lez/HodgeConjecture/blob/main/HodgeConjecture/Definitions/AlgebraicTopology/DerivedSheafSupport.lean#L82}{\texttt{HodgeConjecture/Definitions/AlgebraicTopology/DerivedSheafSupport.lean:82}}.
\end{definition}

\begin{definition}[TopCat.Sheaf.liftSheafSectionsSupportedOutside]\label{def:TopCat.Sheaf.liftSheafSectionsSupportedOutside}
\lean{TopCat.Sheaf.liftSheafSectionsSupportedOutside}\leanok
\uses{def:TopCat.Sheaf.sheafSectionsSupportedOutside}
A morphism of sheaves whose restriction to \texttt{U} is zero factors canonically
through the sheaf of sections supported outside \texttt{U}.

\noindent\emph{Lean:} \texttt{def TopCat.Sheaf.liftSheafSectionsSupportedOutside}, \href{https://github.com/Paul-Lez/HodgeConjecture/blob/main/HodgeConjecture/Definitions/AlgebraicTopology/DerivedSheafSupport.lean#L95}{\texttt{HodgeConjecture/Definitions/AlgebraicTopology/DerivedSheafSupport.lean:95}}.
\end{definition}

\begin{definition}[TopCat.Sheaf.sheafSectionsSupportedOutsideOnOpenIso]\label{def:TopCat.Sheaf.sheafSectionsSupportedOutsideOnOpenIso}
\lean{TopCat.Sheaf.sheafSectionsSupportedOutsideOnOpenIso}\leanok
\uses{def:TopCat.Sheaf.sheafSectionsSupportedOutside}
On every ambient open set, supported sections are exactly the kernel of
restriction to its intersection with \texttt{U}. This is the canonical kernel
comparison, not a supplied equivalence.

\noindent\emph{Lean:} \texttt{def TopCat.Sheaf.sheafSectionsSupportedOutsideOnOpenIso}, \href{https://github.com/Paul-Lez/HodgeConjecture/blob/main/HodgeConjecture/Definitions/AlgebraicTopology/DerivedSheafSupport.lean#L111}{\texttt{HodgeConjecture/Definitions/AlgebraicTopology/DerivedSheafSupport.lean:111}}.
\end{definition}

\section{The localization sequence on injective coefficient complexes}\label{sec:HodgeConjecture.Definitions.AlgebraicTopology.DerivedSheafSupportLocalization}
\noindent\emph{File:} \href{https://github.com/Paul-Lez/HodgeConjecture/blob/main/HodgeConjecture/Definitions/AlgebraicTopology/DerivedSheafSupportLocalization.lean}{\texttt{HodgeConjecture/Definitions/AlgebraicTopology/DerivedSheafSupportLocalization.lean}}.

For an injective coefficient sheaf, the kernel defining sections supported outside \texttt{U} fits
into a short exact sequence with the coefficient sheaf and its restriction-pushforward, with
surjectivity coming from flasqueness of injective sheaves. The sequence stays short exact
after evaluation on an arbitrary open set, giving the complex of supported sections in the
localization calculation. The comparison with the mapping cocone is a quasi-isomorphism
normalized by its projection to the coefficient complex, and the resulting isomorphisms are
built from the \texttt{D⁺} right-derived functors and their canonical units on bounded-below
injective models, then displayed in \texttt{D} along the full inclusion \texttt{D⁺ → D}.

Identifying these fibers with the repository's constant-rational restriction-cone model
requires a normalized comparison between restriction of the ambient injective model and an
independently chosen injective resolution on the complement, which is still to be supplied.

\begin{definition}[TopCat.Sheaf.supportRestrictionShortComplex]\label{def:TopCat.Sheaf.supportRestrictionShortComplex}
\lean{TopCat.Sheaf.supportRestrictionShortComplex}\leanok
\uses{def:TopCat.Sheaf.sheafSectionsSupportedOutsideInclusion}
The actual support inclusion followed by restriction-pushforward.

\noindent\emph{Lean:} \texttt{def TopCat.Sheaf.supportRestrictionShortComplex}, \href{https://github.com/Paul-Lez/HodgeConjecture/blob/main/HodgeConjecture/Definitions/AlgebraicTopology/DerivedSheafSupportLocalization.lean#L37}{\texttt{HodgeConjecture/Definitions/AlgebraicTopology/DerivedSheafSupportLocalization.lean:37}}.
\end{definition}

\begin{definition}[TopCat.Sheaf.supportEvaluation]\label{def:TopCat.Sheaf.supportEvaluation}
\lean{TopCat.Sheaf.supportEvaluation}\leanok
Evaluation of an additive sheaf on an ambient open set.

\noindent\emph{Lean:} \texttt{def TopCat.Sheaf.supportEvaluation}, \href{https://github.com/Paul-Lez/HodgeConjecture/blob/main/HodgeConjecture/Definitions/AlgebraicTopology/DerivedSheafSupportLocalization.lean#L68}{\texttt{HodgeConjecture/Definitions/AlgebraicTopology/DerivedSheafSupportLocalization.lean:68}}.
\end{definition}

\begin{definition}[TopCat.Sheaf.supportRestrictionSectionsShortComplex]\label{def:TopCat.Sheaf.supportRestrictionSectionsShortComplex}
\lean{TopCat.Sheaf.supportRestrictionSectionsShortComplex}\leanok
\uses{def:TopCat.Sheaf.supportEvaluation, def:TopCat.Sheaf.supportRestrictionShortComplex}
The sequence of sections on \texttt{V}, with the actual supported-sections inclusion
and actual restriction map.

\noindent\emph{Lean:} \texttt{def TopCat.Sheaf.supportRestrictionSectionsShortComplex}, \href{https://github.com/Paul-Lez/HodgeConjecture/blob/main/HodgeConjecture/Definitions/AlgebraicTopology/DerivedSheafSupportLocalization.lean#L83}{\texttt{HodgeConjecture/Definitions/AlgebraicTopology/DerivedSheafSupportLocalization.lean:83}}.
\end{definition}

\begin{definition}[TopCat.Sheaf.supportRestrictionComplexShortComplex]\label{def:TopCat.Sheaf.supportRestrictionComplexShortComplex}
\lean{TopCat.Sheaf.supportRestrictionComplexShortComplex}\leanok
\uses{def:TopCat.Sheaf.sheafSectionsSupportedOutsideInclusion}
The actual support/restriction sequence applied termwise to a coefficient
complex. No boundedness is needed for this algebraic sequence.

\noindent\emph{Lean:} \texttt{def TopCat.Sheaf.supportRestrictionComplexShortComplex}, \href{https://github.com/Paul-Lez/HodgeConjecture/blob/main/HodgeConjecture/Definitions/AlgebraicTopology/DerivedSheafSupportLocalization.lean#L106}{\texttt{HodgeConjecture/Definitions/AlgebraicTopology/DerivedSheafSupportLocalization.lean:106}}.
\end{definition}

\begin{definition}[TopCat.Sheaf.supportRestrictionSectionsComplexShortComplex]\label{def:TopCat.Sheaf.supportRestrictionSectionsComplexShortComplex}
\lean{TopCat.Sheaf.supportRestrictionSectionsComplexShortComplex}\leanok
\uses{def:TopCat.Sheaf.supportEvaluation, def:TopCat.Sheaf.supportRestrictionComplexShortComplex}
Evaluate the actual termwise support/restriction sequence on \texttt{V}.

\noindent\emph{Lean:} \texttt{def TopCat.Sheaf.supportRestrictionSectionsComplexShortComplex}, \href{https://github.com/Paul-Lez/HodgeConjecture/blob/main/HodgeConjecture/Definitions/AlgebraicTopology/DerivedSheafSupportLocalization.lean#L128}{\texttt{HodgeConjecture/Definitions/AlgebraicTopology/DerivedSheafSupportLocalization.lean:128}}.
\end{definition}

\section{The singular-cochain sheaf}\label{sec:HodgeConjecture.Definitions.AlgebraicTopology.SingularCochainSheaf}
\noindent\emph{File:} \href{https://github.com/Paul-Lez/HodgeConjecture/blob/main/HodgeConjecture/Definitions/AlgebraicTopology/SingularCochainSheaf.lean}{\texttt{HodgeConjecture/Definitions/AlgebraicTopology/SingularCochainSheaf.lean}}.

This file constructs the singular cochain presheaf on a topological space. In degree \texttt{n}, its
sections over an open set \texttt{U} are the linear dual of the singular \texttt{n}-chains of \texttt{U}. Restriction is
dual to the chain map induced by an inclusion of open sets. The singular boundary dualizes to the
coboundary, giving a cochain complex. Degreewise sheafification gives the singular-cochain sheaf
complex.

These constructions are prerequisites for comparing singular cohomology with constant-sheaf
cohomology. The comparison is proved in degree zero on locally path-connected spaces. In positive
degrees, an explicit neighborhood-wise primitive condition is shown to imply the comparison.
No local acyclicity statement is assumed.

\begin{definition}[AlgebraicTopology.Singular.openSingularChainComplexFunctor]\label{def:AlgebraicTopology.Singular.openSingularChainComplexFunctor}
\lean{AlgebraicTopology.Singular.openSingularChainComplexFunctor}\leanok
The singular chain complex, functorially restricted to the open subsets of \texttt{X}.

\noindent\emph{Lean:} \texttt{def AlgebraicTopology.Singular.openSingularChainComplexFunctor}, \href{https://github.com/Paul-Lez/HodgeConjecture/blob/main/HodgeConjecture/Definitions/AlgebraicTopology/SingularCochainSheaf.lean#L58}{\texttt{HodgeConjecture/Definitions/AlgebraicTopology/SingularCochainSheaf.lean:58}}.
\end{definition}

\begin{definition}[AlgebraicTopology.Singular.OpenChains]\label{def:AlgebraicTopology.Singular.OpenChains}
\lean{AlgebraicTopology.Singular.OpenChains}\leanok
\uses{def:AlgebraicTopology.Singular.openSingularChainComplexFunctor}
Singular chains of an open subset of \texttt{X} in degree \texttt{n}.

\noindent\emph{Lean:} \texttt{abbrev AlgebraicTopology.Singular.OpenChains}, \href{https://github.com/Paul-Lez/HodgeConjecture/blob/main/HodgeConjecture/Definitions/AlgebraicTopology/SingularCochainSheaf.lean#L64}{\texttt{HodgeConjecture/Definitions/AlgebraicTopology/SingularCochainSheaf.lean:64}}.
\end{definition}

\begin{definition}[AlgebraicTopology.Singular.OpenCochains]\label{def:AlgebraicTopology.Singular.OpenCochains}
\lean{AlgebraicTopology.Singular.OpenCochains}\leanok
\uses{def:AlgebraicTopology.Singular.OpenChains}
Singular cochains of an open subset of \texttt{X} in degree \texttt{n}.

\noindent\emph{Lean:} \texttt{abbrev AlgebraicTopology.Singular.OpenCochains}, \href{https://github.com/Paul-Lez/HodgeConjecture/blob/main/HodgeConjecture/Definitions/AlgebraicTopology/SingularCochainSheaf.lean#L68}{\texttt{HodgeConjecture/Definitions/AlgebraicTopology/SingularCochainSheaf.lean:68}}.
\end{definition}

\begin{definition}[AlgebraicTopology.Singular.singularCochainPresheaf]\label{def:AlgebraicTopology.Singular.singularCochainPresheaf}
\lean{AlgebraicTopology.Singular.singularCochainPresheaf}\leanok
\uses{def:AlgebraicTopology.Singular.OpenCochains}
The singular cochain presheaf in a fixed degree.

\noindent\emph{Lean:} \texttt{def AlgebraicTopology.Singular.singularCochainPresheaf}, \href{https://github.com/Paul-Lez/HodgeConjecture/blob/main/HodgeConjecture/Definitions/AlgebraicTopology/SingularCochainSheaf.lean#L72}{\texttt{HodgeConjecture/Definitions/AlgebraicTopology/SingularCochainSheaf.lean:72}}.
\end{definition}

\begin{definition}[AlgebraicTopology.Singular.singularCochainCoboundary]\label{def:AlgebraicTopology.Singular.singularCochainCoboundary}
\lean{AlgebraicTopology.Singular.singularCochainCoboundary}\leanok
\uses{def:AlgebraicTopology.Singular.singularCochainPresheaf}
The singular coboundary, dual to the singular boundary.

\noindent\emph{Lean:} \texttt{def AlgebraicTopology.Singular.singularCochainCoboundary}, \href{https://github.com/Paul-Lez/HodgeConjecture/blob/main/HodgeConjecture/Definitions/AlgebraicTopology/SingularCochainSheaf.lean#L88}{\texttt{HodgeConjecture/Definitions/AlgebraicTopology/SingularCochainSheaf.lean:88}}.
\end{definition}

\begin{definition}[AlgebraicTopology.Singular.singularCochainPresheafComplex]\label{def:AlgebraicTopology.Singular.singularCochainPresheafComplex}
\lean{AlgebraicTopology.Singular.singularCochainPresheafComplex}\leanok
\uses{def:AlgebraicTopology.Singular.singularCochainCoboundary}
The singular cochain complex as a complex of presheaves.

\noindent\emph{Lean:} \texttt{def AlgebraicTopology.Singular.singularCochainPresheafComplex}, \href{https://github.com/Paul-Lez/HodgeConjecture/blob/main/HodgeConjecture/Definitions/AlgebraicTopology/SingularCochainSheaf.lean#L136}{\texttt{HodgeConjecture/Definitions/AlgebraicTopology/SingularCochainSheaf.lean:136}}.
\end{definition}

\begin{definition}[AlgebraicTopology.Singular.singularCochainSheaf]\label{def:AlgebraicTopology.Singular.singularCochainSheaf}
\lean{AlgebraicTopology.Singular.singularCochainSheaf}\leanok
\uses{def:AlgebraicTopology.Singular.singularCochainPresheaf}
Singular cochains in a fixed degree, sheafified as additive groups.

\noindent\emph{Lean:} \texttt{def AlgebraicTopology.Singular.singularCochainSheaf}, \href{https://github.com/Paul-Lez/HodgeConjecture/blob/main/HodgeConjecture/Definitions/AlgebraicTopology/SingularCochainSheaf.lean#L150}{\texttt{HodgeConjecture/Definitions/AlgebraicTopology/SingularCochainSheaf.lean:150}}.
\end{definition}

\begin{definition}[AlgebraicTopology.Singular.singularCochainSheafCoboundary]\label{def:AlgebraicTopology.Singular.singularCochainSheafCoboundary}
\lean{AlgebraicTopology.Singular.singularCochainSheafCoboundary}\leanok
\uses{def:AlgebraicTopology.Singular.singularCochainCoboundary, def:AlgebraicTopology.Singular.singularCochainSheaf}
The sheafified singular coboundary.

\noindent\emph{Lean:} \texttt{def AlgebraicTopology.Singular.singularCochainSheafCoboundary}, \href{https://github.com/Paul-Lez/HodgeConjecture/blob/main/HodgeConjecture/Definitions/AlgebraicTopology/SingularCochainSheaf.lean#L155}{\texttt{HodgeConjecture/Definitions/AlgebraicTopology/SingularCochainSheaf.lean:155}}.
\end{definition}

\begin{definition}[AlgebraicTopology.Singular.singularCochainSheafComplex]\label{def:AlgebraicTopology.Singular.singularCochainSheafComplex}
\lean{AlgebraicTopology.Singular.singularCochainSheafComplex}\leanok
\uses{def:AlgebraicTopology.Singular.singularCochainPresheafComplex}
The degreewise sheafification of the singular cochain complex.

\noindent\emph{Lean:} \texttt{def AlgebraicTopology.Singular.singularCochainSheafComplex}, \href{https://github.com/Paul-Lez/HodgeConjecture/blob/main/HodgeConjecture/Definitions/AlgebraicTopology/SingularCochainSheaf.lean#L161}{\texttt{HodgeConjecture/Definitions/AlgebraicTopology/SingularCochainSheaf.lean:161}}.
\end{definition}

\begin{definition}[AlgebraicTopology.Singular.singularCochainSheafificationUnit]\label{def:AlgebraicTopology.Singular.singularCochainSheafificationUnit}
\lean{AlgebraicTopology.Singular.singularCochainSheafificationUnit}\leanok
\uses{def:AlgebraicTopology.Singular.singularCochainSheafCoboundary, def:AlgebraicTopology.Singular.singularCochainSheafComplex}
The degreewise sheafification unit from singular cochains to the underlying presheaf of the
singular-cochain sheaf complex.

\noindent\emph{Lean:} \texttt{def AlgebraicTopology.Singular.singularCochainSheafificationUnit}, \href{https://github.com/Paul-Lez/HodgeConjecture/blob/main/HodgeConjecture/Definitions/AlgebraicTopology/SingularCochainSheaf.lean#L179}{\texttt{HodgeConjecture/Definitions/AlgebraicTopology/SingularCochainSheaf.lean:179}}.
\end{definition}

\begin{definition}[AlgebraicTopology.Singular.simplicialZeroAugmentation]\label{def:AlgebraicTopology.Singular.simplicialZeroAugmentation}
\lean{AlgebraicTopology.Singular.simplicialZeroAugmentation}\leanok
The augmentation of simplicial zero-chains, sending every vertex to \texttt{1}.

\noindent\emph{Lean:} \texttt{def AlgebraicTopology.Singular.simplicialZeroAugmentation}, \href{https://github.com/Paul-Lez/HodgeConjecture/blob/main/HodgeConjecture/Definitions/AlgebraicTopology/SingularCochainSheaf.lean#L195}{\texttt{HodgeConjecture/Definitions/AlgebraicTopology/SingularCochainSheaf.lean:195}}.
\end{definition}

\begin{definition}[AlgebraicTopology.Singular.openZeroAugmentation]\label{def:AlgebraicTopology.Singular.openZeroAugmentation}
\lean{AlgebraicTopology.Singular.openZeroAugmentation}\leanok
\uses{def:AlgebraicTopology.Singular.OpenChains, def:AlgebraicTopology.Singular.simplicialZeroAugmentation}
The augmentation on the zero-chains of an open subset.

\noindent\emph{Lean:} \texttt{def AlgebraicTopology.Singular.openZeroAugmentation}, \href{https://github.com/Paul-Lez/HodgeConjecture/blob/main/HodgeConjecture/Definitions/AlgebraicTopology/SingularCochainSheaf.lean#L226}{\texttt{HodgeConjecture/Definitions/AlgebraicTopology/SingularCochainSheaf.lean:226}}.
\end{definition}

\begin{definition}[AlgebraicTopology.Singular.constantCoefficientPresheaf]\label{def:AlgebraicTopology.Singular.constantCoefficientPresheaf}
\lean{AlgebraicTopology.Singular.constantCoefficientPresheaf}\leanok
The constant presheaf with value the additive group of \texttt{R}.

\noindent\emph{Lean:} \texttt{def AlgebraicTopology.Singular.constantCoefficientPresheaf}, \href{https://github.com/Paul-Lez/HodgeConjecture/blob/main/HodgeConjecture/Definitions/AlgebraicTopology/SingularCochainSheaf.lean#L246}{\texttt{HodgeConjecture/Definitions/AlgebraicTopology/SingularCochainSheaf.lean:246}}.
\end{definition}

\begin{definition}[AlgebraicTopology.Singular.constantSingularZeroCochain]\label{def:AlgebraicTopology.Singular.constantSingularZeroCochain}
\lean{AlgebraicTopology.Singular.constantSingularZeroCochain}\leanok
\uses{def:AlgebraicTopology.Singular.OpenCochains, def:AlgebraicTopology.Singular.openZeroAugmentation}
A scalar defines the corresponding constant singular zero-cochain.

\noindent\emph{Lean:} \texttt{def AlgebraicTopology.Singular.constantSingularZeroCochain}, \href{https://github.com/Paul-Lez/HodgeConjecture/blob/main/HodgeConjecture/Definitions/AlgebraicTopology/SingularCochainSheaf.lean#L250}{\texttt{HodgeConjecture/Definitions/AlgebraicTopology/SingularCochainSheaf.lean:250}}.
\end{definition}

\begin{definition}[AlgebraicTopology.Singular.constantsToSingularCochainZero]\label{def:AlgebraicTopology.Singular.constantsToSingularCochainZero}
\lean{AlgebraicTopology.Singular.constantsToSingularCochainZero}\leanok
\uses{def:AlgebraicTopology.Singular.constantCoefficientPresheaf, def:AlgebraicTopology.Singular.constantSingularZeroCochain, def:AlgebraicTopology.Singular.singularCochainPresheaf}
Constant functions form singular zero-cochains, naturally under restriction.

\noindent\emph{Lean:} \texttt{def AlgebraicTopology.Singular.constantsToSingularCochainZero}, \href{https://github.com/Paul-Lez/HodgeConjecture/blob/main/HodgeConjecture/Definitions/AlgebraicTopology/SingularCochainSheaf.lean#L257}{\texttt{HodgeConjecture/Definitions/AlgebraicTopology/SingularCochainSheaf.lean:257}}.
\end{definition}

\begin{definition}[AlgebraicTopology.Singular.constantCoefficientSheaf]\label{def:AlgebraicTopology.Singular.constantCoefficientSheaf}
\lean{AlgebraicTopology.Singular.constantCoefficientSheaf}\leanok
The constant sheaf with value the additive group of \texttt{R}.

\noindent\emph{Lean:} \texttt{def AlgebraicTopology.Singular.constantCoefficientSheaf}, \href{https://github.com/Paul-Lez/HodgeConjecture/blob/main/HodgeConjecture/Definitions/AlgebraicTopology/SingularCochainSheaf.lean#L308}{\texttt{HodgeConjecture/Definitions/AlgebraicTopology/SingularCochainSheaf.lean:308}}.
\end{definition}

\begin{definition}[AlgebraicTopology.Singular.constantsToSingularCochainZeroSheaf]\label{def:AlgebraicTopology.Singular.constantsToSingularCochainZeroSheaf}
\lean{AlgebraicTopology.Singular.constantsToSingularCochainZeroSheaf}\leanok
\uses{def:AlgebraicTopology.Singular.constantCoefficientSheaf, def:AlgebraicTopology.Singular.constantsToSingularCochainZero, def:AlgebraicTopology.Singular.singularCochainSheaf}
The sheafified inclusion of constants into singular zero-cochains.

\noindent\emph{Lean:} \texttt{def AlgebraicTopology.Singular.constantsToSingularCochainZeroSheaf}, \href{https://github.com/Paul-Lez/HodgeConjecture/blob/main/HodgeConjecture/Definitions/AlgebraicTopology/SingularCochainSheaf.lean#L313}{\texttt{HodgeConjecture/Definitions/AlgebraicTopology/SingularCochainSheaf.lean:313}}.
\end{definition}

\begin{definition}[AlgebraicTopology.Singular.constantsToSingularCochainPresheafShortComplex]\label{def:AlgebraicTopology.Singular.constantsToSingularCochainPresheafShortComplex}
\lean{AlgebraicTopology.Singular.constantsToSingularCochainPresheafShortComplex}\leanok
\uses{def:AlgebraicTopology.Singular.constantsToSingularCochainZero, def:AlgebraicTopology.Singular.singularCochainCoboundary}
The augmented singular zero-cochain short complex before sheafification.

\noindent\emph{Lean:} \texttt{def AlgebraicTopology.Singular.constantsToSingularCochainPresheafShortComplex}, \href{https://github.com/Paul-Lez/HodgeConjecture/blob/main/HodgeConjecture/Definitions/AlgebraicTopology/SingularCochainSheaf.lean#L331}{\texttt{HodgeConjecture/Definitions/AlgebraicTopology/SingularCochainSheaf.lean:331}}.
\end{definition}

\begin{definition}[AlgebraicTopology.Singular.constantsToSingularCochainSheafShortComplex]\label{def:AlgebraicTopology.Singular.constantsToSingularCochainSheafShortComplex}
\lean{AlgebraicTopology.Singular.constantsToSingularCochainSheafShortComplex}\leanok
\uses{def:AlgebraicTopology.Singular.constantsToSingularCochainZeroSheaf, def:AlgebraicTopology.Singular.singularCochainSheafCoboundary}
The augmented singular zero-cochain short complex after sheafification.

\noindent\emph{Lean:} \texttt{def AlgebraicTopology.Singular.constantsToSingularCochainSheafShortComplex}, \href{https://github.com/Paul-Lez/HodgeConjecture/blob/main/HodgeConjecture/Definitions/AlgebraicTopology/SingularCochainSheaf.lean#L338}{\texttt{HodgeConjecture/Definitions/AlgebraicTopology/SingularCochainSheaf.lean:338}}.
\end{definition}

\begin{definition}[AlgebraicTopology.Singular.constantsToSingularCochainShortComplexSheafificationUnit]\label{def:AlgebraicTopology.Singular.constantsToSingularCochainShortComplexSheafificationUnit}
\lean{AlgebraicTopology.Singular.constantsToSingularCochainShortComplexSheafificationUnit}\leanok
\uses{def:AlgebraicTopology.Singular.constantsToSingularCochainPresheafShortComplex, def:AlgebraicTopology.Singular.constantsToSingularCochainSheafShortComplex}
The sheafification unit between the augmented presheaf and sheaf short complexes.

\noindent\emph{Lean:} \texttt{def AlgebraicTopology.Singular.constantsToSingularCochainShortComplexSheafificationUnit}, \href{https://github.com/Paul-Lez/HodgeConjecture/blob/main/HodgeConjecture/Definitions/AlgebraicTopology/SingularCochainSheaf.lean#L346}{\texttt{HodgeConjecture/Definitions/AlgebraicTopology/SingularCochainSheaf.lean:346}}.
\end{definition}

\begin{definition}[AlgebraicTopology.Singular.constantsToSingularCochainSheafComplex]\label{def:AlgebraicTopology.Singular.constantsToSingularCochainSheafComplex}
\lean{AlgebraicTopology.Singular.constantsToSingularCochainSheafComplex}\leanok
\uses{def:AlgebraicTopology.Singular.constantsToSingularCochainZeroSheaf, def:AlgebraicTopology.Singular.singularCochainSheafCoboundary, def:AlgebraicTopology.Singular.singularCochainSheafComplex}
The canonical comparison from the constant sheaf complex to the singular-cochain sheaf
complex.

\noindent\emph{Lean:} \texttt{def AlgebraicTopology.Singular.constantsToSingularCochainSheafComplex}, \href{https://github.com/Paul-Lez/HodgeConjecture/blob/main/HodgeConjecture/Definitions/AlgebraicTopology/SingularCochainSheaf.lean#L359}{\texttt{HodgeConjecture/Definitions/AlgebraicTopology/SingularCochainSheaf.lean:359}}.
\end{definition}

\section{Global sections of bounded-below exact flasque complexes}\label{sec:HodgeConjecture.Definitions.AlgebraicTopology.BoundedBelowFlasqueComplex}
\noindent\emph{File:} \href{https://github.com/Paul-Lez/HodgeConjecture/blob/main/HodgeConjecture/Definitions/AlgebraicTopology/BoundedBelowFlasqueComplex.lean}{\texttt{HodgeConjecture/Definitions/AlgebraicTopology/BoundedBelowFlasqueComplex.lean}}.

An exact bounded-below cochain complex of flasque sheaves remains exact after taking global
sections.  The proof is elementary: starting at the lower bound, the cycle sheaves are flasque
by induction through their short exact sequences.  Global sections are then exact on each of
those sequences.

This is the acyclic-complex lemma needed to compare a bounded-below flasque resolution with a
termwise-injective replacement.  It does not assume or invoke a hypercohomology spectral
sequence.

\begin{definition}[TopCat.Sheaf.IsFlasque.BoundedBelowComplex.cyclesShortComplex]\label{def:TopCat.Sheaf.IsFlasque.BoundedBelowComplex.cyclesShortComplex}
\lean{TopCat.Sheaf.IsFlasque.BoundedBelowComplex.cyclesShortComplex}\leanok
The short exact sequence from cycles in degree \texttt{i}, through the degree-\texttt{i} term, to cycles
in degree \texttt{i + 1}.

\noindent\emph{Lean:} \texttt{def TopCat.Sheaf.IsFlasque.BoundedBelowComplex.cyclesShortComplex}, \href{https://github.com/Paul-Lez/HodgeConjecture/blob/main/HodgeConjecture/Definitions/AlgebraicTopology/BoundedBelowFlasqueComplex.lean#L48}{\texttt{HodgeConjecture/Definitions/AlgebraicTopology/BoundedBelowFlasqueComplex.lean:48}}.
\end{definition}

\begin{definition}[TopCat.Sheaf.IsFlasque.BoundedBelowComplex.globalSectionsFunctor]\label{def:TopCat.Sheaf.IsFlasque.BoundedBelowComplex.globalSectionsFunctor}
\lean{TopCat.Sheaf.IsFlasque.BoundedBelowComplex.globalSectionsFunctor}\leanok
Evaluation of a sheaf on the top open subset, viewed as a functor. This is
\texttt{CategoryTheory.sheafSections} at the top open, stated with the domain spelled \texttt{TopCat.Sheaf}.

\noindent\emph{Lean:} \texttt{def TopCat.Sheaf.IsFlasque.BoundedBelowComplex.globalSectionsFunctor}, \href{https://github.com/Paul-Lez/HodgeConjecture/blob/main/HodgeConjecture/Definitions/AlgebraicTopology/BoundedBelowFlasqueComplex.lean#L56}{\texttt{HodgeConjecture/Definitions/AlgebraicTopology/BoundedBelowFlasqueComplex.lean:56}}.
\end{definition}

\begin{definition}[TopCat.Sheaf.IsFlasque.BoundedBelowComplex.globalSectionsComplex]\label{def:TopCat.Sheaf.IsFlasque.BoundedBelowComplex.globalSectionsComplex}
\lean{TopCat.Sheaf.IsFlasque.BoundedBelowComplex.globalSectionsComplex}\leanok
\uses{def:TopCat.Sheaf.IsFlasque.BoundedBelowComplex.globalSectionsFunctor}
The integer-indexed cochain complex obtained by evaluating a sheaf complex on the top open
subset.

\noindent\emph{Lean:} \texttt{def TopCat.Sheaf.IsFlasque.BoundedBelowComplex.globalSectionsComplex}, \href{https://github.com/Paul-Lez/HodgeConjecture/blob/main/HodgeConjecture/Definitions/AlgebraicTopology/BoundedBelowFlasqueComplex.lean#L79}{\texttt{HodgeConjecture/Definitions/AlgebraicTopology/BoundedBelowFlasqueComplex.lean:79}}.
\end{definition}

\section{Global and locally defined singular cochains}\label{sec:HodgeConjecture.Definitions.AlgebraicTopology.SingularSubdivisionCochainSheaf}
\noindent\emph{File:} \href{https://github.com/Paul-Lez/HodgeConjecture/blob/main/HodgeConjecture/Definitions/AlgebraicTopology/SingularSubdivisionCochainSheaf.lean}{\texttt{HodgeConjecture/Definitions/AlgebraicTopology/SingularSubdivisionCochainSheaf.lean}}.

This file proves the global comparison used in the singular-to-sheaf cohomology theorem.
Ordinary singular cochains agree with cochains on the top open subset. Subdivision kills the
locally zero kernel of the first plus construction. On paracompact Hausdorff spaces, a locally
finite closed refinement also makes the second plus map surjective on global sections. The
double-plus complex is then identified with Mathlib's chosen degreewise sheafification.

The first plus construction is only an intermediate separated presheaf; no sheaf claim is made
for it. The final comparison uses the double-plus sheaf and the canonical isomorphism to the
chosen singular-cochain sheaf complex.

\begin{definition}[AlgebraicTopology.Singular.coveringSieveOpenFamilyOn]\label{def:AlgebraicTopology.Singular.coveringSieveOpenFamilyOn}
\lean{AlgebraicTopology.Singular.coveringSieveOpenFamilyOn}\leanok
The domains of a covering sieve, regarded as open subsets of the covered open set.

\noindent\emph{Lean:} \texttt{def AlgebraicTopology.Singular.coveringSieveOpenFamilyOn}, \href{https://github.com/Paul-Lez/HodgeConjecture/blob/main/HodgeConjecture/Definitions/AlgebraicTopology/SingularSubdivisionCochainSheaf.lean#L51}{\texttt{HodgeConjecture/Definitions/AlgebraicTopology/SingularSubdivisionCochainSheaf.lean:51}}.
\end{definition}

\begin{definition}[AlgebraicTopology.Singular.TopOpenSingularChainComplex]\label{def:AlgebraicTopology.Singular.TopOpenSingularChainComplex}
\lean{AlgebraicTopology.Singular.TopOpenSingularChainComplex}\leanok
\uses{def:AlgebraicTopology.Singular.openSingularChainComplexFunctor}
The singular chain complex of the top open subset of \texttt{X}.

\noindent\emph{Lean:} \texttt{abbrev AlgebraicTopology.Singular.TopOpenSingularChainComplex}, \href{https://github.com/Paul-Lez/HodgeConjecture/blob/main/HodgeConjecture/Definitions/AlgebraicTopology/SingularSubdivisionCochainSheaf.lean#L118}{\texttt{HodgeConjecture/Definitions/AlgebraicTopology/SingularSubdivisionCochainSheaf.lean:118}}.
\end{definition}

\begin{definition}[AlgebraicTopology.Singular.topOpenSingularChainComplexIso]\label{def:AlgebraicTopology.Singular.topOpenSingularChainComplexIso}
\lean{AlgebraicTopology.Singular.topOpenSingularChainComplexIso}\leanok
\uses{def:AlgebraicTopology.Singular.SingularChainComplex, def:AlgebraicTopology.Singular.TopOpenSingularChainComplex}
The inclusion of the top open subset into \texttt{X} identifies their singular chain complexes.

\noindent\emph{Lean:} \texttt{def AlgebraicTopology.Singular.topOpenSingularChainComplexIso}, \href{https://github.com/Paul-Lez/HodgeConjecture/blob/main/HodgeConjecture/Definitions/AlgebraicTopology/SingularSubdivisionCochainSheaf.lean#L122}{\texttt{HodgeConjecture/Definitions/AlgebraicTopology/SingularSubdivisionCochainSheaf.lean:122}}.
\end{definition}

\begin{definition}[AlgebraicTopology.Singular.singularCochainComplexIsoTopOpen]\label{def:AlgebraicTopology.Singular.singularCochainComplexIsoTopOpen}
\lean{AlgebraicTopology.Singular.singularCochainComplexIsoTopOpen}\leanok
\uses{def:AlgebraicTopology.Singular.topOpenSingularChainComplexIso, def:HomologicalComplex.linearDualIso}
Ordinary singular cochains and top-open singular cochains are isomorphic as cochain
complexes.

\noindent\emph{Lean:} \texttt{def AlgebraicTopology.Singular.singularCochainComplexIsoTopOpen}, \href{https://github.com/Paul-Lez/HodgeConjecture/blob/main/HodgeConjecture/Definitions/AlgebraicTopology/SingularSubdivisionCochainSheaf.lean#L128}{\texttt{HodgeConjecture/Definitions/AlgebraicTopology/SingularSubdivisionCochainSheaf.lean:128}}.
\end{definition}

\begin{definition}[AlgebraicTopology.Singular.globalRawSingularCochainComplex]\label{def:AlgebraicTopology.Singular.globalRawSingularCochainComplex}
\lean{AlgebraicTopology.Singular.globalRawSingularCochainComplex}\leanok
\uses{def:AlgebraicTopology.Singular.singularCochainPresheafComplex}
The singular-cochain presheaf complex evaluated on the top open subset.

\noindent\emph{Lean:} \texttt{def AlgebraicTopology.Singular.globalRawSingularCochainComplex}, \href{https://github.com/Paul-Lez/HodgeConjecture/blob/main/HodgeConjecture/Definitions/AlgebraicTopology/SingularSubdivisionCochainSheaf.lean#L143}{\texttt{HodgeConjecture/Definitions/AlgebraicTopology/SingularSubdivisionCochainSheaf.lean:143}}.
\end{definition}

\begin{definition}[AlgebraicTopology.Singular.topOpenForgottenSingularCochainComplex]\label{def:AlgebraicTopology.Singular.topOpenForgottenSingularCochainComplex}
\lean{AlgebraicTopology.Singular.topOpenForgottenSingularCochainComplex}\leanok
\uses{def:AlgebraicTopology.Singular.TopOpenSingularChainComplex, def:HomologicalComplex.linearDualCochainComplex}
Top-open singular cochains, with their scalar structure forgotten.

\noindent\emph{Lean:} \texttt{def AlgebraicTopology.Singular.topOpenForgottenSingularCochainComplex}, \href{https://github.com/Paul-Lez/HodgeConjecture/blob/main/HodgeConjecture/Definitions/AlgebraicTopology/SingularSubdivisionCochainSheaf.lean#L148}{\texttt{HodgeConjecture/Definitions/AlgebraicTopology/SingularSubdivisionCochainSheaf.lean:148}}.
\end{definition}

\begin{definition}[AlgebraicTopology.Singular.globalRawSingularCochainAddEquiv]\label{def:AlgebraicTopology.Singular.globalRawSingularCochainAddEquiv}
\lean{AlgebraicTopology.Singular.globalRawSingularCochainAddEquiv}\leanok
\uses{def:AlgebraicTopology.Singular.globalRawSingularCochainComplex, def:AlgebraicTopology.Singular.topOpenForgottenSingularCochainComplex}
Evaluation of the cochain presheaf and forgetting the scalar structure on top-open cochains
have the same underlying additive groups.

\noindent\emph{Lean:} \texttt{def AlgebraicTopology.Singular.globalRawSingularCochainAddEquiv}, \href{https://github.com/Paul-Lez/HodgeConjecture/blob/main/HodgeConjecture/Definitions/AlgebraicTopology/SingularSubdivisionCochainSheaf.lean#L155}{\texttt{HodgeConjecture/Definitions/AlgebraicTopology/SingularSubdivisionCochainSheaf.lean:155}}.
\end{definition}

\begin{definition}[AlgebraicTopology.Singular.globalRawSingularCochainXIso]\label{def:AlgebraicTopology.Singular.globalRawSingularCochainXIso}
\lean{AlgebraicTopology.Singular.globalRawSingularCochainXIso}\leanok
\uses{def:AlgebraicTopology.Singular.globalRawSingularCochainAddEquiv}
The degreewise additive-group isomorphism between the two presentations of top-open
cochains.

\noindent\emph{Lean:} \texttt{def AlgebraicTopology.Singular.globalRawSingularCochainXIso}, \href{https://github.com/Paul-Lez/HodgeConjecture/blob/main/HodgeConjecture/Definitions/AlgebraicTopology/SingularSubdivisionCochainSheaf.lean#L166}{\texttt{HodgeConjecture/Definitions/AlgebraicTopology/SingularSubdivisionCochainSheaf.lean:166}}.
\end{definition}

\begin{definition}[AlgebraicTopology.Singular.globalRawSingularCochainComplexIso]\label{def:AlgebraicTopology.Singular.globalRawSingularCochainComplexIso}
\lean{AlgebraicTopology.Singular.globalRawSingularCochainComplexIso}\leanok
\uses{def:AlgebraicTopology.Singular.globalRawSingularCochainXIso}
Evaluating the raw cochain presheaf on the top open agrees, as a complex of additive groups,
with forgetting the scalar structure on the top-open linear-dual complex.

\noindent\emph{Lean:} \texttt{def AlgebraicTopology.Singular.globalRawSingularCochainComplexIso}, \href{https://github.com/Paul-Lez/HodgeConjecture/blob/main/HodgeConjecture/Definitions/AlgebraicTopology/SingularSubdivisionCochainSheaf.lean#L174}{\texttt{HodgeConjecture/Definitions/AlgebraicTopology/SingularSubdivisionCochainSheaf.lean:174}}.
\end{definition}

\begin{definition}[AlgebraicTopology.Singular.singularCochainPlusPresheafComplex]\label{def:AlgebraicTopology.Singular.singularCochainPlusPresheafComplex}
\lean{AlgebraicTopology.Singular.singularCochainPlusPresheafComplex}\leanok
\uses{def:AlgebraicTopology.Singular.singularCochainPresheafComplex}
The singular-cochain presheaf complex after applying the first plus construction degreewise.
This is a complex of presheaves, not a complex of sheaves.

\noindent\emph{Lean:} \texttt{def AlgebraicTopology.Singular.singularCochainPlusPresheafComplex}, \href{https://github.com/Paul-Lez/HodgeConjecture/blob/main/HodgeConjecture/Definitions/AlgebraicTopology/SingularSubdivisionCochainSheaf.lean#L201}{\texttt{HodgeConjecture/Definitions/AlgebraicTopology/SingularSubdivisionCochainSheaf.lean:201}}.
\end{definition}

\begin{definition}[AlgebraicTopology.Singular.singularCochainToPlusPresheafComplex]\label{def:AlgebraicTopology.Singular.singularCochainToPlusPresheafComplex}
\lean{AlgebraicTopology.Singular.singularCochainToPlusPresheafComplex}\leanok
\uses{def:AlgebraicTopology.Singular.singularCochainPlusPresheafComplex}
The degreewise first-plus unit on the singular-cochain presheaf complex.

\noindent\emph{Lean:} \texttt{def AlgebraicTopology.Singular.singularCochainToPlusPresheafComplex}, \href{https://github.com/Paul-Lez/HodgeConjecture/blob/main/HodgeConjecture/Definitions/AlgebraicTopology/SingularSubdivisionCochainSheaf.lean#L208}{\texttt{HodgeConjecture/Definitions/AlgebraicTopology/SingularSubdivisionCochainSheaf.lean:208}}.
\end{definition}

\begin{definition}[AlgebraicTopology.Singular.singularCochainPlusPlusPresheafComplex]\label{def:AlgebraicTopology.Singular.singularCochainPlusPlusPresheafComplex}
\lean{AlgebraicTopology.Singular.singularCochainPlusPlusPresheafComplex}\leanok
\uses{def:AlgebraicTopology.Singular.singularCochainPresheafComplex}
The singular-cochain presheaf complex after applying the plus construction twice
degreewise. Every term of this complex is a sheaf.

\noindent\emph{Lean:} \texttt{def AlgebraicTopology.Singular.singularCochainPlusPlusPresheafComplex}, \href{https://github.com/Paul-Lez/HodgeConjecture/blob/main/HodgeConjecture/Definitions/AlgebraicTopology/SingularSubdivisionCochainSheaf.lean#L221}{\texttt{HodgeConjecture/Definitions/AlgebraicTopology/SingularSubdivisionCochainSheaf.lean:221}}.
\end{definition}

\begin{definition}[AlgebraicTopology.Singular.singularCochainPlusToPlusPlusPresheafComplex]\label{def:AlgebraicTopology.Singular.singularCochainPlusToPlusPlusPresheafComplex}
\lean{AlgebraicTopology.Singular.singularCochainPlusToPlusPlusPresheafComplex}\leanok
\uses{def:AlgebraicTopology.Singular.singularCochainPlusPlusPresheafComplex, def:AlgebraicTopology.Singular.singularCochainPlusPresheafComplex}
The degreewise second-plus unit on the singular-cochain complex.

\noindent\emph{Lean:} \texttt{def AlgebraicTopology.Singular.singularCochainPlusToPlusPlusPresheafComplex}, \href{https://github.com/Paul-Lez/HodgeConjecture/blob/main/HodgeConjecture/Definitions/AlgebraicTopology/SingularSubdivisionCochainSheaf.lean#L229}{\texttt{HodgeConjecture/Definitions/AlgebraicTopology/SingularSubdivisionCochainSheaf.lean:229}}.
\end{definition}

\begin{definition}[AlgebraicTopology.Singular.singularCochainToPlusPlusPresheafComplex]\label{def:AlgebraicTopology.Singular.singularCochainToPlusPlusPresheafComplex}
\lean{AlgebraicTopology.Singular.singularCochainToPlusPlusPresheafComplex}\leanok
\uses{def:AlgebraicTopology.Singular.singularCochainPlusToPlusPlusPresheafComplex, def:AlgebraicTopology.Singular.singularCochainToPlusPresheafComplex}
The degreewise double-plus comparison from raw singular cochains.

\noindent\emph{Lean:} \texttt{def AlgebraicTopology.Singular.singularCochainToPlusPlusPresheafComplex}, \href{https://github.com/Paul-Lez/HodgeConjecture/blob/main/HodgeConjecture/Definitions/AlgebraicTopology/SingularSubdivisionCochainSheaf.lean#L247}{\texttt{HodgeConjecture/Definitions/AlgebraicTopology/SingularSubdivisionCochainSheaf.lean:247}}.
\end{definition}

\begin{definition}[AlgebraicTopology.Singular.forgottenSingularCochainSheafComplex]\label{def:AlgebraicTopology.Singular.forgottenSingularCochainSheafComplex}
\lean{AlgebraicTopology.Singular.forgottenSingularCochainSheafComplex}\leanok
\uses{def:AlgebraicTopology.Singular.singularCochainSheafComplex}
The underlying presheaf complex of the chosen degreewise sheafification.

\noindent\emph{Lean:} \texttt{def AlgebraicTopology.Singular.forgottenSingularCochainSheafComplex}, \href{https://github.com/Paul-Lez/HodgeConjecture/blob/main/HodgeConjecture/Definitions/AlgebraicTopology/SingularSubdivisionCochainSheaf.lean#L269}{\texttt{HodgeConjecture/Definitions/AlgebraicTopology/SingularSubdivisionCochainSheaf.lean:269}}.
\end{definition}

\begin{definition}[AlgebraicTopology.Singular.singularCochainPlusPlusPresheafComplexIsoSheafComplex]\label{def:AlgebraicTopology.Singular.singularCochainPlusPlusPresheafComplexIsoSheafComplex}
\lean{AlgebraicTopology.Singular.singularCochainPlusPlusPresheafComplexIsoSheafComplex}\leanok
\uses{def:AlgebraicTopology.Singular.forgottenSingularCochainSheafComplex, def:AlgebraicTopology.Singular.singularCochainPlusPlusPresheafComplex}
Double-plus sheafification agrees with Mathlib's chosen sheafification, compatibly with the
singular coboundary.

\noindent\emph{Lean:} \texttt{def AlgebraicTopology.Singular.singularCochainPlusPlusPresheafComplexIsoSheafComplex}, \href{https://github.com/Paul-Lez/HodgeConjecture/blob/main/HodgeConjecture/Definitions/AlgebraicTopology/SingularSubdivisionCochainSheaf.lean#L275}{\texttt{HodgeConjecture/Definitions/AlgebraicTopology/SingularSubdivisionCochainSheaf.lean:275}}.
\end{definition}

\begin{definition}[AlgebraicTopology.Singular.globalSingularCochainPlusComplex]\label{def:AlgebraicTopology.Singular.globalSingularCochainPlusComplex}
\lean{AlgebraicTopology.Singular.globalSingularCochainPlusComplex}\leanok
\uses{def:AlgebraicTopology.Singular.singularCochainPlusPresheafComplex}
Evaluation of the first-plus singular-cochain complex on the top open subset.

\noindent\emph{Lean:} \texttt{def AlgebraicTopology.Singular.globalSingularCochainPlusComplex}, \href{https://github.com/Paul-Lez/HodgeConjecture/blob/main/HodgeConjecture/Definitions/AlgebraicTopology/SingularSubdivisionCochainSheaf.lean#L303}{\texttt{HodgeConjecture/Definitions/AlgebraicTopology/SingularSubdivisionCochainSheaf.lean:303}}.
\end{definition}

\begin{definition}[AlgebraicTopology.Singular.topOpenToGlobalSingularCochainPlusComplex]\label{def:AlgebraicTopology.Singular.topOpenToGlobalSingularCochainPlusComplex}
\lean{AlgebraicTopology.Singular.topOpenToGlobalSingularCochainPlusComplex}\leanok
\uses{def:AlgebraicTopology.Singular.globalRawSingularCochainComplex, def:AlgebraicTopology.Singular.globalSingularCochainPlusComplex, def:AlgebraicTopology.Singular.singularCochainToPlusPresheafComplex}
Evaluation on the top open subset of the degreewise first-plus comparison.

\noindent\emph{Lean:} \texttt{def AlgebraicTopology.Singular.topOpenToGlobalSingularCochainPlusComplex}, \href{https://github.com/Paul-Lez/HodgeConjecture/blob/main/HodgeConjecture/Definitions/AlgebraicTopology/SingularSubdivisionCochainSheaf.lean#L308}{\texttt{HodgeConjecture/Definitions/AlgebraicTopology/SingularSubdivisionCochainSheaf.lean:308}}.
\end{definition}

\begin{definition}[AlgebraicTopology.Singular.globalSingularCochainPlusPlusComplex]\label{def:AlgebraicTopology.Singular.globalSingularCochainPlusPlusComplex}
\lean{AlgebraicTopology.Singular.globalSingularCochainPlusPlusComplex}\leanok
\uses{def:AlgebraicTopology.Singular.singularCochainPlusPlusPresheafComplex}
Evaluation of the double-plus singular-cochain complex on the top open subset.

\noindent\emph{Lean:} \texttt{def AlgebraicTopology.Singular.globalSingularCochainPlusPlusComplex}, \href{https://github.com/Paul-Lez/HodgeConjecture/blob/main/HodgeConjecture/Definitions/AlgebraicTopology/SingularSubdivisionCochainSheaf.lean#L315}{\texttt{HodgeConjecture/Definitions/AlgebraicTopology/SingularSubdivisionCochainSheaf.lean:315}}.
\end{definition}

\begin{definition}[AlgebraicTopology.Singular.globalSingularCochainSheafComplex]\label{def:AlgebraicTopology.Singular.globalSingularCochainSheafComplex}
\lean{AlgebraicTopology.Singular.globalSingularCochainSheafComplex}\leanok
\uses{def:AlgebraicTopology.Singular.forgottenSingularCochainSheafComplex}
Global sections of the chosen singular-cochain sheaf complex.

\noindent\emph{Lean:} \texttt{def AlgebraicTopology.Singular.globalSingularCochainSheafComplex}, \href{https://github.com/Paul-Lez/HodgeConjecture/blob/main/HodgeConjecture/Definitions/AlgebraicTopology/SingularSubdivisionCochainSheaf.lean#L320}{\texttt{HodgeConjecture/Definitions/AlgebraicTopology/SingularSubdivisionCochainSheaf.lean:320}}.
\end{definition}

\begin{definition}[AlgebraicTopology.Singular.globalSingularCochainPlusPlusComplexIsoSheafComplex]\label{def:AlgebraicTopology.Singular.globalSingularCochainPlusPlusComplexIsoSheafComplex}
\lean{AlgebraicTopology.Singular.globalSingularCochainPlusPlusComplexIsoSheafComplex}\leanok
\uses{def:AlgebraicTopology.Singular.globalSingularCochainPlusPlusComplex, def:AlgebraicTopology.Singular.globalSingularCochainSheafComplex, def:AlgebraicTopology.Singular.singularCochainPlusPlusPresheafComplexIsoSheafComplex}
The double-plus model and the chosen singular-cochain sheaf complex have isomorphic global
section complexes.

\noindent\emph{Lean:} \texttt{def AlgebraicTopology.Singular.globalSingularCochainPlusPlusComplexIsoSheafComplex}, \href{https://github.com/Paul-Lez/HodgeConjecture/blob/main/HodgeConjecture/Definitions/AlgebraicTopology/SingularSubdivisionCochainSheaf.lean#L325}{\texttt{HodgeConjecture/Definitions/AlgebraicTopology/SingularSubdivisionCochainSheaf.lean:325}}.
\end{definition}

\begin{definition}[AlgebraicTopology.Singular.globalSingularCochainPlusToPlusPlusComplex]\label{def:AlgebraicTopology.Singular.globalSingularCochainPlusToPlusPlusComplex}
\lean{AlgebraicTopology.Singular.globalSingularCochainPlusToPlusPlusComplex}\leanok
\uses{def:AlgebraicTopology.Singular.globalSingularCochainPlusComplex, def:AlgebraicTopology.Singular.globalSingularCochainPlusPlusComplex, def:AlgebraicTopology.Singular.singularCochainPlusToPlusPlusPresheafComplex}
Evaluation on the top open subset of the degreewise second-plus comparison.

\noindent\emph{Lean:} \texttt{def AlgebraicTopology.Singular.globalSingularCochainPlusToPlusPlusComplex}, \href{https://github.com/Paul-Lez/HodgeConjecture/blob/main/HodgeConjecture/Definitions/AlgebraicTopology/SingularSubdivisionCochainSheaf.lean#L334}{\texttt{HodgeConjecture/Definitions/AlgebraicTopology/SingularSubdivisionCochainSheaf.lean:334}}.
\end{definition}

\begin{definition}[AlgebraicTopology.Singular.topOpenToGlobalSingularCochainPlusPlusComplex]\label{def:AlgebraicTopology.Singular.topOpenToGlobalSingularCochainPlusPlusComplex}
\lean{AlgebraicTopology.Singular.topOpenToGlobalSingularCochainPlusPlusComplex}\leanok
\uses{def:AlgebraicTopology.Singular.globalSingularCochainPlusToPlusPlusComplex, def:AlgebraicTopology.Singular.topOpenToGlobalSingularCochainPlusComplex}
Evaluation on the top open subset of the degreewise double-plus comparison.

\noindent\emph{Lean:} \texttt{def AlgebraicTopology.Singular.topOpenToGlobalSingularCochainPlusPlusComplex}, \href{https://github.com/Paul-Lez/HodgeConjecture/blob/main/HodgeConjecture/Definitions/AlgebraicTopology/SingularSubdivisionCochainSheaf.lean#L341}{\texttt{HodgeConjecture/Definitions/AlgebraicTopology/SingularSubdivisionCochainSheaf.lean:341}}.
\end{definition}

\begin{definition}[AlgebraicTopology.Singular.topOpenToGlobalSingularCochainSheafComplex]\label{def:AlgebraicTopology.Singular.topOpenToGlobalSingularCochainSheafComplex}
\lean{AlgebraicTopology.Singular.topOpenToGlobalSingularCochainSheafComplex}\leanok
\uses{def:AlgebraicTopology.Singular.globalSingularCochainPlusPlusComplexIsoSheafComplex, def:AlgebraicTopology.Singular.topOpenToGlobalSingularCochainPlusPlusComplex}
The comparison from ordinary singular cochains to global sections of the chosen
singular-cochain sheaf complex.

\noindent\emph{Lean:} \texttt{def AlgebraicTopology.Singular.topOpenToGlobalSingularCochainSheafComplex}, \href{https://github.com/Paul-Lez/HodgeConjecture/blob/main/HodgeConjecture/Definitions/AlgebraicTopology/SingularSubdivisionCochainSheaf.lean#L348}{\texttt{HodgeConjecture/Definitions/AlgebraicTopology/SingularSubdivisionCochainSheaf.lean:348}}.
\end{definition}

\begin{definition}[AlgebraicTopology.Singular.coveringSieveOpenFamily]\label{def:AlgebraicTopology.Singular.coveringSieveOpenFamily}
\lean{AlgebraicTopology.Singular.coveringSieveOpenFamily}\leanok
The family of open subsets appearing as the domains of the arrows in a covering sieve of the
top open subset.

\noindent\emph{Lean:} \texttt{def AlgebraicTopology.Singular.coveringSieveOpenFamily}, \href{https://github.com/Paul-Lez/HodgeConjecture/blob/main/HodgeConjecture/Definitions/AlgebraicTopology/SingularSubdivisionCochainSheaf.lean#L470}{\texttt{HodgeConjecture/Definitions/AlgebraicTopology/SingularSubdivisionCochainSheaf.lean:470}}.
\end{definition}

\begin{definition}[AlgebraicTopology.Singular.rationalCochainRestrictionToCoverSmall]\label{def:AlgebraicTopology.Singular.rationalCochainRestrictionToCoverSmall}
\lean{AlgebraicTopology.Singular.rationalCochainRestrictionToCoverSmall}\leanok
\uses{def:AlgebraicTopology.Singular.coverSmallRationalSingularChainInclusion, def:HomologicalComplex.linearDualMap}
Restriction of rational singular cochains to chains subordinate to a family of subsets.

\noindent\emph{Lean:} \texttt{def AlgebraicTopology.Singular.rationalCochainRestrictionToCoverSmall}, \href{https://github.com/Paul-Lez/HodgeConjecture/blob/main/HodgeConjecture/Definitions/AlgebraicTopology/SingularSubdivisionCochainSheaf.lean#L476}{\texttt{HodgeConjecture/Definitions/AlgebraicTopology/SingularSubdivisionCochainSheaf.lean:476}}.
\end{definition}

\begin{definition}[AlgebraicTopology.Singular.coverMemberToSmallRationalSingularChains]\label{def:AlgebraicTopology.Singular.coverMemberToSmallRationalSingularChains}
\lean{AlgebraicTopology.Singular.coverMemberToSmallRationalSingularChains}\leanok
\uses{def:AlgebraicTopology.Singular.CoverSmallRationalSingularChainComplex, def:AlgebraicTopology.Singular.coverMemberToSmallSingularSet}
The chain map from one member of a family into the cover-small rational chains.

\noindent\emph{Lean:} \texttt{def AlgebraicTopology.Singular.coverMemberToSmallRationalSingularChains}, \href{https://github.com/Paul-Lez/HodgeConjecture/blob/main/HodgeConjecture/Definitions/AlgebraicTopology/SingularSubdivisionCochainSheaf.lean#L483}{\texttt{HodgeConjecture/Definitions/AlgebraicTopology/SingularSubdivisionCochainSheaf.lean:483}}.
\end{definition}

\begin{definition}[AlgebraicTopology.Singular.rationalCochainHomotopyEquivCoverSmall]\label{def:AlgebraicTopology.Singular.rationalCochainHomotopyEquivCoverSmall}
\lean{AlgebraicTopology.Singular.rationalCochainHomotopyEquivCoverSmall}\leanok
\uses{def:AlgebraicTopology.Singular.coverSmallRationalChainHomotopyEquiv_of_openCover, def:HomologicalComplex.linearDualHomotopyEquiv}
An open cover gives a homotopy equivalence from all rational cochains to its cover-small
cochains.

\noindent\emph{Lean:} \texttt{def AlgebraicTopology.Singular.rationalCochainHomotopyEquivCoverSmall}, \href{https://github.com/Paul-Lez/HodgeConjecture/blob/main/HodgeConjecture/Definitions/AlgebraicTopology/SingularSubdivisionCochainSheaf.lean#L505}{\texttt{HodgeConjecture/Definitions/AlgebraicTopology/SingularSubdivisionCochainSheaf.lean:505}}.
\end{definition}

\begin{definition}[AlgebraicTopology.Singular.ambientOpen]\label{def:AlgebraicTopology.Singular.ambientOpen}
\lean{AlgebraicTopology.Singular.ambientOpen}\leanok
Regard an open subset of an open subspace as an open subset of the ambient space.

\noindent\emph{Lean:} \texttt{def AlgebraicTopology.Singular.ambientOpen}, \href{https://github.com/Paul-Lez/HodgeConjecture/blob/main/HodgeConjecture/Definitions/AlgebraicTopology/SingularSubdivisionCochainSheaf.lean#L570}{\texttt{HodgeConjecture/Definitions/AlgebraicTopology/SingularSubdivisionCochainSheaf.lean:570}}.
\end{definition}

\begin{definition}[AlgebraicTopology.Singular.openFamilySieve]\label{def:AlgebraicTopology.Singular.openFamilySieve}
\lean{AlgebraicTopology.Singular.openFamilySieve}\leanok
The sieve generated by an indexed family of open subsets of a fixed open set.

\noindent\emph{Lean:} \texttt{def AlgebraicTopology.Singular.openFamilySieve}, \href{https://github.com/Paul-Lez/HodgeConjecture/blob/main/HodgeConjecture/Definitions/AlgebraicTopology/SingularSubdivisionCochainSheaf.lean#L588}{\texttt{HodgeConjecture/Definitions/AlgebraicTopology/SingularSubdivisionCochainSheaf.lean:588}}.
\end{definition}

\section{Relative singular cohomology as a cochain mapping cone}\label{sec:HodgeConjecture.Definitions.AlgebraicTopology.RelativeCochainCone}
\noindent\emph{File:} \href{https://github.com/Paul-Lez/HodgeConjecture/blob/main/HodgeConjecture/Definitions/AlgebraicTopology/RelativeCochainCone.lean}{\texttt{HodgeConjecture/Definitions/AlgebraicTopology/RelativeCochainCone.lean}}.

For a topological pair \texttt{A ⊆ X}, restriction of singular cochains is the algebraic dual of
the inclusion \texttt{C\_*(A) ⟶ C\_*(X)}.  This file proves, without a finite-dimensionality
hypothesis, that the cohomology in degree \texttt{n - 1} of its mapping cone is \texttt{H\^{}n(X, A)}, the
degree-\texttt{n} cohomology of the dual relative cochain complex.

The proof dualizes the degreewise short exact sequence
\texttt{C\_*(A) ⟶ C\_*(X) ⟶ C\_*(X, A)}, extends the resulting cochain complexes by zero to
integer degrees, and compares its rotated triangle with Mathlib's mapping-cone triangle.

\begin{definition}[AlgebraicTopology.Singular.relativeChainShortComplex]\label{def:AlgebraicTopology.Singular.relativeChainShortComplex}
\lean{AlgebraicTopology.Singular.relativeChainShortComplex}\leanok
\uses{def:AlgebraicTopology.Singular.relativeChainProjection}
The short exact sequence of subspace, ambient, and relative singular chains.

\noindent\emph{Lean:} \texttt{def AlgebraicTopology.Singular.relativeChainShortComplex}, \href{https://github.com/Paul-Lez/HodgeConjecture/blob/main/HodgeConjecture/Definitions/AlgebraicTopology/RelativeCochainCone.lean#L49}{\texttt{HodgeConjecture/Definitions/AlgebraicTopology/RelativeCochainCone.lean:49}}.
\end{definition}

\begin{definition}[AlgebraicTopology.Singular.relativeDualCochainShortComplexNat]\label{def:AlgebraicTopology.Singular.relativeDualCochainShortComplexNat}
\lean{AlgebraicTopology.Singular.relativeDualCochainShortComplexNat}\leanok
\uses{def:AlgebraicTopology.Singular.relativeChainProjection, def:HomologicalComplex.linearDualMap}
The dual relative, ambient, and subspace cochain complexes in nonnegative degrees.

\noindent\emph{Lean:} \texttt{def AlgebraicTopology.Singular.relativeDualCochainShortComplexNat}, \href{https://github.com/Paul-Lez/HodgeConjecture/blob/main/HodgeConjecture/Definitions/AlgebraicTopology/RelativeCochainCone.lean#L82}{\texttt{HodgeConjecture/Definitions/AlgebraicTopology/RelativeCochainCone.lean:82}}.
\end{definition}

\begin{definition}[AlgebraicTopology.Singular.relativeDualCochainShortComplexInt]\label{def:AlgebraicTopology.Singular.relativeDualCochainShortComplexInt}
\lean{AlgebraicTopology.Singular.relativeDualCochainShortComplexInt}\leanok
\uses{def:AlgebraicTopology.Singular.relativeDualCochainShortComplexNat}
The dual relative, ambient, and subspace cochain complexes, extended by zero to integer
degrees.

\noindent\emph{Lean:} \texttt{def AlgebraicTopology.Singular.relativeDualCochainShortComplexInt}, \href{https://github.com/Paul-Lez/HodgeConjecture/blob/main/HodgeConjecture/Definitions/AlgebraicTopology/RelativeCochainCone.lean#L125}{\texttt{HodgeConjecture/Definitions/AlgebraicTopology/RelativeCochainCone.lean:125}}.
\end{definition}

\begin{definition}[AlgebraicTopology.Singular.relativeCochainRestrictionInt]\label{def:AlgebraicTopology.Singular.relativeCochainRestrictionInt}
\lean{AlgebraicTopology.Singular.relativeCochainRestrictionInt}\leanok
\uses{def:AlgebraicTopology.Singular.SingularChainComplex, def:AlgebraicTopology.Singular.chainPairFunctor, def:HomologicalComplex.linearDualMap}
Restriction from ambient singular cochains to subspace singular cochains, in integer
degrees.

\noindent\emph{Lean:} \texttt{def AlgebraicTopology.Singular.relativeCochainRestrictionInt}, \href{https://github.com/Paul-Lez/HodgeConjecture/blob/main/HodgeConjecture/Definitions/AlgebraicTopology/RelativeCochainCone.lean#L132}{\texttt{HodgeConjecture/Definitions/AlgebraicTopology/RelativeCochainCone.lean:132}}.
\end{definition}

\begin{definition}[AlgebraicTopology.Singular.relativeDualCochainShortComplexIntEvalIso]\label{def:AlgebraicTopology.Singular.relativeDualCochainShortComplexIntEvalIso}
\lean{AlgebraicTopology.Singular.relativeDualCochainShortComplexIntEvalIso}\leanok
\uses{def:AlgebraicTopology.Singular.relativeDualCochainShortComplexInt}
In a nonnegative degree, evaluation of the integer extension recovers evaluation of the
original nonnegative short complex.

\noindent\emph{Lean:} \texttt{def AlgebraicTopology.Singular.relativeDualCochainShortComplexIntEvalIso}, \href{https://github.com/Paul-Lez/HodgeConjecture/blob/main/HodgeConjecture/Definitions/AlgebraicTopology/RelativeCochainCone.lean#L149}{\texttt{HodgeConjecture/Definitions/AlgebraicTopology/RelativeCochainCone.lean:149}}.
\end{definition}

\section{Comparing an ambient resolution with a resolution on an open subspace}\label{sec:HodgeConjecture.Definitions.AlgebraicTopology.OpenInjectiveResolutionComparison}
\noindent\emph{File:} \href{https://github.com/Paul-Lez/HodgeConjecture/blob/main/HodgeConjecture/Definitions/AlgebraicTopology/OpenInjectiveResolutionComparison.lean}{\texttt{HodgeConjecture/Definitions/AlgebraicTopology/OpenInjectiveResolutionComparison.lean}}.

The comparison is lifted on the open subspace across the actual restricted
augmentation. Exact open restriction and the normalized constant-sheaf
comparison prove that this augmentation is a monic quasi-isomorphism.

\begin{definition}[CochainComplex.liftToInjectiveNat]\label{def:CochainComplex.liftToInjectiveNat}
\lean{CochainComplex.liftToInjectiveNat}\leanok
\uses{def:CochainComplex.liftToInjective}
Strict injective lifting for nonnegative cochain complexes, transported
through the fully faithful extension to integer degrees.

\noindent\emph{Lean:} \texttt{def CochainComplex.liftToInjectiveNat}, \href{https://github.com/Paul-Lez/HodgeConjecture/blob/main/HodgeConjecture/Definitions/AlgebraicTopology/OpenInjectiveResolutionComparison.lean#L50}{\texttt{HodgeConjecture/Definitions/AlgebraicTopology/OpenInjectiveResolutionComparison.lean:50}}.
\end{definition}

\begin{definition}[TopCat.Sheaf.ambientConstantInjectiveResolution]\label{def:TopCat.Sheaf.ambientConstantInjectiveResolution}
\lean{TopCat.Sheaf.ambientConstantInjectiveResolution}\leanok
The fixed injective resolution of the constant coefficient sheaf.

\noindent\emph{Lean:} \texttt{def TopCat.Sheaf.ambientConstantInjectiveResolution}, \href{https://github.com/Paul-Lez/HodgeConjecture/blob/main/HodgeConjecture/Definitions/AlgebraicTopology/OpenInjectiveResolutionComparison.lean#L73}{\texttt{HodgeConjecture/Definitions/AlgebraicTopology/OpenInjectiveResolutionComparison.lean:73}}.
\end{definition}

\section{Relative homology near a flattened support}\label{sec:HodgeConjecture.Definitions.AlgebraicTopology.FlattenedSupportLocalHomology}
\noindent\emph{File:} \href{https://github.com/Paul-Lez/HodgeConjecture/blob/main/HodgeConjecture/Definitions/AlgebraicTopology/FlattenedSupportLocalHomology.lean}{\texttt{HodgeConjecture/Definitions/AlgebraicTopology/FlattenedSupportLocalHomology.lean}}.

Radial compression into a small product-norm ball preserves the zero-normal plane.
Composing with the inverse flattening chart gives a pair homeomorphism from the normal-slice
model to a small open neighborhood paired with its support complement, along which the
homology calculation and the normalized class are transported.

\begin{definition}[AlgebraicTopology.Singular.neighborhoodSupportComplementPair]\label{def:AlgebraicTopology.Singular.neighborhoodSupportComplementPair}
\lean{AlgebraicTopology.Singular.neighborhoodSupportComplementPair}\leanok
The actual pair consisting of a neighborhood and the complement of a support in it.

\noindent\emph{Lean:} \texttt{abbrev AlgebraicTopology.Singular.neighborhoodSupportComplementPair}, \href{https://github.com/Paul-Lez/HodgeConjecture/blob/main/HodgeConjecture/Definitions/AlgebraicTopology/FlattenedSupportLocalHomology.lean#L28}{\texttt{HodgeConjecture/Definitions/AlgebraicTopology/FlattenedSupportLocalHomology.lean:28}}.
\end{definition}

\begin{definition}[AlgebraicTopology.Singular.flattenedSupportRadius]\label{def:AlgebraicTopology.Singular.flattenedSupportRadius}
\lean{AlgebraicTopology.Singular.flattenedSupportRadius}\leanok
\noindent\emph{Lean:} \texttt{def AlgebraicTopology.Singular.flattenedSupportRadius}, \href{https://github.com/Paul-Lez/HodgeConjecture/blob/main/HodgeConjecture/Definitions/AlgebraicTopology/FlattenedSupportLocalHomology.lean#L51}{\texttt{HodgeConjecture/Definitions/AlgebraicTopology/FlattenedSupportLocalHomology.lean:51}}.
\end{definition}

\begin{definition}[AlgebraicTopology.Singular.flattenedSupportEmbedding]\label{def:AlgebraicTopology.Singular.flattenedSupportEmbedding}
\lean{AlgebraicTopology.Singular.flattenedSupportEmbedding}\leanok
\uses{def:AlgebraicTopology.Singular.flattenedSupportRadius}
Actual radial compression followed by the inverse chart.

\noindent\emph{Lean:} \texttt{def AlgebraicTopology.Singular.flattenedSupportEmbedding}, \href{https://github.com/Paul-Lez/HodgeConjecture/blob/main/HodgeConjecture/Definitions/AlgebraicTopology/FlattenedSupportLocalHomology.lean#L62}{\texttt{HodgeConjecture/Definitions/AlgebraicTopology/FlattenedSupportLocalHomology.lean:62}}.
\end{definition}

\begin{definition}[AlgebraicTopology.Singular.flattenedSupportNeighborhood]\label{def:AlgebraicTopology.Singular.flattenedSupportNeighborhood}
\lean{AlgebraicTopology.Singular.flattenedSupportNeighborhood}\leanok
\uses{def:AlgebraicTopology.Singular.flattenedSupportEmbedding}
The actual open neighborhood on which the support has the normal-slice pair model.

\noindent\emph{Lean:} \texttt{def AlgebraicTopology.Singular.flattenedSupportNeighborhood}, \href{https://github.com/Paul-Lez/HodgeConjecture/blob/main/HodgeConjecture/Definitions/AlgebraicTopology/FlattenedSupportLocalHomology.lean#L75}{\texttt{HodgeConjecture/Definitions/AlgebraicTopology/FlattenedSupportLocalHomology.lean:75}}.
\end{definition}

\begin{definition}[AlgebraicTopology.Singular.flattenedSupportHomeomorph]\label{def:AlgebraicTopology.Singular.flattenedSupportHomeomorph}
\lean{AlgebraicTopology.Singular.flattenedSupportHomeomorph}\leanok
\uses{def:AlgebraicTopology.Singular.flattenedSupportNeighborhood}
The homeomorphism is constructed from radial compression and the inverse chart.

\noindent\emph{Lean:} \texttt{def AlgebraicTopology.Singular.flattenedSupportHomeomorph}, \href{https://github.com/Paul-Lez/HodgeConjecture/blob/main/HodgeConjecture/Definitions/AlgebraicTopology/FlattenedSupportLocalHomology.lean#L84}{\texttt{HodgeConjecture/Definitions/AlgebraicTopology/FlattenedSupportLocalHomology.lean:84}}.
\end{definition}

\begin{definition}[AlgebraicTopology.Singular.flattenedSupportComplementHomeomorph]\label{def:AlgebraicTopology.Singular.flattenedSupportComplementHomeomorph}
\lean{AlgebraicTopology.Singular.flattenedSupportComplementHomeomorph}\leanok
\uses{def:AlgebraicTopology.Singular.flattenedSupportHomeomorph}
Restriction of the actual homeomorphism to the support complements.

\noindent\emph{Lean:} \texttt{def AlgebraicTopology.Singular.flattenedSupportComplementHomeomorph}, \href{https://github.com/Paul-Lez/HodgeConjecture/blob/main/HodgeConjecture/Definitions/AlgebraicTopology/FlattenedSupportLocalHomology.lean#L113}{\texttt{HodgeConjecture/Definitions/AlgebraicTopology/FlattenedSupportLocalHomology.lean:113}}.
\end{definition}

\begin{definition}[AlgebraicTopology.Singular.flattenedSupportPairIso]\label{def:AlgebraicTopology.Singular.flattenedSupportPairIso}
\lean{AlgebraicTopology.Singular.flattenedSupportPairIso}\leanok
\uses{def:AlgebraicTopology.Singular.flattenedSupportComplementHomeomorph, def:AlgebraicTopology.Singular.neighborhoodSupportComplementPair, def:AlgebraicTopology.Singular.normalSlicePair}
A genuine pair isomorphism from the standard normal model to the local support pair.

\noindent\emph{Lean:} \texttt{def AlgebraicTopology.Singular.flattenedSupportPairIso}, \href{https://github.com/Paul-Lez/HodgeConjecture/blob/main/HodgeConjecture/Definitions/AlgebraicTopology/FlattenedSupportLocalHomology.lean#L120}{\texttt{HodgeConjecture/Definitions/AlgebraicTopology/FlattenedSupportLocalHomology.lean:120}}.
\end{definition}

\begin{definition}[AlgebraicTopology.Singular.flattenedSupportRelativeHomologyIso]\label{def:AlgebraicTopology.Singular.flattenedSupportRelativeHomologyIso}
\lean{AlgebraicTopology.Singular.flattenedSupportRelativeHomologyIso}\leanok
\uses{def:AlgebraicTopology.Singular.flattenedSupportPairIso, def:AlgebraicTopology.Singular.normalSliceRelativeHomologyIso}
Local relative homology is computed by actual pair maps and tangent contraction.

\noindent\emph{Lean:} \texttt{def AlgebraicTopology.Singular.flattenedSupportRelativeHomologyIso}, \href{https://github.com/Paul-Lez/HodgeConjecture/blob/main/HodgeConjecture/Definitions/AlgebraicTopology/FlattenedSupportLocalHomology.lean#L146}{\texttt{HodgeConjecture/Definitions/AlgebraicTopology/FlattenedSupportLocalHomology.lean:146}}.
\end{definition}

\begin{definition}[AlgebraicTopology.Singular.flattenedSupportNormalClass]\label{def:AlgebraicTopology.Singular.flattenedSupportNormalClass}
\lean{AlgebraicTopology.Singular.flattenedSupportNormalClass}\leanok
\uses{def:AlgebraicTopology.Singular.flattenedSupportRelativeHomologyIso, def:AlgebraicTopology.Singular.standardComplexLocalClass}
The local class is the transport of the fixed, exactly normalized complex normal class.

\noindent\emph{Lean:} \texttt{def AlgebraicTopology.Singular.flattenedSupportNormalClass}, \href{https://github.com/Paul-Lez/HodgeConjecture/blob/main/HodgeConjecture/Definitions/AlgebraicTopology/FlattenedSupportLocalHomology.lean#L154}{\texttt{HodgeConjecture/Definitions/AlgebraicTopology/FlattenedSupportLocalHomology.lean:154}}.
\end{definition}

\section{Normal projections and the local support coclass}\label{sec:HodgeConjecture.Definitions.AlgebraicTopology.NormalProjectionCoclass}
\noindent\emph{File:} \href{https://github.com/Paul-Lez/HodgeConjecture/blob/main/HodgeConjecture/Definitions/AlgebraicTopology/NormalProjectionCoclass.lean}{\texttt{HodgeConjecture/Definitions/AlgebraicTopology/NormalProjectionCoclass.lean}}.

The normal projection of a support-flattening chart is a map of pairs on every open subset of
its source, and pulling the fixed normal coclass back along it is compatible with restriction.
The radial-fiber factorization below identifies that coclass with the normalized local coclass.

\begin{definition}[AlgebraicTopology.Singular.neighborhoodSupportInclusionPairMap]\label{def:AlgebraicTopology.Singular.neighborhoodSupportInclusionPairMap}
\lean{AlgebraicTopology.Singular.neighborhoodSupportInclusionPairMap}\leanok
\uses{def:AlgebraicTopology.Singular.neighborhoodSupportComplementPair}
Inclusion of actual neighborhood/support-complement pairs.

\noindent\emph{Lean:} \texttt{def AlgebraicTopology.Singular.neighborhoodSupportInclusionPairMap}, \href{https://github.com/Paul-Lez/HodgeConjecture/blob/main/HodgeConjecture/Definitions/AlgebraicTopology/NormalProjectionCoclass.lean#L26}{\texttt{HodgeConjecture/Definitions/AlgebraicTopology/NormalProjectionCoclass.lean:26}}.
\end{definition}

\begin{definition}[AlgebraicTopology.Singular.chartNormalProjectionPair]\label{def:AlgebraicTopology.Singular.chartNormalProjectionPair}
\lean{AlgebraicTopology.Singular.chartNormalProjectionPair}\leanok
\uses{def:AlgebraicTopology.Singular.neighborhoodSupportComplementPair, def:AlgebraicTopology.Singular.standardComplexPuncturedPair}
The actual normal projection from a neighborhood support pair.

\noindent\emph{Lean:} \texttt{def AlgebraicTopology.Singular.chartNormalProjectionPair}, \href{https://github.com/Paul-Lez/HodgeConjecture/blob/main/HodgeConjecture/Definitions/AlgebraicTopology/NormalProjectionCoclass.lean#L37}{\texttt{HodgeConjecture/Definitions/AlgebraicTopology/NormalProjectionCoclass.lean:37}}.
\end{definition}

\begin{definition}[AlgebraicTopology.Singular.chartNormalProjectionCoclass]\label{def:AlgebraicTopology.Singular.chartNormalProjectionCoclass}
\lean{AlgebraicTopology.Singular.chartNormalProjectionCoclass}\leanok
\uses{def:AlgebraicTopology.Singular.chartNormalProjectionPair, def:AlgebraicTopology.Singular.normalizedRelativeCoclass, def:AlgebraicTopology.Singular.relativeCohomologyMap, def:AlgebraicTopology.Singular.relativeHomologyDegreeCast, def:AlgebraicTopology.Singular.relativeSingularBoundary, def:AlgebraicTopology.Singular.standardAffineBoundaryChainRetractionHomotopy, def:AlgebraicTopology.Singular.standardAffineBoundaryHomologyMap, def:AlgebraicTopology.Singular.standardAffineBoundaryHomologyRetraction, def:AlgebraicTopology.Singular.standardAmbientClassZero, def:AlgebraicTopology.Singular.standardComplexLocalClass, def:AlgebraicTopology.Singular.standardComplexLocalClassMul, def:AlgebraicTopology.Singular.standardLocalRelativeHomologyZeroε, def:AlgebraicTopology.Singular.standardNegativePoint, def:AlgebraicTopology.Singular.standardPositivePoint, def:AlgebraicTopology.Singular.standardPuncturedAffineBoundaryChain, def:AlgebraicTopology.Singular.standardPuncturedBoundaryClass, def:AlgebraicTopology.Singular.standardPuncturedBoundaryDetector, def:AlgebraicTopology.Singular.standardPuncturedFaceCycle, def:AlgebraicTopology.Singular.standardPuncturedRelativeBoundaryIso, def:AlgebraicTopology.Singular.standardSphereSimplicialBoundaryClass, def:AlgebraicTopology.Singular.standardSphereSimplicialNormalizedBoundaryClass, def:AlgebraicTopology.Singular.standardSphereSimplicialNormalizedBoundaryDetector}
The coclass on a whole chart neighborhood is the actual normal-projection pullback.

\noindent\emph{Lean:} \texttt{def AlgebraicTopology.Singular.chartNormalProjectionCoclass}, \href{https://github.com/Paul-Lez/HodgeConjecture/blob/main/HodgeConjecture/Definitions/AlgebraicTopology/NormalProjectionCoclass.lean#L46}{\texttt{HodgeConjecture/Definitions/AlgebraicTopology/NormalProjectionCoclass.lean:46}}.
\end{definition}

\section{Transporting actual neighborhood-support pairs through embeddings}\label{sec:HodgeConjecture.Definitions.AlgebraicTopology.NeighborhoodSupportPairImage}
\noindent\emph{File:} \href{https://github.com/Paul-Lez/HodgeConjecture/blob/main/HodgeConjecture/Definitions/AlgebraicTopology/NeighborhoodSupportPairImage.lean}{\texttt{HodgeConjecture/Definitions/AlgebraicTopology/NeighborhoodSupportPairImage.lean}}.

An embedding identifies a neighborhood and its support complement with their actual
images. Only equality of support membership on that neighborhood is required; no
cohomology comparison is supplied. Open embeddings therefore transport the cofinal
normal neighborhoods of auxiliary algebraic opens to the original ambient space.

\begin{definition}[AlgebraicTopology.Singular.neighborhoodSupportComplementImageHomeomorph]\label{def:AlgebraicTopology.Singular.neighborhoodSupportComplementImageHomeomorph}
\lean{AlgebraicTopology.Singular.neighborhoodSupportComplementImageHomeomorph}\leanok
The actual embedding-induced homeomorphism of the two support complements.

\noindent\emph{Lean:} \texttt{def AlgebraicTopology.Singular.neighborhoodSupportComplementImageHomeomorph}, \href{https://github.com/Paul-Lez/HodgeConjecture/blob/main/HodgeConjecture/Definitions/AlgebraicTopology/NeighborhoodSupportPairImage.lean#L27}{\texttt{HodgeConjecture/Definitions/AlgebraicTopology/NeighborhoodSupportPairImage.lean:27}}.
\end{definition}

\begin{definition}[AlgebraicTopology.Singular.neighborhoodSupportPairImageIso]\label{def:AlgebraicTopology.Singular.neighborhoodSupportPairImageIso}
\lean{AlgebraicTopology.Singular.neighborhoodSupportPairImageIso}\leanok
\uses{def:AlgebraicTopology.Singular.neighborhoodSupportComplementImageHomeomorph, def:AlgebraicTopology.Singular.neighborhoodSupportComplementPair}
The literal pair isomorphism, with the embedding as ambient map.

\noindent\emph{Lean:} \texttt{def AlgebraicTopology.Singular.neighborhoodSupportPairImageIso}, \href{https://github.com/Paul-Lez/HodgeConjecture/blob/main/HodgeConjecture/Definitions/AlgebraicTopology/NeighborhoodSupportPairImage.lean#L32}{\texttt{HodgeConjecture/Definitions/AlgebraicTopology/NeighborhoodSupportPairImage.lean:32}}.
\end{definition}

\begin{definition}[AlgebraicTopology.Singular.neighborhoodSupportPairImageCohomologyIso]\label{def:AlgebraicTopology.Singular.neighborhoodSupportPairImageCohomologyIso}
\lean{AlgebraicTopology.Singular.neighborhoodSupportPairImageCohomologyIso}\leanok
\uses{def:AlgebraicTopology.Singular.RelativeCohomology, def:AlgebraicTopology.Singular.neighborhoodSupportPairImageIso, def:HomologicalComplex.linearDualIso}
Relative cohomology transport, obtained by dualising the actual induced chain map.

\noindent\emph{Lean:} \texttt{def AlgebraicTopology.Singular.neighborhoodSupportPairImageCohomologyIso}, \href{https://github.com/Paul-Lez/HodgeConjecture/blob/main/HodgeConjecture/Definitions/AlgebraicTopology/NeighborhoodSupportPairImage.lean#L59}{\texttt{HodgeConjecture/Definitions/AlgebraicTopology/NeighborhoodSupportPairImage.lean:59}}.
\end{definition}

\begin{definition}[AlgebraicTopology.Singular.neighborhoodSupportPairImageCohomologyEquiv]\label{def:AlgebraicTopology.Singular.neighborhoodSupportPairImageCohomologyEquiv}
\lean{AlgebraicTopology.Singular.neighborhoodSupportPairImageCohomologyEquiv}\leanok
\uses{def:AlgebraicTopology.Singular.neighborhoodSupportPairImageCohomologyIso}
Relative cohomology transport as a linear equivalence.

\noindent\emph{Lean:} \texttt{def AlgebraicTopology.Singular.neighborhoodSupportPairImageCohomologyEquiv}, \href{https://github.com/Paul-Lez/HodgeConjecture/blob/main/HodgeConjecture/Definitions/AlgebraicTopology/NeighborhoodSupportPairImage.lean#L67}{\texttt{HodgeConjecture/Definitions/AlgebraicTopology/NeighborhoodSupportPairImage.lean:67}}.
\end{definition}

\section{Cohomology-sheaf stalk vanishing from cofinal local section calculations}\label{sec:HodgeConjecture.Definitions.AlgebraicTopology.CohomologySheafStalkVanishing}
\noindent\emph{File:} \href{https://github.com/Paul-Lez/HodgeConjecture/blob/main/HodgeConjecture/Definitions/AlgebraicTopology/CohomologySheafStalkVanishing.lean}{\texttt{HodgeConjecture/Definitions/AlgebraicTopology/CohomologySheafStalkVanishing.lean}}.

For a complex of additive sheaves, the stalk of the homology presheaf of its
underlying presheaf complex is canonically the stalk of its homology sheaf.
Both comparisons use exact stalk functors; evaluation on an open set is not
mistakenly treated as exact on sheaves.

Consequently, vanishing of section-complex homology on a cofinal system of
neighborhoods implies vanishing of the cohomology-sheaf stalk. The neighborhood
may depend on the original open set; no single neighborhood is silently
substituted for a cofinal family. This is the generic passage from the local
normal-slice calculation to sheaf support purity.

\begin{definition}[TopCat.Sheaf.sectionCohomologyPresheaf]\label{def:TopCat.Sheaf.sectionCohomologyPresheaf}
\lean{TopCat.Sheaf.sectionCohomologyPresheaf}\leanok
The presheaf homology of the actual underlying complex. This is not defined
as evaluation of the cohomology sheaf, since that evaluation is not exact.

\noindent\emph{Lean:} \texttt{def TopCat.Sheaf.sectionCohomologyPresheaf}, \href{https://github.com/Paul-Lez/HodgeConjecture/blob/main/HodgeConjecture/Definitions/AlgebraicTopology/CohomologySheafStalkVanishing.lean#L45}{\texttt{HodgeConjecture/Definitions/AlgebraicTopology/CohomologySheafStalkVanishing.lean:45}}.
\end{definition}

\begin{definition}[TopCat.Sheaf.sectionCohomologyPresheafOnOpenIso]\label{def:TopCat.Sheaf.sectionCohomologyPresheafOnOpenIso}
\lean{TopCat.Sheaf.sectionCohomologyPresheafOnOpenIso}\leanok
\uses{def:TopCat.Sheaf.sectionCohomologyPresheaf, def:TopCat.Sheaf.supportEvaluation}
Exact evaluation of presheaves identifies section-complex homology with
the homology presheaf on the same open set.

\noindent\emph{Lean:} \texttt{def TopCat.Sheaf.sectionCohomologyPresheafOnOpenIso}, \href{https://github.com/Paul-Lez/HodgeConjecture/blob/main/HodgeConjecture/Definitions/AlgebraicTopology/CohomologySheafStalkVanishing.lean#L50}{\texttt{HodgeConjecture/Definitions/AlgebraicTopology/CohomologySheafStalkVanishing.lean:50}}.
\end{definition}

\begin{definition}[TopCat.Sheaf.sectionCohomologyPresheafStalkIso]\label{def:TopCat.Sheaf.sectionCohomologyPresheafStalkIso}
\lean{TopCat.Sheaf.sectionCohomologyPresheafStalkIso}\leanok
\uses{def:AlgebraicTopology.Singular.additivePresheafStalkSheafificationIso, def:TopCat.Sheaf.sectionCohomologyPresheaf}
The homology presheaf and actual homology sheaf have canonically equal stalks.
The underlying sheaf-forgetful functor is not assumed exact.

\noindent\emph{Lean:} \texttt{def TopCat.Sheaf.sectionCohomologyPresheafStalkIso}, \href{https://github.com/Paul-Lez/HodgeConjecture/blob/main/HodgeConjecture/Definitions/AlgebraicTopology/CohomologySheafStalkVanishing.lean#L59}{\texttt{HodgeConjecture/Definitions/AlgebraicTopology/CohomologySheafStalkVanishing.lean:59}}.
\end{definition}

\section{Singular-cochain models for supported injective resolutions}\label{sec:HodgeConjecture.Definitions.AlgebraicTopology.SingularFlasqueSupportModel}
\noindent\emph{File:} \href{https://github.com/Paul-Lez/HodgeConjecture/blob/main/HodgeConjecture/Definitions/AlgebraicTopology/SingularFlasqueSupportModel.lean}{\texttt{HodgeConjecture/Definitions/AlgebraicTopology/SingularFlasqueSupportModel.lean}}.

On a space with a contractible open basis, the singular-cochain augmentation resolves the
constant rational sheaf, and we extend that augmentation to the fixed injective resolution
with a strict normalization. On a hereditarily paracompact Hausdorff space the singular model
is termwise flasque, so the supported comparison is a quasi-isomorphism on sheaves and on
every open set. The contractible basis is a hypothesis of these generic topological lemmas,
constructed for the smooth complex-point application below.

\begin{definition}[AlgebraicTopology.Singular.rationalSingularCochainComplex]\label{def:AlgebraicTopology.Singular.rationalSingularCochainComplex}
\lean{AlgebraicTopology.Singular.rationalSingularCochainComplex}\leanok
\uses{def:AlgebraicTopology.Singular.singularCochainSheafComplex}
The actual singular-cochain sheaf complex in integer degrees.

\noindent\emph{Lean:} \texttt{def AlgebraicTopology.Singular.rationalSingularCochainComplex}, \href{https://github.com/Paul-Lez/HodgeConjecture/blob/main/HodgeConjecture/Definitions/AlgebraicTopology/SingularFlasqueSupportModel.lean#L30}{\texttt{HodgeConjecture/Definitions/AlgebraicTopology/SingularFlasqueSupportModel.lean:30}}.
\end{definition}

\begin{definition}[AlgebraicTopology.Singular.rationalConstantInjectiveComplex]\label{def:AlgebraicTopology.Singular.rationalConstantInjectiveComplex}
\lean{AlgebraicTopology.Singular.rationalConstantInjectiveComplex}\leanok
\uses{def:TopCat.Sheaf.ambientConstantInjectiveResolution}
The fixed injective resolution of actual rational constants, in integer degrees.

\noindent\emph{Lean:} \texttt{def AlgebraicTopology.Singular.rationalConstantInjectiveComplex}, \href{https://github.com/Paul-Lez/HodgeConjecture/blob/main/HodgeConjecture/Definitions/AlgebraicTopology/SingularFlasqueSupportModel.lean#L34}{\texttt{HodgeConjecture/Definitions/AlgebraicTopology/SingularFlasqueSupportModel.lean:34}}.
\end{definition}

\begin{definition}[AlgebraicTopology.Singular.singularToConstantInjectiveResolution]\label{def:AlgebraicTopology.Singular.singularToConstantInjectiveResolution}
\lean{AlgebraicTopology.Singular.singularToConstantInjectiveResolution}\leanok
\uses{def:AlgebraicTopology.Singular.constantsToSingularCochainSheafComplex, def:AlgebraicTopology.Singular.constantsToSingularCochainShortComplexSheafificationUnit, def:AlgebraicTopology.Singular.singularCochainSheafificationUnit, def:AlgebraicTopology.contractibleSingularChainHomotopyEquiv, def:CochainComplex.liftToInjectiveNat, def:TopCat.Sheaf.ambientConstantInjectiveResolution}
Strictly extend the constant augmentation to the actual injective resolution.

\noindent\emph{Lean:} \texttt{def AlgebraicTopology.Singular.singularToConstantInjectiveResolution}, \href{https://github.com/Paul-Lez/HodgeConjecture/blob/main/HodgeConjecture/Definitions/AlgebraicTopology/SingularFlasqueSupportModel.lean#L57}{\texttt{HodgeConjecture/Definitions/AlgebraicTopology/SingularFlasqueSupportModel.lean:57}}.
\end{definition}

\begin{definition}[AlgebraicTopology.Singular.singularToConstantInjectiveComplex]\label{def:AlgebraicTopology.Singular.singularToConstantInjectiveComplex}
\lean{AlgebraicTopology.Singular.singularToConstantInjectiveComplex}\leanok
\uses{def:AlgebraicTopology.Singular.rationalConstantInjectiveComplex, def:AlgebraicTopology.Singular.rationalSingularCochainComplex, def:AlgebraicTopology.Singular.singularToConstantInjectiveResolution}
The same strictly normalized comparison in integer degrees.

\noindent\emph{Lean:} \texttt{def AlgebraicTopology.Singular.singularToConstantInjectiveComplex}, \href{https://github.com/Paul-Lez/HodgeConjecture/blob/main/HodgeConjecture/Definitions/AlgebraicTopology/SingularFlasqueSupportModel.lean#L91}{\texttt{HodgeConjecture/Definitions/AlgebraicTopology/SingularFlasqueSupportModel.lean:91}}.
\end{definition}

\begin{definition}[AlgebraicTopology.Singular.supportedRationalSingularCochainComplex]\label{def:AlgebraicTopology.Singular.supportedRationalSingularCochainComplex}
\lean{AlgebraicTopology.Singular.supportedRationalSingularCochainComplex}\leanok
\uses{def:AlgebraicTopology.Singular.rationalSingularCochainComplex, def:TopCat.Sheaf.sheafSectionsSupportedOutside}
The actual supported singular-cochain sheaf model. Its terms are kernels
of restriction; its cohomology is not defined to be a desired purity group.

\noindent\emph{Lean:} \texttt{def AlgebraicTopology.Singular.supportedRationalSingularCochainComplex}, \href{https://github.com/Paul-Lez/HodgeConjecture/blob/main/HodgeConjecture/Definitions/AlgebraicTopology/SingularFlasqueSupportModel.lean#L101}{\texttt{HodgeConjecture/Definitions/AlgebraicTopology/SingularFlasqueSupportModel.lean:101}}.
\end{definition}

\begin{definition}[AlgebraicTopology.Singular.supportedSingularToInjectiveComplex]\label{def:AlgebraicTopology.Singular.supportedSingularToInjectiveComplex}
\lean{AlgebraicTopology.Singular.supportedSingularToInjectiveComplex}\leanok
\uses{def:AlgebraicTopology.Singular.singularToConstantInjectiveComplex, def:AlgebraicTopology.Singular.supportedRationalSingularCochainComplex}
Apply actual supported sections to the constructed resolution comparison.

\noindent\emph{Lean:} \texttt{def AlgebraicTopology.Singular.supportedSingularToInjectiveComplex}, \href{https://github.com/Paul-Lez/HodgeConjecture/blob/main/HodgeConjecture/Definitions/AlgebraicTopology/SingularFlasqueSupportModel.lean#L108}{\texttt{HodgeConjecture/Definitions/AlgebraicTopology/SingularFlasqueSupportModel.lean:108}}.
\end{definition}

\begin{definition}[AlgebraicTopology.Singular.supportedSingularInjectiveHomologyIso]\label{def:AlgebraicTopology.Singular.supportedSingularInjectiveHomologyIso}
\lean{AlgebraicTopology.Singular.supportedSingularInjectiveHomologyIso}\leanok
\uses{def:AlgebraicTopology.Singular.ambientOpen, def:AlgebraicTopology.Singular.coveringSieveOpenFamilyOn, def:AlgebraicTopology.Singular.openFamilySieve, def:AlgebraicTopology.Singular.openSimplexMap, def:AlgebraicTopology.Singular.singularCochainEquivSimplexFunction, def:AlgebraicTopology.Singular.supportedSingularToInjectiveComplex, def:TopCat.Sheaf.IsFlasque.BoundedBelowComplex.cyclesShortComplex, def:TopCat.Sheaf.IsFlasque.BoundedBelowComplex.globalSectionsComplex, def:TopCat.Sheaf.freeAbelianYonedaSheafHomEquiv, def:TopCat.Sheaf.openRestrictionGlobalSectionsIso, def:TopCat.Sheaf.supportRestrictionComplexShortComplexMap, def:TopCat.Sheaf.supportRestrictionSectionsComplexShortComplex, def:TopCat.Sheaf.supportRestrictionSectionsShortComplex}
Actual section-complex cohomology agrees through the normalized comparison.

\noindent\emph{Lean:} \texttt{def AlgebraicTopology.Singular.supportedSingularInjectiveHomologyIso}, \href{https://github.com/Paul-Lez/HodgeConjecture/blob/main/HodgeConjecture/Definitions/AlgebraicTopology/SingularFlasqueSupportModel.lean#L140}{\texttt{HodgeConjecture/Definitions/AlgebraicTopology/SingularFlasqueSupportModel.lean:140}}.
\end{definition}

\section{Supported sections preserve flasqueness}\label{sec:HodgeConjecture.Definitions.AlgebraicTopology.FlasqueSupportedSections}
\noindent\emph{File:} \href{https://github.com/Paul-Lez/HodgeConjecture/blob/main/HodgeConjecture/Definitions/AlgebraicTopology/FlasqueSupportedSections.lean}{\texttt{HodgeConjecture/Definitions/AlgebraicTopology/FlasqueSupportedSections.lean}}.

A supported section over \texttt{V} glues with zero on the excluded open \texttt{U} to a section on
\texttt{V ⊔ U}. Flasqueness extends that section to \texttt{W ⊔ U} when \texttt{V ≤ W}, and its restriction to
\texttt{W} vanishes on \texttt{W ⊓ U}, so it lies in the kernel defining supported sections.

\begin{definition}[TopCat.Sheaf.supportedOutsideIntersectionIso]\label{def:TopCat.Sheaf.supportedOutsideIntersectionIso}
\lean{TopCat.Sheaf.supportedOutsideIntersectionIso}\leanok
\uses{def:TopCat.Sheaf.openRestrictionPushforward}
Restriction-pushforward on an ambient open is literally evaluation on its
intersection with the excluded open.

\noindent\emph{Lean:} \texttt{def TopCat.Sheaf.supportedOutsideIntersectionIso}, \href{https://github.com/Paul-Lez/HodgeConjecture/blob/main/HodgeConjecture/Definitions/AlgebraicTopology/FlasqueSupportedSections.lean#L29}{\texttt{HodgeConjecture/Definitions/AlgebraicTopology/FlasqueSupportedSections.lean:29}}.
\end{definition}

\begin{definition}[TopCat.Sheaf.supportedOutsideGlueZero]\label{def:TopCat.Sheaf.supportedOutsideGlueZero}
\lean{TopCat.Sheaf.supportedOutsideGlueZero}\leanok
\uses{def:TopCat.Sheaf.sheafSectionsSupportedOutsideInclusion, def:TopCat.Sheaf.supportedOutsideIntersectionIso}
Pairwise sheaf gluing extends a supported section by zero across \texttt{U}, as an actual
additive morphism from the defining kernel.

\noindent\emph{Lean:} \texttt{def TopCat.Sheaf.supportedOutsideGlueZero}, \href{https://github.com/Paul-Lez/HodgeConjecture/blob/main/HodgeConjecture/Definitions/AlgebraicTopology/FlasqueSupportedSections.lean#L57}{\texttt{HodgeConjecture/Definitions/AlgebraicTopology/FlasqueSupportedSections.lean:57}}.
\end{definition}

\section{Supported-section kernels and open restriction cones}\label{sec:HodgeConjecture.Definitions.AlgebraicTopology.SupportedSectionRestrictionCone}
\noindent\emph{File:} \href{https://github.com/Paul-Lez/HodgeConjecture/blob/main/HodgeConjecture/Definitions/AlgebraicTopology/SupportedSectionRestrictionCone.lean}{\texttt{HodgeConjecture/Definitions/AlgebraicTopology/SupportedSectionRestrictionCone.lean}}.

For a termwise flasque coefficient complex, the kernel-defined supported section complex
computes the cone of restriction from \texttt{V} to \texttt{V ∩ U}, with the conventional degree shift.
The arrow identification displays the open-intersection map.

\begin{definition}[TopCat.Sheaf.supportEvaluationRestriction]\label{def:TopCat.Sheaf.supportEvaluationRestriction}
\lean{TopCat.Sheaf.supportEvaluationRestriction}\leanok
\uses{def:TopCat.Sheaf.supportEvaluation}
Restriction between the actual open-section evaluation functors.

\noindent\emph{Lean:} \texttt{def TopCat.Sheaf.supportEvaluationRestriction}, \href{https://github.com/Paul-Lez/HodgeConjecture/blob/main/HodgeConjecture/Definitions/AlgebraicTopology/SupportedSectionRestrictionCone.lean#L29}{\texttt{HodgeConjecture/Definitions/AlgebraicTopology/SupportedSectionRestrictionCone.lean:29}}.
\end{definition}

\begin{definition}[TopCat.Sheaf.sectionComplexRestriction]\label{def:TopCat.Sheaf.sectionComplexRestriction}
\lean{TopCat.Sheaf.sectionComplexRestriction}\leanok
\uses{def:TopCat.Sheaf.supportEvaluationRestriction}
The actual restriction morphism of section complexes, in any grading.

\noindent\emph{Lean:} \texttt{def TopCat.Sheaf.sectionComplexRestriction}, \href{https://github.com/Paul-Lez/HodgeConjecture/blob/main/HodgeConjecture/Definitions/AlgebraicTopology/SupportedSectionRestrictionCone.lean#L35}{\texttt{HodgeConjecture/Definitions/AlgebraicTopology/SupportedSectionRestrictionCone.lean:35}}.
\end{definition}

\begin{definition}[TopCat.Sheaf.sectionComplexRestrictionExtendConeIso]\label{def:TopCat.Sheaf.sectionComplexRestrictionExtendConeIso}
\lean{TopCat.Sheaf.sectionComplexRestrictionExtendConeIso}\leanok
\uses{def:HomologicalComplex.mapExtendCanonicalIso, def:TopCat.Sheaf.sectionComplexRestriction}
The corresponding actual restriction cones agree under the canonical grading comparison.

\noindent\emph{Lean:} \texttt{def TopCat.Sheaf.sectionComplexRestrictionExtendConeIso}, \href{https://github.com/Paul-Lez/HodgeConjecture/blob/main/HodgeConjecture/Definitions/AlgebraicTopology/SupportedSectionRestrictionCone.lean#L57}{\texttt{HodgeConjecture/Definitions/AlgebraicTopology/SupportedSectionRestrictionCone.lean:57}}.
\end{definition}

\begin{definition}[TopCat.Sheaf.supportRestrictionSectionsIntersectionIso]\label{def:TopCat.Sheaf.supportRestrictionSectionsIntersectionIso}
\lean{TopCat.Sheaf.supportRestrictionSectionsIntersectionIso}\leanok
\uses{def:TopCat.Sheaf.supportRestrictionSectionsComplexShortComplex, def:TopCat.Sheaf.supportedOutsideIntersectionIso}
The third term of the actual support sequence is sections on the actual intersection.

\noindent\emph{Lean:} \texttt{def TopCat.Sheaf.supportRestrictionSectionsIntersectionIso}, \href{https://github.com/Paul-Lez/HodgeConjecture/blob/main/HodgeConjecture/Definitions/AlgebraicTopology/SupportedSectionRestrictionCone.lean#L77}{\texttt{HodgeConjecture/Definitions/AlgebraicTopology/SupportedSectionRestrictionCone.lean:77}}.
\end{definition}

\begin{definition}[TopCat.Sheaf.supportRestrictionSectionsConeIso]\label{def:TopCat.Sheaf.supportRestrictionSectionsConeIso}
\lean{TopCat.Sheaf.supportRestrictionSectionsConeIso}\leanok
\uses{def:TopCat.Sheaf.sectionComplexRestriction, def:TopCat.Sheaf.supportRestrictionSectionsIntersectionIso}
The cone of the actual support-sequence arrow is the cone of actual open restriction.

\noindent\emph{Lean:} \texttt{def TopCat.Sheaf.supportRestrictionSectionsConeIso}, \href{https://github.com/Paul-Lez/HodgeConjecture/blob/main/HodgeConjecture/Definitions/AlgebraicTopology/SupportedSectionRestrictionCone.lean#L96}{\texttt{HodgeConjecture/Definitions/AlgebraicTopology/SupportedSectionRestrictionCone.lean:96}}.
\end{definition}

\begin{definition}[TopCat.Sheaf.supportedSectionHomologyIsoRestrictionCone]\label{def:TopCat.Sheaf.supportedSectionHomologyIsoRestrictionCone}
\lean{TopCat.Sheaf.supportedSectionHomologyIsoRestrictionCone}\leanok
\uses{def:CochainComplex.mappingCocone.shiftedLiftShortComplex, def:TopCat.Sheaf.supportRestrictionSectionsConeIso, def:TopCat.Sheaf.supportRestrictionSectionsShortComplex}
A flasque supported-section kernel computes the true restriction cone, with
supported degree \texttt{n} corresponding to cone degree \texttt{n - 1}.

\noindent\emph{Lean:} \texttt{def TopCat.Sheaf.supportedSectionHomologyIsoRestrictionCone}, \href{https://github.com/Paul-Lez/HodgeConjecture/blob/main/HodgeConjecture/Definitions/AlgebraicTopology/SupportedSectionRestrictionCone.lean#L108}{\texttt{HodgeConjecture/Definitions/AlgebraicTopology/SupportedSectionRestrictionCone.lean:108}}.
\end{definition}

\section{Singular cochains compute sheafified cochains on an arbitrary open}\label{sec:HodgeConjecture.Definitions.AlgebraicTopology.SingularCochainOpenSections}
\noindent\emph{File:} \href{https://github.com/Paul-Lez/HodgeConjecture/blob/main/HodgeConjecture/Definitions/AlgebraicTopology/SingularCochainOpenSections.lean}{\texttt{HodgeConjecture/Definitions/AlgebraicTopology/SingularCochainOpenSections.lean}}.

\begin{definition}[AlgebraicTopology.Singular.openRawSingularCochainComplex]\label{def:AlgebraicTopology.Singular.openRawSingularCochainComplex}
\lean{AlgebraicTopology.Singular.openRawSingularCochainComplex}\leanok
\uses{def:AlgebraicTopology.Singular.singularCochainPresheafComplex}
Raw singular cochains evaluated on an ambient open.

\noindent\emph{Lean:} \texttt{def AlgebraicTopology.Singular.openRawSingularCochainComplex}, \href{https://github.com/Paul-Lez/HodgeConjecture/blob/main/HodgeConjecture/Definitions/AlgebraicTopology/SingularCochainOpenSections.lean#L20}{\texttt{HodgeConjecture/Definitions/AlgebraicTopology/SingularCochainOpenSections.lean:20}}.
\end{definition}

\begin{definition}[AlgebraicTopology.Singular.openSingularCochainSheafComplex]\label{def:AlgebraicTopology.Singular.openSingularCochainSheafComplex}
\lean{AlgebraicTopology.Singular.openSingularCochainSheafComplex}\leanok
\uses{def:AlgebraicTopology.Singular.singularCochainSheafComplex, def:TopCat.Sheaf.supportEvaluation}
Sections on an ambient open of the actual singular cochain sheaf complex.

\noindent\emph{Lean:} \texttt{def AlgebraicTopology.Singular.openSingularCochainSheafComplex}, \href{https://github.com/Paul-Lez/HodgeConjecture/blob/main/HodgeConjecture/Definitions/AlgebraicTopology/SingularCochainOpenSections.lean#L25}{\texttt{HodgeConjecture/Definitions/AlgebraicTopology/SingularCochainOpenSections.lean:25}}.
\end{definition}

\begin{definition}[AlgebraicTopology.Singular.openRawToSingularCochainSheafComplex]\label{def:AlgebraicTopology.Singular.openRawToSingularCochainSheafComplex}
\lean{AlgebraicTopology.Singular.openRawToSingularCochainSheafComplex}\leanok
\uses{def:AlgebraicTopology.Singular.openRawSingularCochainComplex, def:AlgebraicTopology.Singular.openSingularCochainSheafComplex, def:AlgebraicTopology.Singular.singularCochainSheafificationUnit}
The actual singular-cochain sheafification unit evaluated on an ambient open.

\noindent\emph{Lean:} \texttt{def AlgebraicTopology.Singular.openRawToSingularCochainSheafComplex}, \href{https://github.com/Paul-Lez/HodgeConjecture/blob/main/HodgeConjecture/Definitions/AlgebraicTopology/SingularCochainOpenSections.lean#L30}{\texttt{HodgeConjecture/Definitions/AlgebraicTopology/SingularCochainOpenSections.lean:30}}.
\end{definition}

\begin{definition}[AlgebraicTopology.Singular.openSubspaceImageTopIso]\label{def:AlgebraicTopology.Singular.openSubspaceImageTopIso}
\lean{AlgebraicTopology.Singular.openSubspaceImageTopIso}\leanok
The image of the top open of a subspace is that ambient open itself.

\noindent\emph{Lean:} \texttt{def AlgebraicTopology.Singular.openSubspaceImageTopIso}, \href{https://github.com/Paul-Lez/HodgeConjecture/blob/main/HodgeConjecture/Definitions/AlgebraicTopology/SingularCochainOpenSections.lean#L36}{\texttt{HodgeConjecture/Definitions/AlgebraicTopology/SingularCochainOpenSections.lean:36}}.
\end{definition}

\begin{definition}[AlgebraicTopology.Singular.openRawSingularCochainComplexIsoGlobal]\label{def:AlgebraicTopology.Singular.openRawSingularCochainComplexIsoGlobal}
\lean{AlgebraicTopology.Singular.openRawSingularCochainComplexIsoGlobal}\leanok
\uses{def:AlgebraicTopology.Singular.globalRawSingularCochainComplex, def:AlgebraicTopology.Singular.openRawSingularCochainComplex, def:AlgebraicTopology.Singular.openSubspaceImageTopIso, def:AlgebraicTopology.Singular.singularCochainPresheafComplexOpenRestrictionIso}
Ambient raw cochains on \texttt{V} identified with top-open raw cochains of
the space \texttt{V}, through the actual image homeomorphism.

\noindent\emph{Lean:} \texttt{def AlgebraicTopology.Singular.openRawSingularCochainComplexIsoGlobal}, \href{https://github.com/Paul-Lez/HodgeConjecture/blob/main/HodgeConjecture/Definitions/AlgebraicTopology/SingularCochainOpenSections.lean#L42}{\texttt{HodgeConjecture/Definitions/AlgebraicTopology/SingularCochainOpenSections.lean:42}}.
\end{definition}

\begin{definition}[AlgebraicTopology.Singular.openSingularCochainSheafComplexIsoGlobal]\label{def:AlgebraicTopology.Singular.openSingularCochainSheafComplexIsoGlobal}
\lean{AlgebraicTopology.Singular.openSingularCochainSheafComplexIsoGlobal}\leanok
\uses{def:AlgebraicTopology.Singular.globalSingularCochainSheafComplex, def:AlgebraicTopology.Singular.openSingularCochainSheafComplex, def:AlgebraicTopology.Singular.openSubspaceImageTopIso, def:AlgebraicTopology.Singular.singularCochainSheafComplexOpenRestrictionIso}
Sections of the ambient singular sheaf on \texttt{V} identified with global
sections of its intrinsic singular sheaf, using the normalized open restriction.

\noindent\emph{Lean:} \texttt{def AlgebraicTopology.Singular.openSingularCochainSheafComplexIsoGlobal}, \href{https://github.com/Paul-Lez/HodgeConjecture/blob/main/HodgeConjecture/Definitions/AlgebraicTopology/SingularCochainOpenSections.lean#L54}{\texttt{HodgeConjecture/Definitions/AlgebraicTopology/SingularCochainOpenSections.lean:54}}.
\end{definition}

\begin{definition}[AlgebraicTopology.Singular.openRawSingularRestriction]\label{def:AlgebraicTopology.Singular.openRawSingularRestriction}
\lean{AlgebraicTopology.Singular.openRawSingularRestriction}\leanok
\uses{def:AlgebraicTopology.Singular.openRawSingularCochainComplex}
Actual restriction of raw cochains between two ambient opens.

\noindent\emph{Lean:} \texttt{def AlgebraicTopology.Singular.openRawSingularRestriction}, \href{https://github.com/Paul-Lez/HodgeConjecture/blob/main/HodgeConjecture/Definitions/AlgebraicTopology/SingularCochainOpenSections.lean#L75}{\texttt{HodgeConjecture/Definitions/AlgebraicTopology/SingularCochainOpenSections.lean:75}}.
\end{definition}

\begin{definition}[AlgebraicTopology.Singular.openSingularSheafRestriction]\label{def:AlgebraicTopology.Singular.openSingularSheafRestriction}
\lean{AlgebraicTopology.Singular.openSingularSheafRestriction}\leanok
\uses{def:AlgebraicTopology.Singular.openSingularCochainSheafComplex}
Actual restriction of sections of the singular cochain sheaf.

\noindent\emph{Lean:} \texttt{def AlgebraicTopology.Singular.openSingularSheafRestriction}, \href{https://github.com/Paul-Lez/HodgeConjecture/blob/main/HodgeConjecture/Definitions/AlgebraicTopology/SingularCochainOpenSections.lean#L82}{\texttt{HodgeConjecture/Definitions/AlgebraicTopology/SingularCochainOpenSections.lean:82}}.
\end{definition}

\section{Canonical natural relative-cochain cone comparison}\label{sec:HodgeConjecture.Definitions.AlgebraicTopology.RelativeCochainConeNaturality}
\noindent\emph{File:} \href{https://github.com/Paul-Lez/HodgeConjecture/blob/main/HodgeConjecture/Definitions/AlgebraicTopology/RelativeCochainConeNaturality.lean}{\texttt{HodgeConjecture/Definitions/AlgebraicTopology/RelativeCochainConeNaturality.lean}}.

Unlike a completion of a morphism of distinguished triangles, the explicit
short-exact-sequence lift is natural before passage to homology.

The lift is positive \texttt{(inclusion, 0)}; Mathlib's cone connecting morphism
is the negative first projection. Their exact composite is displayed below.
This file does not assert equality with the older
\texttt{relativeCochainConeCohomologyEquiv}, which uses a completed triangle map.
Such equality, in particular on normalized point neighborhoods, is a
separate comparison theorem.

\begin{definition}[AlgebraicTopology.Singular.relativeDualCochainShortComplexNatMap]\label{def:AlgebraicTopology.Singular.relativeDualCochainShortComplexNatMap}
\lean{AlgebraicTopology.Singular.relativeDualCochainShortComplexNatMap}\leanok
\uses{def:AlgebraicTopology.Singular.relativeDualCochainShortComplexNat}
The actual contravariant map of dual relative short exact sequences.

\noindent\emph{Lean:} \texttt{def AlgebraicTopology.Singular.relativeDualCochainShortComplexNatMap}, \href{https://github.com/Paul-Lez/HodgeConjecture/blob/main/HodgeConjecture/Definitions/AlgebraicTopology/RelativeCochainConeNaturality.lean#L35}{\texttt{HodgeConjecture/Definitions/AlgebraicTopology/RelativeCochainConeNaturality.lean:35}}.
\end{definition}

\begin{definition}[AlgebraicTopology.Singular.relativeDualCochainShortComplexIntMap]\label{def:AlgebraicTopology.Singular.relativeDualCochainShortComplexIntMap}
\lean{AlgebraicTopology.Singular.relativeDualCochainShortComplexIntMap}\leanok
\uses{def:AlgebraicTopology.Singular.relativeDualCochainShortComplexInt, def:AlgebraicTopology.Singular.relativeDualCochainShortComplexNatMap}
Extend the actual relative short-complex map to integer degrees.

\noindent\emph{Lean:} \texttt{def AlgebraicTopology.Singular.relativeDualCochainShortComplexIntMap}, \href{https://github.com/Paul-Lez/HodgeConjecture/blob/main/HodgeConjecture/Definitions/AlgebraicTopology/RelativeCochainConeNaturality.lean#L57}{\texttt{HodgeConjecture/Definitions/AlgebraicTopology/RelativeCochainConeNaturality.lean:57}}.
\end{definition}

\begin{definition}[AlgebraicTopology.Singular.relativeCochainConeMap]\label{def:AlgebraicTopology.Singular.relativeCochainConeMap}
\lean{AlgebraicTopology.Singular.relativeCochainConeMap}\leanok
\uses{def:AlgebraicTopology.Singular.relativeCochainRestrictionInt, def:AlgebraicTopology.Singular.relativeDualCochainShortComplexIntMap}
The literal relative restriction-cone map induced by a map of pairs.

\noindent\emph{Lean:} \texttt{def AlgebraicTopology.Singular.relativeCochainConeMap}, \href{https://github.com/Paul-Lez/HodgeConjecture/blob/main/HodgeConjecture/Definitions/AlgebraicTopology/RelativeCochainConeNaturality.lean#L63}{\texttt{HodgeConjecture/Definitions/AlgebraicTopology/RelativeCochainConeNaturality.lean:63}}.
\end{definition}

\begin{definition}[AlgebraicTopology.Singular.relativeDualCochainHomologyIsoCone]\label{def:AlgebraicTopology.Singular.relativeDualCochainHomologyIsoCone}
\lean{AlgebraicTopology.Singular.relativeDualCochainHomologyIsoCone}\leanok
\uses{def:AlgebraicTopology.Singular.relativeChainShortComplex, def:AlgebraicTopology.Singular.relativeCochainRestrictionInt, def:AlgebraicTopology.Singular.relativeDualCochainShortComplexIntEvalIso, def:CochainComplex.mappingCocone.shortExactHomologyIsoCone}
Canonical comparison from integer dual-relative cohomology to its
restriction cone, induced by the positive \texttt{(inclusion, 0)} lift.

\noindent\emph{Lean:} \texttt{def AlgebraicTopology.Singular.relativeDualCochainHomologyIsoCone}, \href{https://github.com/Paul-Lez/HodgeConjecture/blob/main/HodgeConjecture/Definitions/AlgebraicTopology/RelativeCochainConeNaturality.lean#L72}{\texttt{HodgeConjecture/Definitions/AlgebraicTopology/RelativeCochainConeNaturality.lean:72}}.
\end{definition}

\begin{definition}[AlgebraicTopology.Singular.relativeDualCochainCohomologyEquiv]\label{def:AlgebraicTopology.Singular.relativeDualCochainCohomologyEquiv}
\lean{AlgebraicTopology.Singular.relativeDualCochainCohomologyEquiv}\leanok
\uses{def:AlgebraicTopology.Singular.RelativeCohomology, def:AlgebraicTopology.Singular.relativeDualCochainShortComplexInt}
Actual integer dual-relative cohomology computes the repository's
relative cohomology by the evaluation pairing.

\noindent\emph{Lean:} \texttt{def AlgebraicTopology.Singular.relativeDualCochainCohomologyEquiv}, \href{https://github.com/Paul-Lez/HodgeConjecture/blob/main/HodgeConjecture/Definitions/AlgebraicTopology/RelativeCochainConeNaturality.lean#L82}{\texttt{HodgeConjecture/Definitions/AlgebraicTopology/RelativeCochainConeNaturality.lean:82}}.
\end{definition}

\begin{definition}[AlgebraicTopology.Singular.relativeCochainConeCohomologyEquivCanonical]\label{def:AlgebraicTopology.Singular.relativeCochainConeCohomologyEquivCanonical}
\lean{AlgebraicTopology.Singular.relativeCochainConeCohomologyEquivCanonical}\leanok
\uses{def:AlgebraicTopology.Singular.relativeDualCochainCohomologyEquiv, def:AlgebraicTopology.Singular.relativeDualCochainHomologyIsoCone}
Relative cohomology computed from the explicit canonical cone lift.
No arbitrary completion of a triangle map enters this equivalence.

\noindent\emph{Lean:} \texttt{def AlgebraicTopology.Singular.relativeCochainConeCohomologyEquivCanonical}, \href{https://github.com/Paul-Lez/HodgeConjecture/blob/main/HodgeConjecture/Definitions/AlgebraicTopology/RelativeCochainConeNaturality.lean#L90}{\texttt{HodgeConjecture/Definitions/AlgebraicTopology/RelativeCochainConeNaturality.lean:90}}.
\end{definition}

\section{Local singular-cochain restriction cones}\label{sec:HodgeConjecture.Definitions.AlgebraicTopology.SingularCochainOpenCone}
\noindent\emph{File:} \href{https://github.com/Paul-Lez/HodgeConjecture/blob/main/HodgeConjecture/Definitions/AlgebraicTopology/SingularCochainOpenCone.lean}{\texttt{HodgeConjecture/Definitions/AlgebraicTopology/SingularCochainOpenCone.lean}}.

The comparison is the cone map of the actual sheafification units and
actual restriction maps. Cone degree \texttt{n - 1} computes relative degree \texttt{n}.

\begin{definition}[AlgebraicTopology.Singular.singularCochainOpenConeDerivedCategory]\label{def:AlgebraicTopology.Singular.singularCochainOpenConeDerivedCategory}
\lean{AlgebraicTopology.Singular.singularCochainOpenConeDerivedCategory}\leanok
\noindent\emph{Lean:} \texttt{def AlgebraicTopology.Singular.singularCochainOpenConeDerivedCategory}, \href{https://github.com/Paul-Lez/HodgeConjecture/blob/main/HodgeConjecture/Definitions/AlgebraicTopology/SingularCochainOpenCone.lean#L28}{\texttt{HodgeConjecture/Definitions/AlgebraicTopology/SingularCochainOpenCone.lean:28}}.
\end{definition}

\begin{definition}[AlgebraicTopology.Singular.openRawSingularRestrictionCone]\label{def:AlgebraicTopology.Singular.openRawSingularRestrictionCone}
\lean{AlgebraicTopology.Singular.openRawSingularRestrictionCone}\leanok
\uses{def:AlgebraicTopology.Singular.openRawSingularRestriction}
The local raw cochain restriction cone, extended by zero in negative degrees.

\noindent\emph{Lean:} \texttt{def AlgebraicTopology.Singular.openRawSingularRestrictionCone}, \href{https://github.com/Paul-Lez/HodgeConjecture/blob/main/HodgeConjecture/Definitions/AlgebraicTopology/SingularCochainOpenCone.lean#L31}{\texttt{HodgeConjecture/Definitions/AlgebraicTopology/SingularCochainOpenCone.lean:31}}.
\end{definition}

\begin{definition}[AlgebraicTopology.Singular.openSingularSheafRestrictionCone]\label{def:AlgebraicTopology.Singular.openSingularSheafRestrictionCone}
\lean{AlgebraicTopology.Singular.openSingularSheafRestrictionCone}\leanok
\uses{def:AlgebraicTopology.Singular.openSingularSheafRestriction}
The local sheaf-section restriction cone.

\noindent\emph{Lean:} \texttt{def AlgebraicTopology.Singular.openSingularSheafRestrictionCone}, \href{https://github.com/Paul-Lez/HodgeConjecture/blob/main/HodgeConjecture/Definitions/AlgebraicTopology/SingularCochainOpenCone.lean#L38}{\texttt{HodgeConjecture/Definitions/AlgebraicTopology/SingularCochainOpenCone.lean:38}}.
\end{definition}

\begin{definition}[AlgebraicTopology.Singular.openRawToSingularSheafRestrictionCone]\label{def:AlgebraicTopology.Singular.openRawToSingularSheafRestrictionCone}
\lean{AlgebraicTopology.Singular.openRawToSingularSheafRestrictionCone}\leanok
\uses{def:AlgebraicTopology.Singular.openRawSingularRestrictionCone, def:AlgebraicTopology.Singular.openRawToSingularCochainSheafComplex, def:AlgebraicTopology.Singular.openSingularSheafRestrictionCone}
The local cone comparison induced by the actual sheafification units.

\noindent\emph{Lean:} \texttt{def AlgebraicTopology.Singular.openRawToSingularSheafRestrictionCone}, \href{https://github.com/Paul-Lez/HodgeConjecture/blob/main/HodgeConjecture/Definitions/AlgebraicTopology/SingularCochainOpenCone.lean#L58}{\texttt{HodgeConjecture/Definitions/AlgebraicTopology/SingularCochainOpenCone.lean:58}}.
\end{definition}

\begin{definition}[AlgebraicTopology.Singular.openInclusionPair]\label{def:AlgebraicTopology.Singular.openInclusionPair}
\lean{AlgebraicTopology.Singular.openInclusionPair}\leanok
The actual topological pair consisting of two nested ambient opens.

\noindent\emph{Lean:} \texttt{def AlgebraicTopology.Singular.openInclusionPair}, \href{https://github.com/Paul-Lez/HodgeConjecture/blob/main/HodgeConjecture/Definitions/AlgebraicTopology/SingularCochainOpenCone.lean#L77}{\texttt{HodgeConjecture/Definitions/AlgebraicTopology/SingularCochainOpenCone.lean:77}}.
\end{definition}

\begin{definition}[AlgebraicTopology.Singular.openRawSingularCochainComplexIsoDual]\label{def:AlgebraicTopology.Singular.openRawSingularCochainComplexIsoDual}
\lean{AlgebraicTopology.Singular.openRawSingularCochainComplexIsoDual}\leanok
\uses{def:AlgebraicTopology.Singular.SingularChainComplex, def:AlgebraicTopology.Singular.openRawSingularCochainComplex, def:HomologicalComplex.linearDualCochainComplex}
Raw sections are the actual dual singular complex of the open space.

\noindent\emph{Lean:} \texttt{def AlgebraicTopology.Singular.openRawSingularCochainComplexIsoDual}, \href{https://github.com/Paul-Lez/HodgeConjecture/blob/main/HodgeConjecture/Definitions/AlgebraicTopology/SingularCochainOpenCone.lean#L84}{\texttt{HodgeConjecture/Definitions/AlgebraicTopology/SingularCochainOpenCone.lean:84}}.
\end{definition}

\begin{definition}[AlgebraicTopology.Singular.openRawSingularCochainComplexIntIsoDual]\label{def:AlgebraicTopology.Singular.openRawSingularCochainComplexIntIsoDual}
\lean{AlgebraicTopology.Singular.openRawSingularCochainComplexIntIsoDual}\leanok
\uses{def:AlgebraicTopology.Singular.openRawSingularCochainComplexIsoDual, def:HomologicalComplex.mapExtendCanonicalIso}
The actual raw cochains on an ambient open, in the integer-indexed
presentation used by relative cohomology.

\noindent\emph{Lean:} \texttt{def AlgebraicTopology.Singular.openRawSingularCochainComplexIntIsoDual}, \href{https://github.com/Paul-Lez/HodgeConjecture/blob/main/HodgeConjecture/Definitions/AlgebraicTopology/SingularCochainOpenCone.lean#L113}{\texttt{HodgeConjecture/Definitions/AlgebraicTopology/SingularCochainOpenCone.lean:113}}.
\end{definition}

\begin{definition}[AlgebraicTopology.Singular.openRawSingularRestrictionConeIsoRelative]\label{def:AlgebraicTopology.Singular.openRawSingularRestrictionConeIsoRelative}
\lean{AlgebraicTopology.Singular.openRawSingularRestrictionConeIsoRelative}\leanok
\uses{def:AlgebraicTopology.Singular.openInclusionPair, def:AlgebraicTopology.Singular.openRawSingularCochainComplexIntIsoDual, def:AlgebraicTopology.Singular.openRawSingularRestrictionCone, def:AlgebraicTopology.Singular.relativeCochainRestrictionInt}
The raw local cone is the relative-cochain cone, via its actual
restriction square.

\noindent\emph{Lean:} \texttt{def AlgebraicTopology.Singular.openRawSingularRestrictionConeIsoRelative}, \href{https://github.com/Paul-Lez/HodgeConjecture/blob/main/HodgeConjecture/Definitions/AlgebraicTopology/SingularCochainOpenCone.lean#L154}{\texttt{HodgeConjecture/Definitions/AlgebraicTopology/SingularCochainOpenCone.lean:154}}.
\end{definition}

\begin{definition}[AlgebraicTopology.Singular.openRawSingularRestrictionConeCohomologyEquivRelative]\label{def:AlgebraicTopology.Singular.openRawSingularRestrictionConeCohomologyEquivRelative}
\lean{AlgebraicTopology.Singular.openRawSingularRestrictionConeCohomologyEquivRelative}\leanok
\uses{def:AlgebraicTopology.Singular.openRawSingularRestrictionConeIsoRelative, def:AlgebraicTopology.Singular.relativeCochainConeCohomologyEquivCanonical}
The raw local cone computes relative cohomology of the literal
embedded pair \texttt{(V, W)}.

\noindent\emph{Lean:} \texttt{def AlgebraicTopology.Singular.openRawSingularRestrictionConeCohomologyEquivRelative}, \href{https://github.com/Paul-Lez/HodgeConjecture/blob/main/HodgeConjecture/Definitions/AlgebraicTopology/SingularCochainOpenCone.lean#L167}{\texttt{HodgeConjecture/Definitions/AlgebraicTopology/SingularCochainOpenCone.lean:167}}.
\end{definition}

\begin{definition}[AlgebraicTopology.Singular.openSingularSheafRestrictionConeCohomologyEquivRelative]\label{def:AlgebraicTopology.Singular.openSingularSheafRestrictionConeCohomologyEquivRelative}
\lean{AlgebraicTopology.Singular.openSingularSheafRestrictionConeCohomologyEquivRelative}\leanok
\uses{def:AlgebraicTopology.Singular.ambientOpen, def:AlgebraicTopology.Singular.coverMemberToSmallRationalSingularChains, def:AlgebraicTopology.Singular.coveringSieveOpenFamily, def:AlgebraicTopology.Singular.coveringSieveOpenFamilyOn, def:AlgebraicTopology.Singular.globalRawSingularCochainComplexIso, def:AlgebraicTopology.Singular.openFamilySieve, def:AlgebraicTopology.Singular.openRawSingularCochainComplexIsoGlobal, def:AlgebraicTopology.Singular.openRawSingularRestrictionConeCohomologyEquivRelative, def:AlgebraicTopology.Singular.openRawToSingularSheafRestrictionCone, def:AlgebraicTopology.Singular.openSimplexMap, def:AlgebraicTopology.Singular.openSingularCochainSheafComplexIsoGlobal, def:AlgebraicTopology.Singular.rationalCochainHomotopyEquivCoverSmall, def:AlgebraicTopology.Singular.rationalCochainRestrictionToCoverSmall, def:AlgebraicTopology.Singular.singularCochainComplexIsoTopOpen, def:AlgebraicTopology.Singular.singularCochainEquivSimplexFunction, def:AlgebraicTopology.Singular.singularCochainOpenConeDerivedCategory, def:AlgebraicTopology.Singular.singularCochainToPlusPlusPresheafComplex, def:AlgebraicTopology.Singular.topOpenToGlobalSingularCochainSheafComplex}
Sections of the actual singular sheaf restriction cone compute
relative rational cohomology. Apply to \texttt{Opens.infLELeft V U} for \texttt{(V, V ∩ U)}.

\noindent\emph{Lean:} \texttt{def AlgebraicTopology.Singular.openSingularSheafRestrictionConeCohomologyEquivRelative}, \href{https://github.com/Paul-Lez/HodgeConjecture/blob/main/HodgeConjecture/Definitions/AlgebraicTopology/SingularCochainOpenCone.lean#L187}{\texttt{HodgeConjecture/Definitions/AlgebraicTopology/SingularCochainOpenCone.lean:187}}.
\end{definition}

\section{Supported singular-section cohomology on arbitrary opens}\label{sec:HodgeConjecture.Definitions.AlgebraicTopology.SupportedSingularSectionCohomology}
\noindent\emph{File:} \href{https://github.com/Paul-Lez/HodgeConjecture/blob/main/HodgeConjecture/Definitions/AlgebraicTopology/SupportedSingularSectionCohomology.lean}{\texttt{HodgeConjecture/Definitions/AlgebraicTopology/SupportedSingularSectionCohomology.lean}}.

The kernel of restriction on singular cochain sheaves is compared to the relative singular
pair on an arbitrary open neighborhood. The construction composes the flasque kernel/cone
comparison, the canonical grading comparison, and the sheafification-unit cone comparison.

\begin{definition}[AlgebraicTopology.Singular.openIntersectionSupportComplementHomeomorph]\label{def:AlgebraicTopology.Singular.openIntersectionSupportComplementHomeomorph}
\lean{AlgebraicTopology.Singular.openIntersectionSupportComplementHomeomorph}\leanok
The literal intersection subspace and the complement of a support inside an open
are homeomorphic by regrouping subtype witnesses.

\noindent\emph{Lean:} \texttt{def AlgebraicTopology.Singular.openIntersectionSupportComplementHomeomorph}, \href{https://github.com/Paul-Lez/HodgeConjecture/blob/main/HodgeConjecture/Definitions/AlgebraicTopology/SupportedSingularSectionCohomology.lean#L27}{\texttt{HodgeConjecture/Definitions/AlgebraicTopology/SupportedSingularSectionCohomology.lean:27}}.
\end{definition}

\begin{definition}[AlgebraicTopology.Singular.openIntersectionPairIsoSupportComplement]\label{def:AlgebraicTopology.Singular.openIntersectionPairIsoSupportComplement}
\lean{AlgebraicTopology.Singular.openIntersectionPairIsoSupportComplement}\leanok
\uses{def:AlgebraicTopology.Singular.neighborhoodSupportComplementPair, def:AlgebraicTopology.Singular.openInclusionPair, def:AlgebraicTopology.Singular.openIntersectionSupportComplementHomeomorph}
The actual open-inclusion pair is the actual support-complement pair, preserving
the ambient open pointwise.

\noindent\emph{Lean:} \texttt{def AlgebraicTopology.Singular.openIntersectionPairIsoSupportComplement}, \href{https://github.com/Paul-Lez/HodgeConjecture/blob/main/HodgeConjecture/Definitions/AlgebraicTopology/SupportedSingularSectionCohomology.lean#L35}{\texttt{HodgeConjecture/Definitions/AlgebraicTopology/SupportedSingularSectionCohomology.lean:35}}.
\end{definition}

\begin{definition}[AlgebraicTopology.Singular.supportedRationalSingularSectionCohomologyEquivRelative]\label{def:AlgebraicTopology.Singular.supportedRationalSingularSectionCohomologyEquivRelative}
\lean{AlgebraicTopology.Singular.supportedRationalSingularSectionCohomologyEquivRelative}\leanok
\uses{def:AlgebraicTopology.Singular.openSingularSheafRestrictionConeCohomologyEquivRelative, def:AlgebraicTopology.Singular.supportedRationalSingularCochainComplex, def:TopCat.Sheaf.sectionComplexRestrictionExtendConeIso, def:TopCat.Sheaf.supportedSectionHomologyIsoRestrictionCone}
Actual local supported singular cohomology computes the literal relative pair
\texttt{(V, V ∩ U)}, in supported degree \texttt{n}.

\noindent\emph{Lean:} \texttt{def AlgebraicTopology.Singular.supportedRationalSingularSectionCohomologyEquivRelative}, \href{https://github.com/Paul-Lez/HodgeConjecture/blob/main/HodgeConjecture/Definitions/AlgebraicTopology/SupportedSingularSectionCohomology.lean#L53}{\texttt{HodgeConjecture/Definitions/AlgebraicTopology/SupportedSingularSectionCohomology.lean:53}}.
\end{definition}

\begin{definition}[AlgebraicTopology.Singular.supportedRationalSingularSectionCohomologyEquivSupportComplement]\label{def:AlgebraicTopology.Singular.supportedRationalSingularSectionCohomologyEquivSupportComplement}
\lean{AlgebraicTopology.Singular.supportedRationalSingularSectionCohomologyEquivSupportComplement}\leanok
\uses{def:AlgebraicTopology.Singular.openIntersectionPairIsoSupportComplement, def:AlgebraicTopology.Singular.supportedRationalSingularSectionCohomologyEquivRelative}
Actual local supported singular cohomology computes \texttt{(V, V \textbackslash{} S)}, with the pair
homeomorphism displayed explicitly rather than silently replacing an inclusion.

\noindent\emph{Lean:} \texttt{def AlgebraicTopology.Singular.supportedRationalSingularSectionCohomologyEquivSupportComplement}, \href{https://github.com/Paul-Lez/HodgeConjecture/blob/main/HodgeConjecture/Definitions/AlgebraicTopology/SupportedSingularSectionCohomology.lean#L68}{\texttt{HodgeConjecture/Definitions/AlgebraicTopology/SupportedSingularSectionCohomology.lean:68}}.
\end{definition}

\section{Canonical cohomology-sheaf sections from local section-complex classes}\label{sec:HodgeConjecture.Definitions.AlgebraicTopology.CohomologySheafSection}
\noindent\emph{File:} \href{https://github.com/Paul-Lez/HodgeConjecture/blob/main/HodgeConjecture/Definitions/AlgebraicTopology/CohomologySheafSection.lean}{\texttt{HodgeConjecture/Definitions/AlgebraicTopology/CohomologySheafSection.lean}}.

Exact sheafification identifies the sheafification of the presheaf of local
section-complex cohomology with the actual cohomology sheaf. Composing its unit
with that identification gives the canonical local-to-sheaf class map. This
provides the target in which normalized local purity classes can be glued.

No exactness of open-set evaluation on sheaves is assumed. The presheaf
homology is taken before sheafification, and the counit is the actual
sheafification counit on the original coefficient complex.

\begin{definition}[TopCat.Sheaf.sectionCohomologyPresheafSheafificationIso]\label{def:TopCat.Sheaf.sectionCohomologyPresheafSheafificationIso}
\lean{TopCat.Sheaf.sectionCohomologyPresheafSheafificationIso}\leanok
\uses{def:TopCat.Sheaf.sectionCohomologyPresheaf}
Sheafification of the local cohomology presheaf is the actual cohomology
sheaf, by exact sheafification and its counit.

\noindent\emph{Lean:} \texttt{def TopCat.Sheaf.sectionCohomologyPresheafSheafificationIso}, \href{https://github.com/Paul-Lez/HodgeConjecture/blob/main/HodgeConjecture/Definitions/AlgebraicTopology/CohomologySheafSection.lean#L31}{\texttt{HodgeConjecture/Definitions/AlgebraicTopology/CohomologySheafSection.lean:31}}.
\end{definition}

\begin{definition}[TopCat.Sheaf.sectionCohomologyPresheafToSheaf]\label{def:TopCat.Sheaf.sectionCohomologyPresheafToSheaf}
\lean{TopCat.Sheaf.sectionCohomologyPresheafToSheaf}\leanok
\uses{def:TopCat.Sheaf.sectionCohomologyPresheafSheafificationIso}
The canonical map from local cohomology classes to the cohomology sheaf.

\noindent\emph{Lean:} \texttt{def TopCat.Sheaf.sectionCohomologyPresheafToSheaf}, \href{https://github.com/Paul-Lez/HodgeConjecture/blob/main/HodgeConjecture/Definitions/AlgebraicTopology/CohomologySheafSection.lean#L46}{\texttt{HodgeConjecture/Definitions/AlgebraicTopology/CohomologySheafSection.lean:46}}.
\end{definition}

\begin{definition}[TopCat.Sheaf.sectionCohomologyToSheafSection]\label{def:TopCat.Sheaf.sectionCohomologyToSheafSection}
\lean{TopCat.Sheaf.sectionCohomologyToSheafSection}\leanok
\uses{def:TopCat.Sheaf.sectionCohomologyPresheafOnOpenIso, def:TopCat.Sheaf.sectionCohomologyPresheafToSheaf}
An actual cohomology class of sections on \texttt{U} determines a section of the
actual cohomology sheaf on \texttt{U}.

\noindent\emph{Lean:} \texttt{def TopCat.Sheaf.sectionCohomologyToSheafSection}, \href{https://github.com/Paul-Lez/HodgeConjecture/blob/main/HodgeConjecture/Definitions/AlgebraicTopology/CohomologySheafSection.lean#L52}{\texttt{HodgeConjecture/Definitions/AlgebraicTopology/CohomologySheafSection.lean:52}}.
\end{definition}

\section{Lowest-degree cohomology comparison for flasque complexes}\label{sec:HodgeConjecture.Definitions.AlgebraicTopology.LowestFlasqueCohomology}
\noindent\emph{File:} \href{https://github.com/Paul-Lez/HodgeConjecture/blob/main/HodgeConjecture/Definitions/AlgebraicTopology/LowestFlasqueCohomology.lean}{\texttt{HodgeConjecture/Definitions/AlgebraicTopology/LowestFlasqueCohomology.lean}}.

For a bounded-below termwise-flasque sheaf complex with cohomology sheaves zero below \texttt{n},
the canonical map from cohomology of sections on any open \texttt{U} to sections of its degree-\texttt{n}
cohomology sheaf is an isomorphism, exhibited as the sheafification unit/counit comparison.
The proof shows the degree-\texttt{n} cohomology presheaf is already a sheaf, deriving flasqueness
of the preceding cycles from the lower vanishing and cokernel preservation from that.

\begin{definition}[TopCat.Sheaf.kernelBoundaryToCyclesIso]\label{def:TopCat.Sheaf.kernelBoundaryToCyclesIso}
\lean{TopCat.Sheaf.kernelBoundaryToCyclesIso}\leanok
The kernel of the actual boundary-to-cycles map is the preceding cycle sheaf.

\noindent\emph{Lean:} \texttt{def TopCat.Sheaf.kernelBoundaryToCyclesIso}, \href{https://github.com/Paul-Lez/HodgeConjecture/blob/main/HodgeConjecture/Definitions/AlgebraicTopology/LowestFlasqueCohomology.lean#L31}{\texttt{HodgeConjecture/Definitions/AlgebraicTopology/LowestFlasqueCohomology.lean:31}}.
\end{definition}

\begin{definition}[TopCat.Sheaf.lowestSectionCohomologyIso]\label{def:TopCat.Sheaf.lowestSectionCohomologyIso}
\lean{TopCat.Sheaf.lowestSectionCohomologyIso}\leanok
\uses{def:TopCat.Sheaf.IsFlasque.BoundedBelowComplex.cyclesShortComplex, def:TopCat.Sheaf.IsFlasque.imageCokernelShortComplex, def:TopCat.Sheaf.IsFlasque.kernelImageShortComplex, def:TopCat.Sheaf.kernelBoundaryToCyclesIso, def:TopCat.Sheaf.sectionCohomologyToSheafSection}
The constructed lowest-degree isomorphism, with the actual comparison as its
forward map. Taking \texttt{U = ⊤} gives the global-sections isomorphism.

\noindent\emph{Lean:} \texttt{def TopCat.Sheaf.lowestSectionCohomologyIso}, \href{https://github.com/Paul-Lez/HodgeConjecture/blob/main/HodgeConjecture/Definitions/AlgebraicTopology/LowestFlasqueCohomology.lean#L94}{\texttt{HodgeConjecture/Definitions/AlgebraicTopology/LowestFlasqueCohomology.lean:94}}.
\end{definition}

\section{Canonical lowest-degree section comparison on an ambient open}\label{sec:HodgeConjecture.Definitions.AlgebraicTopology.OpenRestrictedLowestCohomology}
\noindent\emph{File:} \href{https://github.com/Paul-Lez/HodgeConjecture/blob/main/HodgeConjecture/Definitions/AlgebraicTopology/OpenRestrictedLowestCohomology.lean}{\texttt{HodgeConjecture/Definitions/AlgebraicTopology/OpenRestrictedLowestCohomology.lean}}.

The coefficient complex stays on the original ambient space. The actual open
restriction, its exact homology comparison, and literal equality of the top open's
image identify the lowest-degree comparison on the restricted space with an
isomorphism between ambient section cohomology and ambient cohomology-sheaf sections.

\begin{definition}[TopCat.Sheaf.openRestrictionTopSectionsIso]\label{def:TopCat.Sheaf.openRestrictionTopSectionsIso}
\lean{TopCat.Sheaf.openRestrictionTopSectionsIso}\leanok
\uses{def:TopCat.Sheaf.supportEvaluation}
Actual global sections of open restriction are ambient sections on the same open.

\noindent\emph{Lean:} \texttt{def TopCat.Sheaf.openRestrictionTopSectionsIso}, \href{https://github.com/Paul-Lez/HodgeConjecture/blob/main/HodgeConjecture/Definitions/AlgebraicTopology/OpenRestrictedLowestCohomology.lean#L28}{\texttt{HodgeConjecture/Definitions/AlgebraicTopology/OpenRestrictedLowestCohomology.lean:28}}.
\end{definition}

\begin{definition}[TopCat.Sheaf.openRestrictionTopSectionComplexIso]\label{def:TopCat.Sheaf.openRestrictionTopSectionComplexIso}
\lean{TopCat.Sheaf.openRestrictionTopSectionComplexIso}\leanok
\uses{def:TopCat.Sheaf.openRestrictionTopSectionsIso}
The same actual open equality applied to coefficient complexes.

\noindent\emph{Lean:} \texttt{def TopCat.Sheaf.openRestrictionTopSectionComplexIso}, \href{https://github.com/Paul-Lez/HodgeConjecture/blob/main/HodgeConjecture/Definitions/AlgebraicTopology/OpenRestrictedLowestCohomology.lean#L40}{\texttt{HodgeConjecture/Definitions/AlgebraicTopology/OpenRestrictedLowestCohomology.lean:40}}.
\end{definition}

\begin{definition}[TopCat.Sheaf.openRestrictionHomologyTopSectionsIso]\label{def:TopCat.Sheaf.openRestrictionHomologyTopSectionsIso}
\lean{TopCat.Sheaf.openRestrictionHomologyTopSectionsIso}\leanok
\uses{def:TopCat.Sheaf.openRestrictionTopSectionsIso}
Exact open restriction identifies sections of the two actual homology sheaves.

\noindent\emph{Lean:} \texttt{def TopCat.Sheaf.openRestrictionHomologyTopSectionsIso}, \href{https://github.com/Paul-Lez/HodgeConjecture/blob/main/HodgeConjecture/Definitions/AlgebraicTopology/OpenRestrictedLowestCohomology.lean#L48}{\texttt{HodgeConjecture/Definitions/AlgebraicTopology/OpenRestrictedLowestCohomology.lean:48}}.
\end{definition}

\begin{definition}[TopCat.Sheaf.openRestrictedLowestSectionCohomologyIso]\label{def:TopCat.Sheaf.openRestrictedLowestSectionCohomologyIso}
\lean{TopCat.Sheaf.openRestrictedLowestSectionCohomologyIso}\leanok
\uses{def:TopCat.Sheaf.lowestSectionCohomologyIso, def:TopCat.Sheaf.openRestrictionHomologyTopSectionsIso, def:TopCat.Sheaf.openRestrictionTopSectionComplexIso}
Actual lowest-degree comparison on an open where lower cohomology sheaves vanish.
Its maps are fixed by restriction, sheafification, and exact homology functoriality.

\noindent\emph{Lean:} \texttt{def TopCat.Sheaf.openRestrictedLowestSectionCohomologyIso}, \href{https://github.com/Paul-Lez/HodgeConjecture/blob/main/HodgeConjecture/Definitions/AlgebraicTopology/OpenRestrictedLowestCohomology.lean#L59}{\texttt{HodgeConjecture/Definitions/AlgebraicTopology/OpenRestrictedLowestCohomology.lean:59}}.
\end{definition}

\section{Assembling constant-sheaf maps from locally represented stalk maps}\label{sec:HodgeConjecture.Definitions.AlgebraicTopology.SheafMapOfLocallyRepresentableStalks}
\noindent\emph{File:} \href{https://github.com/Paul-Lez/HodgeConjecture/blob/main/HodgeConjecture/Definitions/AlgebraicTopology/SheafMapOfLocallyRepresentableStalks.lean}{\texttt{HodgeConjecture/Definitions/AlgebraicTopology/SheafMapOfLocallyRepresentableStalks.lean}}.

A pointwise family of additive maps from a fixed group to the stalks is not automatically
a sheaf map. This module proves the precise assembly theorem: each value must be locally
represented by a section. Unique sheaf gluing then constructs the map. The stalk formula
retains the specified maps exactly, and isomorphisms on stalks give an actual sheaf
isomorphism. Geometric applications must prove the local representability hypothesis.

\begin{definition}[TopCat.Sheaf.sectionOfLocallyRepresentable]\label{def:TopCat.Sheaf.sectionOfLocallyRepresentable}
\lean{TopCat.Sheaf.sectionOfLocallyRepresentable}\leanok
The uniquely specified global section; choice selects only the unique gluing.

\noindent\emph{Lean:} \texttt{def TopCat.Sheaf.sectionOfLocallyRepresentable}, \href{https://github.com/Paul-Lez/HodgeConjecture/blob/main/HodgeConjecture/Definitions/AlgebraicTopology/SheafMapOfLocallyRepresentableStalks.lean#L71}{\texttt{HodgeConjecture/Definitions/AlgebraicTopology/SheafMapOfLocallyRepresentableStalks.lean:71}}.
\end{definition}

\section{The local relative-cohomology presheaf and its sheafification}\label{sec:HodgeConjecture.Definitions.AlgebraicTopology.SupportRelativeCohomologySheaf}
\noindent\emph{File:} \href{https://github.com/Paul-Lez/HodgeConjecture/blob/main/HodgeConjecture/Definitions/AlgebraicTopology/SupportRelativeCohomologySheaf.lean}{\texttt{HodgeConjecture/Definitions/AlgebraicTopology/SupportRelativeCohomologySheaf.lean}}.

For a support \texttt{S ⊆ X}, the value on \texttt{V} is the rational relative cohomology of the pair
\texttt{(V, V \textbackslash{} S)} as an additive group, with restrictions the pullbacks along inclusions of
these pairs, and the sheaf is its sheafification. This module also proves local vanishing
away from a closed support.

\begin{definition}[AlgebraicTopology.Singular.supportRelativeCohomologyPresheaf]\label{def:AlgebraicTopology.Singular.supportRelativeCohomologyPresheaf}
\lean{AlgebraicTopology.Singular.supportRelativeCohomologyPresheaf}\leanok
\uses{def:AlgebraicTopology.Singular.neighborhoodSupportInclusionPairMap, def:AlgebraicTopology.Singular.relativeCohomologyMap}
The literal presheaf of rational relative cohomology of neighborhood/support pairs.
We retain the rational group but forget its scalar structure for the additive sheaf API.

\noindent\emph{Lean:} \texttt{def AlgebraicTopology.Singular.supportRelativeCohomologyPresheaf}, \href{https://github.com/Paul-Lez/HodgeConjecture/blob/main/HodgeConjecture/Definitions/AlgebraicTopology/SupportRelativeCohomologySheaf.lean#L40}{\texttt{HodgeConjecture/Definitions/AlgebraicTopology/SupportRelativeCohomologySheaf.lean:40}}.
\end{definition}

\begin{definition}[AlgebraicTopology.Singular.supportRelativeCohomologySheaf]\label{def:AlgebraicTopology.Singular.supportRelativeCohomologySheaf}
\lean{AlgebraicTopology.Singular.supportRelativeCohomologySheaf}\leanok
\uses{def:AlgebraicTopology.Singular.supportRelativeCohomologyPresheaf}
Actual sheafification, not a presupposed sheaf property of relative cohomology.

\noindent\emph{Lean:} \texttt{def AlgebraicTopology.Singular.supportRelativeCohomologySheaf}, \href{https://github.com/Paul-Lez/HodgeConjecture/blob/main/HodgeConjecture/Definitions/AlgebraicTopology/SupportRelativeCohomologySheaf.lean#L67}{\texttt{HodgeConjecture/Definitions/AlgebraicTopology/SupportRelativeCohomologySheaf.lean:67}}.
\end{definition}

\begin{definition}[AlgebraicTopology.Singular.supportRelativeCohomologyToSheaf]\label{def:AlgebraicTopology.Singular.supportRelativeCohomologyToSheaf}
\lean{AlgebraicTopology.Singular.supportRelativeCohomologyToSheaf}\leanok
\uses{def:AlgebraicTopology.Singular.supportRelativeCohomologySheaf}
The canonical map taking an actual relative class to its sheafified local section.

\noindent\emph{Lean:} \texttt{def AlgebraicTopology.Singular.supportRelativeCohomologyToSheaf}, \href{https://github.com/Paul-Lez/HodgeConjecture/blob/main/HodgeConjecture/Definitions/AlgebraicTopology/SupportRelativeCohomologySheaf.lean#L72}{\texttt{HodgeConjecture/Definitions/AlgebraicTopology/SupportRelativeCohomologySheaf.lean:72}}.
\end{definition}

\begin{definition}[AlgebraicTopology.Singular.supportRelativeCohomologyGerm]\label{def:AlgebraicTopology.Singular.supportRelativeCohomologyGerm}
\lean{AlgebraicTopology.Singular.supportRelativeCohomologyGerm}\leanok
\uses{def:AlgebraicTopology.Singular.supportRelativeCohomologyToSheaf}
Germ of an actual relative coclass in the sheafification.

\noindent\emph{Lean:} \texttt{def AlgebraicTopology.Singular.supportRelativeCohomologyGerm}, \href{https://github.com/Paul-Lez/HodgeConjecture/blob/main/HodgeConjecture/Definitions/AlgebraicTopology/SupportRelativeCohomologySheaf.lean#L77}{\texttt{HodgeConjecture/Definitions/AlgebraicTopology/SupportRelativeCohomologySheaf.lean:77}}.
\end{definition}

\section{Open-embedding transport of the relative-cohomology sheaf}\label{sec:HodgeConjecture.Definitions.AlgebraicTopology.SupportRelativeCohomologyOpenTransport}
\noindent\emph{File:} \href{https://github.com/Paul-Lez/HodgeConjecture/blob/main/HodgeConjecture/Definitions/AlgebraicTopology/SupportRelativeCohomologyOpenTransport.lean}{\texttt{HodgeConjecture/Definitions/AlgebraicTopology/SupportRelativeCohomologyOpenTransport.lean}}.

The embedding homeomorphisms identify neighborhood/support pairs with their images, and
these identifications commute with pair inclusions, giving a presheaf isomorphism.
Sheafification then transports the normalized section, with the sheafification-unit
square displayed.

\begin{definition}[AlgebraicTopology.Singular.supportRelativeCohomologyPresheafOpenIso]\label{def:AlgebraicTopology.Singular.supportRelativeCohomologyPresheafOpenIso}
\lean{AlgebraicTopology.Singular.supportRelativeCohomologyPresheafOpenIso}\leanok
\uses{def:AlgebraicTopology.Singular.neighborhoodSupportPairImageCohomologyEquiv, def:AlgebraicTopology.Singular.supportRelativeCohomologyPresheaf}
The presheaf comparison comes from the actual pair-image homeomorphisms.

\noindent\emph{Lean:} \texttt{def AlgebraicTopology.Singular.supportRelativeCohomologyPresheafOpenIso}, \href{https://github.com/Paul-Lez/HodgeConjecture/blob/main/HodgeConjecture/Definitions/AlgebraicTopology/SupportRelativeCohomologyOpenTransport.lean#L43}{\texttt{HodgeConjecture/Definitions/AlgebraicTopology/SupportRelativeCohomologyOpenTransport.lean:43}}.
\end{definition}

\begin{definition}[AlgebraicTopology.Singular.supportOpenEmbeddingSheafificationIso]\label{def:AlgebraicTopology.Singular.supportOpenEmbeddingSheafificationIso}
\lean{AlgebraicTopology.Singular.supportOpenEmbeddingSheafificationIso}\leanok
The open-image functor commutes with sheafification, through the actual
restricted unit. This is the general open-embedding version of open restriction.

\noindent\emph{Lean:} \texttt{def AlgebraicTopology.Singular.supportOpenEmbeddingSheafificationIso}, \href{https://github.com/Paul-Lez/HodgeConjecture/blob/main/HodgeConjecture/Definitions/AlgebraicTopology/SupportRelativeCohomologyOpenTransport.lean#L69}{\texttt{HodgeConjecture/Definitions/AlgebraicTopology/SupportRelativeCohomologyOpenTransport.lean:69}}.
\end{definition}

\begin{definition}[AlgebraicTopology.Singular.supportRelativeCohomologySheafOpenIso]\label{def:AlgebraicTopology.Singular.supportRelativeCohomologySheafOpenIso}
\lean{AlgebraicTopology.Singular.supportRelativeCohomologySheafOpenIso}\leanok
\uses{def:AlgebraicTopology.Singular.supportOpenEmbeddingSheafificationIso, def:AlgebraicTopology.Singular.supportRelativeCohomologyPresheafOpenIso, def:AlgebraicTopology.Singular.supportRelativeCohomologySheaf}
Actual sheafification of the pair comparison identifies intrinsic local support
cohomology with restriction of the original ambient relative-cohomology sheaf.

\noindent\emph{Lean:} \texttt{def AlgebraicTopology.Singular.supportRelativeCohomologySheafOpenIso}, \href{https://github.com/Paul-Lez/HodgeConjecture/blob/main/HodgeConjecture/Definitions/AlgebraicTopology/SupportRelativeCohomologyOpenTransport.lean#L82}{\texttt{HodgeConjecture/Definitions/AlgebraicTopology/SupportRelativeCohomologyOpenTransport.lean:82}}.
\end{definition}

\begin{definition}[AlgebraicTopology.Singular.supportRelativeCohomologySectionOpenImage]\label{def:AlgebraicTopology.Singular.supportRelativeCohomologySectionOpenImage}
\lean{AlgebraicTopology.Singular.supportRelativeCohomologySectionOpenImage}\leanok
\uses{def:AlgebraicTopology.Singular.supportRelativeCohomologySheafOpenIso}
A genuine section on an auxiliary open ambient space gives a section on its
actual image open in the original ambient space.

\noindent\emph{Lean:} \texttt{def AlgebraicTopology.Singular.supportRelativeCohomologySectionOpenImage}, \href{https://github.com/Paul-Lez/HodgeConjecture/blob/main/HodgeConjecture/Definitions/AlgebraicTopology/SupportRelativeCohomologyOpenTransport.lean#L91}{\texttt{HodgeConjecture/Definitions/AlgebraicTopology/SupportRelativeCohomologyOpenTransport.lean:91}}.
\end{definition}

\begin{definition}[AlgebraicTopology.Singular.supportRelativeCohomologySectionOnOpen]\label{def:AlgebraicTopology.Singular.supportRelativeCohomologySectionOnOpen}
\lean{AlgebraicTopology.Singular.supportRelativeCohomologySectionOnOpen}\leanok
\uses{def:AlgebraicTopology.Singular.supportRelativeCohomologySectionOpenImage}
Transport to a specified ambient open equal to the actual image. The final
identification is the unique open inclusion, not an arbitrary section equivalence.

\noindent\emph{Lean:} \texttt{def AlgebraicTopology.Singular.supportRelativeCohomologySectionOnOpen}, \href{https://github.com/Paul-Lez/HodgeConjecture/blob/main/HodgeConjecture/Definitions/AlgebraicTopology/SupportRelativeCohomologyOpenTransport.lean#L98}{\texttt{HodgeConjecture/Definitions/AlgebraicTopology/SupportRelativeCohomologyOpenTransport.lean:98}}.
\end{definition}

\section{The supported singular cohomology sheaf and local relative cohomology}\label{sec:HodgeConjecture.Definitions.AlgebraicTopology.SupportedSingularCohomologySheafComparison}
\noindent\emph{File:} \href{https://github.com/Paul-Lez/HodgeConjecture/blob/main/HodgeConjecture/Definitions/AlgebraicTopology/SupportedSingularCohomologySheafComparison.lean}{\texttt{HodgeConjecture/Definitions/AlgebraicTopology/SupportedSingularCohomologySheafComparison.lean}}.

The restriction-natural local comparison is an isomorphism of presheaves, and exact
sheafification with its counit identifies the sheafification with the cohomology sheaf of
the kernel-defined supported singular complex. The local class and unit equations fix this
identification, including the positive short-exact-sequence lift normalization.

\begin{definition}[AlgebraicTopology.Singular.supportedSingularCohomologyPresheafIsoRelative]\label{def:AlgebraicTopology.Singular.supportedSingularCohomologyPresheafIsoRelative}
\lean{AlgebraicTopology.Singular.supportedSingularCohomologyPresheafIsoRelative}\leanok
\uses{def:AlgebraicTopology.Singular.forgottenRelativeCochainConeMap, def:AlgebraicTopology.Singular.openInclusionPairMap, def:AlgebraicTopology.Singular.openRawSingularRestrictionConeHomologyIso, def:AlgebraicTopology.Singular.openRawSingularRestrictionConeMap, def:AlgebraicTopology.Singular.openRawToSingularSheafRestrictionConeHomologyIso, def:AlgebraicTopology.Singular.openSingularSheafRestrictionConeMap, def:AlgebraicTopology.Singular.relativeCochainConeMap, def:AlgebraicTopology.Singular.supportRelativeCohomologyPresheaf, def:AlgebraicTopology.Singular.supportedRationalSingularSectionCohomologyEquivSupportComplement, def:AlgebraicTopology.Singular.supportedSingularSectionConeHomologyIso, def:TopCat.Sheaf.sectionCohomologyPresheafOnOpenIso, def:TopCat.Sheaf.sectionComplexRestrictionConeMap, def:TopCat.Sheaf.supportRestrictionSectionsComplexMap}
The actual local cohomology presheaf of supported singular cochains is the
literal relative-cohomology presheaf, with its literal pair restrictions.

\noindent\emph{Lean:} \texttt{def AlgebraicTopology.Singular.supportedSingularCohomologyPresheafIsoRelative}, \href{https://github.com/Paul-Lez/HodgeConjecture/blob/main/HodgeConjecture/Definitions/AlgebraicTopology/SupportedSingularCohomologySheafComparison.lean#L57}{\texttt{HodgeConjecture/Definitions/AlgebraicTopology/SupportedSingularCohomologySheafComparison.lean:57}}.
\end{definition}

\begin{definition}[AlgebraicTopology.Singular.supportedSingularCohomologySheafIsoRelative]\label{def:AlgebraicTopology.Singular.supportedSingularCohomologySheafIsoRelative}
\lean{AlgebraicTopology.Singular.supportedSingularCohomologySheafIsoRelative}\leanok
\uses{def:AlgebraicTopology.Singular.supportRelativeCohomologySheaf, def:AlgebraicTopology.Singular.supportedSingularCohomologyPresheafIsoRelative, def:TopCat.Sheaf.sectionCohomologyPresheafSheafificationIso}
Exact sheafification identifies the actual supported cohomology sheaf with
the sheafification of literal neighborhood/support relative cohomology.

\noindent\emph{Lean:} \texttt{def AlgebraicTopology.Singular.supportedSingularCohomologySheafIsoRelative}, \href{https://github.com/Paul-Lez/HodgeConjecture/blob/main/HodgeConjecture/Definitions/AlgebraicTopology/SupportedSingularCohomologySheafComparison.lean#L86}{\texttt{HodgeConjecture/Definitions/AlgebraicTopology/SupportedSingularCohomologySheafComparison.lean:86}}.
\end{definition}

\chapter{Definitions: Topology}

\section{The dimension of an irreducible topological space}\label{sec:HodgeConjecture.Definitions.Topology.Dimension}
\noindent\emph{File:} \href{https://github.com/Paul-Lez/HodgeConjecture/blob/main/HodgeConjecture/Definitions/Topology/Dimension.lean}{\texttt{HodgeConjecture/Definitions/Topology/Dimension.lean}}.

\begin{definition}[TopologicalSpace.dim]\label{def:TopologicalSpace.dim}
\lean{TopologicalSpace.dim}\leanok
The dimension of an irreducible topological space is its Krull dimension: the supremum of the
lengths of chains of irreducible closed subsets.

Truncating to \texttt{ℕ} collapses the two degenerate values of \texttt{topologicalKrullDim}: the empty space,
which the \texttt{IrreducibleSpace} hypothesis rules out, and infinite-dimensional spaces, which come out
as \texttt{0}. Irreducibility also makes this *the* dimension of the space, rather than the maximum of the
dimensions of its irreducible components.

\noindent\emph{Lean:} \texttt{def TopologicalSpace.dim}, \href{https://github.com/Paul-Lez/HodgeConjecture/blob/main/HodgeConjecture/Definitions/Topology/Dimension.lean#L30}{\texttt{HodgeConjecture/Definitions/Topology/Dimension.lean:30}}.
\end{definition}

\chapter{Lemmas: Algebra}

\section{Naturality of linear duality on homology}\label{sec:HodgeConjecture.Lemmas.Algebra.Homology.LinearDualNaturality}
\noindent\emph{File:} \href{https://github.com/Paul-Lez/HodgeConjecture/blob/main/HodgeConjecture/Lemmas/Algebra/Homology/LinearDualNaturality.lean}{\texttt{HodgeConjecture/Lemmas/Algebra/Homology/LinearDualNaturality.lean}}.

Concrete cycle representatives and the canonical universal-coefficient equivalence are
compatible with actual morphisms of short complexes. These identities are used to descend
geometric cap-product identities in the cohomology variable.

The same holds one level up: the universal-coefficient equivalence
\texttt{HomologicalComplex.linearDualHomologyEquiv} between the cohomology of the linear-dual cochain
complex and the dual of homology intertwines the dualised cochain map with the dual of the
homology map.

\begin{definition}[CategoryTheory.ShortComplex.moduleCatHomologyClass]\label{def:CategoryTheory.ShortComplex.moduleCatHomologyClass}
\lean{CategoryTheory.ShortComplex.moduleCatHomologyClass}\leanok
The canonical class of an explicit cycle in a short complex of vector spaces.

\noindent\emph{Lean:} \texttt{def CategoryTheory.ShortComplex.moduleCatHomologyClass}, \href{https://github.com/Paul-Lez/HodgeConjecture/blob/main/HodgeConjecture/Lemmas/Algebra/Homology/LinearDualNaturality.lean#L43}{\texttt{HodgeConjecture/Lemmas/Algebra/Homology/LinearDualNaturality.lean:43}}.
\end{definition}

\begin{definition}[CategoryTheory.ShortComplex.moduleCatCycleMap]\label{def:CategoryTheory.ShortComplex.moduleCatCycleMap}
\lean{CategoryTheory.ShortComplex.moduleCatCycleMap}\leanok
A morphism of short complexes sends cycles to cycles.

\noindent\emph{Lean:} \texttt{def CategoryTheory.ShortComplex.moduleCatCycleMap}, \href{https://github.com/Paul-Lez/HodgeConjecture/blob/main/HodgeConjecture/Lemmas/Algebra/Homology/LinearDualNaturality.lean#L55}{\texttt{HodgeConjecture/Lemmas/Algebra/Homology/LinearDualNaturality.lean:55}}.
\end{definition}

\begin{definition}[CategoryTheory.ShortComplex.linearDualMap]\label{def:CategoryTheory.ShortComplex.linearDualMap}
\lean{CategoryTheory.ShortComplex.linearDualMap}\leanok
\uses{def:CategoryTheory.ShortComplex.linearDual}
Reversed linear-dual short complexes are contravariantly functorial.

\noindent\emph{Lean:} \texttt{def CategoryTheory.ShortComplex.linearDualMap}, \href{https://github.com/Paul-Lez/HodgeConjecture/blob/main/HodgeConjecture/Lemmas/Algebra/Homology/LinearDualNaturality.lean#L90}{\texttt{HodgeConjecture/Lemmas/Algebra/Homology/LinearDualNaturality.lean:90}}.
\end{definition}

\chapter{Lemmas: AlgebraicGeometry}

\section{Complex points of affine space}\label{sec:HodgeConjecture.Lemmas.AlgebraicGeometry.ComplexAffineSpace}
\noindent\emph{File:} \href{https://github.com/Paul-Lez/HodgeConjecture/blob/main/HodgeConjecture/Lemmas/AlgebraicGeometry/ComplexAffineSpace.lean}{\texttt{HodgeConjecture/Lemmas/AlgebraicGeometry/ComplexAffineSpace.lean}}.

The complex points of algebraic affine space are canonically homeomorphic to tuples of complex
numbers. The proof identifies global regular functions with multivariate polynomials. It then uses
the localization description of sections on principal opens to handle every local regular function
appearing in the analytic topology.

\begin{definition}[AlgebraicGeometry.ComplexPoint.complexAffineSpace]\label{def:AlgebraicGeometry.ComplexPoint.complexAffineSpace}
\lean{AlgebraicGeometry.ComplexPoint.complexAffineSpace}\leanok
Complex affine space as a scheme over \texttt{Spec ℂ}.

\noindent\emph{Lean:} \texttt{abbrev AlgebraicGeometry.ComplexPoint.complexAffineSpace}, \href{https://github.com/Paul-Lez/HodgeConjecture/blob/main/HodgeConjecture/Lemmas/AlgebraicGeometry/ComplexAffineSpace.lean#L62}{\texttt{HodgeConjecture/Lemmas/AlgebraicGeometry/ComplexAffineSpace.lean:62}}.
\end{definition}

\begin{definition}[AlgebraicGeometry.ComplexPoint.affineSpaceEquiv]\label{def:AlgebraicGeometry.ComplexPoint.affineSpaceEquiv}
\lean{AlgebraicGeometry.ComplexPoint.affineSpaceEquiv}\leanok
\uses{def:AlgebraicGeometry.ComplexPoint, def:AlgebraicGeometry.ComplexPoint.complexAffineSpace}
Algebraic coordinates identify complex points of affine space with tuples of complex numbers.

\noindent\emph{Lean:} \texttt{def AlgebraicGeometry.ComplexPoint.affineSpaceEquiv}, \href{https://github.com/Paul-Lez/HodgeConjecture/blob/main/HodgeConjecture/Lemmas/AlgebraicGeometry/ComplexAffineSpace.lean#L66}{\texttt{HodgeConjecture/Lemmas/AlgebraicGeometry/ComplexAffineSpace.lean:66}}.
\end{definition}

\begin{definition}[AlgebraicGeometry.ComplexPoint.affineGlobalPolynomial]\label{def:AlgebraicGeometry.ComplexPoint.affineGlobalPolynomial}
\lean{AlgebraicGeometry.ComplexPoint.affineGlobalPolynomial}\leanok
\uses{def:AlgebraicGeometry.ComplexPoint.complexAffineSpace}
The polynomial represented by a global regular function on complex affine space.

\noindent\emph{Lean:} \texttt{def AlgebraicGeometry.ComplexPoint.affineGlobalPolynomial}, \href{https://github.com/Paul-Lez/HodgeConjecture/blob/main/HodgeConjecture/Lemmas/AlgebraicGeometry/ComplexAffineSpace.lean#L108}{\texttt{HodgeConjecture/Lemmas/AlgebraicGeometry/ComplexAffineSpace.lean:108}}.
\end{definition}

\begin{definition}[AlgebraicGeometry.ComplexPoint.affineSpaceHomeomorph]\label{def:AlgebraicGeometry.ComplexPoint.affineSpaceHomeomorph}
\lean{AlgebraicGeometry.ComplexPoint.affineSpaceHomeomorph}\leanok
\uses{def:AlgebraicGeometry.ComplexPoint.affineGlobalPolynomial, def:AlgebraicGeometry.ComplexPoint.affineSpaceEquiv, def:AlgebraicGeometry.Point.analyticSubbasis, def:AlgebraicGeometry.Point.chartTopology}
Complex affine space with the evaluation topology is ordinary complex affine space.

\noindent\emph{Lean:} \texttt{def AlgebraicGeometry.ComplexPoint.affineSpaceHomeomorph}, \href{https://github.com/Paul-Lez/HodgeConjecture/blob/main/HodgeConjecture/Lemmas/AlgebraicGeometry/ComplexAffineSpace.lean#L271}{\texttt{HodgeConjecture/Lemmas/AlgebraicGeometry/ComplexAffineSpace.lean:271}}.
\end{definition}

\section{Complex points of affine schemes}\label{sec:HodgeConjecture.Lemmas.AlgebraicGeometry.ComplexAffineScheme}
\noindent\emph{File:} \href{https://github.com/Paul-Lez/HodgeConjecture/blob/main/HodgeConjecture/Lemmas/AlgebraicGeometry/ComplexAffineScheme.lean}{\texttt{HodgeConjecture/Lemmas/AlgebraicGeometry/ComplexAffineScheme.lean}}.

The complex points of \texttt{Spec R} are canonically homeomorphic to the space of complex algebra
homomorphisms from \texttt{R} to \texttt{ℂ}, equipped with the topology of pointwise convergence. The proof of
continuity in the reverse direction uses the localization description of regular functions on
principal opens.

\begin{definition}[AlgebraicGeometry.ComplexPoint.affineSpecStructureMap]\label{def:AlgebraicGeometry.ComplexPoint.affineSpecStructureMap}
\lean{AlgebraicGeometry.ComplexPoint.affineSpecStructureMap}\leanok
The structure morphism of an affine complex scheme.

\noindent\emph{Lean:} \texttt{abbrev AlgebraicGeometry.ComplexPoint.affineSpecStructureMap}, \href{https://github.com/Paul-Lez/HodgeConjecture/blob/main/HodgeConjecture/Lemmas/AlgebraicGeometry/ComplexAffineScheme.lean#L48}{\texttt{HodgeConjecture/Lemmas/AlgebraicGeometry/ComplexAffineScheme.lean:48}}.
\end{definition}

\begin{definition}[AlgebraicGeometry.ComplexPoint.affineSpecEquiv]\label{def:AlgebraicGeometry.ComplexPoint.affineSpecEquiv}
\lean{AlgebraicGeometry.ComplexPoint.affineSpecEquiv}\leanok
\uses{def:AlgebraicGeometry.ComplexPoint, def:AlgebraicGeometry.ComplexPoint.affineSpecStructureMap}
Complex points of \texttt{Spec R} correspond to complex algebra homomorphisms from \texttt{R} to \texttt{ℂ}.

\noindent\emph{Lean:} \texttt{def AlgebraicGeometry.ComplexPoint.affineSpecEquiv}, \href{https://github.com/Paul-Lez/HodgeConjecture/blob/main/HodgeConjecture/Lemmas/AlgebraicGeometry/ComplexAffineScheme.lean#L56}{\texttt{HodgeConjecture/Lemmas/AlgebraicGeometry/ComplexAffineScheme.lean:56}}.
\end{definition}

\begin{definition}[AlgebraicGeometry.ComplexPoint.affineSpecHomeomorph]\label{def:AlgebraicGeometry.ComplexPoint.affineSpecHomeomorph}
\lean{AlgebraicGeometry.ComplexPoint.affineSpecHomeomorph}\leanok
\uses{def:AlgebraicGeometry.ComplexPoint.affineSpecEquiv, def:AlgebraicGeometry.Point.analyticSubbasis, def:AlgebraicGeometry.Point.chartTopology}
The complex points of \texttt{Spec R} are homeomorphic to the complex algebra homomorphisms
from \texttt{R} to \texttt{ℂ}, with their topology of pointwise convergence.

\noindent\emph{Lean:} \texttt{def AlgebraicGeometry.ComplexPoint.affineSpecHomeomorph}, \href{https://github.com/Paul-Lez/HodgeConjecture/blob/main/HodgeConjecture/Lemmas/AlgebraicGeometry/ComplexAffineScheme.lean#L214}{\texttt{HodgeConjecture/Lemmas/AlgebraicGeometry/ComplexAffineScheme.lean:214}}.
\end{definition}

\begin{definition}[AlgebraicGeometry.ComplexPoint.affineSpecComplexPointMap]\label{def:AlgebraicGeometry.ComplexPoint.affineSpecComplexPointMap}
\lean{AlgebraicGeometry.ComplexPoint.affineSpecComplexPointMap}\leanok
\uses{def:AlgebraicGeometry.ComplexPoint, def:AlgebraicGeometry.ComplexPoint.affineSpecStructureMap, def:AlgebraicGeometry.Point.map}
The map on complex points contravariantly associated to a complex algebra homomorphism.

\noindent\emph{Lean:} \texttt{def AlgebraicGeometry.ComplexPoint.affineSpecComplexPointMap}, \href{https://github.com/Paul-Lez/HodgeConjecture/blob/main/HodgeConjecture/Lemmas/AlgebraicGeometry/ComplexAffineScheme.lean#L234}{\texttt{HodgeConjecture/Lemmas/AlgebraicGeometry/ComplexAffineScheme.lean:234}}.
\end{definition}

\section{Complex points of a localization}\label{sec:HodgeConjecture.Lemmas.AlgebraicGeometry.ComplexLocalization}
\noindent\emph{File:} \href{https://github.com/Paul-Lez/HodgeConjecture/blob/main/HodgeConjecture/Lemmas/AlgebraicGeometry/ComplexLocalization.lean}{\texttt{HodgeConjecture/Lemmas/AlgebraicGeometry/ComplexLocalization.lean}}.

Complex algebra homomorphisms from \texttt{S[1/f]} are canonically homeomorphic to the open subspace of
complex algebra homomorphisms from \texttt{S} at which \texttt{f} is nonzero. The forward map is restriction.
Continuity of the extension map follows by representing a localized element as a quotient whose
denominator is a power of \texttt{f}.

Nothing here mentions a complex point of a scheme, so it lives in namespace
\texttt{AlgebraicGeometry.ComplexAlgHom}: the algebraic model of the complex points of an affine scheme,
as the space of \texttt{ℂ}-algebra homomorphisms out of its ring of sections.

\begin{definition}[AlgebraicGeometry.ComplexAlgHom.precompAlgEquiv]\label{def:AlgebraicGeometry.ComplexAlgHom.precompAlgEquiv}
\lean{AlgebraicGeometry.ComplexAlgHom.precompAlgEquiv}\leanok
Precomposition by a complex algebra equivalence, as an equivalence of complex-valued algebra
homomorphisms.

\noindent\emph{Lean:} \texttt{def AlgebraicGeometry.ComplexAlgHom.precompAlgEquiv}, \href{https://github.com/Paul-Lez/HodgeConjecture/blob/main/HodgeConjecture/Lemmas/AlgebraicGeometry/ComplexLocalization.lean#L48}{\texttt{HodgeConjecture/Lemmas/AlgebraicGeometry/ComplexLocalization.lean:48}}.
\end{definition}

\begin{definition}[AlgebraicGeometry.ComplexAlgHom.precompAlgEquivHomeomorph]\label{def:AlgebraicGeometry.ComplexAlgHom.precompAlgEquivHomeomorph}
\lean{AlgebraicGeometry.ComplexAlgHom.precompAlgEquivHomeomorph}\leanok
\uses{def:AlgebraicGeometry.ComplexAlgHom.precompAlgEquiv}
Precomposition by a complex algebra equivalence is a homeomorphism for pointwise
convergence.

\noindent\emph{Lean:} \texttt{def AlgebraicGeometry.ComplexAlgHom.precompAlgEquivHomeomorph}, \href{https://github.com/Paul-Lez/HodgeConjecture/blob/main/HodgeConjecture/Lemmas/AlgebraicGeometry/ComplexLocalization.lean#L63}{\texttt{HodgeConjecture/Lemmas/AlgebraicGeometry/ComplexLocalization.lean:63}}.
\end{definition}

\begin{definition}[AlgebraicGeometry.ComplexAlgHom.nonvanishingAlgHom]\label{def:AlgebraicGeometry.ComplexAlgHom.nonvanishingAlgHom}
\lean{AlgebraicGeometry.ComplexAlgHom.nonvanishingAlgHom}\leanok
The open subspace of complex algebra homomorphisms at which \texttt{f} does not vanish.

\noindent\emph{Lean:} \texttt{abbrev AlgebraicGeometry.ComplexAlgHom.nonvanishingAlgHom}, \href{https://github.com/Paul-Lez/HodgeConjecture/blob/main/HodgeConjecture/Lemmas/AlgebraicGeometry/ComplexLocalization.lean#L73}{\texttt{HodgeConjecture/Lemmas/AlgebraicGeometry/ComplexLocalization.lean:73}}.
\end{definition}

\begin{definition}[AlgebraicGeometry.ComplexAlgHom.localizationAwayAlgHomRestriction]\label{def:AlgebraicGeometry.ComplexAlgHom.localizationAwayAlgHomRestriction}
\lean{AlgebraicGeometry.ComplexAlgHom.localizationAwayAlgHomRestriction}\leanok
\uses{def:AlgebraicGeometry.ComplexAlgHom.nonvanishingAlgHom}
Restriction of an algebra homomorphism from \texttt{S[1/f]} to \texttt{S}, bundled with its nonvanishing
property.

\noindent\emph{Lean:} \texttt{def AlgebraicGeometry.ComplexAlgHom.localizationAwayAlgHomRestriction}, \href{https://github.com/Paul-Lez/HodgeConjecture/blob/main/HodgeConjecture/Lemmas/AlgebraicGeometry/ComplexLocalization.lean#L76}{\texttt{HodgeConjecture/Lemmas/AlgebraicGeometry/ComplexLocalization.lean:76}}.
\end{definition}

\begin{definition}[AlgebraicGeometry.ComplexAlgHom.nonvanishingAlgHomExtension]\label{def:AlgebraicGeometry.ComplexAlgHom.nonvanishingAlgHomExtension}
\lean{AlgebraicGeometry.ComplexAlgHom.nonvanishingAlgHomExtension}\leanok
\uses{def:AlgebraicGeometry.ComplexAlgHom.nonvanishingAlgHom}
Extension of an algebra homomorphism on which \texttt{f} is nonzero to \texttt{S[1/f]}.

\noindent\emph{Lean:} \texttt{def AlgebraicGeometry.ComplexAlgHom.nonvanishingAlgHomExtension}, \href{https://github.com/Paul-Lez/HodgeConjecture/blob/main/HodgeConjecture/Lemmas/AlgebraicGeometry/ComplexLocalization.lean#L83}{\texttt{HodgeConjecture/Lemmas/AlgebraicGeometry/ComplexLocalization.lean:83}}.
\end{definition}

\begin{definition}[AlgebraicGeometry.ComplexAlgHom.localizationAwayAlgHomEquiv]\label{def:AlgebraicGeometry.ComplexAlgHom.localizationAwayAlgHomEquiv}
\lean{AlgebraicGeometry.ComplexAlgHom.localizationAwayAlgHomEquiv}\leanok
\uses{def:AlgebraicGeometry.ComplexAlgHom.localizationAwayAlgHomRestriction, def:AlgebraicGeometry.ComplexAlgHom.nonvanishingAlgHomExtension}
The algebraic equivalence between homomorphisms from \texttt{S[1/f]} and homomorphisms from \texttt{S} on
which \texttt{f} does not vanish.

\noindent\emph{Lean:} \texttt{def AlgebraicGeometry.ComplexAlgHom.localizationAwayAlgHomEquiv}, \href{https://github.com/Paul-Lez/HodgeConjecture/blob/main/HodgeConjecture/Lemmas/AlgebraicGeometry/ComplexLocalization.lean#L102}{\texttt{HodgeConjecture/Lemmas/AlgebraicGeometry/ComplexLocalization.lean:102}}.
\end{definition}

\begin{definition}[AlgebraicGeometry.ComplexAlgHom.localizationAwayAlgHomHomeomorph]\label{def:AlgebraicGeometry.ComplexAlgHom.localizationAwayAlgHomHomeomorph}
\lean{AlgebraicGeometry.ComplexAlgHom.localizationAwayAlgHomHomeomorph}\leanok
\uses{def:AlgebraicGeometry.ComplexAlgHom.localizationAwayAlgHomEquiv}
Complex points of \texttt{S[1/f]} are homeomorphic to the nonvanishing locus of \texttt{f} in the complex
points of \texttt{S}.

\noindent\emph{Lean:} \texttt{def AlgebraicGeometry.ComplexAlgHom.localizationAwayAlgHomHomeomorph}, \href{https://github.com/Paul-Lez/HodgeConjecture/blob/main/HodgeConjecture/Lemmas/AlgebraicGeometry/ComplexLocalization.lean#L149}{\texttt{HodgeConjecture/Lemmas/AlgebraicGeometry/ComplexLocalization.lean:149}}.
\end{definition}

\begin{definition}[AlgebraicGeometry.ComplexAlgHom.localizationAwayAlgHomMap]\label{def:AlgebraicGeometry.ComplexAlgHom.localizationAwayAlgHomMap}
\lean{AlgebraicGeometry.ComplexAlgHom.localizationAwayAlgHomMap}\leanok
\uses{def:AlgebraicGeometry.ComplexAlgHom.localizationAwayAlgHomRestriction}
Restriction from complex points of \texttt{S[1/f]} to complex points of \texttt{S}, with the nonvanishing
witness forgotten.

\noindent\emph{Lean:} \texttt{def AlgebraicGeometry.ComplexAlgHom.localizationAwayAlgHomMap}, \href{https://github.com/Paul-Lez/HodgeConjecture/blob/main/HodgeConjecture/Lemmas/AlgebraicGeometry/ComplexLocalization.lean#L157}{\texttt{HodgeConjecture/Lemmas/AlgebraicGeometry/ComplexLocalization.lean:157}}.
\end{definition}

\section{Complex points of standard étale algebras}\label{sec:HodgeConjecture.Lemmas.AlgebraicGeometry.ComplexStandardEtale}
\noindent\emph{File:} \href{https://github.com/Paul-Lez/HodgeConjecture/blob/main/HodgeConjecture/Lemmas/AlgebraicGeometry/ComplexStandardEtale.lean}{\texttt{HodgeConjecture/Lemmas/AlgebraicGeometry/ComplexStandardEtale.lean}}.

A complex point of a standard étale algebra over a polynomial ring is a base point together with
a root of the defining polynomial at which the localization polynomial does not vanish. This file
constructs that correspondence as a homeomorphism. The inverse continuity proof evaluates an
arbitrary representative in the explicit bivariate-polynomial quotient presentation.

Like \texttt{ComplexLocalization}, this works throughout with \texttt{ℂ}-algebra homomorphisms rather than with
complex points of schemes, so it lives in namespace \texttt{AlgebraicGeometry.ComplexAlgHom}.

\begin{definition}[AlgebraicGeometry.ComplexAlgHom.complexPolynomialRing]\label{def:AlgebraicGeometry.ComplexAlgHom.complexPolynomialRing}
\lean{AlgebraicGeometry.ComplexAlgHom.complexPolynomialRing}\leanok
The polynomial coordinate ring of complex affine \texttt{n}-space.

\noindent\emph{Lean:} \texttt{abbrev AlgebraicGeometry.ComplexAlgHom.complexPolynomialRing}, \href{https://github.com/Paul-Lez/HodgeConjecture/blob/main/HodgeConjecture/Lemmas/AlgebraicGeometry/ComplexStandardEtale.lean#L48}{\texttt{HodgeConjecture/Lemmas/AlgebraicGeometry/ComplexStandardEtale.lean:48}}.
\end{definition}

\begin{definition}[AlgebraicGeometry.ComplexAlgHom.mvPolynomialAlgHomEquiv]\label{def:AlgebraicGeometry.ComplexAlgHom.mvPolynomialAlgHomEquiv}
\lean{AlgebraicGeometry.ComplexAlgHom.mvPolynomialAlgHomEquiv}\leanok
\uses{def:AlgebraicGeometry.ComplexAlgHom.complexPolynomialRing}
A complex algebra homomorphism out of a polynomial ring is determined by the images of its
variables.

\noindent\emph{Lean:} \texttt{def AlgebraicGeometry.ComplexAlgHom.mvPolynomialAlgHomEquiv}, \href{https://github.com/Paul-Lez/HodgeConjecture/blob/main/HodgeConjecture/Lemmas/AlgebraicGeometry/ComplexStandardEtale.lean#L51}{\texttt{HodgeConjecture/Lemmas/AlgebraicGeometry/ComplexStandardEtale.lean:51}}.
\end{definition}

\begin{definition}[AlgebraicGeometry.ComplexAlgHom.mvPolynomialAlgHomHomeomorph]\label{def:AlgebraicGeometry.ComplexAlgHom.mvPolynomialAlgHomHomeomorph}
\lean{AlgebraicGeometry.ComplexAlgHom.mvPolynomialAlgHomHomeomorph}\leanok
\uses{def:AlgebraicGeometry.ComplexAlgHom.mvPolynomialAlgHomEquiv}
Pointwise convergence on complex algebra homomorphisms out of a polynomial ring is the usual
Euclidean topology on the tuple of variable values.

\noindent\emph{Lean:} \texttt{def AlgebraicGeometry.ComplexAlgHom.mvPolynomialAlgHomHomeomorph}, \href{https://github.com/Paul-Lez/HodgeConjecture/blob/main/HodgeConjecture/Lemmas/AlgebraicGeometry/ComplexStandardEtale.lean#L71}{\texttt{HodgeConjecture/Lemmas/AlgebraicGeometry/ComplexStandardEtale.lean:71}}.
\end{definition}

\begin{definition}[AlgebraicGeometry.ComplexAlgHom.standardEtalePointSpace]\label{def:AlgebraicGeometry.ComplexAlgHom.standardEtalePointSpace}
\lean{AlgebraicGeometry.ComplexAlgHom.standardEtalePointSpace}\leanok
\uses{def:AlgebraicGeometry.ComplexAlgHom.complexPolynomialRing}
A base point and a root satisfying the equations represented by a standard étale pair.

\noindent\emph{Lean:} \texttt{abbrev AlgebraicGeometry.ComplexAlgHom.standardEtalePointSpace}, \href{https://github.com/Paul-Lez/HodgeConjecture/blob/main/HodgeConjecture/Lemmas/AlgebraicGeometry/ComplexStandardEtale.lean#L90}{\texttt{HodgeConjecture/Lemmas/AlgebraicGeometry/ComplexStandardEtale.lean:90}}.
\end{definition}

\begin{definition}[AlgebraicGeometry.ComplexAlgHom.standardEtalePointToAlgHom]\label{def:AlgebraicGeometry.ComplexAlgHom.standardEtalePointToAlgHom}
\lean{AlgebraicGeometry.ComplexAlgHom.standardEtalePointToAlgHom}\leanok
\uses{def:AlgebraicGeometry.ComplexAlgHom.standardEtalePointSpace}
Construct a complex algebra homomorphism from the root data of a standard étale point.

\noindent\emph{Lean:} \texttt{def AlgebraicGeometry.ComplexAlgHom.standardEtalePointToAlgHom}, \href{https://github.com/Paul-Lez/HodgeConjecture/blob/main/HodgeConjecture/Lemmas/AlgebraicGeometry/ComplexStandardEtale.lean#L95}{\texttt{HodgeConjecture/Lemmas/AlgebraicGeometry/ComplexStandardEtale.lean:95}}.
\end{definition}

\begin{definition}[AlgebraicGeometry.ComplexAlgHom.algHomToStandardEtalePoint]\label{def:AlgebraicGeometry.ComplexAlgHom.algHomToStandardEtalePoint}
\lean{AlgebraicGeometry.ComplexAlgHom.algHomToStandardEtalePoint}\leanok
\uses{def:AlgebraicGeometry.ComplexAlgHom.standardEtalePointSpace}
Read the base point and distinguished root from a complex point of a standard étale algebra.

\noindent\emph{Lean:} \texttt{def AlgebraicGeometry.ComplexAlgHom.algHomToStandardEtalePoint}, \href{https://github.com/Paul-Lez/HodgeConjecture/blob/main/HodgeConjecture/Lemmas/AlgebraicGeometry/ComplexStandardEtale.lean#L107}{\texttt{HodgeConjecture/Lemmas/AlgebraicGeometry/ComplexStandardEtale.lean:107}}.
\end{definition}

\begin{definition}[AlgebraicGeometry.ComplexAlgHom.standardEtalePointEquiv]\label{def:AlgebraicGeometry.ComplexAlgHom.standardEtalePointEquiv}
\lean{AlgebraicGeometry.ComplexAlgHom.standardEtalePointEquiv}\leanok
\uses{def:AlgebraicGeometry.ComplexAlgHom.algHomToStandardEtalePoint, def:AlgebraicGeometry.ComplexAlgHom.standardEtalePointToAlgHom}
Complex algebra homomorphisms out of a standard étale algebra are equivalent to its root
data.

\noindent\emph{Lean:} \texttt{def AlgebraicGeometry.ComplexAlgHom.standardEtalePointEquiv}, \href{https://github.com/Paul-Lez/HodgeConjecture/blob/main/HodgeConjecture/Lemmas/AlgebraicGeometry/ComplexStandardEtale.lean#L124}{\texttt{HodgeConjecture/Lemmas/AlgebraicGeometry/ComplexStandardEtale.lean:124}}.
\end{definition}

\begin{definition}[AlgebraicGeometry.ComplexAlgHom.standardEtalePointHomeomorph]\label{def:AlgebraicGeometry.ComplexAlgHom.standardEtalePointHomeomorph}
\lean{AlgebraicGeometry.ComplexAlgHom.standardEtalePointHomeomorph}\leanok
\uses{def:AlgebraicGeometry.ComplexAlgHom.standardEtalePointEquiv}
The pointwise topology on homomorphisms out of a standard étale algebra is the topology on
its explicit root data.

\noindent\emph{Lean:} \texttt{def AlgebraicGeometry.ComplexAlgHom.standardEtalePointHomeomorph}, \href{https://github.com/Paul-Lez/HodgeConjecture/blob/main/HodgeConjecture/Lemmas/AlgebraicGeometry/ComplexStandardEtale.lean#L279}{\texttt{HodgeConjecture/Lemmas/AlgebraicGeometry/ComplexStandardEtale.lean:279}}.
\end{definition}

\begin{definition}[AlgebraicGeometry.ComplexAlgHom.standardEtaleCoordinateSpace]\label{def:AlgebraicGeometry.ComplexAlgHom.standardEtaleCoordinateSpace}
\lean{AlgebraicGeometry.ComplexAlgHom.standardEtaleCoordinateSpace}\leanok
\uses{def:AlgebraicGeometry.ComplexAlgHom.complexPolynomialRing}
The standard étale equations written in ordinary complex coordinates.

\noindent\emph{Lean:} \texttt{abbrev AlgebraicGeometry.ComplexAlgHom.standardEtaleCoordinateSpace}, \href{https://github.com/Paul-Lez/HodgeConjecture/blob/main/HodgeConjecture/Lemmas/AlgebraicGeometry/ComplexStandardEtale.lean#L287}{\texttt{HodgeConjecture/Lemmas/AlgebraicGeometry/ComplexStandardEtale.lean:287}}.
\end{definition}

\begin{definition}[AlgebraicGeometry.ComplexAlgHom.standardEtalePointCoordinateHomeomorph]\label{def:AlgebraicGeometry.ComplexAlgHom.standardEtalePointCoordinateHomeomorph}
\lean{AlgebraicGeometry.ComplexAlgHom.standardEtalePointCoordinateHomeomorph}\leanok
\uses{def:AlgebraicGeometry.ComplexAlgHom.mvPolynomialAlgHomHomeomorph, def:AlgebraicGeometry.ComplexAlgHom.standardEtaleCoordinateSpace, def:AlgebraicGeometry.ComplexAlgHom.standardEtalePointSpace}
Replace the polynomial-ring homomorphism in standard étale root data by its tuple of variable
values.

\noindent\emph{Lean:} \texttt{def AlgebraicGeometry.ComplexAlgHom.standardEtalePointCoordinateHomeomorph}, \href{https://github.com/Paul-Lez/HodgeConjecture/blob/main/HodgeConjecture/Lemmas/AlgebraicGeometry/ComplexStandardEtale.lean#L293}{\texttt{HodgeConjecture/Lemmas/AlgebraicGeometry/ComplexStandardEtale.lean:293}}.
\end{definition}

\begin{definition}[AlgebraicGeometry.ComplexAlgHom.standardEtaleCoordinateHomeomorph]\label{def:AlgebraicGeometry.ComplexAlgHom.standardEtaleCoordinateHomeomorph}
\lean{AlgebraicGeometry.ComplexAlgHom.standardEtaleCoordinateHomeomorph}\leanok
\uses{def:AlgebraicGeometry.ComplexAlgHom.standardEtalePointCoordinateHomeomorph, def:AlgebraicGeometry.ComplexAlgHom.standardEtalePointHomeomorph}
Complex points of a standard étale algebra as the corresponding equation locus in ordinary
complex coordinates.

\noindent\emph{Lean:} \texttt{def AlgebraicGeometry.ComplexAlgHom.standardEtaleCoordinateHomeomorph}, \href{https://github.com/Paul-Lez/HodgeConjecture/blob/main/HodgeConjecture/Lemmas/AlgebraicGeometry/ComplexStandardEtale.lean#L303}{\texttt{HodgeConjecture/Lemmas/AlgebraicGeometry/ComplexStandardEtale.lean:303}}.
\end{definition}

\section{Complex points of étale morphisms}\label{sec:HodgeConjecture.Lemmas.AlgebraicGeometry.ComplexEtale}
\noindent\emph{File:} \href{https://github.com/Paul-Lez/HodgeConjecture/blob/main/HodgeConjecture/Lemmas/AlgebraicGeometry/ComplexEtale.lean}{\texttt{HodgeConjecture/Lemmas/AlgebraicGeometry/ComplexEtale.lean}}.

The complex points of a standard étale algebra over a polynomial ring form the zero locus of
one polynomial with nonzero derivative in its distinguished variable. The complex implicit
function theorem gives explicit local charts in which projection to the polynomial coordinates
is a homeomorphism. This file constructs those charts and transports the result to the associated
affine schemes.

The charts themselves are statements about \texttt{ℂ}-algebra homomorphisms and live in namespace
\texttt{AlgebraicGeometry.ComplexAlgHom}; only the last section, which transports them to schemes, is in
\texttt{AlgebraicGeometry.ComplexPoint}.

\begin{definition}[AlgebraicGeometry.ComplexAlgHom.standardEtaleEquation]\label{def:AlgebraicGeometry.ComplexAlgHom.standardEtaleEquation}
\lean{AlgebraicGeometry.ComplexAlgHom.standardEtaleEquation}\leanok
\uses{def:AlgebraicGeometry.ComplexAlgHom.complexPolynomialRing}
The equation defining a standard étale algebra, evaluated in ordinary complex
coordinates.

\noindent\emph{Lean:} \texttt{def AlgebraicGeometry.ComplexAlgHom.standardEtaleEquation}, \href{https://github.com/Paul-Lez/HodgeConjecture/blob/main/HodgeConjecture/Lemmas/AlgebraicGeometry/ComplexEtale.lean#L58}{\texttt{HodgeConjecture/Lemmas/AlgebraicGeometry/ComplexEtale.lean:58}}.
\end{definition}

\begin{definition}[AlgebraicGeometry.ComplexAlgHom.standardEtaleFiberPolynomial]\label{def:AlgebraicGeometry.ComplexAlgHom.standardEtaleFiberPolynomial}
\lean{AlgebraicGeometry.ComplexAlgHom.standardEtaleFiberPolynomial}\leanok
\uses{def:AlgebraicGeometry.ComplexAlgHom.complexPolynomialRing}
The defining polynomial after specializing the polynomial-ring coordinates.

\noindent\emph{Lean:} \texttt{def AlgebraicGeometry.ComplexAlgHom.standardEtaleFiberPolynomial}, \href{https://github.com/Paul-Lez/HodgeConjecture/blob/main/HodgeConjecture/Lemmas/AlgebraicGeometry/ComplexEtale.lean#L101}{\texttt{HodgeConjecture/Lemmas/AlgebraicGeometry/ComplexEtale.lean:101}}.
\end{definition}

\begin{definition}[AlgebraicGeometry.ComplexAlgHom.standardEtaleImplicitOpenPartialHomeomorph]\label{def:AlgebraicGeometry.ComplexAlgHom.standardEtaleImplicitOpenPartialHomeomorph}
\lean{AlgebraicGeometry.ComplexAlgHom.standardEtaleImplicitOpenPartialHomeomorph}\leanok
\uses{def:AlgebraicGeometry.ComplexAlgHom.standardEtaleCoordinateSpace, def:AlgebraicGeometry.ComplexAlgHom.standardEtaleEquation, def:AlgebraicGeometry.ComplexAlgHom.standardEtaleFiberPolynomial}
The ambient implicit-function chart sending \texttt{(z, t)} to \texttt{(f(z, t), z)}.

\noindent\emph{Lean:} \texttt{def AlgebraicGeometry.ComplexAlgHom.standardEtaleImplicitOpenPartialHomeomorph}, \href{https://github.com/Paul-Lez/HodgeConjecture/blob/main/HodgeConjecture/Lemmas/AlgebraicGeometry/ComplexEtale.lean#L160}{\texttt{HodgeConjecture/Lemmas/AlgebraicGeometry/ComplexEtale.lean:160}}.
\end{definition}

\begin{definition}[AlgebraicGeometry.ComplexAlgHom.standardEtaleLocalizationEquation]\label{def:AlgebraicGeometry.ComplexAlgHom.standardEtaleLocalizationEquation}
\lean{AlgebraicGeometry.ComplexAlgHom.standardEtaleLocalizationEquation}\leanok
\uses{def:AlgebraicGeometry.ComplexAlgHom.complexPolynomialRing}
The localization polynomial evaluated in ordinary complex coordinates.

\noindent\emph{Lean:} \texttt{def AlgebraicGeometry.ComplexAlgHom.standardEtaleLocalizationEquation}, \href{https://github.com/Paul-Lez/HodgeConjecture/blob/main/HodgeConjecture/Lemmas/AlgebraicGeometry/ComplexEtale.lean#L199}{\texttt{HodgeConjecture/Lemmas/AlgebraicGeometry/ComplexEtale.lean:199}}.
\end{definition}

\begin{definition}[AlgebraicGeometry.ComplexAlgHom.standardEtaleBivariateEvaluation]\label{def:AlgebraicGeometry.ComplexAlgHom.standardEtaleBivariateEvaluation}
\lean{AlgebraicGeometry.ComplexAlgHom.standardEtaleBivariateEvaluation}\leanok
\uses{def:AlgebraicGeometry.ComplexAlgHom.standardEtaleLocalizationEquation}
Evaluate a representative of the double-polynomial presentation of a standard étale
algebra in ordinary complex coordinates. The outer variable is sent to the inverse of the
localization polynomial.

\noindent\emph{Lean:} \texttt{def AlgebraicGeometry.ComplexAlgHom.standardEtaleBivariateEvaluation}, \href{https://github.com/Paul-Lez/HodgeConjecture/blob/main/HodgeConjecture/Lemmas/AlgebraicGeometry/ComplexEtale.lean#L207}{\texttt{HodgeConjecture/Lemmas/AlgebraicGeometry/ComplexEtale.lean:207}}.
\end{definition}

\begin{definition}[AlgebraicGeometry.ComplexAlgHom.standardEtaleChartSource]\label{def:AlgebraicGeometry.ComplexAlgHom.standardEtaleChartSource}
\lean{AlgebraicGeometry.ComplexAlgHom.standardEtaleChartSource}\leanok
\uses{def:AlgebraicGeometry.ComplexAlgHom.standardEtaleImplicitOpenPartialHomeomorph}
The source of the standard étale chart centered at \texttt{z}.

\noindent\emph{Lean:} \texttt{def AlgebraicGeometry.ComplexAlgHom.standardEtaleChartSource}, \href{https://github.com/Paul-Lez/HodgeConjecture/blob/main/HodgeConjecture/Lemmas/AlgebraicGeometry/ComplexEtale.lean#L248}{\texttt{HodgeConjecture/Lemmas/AlgebraicGeometry/ComplexEtale.lean:248}}.
\end{definition}

\begin{definition}[AlgebraicGeometry.ComplexAlgHom.standardEtaleChartTarget]\label{def:AlgebraicGeometry.ComplexAlgHom.standardEtaleChartTarget}
\lean{AlgebraicGeometry.ComplexAlgHom.standardEtaleChartTarget}\leanok
\uses{def:AlgebraicGeometry.ComplexAlgHom.standardEtaleImplicitOpenPartialHomeomorph, def:AlgebraicGeometry.ComplexAlgHom.standardEtaleLocalizationEquation}
The target of the standard étale chart centered at \texttt{z}. The second condition retains the
localization defining the standard étale algebra.

\noindent\emph{Lean:} \texttt{def AlgebraicGeometry.ComplexAlgHom.standardEtaleChartTarget}, \href{https://github.com/Paul-Lez/HodgeConjecture/blob/main/HodgeConjecture/Lemmas/AlgebraicGeometry/ComplexEtale.lean#L253}{\texttt{HodgeConjecture/Lemmas/AlgebraicGeometry/ComplexEtale.lean:253}}.
\end{definition}

\begin{definition}[AlgebraicGeometry.ComplexAlgHom.standardEtaleChartInverse]\label{def:AlgebraicGeometry.ComplexAlgHom.standardEtaleChartInverse}
\lean{AlgebraicGeometry.ComplexAlgHom.standardEtaleChartInverse}\leanok
\uses{def:AlgebraicGeometry.ComplexAlgHom.standardEtaleChartTarget}
The inverse of the standard étale chart on its target. It is assigned the center point
outside the target because \texttt{OpenPartialHomeomorph} stores total functions.

\noindent\emph{Lean:} \texttt{def AlgebraicGeometry.ComplexAlgHom.standardEtaleChartInverse}, \href{https://github.com/Paul-Lez/HodgeConjecture/blob/main/HodgeConjecture/Lemmas/AlgebraicGeometry/ComplexEtale.lean#L283}{\texttt{HodgeConjecture/Lemmas/AlgebraicGeometry/ComplexEtale.lean:283}}.
\end{definition}

\begin{definition}[AlgebraicGeometry.ComplexAlgHom.standardEtaleProjectionChart]\label{def:AlgebraicGeometry.ComplexAlgHom.standardEtaleProjectionChart}
\lean{AlgebraicGeometry.ComplexAlgHom.standardEtaleProjectionChart}\leanok
\uses{def:AlgebraicGeometry.ComplexAlgHom.mvPolynomialAlgHomHomeomorph, def:AlgebraicGeometry.ComplexAlgHom.standardEtaleChartInverse, def:AlgebraicGeometry.ComplexAlgHom.standardEtaleChartSource}
A neighborhood of a standard étale point on which projection to the polynomial-ring
coordinates is a homeomorphism onto an open set.

\noindent\emph{Lean:} \texttt{def AlgebraicGeometry.ComplexAlgHom.standardEtaleProjectionChart}, \href{https://github.com/Paul-Lez/HodgeConjecture/blob/main/HodgeConjecture/Lemmas/AlgebraicGeometry/ComplexEtale.lean#L476}{\texttt{HodgeConjecture/Lemmas/AlgebraicGeometry/ComplexEtale.lean:476}}.
\end{definition}

\begin{definition}[AlgebraicGeometry.ComplexAlgHom.standardEtaleAlgHomProjectionChart]\label{def:AlgebraicGeometry.ComplexAlgHom.standardEtaleAlgHomProjectionChart}
\lean{AlgebraicGeometry.ComplexAlgHom.standardEtaleAlgHomProjectionChart}\leanok
\uses{def:AlgebraicGeometry.ComplexAlgHom.standardEtaleCoordinateHomeomorph, def:AlgebraicGeometry.ComplexAlgHom.standardEtaleProjectionChart}
A projection chart on the complex algebra homomorphisms of a standard étale algebra.

\noindent\emph{Lean:} \texttt{def AlgebraicGeometry.ComplexAlgHom.standardEtaleAlgHomProjectionChart}, \href{https://github.com/Paul-Lez/HodgeConjecture/blob/main/HodgeConjecture/Lemmas/AlgebraicGeometry/ComplexEtale.lean#L503}{\texttt{HodgeConjecture/Lemmas/AlgebraicGeometry/ComplexEtale.lean:503}}.
\end{definition}

\begin{definition}[AlgebraicGeometry.ComplexAlgHom.chosenStandardEtalePresentation]\label{def:AlgebraicGeometry.ComplexAlgHom.chosenStandardEtalePresentation}
\lean{AlgebraicGeometry.ComplexAlgHom.chosenStandardEtalePresentation}\leanok
\uses{def:AlgebraicGeometry.ComplexAlgHom.complexPolynomialRing}
A chosen standard étale presentation of a standard étale algebra.

\noindent\emph{Lean:} \texttt{def AlgebraicGeometry.ComplexAlgHom.chosenStandardEtalePresentation}, \href{https://github.com/Paul-Lez/HodgeConjecture/blob/main/HodgeConjecture/Lemmas/AlgebraicGeometry/ComplexEtale.lean#L568}{\texttt{HodgeConjecture/Lemmas/AlgebraicGeometry/ComplexEtale.lean:568}}.
\end{definition}

\begin{definition}[AlgebraicGeometry.ComplexAlgHom.standardEtalePresentationComplexAlgEquiv]\label{def:AlgebraicGeometry.ComplexAlgHom.standardEtalePresentationComplexAlgEquiv}
\lean{AlgebraicGeometry.ComplexAlgHom.standardEtalePresentationComplexAlgEquiv}\leanok
\uses{def:AlgebraicGeometry.ComplexAlgHom.chosenStandardEtalePresentation}
The chosen presentation is also an equivalence of complex algebras.

\noindent\emph{Lean:} \texttt{def AlgebraicGeometry.ComplexAlgHom.standardEtalePresentationComplexAlgEquiv}, \href{https://github.com/Paul-Lez/HodgeConjecture/blob/main/HodgeConjecture/Lemmas/AlgebraicGeometry/ComplexEtale.lean#L573}{\texttt{HodgeConjecture/Lemmas/AlgebraicGeometry/ComplexEtale.lean:573}}.
\end{definition}

\begin{definition}[AlgebraicGeometry.ComplexAlgHom.isStandardEtaleBaseAlgHom]\label{def:AlgebraicGeometry.ComplexAlgHom.isStandardEtaleBaseAlgHom}
\lean{AlgebraicGeometry.ComplexAlgHom.isStandardEtaleBaseAlgHom}\leanok
\uses{def:AlgebraicGeometry.ComplexAlgHom.complexPolynomialRing}
Restriction of a complex point of a standard étale algebra to its polynomial base.

\noindent\emph{Lean:} \texttt{def AlgebraicGeometry.ComplexAlgHom.isStandardEtaleBaseAlgHom}, \href{https://github.com/Paul-Lez/HodgeConjecture/blob/main/HodgeConjecture/Lemmas/AlgebraicGeometry/ComplexEtale.lean#L578}{\texttt{HodgeConjecture/Lemmas/AlgebraicGeometry/ComplexEtale.lean:578}}.
\end{definition}

\begin{definition}[AlgebraicGeometry.ComplexAlgHom.isStandardEtaleAlgHomProjectionChart]\label{def:AlgebraicGeometry.ComplexAlgHom.isStandardEtaleAlgHomProjectionChart}
\lean{AlgebraicGeometry.ComplexAlgHom.isStandardEtaleAlgHomProjectionChart}\leanok
\uses{def:AlgebraicGeometry.ComplexAlgHom.precompAlgEquivHomeomorph, def:AlgebraicGeometry.ComplexAlgHom.standardEtaleAlgHomProjectionChart, def:AlgebraicGeometry.ComplexAlgHom.standardEtalePresentationComplexAlgEquiv}
A projection chart for an arbitrary standard étale algebra, transported through its chosen
standard étale presentation.

\noindent\emph{Lean:} \texttt{def AlgebraicGeometry.ComplexAlgHom.isStandardEtaleAlgHomProjectionChart}, \href{https://github.com/Paul-Lez/HodgeConjecture/blob/main/HodgeConjecture/Lemmas/AlgebraicGeometry/ComplexEtale.lean#L583}{\texttt{HodgeConjecture/Lemmas/AlgebraicGeometry/ComplexEtale.lean:583}}.
\end{definition}

\begin{definition}[AlgebraicGeometry.ComplexAlgHom.etaleBaseAlgHom]\label{def:AlgebraicGeometry.ComplexAlgHom.etaleBaseAlgHom}
\lean{AlgebraicGeometry.ComplexAlgHom.etaleBaseAlgHom}\leanok
\uses{def:AlgebraicGeometry.ComplexAlgHom.complexPolynomialRing}
Restriction of a complex point of an algebra over a polynomial ring to the polynomial
base.

\noindent\emph{Lean:} \texttt{def AlgebraicGeometry.ComplexAlgHom.etaleBaseAlgHom}, \href{https://github.com/Paul-Lez/HodgeConjecture/blob/main/HodgeConjecture/Lemmas/AlgebraicGeometry/ComplexEtale.lean#L665}{\texttt{HodgeConjecture/Lemmas/AlgebraicGeometry/ComplexEtale.lean:665}}.
\end{definition}

\begin{definition}[AlgebraicGeometry.ComplexAlgHom.EtaleStandardNeighborhood]\label{def:AlgebraicGeometry.ComplexAlgHom.EtaleStandardNeighborhood}
\lean{AlgebraicGeometry.ComplexAlgHom.EtaleStandardNeighborhood}\leanok
\uses{def:AlgebraicGeometry.ComplexAlgHom.complexPolynomialRing}
Standard étale localization data at a complex point of an étale algebra.

\noindent\emph{Lean:} \texttt{structure AlgebraicGeometry.ComplexAlgHom.EtaleStandardNeighborhood}, \href{https://github.com/Paul-Lez/HodgeConjecture/blob/main/HodgeConjecture/Lemmas/AlgebraicGeometry/ComplexEtale.lean#L708}{\texttt{HodgeConjecture/Lemmas/AlgebraicGeometry/ComplexEtale.lean:708}}.
\end{definition}

\begin{definition}[AlgebraicGeometry.ComplexAlgHom.etaleStandardNeighborhood]\label{def:AlgebraicGeometry.ComplexAlgHom.etaleStandardNeighborhood}
\lean{AlgebraicGeometry.ComplexAlgHom.etaleStandardNeighborhood}\leanok
\uses{def:AlgebraicGeometry.ComplexAlgHom.EtaleStandardNeighborhood}
A chosen standard étale localization at a complex point of an étale algebra.

\noindent\emph{Lean:} \texttt{def AlgebraicGeometry.ComplexAlgHom.etaleStandardNeighborhood}, \href{https://github.com/Paul-Lez/HodgeConjecture/blob/main/HodgeConjecture/Lemmas/AlgebraicGeometry/ComplexEtale.lean#L732}{\texttt{HodgeConjecture/Lemmas/AlgebraicGeometry/ComplexEtale.lean:732}}.
\end{definition}

\begin{definition}[AlgebraicGeometry.ComplexAlgHom.pointInEtaleStandardNeighborhood]\label{def:AlgebraicGeometry.ComplexAlgHom.pointInEtaleStandardNeighborhood}
\lean{AlgebraicGeometry.ComplexAlgHom.pointInEtaleStandardNeighborhood}\leanok
\uses{def:AlgebraicGeometry.ComplexAlgHom.etaleStandardNeighborhood, def:AlgebraicGeometry.ComplexAlgHom.localizationAwayAlgHomHomeomorph}
The point of the chosen standard étale localization lying over \texttt{u}.

\noindent\emph{Lean:} \texttt{def AlgebraicGeometry.ComplexAlgHom.pointInEtaleStandardNeighborhood}, \href{https://github.com/Paul-Lez/HodgeConjecture/blob/main/HodgeConjecture/Lemmas/AlgebraicGeometry/ComplexEtale.lean#L737}{\texttt{HodgeConjecture/Lemmas/AlgebraicGeometry/ComplexEtale.lean:737}}.
\end{definition}

\begin{definition}[AlgebraicGeometry.ComplexAlgHom.etaleAlgHomProjectionChart]\label{def:AlgebraicGeometry.ComplexAlgHom.etaleAlgHomProjectionChart}
\lean{AlgebraicGeometry.ComplexAlgHom.etaleAlgHomProjectionChart}\leanok
\uses{def:AlgebraicGeometry.ComplexAlgHom.isStandardEtaleAlgHomProjectionChart, def:AlgebraicGeometry.ComplexAlgHom.localizationAwayAlgHomMap, def:AlgebraicGeometry.ComplexAlgHom.pointInEtaleStandardNeighborhood}
A projection chart on the complex points of an arbitrary étale algebra. It is obtained from a
standard étale chart after restricting to a principal open neighborhood.

\noindent\emph{Lean:} \texttt{def AlgebraicGeometry.ComplexAlgHom.etaleAlgHomProjectionChart}, \href{https://github.com/Paul-Lez/HodgeConjecture/blob/main/HodgeConjecture/Lemmas/AlgebraicGeometry/ComplexEtale.lean#L744}{\texttt{HodgeConjecture/Lemmas/AlgebraicGeometry/ComplexEtale.lean:744}}.
\end{definition}

\begin{definition}[AlgebraicGeometry.ComplexPoint.polynomialSpecToAffineSpacePointMap]\label{def:AlgebraicGeometry.ComplexPoint.polynomialSpecToAffineSpacePointMap}
\lean{AlgebraicGeometry.ComplexPoint.polynomialSpecToAffineSpacePointMap}\leanok
\uses{def:AlgebraicGeometry.ComplexAlgHom.complexPolynomialRing, def:AlgebraicGeometry.ComplexPoint, def:AlgebraicGeometry.ComplexPoint.affineSpecStructureMap, def:AlgebraicGeometry.ComplexPoint.complexAffineSpace, def:AlgebraicGeometry.Point.map}
The complex-point map from the spectrum of a polynomial ring to algebraic affine space.

\noindent\emph{Lean:} \texttt{def AlgebraicGeometry.ComplexPoint.polynomialSpecToAffineSpacePointMap}, \href{https://github.com/Paul-Lez/HodgeConjecture/blob/main/HodgeConjecture/Lemmas/AlgebraicGeometry/ComplexEtale.lean#L818}{\texttt{HodgeConjecture/Lemmas/AlgebraicGeometry/ComplexEtale.lean:818}}.
\end{definition}

\section{Complex points of open subschemes}\label{sec:HodgeConjecture.Lemmas.AlgebraicGeometry.ComplexOpen}
\noindent\emph{File:} \href{https://github.com/Paul-Lez/HodgeConjecture/blob/main/HodgeConjecture/Lemmas/AlgebraicGeometry/ComplexOpen.lean}{\texttt{HodgeConjecture/Lemmas/AlgebraicGeometry/ComplexOpen.lean}}.

Analytification respects restriction to an open subscheme. More precisely, complex points of an
open subscheme are canonically homeomorphic to the corresponding open subspace of the complex
points of the ambient scheme. In particular, an open immersion of this form induces an open
embedding on complex points.

\begin{definition}[AlgebraicGeometry.ComplexPoint.openScheme]\label{def:AlgebraicGeometry.ComplexPoint.openScheme}
\lean{AlgebraicGeometry.ComplexPoint.openScheme}\leanok
An open subscheme with its induced complex structure.

\noindent\emph{Lean:} \texttt{abbrev AlgebraicGeometry.ComplexPoint.openScheme}, \href{https://github.com/Paul-Lez/HodgeConjecture/blob/main/HodgeConjecture/Lemmas/AlgebraicGeometry/ComplexOpen.lean#L45}{\texttt{HodgeConjecture/Lemmas/AlgebraicGeometry/ComplexOpen.lean:45}}.
\end{definition}

\begin{definition}[AlgebraicGeometry.ComplexPoint.openInclusion]\label{def:AlgebraicGeometry.ComplexPoint.openInclusion}
\lean{AlgebraicGeometry.ComplexPoint.openInclusion}\leanok
\uses{def:AlgebraicGeometry.ComplexPoint.openScheme}
The inclusion of an open subscheme as a morphism over \texttt{Spec ℂ}.

\noindent\emph{Lean:} \texttt{abbrev AlgebraicGeometry.ComplexPoint.openInclusion}, \href{https://github.com/Paul-Lez/HodgeConjecture/blob/main/HodgeConjecture/Lemmas/AlgebraicGeometry/ComplexOpen.lean#L49}{\texttt{HodgeConjecture/Lemmas/AlgebraicGeometry/ComplexOpen.lean:49}}.
\end{definition}

\begin{definition}[AlgebraicGeometry.ComplexPoint.liftToOpen]\label{def:AlgebraicGeometry.ComplexPoint.liftToOpen}
\lean{AlgebraicGeometry.ComplexPoint.liftToOpen}\leanok
\uses{def:AlgebraicGeometry.ComplexPoint, def:AlgebraicGeometry.Point.overOpen}
A complex point in an open set factors through the corresponding open subscheme.

\noindent\emph{Lean:} \texttt{def AlgebraicGeometry.ComplexPoint.liftToOpen}, \href{https://github.com/Paul-Lez/HodgeConjecture/blob/main/HodgeConjecture/Lemmas/AlgebraicGeometry/ComplexOpen.lean#L63}{\texttt{HodgeConjecture/Lemmas/AlgebraicGeometry/ComplexOpen.lean:63}}.
\end{definition}

\begin{definition}[AlgebraicGeometry.ComplexPoint.asOpenPoint]\label{def:AlgebraicGeometry.ComplexPoint.asOpenPoint}
\lean{AlgebraicGeometry.ComplexPoint.asOpenPoint}\leanok
\uses{def:AlgebraicGeometry.ComplexPoint.liftToOpen, def:AlgebraicGeometry.ComplexPoint.openScheme}
A point of \texttt{X} lying in \texttt{U}, regarded as a complex point of the open subscheme \texttt{U}.

\noindent\emph{Lean:} \texttt{def AlgebraicGeometry.ComplexPoint.asOpenPoint}, \href{https://github.com/Paul-Lez/HodgeConjecture/blob/main/HodgeConjecture/Lemmas/AlgebraicGeometry/ComplexOpen.lean#L73}{\texttt{HodgeConjecture/Lemmas/AlgebraicGeometry/ComplexOpen.lean:73}}.
\end{definition}

\begin{definition}[AlgebraicGeometry.ComplexPoint.openEquiv]\label{def:AlgebraicGeometry.ComplexPoint.openEquiv}
\lean{AlgebraicGeometry.ComplexPoint.openEquiv}\leanok
\uses{def:AlgebraicGeometry.ComplexPoint.asOpenPoint, def:AlgebraicGeometry.ComplexPoint.openInclusion, def:AlgebraicGeometry.Point.map}
Complex points of an open subscheme are the ambient complex points lying in the open.

\noindent\emph{Lean:} \texttt{def AlgebraicGeometry.ComplexPoint.openEquiv}, \href{https://github.com/Paul-Lez/HodgeConjecture/blob/main/HodgeConjecture/Lemmas/AlgebraicGeometry/ComplexOpen.lean#L81}{\texttt{HodgeConjecture/Lemmas/AlgebraicGeometry/ComplexOpen.lean:81}}.
\end{definition}

\begin{definition}[AlgebraicGeometry.ComplexPoint.openHomeomorph]\label{def:AlgebraicGeometry.ComplexPoint.openHomeomorph}
\lean{AlgebraicGeometry.ComplexPoint.openHomeomorph}\leanok
\uses{def:AlgebraicGeometry.ComplexPoint.openEquiv, def:AlgebraicGeometry.Point.analyticSubbasis, def:AlgebraicGeometry.Point.chartTopology}
Analytification of an open subscheme is the corresponding analytic open subspace.

\noindent\emph{Lean:} \texttt{def AlgebraicGeometry.ComplexPoint.openHomeomorph}, \href{https://github.com/Paul-Lez/HodgeConjecture/blob/main/HodgeConjecture/Lemmas/AlgebraicGeometry/ComplexOpen.lean#L177}{\texttt{HodgeConjecture/Lemmas/AlgebraicGeometry/ComplexOpen.lean:177}}.
\end{definition}

\section{Local coordinates on smooth complex schemes}\label{sec:HodgeConjecture.Lemmas.AlgebraicGeometry.SmoothComplexCoordinates}
\noindent\emph{File:} \href{https://github.com/Paul-Lez/HodgeConjecture/blob/main/HodgeConjecture/Lemmas/AlgebraicGeometry/SmoothComplexCoordinates.lean}{\texttt{HodgeConjecture/Lemmas/AlgebraicGeometry/SmoothComplexCoordinates.lean}}.

Every point of a smooth scheme over \texttt{ℂ} has an affine neighborhood with an étale map to a
finite-dimensional affine space. This file constructs such coordinates from algebraic smoothness,
proves that the map is over \texttt{Spec ℂ}, and records the resulting continuous complex-valued
coordinate tuple on the corresponding analytic open set.

\begin{definition}[AlgebraicGeometry.LocalEtaleCoordinates]\label{def:AlgebraicGeometry.LocalEtaleCoordinates}
\lean{AlgebraicGeometry.LocalEtaleCoordinates}\leanok
\uses{def:AlgebraicGeometry.overSpecAlgebra}
Étale algebraic coordinates of the specified relative dimension around a point of a smooth
complex scheme.

\noindent\emph{Lean:} \texttt{structure AlgebraicGeometry.LocalEtaleCoordinates}, \href{https://github.com/Paul-Lez/HodgeConjecture/blob/main/HodgeConjecture/Lemmas/AlgebraicGeometry/SmoothComplexCoordinates.lean#L108}{\texttt{HodgeConjecture/Lemmas/AlgebraicGeometry/SmoothComplexCoordinates.lean:108}}.
\end{definition}

\begin{definition}[AlgebraicGeometry.localEtaleCoordinates]\label{def:AlgebraicGeometry.localEtaleCoordinates}
\lean{AlgebraicGeometry.localEtaleCoordinates}\leanok
\uses{def:AlgebraicGeometry.LocalEtaleCoordinates}
A choice of the \texttt{d} étale algebraic coordinates constructed from relative-dimensional
smoothness.

\noindent\emph{Lean:} \texttt{def AlgebraicGeometry.localEtaleCoordinates}, \href{https://github.com/Paul-Lez/HodgeConjecture/blob/main/HodgeConjecture/Lemmas/AlgebraicGeometry/SmoothComplexCoordinates.lean#L135}{\texttt{HodgeConjecture/Lemmas/AlgebraicGeometry/SmoothComplexCoordinates.lean:135}}.
\end{definition}

\begin{definition}[AlgebraicGeometry.LocalEtaleCoordinates.coordinateRingHomOnOpen]\label{def:AlgebraicGeometry.LocalEtaleCoordinates.coordinateRingHomOnOpen}
\lean{AlgebraicGeometry.LocalEtaleCoordinates.coordinateRingHomOnOpen}\leanok
\uses{def:AlgebraicGeometry.LocalEtaleCoordinates}
The coordinate ring map after identifying sections on an open with its global sections.

\noindent\emph{Lean:} \texttt{def AlgebraicGeometry.LocalEtaleCoordinates.coordinateRingHomOnOpen}, \href{https://github.com/Paul-Lez/HodgeConjecture/blob/main/HodgeConjecture/Lemmas/AlgebraicGeometry/SmoothComplexCoordinates.lean#L146}{\texttt{HodgeConjecture/Lemmas/AlgebraicGeometry/SmoothComplexCoordinates.lean:146}}.
\end{definition}

\begin{definition}[AlgebraicGeometry.LocalEtaleCoordinates.toAffineSpace]\label{def:AlgebraicGeometry.LocalEtaleCoordinates.toAffineSpace}
\lean{AlgebraicGeometry.LocalEtaleCoordinates.toAffineSpace}\leanok
\uses{def:AlgebraicGeometry.ComplexPoint.complexAffineSpace, def:AlgebraicGeometry.LocalEtaleCoordinates.coordinateRingHomOnOpen}
The scheme morphism defined by the étale coordinates.

\noindent\emph{Lean:} \texttt{def AlgebraicGeometry.LocalEtaleCoordinates.toAffineSpace}, \href{https://github.com/Paul-Lez/HodgeConjecture/blob/main/HodgeConjecture/Lemmas/AlgebraicGeometry/SmoothComplexCoordinates.lean#L171}{\texttt{HodgeConjecture/Lemmas/AlgebraicGeometry/SmoothComplexCoordinates.lean:171}}.
\end{definition}

\begin{definition}[AlgebraicGeometry.LocalEtaleCoordinates.toAffineSpaceOver]\label{def:AlgebraicGeometry.LocalEtaleCoordinates.toAffineSpaceOver}
\lean{AlgebraicGeometry.LocalEtaleCoordinates.toAffineSpaceOver}\leanok
\uses{def:AlgebraicGeometry.ComplexPoint.openScheme, def:AlgebraicGeometry.LocalEtaleCoordinates.toAffineSpace}
The étale coordinate morphism bundled over the complex base.

\noindent\emph{Lean:} \texttt{abbrev AlgebraicGeometry.LocalEtaleCoordinates.toAffineSpaceOver}, \href{https://github.com/Paul-Lez/HodgeConjecture/blob/main/HodgeConjecture/Lemmas/AlgebraicGeometry/SmoothComplexCoordinates.lean#L206}{\texttt{HodgeConjecture/Lemmas/AlgebraicGeometry/SmoothComplexCoordinates.lean:206}}.
\end{definition}

\begin{definition}[AlgebraicGeometry.LocalEtaleCoordinates.pointMap]\label{def:AlgebraicGeometry.LocalEtaleCoordinates.pointMap}
\lean{AlgebraicGeometry.LocalEtaleCoordinates.pointMap}\leanok
\uses{def:AlgebraicGeometry.ComplexPoint, def:AlgebraicGeometry.LocalEtaleCoordinates.toAffineSpaceOver, def:AlgebraicGeometry.Point.map}
The map on complex points induced by the étale coordinate morphism.

\noindent\emph{Lean:} \texttt{def AlgebraicGeometry.LocalEtaleCoordinates.pointMap}, \href{https://github.com/Paul-Lez/HodgeConjecture/blob/main/HodgeConjecture/Lemmas/AlgebraicGeometry/SmoothComplexCoordinates.lean#L212}{\texttt{HodgeConjecture/Lemmas/AlgebraicGeometry/SmoothComplexCoordinates.lean:212}}.
\end{definition}

\begin{definition}[AlgebraicGeometry.LocalEtaleCoordinates.analyticCoordinates]\label{def:AlgebraicGeometry.LocalEtaleCoordinates.analyticCoordinates}
\lean{AlgebraicGeometry.LocalEtaleCoordinates.analyticCoordinates}\leanok
\uses{def:AlgebraicGeometry.ComplexPoint.affineSpaceEquiv, def:AlgebraicGeometry.LocalEtaleCoordinates.pointMap}
The ordinary complex-valued coordinate tuple on the analytic neighborhood.

\noindent\emph{Lean:} \texttt{def AlgebraicGeometry.LocalEtaleCoordinates.analyticCoordinates}, \href{https://github.com/Paul-Lez/HodgeConjecture/blob/main/HodgeConjecture/Lemmas/AlgebraicGeometry/SmoothComplexCoordinates.lean#L226}{\texttt{HodgeConjecture/Lemmas/AlgebraicGeometry/SmoothComplexCoordinates.lean:226}}.
\end{definition}

\begin{definition}[AlgebraicGeometry.LocalEtaleCoordinates.ambientAnalyticCoordinates]\label{def:AlgebraicGeometry.LocalEtaleCoordinates.ambientAnalyticCoordinates}
\lean{AlgebraicGeometry.LocalEtaleCoordinates.ambientAnalyticCoordinates}\leanok
\uses{def:AlgebraicGeometry.ComplexPoint.openEquiv, def:AlgebraicGeometry.LocalEtaleCoordinates.analyticCoordinates}
Étale coordinates written on the corresponding open subset of the ambient complex points.

\noindent\emph{Lean:} \texttt{def AlgebraicGeometry.LocalEtaleCoordinates.ambientAnalyticCoordinates}, \href{https://github.com/Paul-Lez/HodgeConjecture/blob/main/HodgeConjecture/Lemmas/AlgebraicGeometry/SmoothComplexCoordinates.lean#L241}{\texttt{HodgeConjecture/Lemmas/AlgebraicGeometry/SmoothComplexCoordinates.lean:241}}.
\end{definition}

\begin{definition}[AlgebraicGeometry.LocalEtaleCoordinates.coordinateRingAlgebra]\label{def:AlgebraicGeometry.LocalEtaleCoordinates.coordinateRingAlgebra}
\lean{AlgebraicGeometry.LocalEtaleCoordinates.coordinateRingAlgebra}\leanok
\uses{def:AlgebraicGeometry.ComplexAlgHom.complexPolynomialRing, def:AlgebraicGeometry.LocalEtaleCoordinates.coordinateRingHomOnOpen}
\noindent\emph{Lean:} \texttt{def AlgebraicGeometry.LocalEtaleCoordinates.coordinateRingAlgebra}, \href{https://github.com/Paul-Lez/HodgeConjecture/blob/main/HodgeConjecture/Lemmas/AlgebraicGeometry/SmoothComplexCoordinates.lean#L257}{\texttt{HodgeConjecture/Lemmas/AlgebraicGeometry/SmoothComplexCoordinates.lean:257}}.
\end{definition}

\begin{definition}[AlgebraicGeometry.LocalEtaleCoordinates.coordinateRingComplexAlgebra]\label{def:AlgebraicGeometry.LocalEtaleCoordinates.coordinateRingComplexAlgebra}
\lean{AlgebraicGeometry.LocalEtaleCoordinates.coordinateRingComplexAlgebra}\leanok
\uses{def:AlgebraicGeometry.LocalEtaleCoordinates.coordinateRingHomOnOpen}
\noindent\emph{Lean:} \texttt{def AlgebraicGeometry.LocalEtaleCoordinates.coordinateRingComplexAlgebra}, \href{https://github.com/Paul-Lez/HodgeConjecture/blob/main/HodgeConjecture/Lemmas/AlgebraicGeometry/SmoothComplexCoordinates.lean#L262}{\texttt{HodgeConjecture/Lemmas/AlgebraicGeometry/SmoothComplexCoordinates.lean:262}}.
\end{definition}

\begin{definition}[AlgebraicGeometry.LocalEtaleCoordinates.affineSpecPointHomeomorph]\label{def:AlgebraicGeometry.LocalEtaleCoordinates.affineSpecPointHomeomorph}
\lean{AlgebraicGeometry.LocalEtaleCoordinates.affineSpecPointHomeomorph}\leanok
\uses{def:AlgebraicGeometry.ComplexPoint, def:AlgebraicGeometry.ComplexPoint.affineSpecStructureMap, def:AlgebraicGeometry.ComplexPoint.openScheme, def:AlgebraicGeometry.LocalEtaleCoordinates.coordinateRingComplexAlgebra, def:AlgebraicGeometry.Point.isoMapHomeomorph}
The affine neighborhood's complex points are homeomorphic to the complex points of the
spectrum of its global sections.

\noindent\emph{Lean:} \texttt{def AlgebraicGeometry.LocalEtaleCoordinates.affineSpecPointHomeomorph}, \href{https://github.com/Paul-Lez/HodgeConjecture/blob/main/HodgeConjecture/Lemmas/AlgebraicGeometry/SmoothComplexCoordinates.lean#L291}{\texttt{HodgeConjecture/Lemmas/AlgebraicGeometry/SmoothComplexCoordinates.lean:291}}.
\end{definition}

\begin{definition}[AlgebraicGeometry.LocalEtaleCoordinates.pointAlgHomHomeomorph]\label{def:AlgebraicGeometry.LocalEtaleCoordinates.pointAlgHomHomeomorph}
\lean{AlgebraicGeometry.LocalEtaleCoordinates.pointAlgHomHomeomorph}\leanok
\uses{def:AlgebraicGeometry.ComplexPoint.affineSpecHomeomorph, def:AlgebraicGeometry.LocalEtaleCoordinates.affineSpecPointHomeomorph}
Complex points of an affine neighborhood as complex-valued algebra homomorphisms on its
coordinate ring.

\noindent\emph{Lean:} \texttt{def AlgebraicGeometry.LocalEtaleCoordinates.pointAlgHomHomeomorph}, \href{https://github.com/Paul-Lez/HodgeConjecture/blob/main/HodgeConjecture/Lemmas/AlgebraicGeometry/SmoothComplexCoordinates.lean#L316}{\texttt{HodgeConjecture/Lemmas/AlgebraicGeometry/SmoothComplexCoordinates.lean:316}}.
\end{definition}

\begin{definition}[AlgebraicGeometry.LocalEtaleCoordinates.algebraicCoordinates]\label{def:AlgebraicGeometry.LocalEtaleCoordinates.algebraicCoordinates}
\lean{AlgebraicGeometry.LocalEtaleCoordinates.algebraicCoordinates}\leanok
\uses{def:AlgebraicGeometry.ComplexAlgHom.etaleBaseAlgHom, def:AlgebraicGeometry.ComplexAlgHom.mvPolynomialAlgHomHomeomorph, def:AlgebraicGeometry.LocalEtaleCoordinates.coordinateRingAlgebra, def:AlgebraicGeometry.LocalEtaleCoordinates.pointAlgHomHomeomorph}
The analytic coordinate map reconstructed through the affine ring and its étale polynomial
subalgebra.

\noindent\emph{Lean:} \texttt{def AlgebraicGeometry.LocalEtaleCoordinates.algebraicCoordinates}, \href{https://github.com/Paul-Lez/HodgeConjecture/blob/main/HodgeConjecture/Lemmas/AlgebraicGeometry/SmoothComplexCoordinates.lean#L348}{\texttt{HodgeConjecture/Lemmas/AlgebraicGeometry/SmoothComplexCoordinates.lean:348}}.
\end{definition}

\begin{definition}[AlgebraicGeometry.LocalEtaleCoordinates.ambientCoordinateOpen]\label{def:AlgebraicGeometry.LocalEtaleCoordinates.ambientCoordinateOpen}
\lean{AlgebraicGeometry.LocalEtaleCoordinates.ambientCoordinateOpen}\leanok
\uses{def:AlgebraicGeometry.LocalEtaleCoordinates}
The ambient open on which a chosen coordinate section is naturally expressed.

\noindent\emph{Lean:} \texttt{abbrev AlgebraicGeometry.LocalEtaleCoordinates.ambientCoordinateOpen}, \href{https://github.com/Paul-Lez/HodgeConjecture/blob/main/HodgeConjecture/Lemmas/AlgebraicGeometry/SmoothComplexCoordinates.lean#L399}{\texttt{HodgeConjecture/Lemmas/AlgebraicGeometry/SmoothComplexCoordinates.lean:399}}.
\end{definition}

\begin{definition}[AlgebraicGeometry.LocalEtaleCoordinates.ambientCoordinateSection]\label{def:AlgebraicGeometry.LocalEtaleCoordinates.ambientCoordinateSection}
\lean{AlgebraicGeometry.LocalEtaleCoordinates.ambientCoordinateSection}\leanok
\uses{def:AlgebraicGeometry.LocalEtaleCoordinates.ambientCoordinateOpen, def:AlgebraicGeometry.LocalEtaleCoordinates.coordinateRingHomOnOpen}
A coordinate function transported from the affine open subscheme to an open of the ambient
scheme.

\noindent\emph{Lean:} \texttt{def AlgebraicGeometry.LocalEtaleCoordinates.ambientCoordinateSection}, \href{https://github.com/Paul-Lez/HodgeConjecture/blob/main/HodgeConjecture/Lemmas/AlgebraicGeometry/SmoothComplexCoordinates.lean#L402}{\texttt{HodgeConjecture/Lemmas/AlgebraicGeometry/SmoothComplexCoordinates.lean:402}}.
\end{definition}

\begin{definition}[AlgebraicGeometry.LocalEtaleCoordinates.ambientProjectionChart]\label{def:AlgebraicGeometry.LocalEtaleCoordinates.ambientProjectionChart}
\lean{AlgebraicGeometry.LocalEtaleCoordinates.ambientProjectionChart}\leanok
\uses{def:AlgebraicGeometry.ComplexAlgHom.etaleAlgHomProjectionChart, def:AlgebraicGeometry.ComplexPoint.openHomeomorph, def:AlgebraicGeometry.LocalEtaleCoordinates.coordinateRingAlgebra, def:AlgebraicGeometry.LocalEtaleCoordinates.pointAlgHomHomeomorph}
The explicit analytic projection chart centered at a point of the ambient coordinate
neighborhood. Unlike a chart chosen only from local-homeomorphism existence, its inverse retains
the standard étale construction and its analyticity theorem.

\noindent\emph{Lean:} \texttt{def AlgebraicGeometry.LocalEtaleCoordinates.ambientProjectionChart}, \href{https://github.com/Paul-Lez/HodgeConjecture/blob/main/HodgeConjecture/Lemmas/AlgebraicGeometry/SmoothComplexCoordinates.lean#L442}{\texttt{HodgeConjecture/Lemmas/AlgebraicGeometry/SmoothComplexCoordinates.lean:442}}.
\end{definition}

\begin{definition}[AlgebraicGeometry.ComplexPoint.localEtaleCoordinates]\label{def:AlgebraicGeometry.ComplexPoint.localEtaleCoordinates}
\lean{AlgebraicGeometry.ComplexPoint.localEtaleCoordinates}\leanok
\uses{def:AlgebraicGeometry.ComplexPoint, def:AlgebraicGeometry.Point.underlying, def:AlgebraicGeometry.localEtaleCoordinates}
Étale coordinates of the specified relative dimension chosen around the underlying point of a
complex point.

\noindent\emph{Lean:} \texttt{def AlgebraicGeometry.ComplexPoint.localEtaleCoordinates}, \href{https://github.com/Paul-Lez/HodgeConjecture/blob/main/HodgeConjecture/Lemmas/AlgebraicGeometry/SmoothComplexCoordinates.lean#L617}{\texttt{HodgeConjecture/Lemmas/AlgebraicGeometry/SmoothComplexCoordinates.lean:617}}.
\end{definition}

\section{Dimension of a reduced cycle component}\label{sec:HodgeConjecture.Lemmas.AlgebraicGeometry.CycleComponentDimension}
\noindent\emph{File:} \href{https://github.com/Paul-Lez/HodgeConjecture/blob/main/HodgeConjecture/Lemmas/AlgebraicGeometry/CycleComponentDimension.lean}{\texttt{HodgeConjecture/Lemmas/AlgebraicGeometry/CycleComponentDimension.lean}}.

This file relates the order-theoretic dimension of the reduced closure of a scheme point to the
specialization order in the ambient scheme.  These facts separate the purely topological part of
the component dimension calculation from the catenary dimension formula needed for a smooth
variety.

Mathlib currently has no catenary or equidimensional scheme API, and its
\texttt{SmoothOfRelativeDimension} API does not relate relative dimension to \texttt{Order.height} or
\texttt{Order.coheight}.  Consequently this file proves the unconditional identity
\texttt{dim(closure \{x\}) = height(x)}, but does not claim the further smooth-variety formula
\texttt{height(x) + coheight(x) = dim(X)}.

\begin{definition}[AlgebraicGeometry.cycleComponentOrderIsoIic]\label{def:AlgebraicGeometry.cycleComponentOrderIsoIic}
\lean{AlgebraicGeometry.cycleComponentOrderIsoIic}\leanok
\uses{def:AlgebraicGeometry.cycleComponentι}
The points of the reduced closure of \texttt{x} are exactly the specializations below \texttt{x}.

This is an order isomorphism for the specialization preorders. It is the order-theoretic core of
the dimension calculation for a cycle component.

\noindent\emph{Lean:} \texttt{def AlgebraicGeometry.cycleComponentOrderIsoIic}, \href{https://github.com/Paul-Lez/HodgeConjecture/blob/main/HodgeConjecture/Lemmas/AlgebraicGeometry/CycleComponentDimension.lean#L57}{\texttt{HodgeConjecture/Lemmas/AlgebraicGeometry/CycleComponentDimension.lean:57}}.
\end{definition}

\section{The topological manifold of complex points}\label{sec:HodgeConjecture.Lemmas.AlgebraicGeometry.ComplexManifold}
\noindent\emph{File:} \href{https://github.com/Paul-Lez/HodgeConjecture/blob/main/HodgeConjecture/Lemmas/AlgebraicGeometry/ComplexManifold.lean}{\texttt{HodgeConjecture/Lemmas/AlgebraicGeometry/ComplexManifold.lean}}.

If a complex scheme is smooth of relative dimension \texttt{d}, its constructed analytic complex-point
space is locally homeomorphic to \texttt{ℂ\^{}d}. This file turns the chosen algebraic étale coordinates into
an actual \texttt{ChartedSpace (Fin d → ℂ)} structure and proves that its transition maps are
holomorphic. So the complex points form a complex-analytic manifold of complex dimension \texttt{d}, in
particular a topological manifold of real dimension \texttt{2 * d}; both structures are instances.

The charts are not extra input. At a complex point we choose the affine étale coordinates supplied
by relative-dimensional smoothness, take a local inverse for the proven local homeomorphism, and
extend that chart from the corresponding analytic open subset to the ambient space.

\begin{definition}[AlgebraicGeometry.ComplexPoint.coordinateNeighborhood]\label{def:AlgebraicGeometry.ComplexPoint.coordinateNeighborhood}
\lean{AlgebraicGeometry.ComplexPoint.coordinateNeighborhood}\leanok
\uses{def:AlgebraicGeometry.ComplexPoint.localEtaleCoordinates, def:AlgebraicGeometry.Point.overOpen}
The analytic open subset on which the chosen coordinates at \texttt{z} are defined.

\noindent\emph{Lean:} \texttt{abbrev AlgebraicGeometry.ComplexPoint.coordinateNeighborhood}, \href{https://github.com/Paul-Lez/HodgeConjecture/blob/main/HodgeConjecture/Lemmas/AlgebraicGeometry/ComplexManifold.lean#L53}{\texttt{HodgeConjecture/Lemmas/AlgebraicGeometry/ComplexManifold.lean:53}}.
\end{definition}

\begin{definition}[AlgebraicGeometry.ComplexPoint.pointInCoordinateNeighborhood]\label{def:AlgebraicGeometry.ComplexPoint.pointInCoordinateNeighborhood}
\lean{AlgebraicGeometry.ComplexPoint.pointInCoordinateNeighborhood}\leanok
\uses{def:AlgebraicGeometry.ComplexPoint.coordinateNeighborhood}
The point \texttt{z} regarded as a point of its chosen coordinate neighborhood.

\noindent\emph{Lean:} \texttt{def AlgebraicGeometry.ComplexPoint.pointInCoordinateNeighborhood}, \href{https://github.com/Paul-Lez/HodgeConjecture/blob/main/HodgeConjecture/Lemmas/AlgebraicGeometry/ComplexManifold.lean#L59}{\texttt{HodgeConjecture/Lemmas/AlgebraicGeometry/ComplexManifold.lean:59}}.
\end{definition}

\begin{definition}[AlgebraicGeometry.ComplexPoint.coordinateNeighborhoodChart]\label{def:AlgebraicGeometry.ComplexPoint.coordinateNeighborhoodChart}
\lean{AlgebraicGeometry.ComplexPoint.coordinateNeighborhoodChart}\leanok
\uses{def:AlgebraicGeometry.ComplexPoint.pointInCoordinateNeighborhood, def:AlgebraicGeometry.LocalEtaleCoordinates.ambientProjectionChart}
A local coordinate homeomorphism on the chosen analytic open neighborhood.

\noindent\emph{Lean:} \texttt{def AlgebraicGeometry.ComplexPoint.coordinateNeighborhoodChart}, \href{https://github.com/Paul-Lez/HodgeConjecture/blob/main/HodgeConjecture/Lemmas/AlgebraicGeometry/ComplexManifold.lean#L64}{\texttt{HodgeConjecture/Lemmas/AlgebraicGeometry/ComplexManifold.lean:64}}.
\end{definition}

\begin{definition}[AlgebraicGeometry.ComplexPoint.localChart]\label{def:AlgebraicGeometry.ComplexPoint.localChart}
\lean{AlgebraicGeometry.ComplexPoint.localChart}\leanok
\uses{def:AlgebraicGeometry.ComplexPoint.coordinateNeighborhoodChart}
The chosen local chart at a complex point, extended from its analytic open neighborhood to the
whole complex-point space.

\noindent\emph{Lean:} \texttt{def AlgebraicGeometry.ComplexPoint.localChart}, \href{https://github.com/Paul-Lez/HodgeConjecture/blob/main/HodgeConjecture/Lemmas/AlgebraicGeometry/ComplexManifold.lean#L78}{\texttt{HodgeConjecture/Lemmas/AlgebraicGeometry/ComplexManifold.lean:78}}.
\end{definition}

\section{Residue fields and local sections along a closed immersion}\label{sec:HodgeConjecture.Lemmas.AlgebraicGeometry.ClosedImmersionResidueField}
\noindent\emph{File:} \href{https://github.com/Paul-Lez/HodgeConjecture/blob/main/HodgeConjecture/Lemmas/AlgebraicGeometry/ClosedImmersionResidueField.lean}{\texttt{HodgeConjecture/Lemmas/AlgebraicGeometry/ClosedImmersionResidueField.lean}}.

A closed immersion is surjective on stalks of structure sheaves, hence on residue fields; since a
map of fields is injective, the residue fields at corresponding points are canonically isomorphic.

Surjectivity on stalks also says that a section of the closed subscheme is, near each of its
points, the restriction of a section defined on an open set of the ambient scheme. That local
lifting is what transports the analytic topology across a closed immersion.

\begin{definition}[AlgebraicGeometry.IsClosedImmersion.residueFieldIso]\label{def:AlgebraicGeometry.IsClosedImmersion.residueFieldIso}
\lean{AlgebraicGeometry.IsClosedImmersion.residueFieldIso}\leanok
The residue-field isomorphism induced by a closed immersion.

\noindent\emph{Lean:} \texttt{def AlgebraicGeometry.IsClosedImmersion.residueFieldIso}, \href{https://github.com/Paul-Lez/HodgeConjecture/blob/main/HodgeConjecture/Lemmas/AlgebraicGeometry/ClosedImmersionResidueField.lean#L53}{\texttt{HodgeConjecture/Lemmas/AlgebraicGeometry/ClosedImmersionResidueField.lean:53}}.
\end{definition}

\section{Projective analytification}\label{sec:HodgeConjecture.Lemmas.AlgebraicGeometry.ProjectiveAnalytification}
\noindent\emph{File:} \href{https://github.com/Paul-Lez/HodgeConjecture/blob/main/HodgeConjecture/Lemmas/AlgebraicGeometry/ProjectiveAnalytification.lean}{\texttt{HodgeConjecture/Lemmas/AlgebraicGeometry/ProjectiveAnalytification.lean}}.

This file supplies the compact topological model of finite-dimensional complex projective space.
It equips linear-algebraic projectivization with the quotient topology from nonzero coordinate
vectors. Normalizing coordinates gives a continuous surjection from the unit sphere, and hence
proves compactness.

It also constructs the standard scheme-theoretic projective points from homogeneous coordinates.
Standard affine charts prove that this construction is a continuous bijection from linear
projectivization onto the complex points of scheme-theoretic projective space. This proves that
scheme-theoretic projective space is analytically compact.

\begin{definition}[AlgebraicGeometry.ComplexProjectiveSpace.CoordinateSpace]\label{def:AlgebraicGeometry.ComplexProjectiveSpace.CoordinateSpace}
\lean{AlgebraicGeometry.ComplexProjectiveSpace.CoordinateSpace}\leanok
The coordinate vector space underlying complex projective \texttt{n}-space.

\noindent\emph{Lean:} \texttt{abbrev AlgebraicGeometry.ComplexProjectiveSpace.CoordinateSpace}, \href{https://github.com/Paul-Lez/HodgeConjecture/blob/main/HodgeConjecture/Lemmas/AlgebraicGeometry/ProjectiveAnalytification.lean#L50}{\texttt{HodgeConjecture/Lemmas/AlgebraicGeometry/ProjectiveAnalytification.lean:50}}.
\end{definition}

\begin{definition}[AlgebraicGeometry.ComplexProjectiveSpace.UniversalRing]\label{def:AlgebraicGeometry.ComplexProjectiveSpace.UniversalRing}
\lean{AlgebraicGeometry.ComplexProjectiveSpace.UniversalRing}\leanok
The polynomial ring in homogeneous coordinates over the integral model.

\noindent\emph{Lean:} \texttt{abbrev AlgebraicGeometry.ComplexProjectiveSpace.UniversalRing}, \href{https://github.com/Paul-Lez/HodgeConjecture/blob/main/HodgeConjecture/Lemmas/AlgebraicGeometry/ProjectiveAnalytification.lean#L53}{\texttt{HodgeConjecture/Lemmas/AlgebraicGeometry/ProjectiveAnalytification.lean:53}}.
\end{definition}

\begin{definition}[AlgebraicGeometry.ComplexProjectiveSpace.UniversalGrading]\label{def:AlgebraicGeometry.ComplexProjectiveSpace.UniversalGrading}
\lean{AlgebraicGeometry.ComplexProjectiveSpace.UniversalGrading}\leanok
The standard grading on the homogeneous coordinate ring.

\noindent\emph{Lean:} \texttt{abbrev AlgebraicGeometry.ComplexProjectiveSpace.UniversalGrading}, \href{https://github.com/Paul-Lez/HodgeConjecture/blob/main/HodgeConjecture/Lemmas/AlgebraicGeometry/ProjectiveAnalytification.lean#L56}{\texttt{HodgeConjecture/Lemmas/AlgebraicGeometry/ProjectiveAnalytification.lean:56}}.
\end{definition}

\begin{definition}[AlgebraicGeometry.ComplexProjectiveSpace.normalize]\label{def:AlgebraicGeometry.ComplexProjectiveSpace.normalize}
\lean{AlgebraicGeometry.ComplexProjectiveSpace.normalize}\leanok
\uses{def:AlgebraicGeometry.ComplexProjectiveSpace.CoordinateSpace}
Normalize a nonzero vector to a vector on the unit sphere.

\noindent\emph{Lean:} \texttt{def AlgebraicGeometry.ComplexProjectiveSpace.normalize}, \href{https://github.com/Paul-Lez/HodgeConjecture/blob/main/HodgeConjecture/Lemmas/AlgebraicGeometry/ProjectiveAnalytification.lean#L67}{\texttt{HodgeConjecture/Lemmas/AlgebraicGeometry/ProjectiveAnalytification.lean:67}}.
\end{definition}

\begin{definition}[AlgebraicGeometry.ComplexProjectiveSpace.sphereToProjectivization]\label{def:AlgebraicGeometry.ComplexProjectiveSpace.sphereToProjectivization}
\lean{AlgebraicGeometry.ComplexProjectiveSpace.sphereToProjectivization}\leanok
\uses{def:AlgebraicGeometry.ComplexProjectiveSpace.CoordinateSpace}
The quotient map from the unit sphere to complex projective space.

\noindent\emph{Lean:} \texttt{def AlgebraicGeometry.ComplexProjectiveSpace.sphereToProjectivization}, \href{https://github.com/Paul-Lez/HodgeConjecture/blob/main/HodgeConjecture/Lemmas/AlgebraicGeometry/ProjectiveAnalytification.lean#L81}{\texttt{HodgeConjecture/Lemmas/AlgebraicGeometry/ProjectiveAnalytification.lean:81}}.
\end{definition}

\begin{definition}[AlgebraicGeometry.ComplexProjectiveSpace.coordinateGlobalSectionsHom]\label{def:AlgebraicGeometry.ComplexProjectiveSpace.coordinateGlobalSectionsHom}
\lean{AlgebraicGeometry.ComplexProjectiveSpace.coordinateGlobalSectionsHom}\leanok
\uses{def:AlgebraicGeometry.ComplexProjectiveSpace.CoordinateSpace, def:AlgebraicGeometry.ComplexProjectiveSpace.UniversalRing}
Evaluation at a coordinate vector, regarded as a map to the global functions on
\texttt{Spec ℂ}.

\noindent\emph{Lean:} \texttt{def AlgebraicGeometry.ComplexProjectiveSpace.coordinateGlobalSectionsHom}, \href{https://github.com/Paul-Lez/HodgeConjecture/blob/main/HodgeConjecture/Lemmas/AlgebraicGeometry/ProjectiveAnalytification.lean#L117}{\texttt{HodgeConjecture/Lemmas/AlgebraicGeometry/ProjectiveAnalytification.lean:117}}.
\end{definition}

\begin{definition}[AlgebraicGeometry.ComplexProjectiveSpace.coordinateEvaluationHom]\label{def:AlgebraicGeometry.ComplexProjectiveSpace.coordinateEvaluationHom}
\lean{AlgebraicGeometry.ComplexProjectiveSpace.coordinateEvaluationHom}\leanok
\uses{def:AlgebraicGeometry.ComplexProjectiveSpace.coordinateGlobalSectionsHom}
Evaluation of the integral homogeneous coordinate ring at a complex vector, with the global
functions on \texttt{Spec ℂ} identified with \texttt{ℂ}.

\noindent\emph{Lean:} \texttt{def AlgebraicGeometry.ComplexProjectiveSpace.coordinateEvaluationHom}, \href{https://github.com/Paul-Lez/HodgeConjecture/blob/main/HodgeConjecture/Lemmas/AlgebraicGeometry/ProjectiveAnalytification.lean#L177}{\texttt{HodgeConjecture/Lemmas/AlgebraicGeometry/ProjectiveAnalytification.lean:177}}.
\end{definition}

\begin{definition}[AlgebraicGeometry.ComplexProjectiveSpace.awayCoordinateEvaluation]\label{def:AlgebraicGeometry.ComplexProjectiveSpace.awayCoordinateEvaluation}
\lean{AlgebraicGeometry.ComplexProjectiveSpace.awayCoordinateEvaluation}\leanok
\uses{def:AlgebraicGeometry.ComplexProjectiveSpace.UniversalGrading, def:AlgebraicGeometry.ComplexProjectiveSpace.coordinateEvaluationHom}
Evaluation of regular functions on the standard projective chart where the \texttt{i}-th coordinate
is nonzero.

\noindent\emph{Lean:} \texttt{def AlgebraicGeometry.ComplexProjectiveSpace.awayCoordinateEvaluation}, \href{https://github.com/Paul-Lez/HodgeConjecture/blob/main/HodgeConjecture/Lemmas/AlgebraicGeometry/ProjectiveAnalytification.lean#L195}{\texttt{HodgeConjecture/Lemmas/AlgebraicGeometry/ProjectiveAnalytification.lean:195}}.
\end{definition}

\begin{definition}[AlgebraicGeometry.ComplexProjectiveSpace.awayHomogeneousEvaluation]\label{def:AlgebraicGeometry.ComplexProjectiveSpace.awayHomogeneousEvaluation}
\lean{AlgebraicGeometry.ComplexProjectiveSpace.awayHomogeneousEvaluation}\leanok
\uses{def:AlgebraicGeometry.ComplexProjectiveSpace.UniversalGrading, def:AlgebraicGeometry.ComplexProjectiveSpace.coordinateEvaluationHom}
Evaluation on the degree-zero localization away from an arbitrary homogeneous polynomial.

\noindent\emph{Lean:} \texttt{def AlgebraicGeometry.ComplexProjectiveSpace.awayHomogeneousEvaluation}, \href{https://github.com/Paul-Lez/HodgeConjecture/blob/main/HodgeConjecture/Lemmas/AlgebraicGeometry/ProjectiveAnalytification.lean#L251}{\texttt{HodgeConjecture/Lemmas/AlgebraicGeometry/ProjectiveAnalytification.lean:251}}.
\end{definition}

\begin{definition}[AlgebraicGeometry.ComplexProjectiveSpace.chartCoordinate]\label{def:AlgebraicGeometry.ComplexProjectiveSpace.chartCoordinate}
\lean{AlgebraicGeometry.ComplexProjectiveSpace.chartCoordinate}\leanok
\uses{def:AlgebraicGeometry.ComplexProjectiveSpace.UniversalGrading}
The regular function \texttt{X\_j / X\_i} on the standard projective chart where \texttt{X\_i} is
nonzero.

\noindent\emph{Lean:} \texttt{def AlgebraicGeometry.ComplexProjectiveSpace.chartCoordinate}, \href{https://github.com/Paul-Lez/HodgeConjecture/blob/main/HodgeConjecture/Lemmas/AlgebraicGeometry/ProjectiveAnalytification.lean#L279}{\texttt{HodgeConjecture/Lemmas/AlgebraicGeometry/ProjectiveAnalytification.lean:279}}.
\end{definition}

\begin{definition}[AlgebraicGeometry.ComplexProjectiveSpace.vectorOfAwayRingHom]\label{def:AlgebraicGeometry.ComplexProjectiveSpace.vectorOfAwayRingHom}
\lean{AlgebraicGeometry.ComplexProjectiveSpace.vectorOfAwayRingHom}\leanok
\uses{def:AlgebraicGeometry.ComplexProjectiveSpace.CoordinateSpace, def:AlgebraicGeometry.ComplexProjectiveSpace.chartCoordinate}
Homogeneous coordinates reconstructed from a complex point of one standard affine chart.

\noindent\emph{Lean:} \texttt{def AlgebraicGeometry.ComplexProjectiveSpace.vectorOfAwayRingHom}, \href{https://github.com/Paul-Lez/HodgeConjecture/blob/main/HodgeConjecture/Lemmas/AlgebraicGeometry/ProjectiveAnalytification.lean#L397}{\texttt{HodgeConjecture/Lemmas/AlgebraicGeometry/ProjectiveAnalytification.lean:397}}.
\end{definition}

\begin{definition}[AlgebraicGeometry.ComplexProjectiveSpace.coordinateSupport]\label{def:AlgebraicGeometry.ComplexProjectiveSpace.coordinateSupport}
\lean{AlgebraicGeometry.ComplexProjectiveSpace.coordinateSupport}\leanok
\uses{def:AlgebraicGeometry.ComplexProjectiveSpace.CoordinateSpace}
The finite set of indices of the nonzero coordinates of a vector.

\noindent\emph{Lean:} \texttt{def AlgebraicGeometry.ComplexProjectiveSpace.coordinateSupport}, \href{https://github.com/Paul-Lez/HodgeConjecture/blob/main/HodgeConjecture/Lemmas/AlgebraicGeometry/ProjectiveAnalytification.lean#L508}{\texttt{HodgeConjecture/Lemmas/AlgebraicGeometry/ProjectiveAnalytification.lean:508}}.
\end{definition}

\begin{definition}[AlgebraicGeometry.ComplexProjectiveSpace.coordinateIndex]\label{def:AlgebraicGeometry.ComplexProjectiveSpace.coordinateIndex}
\lean{AlgebraicGeometry.ComplexProjectiveSpace.coordinateIndex}\leanok
\uses{def:AlgebraicGeometry.ComplexProjectiveSpace.coordinateSupport}
The least index of a nonzero coordinate. This gives a scale-independent choice of a
standard projective chart.

\noindent\emph{Lean:} \texttt{def AlgebraicGeometry.ComplexProjectiveSpace.coordinateIndex}, \href{https://github.com/Paul-Lez/HodgeConjecture/blob/main/HodgeConjecture/Lemmas/AlgebraicGeometry/ProjectiveAnalytification.lean#L518}{\texttt{HodgeConjecture/Lemmas/AlgebraicGeometry/ProjectiveAnalytification.lean:518}}.
\end{definition}

\begin{definition}[AlgebraicGeometry.ComplexProjectiveSpace.chartIntegralProjAt]\label{def:AlgebraicGeometry.ComplexProjectiveSpace.chartIntegralProjAt}
\lean{AlgebraicGeometry.ComplexProjectiveSpace.chartIntegralProjAt}\leanok
\uses{def:AlgebraicGeometry.ComplexProjectiveSpace.awayCoordinateEvaluation}
The projective point defined directly in one standard affine chart.

\noindent\emph{Lean:} \texttt{def AlgebraicGeometry.ComplexProjectiveSpace.chartIntegralProjAt}, \href{https://github.com/Paul-Lez/HodgeConjecture/blob/main/HodgeConjecture/Lemmas/AlgebraicGeometry/ProjectiveAnalytification.lean#L541}{\texttt{HodgeConjecture/Lemmas/AlgebraicGeometry/ProjectiveAnalytification.lean:541}}.
\end{definition}

\begin{definition}[AlgebraicGeometry.ComplexProjectiveSpace.chartIntegralProj]\label{def:AlgebraicGeometry.ComplexProjectiveSpace.chartIntegralProj}
\lean{AlgebraicGeometry.ComplexProjectiveSpace.chartIntegralProj}\leanok
\uses{def:AlgebraicGeometry.ComplexProjectiveSpace.chartIntegralProjAt, def:AlgebraicGeometry.ComplexProjectiveSpace.coordinateIndex}
The projective point obtained in the least standard chart containing it.

\noindent\emph{Lean:} \texttt{def AlgebraicGeometry.ComplexProjectiveSpace.chartIntegralProj}, \href{https://github.com/Paul-Lez/HodgeConjecture/blob/main/HodgeConjecture/Lemmas/AlgebraicGeometry/ProjectiveAnalytification.lean#L686}{\texttt{HodgeConjecture/Lemmas/AlgebraicGeometry/ProjectiveAnalytification.lean:686}}.
\end{definition}

\begin{definition}[AlgebraicGeometry.ComplexProjectiveSpace.vectorToProjectiveSpace]\label{def:AlgebraicGeometry.ComplexProjectiveSpace.vectorToProjectiveSpace}
\lean{AlgebraicGeometry.ComplexProjectiveSpace.vectorToProjectiveSpace}\leanok
\uses{def:AlgebraicGeometry.ComplexProjectiveSpace.chartIntegralProj, def:ProjectiveSpace}
Nonzero homogeneous coordinates define a point of scheme-theoretic projective space over
\texttt{Spec ℂ}.

\noindent\emph{Lean:} \texttt{def AlgebraicGeometry.ComplexProjectiveSpace.vectorToProjectiveSpace}, \href{https://github.com/Paul-Lez/HodgeConjecture/blob/main/HodgeConjecture/Lemmas/AlgebraicGeometry/ProjectiveAnalytification.lean#L981}{\texttt{HodgeConjecture/Lemmas/AlgebraicGeometry/ProjectiveAnalytification.lean:981}}.
\end{definition}

\begin{definition}[AlgebraicGeometry.ComplexProjectiveSpace.vectorToComplexPoint]\label{def:AlgebraicGeometry.ComplexProjectiveSpace.vectorToComplexPoint}
\lean{AlgebraicGeometry.ComplexProjectiveSpace.vectorToComplexPoint}\leanok
\uses{def:AlgebraicGeometry.ComplexPoint, def:AlgebraicGeometry.ComplexProjectiveSpace.vectorToProjectiveSpace, def:ProjectiveSpace.toBase}
Nonzero homogeneous coordinates define a complex point of scheme-theoretic projective
space.

\noindent\emph{Lean:} \texttt{def AlgebraicGeometry.ComplexProjectiveSpace.vectorToComplexPoint}, \href{https://github.com/Paul-Lez/HodgeConjecture/blob/main/HodgeConjecture/Lemmas/AlgebraicGeometry/ProjectiveAnalytification.lean#L1004}{\texttt{HodgeConjecture/Lemmas/AlgebraicGeometry/ProjectiveAnalytification.lean:1004}}.
\end{definition}

\begin{definition}[AlgebraicGeometry.ComplexProjectiveSpace.projectivizationToComplexPoint]\label{def:AlgebraicGeometry.ComplexProjectiveSpace.projectivizationToComplexPoint}
\lean{AlgebraicGeometry.ComplexProjectiveSpace.projectivizationToComplexPoint}\leanok
\uses{def:AlgebraicGeometry.ComplexProjectiveSpace.vectorToComplexPoint}
Homogeneous coordinates define an actual map from linear projectivization to the complex
points of scheme-theoretic projective space.

\noindent\emph{Lean:} \texttt{def AlgebraicGeometry.ComplexProjectiveSpace.projectivizationToComplexPoint}, \href{https://github.com/Paul-Lez/HodgeConjecture/blob/main/HodgeConjecture/Lemmas/AlgebraicGeometry/ProjectiveAnalytification.lean#L1040}{\texttt{HodgeConjecture/Lemmas/AlgebraicGeometry/ProjectiveAnalytification.lean:1040}}.
\end{definition}

\begin{definition}[AlgebraicGeometry.ComplexProjectiveSpace.chartPolynomialEvaluationHom]\label{def:AlgebraicGeometry.ComplexProjectiveSpace.chartPolynomialEvaluationHom}
\lean{AlgebraicGeometry.ComplexProjectiveSpace.chartPolynomialEvaluationHom}\leanok
\uses{def:AlgebraicGeometry.ComplexProjectiveSpace.UniversalRing}
\noindent\emph{Lean:} \texttt{def AlgebraicGeometry.ComplexProjectiveSpace.chartPolynomialEvaluationHom}, \href{https://github.com/Paul-Lez/HodgeConjecture/blob/main/HodgeConjecture/Lemmas/AlgebraicGeometry/ProjectiveAnalytification.lean#L1090}{\texttt{HodgeConjecture/Lemmas/AlgebraicGeometry/ProjectiveAnalytification.lean:1090}}.
\end{definition}

\begin{definition}[AlgebraicGeometry.ComplexProjectiveSpace.chartAwayPolynomialHom]\label{def:AlgebraicGeometry.ComplexProjectiveSpace.chartAwayPolynomialHom}
\lean{AlgebraicGeometry.ComplexProjectiveSpace.chartAwayPolynomialHom}\leanok
\uses{def:AlgebraicGeometry.ComplexProjectiveSpace.UniversalGrading, def:AlgebraicGeometry.ComplexProjectiveSpace.chartPolynomialEvaluationHom}
\noindent\emph{Lean:} \texttt{def AlgebraicGeometry.ComplexProjectiveSpace.chartAwayPolynomialHom}, \href{https://github.com/Paul-Lez/HodgeConjecture/blob/main/HodgeConjecture/Lemmas/AlgebraicGeometry/ProjectiveAnalytification.lean#L1102}{\texttt{HodgeConjecture/Lemmas/AlgebraicGeometry/ProjectiveAnalytification.lean:1102}}.
\end{definition}

\begin{definition}[AlgebraicGeometry.ComplexProjectiveSpace.chartAffineToProj]\label{def:AlgebraicGeometry.ComplexProjectiveSpace.chartAffineToProj}
\lean{AlgebraicGeometry.ComplexProjectiveSpace.chartAffineToProj}\leanok
\uses{def:AlgebraicGeometry.ComplexPoint.complexAffineSpace, def:AlgebraicGeometry.ComplexProjectiveSpace.chartAwayPolynomialHom}
\noindent\emph{Lean:} \texttt{def AlgebraicGeometry.ComplexProjectiveSpace.chartAffineToProj}, \href{https://github.com/Paul-Lez/HodgeConjecture/blob/main/HodgeConjecture/Lemmas/AlgebraicGeometry/ProjectiveAnalytification.lean#L1128}{\texttt{HodgeConjecture/Lemmas/AlgebraicGeometry/ProjectiveAnalytification.lean:1128}}.
\end{definition}

\begin{definition}[AlgebraicGeometry.ComplexProjectiveSpace.chartAffineToProjectiveSpace]\label{def:AlgebraicGeometry.ComplexProjectiveSpace.chartAffineToProjectiveSpace}
\lean{AlgebraicGeometry.ComplexProjectiveSpace.chartAffineToProjectiveSpace}\leanok
\uses{def:AlgebraicGeometry.ComplexProjectiveSpace.chartAffineToProj, def:ProjectiveSpace}
\noindent\emph{Lean:} \texttt{def AlgebraicGeometry.ComplexProjectiveSpace.chartAffineToProjectiveSpace}, \href{https://github.com/Paul-Lez/HodgeConjecture/blob/main/HodgeConjecture/Lemmas/AlgebraicGeometry/ProjectiveAnalytification.lean#L1135}{\texttt{HodgeConjecture/Lemmas/AlgebraicGeometry/ProjectiveAnalytification.lean:1135}}.
\end{definition}

\begin{definition}[AlgebraicGeometry.ComplexProjectiveSpace.chartAffineComplexPointMap]\label{def:AlgebraicGeometry.ComplexProjectiveSpace.chartAffineComplexPointMap}
\lean{AlgebraicGeometry.ComplexProjectiveSpace.chartAffineComplexPointMap}\leanok
\uses{def:AlgebraicGeometry.ComplexPoint, def:AlgebraicGeometry.ComplexProjectiveSpace.chartAffineToProjectiveSpace, def:AlgebraicGeometry.Point.map, def:ProjectiveSpace.toBase}
\noindent\emph{Lean:} \texttt{def AlgebraicGeometry.ComplexProjectiveSpace.chartAffineComplexPointMap}, \href{https://github.com/Paul-Lez/HodgeConjecture/blob/main/HodgeConjecture/Lemmas/AlgebraicGeometry/ProjectiveAnalytification.lean#L1158}{\texttt{HodgeConjecture/Lemmas/AlgebraicGeometry/ProjectiveAnalytification.lean:1158}}.
\end{definition}

\begin{definition}[AlgebraicGeometry.ComplexProjectiveSpace.vectorChartRatios]\label{def:AlgebraicGeometry.ComplexProjectiveSpace.vectorChartRatios}
\lean{AlgebraicGeometry.ComplexProjectiveSpace.vectorChartRatios}\leanok
\uses{def:AlgebraicGeometry.ComplexProjectiveSpace.CoordinateSpace}
Affine ratio coordinates associated to a vector in the \texttt{i}-th standard chart.

\noindent\emph{Lean:} \texttt{def AlgebraicGeometry.ComplexProjectiveSpace.vectorChartRatios}, \href{https://github.com/Paul-Lez/HodgeConjecture/blob/main/HodgeConjecture/Lemmas/AlgebraicGeometry/ProjectiveAnalytification.lean#L1172}{\texttt{HodgeConjecture/Lemmas/AlgebraicGeometry/ProjectiveAnalytification.lean:1172}}.
\end{definition}

\begin{definition}[AlgebraicGeometry.ComplexProjectiveSpace.vectorChartToAffinePoint]\label{def:AlgebraicGeometry.ComplexProjectiveSpace.vectorChartToAffinePoint}
\lean{AlgebraicGeometry.ComplexProjectiveSpace.vectorChartToAffinePoint}\leanok
\uses{def:AlgebraicGeometry.ComplexPoint.affineSpaceEquiv, def:AlgebraicGeometry.ComplexProjectiveSpace.vectorChartRatios}
\noindent\emph{Lean:} \texttt{def AlgebraicGeometry.ComplexProjectiveSpace.vectorChartToAffinePoint}, \href{https://github.com/Paul-Lez/HodgeConjecture/blob/main/HodgeConjecture/Lemmas/AlgebraicGeometry/ProjectiveAnalytification.lean#L1188}{\texttt{HodgeConjecture/Lemmas/AlgebraicGeometry/ProjectiveAnalytification.lean:1188}}.
\end{definition}

\begin{definition}[AlgebraicGeometry.ComplexProjectiveSpace.nonzeroVectorChart]\label{def:AlgebraicGeometry.ComplexProjectiveSpace.nonzeroVectorChart}
\lean{AlgebraicGeometry.ComplexProjectiveSpace.nonzeroVectorChart}\leanok
\uses{def:AlgebraicGeometry.ComplexProjectiveSpace.CoordinateSpace}
Nonzero vectors whose \texttt{i}-th coordinate is nonzero.

\noindent\emph{Lean:} \texttt{def AlgebraicGeometry.ComplexProjectiveSpace.nonzeroVectorChart}, \href{https://github.com/Paul-Lez/HodgeConjecture/blob/main/HodgeConjecture/Lemmas/AlgebraicGeometry/ProjectiveAnalytification.lean#L1352}{\texttt{HodgeConjecture/Lemmas/AlgebraicGeometry/ProjectiveAnalytification.lean:1352}}.
\end{definition}

\begin{definition}[AlgebraicGeometry.ComplexProjectiveSpace.toVectorChart]\label{def:AlgebraicGeometry.ComplexProjectiveSpace.toVectorChart}
\lean{AlgebraicGeometry.ComplexProjectiveSpace.toVectorChart}\leanok
\uses{def:AlgebraicGeometry.ComplexProjectiveSpace.nonzeroVectorChart}
Forget the redundant global nonvanishing proof on a standard coordinate chart.

\noindent\emph{Lean:} \texttt{def AlgebraicGeometry.ComplexProjectiveSpace.toVectorChart}, \href{https://github.com/Paul-Lez/HodgeConjecture/blob/main/HodgeConjecture/Lemmas/AlgebraicGeometry/ProjectiveAnalytification.lean#L1357}{\texttt{HodgeConjecture/Lemmas/AlgebraicGeometry/ProjectiveAnalytification.lean:1357}}.
\end{definition}

\section{Analytic compactness of a projective complex scheme}\label{sec:HodgeConjecture.Lemmas.AlgebraicGeometry.ProjectiveAnalyticImmersion}
\noindent\emph{File:} \href{https://github.com/Paul-Lez/HodgeConjecture/blob/main/HodgeConjecture/Lemmas/AlgebraicGeometry/ProjectiveAnalyticImmersion.lean}{\texttt{HodgeConjecture/Lemmas/AlgebraicGeometry/ProjectiveAnalyticImmersion.lean}}.

A projective presentation of a complex scheme is a closed immersion into scheme-theoretic
projective space. On complex points it therefore induces a closed topological embedding into the
complex points of projective space, which are analytically compact.

Consequently the analytic complex points of any projective complex scheme form a compact space.

\begin{definition}[AlgebraicGeometry.ProjectiveSpace.Presentation.overImmersion]\label{def:AlgebraicGeometry.ProjectiveSpace.Presentation.overImmersion}
\lean{AlgebraicGeometry.ProjectiveSpace.Presentation.overImmersion}\leanok
\uses{def:ProjectiveSpace.Presentation}
The immersion of a projective presentation, bundled over the complex base.

\noindent\emph{Lean:} \texttt{abbrev AlgebraicGeometry.ProjectiveSpace.Presentation.overImmersion}, \href{https://github.com/Paul-Lez/HodgeConjecture/blob/main/HodgeConjecture/Lemmas/AlgebraicGeometry/ProjectiveAnalyticImmersion.lean#L58}{\texttt{HodgeConjecture/Lemmas/AlgebraicGeometry/ProjectiveAnalyticImmersion.lean:58}}.
\end{definition}

\begin{definition}[AlgebraicGeometry.ProjectiveSpace.Presentation.analyticImmersion]\label{def:AlgebraicGeometry.ProjectiveSpace.Presentation.analyticImmersion}
\lean{AlgebraicGeometry.ProjectiveSpace.Presentation.analyticImmersion}\leanok
\uses{def:AlgebraicGeometry.ComplexPoint, def:AlgebraicGeometry.Point.continuousMap, def:AlgebraicGeometry.ProjectiveSpace.Presentation.overImmersion}
The continuous map on complex points induced by an explicit projective presentation.

\noindent\emph{Lean:} \texttt{def AlgebraicGeometry.ProjectiveSpace.Presentation.analyticImmersion}, \href{https://github.com/Paul-Lez/HodgeConjecture/blob/main/HodgeConjecture/Lemmas/AlgebraicGeometry/ProjectiveAnalyticImmersion.lean#L65}{\texttt{HodgeConjecture/Lemmas/AlgebraicGeometry/ProjectiveAnalyticImmersion.lean:65}}.
\end{definition}

\section{Hausdorff analytifications of projective complex schemes}\label{sec:HodgeConjecture.Lemmas.AlgebraicGeometry.ProjectiveAnalytificationHausdorff}
\noindent\emph{File:} \href{https://github.com/Paul-Lez/HodgeConjecture/blob/main/HodgeConjecture/Lemmas/AlgebraicGeometry/ProjectiveAnalytificationHausdorff.lean}{\texttt{HodgeConjecture/Lemmas/AlgebraicGeometry/ProjectiveAnalytificationHausdorff.lean}}.

The analytification of an affine complex scheme is Hausdorff because global regular functions
separate its complex points. Any two points of finite-dimensional projective space lie in a
common affine basic open: an explicit integral linear form can be chosen not to vanish at either
of two homogeneous coordinate vectors. Pulling this open back along a closed projective
presentation proves that every projective complex scheme has Hausdorff analytification.

\begin{definition}[AlgebraicGeometry.ComplexProjectiveSpace.projectiveSpaceBasicOpen]\label{def:AlgebraicGeometry.ComplexProjectiveSpace.projectiveSpaceBasicOpen}
\lean{AlgebraicGeometry.ComplexProjectiveSpace.projectiveSpaceBasicOpen}\leanok
\uses{def:AlgebraicGeometry.ComplexProjectiveSpace.UniversalGrading, def:AlgebraicGeometry.ComplexProjectiveSpace.UniversalRing, def:ProjectiveSpace}
The inverse image of a projective-spectrum basic open in complex projective space.

\noindent\emph{Lean:} \texttt{def AlgebraicGeometry.ComplexProjectiveSpace.projectiveSpaceBasicOpen}, \href{https://github.com/Paul-Lez/HodgeConjecture/blob/main/HodgeConjecture/Lemmas/AlgebraicGeometry/ProjectiveAnalytificationHausdorff.lean#L188}{\texttt{HodgeConjecture/Lemmas/AlgebraicGeometry/ProjectiveAnalytificationHausdorff.lean:188}}.
\end{definition}

\section{The holomorphic Poincare operator on a complex-point space}\label{sec:HodgeConjecture.Lemmas.AlgebraicGeometry.HolomorphicPoincare}
\noindent\emph{File:} \href{https://github.com/Paul-Lez/HodgeConjecture/blob/main/HodgeConjecture/Lemmas/AlgebraicGeometry/HolomorphicPoincare.lean}{\texttt{HodgeConjecture/Lemmas/AlgebraicGeometry/HolomorphicPoincare.lean}}.

This file applies the holomorphic Poincaré lemma of
\texttt{HodgeConjecture.Lemmas.Analysis.Calculus.DifferentialForm.HolomorphicPoincare} to fixed-chart
evaluations of raw holomorphic forms on the analytic complex-point space of a complex scheme.

\begin{definition}[AlgebraicGeometry.ComplexPoint.fixedChartTransition]\label{def:AlgebraicGeometry.ComplexPoint.fixedChartTransition}
\lean{AlgebraicGeometry.ComplexPoint.fixedChartTransition}\leanok
\uses{def:AlgebraicGeometry.ComplexPoint.localChart}
The change of coordinates from the chart at \texttt{z'} to the chart at \texttt{z}.

\noindent\emph{Lean:} \texttt{def AlgebraicGeometry.ComplexPoint.fixedChartTransition}, \href{https://github.com/Paul-Lez/HodgeConjecture/blob/main/HodgeConjecture/Lemmas/AlgebraicGeometry/HolomorphicPoincare.lean#L53}{\texttt{HodgeConjecture/Lemmas/AlgebraicGeometry/HolomorphicPoincare.lean:53}}.
\end{definition}

\begin{definition}[AlgebraicGeometry.ComplexPoint.holomorphicSectionOfChart]\label{def:AlgebraicGeometry.ComplexPoint.holomorphicSectionOfChart}
\lean{AlgebraicGeometry.ComplexPoint.holomorphicSectionOfChart}\leanok
\uses{def:AlgebraicGeometry.ComplexAlgHom.isStandardEtaleBaseAlgHom, def:AlgebraicGeometry.ComplexAlgHom.standardEtaleBivariateEvaluation, def:AlgebraicGeometry.ComplexPoint.OpenHolomorphicFunctions, def:AlgebraicGeometry.ComplexPoint.affineSpecComplexPointMap, def:AlgebraicGeometry.ComplexPoint.chartSectionDomain, def:AlgebraicGeometry.ComplexPoint.polynomialSpecToAffineSpacePointMap, def:AlgebraicGeometry.LocalEtaleCoordinates.algebraicCoordinates, def:AlgebraicGeometry.LocalEtaleCoordinates.ambientAnalyticCoordinates, def:AlgebraicGeometry.LocalEtaleCoordinates.ambientCoordinateSection}
An analytic scalar function in one chart, regarded as a holomorphic section on the chart
source.

\noindent\emph{Lean:} \texttt{def AlgebraicGeometry.ComplexPoint.holomorphicSectionOfChart}, \href{https://github.com/Paul-Lez/HodgeConjecture/blob/main/HodgeConjecture/Lemmas/AlgebraicGeometry/HolomorphicPoincare.lean#L270}{\texttt{HodgeConjecture/Lemmas/AlgebraicGeometry/HolomorphicPoincare.lean:270}}.
\end{definition}

\begin{definition}[AlgebraicGeometry.ComplexPoint.chartCoordinateSection]\label{def:AlgebraicGeometry.ComplexPoint.chartCoordinateSection}
\lean{AlgebraicGeometry.ComplexPoint.chartCoordinateSection}\leanok
\uses{def:AlgebraicGeometry.ComplexPoint.holomorphicSectionOfChart}
The \texttt{i}-th fixed-chart coordinate as a holomorphic section.

\noindent\emph{Lean:} \texttt{def AlgebraicGeometry.ComplexPoint.chartCoordinateSection}, \href{https://github.com/Paul-Lez/HodgeConjecture/blob/main/HodgeConjecture/Lemmas/AlgebraicGeometry/HolomorphicPoincare.lean#L315}{\texttt{HodgeConjecture/Lemmas/AlgebraicGeometry/HolomorphicPoincare.lean:315}}.
\end{definition}

\begin{definition}[AlgebraicGeometry.ComplexPoint.formOfAnalyticField]\label{def:AlgebraicGeometry.ComplexPoint.formOfAnalyticField}
\lean{AlgebraicGeometry.ComplexPoint.formOfAnalyticField}\leanok
\uses{def:Algebra.DeRham.mk, def:AlgebraicGeometry.ComplexPoint.chartCoordinateSection, def:ContinuousAlternatingMap.alternatingFormEvaluation, def:ContinuousAlternatingMap.coordinateCoefficient}
A finite differential form whose fixed-chart evaluation is a given analytic
alternating-form field.

\noindent\emph{Lean:} \texttt{def AlgebraicGeometry.ComplexPoint.formOfAnalyticField}, \href{https://github.com/Paul-Lez/HodgeConjecture/blob/main/HodgeConjecture/Lemmas/AlgebraicGeometry/HolomorphicPoincare.lean#L352}{\texttt{HodgeConjecture/Lemmas/AlgebraicGeometry/HolomorphicPoincare.lean:352}}.
\end{definition}

\section{Constructed local relative homology for smooth closed supports}\label{sec:HodgeConjecture.Lemmas.AlgebraicGeometry.SmoothClosedSupportLocalHomology}
\noindent\emph{File:} \href{https://github.com/Paul-Lez/HodgeConjecture/blob/main/HodgeConjecture/Lemmas/AlgebraicGeometry/SmoothClosedSupportLocalHomology.lean}{\texttt{HodgeConjecture/Lemmas/AlgebraicGeometry/SmoothClosedSupportLocalHomology.lean}}.

Every prescribed open neighborhood of a point on a smooth closed complex subvariety contains
an explicitly constructed smaller open neighborhood whose support-complement pair has rational
relative homology and cohomology only in degree twice the complex codimension. The comparison
and the exactly normalized normal class come from actual flattening, radial compression, and
tangent contraction. No purity, derivative, or comparison equivalence is supplied.

This is a cofinal local singular calculation. The identification with derived sheaf sections
with support, and gluing its normalizations across overlapping charts, are separate theorems.

\begin{definition}[AlgebraicGeometry.ComplexPoint.smoothClosedSupportRestrictionChart]\label{def:AlgebraicGeometry.ComplexPoint.smoothClosedSupportRestrictionChart}
\lean{AlgebraicGeometry.ComplexPoint.smoothClosedSupportRestrictionChart}\leanok
\uses{def:AlgebraicGeometry.ComplexPoint.closedImmersionStandardFlatteningChart}
Restrict the constructed flattening chart by the prescribed ambient open.

\noindent\emph{Lean:} \texttt{def AlgebraicGeometry.ComplexPoint.smoothClosedSupportRestrictionChart}, \href{https://github.com/Paul-Lez/HodgeConjecture/blob/main/HodgeConjecture/Lemmas/AlgebraicGeometry/SmoothClosedSupportLocalHomology.lean#L37}{\texttt{HodgeConjecture/Lemmas/AlgebraicGeometry/SmoothClosedSupportLocalHomology.lean:37}}.
\end{definition}

\begin{definition}[AlgebraicGeometry.ComplexPoint.smoothClosedSupportNeighborhood]\label{def:AlgebraicGeometry.ComplexPoint.smoothClosedSupportNeighborhood}
\lean{AlgebraicGeometry.ComplexPoint.smoothClosedSupportNeighborhood}\leanok
\uses{def:AlgebraicGeometry.ComplexPoint.smoothClosedSupportRestrictionChart, def:AlgebraicTopology.Singular.flattenedSupportNeighborhood}
An actual small open support-model neighborhood inside the prescribed open.

\noindent\emph{Lean:} \texttt{def AlgebraicGeometry.ComplexPoint.smoothClosedSupportNeighborhood}, \href{https://github.com/Paul-Lez/HodgeConjecture/blob/main/HodgeConjecture/Lemmas/AlgebraicGeometry/SmoothClosedSupportLocalHomology.lean#L60}{\texttt{HodgeConjecture/Lemmas/AlgebraicGeometry/SmoothClosedSupportLocalHomology.lean:60}}.
\end{definition}

\begin{definition}[AlgebraicGeometry.ComplexPoint.smoothClosedSupportNeighborhoodPair]\label{def:AlgebraicGeometry.ComplexPoint.smoothClosedSupportNeighborhoodPair}
\lean{AlgebraicGeometry.ComplexPoint.smoothClosedSupportNeighborhoodPair}\leanok
\uses{def:AlgebraicGeometry.ComplexPoint.smoothClosedSupportNeighborhood, def:AlgebraicTopology.Singular.neighborhoodSupportComplementPair}
The local support-complement pair, with the actual image as support.

\noindent\emph{Lean:} \texttt{abbrev AlgebraicGeometry.ComplexPoint.smoothClosedSupportNeighborhoodPair}, \href{https://github.com/Paul-Lez/HodgeConjecture/blob/main/HodgeConjecture/Lemmas/AlgebraicGeometry/SmoothClosedSupportLocalHomology.lean#L79}{\texttt{HodgeConjecture/Lemmas/AlgebraicGeometry/SmoothClosedSupportLocalHomology.lean:79}}.
\end{definition}

\begin{definition}[AlgebraicGeometry.ComplexPoint.smoothClosedSupportNeighborhoodPairIso]\label{def:AlgebraicGeometry.ComplexPoint.smoothClosedSupportNeighborhoodPairIso}
\lean{AlgebraicGeometry.ComplexPoint.smoothClosedSupportNeighborhoodPairIso}\leanok
\uses{def:AlgebraicGeometry.ComplexPoint.smoothClosedSupportNeighborhoodPair, def:AlgebraicTopology.Singular.flattenedSupportPairIso}
The actual pair homeomorphism to the product normal model.

\noindent\emph{Lean:} \texttt{def AlgebraicGeometry.ComplexPoint.smoothClosedSupportNeighborhoodPairIso}, \href{https://github.com/Paul-Lez/HodgeConjecture/blob/main/HodgeConjecture/Lemmas/AlgebraicGeometry/SmoothClosedSupportLocalHomology.lean#L85}{\texttt{HodgeConjecture/Lemmas/AlgebraicGeometry/SmoothClosedSupportLocalHomology.lean:85}}.
\end{definition}

\begin{definition}[AlgebraicGeometry.ComplexPoint.smoothClosedSupportRelativeHomologyIso]\label{def:AlgebraicGeometry.ComplexPoint.smoothClosedSupportRelativeHomologyIso}
\lean{AlgebraicGeometry.ComplexPoint.smoothClosedSupportRelativeHomologyIso}\leanok
\uses{def:AlgebraicGeometry.ComplexPoint.smoothClosedSupportNeighborhoodPairIso, def:AlgebraicTopology.Singular.normalSliceRelativeHomologyIso}
All-degree local relative homology, computed from the actual pair maps.

\noindent\emph{Lean:} \texttt{def AlgebraicGeometry.ComplexPoint.smoothClosedSupportRelativeHomologyIso}, \href{https://github.com/Paul-Lez/HodgeConjecture/blob/main/HodgeConjecture/Lemmas/AlgebraicGeometry/SmoothClosedSupportLocalHomology.lean#L97}{\texttt{HodgeConjecture/Lemmas/AlgebraicGeometry/SmoothClosedSupportLocalHomology.lean:97}}.
\end{definition}

\section{Local support vanishing on canonical singular-filtration layers}\label{sec:HodgeConjecture.Lemmas.AlgebraicGeometry.SingularFiltrationLocalSupportVanishing}
\noindent\emph{File:} \href{https://github.com/Paul-Lez/HodgeConjecture/blob/main/HodgeConjecture/Lemmas/AlgebraicGeometry/SingularFiltrationLocalSupportVanishing.lean}{\texttt{HodgeConjecture/Lemmas/AlgebraicGeometry/SingularFiltrationLocalSupportVanishing.lean}}.

Every point of a smooth layer has a fixed-dimensional affine source neighborhood, of
dimension strictly less than that of the original cycle component. Deleting the image of
the discarded source complement makes that neighborhood closed in a smaller ambient open,
and normal neighborhoods there transport to the original projective ambient space and its
fixed supported injective resolution. The local lower vanishing follows from those maps
and the dimension bounds.

\begin{definition}[AlgebraicGeometry.ComplexPoint.singularFiltrationLocalSupportVanishingAnalyticTopology]\label{def:AlgebraicGeometry.ComplexPoint.singularFiltrationLocalSupportVanishingAnalyticTopology}
\lean{AlgebraicGeometry.ComplexPoint.singularFiltrationLocalSupportVanishingAnalyticTopology}\leanok
\uses{def:AlgebraicGeometry.ComplexPoint, def:AlgebraicGeometry.Point.chartTopology}
\noindent\emph{Lean:} \texttt{def AlgebraicGeometry.ComplexPoint.singularFiltrationLocalSupportVanishingAnalyticTopology}, \href{https://github.com/Paul-Lez/HodgeConjecture/blob/main/HodgeConjecture/Lemmas/AlgebraicGeometry/SingularFiltrationLocalSupportVanishing.lean#L35}{\texttt{HodgeConjecture/Lemmas/AlgebraicGeometry/SingularFiltrationLocalSupportVanishing.lean:35}}.
\end{definition}

\chapter{Lemmas: AlgebraicTopology}

\section{Local fundamental classes from complex charts}\label{sec:HodgeConjecture.Lemmas.AlgebraicTopology.ChartLocalFundamentalClass}
\noindent\emph{File:} \href{https://github.com/Paul-Lez/HodgeConjecture/blob/main/HodgeConjecture/Lemmas/AlgebraicTopology/ChartLocalFundamentalClass.lean}{\texttt{HodgeConjecture/Lemmas/AlgebraicTopology/ChartLocalFundamentalClass.lean}}.

This file transports the explicit standard class of \texttt{ℂ\^{}d} to a point in a topological space with
a complex coordinate chart. An open ball around the coordinate of the point is chosen inside the
chart target. Mathlib's radial homeomorphism compresses all of \texttt{ℂ\^{}d} into that ball, after which
the inverse chart gives an open embedding into the space. The induced map of punctured pairs sends
the standard complex local class to local homology at the point.

The construction fixes both the support and the complex orientation. It requires only an actual
chart and a proof that the point lies in its source; no fundamental or orientation class is passed
as data.

\begin{definition}[AlgebraicTopology.Singular.pointComplementPair]\label{def:AlgebraicTopology.Singular.pointComplementPair}
\lean{AlgebraicTopology.Singular.pointComplementPair}\leanok
The pair \texttt{(M, M ∖ \{x\})} used for local homology at \texttt{x}.

\noindent\emph{Lean:} \texttt{abbrev AlgebraicTopology.Singular.pointComplementPair}, \href{https://github.com/Paul-Lez/HodgeConjecture/blob/main/HodgeConjecture/Lemmas/AlgebraicTopology/ChartLocalFundamentalClass.lean#L43}{\texttt{HodgeConjecture/Lemmas/AlgebraicTopology/ChartLocalFundamentalClass.lean:43}}.
\end{definition}

\section{SingularCoverSmall}\label{sec:HodgeConjecture.Lemmas.AlgebraicTopology.SingularCoverSmall}
\noindent\emph{File:} \href{https://github.com/Paul-Lez/HodgeConjecture/blob/main/HodgeConjecture/Lemmas/AlgebraicTopology/SingularCoverSmall.lean}{\texttt{HodgeConjecture/Lemmas/AlgebraicTopology/SingularCoverSmall.lean}}.

This module is ported from Paul Lezeau's corresponding file in
\texttt{sphere-six-complex} pull request \#49, under the Apache-2.0 license.

\paragraph{Small singular chains for excision}

This file builds the first chain-level layer of the singular excision argument.  Given a family
of subsets of a space, the small singular subcomplex consists of the singular simplices which
factor through one member of the family.  Its chain complex maps canonically and monomorphically
to the full singular chain complex, and every cover-member chain map factors through it.

The classical subdivision theorem says that, for an open cover, this inclusion is a chain-homotopy
equivalence.  The definitions below state that next step using mathlib's actual \texttt{HomotopyEquiv}
API and prove its full homological consequence.  Mathlib's current simplicial subdivision functor
does not yet provide a last-vertex map, a subdivision chain map, or its chain homotopy to the
identity, so that theorem cannot yet be constructed from library primitives.

\begin{definition}[AlgebraicTopology.Singular.IntegralSingularChainComplexObj]\label{def:AlgebraicTopology.Singular.IntegralSingularChainComplexObj}
\lean{AlgebraicTopology.Singular.IntegralSingularChainComplexObj}\leanok
Integral singular chains of a categorical topological space.

\noindent\emph{Lean:} \texttt{abbrev AlgebraicTopology.Singular.IntegralSingularChainComplexObj}, \href{https://github.com/Paul-Lez/HodgeConjecture/blob/main/HodgeConjecture/Lemmas/AlgebraicTopology/SingularCoverSmall.lean#L47}{\texttt{HodgeConjecture/Lemmas/AlgebraicTopology/SingularCoverSmall.lean:47}}.
\end{definition}

\begin{definition}[AlgebraicTopology.Singular.integralSingularChainMapObj]\label{def:AlgebraicTopology.Singular.integralSingularChainMapObj}
\lean{AlgebraicTopology.Singular.integralSingularChainMapObj}\leanok
\uses{def:AlgebraicTopology.Singular.IntegralSingularChainComplexObj}
The singular-chain map induced by a continuous map.

\noindent\emph{Lean:} \texttt{def AlgebraicTopology.Singular.integralSingularChainMapObj}, \href{https://github.com/Paul-Lez/HodgeConjecture/blob/main/HodgeConjecture/Lemmas/AlgebraicTopology/SingularCoverSmall.lean#L52}{\texttt{HodgeConjecture/Lemmas/AlgebraicTopology/SingularCoverSmall.lean:52}}.
\end{definition}

\begin{definition}[AlgebraicTopology.Singular.topologicalSubsetInclusion]\label{def:AlgebraicTopology.Singular.topologicalSubsetInclusion}
\lean{AlgebraicTopology.Singular.topologicalSubsetInclusion}\leanok
The categorical inclusion of a topological subspace.

\noindent\emph{Lean:} \texttt{def AlgebraicTopology.Singular.topologicalSubsetInclusion}, \href{https://github.com/Paul-Lez/HodgeConjecture/blob/main/HodgeConjecture/Lemmas/AlgebraicTopology/SingularCoverSmall.lean#L61}{\texttt{HodgeConjecture/Lemmas/AlgebraicTopology/SingularCoverSmall.lean:61}}.
\end{definition}

\begin{definition}[AlgebraicTopology.Singular.coverSmallSingularSubcomplex]\label{def:AlgebraicTopology.Singular.coverSmallSingularSubcomplex}
\lean{AlgebraicTopology.Singular.coverSmallSingularSubcomplex}\leanok
\uses{def:AlgebraicTopology.Singular.topologicalSubsetInclusion}
Singular simplices which factor through one member of \texttt{U}.

\noindent\emph{Lean:} \texttt{def AlgebraicTopology.Singular.coverSmallSingularSubcomplex}, \href{https://github.com/Paul-Lez/HodgeConjecture/blob/main/HodgeConjecture/Lemmas/AlgebraicTopology/SingularCoverSmall.lean#L66}{\texttt{HodgeConjecture/Lemmas/AlgebraicTopology/SingularCoverSmall.lean:66}}.
\end{definition}

\begin{definition}[AlgebraicTopology.Singular.CoverSmallIntegralSingularChainComplex]\label{def:AlgebraicTopology.Singular.CoverSmallIntegralSingularChainComplex}
\lean{AlgebraicTopology.Singular.CoverSmallIntegralSingularChainComplex}\leanok
\uses{def:AlgebraicTopology.Singular.coverSmallSingularSubcomplex}
Integral chains on the cover-small singular simplicial set.

\noindent\emph{Lean:} \texttt{abbrev AlgebraicTopology.Singular.CoverSmallIntegralSingularChainComplex}, \href{https://github.com/Paul-Lez/HodgeConjecture/blob/main/HodgeConjecture/Lemmas/AlgebraicTopology/SingularCoverSmall.lean#L88}{\texttt{HodgeConjecture/Lemmas/AlgebraicTopology/SingularCoverSmall.lean:88}}.
\end{definition}

\begin{definition}[AlgebraicTopology.Singular.coverSmallIntegralSingularChainInclusion]\label{def:AlgebraicTopology.Singular.coverSmallIntegralSingularChainInclusion}
\lean{AlgebraicTopology.Singular.coverSmallIntegralSingularChainInclusion}\leanok
\uses{def:AlgebraicTopology.Singular.CoverSmallIntegralSingularChainComplex, def:AlgebraicTopology.Singular.IntegralSingularChainComplexObj}
Inclusion of cover-small integral singular chains into all integral singular chains.

\noindent\emph{Lean:} \texttt{def AlgebraicTopology.Singular.coverSmallIntegralSingularChainInclusion}, \href{https://github.com/Paul-Lez/HodgeConjecture/blob/main/HodgeConjecture/Lemmas/AlgebraicTopology/SingularCoverSmall.lean#L93}{\texttt{HodgeConjecture/Lemmas/AlgebraicTopology/SingularCoverSmall.lean:93}}.
\end{definition}

\begin{definition}[AlgebraicTopology.Singular.coverMemberToSmallSingularSet]\label{def:AlgebraicTopology.Singular.coverMemberToSmallSingularSet}
\lean{AlgebraicTopology.Singular.coverMemberToSmallSingularSet}\leanok
\uses{def:AlgebraicTopology.Singular.coverSmallSingularSubcomplex}
The singular set of each cover member factors through the small singular subcomplex.

\noindent\emph{Lean:} \texttt{def AlgebraicTopology.Singular.coverMemberToSmallSingularSet}, \href{https://github.com/Paul-Lez/HodgeConjecture/blob/main/HodgeConjecture/Lemmas/AlgebraicTopology/SingularCoverSmall.lean#L107}{\texttt{HodgeConjecture/Lemmas/AlgebraicTopology/SingularCoverSmall.lean:107}}.
\end{definition}

\begin{definition}[AlgebraicTopology.Singular.coverMemberToSmallIntegralSingularChains]\label{def:AlgebraicTopology.Singular.coverMemberToSmallIntegralSingularChains}
\lean{AlgebraicTopology.Singular.coverMemberToSmallIntegralSingularChains}\leanok
\uses{def:AlgebraicTopology.Singular.CoverSmallIntegralSingularChainComplex, def:AlgebraicTopology.Singular.IntegralSingularChainComplexObj, def:AlgebraicTopology.Singular.coverMemberToSmallSingularSet}
The chain map from a cover member into the small singular chains.

\noindent\emph{Lean:} \texttt{def AlgebraicTopology.Singular.coverMemberToSmallIntegralSingularChains}, \href{https://github.com/Paul-Lez/HodgeConjecture/blob/main/HodgeConjecture/Lemmas/AlgebraicTopology/SingularCoverSmall.lean#L121}{\texttt{HodgeConjecture/Lemmas/AlgebraicTopology/SingularCoverSmall.lean:121}}.
\end{definition}

\section{SingularSubdivision}\label{sec:HodgeConjecture.Lemmas.AlgebraicTopology.SingularSubdivision}
\noindent\emph{File:} \href{https://github.com/Paul-Lez/HodgeConjecture/blob/main/HodgeConjecture/Lemmas/AlgebraicTopology/SingularSubdivision.lean}{\texttt{HodgeConjecture/Lemmas/AlgebraicTopology/SingularSubdivision.lean}}.

This module is ported from Paul Lezeau's corresponding file in
\texttt{sphere-six-complex} pull request \#49, under the Apache-2.0 license.

\paragraph{The last-vertex map for simplicial subdivision}

Mathlib defines barycentric subdivision as a left Kan extension, but currently does not include
the classical last-vertex natural transformation.  This file begins that construction at its
finite-poset source: a nonempty finite chain is sent to its greatest vertex.  The maximum map is
monotone and natural in monotone maps between finite linear orders, so taking nerves gives the
last-vertex map on the subdivided standard simplices.

\begin{definition}[AlgebraicTopology.Singular.nonemptyFiniteChainMaximum]\label{def:AlgebraicTopology.Singular.nonemptyFiniteChainMaximum}
\lean{AlgebraicTopology.Singular.nonemptyFiniteChainMaximum}\leanok
A nonempty finite chain in a linear order is sent monotonically to its greatest element.

\noindent\emph{Lean:} \texttt{def AlgebraicTopology.Singular.nonemptyFiniteChainMaximum}, \href{https://github.com/Paul-Lez/HodgeConjecture/blob/main/HodgeConjecture/Lemmas/AlgebraicTopology/SingularSubdivision.lean#L44}{\texttt{HodgeConjecture/Lemmas/AlgebraicTopology/SingularSubdivision.lean:44}}.
\end{definition}

\begin{definition}[AlgebraicTopology.Singular.simplexSubdivisionMaximum]\label{def:AlgebraicTopology.Singular.simplexSubdivisionMaximum}
\lean{AlgebraicTopology.Singular.simplexSubdivisionMaximum}\leanok
\uses{def:AlgebraicTopology.Singular.nonemptyFiniteChainMaximum}
The maximum-vertex morphism from chains in the vertex poset of a standard simplex.

\noindent\emph{Lean:} \texttt{def AlgebraicTopology.Singular.simplexSubdivisionMaximum}, \href{https://github.com/Paul-Lez/HodgeConjecture/blob/main/HodgeConjecture/Lemmas/AlgebraicTopology/SingularSubdivision.lean#L71}{\texttt{HodgeConjecture/Lemmas/AlgebraicTopology/SingularSubdivision.lean:71}}.
\end{definition}

\begin{definition}[AlgebraicTopology.Singular.simplexSubdivisionMaximumNatTrans]\label{def:AlgebraicTopology.Singular.simplexSubdivisionMaximumNatTrans}
\lean{AlgebraicTopology.Singular.simplexSubdivisionMaximumNatTrans}\leanok
\uses{def:AlgebraicTopology.Singular.simplexSubdivisionMaximum}
Maximum vertex is natural in morphisms of the simplex category.

\noindent\emph{Lean:} \texttt{def AlgebraicTopology.Singular.simplexSubdivisionMaximumNatTrans}, \href{https://github.com/Paul-Lez/HodgeConjecture/blob/main/HodgeConjecture/Lemmas/AlgebraicTopology/SingularSubdivision.lean#L90}{\texttt{HodgeConjecture/Lemmas/AlgebraicTopology/SingularSubdivision.lean:90}}.
\end{definition}

\begin{definition}[AlgebraicTopology.Singular.simplexSubdivisionLastVertexToNerve]\label{def:AlgebraicTopology.Singular.simplexSubdivisionLastVertexToNerve}
\lean{AlgebraicTopology.Singular.simplexSubdivisionLastVertexToNerve}\leanok
\uses{def:AlgebraicTopology.Singular.simplexSubdivisionMaximumNatTrans}
Taking nerves gives the last-vertex map on the subdivision model of standard simplices.

\noindent\emph{Lean:} \texttt{def AlgebraicTopology.Singular.simplexSubdivisionLastVertexToNerve}, \href{https://github.com/Paul-Lez/HodgeConjecture/blob/main/HodgeConjecture/Lemmas/AlgebraicTopology/SingularSubdivision.lean#L107}{\texttt{HodgeConjecture/Lemmas/AlgebraicTopology/SingularSubdivision.lean:107}}.
\end{definition}

\begin{definition}[AlgebraicTopology.Singular.standardSimplexNerveIso]\label{def:AlgebraicTopology.Singular.standardSimplexNerveIso}
\lean{AlgebraicTopology.Singular.standardSimplexNerveIso}\leanok
The objectwise identification of the standard simplex with the nerve of its vertex order is
natural in the simplex category.

\noindent\emph{Lean:} \texttt{def AlgebraicTopology.Singular.standardSimplexNerveIso}, \href{https://github.com/Paul-Lez/HodgeConjecture/blob/main/HodgeConjecture/Lemmas/AlgebraicTopology/SingularSubdivision.lean#L113}{\texttt{HodgeConjecture/Lemmas/AlgebraicTopology/SingularSubdivision.lean:113}}.
\end{definition}

\begin{definition}[AlgebraicTopology.Singular.simplexSubdivisionLastVertex]\label{def:AlgebraicTopology.Singular.simplexSubdivisionLastVertex}
\lean{AlgebraicTopology.Singular.simplexSubdivisionLastVertex}\leanok
\uses{def:AlgebraicTopology.Singular.simplexSubdivisionLastVertexToNerve, def:AlgebraicTopology.Singular.standardSimplexNerveIso}
The last-vertex map from the subdivision model of standard simplices to standard simplices.

\noindent\emph{Lean:} \texttt{def AlgebraicTopology.Singular.simplexSubdivisionLastVertex}, \href{https://github.com/Paul-Lez/HodgeConjecture/blob/main/HodgeConjecture/Lemmas/AlgebraicTopology/SingularSubdivision.lean#L120}{\texttt{HodgeConjecture/Lemmas/AlgebraicTopology/SingularSubdivision.lean:120}}.
\end{definition}

\begin{definition}[AlgebraicTopology.Singular.simplexSubdivisionLastVertexExtension]\label{def:AlgebraicTopology.Singular.simplexSubdivisionLastVertexExtension}
\lean{AlgebraicTopology.Singular.simplexSubdivisionLastVertexExtension}\leanok
\uses{def:AlgebraicTopology.Singular.simplexSubdivisionLastVertex}
The last-vertex map, viewed with the codomain required by the left Kan extension universal
property.

\noindent\emph{Lean:} \texttt{def AlgebraicTopology.Singular.simplexSubdivisionLastVertexExtension}, \href{https://github.com/Paul-Lez/HodgeConjecture/blob/main/HodgeConjecture/Lemmas/AlgebraicTopology/SingularSubdivision.lean#L125}{\texttt{HodgeConjecture/Lemmas/AlgebraicTopology/SingularSubdivision.lean:125}}.
\end{definition}

\begin{definition}[AlgebraicTopology.Singular.subdivisionLastVertex]\label{def:AlgebraicTopology.Singular.subdivisionLastVertex}
\lean{AlgebraicTopology.Singular.subdivisionLastVertex}\leanok
\uses{def:AlgebraicTopology.Singular.simplexSubdivisionLastVertexExtension}
The global last-vertex natural transformation \texttt{sd X ⟶ X}, obtained from the defining left
Kan extension universal property of simplicial subdivision.

\noindent\emph{Lean:} \texttt{def AlgebraicTopology.Singular.subdivisionLastVertex}, \href{https://github.com/Paul-Lez/HodgeConjecture/blob/main/HodgeConjecture/Lemmas/AlgebraicTopology/SingularSubdivision.lean#L133}{\texttt{HodgeConjecture/Lemmas/AlgebraicTopology/SingularSubdivision.lean:133}}.
\end{definition}

\begin{definition}[AlgebraicTopology.Singular.subdivisionLastVertexChainMap]\label{def:AlgebraicTopology.Singular.subdivisionLastVertexChainMap}
\lean{AlgebraicTopology.Singular.subdivisionLastVertexChainMap}\leanok
\uses{def:AlgebraicTopology.Singular.subdivisionLastVertex}
The chain map induced by the last-vertex map of a simplicial set.

\noindent\emph{Lean:} \texttt{def AlgebraicTopology.Singular.subdivisionLastVertexChainMap}, \href{https://github.com/Paul-Lez/HodgeConjecture/blob/main/HodgeConjecture/Lemmas/AlgebraicTopology/SingularSubdivision.lean#L150}{\texttt{HodgeConjecture/Lemmas/AlgebraicTopology/SingularSubdivision.lean:150}}.
\end{definition}

\section{SingularBarycentricChains}\label{sec:HodgeConjecture.Lemmas.AlgebraicTopology.SingularBarycentricChains}
\noindent\emph{File:} \href{https://github.com/Paul-Lez/HodgeConjecture/blob/main/HodgeConjecture/Lemmas/AlgebraicTopology/SingularBarycentricChains.lean}{\texttt{HodgeConjecture/Lemmas/AlgebraicTopology/SingularBarycentricChains.lean}}.

This module is ported from Paul Lezeau's corresponding file in
\texttt{sphere-six-complex} pull request \#49, under the Apache-2.0 license.

\paragraph{Low-dimensional barycentric fundamental chains}

This file constructs the first actual components of the barycentric subdivision chain operator.
The construction starts from the vertex of \texttt{sd Δ[0]} represented by the singleton chain and
uses Yoneda naturality to associate a subdivided vertex to every vertex of a simplicial set.

\begin{definition}[AlgebraicTopology.Singular.nonemptyFiniteChainSingleton]\label{def:AlgebraicTopology.Singular.nonemptyFiniteChainSingleton}
\lean{AlgebraicTopology.Singular.nonemptyFiniteChainSingleton}\leanok
A singleton, regarded as a nonempty finite chain.

\noindent\emph{Lean:} \texttt{def AlgebraicTopology.Singular.nonemptyFiniteChainSingleton}, \href{https://github.com/Paul-Lez/HodgeConjecture/blob/main/HodgeConjecture/Lemmas/AlgebraicTopology/SingularBarycentricChains.lean#L39}{\texttt{HodgeConjecture/Lemmas/AlgebraicTopology/SingularBarycentricChains.lean:39}}.
\end{definition}

\begin{definition}[AlgebraicTopology.Singular.nonemptyFiniteChainOfFinset]\label{def:AlgebraicTopology.Singular.nonemptyFiniteChainOfFinset}
\lean{AlgebraicTopology.Singular.nonemptyFiniteChainOfFinset}\leanok
A specified nonempty finset in a linear order, regarded as a finite chain.

This constructor is shared by the all-degree barycentric development; unlike the removed
degree-one scaffolding, it is part of that file's live API.

\noindent\emph{Lean:} \texttt{def AlgebraicTopology.Singular.nonemptyFiniteChainOfFinset}, \href{https://github.com/Paul-Lez/HodgeConjecture/blob/main/HodgeConjecture/Lemmas/AlgebraicTopology/SingularBarycentricChains.lean#L51}{\texttt{HodgeConjecture/Lemmas/AlgebraicTopology/SingularBarycentricChains.lean:51}}.
\end{definition}

\begin{definition}[AlgebraicTopology.Singular.subdividedZeroSimplexVertex]\label{def:AlgebraicTopology.Singular.subdividedZeroSimplexVertex}
\lean{AlgebraicTopology.Singular.subdividedZeroSimplexVertex}\leanok
\uses{def:AlgebraicTopology.Singular.nonemptyFiniteChainSingleton}
The vertex of the subdivision model of \texttt{Δ[0]} represented by its singleton vertex chain.

\noindent\emph{Lean:} \texttt{def AlgebraicTopology.Singular.subdividedZeroSimplexVertex}, \href{https://github.com/Paul-Lez/HodgeConjecture/blob/main/HodgeConjecture/Lemmas/AlgebraicTopology/SingularBarycentricChains.lean#L62}{\texttt{HodgeConjecture/Lemmas/AlgebraicTopology/SingularBarycentricChains.lean:62}}.
\end{definition}

\begin{definition}[AlgebraicTopology.Singular.subdividedStandardZeroVertex]\label{def:AlgebraicTopology.Singular.subdividedStandardZeroVertex}
\lean{AlgebraicTopology.Singular.subdividedStandardZeroVertex}\leanok
\uses{def:AlgebraicTopology.Singular.subdividedZeroSimplexVertex}
The corresponding vertex of Mathlib's left-Kan-extension model \texttt{sd Δ[0]}.

\noindent\emph{Lean:} \texttt{def AlgebraicTopology.Singular.subdividedStandardZeroVertex}, \href{https://github.com/Paul-Lez/HodgeConjecture/blob/main/HodgeConjecture/Lemmas/AlgebraicTopology/SingularBarycentricChains.lean#L68}{\texttt{HodgeConjecture/Lemmas/AlgebraicTopology/SingularBarycentricChains.lean:68}}.
\end{definition}

\begin{definition}[AlgebraicTopology.Singular.barycentricSubdivisionVertex]\label{def:AlgebraicTopology.Singular.barycentricSubdivisionVertex}
\lean{AlgebraicTopology.Singular.barycentricSubdivisionVertex}\leanok
\uses{def:AlgebraicTopology.Singular.subdividedStandardZeroVertex}
Every vertex of a simplicial set determines its canonical vertex after subdivision.

\noindent\emph{Lean:} \texttt{def AlgebraicTopology.Singular.barycentricSubdivisionVertex}, \href{https://github.com/Paul-Lez/HodgeConjecture/blob/main/HodgeConjecture/Lemmas/AlgebraicTopology/SingularBarycentricChains.lean#L75}{\texttt{HodgeConjecture/Lemmas/AlgebraicTopology/SingularBarycentricChains.lean:75}}.
\end{definition}

\begin{definition}[AlgebraicTopology.Singular.barycentricSubdivisionChainMapZero]\label{def:AlgebraicTopology.Singular.barycentricSubdivisionChainMapZero}
\lean{AlgebraicTopology.Singular.barycentricSubdivisionChainMapZero}\leanok
\uses{def:AlgebraicTopology.Singular.barycentricSubdivisionVertex}
The degree-zero component of barycentric subdivision on integral simplicial chains.

\noindent\emph{Lean:} \texttt{def AlgebraicTopology.Singular.barycentricSubdivisionChainMapZero}, \href{https://github.com/Paul-Lez/HodgeConjecture/blob/main/HodgeConjecture/Lemmas/AlgebraicTopology/SingularBarycentricChains.lean#L123}{\texttt{HodgeConjecture/Lemmas/AlgebraicTopology/SingularBarycentricChains.lean:123}}.
\end{definition}

\section{SingularBarycentricAllDegrees}\label{sec:HodgeConjecture.Lemmas.AlgebraicTopology.SingularBarycentricAllDegrees}
\noindent\emph{File:} \href{https://github.com/Paul-Lez/HodgeConjecture/blob/main/HodgeConjecture/Lemmas/AlgebraicTopology/SingularBarycentricAllDegrees.lean}{\texttt{HodgeConjecture/Lemmas/AlgebraicTopology/SingularBarycentricAllDegrees.lean}}.

This module is ported from Paul Lezeau's corresponding file in
\texttt{sphere-six-complex} pull request \#49, under the Apache-2.0 license.

\paragraph{Barycentric fundamental chains in all degrees}

For each permutation of the vertices of the standard \texttt{n}-simplex, its successive nonempty
prefixes form a maximal flag in the poset of nonempty subsets.  The signed sum of these flags is
the oriented barycentric fundamental chain.  This file constructs that chain in Mathlib's actual
\texttt{SimplexCategory.sd} nerve model and transports it naturally to every simplex of every simplicial
set.

\begin{definition}[AlgebraicTopology.Singular.permutationPrefixFinset]\label{def:AlgebraicTopology.Singular.permutationPrefixFinset}
\lean{AlgebraicTopology.Singular.permutationPrefixFinset}\leanok
The set of the first \texttt{k+1} vertices in the ordering specified by \texttt{σ}.

\noindent\emph{Lean:} \texttt{def AlgebraicTopology.Singular.permutationPrefixFinset}, \href{https://github.com/Paul-Lez/HodgeConjecture/blob/main/HodgeConjecture/Lemmas/AlgebraicTopology/SingularBarycentricAllDegrees.lean#L42}{\texttt{HodgeConjecture/Lemmas/AlgebraicTopology/SingularBarycentricAllDegrees.lean:42}}.
\end{definition}

\begin{definition}[AlgebraicTopology.Singular.permutationPrefixChain]\label{def:AlgebraicTopology.Singular.permutationPrefixChain}
\lean{AlgebraicTopology.Singular.permutationPrefixChain}\leanok
\uses{def:AlgebraicTopology.Singular.nonemptyFiniteChainOfFinset, def:AlgebraicTopology.Singular.permutationPrefixFinset}
A permutation prefix, regarded as a nonempty finite chain in the linearly ordered vertex
set.

\noindent\emph{Lean:} \texttt{def AlgebraicTopology.Singular.permutationPrefixChain}, \href{https://github.com/Paul-Lez/HodgeConjecture/blob/main/HodgeConjecture/Lemmas/AlgebraicTopology/SingularBarycentricAllDegrees.lean#L64}{\texttt{HodgeConjecture/Lemmas/AlgebraicTopology/SingularBarycentricAllDegrees.lean:64}}.
\end{definition}

\begin{definition}[AlgebraicTopology.Singular.permutationMaximalFlagOrderHom]\label{def:AlgebraicTopology.Singular.permutationMaximalFlagOrderHom}
\lean{AlgebraicTopology.Singular.permutationMaximalFlagOrderHom}\leanok
\uses{def:AlgebraicTopology.Singular.permutationPrefixChain}
The monotone maximal flag of prefix chains associated to a permutation.

\noindent\emph{Lean:} \texttt{def AlgebraicTopology.Singular.permutationMaximalFlagOrderHom}, \href{https://github.com/Paul-Lez/HodgeConjecture/blob/main/HodgeConjecture/Lemmas/AlgebraicTopology/SingularBarycentricAllDegrees.lean#L72}{\texttt{HodgeConjecture/Lemmas/AlgebraicTopology/SingularBarycentricAllDegrees.lean:72}}.
\end{definition}

\begin{definition}[AlgebraicTopology.Singular.permutationMaximalFlagSimplex]\label{def:AlgebraicTopology.Singular.permutationMaximalFlagSimplex}
\lean{AlgebraicTopology.Singular.permutationMaximalFlagSimplex}\leanok
\uses{def:AlgebraicTopology.Singular.permutationMaximalFlagOrderHom}
The \texttt{n}-simplex of the subdivision model corresponding to a maximal permutation flag.

\noindent\emph{Lean:} \texttt{def AlgebraicTopology.Singular.permutationMaximalFlagSimplex}, \href{https://github.com/Paul-Lez/HodgeConjecture/blob/main/HodgeConjecture/Lemmas/AlgebraicTopology/SingularBarycentricAllDegrees.lean#L79}{\texttt{HodgeConjecture/Lemmas/AlgebraicTopology/SingularBarycentricAllDegrees.lean:79}}.
\end{definition}

\begin{definition}[AlgebraicTopology.Singular.permutationSignInteger]\label{def:AlgebraicTopology.Singular.permutationSignInteger}
\lean{AlgebraicTopology.Singular.permutationSignInteger}\leanok
The sign of a permutation, as an integer coefficient.

\noindent\emph{Lean:} \texttt{def AlgebraicTopology.Singular.permutationSignInteger}, \href{https://github.com/Paul-Lez/HodgeConjecture/blob/main/HodgeConjecture/Lemmas/AlgebraicTopology/SingularBarycentricAllDegrees.lean#L86}{\texttt{HodgeConjecture/Lemmas/AlgebraicTopology/SingularBarycentricAllDegrees.lean:86}}.
\end{definition}

\begin{definition}[AlgebraicTopology.Singular.subdividedSimplexFundamentalChain]\label{def:AlgebraicTopology.Singular.subdividedSimplexFundamentalChain}
\lean{AlgebraicTopology.Singular.subdividedSimplexFundamentalChain}\leanok
\uses{def:AlgebraicTopology.Singular.permutationMaximalFlagSimplex, def:AlgebraicTopology.Singular.permutationSignInteger}
The oriented barycentric fundamental \texttt{n}-chain in the explicit subdivision model of
\texttt{Δ[n]}.

\noindent\emph{Lean:} \texttt{def AlgebraicTopology.Singular.subdividedSimplexFundamentalChain}, \href{https://github.com/Paul-Lez/HodgeConjecture/blob/main/HodgeConjecture/Lemmas/AlgebraicTopology/SingularBarycentricAllDegrees.lean#L173}{\texttt{HodgeConjecture/Lemmas/AlgebraicTopology/SingularBarycentricAllDegrees.lean:173}}.
\end{definition}

\begin{definition}[AlgebraicTopology.Singular.subdividedStandardSimplexFundamentalChain]\label{def:AlgebraicTopology.Singular.subdividedStandardSimplexFundamentalChain}
\lean{AlgebraicTopology.Singular.subdividedStandardSimplexFundamentalChain}\leanok
\uses{def:AlgebraicTopology.Singular.subdividedSimplexFundamentalChain}
The all-degree fundamental chain transported to Mathlib's left-Kan-extension model
\texttt{sd Δ[n]}.

\noindent\emph{Lean:} \texttt{def AlgebraicTopology.Singular.subdividedStandardSimplexFundamentalChain}, \href{https://github.com/Paul-Lez/HodgeConjecture/blob/main/HodgeConjecture/Lemmas/AlgebraicTopology/SingularBarycentricAllDegrees.lean#L184}{\texttt{HodgeConjecture/Lemmas/AlgebraicTopology/SingularBarycentricAllDegrees.lean:184}}.
\end{definition}

\begin{definition}[AlgebraicTopology.Singular.barycentricSubdivisionSimplexChain]\label{def:AlgebraicTopology.Singular.barycentricSubdivisionSimplexChain}
\lean{AlgebraicTopology.Singular.barycentricSubdivisionSimplexChain}\leanok
\uses{def:AlgebraicTopology.Singular.subdividedStandardSimplexFundamentalChain}
The barycentric fundamental chain associated to an arbitrary \texttt{n}-simplex.

\noindent\emph{Lean:} \texttt{def AlgebraicTopology.Singular.barycentricSubdivisionSimplexChain}, \href{https://github.com/Paul-Lez/HodgeConjecture/blob/main/HodgeConjecture/Lemmas/AlgebraicTopology/SingularBarycentricAllDegrees.lean#L194}{\texttt{HodgeConjecture/Lemmas/AlgebraicTopology/SingularBarycentricAllDegrees.lean:194}}.
\end{definition}

\begin{definition}[AlgebraicTopology.Singular.barycentricSubdivisionComponent]\label{def:AlgebraicTopology.Singular.barycentricSubdivisionComponent}
\lean{AlgebraicTopology.Singular.barycentricSubdivisionComponent}\leanok
\uses{def:AlgebraicTopology.Singular.barycentricSubdivisionSimplexChain}
The degree-\texttt{n} component of barycentric subdivision, defined on every simplex by its signed
maximal-flag fundamental chain.

\noindent\emph{Lean:} \texttt{def AlgebraicTopology.Singular.barycentricSubdivisionComponent}, \href{https://github.com/Paul-Lez/HodgeConjecture/blob/main/HodgeConjecture/Lemmas/AlgebraicTopology/SingularBarycentricAllDegrees.lean#L204}{\texttt{HodgeConjecture/Lemmas/AlgebraicTopology/SingularBarycentricAllDegrees.lean:204}}.
\end{definition}

\begin{definition}[AlgebraicTopology.Singular.subdividedSimplexAlternatingFaceChain]\label{def:AlgebraicTopology.Singular.subdividedSimplexAlternatingFaceChain}
\lean{AlgebraicTopology.Singular.subdividedSimplexAlternatingFaceChain}\leanok
\uses{def:AlgebraicTopology.Singular.subdividedSimplexFundamentalChain}
The alternating sum of the subdivided fundamental chains of the codimension-one faces of
the standard \texttt{(n+1)}-simplex, in the explicit nerve model.

\noindent\emph{Lean:} \texttt{def AlgebraicTopology.Singular.subdividedSimplexAlternatingFaceChain}, \href{https://github.com/Paul-Lez/HodgeConjecture/blob/main/HodgeConjecture/Lemmas/AlgebraicTopology/SingularBarycentricAllDegrees.lean#L251}{\texttt{HodgeConjecture/Lemmas/AlgebraicTopology/SingularBarycentricAllDegrees.lean:251}}.
\end{definition}

\begin{definition}[AlgebraicTopology.Singular.subdividedSimplexInteriorFaceSum]\label{def:AlgebraicTopology.Singular.subdividedSimplexInteriorFaceSum}
\lean{AlgebraicTopology.Singular.subdividedSimplexInteriorFaceSum}\leanok
\uses{def:AlgebraicTopology.Singular.permutationMaximalFlagSimplex, def:AlgebraicTopology.Singular.permutationSignInteger}
The contribution to the boundary in which the prefix at the interior position \texttt{r} is
deleted.

\noindent\emph{Lean:} \texttt{def AlgebraicTopology.Singular.subdividedSimplexInteriorFaceSum}, \href{https://github.com/Paul-Lez/HodgeConjecture/blob/main/HodgeConjecture/Lemmas/AlgebraicTopology/SingularBarycentricAllDegrees.lean#L263}{\texttt{HodgeConjecture/Lemmas/AlgebraicTopology/SingularBarycentricAllDegrees.lean:263}}.
\end{definition}

\begin{definition}[AlgebraicTopology.Singular.subdividedSimplexOuterFaceSum]\label{def:AlgebraicTopology.Singular.subdividedSimplexOuterFaceSum}
\lean{AlgebraicTopology.Singular.subdividedSimplexOuterFaceSum}\leanok
\uses{def:AlgebraicTopology.Singular.permutationMaximalFlagSimplex, def:AlgebraicTopology.Singular.permutationSignInteger}
The surviving last-prefix faces in the boundary of the signed maximal-flag chain.

\noindent\emph{Lean:} \texttt{def AlgebraicTopology.Singular.subdividedSimplexOuterFaceSum}, \href{https://github.com/Paul-Lez/HodgeConjecture/blob/main/HodgeConjecture/Lemmas/AlgebraicTopology/SingularBarycentricAllDegrees.lean#L292}{\texttt{HodgeConjecture/Lemmas/AlgebraicTopology/SingularBarycentricAllDegrees.lean:292}}.
\end{definition}

\begin{definition}[AlgebraicTopology.Singular.subdividedStandardSimplexAlternatingFaceChain]\label{def:AlgebraicTopology.Singular.subdividedStandardSimplexAlternatingFaceChain}
\lean{AlgebraicTopology.Singular.subdividedStandardSimplexAlternatingFaceChain}\leanok
\uses{def:AlgebraicTopology.Singular.subdividedStandardSimplexFundamentalChain}
The alternating face chain after transport to \texttt{sd Δ[n+1]}.

\noindent\emph{Lean:} \texttt{def AlgebraicTopology.Singular.subdividedStandardSimplexAlternatingFaceChain}, \href{https://github.com/Paul-Lez/HodgeConjecture/blob/main/HodgeConjecture/Lemmas/AlgebraicTopology/SingularBarycentricAllDegrees.lean#L400}{\texttt{HodgeConjecture/Lemmas/AlgebraicTopology/SingularBarycentricAllDegrees.lean:400}}.
\end{definition}

\begin{definition}[AlgebraicTopology.Singular.barycentricSubdivisionAlternatingFaceChain]\label{def:AlgebraicTopology.Singular.barycentricSubdivisionAlternatingFaceChain}
\lean{AlgebraicTopology.Singular.barycentricSubdivisionAlternatingFaceChain}\leanok
\uses{def:AlgebraicTopology.Singular.barycentricSubdivisionSimplexChain}
The alternating sum of the barycentrically subdivided faces of an arbitrary simplex.

\noindent\emph{Lean:} \texttt{def AlgebraicTopology.Singular.barycentricSubdivisionAlternatingFaceChain}, \href{https://github.com/Paul-Lez/HodgeConjecture/blob/main/HodgeConjecture/Lemmas/AlgebraicTopology/SingularBarycentricAllDegrees.lean#L431}{\texttt{HodgeConjecture/Lemmas/AlgebraicTopology/SingularBarycentricAllDegrees.lean:431}}.
\end{definition}

\section{SingularBarycentricOuterFaces}\label{sec:HodgeConjecture.Lemmas.AlgebraicTopology.SingularBarycentricOuterFaces}
\noindent\emph{File:} \href{https://github.com/Paul-Lez/HodgeConjecture/blob/main/HodgeConjecture/Lemmas/AlgebraicTopology/SingularBarycentricOuterFaces.lean}{\texttt{HodgeConjecture/Lemmas/AlgebraicTopology/SingularBarycentricOuterFaces.lean}}.

This module is ported from Paul Lezeau's corresponding file in
\texttt{sphere-six-complex} pull request \#49, under the Apache-2.0 license.

\paragraph{The outer-face identity for barycentric fundamental chains}

This file reindexes the surviving outer faces of the signed maximal-flag sum by the omitted
vertex and the induced ordering of its complement.  It completes the boundary calculation and
therefore supplies the unconditional natural barycentric subdivision chain morphism.

\begin{definition}[AlgebraicTopology.Singular.insertOmittedVertexLast]\label{def:AlgebraicTopology.Singular.insertOmittedVertexLast}
\lean{AlgebraicTopology.Singular.insertOmittedVertexLast}\leanok
Insert the omitted vertex \texttt{p} at the final position of an ordering of its complement.

\noindent\emph{Lean:} \texttt{def AlgebraicTopology.Singular.insertOmittedVertexLast}, \href{https://github.com/Paul-Lez/HodgeConjecture/blob/main/HodgeConjecture/Lemmas/AlgebraicTopology/SingularBarycentricOuterFaces.lean#L48}{\texttt{HodgeConjecture/Lemmas/AlgebraicTopology/SingularBarycentricOuterFaces.lean:48}}.
\end{definition}

\begin{definition}[AlgebraicTopology.Singular.insertOmittedVertexLastEquiv]\label{def:AlgebraicTopology.Singular.insertOmittedVertexLastEquiv}
\lean{AlgebraicTopology.Singular.insertOmittedVertexLastEquiv}\leanok
\uses{def:AlgebraicTopology.Singular.insertOmittedVertexLast}
Every permutation is uniquely an omitted final vertex together with an ordering of its
complement.

\noindent\emph{Lean:} \texttt{def AlgebraicTopology.Singular.insertOmittedVertexLastEquiv}, \href{https://github.com/Paul-Lez/HodgeConjecture/blob/main/HodgeConjecture/Lemmas/AlgebraicTopology/SingularBarycentricOuterFaces.lean#L70}{\texttt{HodgeConjecture/Lemmas/AlgebraicTopology/SingularBarycentricOuterFaces.lean:70}}.
\end{definition}

\begin{definition}[AlgebraicTopology.Singular.finCastSuccEquivNotLast]\label{def:AlgebraicTopology.Singular.finCastSuccEquivNotLast}
\lean{AlgebraicTopology.Singular.finCastSuccEquivNotLast}\leanok
\texttt{Fin.castSucc} identifies \texttt{Fin (n+1)} with all vertices except the last one.

\noindent\emph{Lean:} \texttt{def AlgebraicTopology.Singular.finCastSuccEquivNotLast}, \href{https://github.com/Paul-Lez/HodgeConjecture/blob/main/HodgeConjecture/Lemmas/AlgebraicTopology/SingularBarycentricOuterFaces.lean#L144}{\texttt{HodgeConjecture/Lemmas/AlgebraicTopology/SingularBarycentricOuterFaces.lean:144}}.
\end{definition}

\begin{definition}[AlgebraicTopology.Singular.barycentricSubdivisionChainMapCanonical]\label{def:AlgebraicTopology.Singular.barycentricSubdivisionChainMapCanonical}
\lean{AlgebraicTopology.Singular.barycentricSubdivisionChainMapCanonical}\leanok
\uses{def:AlgebraicTopology.Singular.barycentricSubdivisionAlternatingFaceChain, def:AlgebraicTopology.Singular.barycentricSubdivisionComponent, def:AlgebraicTopology.Singular.finCastSuccEquivNotLast, def:AlgebraicTopology.Singular.insertOmittedVertexLastEquiv, def:AlgebraicTopology.Singular.subdividedSimplexAlternatingFaceChain, def:AlgebraicTopology.Singular.subdividedSimplexInteriorFaceSum, def:AlgebraicTopology.Singular.subdividedSimplexOuterFaceSum, def:AlgebraicTopology.Singular.subdividedStandardSimplexAlternatingFaceChain}
The unconditional natural barycentric subdivision chain morphism.

\noindent\emph{Lean:} \texttt{def AlgebraicTopology.Singular.barycentricSubdivisionChainMapCanonical}, \href{https://github.com/Paul-Lez/HodgeConjecture/blob/main/HodgeConjecture/Lemmas/AlgebraicTopology/SingularBarycentricOuterFaces.lean#L381}{\texttt{HodgeConjecture/Lemmas/AlgebraicTopology/SingularBarycentricOuterFaces.lean:381}}.
\end{definition}

\section{SingularAffineSubdivision}\label{sec:HodgeConjecture.Lemmas.AlgebraicTopology.SingularAffineSubdivision}
\noindent\emph{File:} \href{https://github.com/Paul-Lez/HodgeConjecture/blob/main/HodgeConjecture/Lemmas/AlgebraicTopology/SingularAffineSubdivision.lean}{\texttt{HodgeConjecture/Lemmas/AlgebraicTopology/SingularAffineSubdivision.lean}}.

This module is ported from Paul Lezeau's corresponding file in
\texttt{sphere-six-complex} pull request \#49, under the Apache-2.0 license.

\paragraph{Affine realization of barycentric flags as singular simplices}

A simplex in the nerve model of barycentric subdivision is a flag of nonempty faces.  This
file sends such a flag to the affine singular simplex whose vertices are the barycenters of
those faces.  The construction is deliberately degreewise: the codimension-one face formula
is the piece needed to transport the combinatorial subdivision-chain boundary calculation to
ordinary singular chains.

\begin{definition}[AlgebraicTopology.Singular.stdSimplexAffineCombination]\label{def:AlgebraicTopology.Singular.stdSimplexAffineCombination}
\lean{AlgebraicTopology.Singular.stdSimplexAffineCombination}\leanok
The affine combination of a finite family of points in a standard simplex.

\noindent\emph{Lean:} \texttt{def AlgebraicTopology.Singular.stdSimplexAffineCombination}, \href{https://github.com/Paul-Lez/HodgeConjecture/blob/main/HodgeConjecture/Lemmas/AlgebraicTopology/SingularAffineSubdivision.lean#L44}{\texttt{HodgeConjecture/Lemmas/AlgebraicTopology/SingularAffineSubdivision.lean:44}}.
\end{definition}

\begin{definition}[AlgebraicTopology.Singular.nonemptyFiniteChainBarycenter]\label{def:AlgebraicTopology.Singular.nonemptyFiniteChainBarycenter}
\lean{AlgebraicTopology.Singular.nonemptyFiniteChainBarycenter}\leanok
The barycenter of a nonempty face of the standard \texttt{n}-simplex, embedded in the ambient
standard simplex.

\noindent\emph{Lean:} \texttt{def AlgebraicTopology.Singular.nonemptyFiniteChainBarycenter}, \href{https://github.com/Paul-Lez/HodgeConjecture/blob/main/HodgeConjecture/Lemmas/AlgebraicTopology/SingularAffineSubdivision.lean#L138}{\texttt{HodgeConjecture/Lemmas/AlgebraicTopology/SingularAffineSubdivision.lean:138}}.
\end{definition}

\begin{definition}[AlgebraicTopology.Singular.standardSimplexFaceVertexOrderHom]\label{def:AlgebraicTopology.Singular.standardSimplexFaceVertexOrderHom}
\lean{AlgebraicTopology.Singular.standardSimplexFaceVertexOrderHom}\leanok
The order embedding of vertices associated to a standard-simplex face inclusion.

\noindent\emph{Lean:} \texttt{def AlgebraicTopology.Singular.standardSimplexFaceVertexOrderHom}, \href{https://github.com/Paul-Lez/HodgeConjecture/blob/main/HodgeConjecture/Lemmas/AlgebraicTopology/SingularAffineSubdivision.lean#L186}{\texttt{HodgeConjecture/Lemmas/AlgebraicTopology/SingularAffineSubdivision.lean:186}}.
\end{definition}

\begin{definition}[AlgebraicTopology.Singular.affineFlagContinuousMap]\label{def:AlgebraicTopology.Singular.affineFlagContinuousMap}
\lean{AlgebraicTopology.Singular.affineFlagContinuousMap}\leanok
\uses{def:AlgebraicTopology.Singular.nonemptyFiniteChainBarycenter, def:AlgebraicTopology.Singular.stdSimplexAffineCombination}
The affine map associated to a flag of nonempty faces.

\noindent\emph{Lean:} \texttt{def AlgebraicTopology.Singular.affineFlagContinuousMap}, \href{https://github.com/Paul-Lez/HodgeConjecture/blob/main/HodgeConjecture/Lemmas/AlgebraicTopology/SingularAffineSubdivision.lean#L256}{\texttt{HodgeConjecture/Lemmas/AlgebraicTopology/SingularAffineSubdivision.lean:256}}.
\end{definition}

\begin{definition}[AlgebraicTopology.Singular.standardSimplexFaceContinuousMap]\label{def:AlgebraicTopology.Singular.standardSimplexFaceContinuousMap}
\lean{AlgebraicTopology.Singular.standardSimplexFaceContinuousMap}\leanok
The continuous affine inclusion of a codimension-one face of a standard simplex.

\noindent\emph{Lean:} \texttt{def AlgebraicTopology.Singular.standardSimplexFaceContinuousMap}, \href{https://github.com/Paul-Lez/HodgeConjecture/blob/main/HodgeConjecture/Lemmas/AlgebraicTopology/SingularAffineSubdivision.lean#L265}{\texttt{HodgeConjecture/Lemmas/AlgebraicTopology/SingularAffineSubdivision.lean:265}}.
\end{definition}

\begin{definition}[AlgebraicTopology.Singular.standardSimplexFaceTopCatMap]\label{def:AlgebraicTopology.Singular.standardSimplexFaceTopCatMap}
\lean{AlgebraicTopology.Singular.standardSimplexFaceTopCatMap}\leanok
\uses{def:AlgebraicTopology.Singular.standardSimplexFaceContinuousMap}
The face inclusion as a morphism of topological spaces.

\noindent\emph{Lean:} \texttt{def AlgebraicTopology.Singular.standardSimplexFaceTopCatMap}, \href{https://github.com/Paul-Lez/HodgeConjecture/blob/main/HodgeConjecture/Lemmas/AlgebraicTopology/SingularAffineSubdivision.lean#L297}{\texttt{HodgeConjecture/Lemmas/AlgebraicTopology/SingularAffineSubdivision.lean:297}}.
\end{definition}

\begin{definition}[AlgebraicTopology.Singular.affineFlagSingularSimplex]\label{def:AlgebraicTopology.Singular.affineFlagSingularSimplex}
\lean{AlgebraicTopology.Singular.affineFlagSingularSimplex}\leanok
\uses{def:AlgebraicTopology.Singular.affineFlagContinuousMap}
A flag in the barycentric subdivision nerve, regarded as an ordinary singular simplex of
the ambient topological standard simplex.

\noindent\emph{Lean:} \texttt{def AlgebraicTopology.Singular.affineFlagSingularSimplex}, \href{https://github.com/Paul-Lez/HodgeConjecture/blob/main/HodgeConjecture/Lemmas/AlgebraicTopology/SingularAffineSubdivision.lean#L304}{\texttt{HodgeConjecture/Lemmas/AlgebraicTopology/SingularAffineSubdivision.lean:304}}.
\end{definition}

\begin{definition}[AlgebraicTopology.Singular.affineFlagChainComponent]\label{def:AlgebraicTopology.Singular.affineFlagChainComponent}
\lean{AlgebraicTopology.Singular.affineFlagChainComponent}\leanok
\uses{def:AlgebraicTopology.Singular.affineFlagSingularSimplex}
The degree-\texttt{k} homomorphism which realizes every flag as its affine singular simplex.

\noindent\emph{Lean:} \texttt{def AlgebraicTopology.Singular.affineFlagChainComponent}, \href{https://github.com/Paul-Lez/HodgeConjecture/blob/main/HodgeConjecture/Lemmas/AlgebraicTopology/SingularAffineSubdivision.lean#L368}{\texttt{HodgeConjecture/Lemmas/AlgebraicTopology/SingularAffineSubdivision.lean:368}}.
\end{definition}

\begin{definition}[AlgebraicTopology.Singular.affineFlagChainMap]\label{def:AlgebraicTopology.Singular.affineFlagChainMap}
\lean{AlgebraicTopology.Singular.affineFlagChainMap}\leanok
\uses{def:AlgebraicTopology.Singular.affineFlagChainComponent}
The all-degree affine realization from the barycentric nerve model to the ordinary singular
chain complex of the topological standard simplex.

\noindent\emph{Lean:} \texttt{def AlgebraicTopology.Singular.affineFlagChainMap}, \href{https://github.com/Paul-Lez/HodgeConjecture/blob/main/HodgeConjecture/Lemmas/AlgebraicTopology/SingularAffineSubdivision.lean#L406}{\texttt{HodgeConjecture/Lemmas/AlgebraicTopology/SingularAffineSubdivision.lean:406}}.
\end{definition}

\begin{definition}[AlgebraicTopology.Singular.affineSubdividedSimplexFundamentalChain]\label{def:AlgebraicTopology.Singular.affineSubdividedSimplexFundamentalChain}
\lean{AlgebraicTopology.Singular.affineSubdividedSimplexFundamentalChain}\leanok
\uses{def:AlgebraicTopology.Singular.affineFlagChainMap, def:AlgebraicTopology.Singular.subdividedSimplexFundamentalChain}
The signed affine barycentric subdivision of the topological standard \texttt{n}-simplex, now as
an actual integral singular chain.

\noindent\emph{Lean:} \texttt{def AlgebraicTopology.Singular.affineSubdividedSimplexFundamentalChain}, \href{https://github.com/Paul-Lez/HodgeConjecture/blob/main/HodgeConjecture/Lemmas/AlgebraicTopology/SingularAffineSubdivision.lean#L454}{\texttt{HodgeConjecture/Lemmas/AlgebraicTopology/SingularAffineSubdivision.lean:454}}.
\end{definition}

\begin{definition}[AlgebraicTopology.Singular.affineSubdividedSimplexAlternatingFaceChain]\label{def:AlgebraicTopology.Singular.affineSubdividedSimplexAlternatingFaceChain}
\lean{AlgebraicTopology.Singular.affineSubdividedSimplexAlternatingFaceChain}\leanok
\uses{def:AlgebraicTopology.Singular.affineFlagChainMap, def:AlgebraicTopology.Singular.subdividedSimplexAlternatingFaceChain}
The realized alternating boundary chain: each subdivided face flag is interpreted inside
the ambient topological \texttt{(n+1)}-simplex.

\noindent\emph{Lean:} \texttt{def AlgebraicTopology.Singular.affineSubdividedSimplexAlternatingFaceChain}, \href{https://github.com/Paul-Lez/HodgeConjecture/blob/main/HodgeConjecture/Lemmas/AlgebraicTopology/SingularAffineSubdivision.lean#L486}{\texttt{HodgeConjecture/Lemmas/AlgebraicTopology/SingularAffineSubdivision.lean:486}}.
\end{definition}

\begin{definition}[AlgebraicTopology.Singular.affineSubdividedSimplexExpectedBoundaryChain]\label{def:AlgebraicTopology.Singular.affineSubdividedSimplexExpectedBoundaryChain}
\lean{AlgebraicTopology.Singular.affineSubdividedSimplexExpectedBoundaryChain}\leanok
\uses{def:AlgebraicTopology.Singular.affineSubdividedSimplexFundamentalChain, def:AlgebraicTopology.Singular.standardSimplexFaceTopCatMap}
The alternating sum of the affine subdivided faces, each included into the ambient
topological standard simplex.

\noindent\emph{Lean:} \texttt{def AlgebraicTopology.Singular.affineSubdividedSimplexExpectedBoundaryChain}, \href{https://github.com/Paul-Lez/HodgeConjecture/blob/main/HodgeConjecture/Lemmas/AlgebraicTopology/SingularAffineSubdivision.lean#L494}{\texttt{HodgeConjecture/Lemmas/AlgebraicTopology/SingularAffineSubdivision.lean:494}}.
\end{definition}

\begin{definition}[AlgebraicTopology.Singular.singularSimplexTopCatMap]\label{def:AlgebraicTopology.Singular.singularSimplexTopCatMap}
\lean{AlgebraicTopology.Singular.singularSimplexTopCatMap}\leanok
The topological map represented by a singular simplex.

\noindent\emph{Lean:} \texttt{def AlgebraicTopology.Singular.singularSimplexTopCatMap}, \href{https://github.com/Paul-Lez/HodgeConjecture/blob/main/HodgeConjecture/Lemmas/AlgebraicTopology/SingularAffineSubdivision.lean#L553}{\texttt{HodgeConjecture/Lemmas/AlgebraicTopology/SingularAffineSubdivision.lean:553}}.
\end{definition}

\begin{definition}[AlgebraicTopology.Singular.affineSubdivisionSingularSimplexChain]\label{def:AlgebraicTopology.Singular.affineSubdivisionSingularSimplexChain}
\lean{AlgebraicTopology.Singular.affineSubdivisionSingularSimplexChain}\leanok
\uses{def:AlgebraicTopology.Singular.affineSubdividedSimplexFundamentalChain, def:AlgebraicTopology.Singular.singularSimplexTopCatMap}
The affine barycentric subdivision chain associated to one singular simplex.

\noindent\emph{Lean:} \texttt{def AlgebraicTopology.Singular.affineSubdivisionSingularSimplexChain}, \href{https://github.com/Paul-Lez/HodgeConjecture/blob/main/HodgeConjecture/Lemmas/AlgebraicTopology/SingularAffineSubdivision.lean#L584}{\texttt{HodgeConjecture/Lemmas/AlgebraicTopology/SingularAffineSubdivision.lean:584}}.
\end{definition}

\begin{definition}[AlgebraicTopology.Singular.affineSubdivisionSingularSimplexAlternatingFaceChain]\label{def:AlgebraicTopology.Singular.affineSubdivisionSingularSimplexAlternatingFaceChain}
\lean{AlgebraicTopology.Singular.affineSubdivisionSingularSimplexAlternatingFaceChain}\leanok
\uses{def:AlgebraicTopology.Singular.affineSubdivisionSingularSimplexChain}
The alternating sum of the affine subdivisions of the faces of a singular simplex.

\noindent\emph{Lean:} \texttt{def AlgebraicTopology.Singular.affineSubdivisionSingularSimplexAlternatingFaceChain}, \href{https://github.com/Paul-Lez/HodgeConjecture/blob/main/HodgeConjecture/Lemmas/AlgebraicTopology/SingularAffineSubdivision.lean#L638}{\texttt{HodgeConjecture/Lemmas/AlgebraicTopology/SingularAffineSubdivision.lean:638}}.
\end{definition}

\begin{definition}[AlgebraicTopology.Singular.affineSingularSubdivisionComponent]\label{def:AlgebraicTopology.Singular.affineSingularSubdivisionComponent}
\lean{AlgebraicTopology.Singular.affineSingularSubdivisionComponent}\leanok
\uses{def:AlgebraicTopology.Singular.affineSubdivisionSingularSimplexChain}
The degreewise affine barycentric subdivision operator on integral singular chains.

\noindent\emph{Lean:} \texttt{def AlgebraicTopology.Singular.affineSingularSubdivisionComponent}, \href{https://github.com/Paul-Lez/HodgeConjecture/blob/main/HodgeConjecture/Lemmas/AlgebraicTopology/SingularAffineSubdivision.lean#L707}{\texttt{HodgeConjecture/Lemmas/AlgebraicTopology/SingularAffineSubdivision.lean:707}}.
\end{definition}

\begin{definition}[AlgebraicTopology.Singular.affineSingularSubdivisionChainMap]\label{def:AlgebraicTopology.Singular.affineSingularSubdivisionChainMap}
\lean{AlgebraicTopology.Singular.affineSingularSubdivisionChainMap}\leanok
\uses{def:AlgebraicTopology.Singular.affineSingularSubdivisionComponent, def:AlgebraicTopology.Singular.affineSubdividedSimplexAlternatingFaceChain, def:AlgebraicTopology.Singular.affineSubdividedSimplexExpectedBoundaryChain, def:AlgebraicTopology.Singular.affineSubdivisionSingularSimplexAlternatingFaceChain, def:AlgebraicTopology.Singular.finCastSuccEquivNotLast, def:AlgebraicTopology.Singular.insertOmittedVertexLastEquiv, def:AlgebraicTopology.Singular.standardSimplexFaceVertexOrderHom, def:AlgebraicTopology.Singular.subdividedSimplexInteriorFaceSum, def:AlgebraicTopology.Singular.subdividedSimplexOuterFaceSum}
The genuine affine barycentric subdivision endomorphism of the integral singular chain
complex of a topological space.

\noindent\emph{Lean:} \texttt{def AlgebraicTopology.Singular.affineSingularSubdivisionChainMap}, \href{https://github.com/Paul-Lez/HodgeConjecture/blob/main/HodgeConjecture/Lemmas/AlgebraicTopology/SingularAffineSubdivision.lean#L757}{\texttt{HodgeConjecture/Lemmas/AlgebraicTopology/SingularAffineSubdivision.lean:757}}.
\end{definition}

\section{SingularAffineSubdivisionPrism}\label{sec:HodgeConjecture.Lemmas.AlgebraicTopology.SingularAffineSubdivisionPrism}
\noindent\emph{File:} \href{https://github.com/Paul-Lez/HodgeConjecture/blob/main/HodgeConjecture/Lemmas/AlgebraicTopology/SingularAffineSubdivisionPrism.lean}{\texttt{HodgeConjecture/Lemmas/AlgebraicTopology/SingularAffineSubdivisionPrism.lean}}.

This module is ported from Paul Lezeau's corresponding file in
\texttt{sphere-six-complex} pull request \#49, under the Apache-2.0 license.

\paragraph{A universal prism for affine singular subdivision}

The singular chain complex of every topological standard simplex is exact in positive degrees.
This supplies fillers for the universal subdivision discrepancies and hence the acyclic-models
prism between affine barycentric subdivision and the identity.  The first part of this file
establishes the exactness and packages a selected filler for every positive-degree cycle.

\begin{definition}[AlgebraicTopology.Singular.IntegralSingularHomology]\label{def:AlgebraicTopology.Singular.IntegralSingularHomology}
\lean{AlgebraicTopology.Singular.IntegralSingularHomology}\leanok
Integral singular homology with coefficients in \texttt{ℤ}.

\noindent\emph{Lean:} \texttt{abbrev AlgebraicTopology.Singular.IntegralSingularHomology}, \href{https://github.com/Paul-Lez/HodgeConjecture/blob/main/HodgeConjecture/Lemmas/AlgebraicTopology/SingularAffineSubdivisionPrism.lean#L47}{\texttt{HodgeConjecture/Lemmas/AlgebraicTopology/SingularAffineSubdivisionPrism.lean:47}}.
\end{definition}

\begin{definition}[AlgebraicTopology.Singular.integralSingularHomologyEquivOfHomotopyEquiv]\label{def:AlgebraicTopology.Singular.integralSingularHomologyEquivOfHomotopyEquiv}
\lean{AlgebraicTopology.Singular.integralSingularHomologyEquivOfHomotopyEquiv}\leanok
\uses{def:AlgebraicTopology.Singular.IntegralSingularHomology}
Integral singular homology is invariant under a homotopy equivalence.

\noindent\emph{Lean:} \texttt{def AlgebraicTopology.Singular.integralSingularHomologyEquivOfHomotopyEquiv}, \href{https://github.com/Paul-Lez/HodgeConjecture/blob/main/HodgeConjecture/Lemmas/AlgebraicTopology/SingularAffineSubdivisionPrism.lean#L51}{\texttt{HodgeConjecture/Lemmas/AlgebraicTopology/SingularAffineSubdivisionPrism.lean:51}}.
\end{definition}

\begin{definition}[AlgebraicTopology.Singular.standardTopologicalSimplexCycleFiller]\label{def:AlgebraicTopology.Singular.standardTopologicalSimplexCycleFiller}
\lean{AlgebraicTopology.Singular.standardTopologicalSimplexCycleFiller}\leanok
\uses{def:AlgebraicTopology.Singular.integralSingularHomologyEquivOfHomotopyEquiv}
A selected filler of a positive-degree cycle in a topological standard simplex.

\noindent\emph{Lean:} \texttt{def AlgebraicTopology.Singular.standardTopologicalSimplexCycleFiller}, \href{https://github.com/Paul-Lez/HodgeConjecture/blob/main/HodgeConjecture/Lemmas/AlgebraicTopology/SingularAffineSubdivisionPrism.lean#L149}{\texttt{HodgeConjecture/Lemmas/AlgebraicTopology/SingularAffineSubdivisionPrism.lean:149}}.
\end{definition}

\begin{definition}[AlgebraicTopology.Singular.standardTopologicalSimplexTotalCycleFiller]\label{def:AlgebraicTopology.Singular.standardTopologicalSimplexTotalCycleFiller}
\lean{AlgebraicTopology.Singular.standardTopologicalSimplexTotalCycleFiller}\leanok
\uses{def:AlgebraicTopology.Singular.standardTopologicalSimplexCycleFiller}
A total version of the positive-degree filler.  It is the selected filler when its input
is a cycle and zero otherwise; this form can be used in a structurally recursive definition
before the cycle condition has been proved.

\noindent\emph{Lean:} \texttt{def AlgebraicTopology.Singular.standardTopologicalSimplexTotalCycleFiller}, \href{https://github.com/Paul-Lez/HodgeConjecture/blob/main/HodgeConjecture/Lemmas/AlgebraicTopology/SingularAffineSubdivisionPrism.lean#L183}{\texttt{HodgeConjecture/Lemmas/AlgebraicTopology/SingularAffineSubdivisionPrism.lean:183}}.
\end{definition}

\begin{definition}[AlgebraicTopology.Singular.standardTopologicalSimplexIdentitySimplex]\label{def:AlgebraicTopology.Singular.standardTopologicalSimplexIdentitySimplex}
\lean{AlgebraicTopology.Singular.standardTopologicalSimplexIdentitySimplex}\leanok
The identity map of the topological standard simplex, regarded as its universal singular
simplex.

\noindent\emph{Lean:} \texttt{def AlgebraicTopology.Singular.standardTopologicalSimplexIdentitySimplex}, \href{https://github.com/Paul-Lez/HodgeConjecture/blob/main/HodgeConjecture/Lemmas/AlgebraicTopology/SingularAffineSubdivisionPrism.lean#L222}{\texttt{HodgeConjecture/Lemmas/AlgebraicTopology/SingularAffineSubdivisionPrism.lean:222}}.
\end{definition}

\begin{definition}[AlgebraicTopology.Singular.affineSingularSubdivisionDiscrepancyChainMap]\label{def:AlgebraicTopology.Singular.affineSingularSubdivisionDiscrepancyChainMap}
\lean{AlgebraicTopology.Singular.affineSingularSubdivisionDiscrepancyChainMap}\leanok
\uses{def:AlgebraicTopology.Singular.affineSingularSubdivisionChainMap}
Affine singular subdivision minus the identity, as a natural chain endomorphism.

\noindent\emph{Lean:} \texttt{def AlgebraicTopology.Singular.affineSingularSubdivisionDiscrepancyChainMap}, \href{https://github.com/Paul-Lez/HodgeConjecture/blob/main/HodgeConjecture/Lemmas/AlgebraicTopology/SingularAffineSubdivisionPrism.lean#L243}{\texttt{HodgeConjecture/Lemmas/AlgebraicTopology/SingularAffineSubdivisionPrism.lean:243}}.
\end{definition}

\begin{definition}[AlgebraicTopology.Singular.standardAffineSubdivisionDiscrepancy]\label{def:AlgebraicTopology.Singular.standardAffineSubdivisionDiscrepancy}
\lean{AlgebraicTopology.Singular.standardAffineSubdivisionDiscrepancy}\leanok
\uses{def:AlgebraicTopology.Singular.affineSingularSubdivisionDiscrepancyChainMap, def:AlgebraicTopology.Singular.standardTopologicalSimplexIdentitySimplex}
The universal degree-\texttt{n} affine-subdivision discrepancy on the topological standard
\texttt{n}-simplex.

\noindent\emph{Lean:} \texttt{def AlgebraicTopology.Singular.standardAffineSubdivisionDiscrepancy}, \href{https://github.com/Paul-Lez/HodgeConjecture/blob/main/HodgeConjecture/Lemmas/AlgebraicTopology/SingularAffineSubdivisionPrism.lean#L250}{\texttt{HodgeConjecture/Lemmas/AlgebraicTopology/SingularAffineSubdivisionPrism.lean:250}}.
\end{definition}

\begin{definition}[AlgebraicTopology.Singular.universalAffinePrismComponent]\label{def:AlgebraicTopology.Singular.universalAffinePrismComponent}
\lean{AlgebraicTopology.Singular.universalAffinePrismComponent}\leanok
\uses{def:AlgebraicTopology.Singular.singularSimplexTopCatMap}
Transport one universal degree-raising singular chain along every singular simplex of an
arbitrary topological space.

\noindent\emph{Lean:} \texttt{def AlgebraicTopology.Singular.universalAffinePrismComponent}, \href{https://github.com/Paul-Lez/HodgeConjecture/blob/main/HodgeConjecture/Lemmas/AlgebraicTopology/SingularAffineSubdivisionPrism.lean#L394}{\texttt{HodgeConjecture/Lemmas/AlgebraicTopology/SingularAffineSubdivisionPrism.lean:394}}.
\end{definition}

\begin{definition}[AlgebraicTopology.Singular.standardAffinePrismFaceChain]\label{def:AlgebraicTopology.Singular.standardAffinePrismFaceChain}
\lean{AlgebraicTopology.Singular.standardAffinePrismFaceChain}\leanok
\uses{def:AlgebraicTopology.Singular.standardSimplexFaceTopCatMap}
The raw sum of the already constructed prisms over the faces of a standard simplex.

\noindent\emph{Lean:} \texttt{def AlgebraicTopology.Singular.standardAffinePrismFaceChain}, \href{https://github.com/Paul-Lez/HodgeConjecture/blob/main/HodgeConjecture/Lemmas/AlgebraicTopology/SingularAffineSubdivisionPrism.lean#L427}{\texttt{HodgeConjecture/Lemmas/AlgebraicTopology/SingularAffineSubdivisionPrism.lean:427}}.
\end{definition}

\begin{definition}[AlgebraicTopology.Singular.canonicalAffineSubdivisionPrism]\label{def:AlgebraicTopology.Singular.canonicalAffineSubdivisionPrism}
\lean{AlgebraicTopology.Singular.canonicalAffineSubdivisionPrism}\leanok
\uses{def:AlgebraicTopology.Singular.standardAffineSubdivisionDiscrepancy, def:AlgebraicTopology.Singular.standardTopologicalSimplexTotalCycleFiller}
The universal affine prisms, defined recursively by filling the discrepancy left after
the already constructed face prisms.

\noindent\emph{Lean:} \texttt{def AlgebraicTopology.Singular.canonicalAffineSubdivisionPrism}, \href{https://github.com/Paul-Lez/HodgeConjecture/blob/main/HodgeConjecture/Lemmas/AlgebraicTopology/SingularAffineSubdivisionPrism.lean#L639}{\texttt{HodgeConjecture/Lemmas/AlgebraicTopology/SingularAffineSubdivisionPrism.lean:639}}.
\end{definition}

\begin{definition}[AlgebraicTopology.Singular.affineSingularSubdivisionPrismComponent]\label{def:AlgebraicTopology.Singular.affineSingularSubdivisionPrismComponent}
\lean{AlgebraicTopology.Singular.affineSingularSubdivisionPrismComponent}\leanok
\uses{def:AlgebraicTopology.Singular.canonicalAffineSubdivisionPrism, def:AlgebraicTopology.Singular.universalAffinePrismComponent}
The canonical universal prism transported to every singular simplex in degree \texttt{n}.

\noindent\emph{Lean:} \texttt{def AlgebraicTopology.Singular.affineSingularSubdivisionPrismComponent}, \href{https://github.com/Paul-Lez/HodgeConjecture/blob/main/HodgeConjecture/Lemmas/AlgebraicTopology/SingularAffineSubdivisionPrism.lean#L717}{\texttt{HodgeConjecture/Lemmas/AlgebraicTopology/SingularAffineSubdivisionPrism.lean:717}}.
\end{definition}

\begin{definition}[AlgebraicTopology.Singular.affineSingularSubdivisionPrismHom]\label{def:AlgebraicTopology.Singular.affineSingularSubdivisionPrismHom}
\lean{AlgebraicTopology.Singular.affineSingularSubdivisionPrismHom}\leanok
\uses{def:AlgebraicTopology.Singular.affineSingularSubdivisionPrismComponent}
The homotopy family, supported from degree \texttt{i} to degree \texttt{i+1}.

\noindent\emph{Lean:} \texttt{def AlgebraicTopology.Singular.affineSingularSubdivisionPrismHom}, \href{https://github.com/Paul-Lez/HodgeConjecture/blob/main/HodgeConjecture/Lemmas/AlgebraicTopology/SingularAffineSubdivisionPrism.lean#L833}{\texttt{HodgeConjecture/Lemmas/AlgebraicTopology/SingularAffineSubdivisionPrism.lean:833}}.
\end{definition}

\begin{definition}[AlgebraicTopology.Singular.affineSingularSubdivisionHomotopy]\label{def:AlgebraicTopology.Singular.affineSingularSubdivisionHomotopy}
\lean{AlgebraicTopology.Singular.affineSingularSubdivisionHomotopy}\leanok
\uses{def:AlgebraicTopology.Singular.affineSingularSubdivisionPrismHom, def:AlgebraicTopology.Singular.standardAffinePrismFaceChain}
Affine barycentric subdivision is chain homotopic to the identity on the integral singular
chain complex of every topological space.

\noindent\emph{Lean:} \texttt{def AlgebraicTopology.Singular.affineSingularSubdivisionHomotopy}, \href{https://github.com/Paul-Lez/HodgeConjecture/blob/main/HodgeConjecture/Lemmas/AlgebraicTopology/SingularAffineSubdivisionPrism.lean#L850}{\texttt{HodgeConjecture/Lemmas/AlgebraicTopology/SingularAffineSubdivisionPrism.lean:850}}.
\end{definition}

\section{SingularAffineSubdivisionSmall}\label{sec:HodgeConjecture.Lemmas.AlgebraicTopology.SingularAffineSubdivisionSmall}
\noindent\emph{File:} \href{https://github.com/Paul-Lez/HodgeConjecture/blob/main/HodgeConjecture/Lemmas/AlgebraicTopology/SingularAffineSubdivisionSmall.lean}{\texttt{HodgeConjecture/Lemmas/AlgebraicTopology/SingularAffineSubdivisionSmall.lean}}.

This module is ported from Paul Lezeau's corresponding file in
\texttt{sphere-six-complex} pull request \#49, under the Apache-2.0 license.

\paragraph{Affine subdivision on cover-small singular chains}

Affine barycentric subdivision of a simplex only precomposes that simplex.  Consequently a
simplex which factors through one member of a cover subdivides entirely inside the same member.
This file packages that range-preservation statement as an endomorphism of the cover-small chain
complex and proves compatibility with its inclusion into all singular chains.

\begin{definition}[AlgebraicTopology.Singular.coverSmallSimplexPreimageIndex]\label{def:AlgebraicTopology.Singular.coverSmallSimplexPreimageIndex}
\lean{AlgebraicTopology.Singular.coverSmallSimplexPreimageIndex}\leanok
\uses{def:AlgebraicTopology.Singular.coverSmallSingularSubcomplex}
A selected cover member through which a cover-small simplex factors.

\noindent\emph{Lean:} \texttt{def AlgebraicTopology.Singular.coverSmallSimplexPreimageIndex}, \href{https://github.com/Paul-Lez/HodgeConjecture/blob/main/HodgeConjecture/Lemmas/AlgebraicTopology/SingularAffineSubdivisionSmall.lean#L44}{\texttt{HodgeConjecture/Lemmas/AlgebraicTopology/SingularAffineSubdivisionSmall.lean:44}}.
\end{definition}

\begin{definition}[AlgebraicTopology.Singular.coverSmallSimplexPreimage]\label{def:AlgebraicTopology.Singular.coverSmallSimplexPreimage}
\lean{AlgebraicTopology.Singular.coverSmallSimplexPreimage}\leanok
\uses{def:AlgebraicTopology.Singular.coverSmallSimplexPreimageIndex}
A selected lift of a cover-small simplex to its selected cover member.

\noindent\emph{Lean:} \texttt{def AlgebraicTopology.Singular.coverSmallSimplexPreimage}, \href{https://github.com/Paul-Lez/HodgeConjecture/blob/main/HodgeConjecture/Lemmas/AlgebraicTopology/SingularAffineSubdivisionSmall.lean#L50}{\texttt{HodgeConjecture/Lemmas/AlgebraicTopology/SingularAffineSubdivisionSmall.lean:50}}.
\end{definition}

\begin{definition}[AlgebraicTopology.Singular.coverSmallAffineSubdivisionSimplexChain]\label{def:AlgebraicTopology.Singular.coverSmallAffineSubdivisionSimplexChain}
\lean{AlgebraicTopology.Singular.coverSmallAffineSubdivisionSimplexChain}\leanok
\uses{def:AlgebraicTopology.Singular.affineSubdivisionSingularSimplexChain, def:AlgebraicTopology.Singular.coverMemberToSmallIntegralSingularChains, def:AlgebraicTopology.Singular.coverSmallSimplexPreimage}
Affine subdivision of a cover-small generator, performed in a cover member and then mapped
back into the cover-small complex.

\noindent\emph{Lean:} \texttt{def AlgebraicTopology.Singular.coverSmallAffineSubdivisionSimplexChain}, \href{https://github.com/Paul-Lez/HodgeConjecture/blob/main/HodgeConjecture/Lemmas/AlgebraicTopology/SingularAffineSubdivisionSmall.lean#L70}{\texttt{HodgeConjecture/Lemmas/AlgebraicTopology/SingularAffineSubdivisionSmall.lean:70}}.
\end{definition}

\begin{definition}[AlgebraicTopology.Singular.coverSmallAffineSubdivisionComponent]\label{def:AlgebraicTopology.Singular.coverSmallAffineSubdivisionComponent}
\lean{AlgebraicTopology.Singular.coverSmallAffineSubdivisionComponent}\leanok
\uses{def:AlgebraicTopology.Singular.coverSmallAffineSubdivisionSimplexChain}
The degreewise affine subdivision endomorphism of cover-small chains.

\noindent\emph{Lean:} \texttt{def AlgebraicTopology.Singular.coverSmallAffineSubdivisionComponent}, \href{https://github.com/Paul-Lez/HodgeConjecture/blob/main/HodgeConjecture/Lemmas/AlgebraicTopology/SingularAffineSubdivisionSmall.lean#L84}{\texttt{HodgeConjecture/Lemmas/AlgebraicTopology/SingularAffineSubdivisionSmall.lean:84}}.
\end{definition}

\begin{definition}[AlgebraicTopology.Singular.coverSmallAffineSubdivisionChainMap]\label{def:AlgebraicTopology.Singular.coverSmallAffineSubdivisionChainMap}
\lean{AlgebraicTopology.Singular.coverSmallAffineSubdivisionChainMap}\leanok
\uses{def:AlgebraicTopology.Singular.affineSingularSubdivisionComponent, def:AlgebraicTopology.Singular.affineSubdividedSimplexAlternatingFaceChain, def:AlgebraicTopology.Singular.affineSubdividedSimplexExpectedBoundaryChain, def:AlgebraicTopology.Singular.affineSubdivisionSingularSimplexAlternatingFaceChain, def:AlgebraicTopology.Singular.coverSmallAffineSubdivisionComponent, def:AlgebraicTopology.Singular.coverSmallIntegralSingularChainInclusion, def:AlgebraicTopology.Singular.finCastSuccEquivNotLast, def:AlgebraicTopology.Singular.insertOmittedVertexLastEquiv, def:AlgebraicTopology.Singular.integralSingularChainMapObj, def:AlgebraicTopology.Singular.standardSimplexFaceVertexOrderHom, def:AlgebraicTopology.Singular.subdividedSimplexInteriorFaceSum, def:AlgebraicTopology.Singular.subdividedSimplexOuterFaceSum}
Affine barycentric subdivision as an endomorphism of the cover-small chain complex.

\noindent\emph{Lean:} \texttt{def AlgebraicTopology.Singular.coverSmallAffineSubdivisionChainMap}, \href{https://github.com/Paul-Lez/HodgeConjecture/blob/main/HodgeConjecture/Lemmas/AlgebraicTopology/SingularAffineSubdivisionSmall.lean#L182}{\texttt{HodgeConjecture/Lemmas/AlgebraicTopology/SingularAffineSubdivisionSmall.lean:182}}.
\end{definition}

\section{SingularAffineSubdivisionSmallPrism}\label{sec:HodgeConjecture.Lemmas.AlgebraicTopology.SingularAffineSubdivisionSmallPrism}
\noindent\emph{File:} \href{https://github.com/Paul-Lez/HodgeConjecture/blob/main/HodgeConjecture/Lemmas/AlgebraicTopology/SingularAffineSubdivisionSmallPrism.lean}{\texttt{HodgeConjecture/Lemmas/AlgebraicTopology/SingularAffineSubdivisionSmallPrism.lean}}.

This module is ported from Paul Lezeau's corresponding file in
\texttt{sphere-six-complex} pull request \#49, under the Apache-2.0 license.

\paragraph{The affine subdivision prism on cover-small singular chains}

The universal affine prism only precomposes a singular simplex.  It therefore preserves the
range of any cover member, just as affine subdivision itself does.  This file restricts the
prism to the cover-small chain complex and proves that cover-small affine subdivision is chain
homotopic to the identity there.

\begin{definition}[AlgebraicTopology.Singular.coverSmallAffinePrismSimplexChain]\label{def:AlgebraicTopology.Singular.coverSmallAffinePrismSimplexChain}
\lean{AlgebraicTopology.Singular.coverSmallAffinePrismSimplexChain}\leanok
\uses{def:AlgebraicTopology.Singular.affineSingularSubdivisionPrismComponent, def:AlgebraicTopology.Singular.coverMemberToSmallIntegralSingularChains, def:AlgebraicTopology.Singular.coverSmallSimplexPreimage}
The affine prism on a selected lift of a cover-small generator, mapped back into the small
chain complex.

\noindent\emph{Lean:} \texttt{def AlgebraicTopology.Singular.coverSmallAffinePrismSimplexChain}, \href{https://github.com/Paul-Lez/HodgeConjecture/blob/main/HodgeConjecture/Lemmas/AlgebraicTopology/SingularAffineSubdivisionSmallPrism.lean#L46}{\texttt{HodgeConjecture/Lemmas/AlgebraicTopology/SingularAffineSubdivisionSmallPrism.lean:46}}.
\end{definition}

\begin{definition}[AlgebraicTopology.Singular.coverSmallAffinePrismComponent]\label{def:AlgebraicTopology.Singular.coverSmallAffinePrismComponent}
\lean{AlgebraicTopology.Singular.coverSmallAffinePrismComponent}\leanok
\uses{def:AlgebraicTopology.Singular.coverSmallAffinePrismSimplexChain}
The degree-raising affine prism operator on cover-small chains.

\noindent\emph{Lean:} \texttt{def AlgebraicTopology.Singular.coverSmallAffinePrismComponent}, \href{https://github.com/Paul-Lez/HodgeConjecture/blob/main/HodgeConjecture/Lemmas/AlgebraicTopology/SingularAffineSubdivisionSmallPrism.lean#L62}{\texttt{HodgeConjecture/Lemmas/AlgebraicTopology/SingularAffineSubdivisionSmallPrism.lean:62}}.
\end{definition}

\begin{definition}[AlgebraicTopology.Singular.coverSmallAffinePrismHom]\label{def:AlgebraicTopology.Singular.coverSmallAffinePrismHom}
\lean{AlgebraicTopology.Singular.coverSmallAffinePrismHom}\leanok
\uses{def:AlgebraicTopology.Singular.coverSmallAffinePrismComponent}
The degree-raising small-prism family.

\noindent\emph{Lean:} \texttt{def AlgebraicTopology.Singular.coverSmallAffinePrismHom}, \href{https://github.com/Paul-Lez/HodgeConjecture/blob/main/HodgeConjecture/Lemmas/AlgebraicTopology/SingularAffineSubdivisionSmallPrism.lean#L206}{\texttt{HodgeConjecture/Lemmas/AlgebraicTopology/SingularAffineSubdivisionSmallPrism.lean:206}}.
\end{definition}

\begin{definition}[AlgebraicTopology.Singular.coverSmallAffineSubdivisionHomotopy]\label{def:AlgebraicTopology.Singular.coverSmallAffineSubdivisionHomotopy}
\lean{AlgebraicTopology.Singular.coverSmallAffineSubdivisionHomotopy}\leanok
\uses{def:AlgebraicTopology.Singular.coverSmallAffinePrismHom, def:AlgebraicTopology.Singular.coverSmallAffineSubdivisionChainMap, def:AlgebraicTopology.Singular.standardAffinePrismFaceChain}
Cover-small affine subdivision is chain homotopic to the identity.

\noindent\emph{Lean:} \texttt{def AlgebraicTopology.Singular.coverSmallAffineSubdivisionHomotopy}, \href{https://github.com/Paul-Lez/HodgeConjecture/blob/main/HodgeConjecture/Lemmas/AlgebraicTopology/SingularAffineSubdivisionSmallPrism.lean#L221}{\texttt{HodgeConjecture/Lemmas/AlgebraicTopology/SingularAffineSubdivisionSmallPrism.lean:221}}.
\end{definition}

\section{SingularAffineSubdivisionIteration}\label{sec:HodgeConjecture.Lemmas.AlgebraicTopology.SingularAffineSubdivisionIteration}
\noindent\emph{File:} \href{https://github.com/Paul-Lez/HodgeConjecture/blob/main/HodgeConjecture/Lemmas/AlgebraicTopology/SingularAffineSubdivisionIteration.lean}{\texttt{HodgeConjecture/Lemmas/AlgebraicTopology/SingularAffineSubdivisionIteration.lean}}.

This module is ported from Paul Lezeau's corresponding file in
\texttt{sphere-six-complex} pull request \#49, under the Apache-2.0 license.

\paragraph{Iterated affine subdivision on singular chains}

This file packages finite iterates of geometric affine subdivision on full and cover-small
singular chains.  Every iterate is chain homotopic to the identity, and the small/full iterates
commute with the canonical inclusion.

\begin{definition}[AlgebraicTopology.Singular.affineSingularSubdivisionIterate]\label{def:AlgebraicTopology.Singular.affineSingularSubdivisionIterate}
\lean{AlgebraicTopology.Singular.affineSingularSubdivisionIterate}\leanok
\uses{def:AlgebraicTopology.Singular.affineSingularSubdivisionChainMap}
Finite iteration of affine subdivision on all integral singular chains.

\noindent\emph{Lean:} \texttt{def AlgebraicTopology.Singular.affineSingularSubdivisionIterate}, \href{https://github.com/Paul-Lez/HodgeConjecture/blob/main/HodgeConjecture/Lemmas/AlgebraicTopology/SingularAffineSubdivisionIteration.lean#L39}{\texttt{HodgeConjecture/Lemmas/AlgebraicTopology/SingularAffineSubdivisionIteration.lean:39}}.
\end{definition}

\begin{definition}[AlgebraicTopology.Singular.affineSingularSubdivisionIterateHomotopy]\label{def:AlgebraicTopology.Singular.affineSingularSubdivisionIterateHomotopy}
\lean{AlgebraicTopology.Singular.affineSingularSubdivisionIterateHomotopy}\leanok
\uses{def:AlgebraicTopology.Singular.affineSingularSubdivisionHomotopy, def:AlgebraicTopology.Singular.affineSingularSubdivisionIterate}
Every finite affine-subdivision iterate is chain homotopic to the identity.

\noindent\emph{Lean:} \texttt{def AlgebraicTopology.Singular.affineSingularSubdivisionIterateHomotopy}, \href{https://github.com/Paul-Lez/HodgeConjecture/blob/main/HodgeConjecture/Lemmas/AlgebraicTopology/SingularAffineSubdivisionIteration.lean#L74}{\texttt{HodgeConjecture/Lemmas/AlgebraicTopology/SingularAffineSubdivisionIteration.lean:74}}.
\end{definition}

\begin{definition}[AlgebraicTopology.Singular.coverSmallAffineSubdivisionIterate]\label{def:AlgebraicTopology.Singular.coverSmallAffineSubdivisionIterate}
\lean{AlgebraicTopology.Singular.coverSmallAffineSubdivisionIterate}\leanok
\uses{def:AlgebraicTopology.Singular.coverSmallAffineSubdivisionChainMap}
Finite iteration of affine subdivision on cover-small chains.

\noindent\emph{Lean:} \texttt{def AlgebraicTopology.Singular.coverSmallAffineSubdivisionIterate}, \href{https://github.com/Paul-Lez/HodgeConjecture/blob/main/HodgeConjecture/Lemmas/AlgebraicTopology/SingularAffineSubdivisionIteration.lean#L99}{\texttt{HodgeConjecture/Lemmas/AlgebraicTopology/SingularAffineSubdivisionIteration.lean:99}}.
\end{definition}

\begin{definition}[AlgebraicTopology.Singular.coverSmallAffineSubdivisionIterateHomotopy]\label{def:AlgebraicTopology.Singular.coverSmallAffineSubdivisionIterateHomotopy}
\lean{AlgebraicTopology.Singular.coverSmallAffineSubdivisionIterateHomotopy}\leanok
\uses{def:AlgebraicTopology.Singular.coverSmallAffineSubdivisionHomotopy, def:AlgebraicTopology.Singular.coverSmallAffineSubdivisionIterate}
Every cover-small affine-subdivision iterate is chain homotopic to the identity.

\noindent\emph{Lean:} \texttt{def AlgebraicTopology.Singular.coverSmallAffineSubdivisionIterateHomotopy}, \href{https://github.com/Paul-Lez/HodgeConjecture/blob/main/HodgeConjecture/Lemmas/AlgebraicTopology/SingularAffineSubdivisionIteration.lean#L118}{\texttt{HodgeConjecture/Lemmas/AlgebraicTopology/SingularAffineSubdivisionIteration.lean:118}}.
\end{definition}

\section{SingularAffineSubdivisionMesh}\label{sec:HodgeConjecture.Lemmas.AlgebraicTopology.SingularAffineSubdivisionMesh}
\noindent\emph{File:} \href{https://github.com/Paul-Lez/HodgeConjecture/blob/main/HodgeConjecture/Lemmas/AlgebraicTopology/SingularAffineSubdivisionMesh.lean}{\texttt{HodgeConjecture/Lemmas/AlgebraicTopology/SingularAffineSubdivisionMesh.lean}}.

This module is ported from Paul Lezeau's corresponding file in
\texttt{sphere-six-complex} pull request \#49, under the Apache-2.0 license.

\paragraph{Mesh estimates for affine barycentric subdivision}

This file proves the quantitative geometric estimate behind the small-chain argument.  In the
sup metric on the standard \texttt{n}-simplex, the affine simplex associated to every flag of nonempty
faces has diameter at most \texttt{n / (n + 1)}.  The proof uses the exact face-barycenter coordinates
from \texttt{SingularAffineSubdivision}.

The final theorem packages the metric conclusion needed for iterated subdivision: any family of
nonempty cells whose diameters are bounded by successive powers of this factor is eventually
subordinate to an arbitrary open cover of a compact metric space.  Connecting that abstract
family to the chain-level iterate requires keeping the ancestry of every iterated affine flag;
that combinatorial bookkeeping is intentionally separate from the metric argument here.

\begin{definition}[AlgebraicTopology.Singular.barycentricContractionFactor]\label{def:AlgebraicTopology.Singular.barycentricContractionFactor}
\lean{AlgebraicTopology.Singular.barycentricContractionFactor}\leanok
The classical mesh-contraction factor for barycentric subdivision of an \texttt{n}-simplex.

\noindent\emph{Lean:} \texttt{def AlgebraicTopology.Singular.barycentricContractionFactor}, \href{https://github.com/Paul-Lez/HodgeConjecture/blob/main/HodgeConjecture/Lemmas/AlgebraicTopology/SingularAffineSubdivisionMesh.lean#L49}{\texttt{HodgeConjecture/Lemmas/AlgebraicTopology/SingularAffineSubdivisionMesh.lean:49}}.
\end{definition}

\begin{definition}[AlgebraicTopology.Singular.TopAffineFlag]\label{def:AlgebraicTopology.Singular.TopAffineFlag}
\lean{AlgebraicTopology.Singular.TopAffineFlag}\leanok
A top-dimensional flag in the barycentric subdivision of the standard \texttt{n}-simplex.

\noindent\emph{Lean:} \texttt{abbrev AlgebraicTopology.Singular.TopAffineFlag}, \href{https://github.com/Paul-Lez/HodgeConjecture/blob/main/HodgeConjecture/Lemmas/AlgebraicTopology/SingularAffineSubdivisionMesh.lean#L283}{\texttt{HodgeConjecture/Lemmas/AlgebraicTopology/SingularAffineSubdivisionMesh.lean:283}}.
\end{definition}

\begin{definition}[AlgebraicTopology.Singular.iteratedAffineCellMap]\label{def:AlgebraicTopology.Singular.iteratedAffineCellMap}
\lean{AlgebraicTopology.Singular.iteratedAffineCellMap}\leanok
\uses{def:AlgebraicTopology.Singular.TopAffineFlag, def:AlgebraicTopology.Singular.affineFlagContinuousMap}
The affine cell map described by a finite ancestry of top-dimensional flags.  The head of
the list is the newest (innermost) subdivision cell and the tail records its parent ancestry.

\noindent\emph{Lean:} \texttt{def AlgebraicTopology.Singular.iteratedAffineCellMap}, \href{https://github.com/Paul-Lez/HodgeConjecture/blob/main/HodgeConjecture/Lemmas/AlgebraicTopology/SingularAffineSubdivisionMesh.lean#L288}{\texttt{HodgeConjecture/Lemmas/AlgebraicTopology/SingularAffineSubdivisionMesh.lean:288}}.
\end{definition}

\section{SingularAffineSubdivisionRelativeMesh}\label{sec:HodgeConjecture.Lemmas.AlgebraicTopology.SingularAffineSubdivisionRelativeMesh}
\noindent\emph{File:} \href{https://github.com/Paul-Lez/HodgeConjecture/blob/main/HodgeConjecture/Lemmas/AlgebraicTopology/SingularAffineSubdivisionRelativeMesh.lean}{\texttt{HodgeConjecture/Lemmas/AlgebraicTopology/SingularAffineSubdivisionRelativeMesh.lean}}.

This module is ported from Paul Lezeau's corresponding file in
\texttt{sphere-six-complex} pull request \#49, under the Apache-2.0 license.

\paragraph{Relative mesh estimates for affine barycentric subdivision}

This file proves the relative (rather than merely intrinsic) mesh estimate.  If the vertices of
an affine parent simplex are \texttt{p i}, then the images of the barycenters in any flag of nested
nonempty faces have diameter at most \texttt{n / (n + 1)} times the diameter of the parent vertices.
The result is stated first in an arbitrary real normed vector space and then specialized to the
affine self-maps used by iterated singular subdivision.

\begin{definition}[AlgebraicTopology.Singular.affineParentFaceCentroid]\label{def:AlgebraicTopology.Singular.affineParentFaceCentroid}
\lean{AlgebraicTopology.Singular.affineParentFaceCentroid}\leanok
The centroid of the vertices indexed by a nonempty face.

\noindent\emph{Lean:} \texttt{def AlgebraicTopology.Singular.affineParentFaceCentroid}, \href{https://github.com/Paul-Lez/HodgeConjecture/blob/main/HodgeConjecture/Lemmas/AlgebraicTopology/SingularAffineSubdivisionRelativeMesh.lean#L45}{\texttt{HodgeConjecture/Lemmas/AlgebraicTopology/SingularAffineSubdivisionRelativeMesh.lean:45}}.
\end{definition}

\begin{definition}[AlgebraicTopology.Singular.normedAffineCombination]\label{def:AlgebraicTopology.Singular.normedAffineCombination}
\lean{AlgebraicTopology.Singular.normedAffineCombination}\leanok
The affine combination of a family in a real normed vector space.

\noindent\emph{Lean:} \texttt{def AlgebraicTopology.Singular.normedAffineCombination}, \href{https://github.com/Paul-Lez/HodgeConjecture/blob/main/HodgeConjecture/Lemmas/AlgebraicTopology/SingularAffineSubdivisionRelativeMesh.lean#L192}{\texttt{HodgeConjecture/Lemmas/AlgebraicTopology/SingularAffineSubdivisionRelativeMesh.lean:192}}.
\end{definition}

\begin{definition}[AlgebraicTopology.Singular.affineParentFlagVertexSet]\label{def:AlgebraicTopology.Singular.affineParentFlagVertexSet}
\lean{AlgebraicTopology.Singular.affineParentFlagVertexSet}\leanok
\uses{def:AlgebraicTopology.Singular.affineParentFaceCentroid}
The parent images of the barycenter vertices in an affine flag.

\noindent\emph{Lean:} \texttt{def AlgebraicTopology.Singular.affineParentFlagVertexSet}, \href{https://github.com/Paul-Lez/HodgeConjecture/blob/main/HodgeConjecture/Lemmas/AlgebraicTopology/SingularAffineSubdivisionRelativeMesh.lean#L248}{\texttt{HodgeConjecture/Lemmas/AlgebraicTopology/SingularAffineSubdivisionRelativeMesh.lean:248}}.
\end{definition}

\begin{definition}[AlgebraicTopology.Singular.affineParentContinuousMap]\label{def:AlgebraicTopology.Singular.affineParentContinuousMap}
\lean{AlgebraicTopology.Singular.affineParentContinuousMap}\leanok
\uses{def:AlgebraicTopology.Singular.stdSimplexAffineCombination}
The affine map determined by vertices in a standard simplex.

\noindent\emph{Lean:} \texttt{def AlgebraicTopology.Singular.affineParentContinuousMap}, \href{https://github.com/Paul-Lez/HodgeConjecture/blob/main/HodgeConjecture/Lemmas/AlgebraicTopology/SingularAffineSubdivisionRelativeMesh.lean#L337}{\texttt{HodgeConjecture/Lemmas/AlgebraicTopology/SingularAffineSubdivisionRelativeMesh.lean:337}}.
\end{definition}

\section{SingularAffineSubdivisionSupport}\label{sec:HodgeConjecture.Lemmas.AlgebraicTopology.SingularAffineSubdivisionSupport}
\noindent\emph{File:} \href{https://github.com/Paul-Lez/HodgeConjecture/blob/main/HodgeConjecture/Lemmas/AlgebraicTopology/SingularAffineSubdivisionSupport.lean}{\texttt{HodgeConjecture/Lemmas/AlgebraicTopology/SingularAffineSubdivisionSupport.lean}}.

This module is ported from Paul Lezeau's corresponding file in
\texttt{sphere-six-complex} pull request \#49, under the Apache-2.0 license.

\paragraph{Support of iterated affine subdivision}

This file connects the permutation expansion of affine barycentric subdivision with the
geometric ancestry cells used in the mesh estimate.  It then lifts an iterate whose ancestry
cells are subordinate to a cover into the cover-small singular chain complex.

\begin{definition}[AlgebraicTopology.Singular.iteratedAffineCellSingularSimplex]\label{def:AlgebraicTopology.Singular.iteratedAffineCellSingularSimplex}
\lean{AlgebraicTopology.Singular.iteratedAffineCellSingularSimplex}\leanok
\uses{def:AlgebraicTopology.Singular.iteratedAffineCellMap}
The singular simplex obtained by restricting \texttt{x} to one iterated affine ancestry cell.

\noindent\emph{Lean:} \texttt{def AlgebraicTopology.Singular.iteratedAffineCellSingularSimplex}, \href{https://github.com/Paul-Lez/HodgeConjecture/blob/main/HodgeConjecture/Lemmas/AlgebraicTopology/SingularAffineSubdivisionSupport.lean#L41}{\texttt{HodgeConjecture/Lemmas/AlgebraicTopology/SingularAffineSubdivisionSupport.lean:41}}.
\end{definition}

\begin{definition}[AlgebraicTopology.Singular.AffinePermutationAncestry]\label{def:AlgebraicTopology.Singular.AffinePermutationAncestry}
\lean{AlgebraicTopology.Singular.AffinePermutationAncestry}\leanok
A depth-\texttt{m} choice of one permutation maximal flag at every subdivision stage.

\noindent\emph{Lean:} \texttt{abbrev AlgebraicTopology.Singular.AffinePermutationAncestry}, \href{https://github.com/Paul-Lez/HodgeConjecture/blob/main/HodgeConjecture/Lemmas/AlgebraicTopology/SingularAffineSubdivisionSupport.lean#L137}{\texttt{HodgeConjecture/Lemmas/AlgebraicTopology/SingularAffineSubdivisionSupport.lean:137}}.
\end{definition}

\begin{definition}[AlgebraicTopology.Singular.affinePermutationAncestryFlags]\label{def:AlgebraicTopology.Singular.affinePermutationAncestryFlags}
\lean{AlgebraicTopology.Singular.affinePermutationAncestryFlags}\leanok
\uses{def:AlgebraicTopology.Singular.AffinePermutationAncestry, def:AlgebraicTopology.Singular.TopAffineFlag, def:AlgebraicTopology.Singular.permutationMaximalFlagSimplex}
The list of affine flags encoded by a permutation ancestry, newest first.

\noindent\emph{Lean:} \texttt{def AlgebraicTopology.Singular.affinePermutationAncestryFlags}, \href{https://github.com/Paul-Lez/HodgeConjecture/blob/main/HodgeConjecture/Lemmas/AlgebraicTopology/SingularAffineSubdivisionSupport.lean#L141}{\texttt{HodgeConjecture/Lemmas/AlgebraicTopology/SingularAffineSubdivisionSupport.lean:141}}.
\end{definition}

\begin{definition}[AlgebraicTopology.Singular.affinePermutationAncestrySign]\label{def:AlgebraicTopology.Singular.affinePermutationAncestrySign}
\lean{AlgebraicTopology.Singular.affinePermutationAncestrySign}\leanok
\uses{def:AlgebraicTopology.Singular.AffinePermutationAncestry, def:AlgebraicTopology.Singular.permutationSignInteger}
The product of the orientation signs along a permutation ancestry.

\noindent\emph{Lean:} \texttt{def AlgebraicTopology.Singular.affinePermutationAncestrySign}, \href{https://github.com/Paul-Lez/HodgeConjecture/blob/main/HodgeConjecture/Lemmas/AlgebraicTopology/SingularAffineSubdivisionSupport.lean#L147}{\texttt{HodgeConjecture/Lemmas/AlgebraicTopology/SingularAffineSubdivisionSupport.lean:147}}.
\end{definition}

\begin{definition}[AlgebraicTopology.Singular.coverSmallIteratedAffineCellSimplex]\label{def:AlgebraicTopology.Singular.coverSmallIteratedAffineCellSimplex}
\lean{AlgebraicTopology.Singular.coverSmallIteratedAffineCellSimplex}\leanok
\uses{def:AlgebraicTopology.Singular.coverSmallSingularSubcomplex, def:AlgebraicTopology.Singular.iteratedAffineCellSingularSimplex}
The cover-small simplex corresponding to a subordinate ancestry cell.

\noindent\emph{Lean:} \texttt{def AlgebraicTopology.Singular.coverSmallIteratedAffineCellSimplex}, \href{https://github.com/Paul-Lez/HodgeConjecture/blob/main/HodgeConjecture/Lemmas/AlgebraicTopology/SingularAffineSubdivisionSupport.lean#L260}{\texttt{HodgeConjecture/Lemmas/AlgebraicTopology/SingularAffineSubdivisionSupport.lean:260}}.
\end{definition}

\begin{definition}[AlgebraicTopology.Singular.coverSmallAffineAncestryLiftChain]\label{def:AlgebraicTopology.Singular.coverSmallAffineAncestryLiftChain}
\lean{AlgebraicTopology.Singular.coverSmallAffineAncestryLiftChain}\leanok
\uses{def:AlgebraicTopology.Singular.CoverSmallIntegralSingularChainComplex, def:AlgebraicTopology.Singular.affinePermutationAncestryFlags, def:AlgebraicTopology.Singular.affinePermutationAncestrySign, def:AlgebraicTopology.Singular.coverSmallIteratedAffineCellSimplex}
The signed ancestry expansion, lifted generator by generator to cover-small chains.

\noindent\emph{Lean:} \texttt{def AlgebraicTopology.Singular.coverSmallAffineAncestryLiftChain}, \href{https://github.com/Paul-Lez/HodgeConjecture/blob/main/HodgeConjecture/Lemmas/AlgebraicTopology/SingularAffineSubdivisionSupport.lean#L285}{\texttt{HodgeConjecture/Lemmas/AlgebraicTopology/SingularAffineSubdivisionSupport.lean:285}}.
\end{definition}

\begin{definition}[AlgebraicTopology.Singular.affineSubdivisionEventuallySmallAddSubgroup]\label{def:AlgebraicTopology.Singular.affineSubdivisionEventuallySmallAddSubgroup}
\lean{AlgebraicTopology.Singular.affineSubdivisionEventuallySmallAddSubgroup}\leanok
\uses{def:AlgebraicTopology.Singular.affineSingularSubdivisionIterate, def:AlgebraicTopology.Singular.coverSmallAffineSubdivisionIterate}
Chains which become cover-small after some (chain-dependent) number of affine subdivisions
form an additive subgroup.

\noindent\emph{Lean:} \texttt{def AlgebraicTopology.Singular.affineSubdivisionEventuallySmallAddSubgroup}, \href{https://github.com/Paul-Lez/HodgeConjecture/blob/main/HodgeConjecture/Lemmas/AlgebraicTopology/SingularAffineSubdivisionSupport.lean#L375}{\texttt{HodgeConjecture/Lemmas/AlgebraicTopology/SingularAffineSubdivisionSupport.lean:375}}.
\end{definition}

\section{SingularExcisionOpenCover}\label{sec:HodgeConjecture.Lemmas.AlgebraicTopology.SingularExcisionOpenCover}
\noindent\emph{File:} \href{https://github.com/Paul-Lez/HodgeConjecture/blob/main/HodgeConjecture/Lemmas/AlgebraicTopology/SingularExcisionOpenCover.lean}{\texttt{HodgeConjecture/Lemmas/AlgebraicTopology/SingularExcisionOpenCover.lean}}.

This module is ported from Paul Lezeau's corresponding file in
\texttt{sphere-six-complex} pull request \#49, under the Apache-2.0 license.

\paragraph{Small-chain approximation for open covers}

This file assembles the geometric affine-mesh theorem, the exact permutation-ancestry expansion,
the adaptive quasi-isomorphism argument, and projectivity.  The result is the classical
small-singular-chain theorem for every open cover, with no additional hypothesis.

\begin{definition}[AlgebraicTopology.Singular.coverSmallChainHomotopyEquiv\_of\_openCover]\label{def:AlgebraicTopology.Singular.coverSmallChainHomotopyEquiv_of_openCover}
\lean{AlgebraicTopology.Singular.coverSmallChainHomotopyEquiv_of_openCover}\leanok
\uses{def:AlgebraicTopology.Singular.affineParentContinuousMap, def:AlgebraicTopology.Singular.affineParentFlagVertexSet, def:AlgebraicTopology.Singular.affineSingularSubdivisionIterateHomotopy, def:AlgebraicTopology.Singular.affineSubdivisionEventuallySmallAddSubgroup, def:AlgebraicTopology.Singular.barycentricContractionFactor, def:AlgebraicTopology.Singular.coverSmallAffineAncestryLiftChain, def:AlgebraicTopology.Singular.coverSmallAffineSubdivisionIterateHomotopy, def:AlgebraicTopology.Singular.normedAffineCombination}
A selected chain-homotopy equivalence witnessing small-chain approximation.

\noindent\emph{Lean:} \texttt{def AlgebraicTopology.Singular.coverSmallChainHomotopyEquiv\_of\_openCover}, \href{https://github.com/Paul-Lez/HodgeConjecture/blob/main/HodgeConjecture/Lemmas/AlgebraicTopology/SingularExcisionOpenCover.lean#L60}{\texttt{HodgeConjecture/Lemmas/AlgebraicTopology/SingularExcisionOpenCover.lean:60}}.
\end{definition}

\section{Small singular chains with rational coefficients}\label{sec:HodgeConjecture.Lemmas.AlgebraicTopology.SingularExcisionField}
\noindent\emph{File:} \href{https://github.com/Paul-Lez/HodgeConjecture/blob/main/HodgeConjecture/Lemmas/AlgebraicTopology/SingularExcisionField.lean}{\texttt{HodgeConjecture/Lemmas/AlgebraicTopology/SingularExcisionField.lean}}.

This file transports the integral subdivision homotopy to rational simplicial chains.  The
transport is constructed directly from the universal bases of the two simplicial chain
complexes.  In particular, it does not assume a universal-coefficient theorem.

\begin{definition}[AlgebraicTopology.Singular.RationalSimplicialChainComplex]\label{def:AlgebraicTopology.Singular.RationalSimplicialChainComplex}
\lean{AlgebraicTopology.Singular.RationalSimplicialChainComplex}\leanok
Rational chains on a simplicial set.

\noindent\emph{Lean:} \texttt{abbrev AlgebraicTopology.Singular.RationalSimplicialChainComplex}, \href{https://github.com/Paul-Lez/HodgeConjecture/blob/main/HodgeConjecture/Lemmas/AlgebraicTopology/SingularExcisionField.lean#L38}{\texttt{HodgeConjecture/Lemmas/AlgebraicTopology/SingularExcisionField.lean:38}}.
\end{definition}

\begin{definition}[AlgebraicTopology.Singular.integralToRationalChainComponent]\label{def:AlgebraicTopology.Singular.integralToRationalChainComponent}
\lean{AlgebraicTopology.Singular.integralToRationalChainComponent}\leanok
\uses{def:AlgebraicTopology.Singular.RationalSimplicialChainComplex}
The coefficient map from integral to rational chains in one degree.

\noindent\emph{Lean:} \texttt{def AlgebraicTopology.Singular.integralToRationalChainComponent}, \href{https://github.com/Paul-Lez/HodgeConjecture/blob/main/HodgeConjecture/Lemmas/AlgebraicTopology/SingularExcisionField.lean#L49}{\texttt{HodgeConjecture/Lemmas/AlgebraicTopology/SingularExcisionField.lean:49}}.
\end{definition}

\begin{definition}[AlgebraicTopology.Singular.rationalizeSimplicialChainComponent]\label{def:AlgebraicTopology.Singular.rationalizeSimplicialChainComponent}
\lean{AlgebraicTopology.Singular.rationalizeSimplicialChainComponent}\leanok
\uses{def:AlgebraicTopology.Singular.integralToRationalChainComponent}
Extend an integral map between free simplicial-chain groups rational-linearly.

\noindent\emph{Lean:} \texttt{def AlgebraicTopology.Singular.rationalizeSimplicialChainComponent}, \href{https://github.com/Paul-Lez/HodgeConjecture/blob/main/HodgeConjecture/Lemmas/AlgebraicTopology/SingularExcisionField.lean#L67}{\texttt{HodgeConjecture/Lemmas/AlgebraicTopology/SingularExcisionField.lean:67}}.
\end{definition}

\begin{definition}[AlgebraicTopology.Singular.rationalizeSimplicialChainMap]\label{def:AlgebraicTopology.Singular.rationalizeSimplicialChainMap}
\lean{AlgebraicTopology.Singular.rationalizeSimplicialChainMap}\leanok
\uses{def:AlgebraicTopology.Singular.rationalizeSimplicialChainComponent}
Rationalization of a chain map between integral simplicial chains.

\noindent\emph{Lean:} \texttt{def AlgebraicTopology.Singular.rationalizeSimplicialChainMap}, \href{https://github.com/Paul-Lez/HodgeConjecture/blob/main/HodgeConjecture/Lemmas/AlgebraicTopology/SingularExcisionField.lean#L236}{\texttt{HodgeConjecture/Lemmas/AlgebraicTopology/SingularExcisionField.lean:236}}.
\end{definition}

\begin{definition}[AlgebraicTopology.Singular.rationalizeSimplicialChainHomotopy]\label{def:AlgebraicTopology.Singular.rationalizeSimplicialChainHomotopy}
\lean{AlgebraicTopology.Singular.rationalizeSimplicialChainHomotopy}\leanok
\uses{def:AlgebraicTopology.Singular.rationalizeSimplicialChainMap}
Rationalization preserves a chain homotopy between integral simplicial chain maps.

\noindent\emph{Lean:} \texttt{def AlgebraicTopology.Singular.rationalizeSimplicialChainHomotopy}, \href{https://github.com/Paul-Lez/HodgeConjecture/blob/main/HodgeConjecture/Lemmas/AlgebraicTopology/SingularExcisionField.lean#L287}{\texttt{HodgeConjecture/Lemmas/AlgebraicTopology/SingularExcisionField.lean:287}}.
\end{definition}

\begin{definition}[AlgebraicTopology.Singular.rationalizeSimplicialChainHomotopyEquiv]\label{def:AlgebraicTopology.Singular.rationalizeSimplicialChainHomotopyEquiv}
\lean{AlgebraicTopology.Singular.rationalizeSimplicialChainHomotopyEquiv}\leanok
\uses{def:AlgebraicTopology.Singular.rationalizeSimplicialChainHomotopy}
Rationalization preserves a chain-homotopy equivalence of simplicial chain complexes.

\noindent\emph{Lean:} \texttt{def AlgebraicTopology.Singular.rationalizeSimplicialChainHomotopyEquiv}, \href{https://github.com/Paul-Lez/HodgeConjecture/blob/main/HodgeConjecture/Lemmas/AlgebraicTopology/SingularExcisionField.lean#L322}{\texttt{HodgeConjecture/Lemmas/AlgebraicTopology/SingularExcisionField.lean:322}}.
\end{definition}

\begin{definition}[AlgebraicTopology.Singular.CoverSmallRationalSingularChainComplex]\label{def:AlgebraicTopology.Singular.CoverSmallRationalSingularChainComplex}
\lean{AlgebraicTopology.Singular.CoverSmallRationalSingularChainComplex}\leanok
\uses{def:AlgebraicTopology.Singular.coverSmallSingularSubcomplex}
Rational chains generated by singular simplices subordinate to one member of a cover.

\noindent\emph{Lean:} \texttt{abbrev AlgebraicTopology.Singular.CoverSmallRationalSingularChainComplex}, \href{https://github.com/Paul-Lez/HodgeConjecture/blob/main/HodgeConjecture/Lemmas/AlgebraicTopology/SingularExcisionField.lean#L376}{\texttt{HodgeConjecture/Lemmas/AlgebraicTopology/SingularExcisionField.lean:376}}.
\end{definition}

\begin{definition}[AlgebraicTopology.Singular.coverSmallRationalSingularChainInclusion]\label{def:AlgebraicTopology.Singular.coverSmallRationalSingularChainInclusion}
\lean{AlgebraicTopology.Singular.coverSmallRationalSingularChainInclusion}\leanok
\uses{def:AlgebraicTopology.Singular.CoverSmallRationalSingularChainComplex}
Inclusion of cover-small rational singular chains into all rational singular chains.

\noindent\emph{Lean:} \texttt{def AlgebraicTopology.Singular.coverSmallRationalSingularChainInclusion}, \href{https://github.com/Paul-Lez/HodgeConjecture/blob/main/HodgeConjecture/Lemmas/AlgebraicTopology/SingularExcisionField.lean#L381}{\texttt{HodgeConjecture/Lemmas/AlgebraicTopology/SingularExcisionField.lean:381}}.
\end{definition}

\begin{definition}[AlgebraicTopology.Singular.coverSmallRationalChainHomotopyEquiv\_of\_openCover]\label{def:AlgebraicTopology.Singular.coverSmallRationalChainHomotopyEquiv_of_openCover}
\lean{AlgebraicTopology.Singular.coverSmallRationalChainHomotopyEquiv_of_openCover}\leanok
\uses{def:AlgebraicTopology.Singular.coverSmallChainHomotopyEquiv_of_openCover, def:AlgebraicTopology.Singular.coverSmallRationalSingularChainInclusion, def:AlgebraicTopology.Singular.rationalizeSimplicialChainHomotopyEquiv}
A selected rational small-chain homotopy equivalence for an open cover.

\noindent\emph{Lean:} \texttt{def AlgebraicTopology.Singular.coverSmallRationalChainHomotopyEquiv\_of\_openCover}, \href{https://github.com/Paul-Lez/HodgeConjecture/blob/main/HodgeConjecture/Lemmas/AlgebraicTopology/SingularExcisionField.lean#L413}{\texttt{HodgeConjecture/Lemmas/AlgebraicTopology/SingularExcisionField.lean:413}}.
\end{definition}

\section{Relative excision for an open neighborhood of a point}\label{sec:HodgeConjecture.Lemmas.AlgebraicTopology.RelativePairExcision}
\noindent\emph{File:} \href{https://github.com/Paul-Lez/HodgeConjecture/blob/main/HodgeConjecture/Lemmas/AlgebraicTopology/RelativePairExcision.lean}{\texttt{HodgeConjecture/Lemmas/AlgebraicTopology/RelativePairExcision.lean}}.

This file proves the point-complement form of singular excision needed for local fundamental
classes.  If \texttt{U} is an open neighborhood of \texttt{x} in a \texttt{T₁} space \texttt{X}, inclusion induces an
isomorphism

\texttt{Hₙ(U, U ∖ \{x\}; ℚ) ≅ Hₙ(X, X ∖ \{x\}; ℚ)}.

The proof uses the already established affine-subdivision small-chain theorem for the two-set
cover \texttt{U, X ∖ \{x\}}.  No excision or homology-comparison hypothesis is assumed.

\begin{definition}[AlgebraicTopology.Singular.neighborhoodPointComplementPair]\label{def:AlgebraicTopology.Singular.neighborhoodPointComplementPair}
\lean{AlgebraicTopology.Singular.neighborhoodPointComplementPair}\leanok
The pair \texttt{(U, U ∖ \{x\})} for a subset \texttt{U} containing \texttt{x}.

\noindent\emph{Lean:} \texttt{abbrev AlgebraicTopology.Singular.neighborhoodPointComplementPair}, \href{https://github.com/Paul-Lez/HodgeConjecture/blob/main/HodgeConjecture/Lemmas/AlgebraicTopology/RelativePairExcision.lean#L47}{\texttt{HodgeConjecture/Lemmas/AlgebraicTopology/RelativePairExcision.lean:47}}.
\end{definition}

\begin{definition}[AlgebraicTopology.Singular.neighborhoodPointComplementPairMap]\label{def:AlgebraicTopology.Singular.neighborhoodPointComplementPairMap}
\lean{AlgebraicTopology.Singular.neighborhoodPointComplementPairMap}\leanok
\uses{def:AlgebraicTopology.Singular.neighborhoodPointComplementPair, def:AlgebraicTopology.Singular.pointComplementPair}
The inclusion from a neighborhood point-complement pair into the ambient point-complement
pair.

\noindent\emph{Lean:} \texttt{def AlgebraicTopology.Singular.neighborhoodPointComplementPairMap}, \href{https://github.com/Paul-Lez/HodgeConjecture/blob/main/HodgeConjecture/Lemmas/AlgebraicTopology/RelativePairExcision.lean#L51}{\texttt{HodgeConjecture/Lemmas/AlgebraicTopology/RelativePairExcision.lean:51}}.
\end{definition}

\begin{definition}[AlgebraicTopology.Singular.pointExcisionCover]\label{def:AlgebraicTopology.Singular.pointExcisionCover}
\lean{AlgebraicTopology.Singular.pointExcisionCover}\leanok
The two-set cover used in point excision: the neighborhood and the complement of the
point.

\noindent\emph{Lean:} \texttt{def AlgebraicTopology.Singular.pointExcisionCover}, \href{https://github.com/Paul-Lez/HodgeConjecture/blob/main/HodgeConjecture/Lemmas/AlgebraicTopology/RelativePairExcision.lean#L61}{\texttt{HodgeConjecture/Lemmas/AlgebraicTopology/RelativePairExcision.lean:61}}.
\end{definition}

\begin{definition}[AlgebraicTopology.Singular.PointExcisionSmallChainComplex]\label{def:AlgebraicTopology.Singular.PointExcisionSmallChainComplex}
\lean{AlgebraicTopology.Singular.PointExcisionSmallChainComplex}\leanok
\uses{def:AlgebraicTopology.Singular.CoverSmallRationalSingularChainComplex, def:AlgebraicTopology.Singular.pointExcisionCover}
Rational chains on the small singular subcomplex for the point-excision cover.

\noindent\emph{Lean:} \texttt{abbrev AlgebraicTopology.Singular.PointExcisionSmallChainComplex}, \href{https://github.com/Paul-Lez/HodgeConjecture/blob/main/HodgeConjecture/Lemmas/AlgebraicTopology/RelativePairExcision.lean#L82}{\texttt{HodgeConjecture/Lemmas/AlgebraicTopology/RelativePairExcision.lean:82}}.
\end{definition}

\begin{definition}[AlgebraicTopology.Singular.pointComplementAmbientInclusion]\label{def:AlgebraicTopology.Singular.pointComplementAmbientInclusion}
\lean{AlgebraicTopology.Singular.pointComplementAmbientInclusion}\leanok
\uses{def:AlgebraicTopology.Singular.pointComplementPair}
The inclusion of the point complement into the ambient space, with its codomain fixed
explicitly to avoid transports between definitionally equal presentations of the subspace.

\noindent\emph{Lean:} \texttt{def AlgebraicTopology.Singular.pointComplementAmbientInclusion}, \href{https://github.com/Paul-Lez/HodgeConjecture/blob/main/HodgeConjecture/Lemmas/AlgebraicTopology/RelativePairExcision.lean#L87}{\texttt{HodgeConjecture/Lemmas/AlgebraicTopology/RelativePairExcision.lean:87}}.
\end{definition}

\begin{definition}[AlgebraicTopology.Singular.pointComplementToFalseCoverMember]\label{def:AlgebraicTopology.Singular.pointComplementToFalseCoverMember}
\lean{AlgebraicTopology.Singular.pointComplementToFalseCoverMember}\leanok
\uses{def:AlgebraicTopology.Singular.pointComplementAmbientInclusion, def:AlgebraicTopology.Singular.pointExcisionCover}
The identity-on-points map from the point complement to the \texttt{false} member of the
point-excision cover.

\noindent\emph{Lean:} \texttt{def AlgebraicTopology.Singular.pointComplementToFalseCoverMember}, \href{https://github.com/Paul-Lez/HodgeConjecture/blob/main/HodgeConjecture/Lemmas/AlgebraicTopology/RelativePairExcision.lean#L93}{\texttt{HodgeConjecture/Lemmas/AlgebraicTopology/RelativePairExcision.lean:93}}.
\end{definition}

\begin{definition}[AlgebraicTopology.Singular.pointComplementToExcisionSmallSingularSet]\label{def:AlgebraicTopology.Singular.pointComplementToExcisionSmallSingularSet}
\lean{AlgebraicTopology.Singular.pointComplementToExcisionSmallSingularSet}\leanok
\uses{def:AlgebraicTopology.Singular.coverSmallSingularSubcomplex, def:AlgebraicTopology.Singular.pointComplementToFalseCoverMember}
The chain map from the point complement into the small chains.

\noindent\emph{Lean:} \texttt{def AlgebraicTopology.Singular.pointComplementToExcisionSmallSingularSet}, \href{https://github.com/Paul-Lez/HodgeConjecture/blob/main/HodgeConjecture/Lemmas/AlgebraicTopology/RelativePairExcision.lean#L101}{\texttt{HodgeConjecture/Lemmas/AlgebraicTopology/RelativePairExcision.lean:101}}.
\end{definition}

\begin{definition}[AlgebraicTopology.Singular.pointComplementToExcisionSmallChains]\label{def:AlgebraicTopology.Singular.pointComplementToExcisionSmallChains}
\lean{AlgebraicTopology.Singular.pointComplementToExcisionSmallChains}\leanok
\uses{def:AlgebraicTopology.Singular.PointExcisionSmallChainComplex, def:AlgebraicTopology.Singular.pointComplementToExcisionSmallSingularSet}
\noindent\emph{Lean:} \texttt{def AlgebraicTopology.Singular.pointComplementToExcisionSmallChains}, \href{https://github.com/Paul-Lez/HodgeConjecture/blob/main/HodgeConjecture/Lemmas/AlgebraicTopology/RelativePairExcision.lean#L127}{\texttt{HodgeConjecture/Lemmas/AlgebraicTopology/RelativePairExcision.lean:127}}.
\end{definition}

\begin{definition}[AlgebraicTopology.Singular.pointExcisionSmallChainArrow]\label{def:AlgebraicTopology.Singular.pointExcisionSmallChainArrow}
\lean{AlgebraicTopology.Singular.pointExcisionSmallChainArrow}\leanok
\uses{def:AlgebraicTopology.Singular.ChainCategory, def:AlgebraicTopology.Singular.pointComplementToExcisionSmallChains}
The arrow whose cokernel is the small relative chain complex for point excision.

\noindent\emph{Lean:} \texttt{def AlgebraicTopology.Singular.pointExcisionSmallChainArrow}, \href{https://github.com/Paul-Lez/HodgeConjecture/blob/main/HodgeConjecture/Lemmas/AlgebraicTopology/RelativePairExcision.lean#L135}{\texttt{HodgeConjecture/Lemmas/AlgebraicTopology/RelativePairExcision.lean:135}}.
\end{definition}

\begin{definition}[AlgebraicTopology.Singular.PointExcisionSmallRelativeChainComplex]\label{def:AlgebraicTopology.Singular.PointExcisionSmallRelativeChainComplex}
\lean{AlgebraicTopology.Singular.PointExcisionSmallRelativeChainComplex}\leanok
\uses{def:AlgebraicTopology.Singular.pointExcisionSmallChainArrow}
The relative quotient of point-excision-small chains by chains avoiding the point.

\noindent\emph{Lean:} \texttt{abbrev AlgebraicTopology.Singular.PointExcisionSmallRelativeChainComplex}, \href{https://github.com/Paul-Lez/HodgeConjecture/blob/main/HodgeConjecture/Lemmas/AlgebraicTopology/RelativePairExcision.lean#L140}{\texttt{HodgeConjecture/Lemmas/AlgebraicTopology/RelativePairExcision.lean:140}}.
\end{definition}

\begin{definition}[AlgebraicTopology.Singular.PointExcisionSmallSimplex.AvoidsPoint]\label{def:AlgebraicTopology.Singular.PointExcisionSmallSimplex.AvoidsPoint}
\lean{AlgebraicTopology.Singular.PointExcisionSmallSimplex.AvoidsPoint}\leanok
\uses{def:AlgebraicTopology.Singular.coverSmallSingularSubcomplex, def:AlgebraicTopology.Singular.pointComplementAmbientInclusion, def:AlgebraicTopology.Singular.pointExcisionCover}
A small simplex lies away from \texttt{x} when it comes from the point-complement member of the
cover.

\noindent\emph{Lean:} \texttt{def AlgebraicTopology.Singular.PointExcisionSmallSimplex.AvoidsPoint}, \href{https://github.com/Paul-Lez/HodgeConjecture/blob/main/HodgeConjecture/Lemmas/AlgebraicTopology/RelativePairExcision.lean#L145}{\texttt{HodgeConjecture/Lemmas/AlgebraicTopology/RelativePairExcision.lean:145}}.
\end{definition}

\begin{definition}[AlgebraicTopology.Singular.PointExcisionSmallSimplex.neighborhoodPreimage]\label{def:AlgebraicTopology.Singular.PointExcisionSmallSimplex.neighborhoodPreimage}
\lean{AlgebraicTopology.Singular.PointExcisionSmallSimplex.neighborhoodPreimage}\leanok
\uses{def:AlgebraicTopology.Singular.PointExcisionSmallSimplex.AvoidsPoint}
The chosen neighborhood lift of a small simplex which does not avoid \texttt{x}.

\noindent\emph{Lean:} \texttt{def AlgebraicTopology.Singular.PointExcisionSmallSimplex.neighborhoodPreimage}, \href{https://github.com/Paul-Lez/HodgeConjecture/blob/main/HodgeConjecture/Lemmas/AlgebraicTopology/RelativePairExcision.lean#L189}{\texttt{HodgeConjecture/Lemmas/AlgebraicTopology/RelativePairExcision.lean:189}}.
\end{definition}

\begin{definition}[AlgebraicTopology.Singular.neighborhoodPointComplementSimplexOfAvoids]\label{def:AlgebraicTopology.Singular.neighborhoodPointComplementSimplexOfAvoids}
\lean{AlgebraicTopology.Singular.neighborhoodPointComplementSimplexOfAvoids}\leanok
\uses{def:AlgebraicTopology.Singular.neighborhoodPointComplementPair}
A singular simplex of \texttt{U} which avoids \texttt{x} lifts to the subspace \texttt{U ∖ \{x\}}.

\noindent\emph{Lean:} \texttt{def AlgebraicTopology.Singular.neighborhoodPointComplementSimplexOfAvoids}, \href{https://github.com/Paul-Lez/HodgeConjecture/blob/main/HodgeConjecture/Lemmas/AlgebraicTopology/RelativePairExcision.lean#L206}{\texttt{HodgeConjecture/Lemmas/AlgebraicTopology/RelativePairExcision.lean:206}}.
\end{definition}

\begin{definition}[AlgebraicTopology.Singular.neighborhoodSubspaceChainMap]\label{def:AlgebraicTopology.Singular.neighborhoodSubspaceChainMap}
\lean{AlgebraicTopology.Singular.neighborhoodSubspaceChainMap}\leanok
\uses{def:AlgebraicTopology.Singular.neighborhoodPointComplementPair}
A neighborhood simplex which is also represented in the point complement vanishes after
projection to the neighborhood relative chain complex.

\noindent\emph{Lean:} \texttt{abbrev AlgebraicTopology.Singular.neighborhoodSubspaceChainMap}, \href{https://github.com/Paul-Lez/HodgeConjecture/blob/main/HodgeConjecture/Lemmas/AlgebraicTopology/RelativePairExcision.lean#L250}{\texttt{HodgeConjecture/Lemmas/AlgebraicTopology/RelativePairExcision.lean:250}}.
\end{definition}

\begin{definition}[AlgebraicTopology.Singular.neighborhoodRelativeProjection]\label{def:AlgebraicTopology.Singular.neighborhoodRelativeProjection}
\lean{AlgebraicTopology.Singular.neighborhoodRelativeProjection}\leanok
\uses{def:AlgebraicTopology.Singular.neighborhoodPointComplementPair, def:AlgebraicTopology.Singular.relativeChainProjection}
The relative-chain projection for the neighborhood pair, with the absolute neighborhood
chain complex exposed as its source.

\noindent\emph{Lean:} \texttt{abbrev AlgebraicTopology.Singular.neighborhoodRelativeProjection}, \href{https://github.com/Paul-Lez/HodgeConjecture/blob/main/HodgeConjecture/Lemmas/AlgebraicTopology/RelativePairExcision.lean#L259}{\texttt{HodgeConjecture/Lemmas/AlgebraicTopology/RelativePairExcision.lean:259}}.
\end{definition}

\begin{definition}[AlgebraicTopology.Singular.neighborhoodRelativeProjectionComponent]\label{def:AlgebraicTopology.Singular.neighborhoodRelativeProjectionComponent}
\lean{AlgebraicTopology.Singular.neighborhoodRelativeProjectionComponent}\leanok
\uses{def:AlgebraicTopology.Singular.neighborhoodRelativeProjection}
\noindent\emph{Lean:} \texttt{def AlgebraicTopology.Singular.neighborhoodRelativeProjectionComponent}, \href{https://github.com/Paul-Lez/HodgeConjecture/blob/main/HodgeConjecture/Lemmas/AlgebraicTopology/RelativePairExcision.lean#L266}{\texttt{HodgeConjecture/Lemmas/AlgebraicTopology/RelativePairExcision.lean:266}}.
\end{definition}

\begin{definition}[AlgebraicTopology.Singular.pointExcisionSmallToNeighborhoodRelativeComponent]\label{def:AlgebraicTopology.Singular.pointExcisionSmallToNeighborhoodRelativeComponent}
\lean{AlgebraicTopology.Singular.pointExcisionSmallToNeighborhoodRelativeComponent}\leanok
\uses{def:AlgebraicTopology.Singular.PointExcisionSmallChainComplex, def:AlgebraicTopology.Singular.PointExcisionSmallSimplex.neighborhoodPreimage, def:AlgebraicTopology.Singular.neighborhoodRelativeProjectionComponent}
On a basis simplex, send a point-avoiding simplex to zero and otherwise lift it to the
neighborhood before projecting to relative chains.

\noindent\emph{Lean:} \texttt{def AlgebraicTopology.Singular.pointExcisionSmallToNeighborhoodRelativeComponent}, \href{https://github.com/Paul-Lez/HodgeConjecture/blob/main/HodgeConjecture/Lemmas/AlgebraicTopology/RelativePairExcision.lean#L325}{\texttt{HodgeConjecture/Lemmas/AlgebraicTopology/RelativePairExcision.lean:325}}.
\end{definition}

\begin{definition}[AlgebraicTopology.Singular.pointExcisionSmallToNeighborhoodRelativeChainMap]\label{def:AlgebraicTopology.Singular.pointExcisionSmallToNeighborhoodRelativeChainMap}
\lean{AlgebraicTopology.Singular.pointExcisionSmallToNeighborhoodRelativeChainMap}\leanok
\uses{def:AlgebraicTopology.Singular.neighborhoodPointComplementSimplexOfAvoids, def:AlgebraicTopology.Singular.neighborhoodSubspaceChainMap, def:AlgebraicTopology.Singular.pointExcisionSmallToNeighborhoodRelativeComponent}
The chain map from small relative chains to neighborhood relative chains obtained by
discarding simplices which avoid the distinguished point.

\noindent\emph{Lean:} \texttt{def AlgebraicTopology.Singular.pointExcisionSmallToNeighborhoodRelativeChainMap}, \href{https://github.com/Paul-Lez/HodgeConjecture/blob/main/HodgeConjecture/Lemmas/AlgebraicTopology/RelativePairExcision.lean#L402}{\texttt{HodgeConjecture/Lemmas/AlgebraicTopology/RelativePairExcision.lean:402}}.
\end{definition}

\begin{definition}[AlgebraicTopology.Singular.pointExcisionSmallRelativeProjection]\label{def:AlgebraicTopology.Singular.pointExcisionSmallRelativeProjection}
\lean{AlgebraicTopology.Singular.pointExcisionSmallRelativeProjection}\leanok
\uses{def:AlgebraicTopology.Singular.PointExcisionSmallRelativeChainComplex}
Projection from small chains to the small relative quotient.

\noindent\emph{Lean:} \texttt{abbrev AlgebraicTopology.Singular.pointExcisionSmallRelativeProjection}, \href{https://github.com/Paul-Lez/HodgeConjecture/blob/main/HodgeConjecture/Lemmas/AlgebraicTopology/RelativePairExcision.lean#L444}{\texttt{HodgeConjecture/Lemmas/AlgebraicTopology/RelativePairExcision.lean:444}}.
\end{definition}

\begin{definition}[AlgebraicTopology.Singular.pointExcisionSmallRelativeToNeighborhoodRelativeChainMap]\label{def:AlgebraicTopology.Singular.pointExcisionSmallRelativeToNeighborhoodRelativeChainMap}
\lean{AlgebraicTopology.Singular.pointExcisionSmallRelativeToNeighborhoodRelativeChainMap}\leanok
\uses{def:AlgebraicTopology.Singular.PointExcisionSmallRelativeChainComplex, def:AlgebraicTopology.Singular.pointExcisionSmallToNeighborhoodRelativeChainMap}
The chain map from the small relative quotient back to neighborhood relative chains.

\noindent\emph{Lean:} \texttt{def AlgebraicTopology.Singular.pointExcisionSmallRelativeToNeighborhoodRelativeChainMap}, \href{https://github.com/Paul-Lez/HodgeConjecture/blob/main/HodgeConjecture/Lemmas/AlgebraicTopology/RelativePairExcision.lean#L450}{\texttt{HodgeConjecture/Lemmas/AlgebraicTopology/RelativePairExcision.lean:450}}.
\end{definition}

\begin{definition}[AlgebraicTopology.Singular.neighborhoodPointComplementToPointComplement]\label{def:AlgebraicTopology.Singular.neighborhoodPointComplementToPointComplement}
\lean{AlgebraicTopology.Singular.neighborhoodPointComplementToPointComplement}\leanok
\uses{def:AlgebraicTopology.Singular.neighborhoodPointComplementPair, def:AlgebraicTopology.Singular.pointComplementPair}
The inclusion from \texttt{U ∖ \{x\}} to \texttt{X ∖ \{x\}}.

\noindent\emph{Lean:} \texttt{def AlgebraicTopology.Singular.neighborhoodPointComplementToPointComplement}, \href{https://github.com/Paul-Lez/HodgeConjecture/blob/main/HodgeConjecture/Lemmas/AlgebraicTopology/RelativePairExcision.lean#L467}{\texttt{HodgeConjecture/Lemmas/AlgebraicTopology/RelativePairExcision.lean:467}}.
\end{definition}

\begin{definition}[AlgebraicTopology.Singular.neighborhoodToTrueCoverMember]\label{def:AlgebraicTopology.Singular.neighborhoodToTrueCoverMember}
\lean{AlgebraicTopology.Singular.neighborhoodToTrueCoverMember}\leanok
\uses{def:AlgebraicTopology.Singular.pointExcisionCover, def:AlgebraicTopology.Singular.topologicalSubsetInclusion}
The identity-on-points map from the neighborhood to the \texttt{true} member of the
point-excision cover.

\noindent\emph{Lean:} \texttt{def AlgebraicTopology.Singular.neighborhoodToTrueCoverMember}, \href{https://github.com/Paul-Lez/HodgeConjecture/blob/main/HodgeConjecture/Lemmas/AlgebraicTopology/RelativePairExcision.lean#L473}{\texttt{HodgeConjecture/Lemmas/AlgebraicTopology/RelativePairExcision.lean:473}}.
\end{definition}

\begin{definition}[AlgebraicTopology.Singular.neighborhoodToPointExcisionSmallSingularSet]\label{def:AlgebraicTopology.Singular.neighborhoodToPointExcisionSmallSingularSet}
\lean{AlgebraicTopology.Singular.neighborhoodToPointExcisionSmallSingularSet}\leanok
\uses{def:AlgebraicTopology.Singular.coverSmallSingularSubcomplex, def:AlgebraicTopology.Singular.neighborhoodToTrueCoverMember}
The chain map from neighborhood chains to the point-excision-small chains.

\noindent\emph{Lean:} \texttt{def AlgebraicTopology.Singular.neighborhoodToPointExcisionSmallSingularSet}, \href{https://github.com/Paul-Lez/HodgeConjecture/blob/main/HodgeConjecture/Lemmas/AlgebraicTopology/RelativePairExcision.lean#L481}{\texttt{HodgeConjecture/Lemmas/AlgebraicTopology/RelativePairExcision.lean:481}}.
\end{definition}

\begin{definition}[AlgebraicTopology.Singular.neighborhoodToPointExcisionSmallChains]\label{def:AlgebraicTopology.Singular.neighborhoodToPointExcisionSmallChains}
\lean{AlgebraicTopology.Singular.neighborhoodToPointExcisionSmallChains}\leanok
\uses{def:AlgebraicTopology.Singular.PointExcisionSmallChainComplex, def:AlgebraicTopology.Singular.neighborhoodToPointExcisionSmallSingularSet}
\noindent\emph{Lean:} \texttt{def AlgebraicTopology.Singular.neighborhoodToPointExcisionSmallChains}, \href{https://github.com/Paul-Lez/HodgeConjecture/blob/main/HodgeConjecture/Lemmas/AlgebraicTopology/RelativePairExcision.lean#L507}{\texttt{HodgeConjecture/Lemmas/AlgebraicTopology/RelativePairExcision.lean:507}}.
\end{definition}

\begin{definition}[AlgebraicTopology.Singular.neighborhoodPointComplementToPointComplementChains]\label{def:AlgebraicTopology.Singular.neighborhoodPointComplementToPointComplementChains}
\lean{AlgebraicTopology.Singular.neighborhoodPointComplementToPointComplementChains}\leanok
\uses{def:AlgebraicTopology.Singular.neighborhoodPointComplementToPointComplement}
The chain map from \texttt{U ∖ \{x\}} to \texttt{X ∖ \{x\}}.

\noindent\emph{Lean:} \texttt{def AlgebraicTopology.Singular.neighborhoodPointComplementToPointComplementChains}, \href{https://github.com/Paul-Lez/HodgeConjecture/blob/main/HodgeConjecture/Lemmas/AlgebraicTopology/RelativePairExcision.lean#L514}{\texttt{HodgeConjecture/Lemmas/AlgebraicTopology/RelativePairExcision.lean:514}}.
\end{definition}

\begin{definition}[AlgebraicTopology.Singular.neighborhoodPairChainsToPointExcisionSmallArrow]\label{def:AlgebraicTopology.Singular.neighborhoodPairChainsToPointExcisionSmallArrow}
\lean{AlgebraicTopology.Singular.neighborhoodPairChainsToPointExcisionSmallArrow}\leanok
\uses{def:AlgebraicTopology.Singular.chainPairFunctor, def:AlgebraicTopology.Singular.neighborhoodPointComplementToPointComplementChains, def:AlgebraicTopology.Singular.neighborhoodSubspaceChainMap, def:AlgebraicTopology.Singular.neighborhoodToPointExcisionSmallChains, def:AlgebraicTopology.Singular.pointExcisionSmallChainArrow}
The morphism of chain-complex arrows from the neighborhood pair to the small pair.

\noindent\emph{Lean:} \texttt{def AlgebraicTopology.Singular.neighborhoodPairChainsToPointExcisionSmallArrow}, \href{https://github.com/Paul-Lez/HodgeConjecture/blob/main/HodgeConjecture/Lemmas/AlgebraicTopology/RelativePairExcision.lean#L554}{\texttt{HodgeConjecture/Lemmas/AlgebraicTopology/RelativePairExcision.lean:554}}.
\end{definition}

\begin{definition}[AlgebraicTopology.Singular.neighborhoodRelativeToPointExcisionSmallRelativeChainMap]\label{def:AlgebraicTopology.Singular.neighborhoodRelativeToPointExcisionSmallRelativeChainMap}
\lean{AlgebraicTopology.Singular.neighborhoodRelativeToPointExcisionSmallRelativeChainMap}\leanok
\uses{def:AlgebraicTopology.Singular.PointExcisionSmallRelativeChainComplex, def:AlgebraicTopology.Singular.neighborhoodPairChainsToPointExcisionSmallArrow, def:AlgebraicTopology.Singular.relativeChainFunctor}
Inclusion of the neighborhood pair into the small relative quotient.

\noindent\emph{Lean:} \texttt{def AlgebraicTopology.Singular.neighborhoodRelativeToPointExcisionSmallRelativeChainMap}, \href{https://github.com/Paul-Lez/HodgeConjecture/blob/main/HodgeConjecture/Lemmas/AlgebraicTopology/RelativePairExcision.lean#L563}{\texttt{HodgeConjecture/Lemmas/AlgebraicTopology/RelativePairExcision.lean:563}}.
\end{definition}

\begin{definition}[AlgebraicTopology.Singular.neighborhoodRelativePointExcisionSmallIso]\label{def:AlgebraicTopology.Singular.neighborhoodRelativePointExcisionSmallIso}
\lean{AlgebraicTopology.Singular.neighborhoodRelativePointExcisionSmallIso}\leanok
\uses{def:AlgebraicTopology.Singular.neighborhoodRelativeToPointExcisionSmallRelativeChainMap, def:AlgebraicTopology.Singular.pointExcisionSmallRelativeProjection, def:AlgebraicTopology.Singular.pointExcisionSmallRelativeToNeighborhoodRelativeChainMap}
Neighborhood-relative chains are isomorphic to the point-excision-small relative
quotient.

\noindent\emph{Lean:} \texttt{def AlgebraicTopology.Singular.neighborhoodRelativePointExcisionSmallIso}, \href{https://github.com/Paul-Lez/HodgeConjecture/blob/main/HodgeConjecture/Lemmas/AlgebraicTopology/RelativePairExcision.lean#L746}{\texttt{HodgeConjecture/Lemmas/AlgebraicTopology/RelativePairExcision.lean:746}}.
\end{definition}

\begin{definition}[AlgebraicTopology.Singular.ambientPointComplementSubspaceChainMap]\label{def:AlgebraicTopology.Singular.ambientPointComplementSubspaceChainMap}
\lean{AlgebraicTopology.Singular.ambientPointComplementSubspaceChainMap}\leanok
\uses{def:AlgebraicTopology.Singular.chainPairFunctor, def:AlgebraicTopology.Singular.pointComplementPair}
The chain inclusion for the ambient point-complement pair, with both singular chain
complexes exposed.

\noindent\emph{Lean:} \texttt{abbrev AlgebraicTopology.Singular.ambientPointComplementSubspaceChainMap}, \href{https://github.com/Paul-Lez/HodgeConjecture/blob/main/HodgeConjecture/Lemmas/AlgebraicTopology/RelativePairExcision.lean#L762}{\texttt{HodgeConjecture/Lemmas/AlgebraicTopology/RelativePairExcision.lean:762}}.
\end{definition}

\begin{definition}[AlgebraicTopology.Singular.pointComplementAmbientChainMap]\label{def:AlgebraicTopology.Singular.pointComplementAmbientChainMap}
\lean{AlgebraicTopology.Singular.pointComplementAmbientChainMap}\leanok
\uses{def:AlgebraicTopology.Singular.pointComplementAmbientInclusion}
The explicit singular-chain map underlying the ambient point-complement inclusion.

\noindent\emph{Lean:} \texttt{def AlgebraicTopology.Singular.pointComplementAmbientChainMap}, \href{https://github.com/Paul-Lez/HodgeConjecture/blob/main/HodgeConjecture/Lemmas/AlgebraicTopology/RelativePairExcision.lean#L770}{\texttt{HodgeConjecture/Lemmas/AlgebraicTopology/RelativePairExcision.lean:770}}.
\end{definition}

\begin{definition}[AlgebraicTopology.Singular.ambientPointComplementRelativeProjection]\label{def:AlgebraicTopology.Singular.ambientPointComplementRelativeProjection}
\lean{AlgebraicTopology.Singular.ambientPointComplementRelativeProjection}\leanok
\uses{def:AlgebraicTopology.Singular.pointComplementPair, def:AlgebraicTopology.Singular.relativeChainProjection}
The ambient relative-chain projection, with its absolute chain source exposed.

\noindent\emph{Lean:} \texttt{abbrev AlgebraicTopology.Singular.ambientPointComplementRelativeProjection}, \href{https://github.com/Paul-Lez/HodgeConjecture/blob/main/HodgeConjecture/Lemmas/AlgebraicTopology/RelativePairExcision.lean#L802}{\texttt{HodgeConjecture/Lemmas/AlgebraicTopology/RelativePairExcision.lean:802}}.
\end{definition}

\begin{definition}[AlgebraicTopology.Singular.pointComplementChainsToAmbientPairLeft]\label{def:AlgebraicTopology.Singular.pointComplementChainsToAmbientPairLeft}
\lean{AlgebraicTopology.Singular.pointComplementChainsToAmbientPairLeft}\leanok
\uses{def:AlgebraicTopology.Singular.chainPairFunctor, def:AlgebraicTopology.Singular.pointComplementPair}
The identity map on point-complement chains, with the codomain presented as the left
object of the ambient chain-pair arrow.

\noindent\emph{Lean:} \texttt{def AlgebraicTopology.Singular.pointComplementChainsToAmbientPairLeft}, \href{https://github.com/Paul-Lez/HodgeConjecture/blob/main/HodgeConjecture/Lemmas/AlgebraicTopology/RelativePairExcision.lean#L836}{\texttt{HodgeConjecture/Lemmas/AlgebraicTopology/RelativePairExcision.lean:836}}.
\end{definition}

\begin{definition}[AlgebraicTopology.Singular.ambientPairLeftToPointComplementChains]\label{def:AlgebraicTopology.Singular.ambientPairLeftToPointComplementChains}
\lean{AlgebraicTopology.Singular.ambientPairLeftToPointComplementChains}\leanok
\uses{def:AlgebraicTopology.Singular.chainPairFunctor, def:AlgebraicTopology.Singular.pointComplementPair}
The inverse identity map for \texttt{pointComplementChainsToAmbientPairLeft}.

\noindent\emph{Lean:} \texttt{def AlgebraicTopology.Singular.ambientPairLeftToPointComplementChains}, \href{https://github.com/Paul-Lez/HodgeConjecture/blob/main/HodgeConjecture/Lemmas/AlgebraicTopology/RelativePairExcision.lean#L844}{\texttt{HodgeConjecture/Lemmas/AlgebraicTopology/RelativePairExcision.lean:844}}.
\end{definition}

\begin{definition}[AlgebraicTopology.Singular.pointExcisionSmallToAmbientPairArrow]\label{def:AlgebraicTopology.Singular.pointExcisionSmallToAmbientPairArrow}
\lean{AlgebraicTopology.Singular.pointExcisionSmallToAmbientPairArrow}\leanok
\uses{def:AlgebraicTopology.Singular.ambientPointComplementSubspaceChainMap, def:AlgebraicTopology.Singular.coverSmallRationalSingularChainInclusion, def:AlgebraicTopology.Singular.pointComplementChainsToAmbientPairLeft, def:AlgebraicTopology.Singular.pointExcisionSmallChainArrow}
The morphism of chain-pair arrows from the point-excision-small pair to the ambient
point-complement pair.

\noindent\emph{Lean:} \texttt{def AlgebraicTopology.Singular.pointExcisionSmallToAmbientPairArrow}, \href{https://github.com/Paul-Lez/HodgeConjecture/blob/main/HodgeConjecture/Lemmas/AlgebraicTopology/RelativePairExcision.lean#L867}{\texttt{HodgeConjecture/Lemmas/AlgebraicTopology/RelativePairExcision.lean:867}}.
\end{definition}

\begin{definition}[AlgebraicTopology.Singular.pointExcisionSmallRelativeToAmbientRelativeChainMap]\label{def:AlgebraicTopology.Singular.pointExcisionSmallRelativeToAmbientRelativeChainMap}
\lean{AlgebraicTopology.Singular.pointExcisionSmallRelativeToAmbientRelativeChainMap}\leanok
\uses{def:AlgebraicTopology.Singular.PointExcisionSmallRelativeChainComplex, def:AlgebraicTopology.Singular.pointExcisionSmallToAmbientPairArrow, def:AlgebraicTopology.Singular.relativeChainFunctor}
The map from the point-excision-small relative quotient to ambient relative chains.

\noindent\emph{Lean:} \texttt{def AlgebraicTopology.Singular.pointExcisionSmallRelativeToAmbientRelativeChainMap}, \href{https://github.com/Paul-Lez/HodgeConjecture/blob/main/HodgeConjecture/Lemmas/AlgebraicTopology/RelativePairExcision.lean#L878}{\texttt{HodgeConjecture/Lemmas/AlgebraicTopology/RelativePairExcision.lean:878}}.
\end{definition}

\begin{definition}[AlgebraicTopology.Singular.pointExcisionSmallToAmbientCokernelSequenceHom]\label{def:AlgebraicTopology.Singular.pointExcisionSmallToAmbientCokernelSequenceHom}
\lean{AlgebraicTopology.Singular.pointExcisionSmallToAmbientCokernelSequenceHom}\leanok
\uses{def:AlgebraicTopology.Singular.ambientPointComplementRelativeProjection, def:AlgebraicTopology.Singular.pointExcisionSmallRelativeProjection, def:AlgebraicTopology.Singular.pointExcisionSmallRelativeToAmbientRelativeChainMap}
The morphism between the two cokernel short complexes used in point excision.

\noindent\emph{Lean:} \texttt{def AlgebraicTopology.Singular.pointExcisionSmallToAmbientCokernelSequenceHom}, \href{https://github.com/Paul-Lez/HodgeConjecture/blob/main/HodgeConjecture/Lemmas/AlgebraicTopology/RelativePairExcision.lean#L897}{\texttt{HodgeConjecture/Lemmas/AlgebraicTopology/RelativePairExcision.lean:897}}.
\end{definition}

\begin{definition}[AlgebraicTopology.Singular.neighborhoodAmbientChainMap]\label{def:AlgebraicTopology.Singular.neighborhoodAmbientChainMap}
\lean{AlgebraicTopology.Singular.neighborhoodAmbientChainMap}\leanok
\uses{def:AlgebraicTopology.Singular.topologicalSubsetInclusion}
The ordinary chain map induced by inclusion of the neighborhood into the ambient space.

\noindent\emph{Lean:} \texttt{def AlgebraicTopology.Singular.neighborhoodAmbientChainMap}, \href{https://github.com/Paul-Lez/HodgeConjecture/blob/main/HodgeConjecture/Lemmas/AlgebraicTopology/RelativePairExcision.lean#L952}{\texttt{HodgeConjecture/Lemmas/AlgebraicTopology/RelativePairExcision.lean:952}}.
\end{definition}

\section{Generator properties of the standard complex local class}\label{sec:HodgeConjecture.Lemmas.AlgebraicTopology.LocalFundamentalClassGenerator}
\noindent\emph{File:} \href{https://github.com/Paul-Lez/HodgeConjecture/blob/main/HodgeConjecture/Lemmas/AlgebraicTopology/LocalFundamentalClassGenerator.lean}{\texttt{HodgeConjecture/Lemmas/AlgebraicTopology/LocalFundamentalClassGenerator.lean}}.

The ordered real and imaginary coordinate homeomorphism identifies the standard complex local
class with the explicit real local class. This file records that this identification preserves
nonvanishing, linear independence, and generation of local homology.

These results isolate the remaining topological input for local purity: it suffices to prove that
the explicit real affine cycle generates the corresponding real relative homology group.

\begin{definition}[AlgebraicTopology.Singular.standardComplexLocalClassMul]\label{def:AlgebraicTopology.Singular.standardComplexLocalClassMul}
\lean{AlgebraicTopology.Singular.standardComplexLocalClassMul}\leanok
\uses{def:AlgebraicTopology.Singular.standardComplexRealRelativeHomologyIso, def:AlgebraicTopology.Singular.standardLocalClass}
\noindent\emph{Lean:} \texttt{def AlgebraicTopology.Singular.standardComplexLocalClassMul}, \href{https://github.com/Paul-Lez/HodgeConjecture/blob/main/HodgeConjecture/Lemmas/AlgebraicTopology/LocalFundamentalClassGenerator.lean#L39}{\texttt{HodgeConjecture/Lemmas/AlgebraicTopology/LocalFundamentalClassGenerator.lean:39}}.
\end{definition}

\begin{definition}[AlgebraicTopology.Singular.relativeHomologyDegreeCast]\label{def:AlgebraicTopology.Singular.relativeHomologyDegreeCast}
\lean{AlgebraicTopology.Singular.relativeHomologyDegreeCast}\leanok
\uses{def:AlgebraicTopology.Singular.RelativeHomology}
Transport a relative homology class along an equality of degrees.

\noindent\emph{Lean:} \texttt{def AlgebraicTopology.Singular.relativeHomologyDegreeCast}, \href{https://github.com/Paul-Lez/HodgeConjecture/blob/main/HodgeConjecture/Lemmas/AlgebraicTopology/LocalFundamentalClassGenerator.lean#L43}{\texttt{HodgeConjecture/Lemmas/AlgebraicTopology/LocalFundamentalClassGenerator.lean:43}}.
\end{definition}

\begin{definition}[AlgebraicTopology.Singular.standardRelativeChainProjectionZero]\label{def:AlgebraicTopology.Singular.standardRelativeChainProjectionZero}
\lean{AlgebraicTopology.Singular.standardRelativeChainProjectionZero}\leanok
\uses{def:AlgebraicTopology.Singular.relativeChainProjection, def:AlgebraicTopology.Singular.standardLocalRelativeChainComplex}
The relative-chain projection with its ambient source exposed through the pair functor.

\noindent\emph{Lean:} \texttt{def AlgebraicTopology.Singular.standardRelativeChainProjectionZero}, \href{https://github.com/Paul-Lez/HodgeConjecture/blob/main/HodgeConjecture/Lemmas/AlgebraicTopology/LocalFundamentalClassGenerator.lean#L141}{\texttt{HodgeConjecture/Lemmas/AlgebraicTopology/LocalFundamentalClassGenerator.lean:141}}.
\end{definition}

\begin{definition}[AlgebraicTopology.Singular.standardRelativeHomologyProjectionZeroIso]\label{def:AlgebraicTopology.Singular.standardRelativeHomologyProjectionZeroIso}
\lean{AlgebraicTopology.Singular.standardRelativeHomologyProjectionZeroIso}\leanok
\uses{def:AlgebraicTopology.Singular.RelativeHomology, def:AlgebraicTopology.Singular.standardRelativeChainProjectionZero}
\noindent\emph{Lean:} \texttt{def AlgebraicTopology.Singular.standardRelativeHomologyProjectionZeroIso}, \href{https://github.com/Paul-Lez/HodgeConjecture/blob/main/HodgeConjecture/Lemmas/AlgebraicTopology/LocalFundamentalClassGenerator.lean#L151}{\texttt{HodgeConjecture/Lemmas/AlgebraicTopology/LocalFundamentalClassGenerator.lean:151}}.
\end{definition}

\begin{definition}[AlgebraicTopology.Singular.standardLocalRelativeHomologyZeroIso]\label{def:AlgebraicTopology.Singular.standardLocalRelativeHomologyZeroIso}
\lean{AlgebraicTopology.Singular.standardLocalRelativeHomologyZeroIso}\leanok
\uses{def:AlgebraicTopology.Singular.standardRelativeHomologyProjectionZeroIso}
The zero-dimensional standard relative local homology group is canonically one-dimensional.

\noindent\emph{Lean:} \texttt{def AlgebraicTopology.Singular.standardLocalRelativeHomologyZeroIso}, \href{https://github.com/Paul-Lez/HodgeConjecture/blob/main/HodgeConjecture/Lemmas/AlgebraicTopology/LocalFundamentalClassGenerator.lean#L157}{\texttt{HodgeConjecture/Lemmas/AlgebraicTopology/LocalFundamentalClassGenerator.lean:157}}.
\end{definition}

\begin{definition}[AlgebraicTopology.Singular.standardAmbientCycleZero]\label{def:AlgebraicTopology.Singular.standardAmbientCycleZero}
\lean{AlgebraicTopology.Singular.standardAmbientCycleZero}\leanok
\uses{def:AlgebraicTopology.Singular.standardAmbientSimplexChain}
The ambient zero-cycle represented by the unique singular point of \texttt{ℝ⁰}.

\noindent\emph{Lean:} \texttt{def AlgebraicTopology.Singular.standardAmbientCycleZero}, \href{https://github.com/Paul-Lez/HodgeConjecture/blob/main/HodgeConjecture/Lemmas/AlgebraicTopology/LocalFundamentalClassGenerator.lean#L169}{\texttt{HodgeConjecture/Lemmas/AlgebraicTopology/LocalFundamentalClassGenerator.lean:169}}.
\end{definition}

\begin{definition}[AlgebraicTopology.Singular.standardAmbientClassZero]\label{def:AlgebraicTopology.Singular.standardAmbientClassZero}
\lean{AlgebraicTopology.Singular.standardAmbientClassZero}\leanok
\uses{def:AlgebraicTopology.Singular.standardAmbientCycleZero}
The class of the unique singular point in the ambient homology of \texttt{ℝ⁰}.

\noindent\emph{Lean:} \texttt{def AlgebraicTopology.Singular.standardAmbientClassZero}, \href{https://github.com/Paul-Lez/HodgeConjecture/blob/main/HodgeConjecture/Lemmas/AlgebraicTopology/LocalFundamentalClassGenerator.lean#L178}{\texttt{HodgeConjecture/Lemmas/AlgebraicTopology/LocalFundamentalClassGenerator.lean:178}}.
\end{definition}

\begin{definition}[AlgebraicTopology.Singular.standardAmbientHomologyZeroε]\label{def:AlgebraicTopology.Singular.standardAmbientHomologyZeroε}
\lean{AlgebraicTopology.Singular.standardAmbientHomologyZeroε}\leanok
\uses{def:AlgebraicTopology.Singular.chainPairFunctor, def:AlgebraicTopology.Singular.standardPuncturedPair}
The ordinary \texttt{H₀} augmentation, with its source exposed through the pair functor.

\noindent\emph{Lean:} \texttt{def AlgebraicTopology.Singular.standardAmbientHomologyZeroε}, \href{https://github.com/Paul-Lez/HodgeConjecture/blob/main/HodgeConjecture/Lemmas/AlgebraicTopology/LocalFundamentalClassGenerator.lean#L215}{\texttt{HodgeConjecture/Lemmas/AlgebraicTopology/LocalFundamentalClassGenerator.lean:215}}.
\end{definition}

\begin{definition}[AlgebraicTopology.Singular.standardLocalRelativeHomologyZeroε]\label{def:AlgebraicTopology.Singular.standardLocalRelativeHomologyZeroε}
\lean{AlgebraicTopology.Singular.standardLocalRelativeHomologyZeroε}\leanok
\uses{def:AlgebraicTopology.Singular.standardAmbientHomologyZeroε, def:AlgebraicTopology.Singular.standardRelativeHomologyProjectionZeroIso}
The augmentation on zero-dimensional relative local homology.

\noindent\emph{Lean:} \texttt{def AlgebraicTopology.Singular.standardLocalRelativeHomologyZeroε}, \href{https://github.com/Paul-Lez/HodgeConjecture/blob/main/HodgeConjecture/Lemmas/AlgebraicTopology/LocalFundamentalClassGenerator.lean#L221}{\texttt{HodgeConjecture/Lemmas/AlgebraicTopology/LocalFundamentalClassGenerator.lean:221}}.
\end{definition}

\section{Euclidean local homology}\label{sec:HodgeConjecture.Lemmas.AlgebraicTopology.EuclideanLocalHomology}
\noindent\emph{File:} \href{https://github.com/Paul-Lez/HodgeConjecture/blob/main/HodgeConjecture/Lemmas/AlgebraicTopology/EuclideanLocalHomology.lean}{\texttt{HodgeConjecture/Lemmas/AlgebraicTopology/EuclideanLocalHomology.lean}}.

This file develops the long exact sequence of the standard punctured Euclidean pair and computes
the connecting morphism on the explicit standard local class. These are the chain-level
prerequisites for proving that the standard class generates local homology.

\begin{definition}[AlgebraicTopology.Singular.relativeSingularChainShortComplex]\label{def:AlgebraicTopology.Singular.relativeSingularChainShortComplex}
\lean{AlgebraicTopology.Singular.relativeSingularChainShortComplex}\leanok
\uses{def:AlgebraicTopology.Singular.relativeChainProjection}
The short complex of subspace, ambient, and relative rational singular chains.

\noindent\emph{Lean:} \texttt{def AlgebraicTopology.Singular.relativeSingularChainShortComplex}, \href{https://github.com/Paul-Lez/HodgeConjecture/blob/main/HodgeConjecture/Lemmas/AlgebraicTopology/EuclideanLocalHomology.lean#L44}{\texttt{HodgeConjecture/Lemmas/AlgebraicTopology/EuclideanLocalHomology.lean:44}}.
\end{definition}

\begin{definition}[AlgebraicTopology.Singular.relativeSingularBoundary]\label{def:AlgebraicTopology.Singular.relativeSingularBoundary}
\lean{AlgebraicTopology.Singular.relativeSingularBoundary}\leanok
\uses{def:AlgebraicTopology.Singular.Homology, def:AlgebraicTopology.Singular.RelativeHomology, def:AlgebraicTopology.Singular.relativeSingularChainShortComplex}
The connecting map from relative homology in degree \texttt{n + 1} to subspace homology in degree
\texttt{n}.

\noindent\emph{Lean:} \texttt{def AlgebraicTopology.Singular.relativeSingularBoundary}, \href{https://github.com/Paul-Lez/HodgeConjecture/blob/main/HodgeConjecture/Lemmas/AlgebraicTopology/EuclideanLocalHomology.lean#L77}{\texttt{HodgeConjecture/Lemmas/AlgebraicTopology/EuclideanLocalHomology.lean:77}}.
\end{definition}

\begin{definition}[AlgebraicTopology.Singular.standardSubspaceBoundaryChain]\label{def:AlgebraicTopology.Singular.standardSubspaceBoundaryChain}
\lean{AlgebraicTopology.Singular.standardSubspaceBoundaryChain}\leanok
\uses{def:AlgebraicTopology.Singular.standardSubspaceFaceChain}
The alternating sum of the faces of the standard affine \texttt{(n + 1)}-simplex, regarded as a
chain in punctured Euclidean space.

\noindent\emph{Lean:} \texttt{def AlgebraicTopology.Singular.standardSubspaceBoundaryChain}, \href{https://github.com/Paul-Lez/HodgeConjecture/blob/main/HodgeConjecture/Lemmas/AlgebraicTopology/EuclideanLocalHomology.lean#L120}{\texttt{HodgeConjecture/Lemmas/AlgebraicTopology/EuclideanLocalHomology.lean:120}}.
\end{definition}

\begin{definition}[AlgebraicTopology.Singular.standardPuncturedBoundaryCycle]\label{def:AlgebraicTopology.Singular.standardPuncturedBoundaryCycle}
\lean{AlgebraicTopology.Singular.standardPuncturedBoundaryCycle}\leanok
\uses{def:AlgebraicTopology.Singular.standardAmbientFaceChain, def:AlgebraicTopology.Singular.standardAmbientSimplexChain, def:AlgebraicTopology.Singular.standardSubspaceBoundaryChain}
The boundary of the standard affine simplex as a cycle in punctured Euclidean space.

\noindent\emph{Lean:} \texttt{def AlgebraicTopology.Singular.standardPuncturedBoundaryCycle}, \href{https://github.com/Paul-Lez/HodgeConjecture/blob/main/HodgeConjecture/Lemmas/AlgebraicTopology/EuclideanLocalHomology.lean#L151}{\texttt{HodgeConjecture/Lemmas/AlgebraicTopology/EuclideanLocalHomology.lean:151}}.
\end{definition}

\begin{definition}[AlgebraicTopology.Singular.standardPuncturedBoundaryClass]\label{def:AlgebraicTopology.Singular.standardPuncturedBoundaryClass}
\lean{AlgebraicTopology.Singular.standardPuncturedBoundaryClass}\leanok
\uses{def:AlgebraicTopology.Singular.Homology, def:AlgebraicTopology.Singular.standardPuncturedBoundaryCycle}
The homology class of the boundary of the standard affine simplex in punctured Euclidean
space.

\noindent\emph{Lean:} \texttt{def AlgebraicTopology.Singular.standardPuncturedBoundaryClass}, \href{https://github.com/Paul-Lez/HodgeConjecture/blob/main/HodgeConjecture/Lemmas/AlgebraicTopology/EuclideanLocalHomology.lean#L159}{\texttt{HodgeConjecture/Lemmas/AlgebraicTopology/EuclideanLocalHomology.lean:159}}.
\end{definition}

\begin{definition}[AlgebraicTopology.Singular.rationalSingularHomologyIsoOfHomotopyEquiv]\label{def:AlgebraicTopology.Singular.rationalSingularHomologyIsoOfHomotopyEquiv}
\lean{AlgebraicTopology.Singular.rationalSingularHomologyIsoOfHomotopyEquiv}\leanok
\uses{def:AlgebraicTopology.Singular.Homology}
A specified homotopy equivalence induces an isomorphism on rational singular homology.

\noindent\emph{Lean:} \texttt{def AlgebraicTopology.Singular.rationalSingularHomologyIsoOfHomotopyEquiv}, \href{https://github.com/Paul-Lez/HodgeConjecture/blob/main/HodgeConjecture/Lemmas/AlgebraicTopology/EuclideanLocalHomology.lean#L193}{\texttt{HodgeConjecture/Lemmas/AlgebraicTopology/EuclideanLocalHomology.lean:193}}.
\end{definition}

\begin{definition}[AlgebraicTopology.Singular.standardPuncturedRelativeBoundaryIso]\label{def:AlgebraicTopology.Singular.standardPuncturedRelativeBoundaryIso}
\lean{AlgebraicTopology.Singular.standardPuncturedRelativeBoundaryIso}\leanok
\uses{def:AlgebraicTopology.Singular.RelativeHomology, def:AlgebraicTopology.Singular.rationalSingularHomologyIsoOfHomotopyEquiv, def:AlgebraicTopology.Singular.relativeSingularChainShortComplex, def:AlgebraicTopology.Singular.standardPuncturedPair}
In positive degree, the connecting map identifies the local homology of Euclidean space with
the preceding homology of its punctured complement.

\noindent\emph{Lean:} \texttt{def AlgebraicTopology.Singular.standardPuncturedRelativeBoundaryIso}, \href{https://github.com/Paul-Lez/HodgeConjecture/blob/main/HodgeConjecture/Lemmas/AlgebraicTopology/EuclideanLocalHomology.lean#L217}{\texttt{HodgeConjecture/Lemmas/AlgebraicTopology/EuclideanLocalHomology.lean:217}}.
\end{definition}

\begin{definition}[AlgebraicTopology.Singular.standardNegativePoint]\label{def:AlgebraicTopology.Singular.standardNegativePoint}
\lean{AlgebraicTopology.Singular.standardNegativePoint}\leanok
\uses{def:AlgebraicTopology.Singular.StandardRealModel}
The point \texttt{-1} of the standard punctured real line.

\noindent\emph{Lean:} \texttt{def AlgebraicTopology.Singular.standardNegativePoint}, \href{https://github.com/Paul-Lez/HodgeConjecture/blob/main/HodgeConjecture/Lemmas/AlgebraicTopology/EuclideanLocalHomology.lean#L267}{\texttt{HodgeConjecture/Lemmas/AlgebraicTopology/EuclideanLocalHomology.lean:267}}.
\end{definition}

\begin{definition}[AlgebraicTopology.Singular.standardPositivePoint]\label{def:AlgebraicTopology.Singular.standardPositivePoint}
\lean{AlgebraicTopology.Singular.standardPositivePoint}\leanok
\uses{def:AlgebraicTopology.Singular.StandardRealModel}
The point \texttt{1} of the standard punctured real line.

\noindent\emph{Lean:} \texttt{def AlgebraicTopology.Singular.standardPositivePoint}, \href{https://github.com/Paul-Lez/HodgeConjecture/blob/main/HodgeConjecture/Lemmas/AlgebraicTopology/EuclideanLocalHomology.lean#L274}{\texttt{HodgeConjecture/Lemmas/AlgebraicTopology/EuclideanLocalHomology.lean:274}}.
\end{definition}

\begin{definition}[AlgebraicTopology.Singular.standardPuncturedFaceCycle]\label{def:AlgebraicTopology.Singular.standardPuncturedFaceCycle}
\lean{AlgebraicTopology.Singular.standardPuncturedFaceCycle}\leanok
\uses{def:AlgebraicTopology.Singular.standardSubspaceFaceChain}
A vertex of the standard one-simplex, regarded as a zero-cycle in the punctured line.

\noindent\emph{Lean:} \texttt{def AlgebraicTopology.Singular.standardPuncturedFaceCycle}, \href{https://github.com/Paul-Lez/HodgeConjecture/blob/main/HodgeConjecture/Lemmas/AlgebraicTopology/EuclideanLocalHomology.lean#L341}{\texttt{HodgeConjecture/Lemmas/AlgebraicTopology/EuclideanLocalHomology.lean:341}}.
\end{definition}

\begin{definition}[AlgebraicTopology.Singular.standardPuncturedBoundaryDetector]\label{def:AlgebraicTopology.Singular.standardPuncturedBoundaryDetector}
\lean{AlgebraicTopology.Singular.standardPuncturedBoundaryDetector}\leanok
\uses{def:AlgebraicTopology.Singular.standardFaceSimplex}
A linear functional detecting the negative component of the punctured line.

\noindent\emph{Lean:} \texttt{def AlgebraicTopology.Singular.standardPuncturedBoundaryDetector}, \href{https://github.com/Paul-Lez/HodgeConjecture/blob/main/HodgeConjecture/Lemmas/AlgebraicTopology/EuclideanLocalHomology.lean#L387}{\texttt{HodgeConjecture/Lemmas/AlgebraicTopology/EuclideanLocalHomology.lean:387}}.
\end{definition}

\section{Singular chains of a subspace}\label{sec:HodgeConjecture.Lemmas.AlgebraicTopology.SingularSubsetChains}
\noindent\emph{File:} \href{https://github.com/Paul-Lez/HodgeConjecture/blob/main/HodgeConjecture/Lemmas/AlgebraicTopology/SingularSubsetChains.lean}{\texttt{HodgeConjecture/Lemmas/AlgebraicTopology/SingularSubsetChains.lean}}.

A singular simplex of a space factors through a subspace exactly when its topological image is
contained in that subspace, and the factorization is unique because the inclusion is injective.
This file turns that remark into the maps it is used through: a simplicial map into the singular
simplicial set lifts to the subspace whenever every represented simplex has image there, and a
single simplex lifts likewise.

Degreewise, the resulting lift retracts the integral chain inclusion of the subspace, so that
inclusion is a split monomorphism in every degree and is therefore cancellable on the left.

\begin{definition}[AlgebraicTopology.Singular.singularSimplicialMapLiftToSubset]\label{def:AlgebraicTopology.Singular.singularSimplicialMapLiftToSubset}
\lean{AlgebraicTopology.Singular.singularSimplicialMapLiftToSubset}\leanok
A simplicial map to a singular simplicial set lifts to a topological subspace when every
represented continuous simplex has image in that subspace.

\noindent\emph{Lean:} \texttt{def AlgebraicTopology.Singular.singularSimplicialMapLiftToSubset}, \href{https://github.com/Paul-Lez/HodgeConjecture/blob/main/HodgeConjecture/Lemmas/AlgebraicTopology/SingularSubsetChains.lean#L61}{\texttt{HodgeConjecture/Lemmas/AlgebraicTopology/SingularSubsetChains.lean:61}}.
\end{definition}

\begin{definition}[AlgebraicTopology.Singular.singularSimplexLiftToSubset]\label{def:AlgebraicTopology.Singular.singularSimplexLiftToSubset}
\lean{AlgebraicTopology.Singular.singularSimplexLiftToSubset}\leanok
A singular simplex whose image lies in a subspace, regarded as a simplex of that
subspace.

\noindent\emph{Lean:} \texttt{def AlgebraicTopology.Singular.singularSimplexLiftToSubset}, \href{https://github.com/Paul-Lez/HodgeConjecture/blob/main/HodgeConjecture/Lemmas/AlgebraicTopology/SingularSubsetChains.lean#L126}{\texttt{HodgeConjecture/Lemmas/AlgebraicTopology/SingularSubsetChains.lean:126}}.
\end{definition}

\begin{definition}[AlgebraicTopology.Singular.singularSubsetIntegralChainRetractionComponent]\label{def:AlgebraicTopology.Singular.singularSubsetIntegralChainRetractionComponent}
\lean{AlgebraicTopology.Singular.singularSubsetIntegralChainRetractionComponent}\leanok
\uses{def:AlgebraicTopology.Singular.singularSimplexLiftToSubset}
A degreewise retraction of integral singular chains along a topological subspace
inclusion.  It sends a simplex outside the subspace to zero.

\noindent\emph{Lean:} \texttt{def AlgebraicTopology.Singular.singularSubsetIntegralChainRetractionComponent}, \href{https://github.com/Paul-Lez/HodgeConjecture/blob/main/HodgeConjecture/Lemmas/AlgebraicTopology/SingularSubsetChains.lean#L148}{\texttt{HodgeConjecture/Lemmas/AlgebraicTopology/SingularSubsetChains.lean:148}}.
\end{definition}

\section{SingularBarycentricHomotopy}\label{sec:HodgeConjecture.Lemmas.AlgebraicTopology.SingularBarycentricHomotopy}
\noindent\emph{File:} \href{https://github.com/Paul-Lez/HodgeConjecture/blob/main/HodgeConjecture/Lemmas/AlgebraicTopology/SingularBarycentricHomotopy.lean}{\texttt{HodgeConjecture/Lemmas/AlgebraicTopology/SingularBarycentricHomotopy.lean}}.

This module is ported from Paul Lezeau's corresponding file in
\texttt{sphere-six-complex} pull request \#49, under the Apache-2.0 license.

\paragraph{The prism identity for barycentric subdivision}

The composite of barycentric subdivision with the last-vertex map is expected to be chain
homotopic to the identity.  Since simplicial chains are freely generated by simplices, the
remaining content is one coherent prism chain on each standard simplex.  This file isolates that
exact acyclic-model identity and constructs its natural degreewise operators on every simplicial
set.

\begin{definition}[AlgebraicTopology.Singular.standardSimplexTopSimplex]\label{def:AlgebraicTopology.Singular.standardSimplexTopSimplex}
\lean{AlgebraicTopology.Singular.standardSimplexTopSimplex}\leanok
The universal top-dimensional simplex of \texttt{Δ[n]}.

\noindent\emph{Lean:} \texttt{def AlgebraicTopology.Singular.standardSimplexTopSimplex}, \href{https://github.com/Paul-Lez/HodgeConjecture/blob/main/HodgeConjecture/Lemmas/AlgebraicTopology/SingularBarycentricHomotopy.lean#L41}{\texttt{HodgeConjecture/Lemmas/AlgebraicTopology/SingularBarycentricHomotopy.lean:41}}.
\end{definition}

\begin{definition}[AlgebraicTopology.Singular.standardBarycentricLastVertexDiscrepancy]\label{def:AlgebraicTopology.Singular.standardBarycentricLastVertexDiscrepancy}
\lean{AlgebraicTopology.Singular.standardBarycentricLastVertexDiscrepancy}\leanok
\uses{def:AlgebraicTopology.Singular.barycentricSubdivisionSimplexChain, def:AlgebraicTopology.Singular.standardSimplexTopSimplex, def:AlgebraicTopology.Singular.subdivisionLastVertexChainMap}
On the universal simplex, subdivision followed by last vertex minus the identity.

\noindent\emph{Lean:} \texttt{def AlgebraicTopology.Singular.standardBarycentricLastVertexDiscrepancy}, \href{https://github.com/Paul-Lez/HodgeConjecture/blob/main/HodgeConjecture/Lemmas/AlgebraicTopology/SingularBarycentricHomotopy.lean#L46}{\texttt{HodgeConjecture/Lemmas/AlgebraicTopology/SingularBarycentricHomotopy.lean:46}}.
\end{definition}

\begin{definition}[AlgebraicTopology.Singular.standardPrismFaceChain]\label{def:AlgebraicTopology.Singular.standardPrismFaceChain}
\lean{AlgebraicTopology.Singular.standardPrismFaceChain}\leanok
The face contribution made by lower-dimensional prism chains.

\noindent\emph{Lean:} \texttt{def AlgebraicTopology.Singular.standardPrismFaceChain}, \href{https://github.com/Paul-Lez/HodgeConjecture/blob/main/HodgeConjecture/Lemmas/AlgebraicTopology/SingularBarycentricHomotopy.lean#L90}{\texttt{HodgeConjecture/Lemmas/AlgebraicTopology/SingularBarycentricHomotopy.lean:90}}.
\end{definition}

\begin{definition}[AlgebraicTopology.Singular.BarycentricLastVertexPrismData]\label{def:AlgebraicTopology.Singular.BarycentricLastVertexPrismData}
\lean{AlgebraicTopology.Singular.BarycentricLastVertexPrismData}\leanok
\uses{def:AlgebraicTopology.Singular.standardBarycentricLastVertexDiscrepancy, def:AlgebraicTopology.Singular.standardPrismFaceChain}
The smallest remaining acyclic-model datum.  Its equation is the literal universal-simplex
prism identity: discrepancy equals the prism boundary plus the transported prisms of all faces.

\noindent\emph{Lean:} \texttt{structure AlgebraicTopology.Singular.BarycentricLastVertexPrismData}, \href{https://github.com/Paul-Lez/HodgeConjecture/blob/main/HodgeConjecture/Lemmas/AlgebraicTopology/SingularBarycentricHomotopy.lean#L106}{\texttt{HodgeConjecture/Lemmas/AlgebraicTopology/SingularBarycentricHomotopy.lean:106}}.
\end{definition}

\begin{definition}[AlgebraicTopology.Singular.barycentricLastVertexPrismSimplexChain]\label{def:AlgebraicTopology.Singular.barycentricLastVertexPrismSimplexChain}
\lean{AlgebraicTopology.Singular.barycentricLastVertexPrismSimplexChain}\leanok
\uses{def:AlgebraicTopology.Singular.BarycentricLastVertexPrismData}
Transport a universal prism chain along the simplex represented by \texttt{x}.

\noindent\emph{Lean:} \texttt{def AlgebraicTopology.Singular.barycentricLastVertexPrismSimplexChain}, \href{https://github.com/Paul-Lez/HodgeConjecture/blob/main/HodgeConjecture/Lemmas/AlgebraicTopology/SingularBarycentricHomotopy.lean#L120}{\texttt{HodgeConjecture/Lemmas/AlgebraicTopology/SingularBarycentricHomotopy.lean:120}}.
\end{definition}

\begin{definition}[AlgebraicTopology.Singular.barycentricLastVertexPrismComponent]\label{def:AlgebraicTopology.Singular.barycentricLastVertexPrismComponent}
\lean{AlgebraicTopology.Singular.barycentricLastVertexPrismComponent}\leanok
\uses{def:AlgebraicTopology.Singular.barycentricLastVertexPrismSimplexChain}
The degree-raising prism operator on the free integral simplicial chains of \texttt{X}.

\noindent\emph{Lean:} \texttt{def AlgebraicTopology.Singular.barycentricLastVertexPrismComponent}, \href{https://github.com/Paul-Lez/HodgeConjecture/blob/main/HodgeConjecture/Lemmas/AlgebraicTopology/SingularBarycentricHomotopy.lean#L130}{\texttt{HodgeConjecture/Lemmas/AlgebraicTopology/SingularBarycentricHomotopy.lean:130}}.
\end{definition}

\begin{definition}[AlgebraicTopology.Singular.barycentricLastVertexPrismHom]\label{def:AlgebraicTopology.Singular.barycentricLastVertexPrismHom}
\lean{AlgebraicTopology.Singular.barycentricLastVertexPrismHom}\leanok
\uses{def:AlgebraicTopology.Singular.barycentricLastVertexPrismComponent}
The homotopy family, supported exactly from degree \texttt{i} to degree \texttt{i+1}.

\noindent\emph{Lean:} \texttt{def AlgebraicTopology.Singular.barycentricLastVertexPrismHom}, \href{https://github.com/Paul-Lez/HodgeConjecture/blob/main/HodgeConjecture/Lemmas/AlgebraicTopology/SingularBarycentricHomotopy.lean#L284}{\texttt{HodgeConjecture/Lemmas/AlgebraicTopology/SingularBarycentricHomotopy.lean:284}}.
\end{definition}

\begin{definition}[AlgebraicTopology.Singular.barycentricSubdivisionLastVertexHomotopy]\label{def:AlgebraicTopology.Singular.barycentricSubdivisionLastVertexHomotopy}
\lean{AlgebraicTopology.Singular.barycentricSubdivisionLastVertexHomotopy}\leanok
\uses{def:AlgebraicTopology.Singular.barycentricLastVertexPrismHom, def:AlgebraicTopology.Singular.barycentricSubdivisionChainMapCanonical}
A coherent family of universal simplex prisms produces the desired chain homotopy.

\noindent\emph{Lean:} \texttt{def AlgebraicTopology.Singular.barycentricSubdivisionLastVertexHomotopy}, \href{https://github.com/Paul-Lez/HodgeConjecture/blob/main/HodgeConjecture/Lemmas/AlgebraicTopology/SingularBarycentricHomotopy.lean#L302}{\texttt{HodgeConjecture/Lemmas/AlgebraicTopology/SingularBarycentricHomotopy.lean:302}}.
\end{definition}

\section{SingularStandardSimplexCone}\label{sec:HodgeConjecture.Lemmas.AlgebraicTopology.SingularStandardSimplexCone}
\noindent\emph{File:} \href{https://github.com/Paul-Lez/HodgeConjecture/blob/main/HodgeConjecture/Lemmas/AlgebraicTopology/SingularStandardSimplexCone.lean}{\texttt{HodgeConjecture/Lemmas/AlgebraicTopology/SingularStandardSimplexCone.lean}}.

This module is ported from Paul Lezeau's corresponding file in
\texttt{sphere-six-complex} pull request \#49, under the Apache-2.0 license.

\paragraph{The cone contraction on standard-simplex chains}

Prepending the zero vertex defines an extra degeneracy on every standard simplex.  On integral
simplicial chains this is the classical cone operator.  This file constructs it on the actual
coproduct basis used by \texttt{SSet.chainComplex} and proves its positive-degree contraction identity.

\begin{definition}[AlgebraicTopology.Singular.standardSimplexZeroConeSimplex]\label{def:AlgebraicTopology.Singular.standardSimplexZeroConeSimplex}
\lean{AlgebraicTopology.Singular.standardSimplexZeroConeSimplex}\leanok
Prepend the zero vertex to a simplex of \texttt{Δ[m]}.

\noindent\emph{Lean:} \texttt{def AlgebraicTopology.Singular.standardSimplexZeroConeSimplex}, \href{https://github.com/Paul-Lez/HodgeConjecture/blob/main/HodgeConjecture/Lemmas/AlgebraicTopology/SingularStandardSimplexCone.lean#L39}{\texttt{HodgeConjecture/Lemmas/AlgebraicTopology/SingularStandardSimplexCone.lean:39}}.
\end{definition}

\begin{definition}[AlgebraicTopology.Singular.standardSimplexZeroConeComponent]\label{def:AlgebraicTopology.Singular.standardSimplexZeroConeComponent}
\lean{AlgebraicTopology.Singular.standardSimplexZeroConeComponent}\leanok
\uses{def:AlgebraicTopology.Singular.standardSimplexZeroConeSimplex}
The integral cone operator on all basis simplices in degree \texttt{n}.

\noindent\emph{Lean:} \texttt{def AlgebraicTopology.Singular.standardSimplexZeroConeComponent}, \href{https://github.com/Paul-Lez/HodgeConjecture/blob/main/HodgeConjecture/Lemmas/AlgebraicTopology/SingularStandardSimplexCone.lean#L66}{\texttt{HodgeConjecture/Lemmas/AlgebraicTopology/SingularStandardSimplexCone.lean:66}}.
\end{definition}

\begin{definition}[AlgebraicTopology.Singular.barycentricLastVertexDiscrepancyChainMap]\label{def:AlgebraicTopology.Singular.barycentricLastVertexDiscrepancyChainMap}
\lean{AlgebraicTopology.Singular.barycentricLastVertexDiscrepancyChainMap}\leanok
\uses{def:AlgebraicTopology.Singular.barycentricSubdivisionChainMapCanonical, def:AlgebraicTopology.Singular.subdivisionLastVertexChainMap}
Subdivision followed by last vertex, minus the identity, as an actual chain map.

\noindent\emph{Lean:} \texttt{def AlgebraicTopology.Singular.barycentricLastVertexDiscrepancyChainMap}, \href{https://github.com/Paul-Lez/HodgeConjecture/blob/main/HodgeConjecture/Lemmas/AlgebraicTopology/SingularStandardSimplexCone.lean#L127}{\texttt{HodgeConjecture/Lemmas/AlgebraicTopology/SingularStandardSimplexCone.lean:127}}.
\end{definition}

\begin{definition}[AlgebraicTopology.Singular.universalStandardSimplexPrismComponent]\label{def:AlgebraicTopology.Singular.universalStandardSimplexPrismComponent}
\lean{AlgebraicTopology.Singular.universalStandardSimplexPrismComponent}\leanok
Transport one universal degree-raising chain along every simplex of an arbitrary simplicial
set.

\noindent\emph{Lean:} \texttt{def AlgebraicTopology.Singular.universalStandardSimplexPrismComponent}, \href{https://github.com/Paul-Lez/HodgeConjecture/blob/main/HodgeConjecture/Lemmas/AlgebraicTopology/SingularStandardSimplexCone.lean#L260}{\texttt{HodgeConjecture/Lemmas/AlgebraicTopology/SingularStandardSimplexCone.lean:260}}.
\end{definition}

\begin{definition}[AlgebraicTopology.Singular.canonicalBarycentricLastVertexPrism]\label{def:AlgebraicTopology.Singular.canonicalBarycentricLastVertexPrism}
\lean{AlgebraicTopology.Singular.canonicalBarycentricLastVertexPrism}\leanok
\uses{def:AlgebraicTopology.Singular.standardBarycentricLastVertexDiscrepancy, def:AlgebraicTopology.Singular.standardSimplexZeroConeComponent}
The universal prism chains obtained recursively by filling each residual with the
zero-vertex cone.

\noindent\emph{Lean:} \texttt{def AlgebraicTopology.Singular.canonicalBarycentricLastVertexPrism}, \href{https://github.com/Paul-Lez/HodgeConjecture/blob/main/HodgeConjecture/Lemmas/AlgebraicTopology/SingularStandardSimplexCone.lean#L441}{\texttt{HodgeConjecture/Lemmas/AlgebraicTopology/SingularStandardSimplexCone.lean:441}}.
\end{definition}

\begin{definition}[AlgebraicTopology.Singular.barycentricLastVertexPrismDataCanonical]\label{def:AlgebraicTopology.Singular.barycentricLastVertexPrismDataCanonical}
\lean{AlgebraicTopology.Singular.barycentricLastVertexPrismDataCanonical}\leanok
\uses{def:AlgebraicTopology.Singular.BarycentricLastVertexPrismData, def:AlgebraicTopology.Singular.barycentricLastVertexDiscrepancyChainMap, def:AlgebraicTopology.Singular.barycentricSubdivisionChainMapZero, def:AlgebraicTopology.Singular.canonicalBarycentricLastVertexPrism, def:AlgebraicTopology.Singular.universalStandardSimplexPrismComponent}
The explicit recursively coned universal prisms supply the formerly missing acyclic-model
datum.

\noindent\emph{Lean:} \texttt{def AlgebraicTopology.Singular.barycentricLastVertexPrismDataCanonical}, \href{https://github.com/Paul-Lez/HodgeConjecture/blob/main/HodgeConjecture/Lemmas/AlgebraicTopology/SingularStandardSimplexCone.lean#L517}{\texttt{HodgeConjecture/Lemmas/AlgebraicTopology/SingularStandardSimplexCone.lean:517}}.
\end{definition}

\begin{definition}[AlgebraicTopology.Singular.barycentricSubdivisionLastVertexHomotopyCanonical]\label{def:AlgebraicTopology.Singular.barycentricSubdivisionLastVertexHomotopyCanonical}
\lean{AlgebraicTopology.Singular.barycentricSubdivisionLastVertexHomotopyCanonical}\leanok
\uses{def:AlgebraicTopology.Singular.barycentricLastVertexPrismDataCanonical, def:AlgebraicTopology.Singular.barycentricSubdivisionLastVertexHomotopy}
The explicit canonical chain homotopy from subdivision followed by last vertex to the
identity.

\noindent\emph{Lean:} \texttt{def AlgebraicTopology.Singular.barycentricSubdivisionLastVertexHomotopyCanonical}, \href{https://github.com/Paul-Lez/HodgeConjecture/blob/main/HodgeConjecture/Lemmas/AlgebraicTopology/SingularStandardSimplexCone.lean#L524}{\texttt{HodgeConjecture/Lemmas/AlgebraicTopology/SingularStandardSimplexCone.lean:524}}.
\end{definition}

\section{Rational top homology of simplicial spheres}\label{sec:HodgeConjecture.Lemmas.AlgebraicTopology.StandardSphereSimplicialHomology}
\noindent\emph{File:} \href{https://github.com/Paul-Lez/HodgeConjecture/blob/main/HodgeConjecture/Lemmas/AlgebraicTopology/StandardSphereSimplicialHomology.lean}{\texttt{HodgeConjecture/Lemmas/AlgebraicTopology/StandardSphereSimplicialHomology.lean}}.

This file computes the top rational simplicial homology of the boundary of every positive-
dimensional standard simplex. The proof compares normalized chains below the top dimension with
the full simplex and uses its extra degeneracy.

\begin{definition}[AlgebraicTopology.Singular.standardSimplexSuccNormalizedRationalChains]\label{def:AlgebraicTopology.Singular.standardSimplexSuccNormalizedRationalChains}
\lean{AlgebraicTopology.Singular.standardSimplexSuccNormalizedRationalChains}\leanok
Normalized rational chains of the standard \texttt{(n + 2)}-simplex.

\noindent\emph{Lean:} \texttt{abbrev AlgebraicTopology.Singular.standardSimplexSuccNormalizedRationalChains}, \href{https://github.com/Paul-Lez/HodgeConjecture/blob/main/HodgeConjecture/Lemmas/AlgebraicTopology/StandardSphereSimplicialHomology.lean#L41}{\texttt{HodgeConjecture/Lemmas/AlgebraicTopology/StandardSphereSimplicialHomology.lean:41}}.
\end{definition}

\begin{definition}[AlgebraicTopology.Singular.standardSphereSuccNormalizedRationalChains]\label{def:AlgebraicTopology.Singular.standardSphereSuccNormalizedRationalChains}
\lean{AlgebraicTopology.Singular.standardSphereSuccNormalizedRationalChains}\leanok
Normalized rational chains of the boundary of the standard \texttt{(n + 2)}-simplex. Its
realization is an \texttt{(n + 1)}-sphere.

\noindent\emph{Lean:} \texttt{abbrev AlgebraicTopology.Singular.standardSphereSuccNormalizedRationalChains}, \href{https://github.com/Paul-Lez/HodgeConjecture/blob/main/HodgeConjecture/Lemmas/AlgebraicTopology/StandardSphereSimplicialHomology.lean#L46}{\texttt{HodgeConjecture/Lemmas/AlgebraicTopology/StandardSphereSimplicialHomology.lean:46}}.
\end{definition}

\begin{definition}[AlgebraicTopology.Singular.standardSimplexSuccNormalizedChainsToBoundary]\label{def:AlgebraicTopology.Singular.standardSimplexSuccNormalizedChainsToBoundary}
\lean{AlgebraicTopology.Singular.standardSimplexSuccNormalizedChainsToBoundary}\leanok
\uses{def:AlgebraicTopology.Singular.standardSimplexSuccNormalizedRationalChains, def:AlgebraicTopology.Singular.standardSphereSuccNormalizedRationalChains}
Below dimension \texttt{n + 2}, regard a nondegenerate simplex of the full standard simplex as a
simplex of its boundary.

\noindent\emph{Lean:} \texttt{def AlgebraicTopology.Singular.standardSimplexSuccNormalizedChainsToBoundary}, \href{https://github.com/Paul-Lez/HodgeConjecture/blob/main/HodgeConjecture/Lemmas/AlgebraicTopology/StandardSphereSimplicialHomology.lean#L64}{\texttt{HodgeConjecture/Lemmas/AlgebraicTopology/StandardSphereSimplicialHomology.lean:64}}.
\end{definition}

\begin{definition}[AlgebraicTopology.Singular.standardSphereSuccNormalizedChainsXIsoStandard]\label{def:AlgebraicTopology.Singular.standardSphereSuccNormalizedChainsXIsoStandard}
\lean{AlgebraicTopology.Singular.standardSphereSuccNormalizedChainsXIsoStandard}\leanok
\uses{def:AlgebraicTopology.Singular.standardSimplexSuccNormalizedChainsToBoundary}
Below its top simplex, the normalized groups of a standard simplex and its boundary agree.

\noindent\emph{Lean:} \texttt{def AlgebraicTopology.Singular.standardSphereSuccNormalizedChainsXIsoStandard}, \href{https://github.com/Paul-Lez/HodgeConjecture/blob/main/HodgeConjecture/Lemmas/AlgebraicTopology/StandardSphereSimplicialHomology.lean#L89}{\texttt{HodgeConjecture/Lemmas/AlgebraicTopology/StandardSphereSimplicialHomology.lean:89}}.
\end{definition}

\begin{definition}[AlgebraicTopology.Singular.standardSphereSuccTopCyclesIsoStandardSimplexSuccTopCycles]\label{def:AlgebraicTopology.Singular.standardSphereSuccTopCyclesIsoStandardSimplexSuccTopCycles}
\lean{AlgebraicTopology.Singular.standardSphereSuccTopCyclesIsoStandardSimplexSuccTopCycles}\leanok
\uses{def:AlgebraicTopology.Singular.standardSphereSuccNormalizedChainsXIsoStandard}
The top cycle kernels of the boundary and full standard simplex agree.

\noindent\emph{Lean:} \texttt{def AlgebraicTopology.Singular.standardSphereSuccTopCyclesIsoStandardSimplexSuccTopCycles}, \href{https://github.com/Paul-Lez/HodgeConjecture/blob/main/HodgeConjecture/Lemmas/AlgebraicTopology/StandardSphereSimplicialHomology.lean#L120}{\texttt{HodgeConjecture/Lemmas/AlgebraicTopology/StandardSphereSimplicialHomology.lean:120}}.
\end{definition}

\begin{definition}[AlgebraicTopology.Singular.standardSimplexSuccNormalizedChainsXTopIsoRat]\label{def:AlgebraicTopology.Singular.standardSimplexSuccNormalizedChainsXTopIsoRat}
\lean{AlgebraicTopology.Singular.standardSimplexSuccNormalizedChainsXTopIsoRat}\leanok
\uses{def:AlgebraicTopology.Singular.standardSimplexSuccNormalizedRationalChains}
The top normalized group of a standard simplex is canonically one-dimensional.

\noindent\emph{Lean:} \texttt{def AlgebraicTopology.Singular.standardSimplexSuccNormalizedChainsXTopIsoRat}, \href{https://github.com/Paul-Lez/HodgeConjecture/blob/main/HodgeConjecture/Lemmas/AlgebraicTopology/StandardSphereSimplicialHomology.lean#L145}{\texttt{HodgeConjecture/Lemmas/AlgebraicTopology/StandardSphereSimplicialHomology.lean:145}}.
\end{definition}

\begin{definition}[AlgebraicTopology.Singular.standardSimplexSuccNormalizedChainsXTopIsoTopCycles]\label{def:AlgebraicTopology.Singular.standardSimplexSuccNormalizedChainsXTopIsoTopCycles}
\lean{AlgebraicTopology.Singular.standardSimplexSuccNormalizedChainsXTopIsoTopCycles}\leanok
\uses{def:AlgebraicTopology.Singular.standardSimplexSuccNormalizedRationalChains}
Exactness identifies the top group of a full standard simplex with the cycle kernel one
degree below.

\noindent\emph{Lean:} \texttt{def AlgebraicTopology.Singular.standardSimplexSuccNormalizedChainsXTopIsoTopCycles}, \href{https://github.com/Paul-Lez/HodgeConjecture/blob/main/HodgeConjecture/Lemmas/AlgebraicTopology/StandardSphereSimplicialHomology.lean#L182}{\texttt{HodgeConjecture/Lemmas/AlgebraicTopology/StandardSphereSimplicialHomology.lean:182}}.
\end{definition}

\begin{definition}[AlgebraicTopology.Singular.standardSphereSuccTopCyclesIsoRat]\label{def:AlgebraicTopology.Singular.standardSphereSuccTopCyclesIsoRat}
\lean{AlgebraicTopology.Singular.standardSphereSuccTopCyclesIsoRat}\leanok
\uses{def:AlgebraicTopology.Singular.standardSimplexSuccNormalizedChainsXTopIsoRat, def:AlgebraicTopology.Singular.standardSimplexSuccNormalizedChainsXTopIsoTopCycles, def:AlgebraicTopology.Singular.standardSphereSuccTopCyclesIsoStandardSimplexSuccTopCycles}
The top normalized cycle kernel of a standard simplicial sphere is one-dimensional.

\noindent\emph{Lean:} \texttt{def AlgebraicTopology.Singular.standardSphereSuccTopCyclesIsoRat}, \href{https://github.com/Paul-Lez/HodgeConjecture/blob/main/HodgeConjecture/Lemmas/AlgebraicTopology/StandardSphereSimplicialHomology.lean#L201}{\texttt{HodgeConjecture/Lemmas/AlgebraicTopology/StandardSphereSimplicialHomology.lean:201}}.
\end{definition}

\begin{definition}[AlgebraicTopology.Singular.standardSphereSuccNormalizedHomologyTopIsoTopCycles]\label{def:AlgebraicTopology.Singular.standardSphereSuccNormalizedHomologyTopIsoTopCycles}
\lean{AlgebraicTopology.Singular.standardSphereSuccNormalizedHomologyTopIsoTopCycles}\leanok
\uses{def:AlgebraicTopology.Singular.standardSphereSuccNormalizedRationalChains}
Top normalized homology of the boundary of a standard simplex is its top cycle kernel.

\noindent\emph{Lean:} \texttt{def AlgebraicTopology.Singular.standardSphereSuccNormalizedHomologyTopIsoTopCycles}, \href{https://github.com/Paul-Lez/HodgeConjecture/blob/main/HodgeConjecture/Lemmas/AlgebraicTopology/StandardSphereSimplicialHomology.lean#L215}{\texttt{HodgeConjecture/Lemmas/AlgebraicTopology/StandardSphereSimplicialHomology.lean:215}}.
\end{definition}

\begin{definition}[AlgebraicTopology.Singular.standardSphereSuccNormalizedHomologyTopIsoRat]\label{def:AlgebraicTopology.Singular.standardSphereSuccNormalizedHomologyTopIsoRat}
\lean{AlgebraicTopology.Singular.standardSphereSuccNormalizedHomologyTopIsoRat}\leanok
\uses{def:AlgebraicTopology.Singular.standardSphereSuccNormalizedHomologyTopIsoTopCycles, def:AlgebraicTopology.Singular.standardSphereSuccTopCyclesIsoRat}
Normalized rational top homology of a positive-dimensional standard simplicial sphere is
one-dimensional.

\noindent\emph{Lean:} \texttt{def AlgebraicTopology.Singular.standardSphereSuccNormalizedHomologyTopIsoRat}, \href{https://github.com/Paul-Lez/HodgeConjecture/blob/main/HodgeConjecture/Lemmas/AlgebraicTopology/StandardSphereSimplicialHomology.lean#L226}{\texttt{HodgeConjecture/Lemmas/AlgebraicTopology/StandardSphereSimplicialHomology.lean:226}}.
\end{definition}

\begin{definition}[AlgebraicTopology.Singular.standardSphereSuccSimplicialHomologyTopIsoRat]\label{def:AlgebraicTopology.Singular.standardSphereSuccSimplicialHomologyTopIsoRat}
\lean{AlgebraicTopology.Singular.standardSphereSuccSimplicialHomologyTopIsoRat}\leanok
\uses{def:AlgebraicTopology.Singular.standardSphereSuccNormalizedHomologyTopIsoRat}
Rational simplicial homology of the boundary of the standard \texttt{(n + 2)}-simplex in degree
\texttt{n + 1} is one-dimensional.

\noindent\emph{Lean:} \texttt{def AlgebraicTopology.Singular.standardSphereSuccSimplicialHomologyTopIsoRat}, \href{https://github.com/Paul-Lez/HodgeConjecture/blob/main/HodgeConjecture/Lemmas/AlgebraicTopology/StandardSphereSimplicialHomology.lean#L234}{\texttt{HodgeConjecture/Lemmas/AlgebraicTopology/StandardSphereSimplicialHomology.lean:234}}.
\end{definition}

\section{The affine boundary map into punctured Euclidean space}\label{sec:HodgeConjecture.Lemmas.AlgebraicTopology.StandardSphereAffineBoundary}
\noindent\emph{File:} \href{https://github.com/Paul-Lez/HodgeConjecture/blob/main/HodgeConjecture/Lemmas/AlgebraicTopology/StandardSphereAffineBoundary.lean}{\texttt{HodgeConjecture/Lemmas/AlgebraicTopology/StandardSphereAffineBoundary.lean}}.

The universal affine simplex meets the origin only at its barycenter. Restricting it to the
simplicial boundary therefore gives a simplicial map into punctured Euclidean space. This file
constructs that map in every dimension.

\begin{definition}[AlgebraicTopology.Singular.standardAffineBoundarySimplex]\label{def:AlgebraicTopology.Singular.standardAffineBoundarySimplex}
\lean{AlgebraicTopology.Singular.standardAffineBoundarySimplex}\leanok
\uses{def:AlgebraicTopology.Singular.standardAffineSimplex}
A simplex in the simplicial boundary, realized affinely in punctured coordinate space.

\noindent\emph{Lean:} \texttt{def AlgebraicTopology.Singular.standardAffineBoundarySimplex}, \href{https://github.com/Paul-Lez/HodgeConjecture/blob/main/HodgeConjecture/Lemmas/AlgebraicTopology/StandardSphereAffineBoundary.lean#L49}{\texttt{HodgeConjecture/Lemmas/AlgebraicTopology/StandardSphereAffineBoundary.lean:49}}.
\end{definition}

\begin{definition}[AlgebraicTopology.Singular.standardAffineBoundarySimplicialMap]\label{def:AlgebraicTopology.Singular.standardAffineBoundarySimplicialMap}
\lean{AlgebraicTopology.Singular.standardAffineBoundarySimplicialMap}\leanok
\uses{def:AlgebraicTopology.Singular.standardAffineBoundarySimplex, def:AlgebraicTopology.Singular.standardPuncturedPair}
The boundary of the universal affine \texttt{d}-simplex as a map of simplicial sets into the
singular simplicial set of punctured \texttt{ℝ\^{}d}.

\noindent\emph{Lean:} \texttt{def AlgebraicTopology.Singular.standardAffineBoundarySimplicialMap}, \href{https://github.com/Paul-Lez/HodgeConjecture/blob/main/HodgeConjecture/Lemmas/AlgebraicTopology/StandardSphereAffineBoundary.lean#L59}{\texttt{HodgeConjecture/Lemmas/AlgebraicTopology/StandardSphereAffineBoundary.lean:59}}.
\end{definition}

\begin{definition}[AlgebraicTopology.Singular.standardSimplexTopSimplexForLocalClass]\label{def:AlgebraicTopology.Singular.standardSimplexTopSimplexForLocalClass}
\lean{AlgebraicTopology.Singular.standardSimplexTopSimplexForLocalClass}\leanok
The unique nondegenerate simplex in the top dimension of a standard simplex.

\noindent\emph{Lean:} \texttt{def AlgebraicTopology.Singular.standardSimplexTopSimplexForLocalClass}, \href{https://github.com/Paul-Lez/HodgeConjecture/blob/main/HodgeConjecture/Lemmas/AlgebraicTopology/StandardSphereAffineBoundary.lean#L84}{\texttt{HodgeConjecture/Lemmas/AlgebraicTopology/StandardSphereAffineBoundary.lean:84}}.
\end{definition}

\begin{definition}[AlgebraicTopology.Singular.standardSphereBoundaryFaceSimplex]\label{def:AlgebraicTopology.Singular.standardSphereBoundaryFaceSimplex}
\lean{AlgebraicTopology.Singular.standardSphereBoundaryFaceSimplex}\leanok
\uses{def:AlgebraicTopology.Singular.standardSimplexTopSimplexForLocalClass}
A top-dimensional face of the simplicial boundary of the standard \texttt{(n + 1)}-simplex.

\noindent\emph{Lean:} \texttt{def AlgebraicTopology.Singular.standardSphereBoundaryFaceSimplex}, \href{https://github.com/Paul-Lez/HodgeConjecture/blob/main/HodgeConjecture/Lemmas/AlgebraicTopology/StandardSphereAffineBoundary.lean#L89}{\texttt{HodgeConjecture/Lemmas/AlgebraicTopology/StandardSphereAffineBoundary.lean:89}}.
\end{definition}

\begin{definition}[AlgebraicTopology.Singular.standardAffineBoundaryChainMap]\label{def:AlgebraicTopology.Singular.standardAffineBoundaryChainMap}
\lean{AlgebraicTopology.Singular.standardAffineBoundaryChainMap}\leanok
\uses{def:AlgebraicTopology.Singular.standardAffineBoundarySimplicialMap}
The chain map from a simplicial boundary to singular chains of punctured coordinate space.

\noindent\emph{Lean:} \texttt{def AlgebraicTopology.Singular.standardAffineBoundaryChainMap}, \href{https://github.com/Paul-Lez/HodgeConjecture/blob/main/HodgeConjecture/Lemmas/AlgebraicTopology/StandardSphereAffineBoundary.lean#L131}{\texttt{HodgeConjecture/Lemmas/AlgebraicTopology/StandardSphereAffineBoundary.lean:131}}.
\end{definition}

\begin{definition}[AlgebraicTopology.Singular.standardSphereBoundaryChainRetractionComponent]\label{def:AlgebraicTopology.Singular.standardSphereBoundaryChainRetractionComponent}
\lean{AlgebraicTopology.Singular.standardSphereBoundaryChainRetractionComponent}\leanok
A degreewise retraction from chains of the full simplex to chains of its boundary. Simplices
outside the boundary are sent to zero.

\noindent\emph{Lean:} \texttt{def AlgebraicTopology.Singular.standardSphereBoundaryChainRetractionComponent}, \href{https://github.com/Paul-Lez/HodgeConjecture/blob/main/HodgeConjecture/Lemmas/AlgebraicTopology/StandardSphereAffineBoundary.lean#L138}{\texttt{HodgeConjecture/Lemmas/AlgebraicTopology/StandardSphereAffineBoundary.lean:138}}.
\end{definition}

\begin{definition}[AlgebraicTopology.Singular.standardSphereSimplicialBoundaryChain]\label{def:AlgebraicTopology.Singular.standardSphereSimplicialBoundaryChain}
\lean{AlgebraicTopology.Singular.standardSphereSimplicialBoundaryChain}\leanok
\uses{def:AlgebraicTopology.Singular.standardSphereBoundaryFaceSimplex}
The alternating top-face chain in the simplicial boundary.

\noindent\emph{Lean:} \texttt{def AlgebraicTopology.Singular.standardSphereSimplicialBoundaryChain}, \href{https://github.com/Paul-Lez/HodgeConjecture/blob/main/HodgeConjecture/Lemmas/AlgebraicTopology/StandardSphereAffineBoundary.lean#L181}{\texttt{HodgeConjecture/Lemmas/AlgebraicTopology/StandardSphereAffineBoundary.lean:181}}.
\end{definition}

\begin{definition}[AlgebraicTopology.Singular.standardSphereSimplicialBoundaryCycle]\label{def:AlgebraicTopology.Singular.standardSphereSimplicialBoundaryCycle}
\lean{AlgebraicTopology.Singular.standardSphereSimplicialBoundaryCycle}\leanok
\uses{def:AlgebraicTopology.Singular.standardSphereBoundaryChainRetractionComponent, def:AlgebraicTopology.Singular.standardSphereSimplicialBoundaryChain}
The alternating facets as a cycle in every positive-dimensional standard simplicial
sphere.

\noindent\emph{Lean:} \texttt{def AlgebraicTopology.Singular.standardSphereSimplicialBoundaryCycle}, \href{https://github.com/Paul-Lez/HodgeConjecture/blob/main/HodgeConjecture/Lemmas/AlgebraicTopology/StandardSphereAffineBoundary.lean#L220}{\texttt{HodgeConjecture/Lemmas/AlgebraicTopology/StandardSphereAffineBoundary.lean:220}}.
\end{definition}

\begin{definition}[AlgebraicTopology.Singular.standardSphereSimplicialBoundaryClass]\label{def:AlgebraicTopology.Singular.standardSphereSimplicialBoundaryClass}
\lean{AlgebraicTopology.Singular.standardSphereSimplicialBoundaryClass}\leanok
\uses{def:AlgebraicTopology.Singular.standardSphereSimplicialBoundaryCycle}
The simplicial homology class of the alternating facets.

\noindent\emph{Lean:} \texttt{def AlgebraicTopology.Singular.standardSphereSimplicialBoundaryClass}, \href{https://github.com/Paul-Lez/HodgeConjecture/blob/main/HodgeConjecture/Lemmas/AlgebraicTopology/StandardSphereAffineBoundary.lean#L234}{\texttt{HodgeConjecture/Lemmas/AlgebraicTopology/StandardSphereAffineBoundary.lean:234}}.
\end{definition}

\begin{definition}[AlgebraicTopology.Singular.standardSphereSimplicialNormalizedBoundaryChain]\label{def:AlgebraicTopology.Singular.standardSphereSimplicialNormalizedBoundaryChain}
\lean{AlgebraicTopology.Singular.standardSphereSimplicialNormalizedBoundaryChain}\leanok
\uses{def:AlgebraicTopology.Singular.standardSphereSimplicialBoundaryChain, def:AlgebraicTopology.Singular.standardSphereSuccNormalizedRationalChains}
The alternating facet chain after passage to normalized chains.

\noindent\emph{Lean:} \texttt{def AlgebraicTopology.Singular.standardSphereSimplicialNormalizedBoundaryChain}, \href{https://github.com/Paul-Lez/HodgeConjecture/blob/main/HodgeConjecture/Lemmas/AlgebraicTopology/StandardSphereAffineBoundary.lean#L242}{\texttt{HodgeConjecture/Lemmas/AlgebraicTopology/StandardSphereAffineBoundary.lean:242}}.
\end{definition}

\begin{definition}[AlgebraicTopology.Singular.standardSphereSimplicialNormalizedBoundaryDetector]\label{def:AlgebraicTopology.Singular.standardSphereSimplicialNormalizedBoundaryDetector}
\lean{AlgebraicTopology.Singular.standardSphereSimplicialNormalizedBoundaryDetector}\leanok
\uses{def:AlgebraicTopology.Singular.standardSphereBoundaryFaceSimplex, def:AlgebraicTopology.Singular.standardSphereSuccNormalizedRationalChains}
The coefficient functional of the zeroth oriented facet in normalized top chains.

\noindent\emph{Lean:} \texttt{def AlgebraicTopology.Singular.standardSphereSimplicialNormalizedBoundaryDetector}, \href{https://github.com/Paul-Lez/HodgeConjecture/blob/main/HodgeConjecture/Lemmas/AlgebraicTopology/StandardSphereAffineBoundary.lean#L250}{\texttt{HodgeConjecture/Lemmas/AlgebraicTopology/StandardSphereAffineBoundary.lean:250}}.
\end{definition}

\begin{definition}[AlgebraicTopology.Singular.standardSphereSimplicialNormalizedBoundaryCycle]\label{def:AlgebraicTopology.Singular.standardSphereSimplicialNormalizedBoundaryCycle}
\lean{AlgebraicTopology.Singular.standardSphereSimplicialNormalizedBoundaryCycle}\leanok
\uses{def:AlgebraicTopology.Singular.standardSphereBoundaryChainRetractionComponent, def:AlgebraicTopology.Singular.standardSphereSimplicialNormalizedBoundaryChain}
The normalized alternating facets as a top cycle.

\noindent\emph{Lean:} \texttt{def AlgebraicTopology.Singular.standardSphereSimplicialNormalizedBoundaryCycle}, \href{https://github.com/Paul-Lez/HodgeConjecture/blob/main/HodgeConjecture/Lemmas/AlgebraicTopology/StandardSphereAffineBoundary.lean#L328}{\texttt{HodgeConjecture/Lemmas/AlgebraicTopology/StandardSphereAffineBoundary.lean:328}}.
\end{definition}

\begin{definition}[AlgebraicTopology.Singular.standardSphereSimplicialNormalizedBoundaryClass]\label{def:AlgebraicTopology.Singular.standardSphereSimplicialNormalizedBoundaryClass}
\lean{AlgebraicTopology.Singular.standardSphereSimplicialNormalizedBoundaryClass}\leanok
\uses{def:AlgebraicTopology.Singular.standardSphereSimplicialNormalizedBoundaryCycle}
The normalized homology class of the alternating facets.

\noindent\emph{Lean:} \texttt{def AlgebraicTopology.Singular.standardSphereSimplicialNormalizedBoundaryClass}, \href{https://github.com/Paul-Lez/HodgeConjecture/blob/main/HodgeConjecture/Lemmas/AlgebraicTopology/StandardSphereAffineBoundary.lean#L346}{\texttt{HodgeConjecture/Lemmas/AlgebraicTopology/StandardSphereAffineBoundary.lean:346}}.
\end{definition}

\begin{definition}[AlgebraicTopology.Singular.standardPuncturedAffineBoundaryChain]\label{def:AlgebraicTopology.Singular.standardPuncturedAffineBoundaryChain}
\lean{AlgebraicTopology.Singular.standardPuncturedAffineBoundaryChain}\leanok
\uses{def:AlgebraicTopology.Singular.standardSubspaceFaceChain}
The alternating affine face chain in punctured coordinate space.

\noindent\emph{Lean:} \texttt{def AlgebraicTopology.Singular.standardPuncturedAffineBoundaryChain}, \href{https://github.com/Paul-Lez/HodgeConjecture/blob/main/HodgeConjecture/Lemmas/AlgebraicTopology/StandardSphereAffineBoundary.lean#L456}{\texttt{HodgeConjecture/Lemmas/AlgebraicTopology/StandardSphereAffineBoundary.lean:456}}.
\end{definition}

\section{Singular chains of contractible spaces}\label{sec:HodgeConjecture.Lemmas.AlgebraicTopology.SingularContractible}
\noindent\emph{File:} \href{https://github.com/Paul-Lez/HodgeConjecture/blob/main/HodgeConjecture/Lemmas/AlgebraicTopology/SingularContractible.lean}{\texttt{HodgeConjecture/Lemmas/AlgebraicTopology/SingularContractible.lean}}.

A topological homotopy equivalence induces a homotopy equivalence of singular chain complexes.
Consequently, the singular chain complex of a contractible space is exact in every positive
degree. The comparison with a point uses Mathlib's explicit calculation for totally disconnected
spaces.

\begin{definition}[AlgebraicTopology.singularChainHomotopyEquivOfHomotopyEquiv]\label{def:AlgebraicTopology.singularChainHomotopyEquivOfHomotopyEquiv}
\lean{AlgebraicTopology.singularChainHomotopyEquivOfHomotopyEquiv}\leanok
A topological homotopy equivalence induces a homotopy equivalence of singular chain
complexes.

\noindent\emph{Lean:} \texttt{def AlgebraicTopology.singularChainHomotopyEquivOfHomotopyEquiv}, \href{https://github.com/Paul-Lez/HodgeConjecture/blob/main/HodgeConjecture/Lemmas/AlgebraicTopology/SingularContractible.lean#L45}{\texttt{HodgeConjecture/Lemmas/AlgebraicTopology/SingularContractible.lean:45}}.
\end{definition}

\begin{definition}[AlgebraicTopology.contractibleSingularChainHomotopyEquiv]\label{def:AlgebraicTopology.contractibleSingularChainHomotopyEquiv}
\lean{AlgebraicTopology.contractibleSingularChainHomotopyEquiv}\leanok
\uses{def:AlgebraicTopology.singularChainHomotopyEquivOfHomotopyEquiv}
The singular chain complex of a contractible space is homotopy equivalent to that of a
point.

\noindent\emph{Lean:} \texttt{def AlgebraicTopology.contractibleSingularChainHomotopyEquiv}, \href{https://github.com/Paul-Lez/HodgeConjecture/blob/main/HodgeConjecture/Lemmas/AlgebraicTopology/SingularContractible.lean#L74}{\texttt{HodgeConjecture/Lemmas/AlgebraicTopology/SingularContractible.lean:74}}.
\end{definition}

\section{The affine sphere class in punctured Euclidean space}\label{sec:HodgeConjecture.Lemmas.AlgebraicTopology.PuncturedEuclideanFundamentalClass}
\noindent\emph{File:} \href{https://github.com/Paul-Lez/HodgeConjecture/blob/main/HodgeConjecture/Lemmas/AlgebraicTopology/PuncturedEuclideanFundamentalClass.lean}{\texttt{HodgeConjecture/Lemmas/AlgebraicTopology/PuncturedEuclideanFundamentalClass.lean}}.

The recursive acyclic-filler construction is adapted from Paul Lezeau's
\texttt{SphereSixComplex/Topology/SingularAffineSubdivisionPrism.lean} in \texttt{sphere-six-complex}, commit
\texttt{4a00b63d81550972eb77cecb0eff28ef142e40a0}.

The affine boundary map sends the explicit simplicial fundamental cycle to the oriented
boundary cycle used in the standard Euclidean local class.  A finite open-facet cover of
punctured Euclidean space supplies a face-coherent carrier back to the standard simplicial
sphere.  An acyclic-carrier prism and the rational small-chain equivalence upgrade these maps to
a chain-homotopy equivalence.  Consequently the affine boundary map is an isomorphism on
homology, and the explicit punctured and positive-even-dimensional local classes generate.

\begin{definition}[AlgebraicTopology.Singular.standardExtendedCoordinate]\label{def:AlgebraicTopology.Singular.standardExtendedCoordinate}
\lean{AlgebraicTopology.Singular.standardExtendedCoordinate}\leanok
\uses{def:AlgebraicTopology.Singular.StandardRealModel}
The affine barycentric coordinate indexed by the last vertex is normalized to zero; the
other coordinates are the ordinary coordinates on \texttt{ℝ\^{}d}.  Only comparisons between these
coordinates are used below, so this common-translation normalization is canonical.

\noindent\emph{Lean:} \texttt{def AlgebraicTopology.Singular.standardExtendedCoordinate}, \href{https://github.com/Paul-Lez/HodgeConjecture/blob/main/HodgeConjecture/Lemmas/AlgebraicTopology/PuncturedEuclideanFundamentalClass.lean#L56}{\texttt{HodgeConjecture/Lemmas/AlgebraicTopology/PuncturedEuclideanFundamentalClass.lean:56}}.
\end{definition}

\begin{definition}[AlgebraicTopology.Singular.standardPuncturedFacetCover]\label{def:AlgebraicTopology.Singular.standardPuncturedFacetCover}
\lean{AlgebraicTopology.Singular.standardPuncturedFacetCover}\leanok
\uses{def:AlgebraicTopology.Singular.standardExtendedCoordinate, def:AlgebraicTopology.Singular.standardPuncturedPair}
The \texttt{i}th member of the finite facet cover consists of points for which the \texttt{i}th extended
coordinate is not maximal.

\noindent\emph{Lean:} \texttt{def AlgebraicTopology.Singular.standardPuncturedFacetCover}, \href{https://github.com/Paul-Lez/HodgeConjecture/blob/main/HodgeConjecture/Lemmas/AlgebraicTopology/PuncturedEuclideanFundamentalClass.lean#L102}{\texttt{HodgeConjecture/Lemmas/AlgebraicTopology/PuncturedEuclideanFundamentalClass.lean:102}}.
\end{definition}

\begin{definition}[AlgebraicTopology.Singular.standardFacetIntersectionCenter]\label{def:AlgebraicTopology.Singular.standardFacetIntersectionCenter}
\lean{AlgebraicTopology.Singular.standardFacetIntersectionCenter}\leanok
\uses{def:AlgebraicTopology.Singular.StandardRealModel}
The coordinate vector whose extended coordinates are \texttt{-1} on \texttt{I} and \texttt{0} off \texttt{I}, up to
the common translation that makes the last extended coordinate zero.

\noindent\emph{Lean:} \texttt{def AlgebraicTopology.Singular.standardFacetIntersectionCenter}, \href{https://github.com/Paul-Lez/HodgeConjecture/blob/main/HodgeConjecture/Lemmas/AlgebraicTopology/PuncturedEuclideanFundamentalClass.lean#L144}{\texttt{HodgeConjecture/Lemmas/AlgebraicTopology/PuncturedEuclideanFundamentalClass.lean:144}}.
\end{definition}

\begin{definition}[AlgebraicTopology.Singular.standardPuncturedFacetIntersectionSet]\label{def:AlgebraicTopology.Singular.standardPuncturedFacetIntersectionSet}
\lean{AlgebraicTopology.Singular.standardPuncturedFacetIntersectionSet}\leanok
\uses{def:AlgebraicTopology.Singular.standardExtendedCoordinate}
The intersection of a finite collection of facet-cover conditions, expressed in the
ambient coordinate vector space.

\noindent\emph{Lean:} \texttt{def AlgebraicTopology.Singular.standardPuncturedFacetIntersectionSet}, \href{https://github.com/Paul-Lez/HodgeConjecture/blob/main/HodgeConjecture/Lemmas/AlgebraicTopology/PuncturedEuclideanFundamentalClass.lean#L169}{\texttt{HodgeConjecture/Lemmas/AlgebraicTopology/PuncturedEuclideanFundamentalClass.lean:169}}.
\end{definition}

\begin{definition}[AlgebraicTopology.Singular.standardPuncturedFacetIntersectionSubspace]\label{def:AlgebraicTopology.Singular.standardPuncturedFacetIntersectionSubspace}
\lean{AlgebraicTopology.Singular.standardPuncturedFacetIntersectionSubspace}\leanok
\uses{def:AlgebraicTopology.Singular.standardPuncturedFacetIntersectionSet, def:AlgebraicTopology.Singular.standardPuncturedPair}
The same facet intersection, now presented as a subspace of punctured coordinate space.

\noindent\emph{Lean:} \texttt{def AlgebraicTopology.Singular.standardPuncturedFacetIntersectionSubspace}, \href{https://github.com/Paul-Lez/HodgeConjecture/blob/main/HodgeConjecture/Lemmas/AlgebraicTopology/PuncturedEuclideanFundamentalClass.lean#L239}{\texttt{HodgeConjecture/Lemmas/AlgebraicTopology/PuncturedEuclideanFundamentalClass.lean:239}}.
\end{definition}

\begin{definition}[AlgebraicTopology.Singular.standardPuncturedFacetIntersectionHomeomorph]\label{def:AlgebraicTopology.Singular.standardPuncturedFacetIntersectionHomeomorph}
\lean{AlgebraicTopology.Singular.standardPuncturedFacetIntersectionHomeomorph}\leanok
\uses{def:AlgebraicTopology.Singular.standardPuncturedFacetIntersectionSubspace}
The ambient and nested-subspace presentations of a nonempty facet intersection are
homeomorphic.

\noindent\emph{Lean:} \texttt{def AlgebraicTopology.Singular.standardPuncturedFacetIntersectionHomeomorph}, \href{https://github.com/Paul-Lez/HodgeConjecture/blob/main/HodgeConjecture/Lemmas/AlgebraicTopology/PuncturedEuclideanFundamentalClass.lean#L244}{\texttt{HodgeConjecture/Lemmas/AlgebraicTopology/PuncturedEuclideanFundamentalClass.lean:244}}.
\end{definition}

\begin{definition}[AlgebraicTopology.Singular.standardPuncturedFacetIntersectionMapOfSubset]\label{def:AlgebraicTopology.Singular.standardPuncturedFacetIntersectionMapOfSubset}
\lean{AlgebraicTopology.Singular.standardPuncturedFacetIntersectionMapOfSubset}\leanok
\uses{def:AlgebraicTopology.Singular.standardPuncturedFacetIntersectionSubspace}
Inclusion of facet intersections, contravariant in their index sets.

\noindent\emph{Lean:} \texttt{def AlgebraicTopology.Singular.standardPuncturedFacetIntersectionMapOfSubset}, \href{https://github.com/Paul-Lez/HodgeConjecture/blob/main/HodgeConjecture/Lemmas/AlgebraicTopology/PuncturedEuclideanFundamentalClass.lean#L266}{\texttt{HodgeConjecture/Lemmas/AlgebraicTopology/PuncturedEuclideanFundamentalClass.lean:266}}.
\end{definition}

\begin{definition}[AlgebraicTopology.Singular.standardFacetCarrier]\label{def:AlgebraicTopology.Singular.standardFacetCarrier}
\lean{AlgebraicTopology.Singular.standardFacetCarrier}\leanok
\uses{def:AlgebraicTopology.Singular.coverSmallSingularSubcomplex, def:AlgebraicTopology.Singular.standardPuncturedFacetCover}
Cover indices through which a cover-small singular simplex factors.

\noindent\emph{Lean:} \texttt{def AlgebraicTopology.Singular.standardFacetCarrier}, \href{https://github.com/Paul-Lez/HodgeConjecture/blob/main/HodgeConjecture/Lemmas/AlgebraicTopology/PuncturedEuclideanFundamentalClass.lean#L407}{\texttt{HodgeConjecture/Lemmas/AlgebraicTopology/PuncturedEuclideanFundamentalClass.lean:407}}.
\end{definition}

\begin{definition}[AlgebraicTopology.Singular.standardFacetCarrierSourceIntersectionMap]\label{def:AlgebraicTopology.Singular.standardFacetCarrierSourceIntersectionMap}
\lean{AlgebraicTopology.Singular.standardFacetCarrierSourceIntersectionMap}\leanok
\uses{def:AlgebraicTopology.Singular.singularSimplicialMapLiftToSubset, def:AlgebraicTopology.Singular.standardFacetCarrier, def:AlgebraicTopology.Singular.standardPuncturedFacetIntersectionSubspace}
The source singular simplex lifted to the contractible intersection assigned by its
carrier.

\noindent\emph{Lean:} \texttt{def AlgebraicTopology.Singular.standardFacetCarrierSourceIntersectionMap}, \href{https://github.com/Paul-Lez/HodgeConjecture/blob/main/HodgeConjecture/Lemmas/AlgebraicTopology/PuncturedEuclideanFundamentalClass.lean#L493}{\texttt{HodgeConjecture/Lemmas/AlgebraicTopology/PuncturedEuclideanFundamentalClass.lean:493}}.
\end{definition}

\begin{definition}[AlgebraicTopology.Singular.standardFacetCarrierSourceSubdivisionIntersectionMap]\label{def:AlgebraicTopology.Singular.standardFacetCarrierSourceSubdivisionIntersectionMap}
\lean{AlgebraicTopology.Singular.standardFacetCarrierSourceSubdivisionIntersectionMap}\leanok
\uses{def:AlgebraicTopology.Singular.simplexSubdivisionLastVertex, def:AlgebraicTopology.Singular.standardFacetCarrierSourceIntersectionMap}
Subdivision followed by last vertex, transported along the lifted source simplex.

\noindent\emph{Lean:} \texttt{def AlgebraicTopology.Singular.standardFacetCarrierSourceSubdivisionIntersectionMap}, \href{https://github.com/Paul-Lez/HodgeConjecture/blob/main/HodgeConjecture/Lemmas/AlgebraicTopology/PuncturedEuclideanFundamentalClass.lean#L559}{\texttt{HodgeConjecture/Lemmas/AlgebraicTopology/PuncturedEuclideanFundamentalClass.lean:559}}.
\end{definition}

\begin{definition}[AlgebraicTopology.Singular.standardFacetCarrierAtFace]\label{def:AlgebraicTopology.Singular.standardFacetCarrierAtFace}
\lean{AlgebraicTopology.Singular.standardFacetCarrierAtFace}\leanok
\uses{def:AlgebraicTopology.Singular.coverSmallSingularSubcomplex, def:AlgebraicTopology.Singular.standardPuncturedFacetCover}
Cover indices which contain the images of every vertex in \texttt{A}.

\noindent\emph{Lean:} \texttt{def AlgebraicTopology.Singular.standardFacetCarrierAtFace}, \href{https://github.com/Paul-Lez/HodgeConjecture/blob/main/HodgeConjecture/Lemmas/AlgebraicTopology/PuncturedEuclideanFundamentalClass.lean#L587}{\texttt{HodgeConjecture/Lemmas/AlgebraicTopology/PuncturedEuclideanFundamentalClass.lean:587}}.
\end{definition}

\begin{definition}[AlgebraicTopology.Singular.standardFacetComplementCarrier]\label{def:AlgebraicTopology.Singular.standardFacetComplementCarrier}
\lean{AlgebraicTopology.Singular.standardFacetComplementCarrier}\leanok
\uses{def:AlgebraicTopology.Singular.standardFacetCarrierAtFace}
The complementary carrier is a nonempty proper set of vertices and grows with the face.

\noindent\emph{Lean:} \texttt{def AlgebraicTopology.Singular.standardFacetComplementCarrier}, \href{https://github.com/Paul-Lez/HodgeConjecture/blob/main/HodgeConjecture/Lemmas/AlgebraicTopology/PuncturedEuclideanFundamentalClass.lean#L654}{\texttt{HodgeConjecture/Lemmas/AlgebraicTopology/PuncturedEuclideanFundamentalClass.lean:654}}.
\end{definition}

\begin{definition}[AlgebraicTopology.Singular.standardFacetComplementCarrierMax]\label{def:AlgebraicTopology.Singular.standardFacetComplementCarrierMax}
\lean{AlgebraicTopology.Singular.standardFacetComplementCarrierMax}\leanok
\uses{def:AlgebraicTopology.Singular.standardFacetComplementCarrier}
The maximum vertex in the complementary carrier.

\noindent\emph{Lean:} \texttt{def AlgebraicTopology.Singular.standardFacetComplementCarrierMax}, \href{https://github.com/Paul-Lez/HodgeConjecture/blob/main/HodgeConjecture/Lemmas/AlgebraicTopology/PuncturedEuclideanFundamentalClass.lean#L682}{\texttt{HodgeConjecture/Lemmas/AlgebraicTopology/PuncturedEuclideanFundamentalClass.lean:682}}.
\end{definition}

\begin{definition}[AlgebraicTopology.Singular.standardFacetCarrierVertexOrderHom]\label{def:AlgebraicTopology.Singular.standardFacetCarrierVertexOrderHom}
\lean{AlgebraicTopology.Singular.standardFacetCarrierVertexOrderHom}\leanok
\uses{def:AlgebraicTopology.Singular.standardFacetComplementCarrierMax}
The largest vertex outside the face carrier, monotone as the face grows.

\noindent\emph{Lean:} \texttt{def AlgebraicTopology.Singular.standardFacetCarrierVertexOrderHom}, \href{https://github.com/Paul-Lez/HodgeConjecture/blob/main/HodgeConjecture/Lemmas/AlgebraicTopology/PuncturedEuclideanFundamentalClass.lean#L699}{\texttt{HodgeConjecture/Lemmas/AlgebraicTopology/PuncturedEuclideanFundamentalClass.lean:699}}.
\end{definition}

\begin{definition}[AlgebraicTopology.Singular.standardFacetCarrierSimplexMap]\label{def:AlgebraicTopology.Singular.standardFacetCarrierSimplexMap}
\lean{AlgebraicTopology.Singular.standardFacetCarrierSimplexMap}\leanok
\uses{def:AlgebraicTopology.Singular.standardFacetCarrierVertexOrderHom}
The carrier map from the subdivision of a simplex to the standard simplex on cover
indices.

\noindent\emph{Lean:} \texttt{def AlgebraicTopology.Singular.standardFacetCarrierSimplexMap}, \href{https://github.com/Paul-Lez/HodgeConjecture/blob/main/HodgeConjecture/Lemmas/AlgebraicTopology/PuncturedEuclideanFundamentalClass.lean#L708}{\texttt{HodgeConjecture/Lemmas/AlgebraicTopology/PuncturedEuclideanFundamentalClass.lean:708}}.
\end{definition}

\begin{definition}[AlgebraicTopology.Singular.standardFacetCarrierBoundarySimplexMap]\label{def:AlgebraicTopology.Singular.standardFacetCarrierBoundarySimplexMap}
\lean{AlgebraicTopology.Singular.standardFacetCarrierBoundarySimplexMap}\leanok
\uses{def:AlgebraicTopology.Singular.standardFacetCarrier, def:AlgebraicTopology.Singular.standardFacetCarrierSimplexMap}
The carrier map lifted to the simplicial boundary of the cover-index simplex.

\noindent\emph{Lean:} \texttt{def AlgebraicTopology.Singular.standardFacetCarrierBoundarySimplexMap}, \href{https://github.com/Paul-Lez/HodgeConjecture/blob/main/HodgeConjecture/Lemmas/AlgebraicTopology/PuncturedEuclideanFundamentalClass.lean#L768}{\texttt{HodgeConjecture/Lemmas/AlgebraicTopology/PuncturedEuclideanFundamentalClass.lean:768}}.
\end{definition}

\begin{definition}[AlgebraicTopology.Singular.standardFaceVertexOrderHom]\label{def:AlgebraicTopology.Singular.standardFaceVertexOrderHom}
\lean{AlgebraicTopology.Singular.standardFaceVertexOrderHom}\leanok
The vertex injection underlying a codimension-one face.

\noindent\emph{Lean:} \texttt{def AlgebraicTopology.Singular.standardFaceVertexOrderHom}, \href{https://github.com/Paul-Lez/HodgeConjecture/blob/main/HodgeConjecture/Lemmas/AlgebraicTopology/PuncturedEuclideanFundamentalClass.lean#L814}{\texttt{HodgeConjecture/Lemmas/AlgebraicTopology/PuncturedEuclideanFundamentalClass.lean:814}}.
\end{definition}

\begin{definition}[AlgebraicTopology.Singular.standardFacetCarrierSimplexChain]\label{def:AlgebraicTopology.Singular.standardFacetCarrierSimplexChain}
\lean{AlgebraicTopology.Singular.standardFacetCarrierSimplexChain}\leanok
\uses{def:AlgebraicTopology.Singular.standardFacetCarrierBoundarySimplexMap, def:AlgebraicTopology.Singular.subdividedSimplexFundamentalChain}
The signed barycentric carrier chain attached to one cover-small simplex.

\noindent\emph{Lean:} \texttt{def AlgebraicTopology.Singular.standardFacetCarrierSimplexChain}, \href{https://github.com/Paul-Lez/HodgeConjecture/blob/main/HodgeConjecture/Lemmas/AlgebraicTopology/PuncturedEuclideanFundamentalClass.lean#L933}{\texttt{HodgeConjecture/Lemmas/AlgebraicTopology/PuncturedEuclideanFundamentalClass.lean:933}}.
\end{definition}

\begin{definition}[AlgebraicTopology.Singular.standardFacetCarrierComponent]\label{def:AlgebraicTopology.Singular.standardFacetCarrierComponent}
\lean{AlgebraicTopology.Singular.standardFacetCarrierComponent}\leanok
\uses{def:AlgebraicTopology.Singular.CoverSmallIntegralSingularChainComplex, def:AlgebraicTopology.Singular.standardFacetCarrierSimplexChain}
The degreewise carrier map on integral cover-small chains.

\noindent\emph{Lean:} \texttt{def AlgebraicTopology.Singular.standardFacetCarrierComponent}, \href{https://github.com/Paul-Lez/HodgeConjecture/blob/main/HodgeConjecture/Lemmas/AlgebraicTopology/PuncturedEuclideanFundamentalClass.lean#L943}{\texttt{HodgeConjecture/Lemmas/AlgebraicTopology/PuncturedEuclideanFundamentalClass.lean:943}}.
\end{definition}

\begin{definition}[AlgebraicTopology.Singular.standardFacetCarrierChainMap]\label{def:AlgebraicTopology.Singular.standardFacetCarrierChainMap}
\lean{AlgebraicTopology.Singular.standardFacetCarrierChainMap}\leanok
\uses{def:AlgebraicTopology.Singular.finCastSuccEquivNotLast, def:AlgebraicTopology.Singular.insertOmittedVertexLastEquiv, def:AlgebraicTopology.Singular.standardFaceVertexOrderHom, def:AlgebraicTopology.Singular.standardFacetCarrierComponent, def:AlgebraicTopology.Singular.subdividedSimplexAlternatingFaceChain, def:AlgebraicTopology.Singular.subdividedSimplexInteriorFaceSum, def:AlgebraicTopology.Singular.subdividedSimplexOuterFaceSum}
The integral finite-cover carrier chain map into the simplicial sphere.

\noindent\emph{Lean:} \texttt{def AlgebraicTopology.Singular.standardFacetCarrierChainMap}, \href{https://github.com/Paul-Lez/HodgeConjecture/blob/main/HodgeConjecture/Lemmas/AlgebraicTopology/PuncturedEuclideanFundamentalClass.lean#L1000}{\texttt{HodgeConjecture/Lemmas/AlgebraicTopology/PuncturedEuclideanFundamentalClass.lean:1000}}.
\end{definition}

\begin{definition}[AlgebraicTopology.Singular.standardFacetCarrierAffineIntersectionMap]\label{def:AlgebraicTopology.Singular.standardFacetCarrierAffineIntersectionMap}
\lean{AlgebraicTopology.Singular.standardFacetCarrierAffineIntersectionMap}\leanok
\uses{def:AlgebraicTopology.Singular.singularSimplicialMapLiftToSubset, def:AlgebraicTopology.Singular.standardAffineBoundarySimplicialMap, def:AlgebraicTopology.Singular.standardFacetCarrierBoundarySimplexMap, def:AlgebraicTopology.Singular.standardPuncturedFacetIntersectionSubspace}
The affine carrier map lifted to the contractible intersection assigned to its source
simplex.

\noindent\emph{Lean:} \texttt{def AlgebraicTopology.Singular.standardFacetCarrierAffineIntersectionMap}, \href{https://github.com/Paul-Lez/HodgeConjecture/blob/main/HodgeConjecture/Lemmas/AlgebraicTopology/PuncturedEuclideanFundamentalClass.lean#L1051}{\texttt{HodgeConjecture/Lemmas/AlgebraicTopology/PuncturedEuclideanFundamentalClass.lean:1051}}.
\end{definition}

\begin{definition}[AlgebraicTopology.Singular.standardFacetCarrierIntersectionDiscrepancy]\label{def:AlgebraicTopology.Singular.standardFacetCarrierIntersectionDiscrepancy}
\lean{AlgebraicTopology.Singular.standardFacetCarrierIntersectionDiscrepancy}\leanok
\uses{def:AlgebraicTopology.Singular.standardFacetCarrierAffineIntersectionMap, def:AlgebraicTopology.Singular.standardFacetCarrierSourceSubdivisionIntersectionMap, def:AlgebraicTopology.Singular.subdividedSimplexFundamentalChain}
The difference, inside its assigned facet intersection, between the affine carrier and
subdivision followed by the last vertex of the original singular simplex.

\noindent\emph{Lean:} \texttt{def AlgebraicTopology.Singular.standardFacetCarrierIntersectionDiscrepancy}, \href{https://github.com/Paul-Lez/HodgeConjecture/blob/main/HodgeConjecture/Lemmas/AlgebraicTopology/PuncturedEuclideanFundamentalClass.lean#L1120}{\texttt{HodgeConjecture/Lemmas/AlgebraicTopology/PuncturedEuclideanFundamentalClass.lean:1120}}.
\end{definition}

\begin{definition}[AlgebraicTopology.Singular.standardFacetCarrierIntersectionDiscrepancyFaces]\label{def:AlgebraicTopology.Singular.standardFacetCarrierIntersectionDiscrepancyFaces}
\lean{AlgebraicTopology.Singular.standardFacetCarrierIntersectionDiscrepancyFaces}\leanok
\uses{def:AlgebraicTopology.Singular.standardFacetCarrierIntersectionDiscrepancy, def:AlgebraicTopology.Singular.standardPuncturedFacetIntersectionMapOfSubset}
The alternating sum of face discrepancies, all restricted into the intersection assigned
to the original simplex.

\noindent\emph{Lean:} \texttt{def AlgebraicTopology.Singular.standardFacetCarrierIntersectionDiscrepancyFaces}, \href{https://github.com/Paul-Lez/HodgeConjecture/blob/main/HodgeConjecture/Lemmas/AlgebraicTopology/PuncturedEuclideanFundamentalClass.lean#L1139}{\texttt{HodgeConjecture/Lemmas/AlgebraicTopology/PuncturedEuclideanFundamentalClass.lean:1139}}.
\end{definition}

\begin{definition}[AlgebraicTopology.Singular.standardFacetCarrierIntersectionChainInclusion]\label{def:AlgebraicTopology.Singular.standardFacetCarrierIntersectionChainInclusion}
\lean{AlgebraicTopology.Singular.standardFacetCarrierIntersectionChainInclusion}\leanok
\uses{def:AlgebraicTopology.Singular.standardFacetCarrier, def:AlgebraicTopology.Singular.standardPuncturedFacetIntersectionSubspace}
The chain map which includes an assigned carrier intersection into punctured coordinate
space.

\noindent\emph{Lean:} \texttt{def AlgebraicTopology.Singular.standardFacetCarrierIntersectionChainInclusion}, \href{https://github.com/Paul-Lez/HodgeConjecture/blob/main/HodgeConjecture/Lemmas/AlgebraicTopology/PuncturedEuclideanFundamentalClass.lean#L1251}{\texttt{HodgeConjecture/Lemmas/AlgebraicTopology/PuncturedEuclideanFundamentalClass.lean:1251}}.
\end{definition}

\begin{definition}[AlgebraicTopology.Singular.standardFacetCarrierDiscrepancySimplexChain]\label{def:AlgebraicTopology.Singular.standardFacetCarrierDiscrepancySimplexChain}
\lean{AlgebraicTopology.Singular.standardFacetCarrierDiscrepancySimplexChain}\leanok
\uses{def:AlgebraicTopology.Singular.standardFacetCarrierIntersectionChainInclusion, def:AlgebraicTopology.Singular.standardFacetCarrierIntersectionDiscrepancy}
The carrier discrepancy, now regarded as a chain in the ambient punctured space.

\noindent\emph{Lean:} \texttt{def AlgebraicTopology.Singular.standardFacetCarrierDiscrepancySimplexChain}, \href{https://github.com/Paul-Lez/HodgeConjecture/blob/main/HodgeConjecture/Lemmas/AlgebraicTopology/PuncturedEuclideanFundamentalClass.lean#L1270}{\texttt{HodgeConjecture/Lemmas/AlgebraicTopology/PuncturedEuclideanFundamentalClass.lean:1270}}.
\end{definition}

\begin{definition}[AlgebraicTopology.Singular.standardFacetCarrierDiscrepancyComponent]\label{def:AlgebraicTopology.Singular.standardFacetCarrierDiscrepancyComponent}
\lean{AlgebraicTopology.Singular.standardFacetCarrierDiscrepancyComponent}\leanok
\uses{def:AlgebraicTopology.Singular.CoverSmallIntegralSingularChainComplex, def:AlgebraicTopology.Singular.standardFacetCarrierDiscrepancySimplexChain}
The degreewise ambient discrepancy operator on cover-small integral chains.

\noindent\emph{Lean:} \texttt{def AlgebraicTopology.Singular.standardFacetCarrierDiscrepancyComponent}, \href{https://github.com/Paul-Lez/HodgeConjecture/blob/main/HodgeConjecture/Lemmas/AlgebraicTopology/PuncturedEuclideanFundamentalClass.lean#L1281}{\texttt{HodgeConjecture/Lemmas/AlgebraicTopology/PuncturedEuclideanFundamentalClass.lean:1281}}.
\end{definition}

\begin{definition}[AlgebraicTopology.Singular.standardFacetCarrierDiscrepancyChainMap]\label{def:AlgebraicTopology.Singular.standardFacetCarrierDiscrepancyChainMap}
\lean{AlgebraicTopology.Singular.standardFacetCarrierDiscrepancyChainMap}\leanok
\uses{def:AlgebraicTopology.Singular.finCastSuccEquivNotLast, def:AlgebraicTopology.Singular.insertOmittedVertexLastEquiv, def:AlgebraicTopology.Singular.standardFaceVertexOrderHom, def:AlgebraicTopology.Singular.standardFacetCarrierDiscrepancyComponent, def:AlgebraicTopology.Singular.standardFacetCarrierIntersectionDiscrepancyFaces, def:AlgebraicTopology.Singular.subdividedSimplexAlternatingFaceChain, def:AlgebraicTopology.Singular.subdividedSimplexInteriorFaceSum, def:AlgebraicTopology.Singular.subdividedSimplexOuterFaceSum}
The carrier discrepancy as a chain map from cover-small chains to ambient singular
chains.

\noindent\emph{Lean:} \texttt{def AlgebraicTopology.Singular.standardFacetCarrierDiscrepancyChainMap}, \href{https://github.com/Paul-Lez/HodgeConjecture/blob/main/HodgeConjecture/Lemmas/AlgebraicTopology/PuncturedEuclideanFundamentalClass.lean#L1365}{\texttt{HodgeConjecture/Lemmas/AlgebraicTopology/PuncturedEuclideanFundamentalClass.lean:1365}}.
\end{definition}

\begin{definition}[AlgebraicTopology.Singular.StandardFacetCarrierPrismFamily]\label{def:AlgebraicTopology.Singular.StandardFacetCarrierPrismFamily}
\lean{AlgebraicTopology.Singular.StandardFacetCarrierPrismFamily}\leanok
\uses{def:AlgebraicTopology.Singular.standardFacetCarrier, def:AlgebraicTopology.Singular.standardPuncturedFacetIntersectionSubspace}
A family of local carrier prisms, one degree-raising chain for every cover-small
simplex.

\noindent\emph{Lean:} \texttt{abbrev AlgebraicTopology.Singular.StandardFacetCarrierPrismFamily}, \href{https://github.com/Paul-Lez/HodgeConjecture/blob/main/HodgeConjecture/Lemmas/AlgebraicTopology/PuncturedEuclideanFundamentalClass.lean#L1381}{\texttt{HodgeConjecture/Lemmas/AlgebraicTopology/PuncturedEuclideanFundamentalClass.lean:1381}}.
\end{definition}

\begin{definition}[AlgebraicTopology.Singular.standardFacetCarrierPrismFaces]\label{def:AlgebraicTopology.Singular.standardFacetCarrierPrismFaces}
\lean{AlgebraicTopology.Singular.standardFacetCarrierPrismFaces}\leanok
\uses{def:AlgebraicTopology.Singular.StandardFacetCarrierPrismFamily, def:AlgebraicTopology.Singular.standardPuncturedFacetIntersectionMapOfSubset}
The alternating sum of the already constructed face prisms, restricted into the carrier
intersection assigned to the original simplex.

\noindent\emph{Lean:} \texttt{def AlgebraicTopology.Singular.standardFacetCarrierPrismFaces}, \href{https://github.com/Paul-Lez/HodgeConjecture/blob/main/HodgeConjecture/Lemmas/AlgebraicTopology/PuncturedEuclideanFundamentalClass.lean#L1393}{\texttt{HodgeConjecture/Lemmas/AlgebraicTopology/PuncturedEuclideanFundamentalClass.lean:1393}}.
\end{definition}

\begin{definition}[AlgebraicTopology.Singular.standardFacetCarrierPrismSimplexChain]\label{def:AlgebraicTopology.Singular.standardFacetCarrierPrismSimplexChain}
\lean{AlgebraicTopology.Singular.standardFacetCarrierPrismSimplexChain}\leanok
\uses{def:AlgebraicTopology.Singular.StandardFacetCarrierPrismFamily, def:AlgebraicTopology.Singular.standardFacetCarrierIntersectionChainInclusion}
One local prism chain, included into ambient punctured coordinate space.

\noindent\emph{Lean:} \texttt{def AlgebraicTopology.Singular.standardFacetCarrierPrismSimplexChain}, \href{https://github.com/Paul-Lez/HodgeConjecture/blob/main/HodgeConjecture/Lemmas/AlgebraicTopology/PuncturedEuclideanFundamentalClass.lean#L1413}{\texttt{HodgeConjecture/Lemmas/AlgebraicTopology/PuncturedEuclideanFundamentalClass.lean:1413}}.
\end{definition}

\begin{definition}[AlgebraicTopology.Singular.standardFacetCarrierPrismComponent]\label{def:AlgebraicTopology.Singular.standardFacetCarrierPrismComponent}
\lean{AlgebraicTopology.Singular.standardFacetCarrierPrismComponent}\leanok
\uses{def:AlgebraicTopology.Singular.CoverSmallIntegralSingularChainComplex, def:AlgebraicTopology.Singular.standardFacetCarrierPrismSimplexChain}
The global degree-raising carrier-prism operator induced by local prism chains.

\noindent\emph{Lean:} \texttt{def AlgebraicTopology.Singular.standardFacetCarrierPrismComponent}, \href{https://github.com/Paul-Lez/HodgeConjecture/blob/main/HodgeConjecture/Lemmas/AlgebraicTopology/PuncturedEuclideanFundamentalClass.lean#L1423}{\texttt{HodgeConjecture/Lemmas/AlgebraicTopology/PuncturedEuclideanFundamentalClass.lean:1423}}.
\end{definition}

\begin{definition}[AlgebraicTopology.Singular.standardFacetCarrierPrismResidual]\label{def:AlgebraicTopology.Singular.standardFacetCarrierPrismResidual}
\lean{AlgebraicTopology.Singular.standardFacetCarrierPrismResidual}\leanok
\uses{def:AlgebraicTopology.Singular.standardFacetCarrierIntersectionDiscrepancy, def:AlgebraicTopology.Singular.standardFacetCarrierPrismFaces}
The local residual left after subtracting the already constructed face prisms.

\noindent\emph{Lean:} \texttt{def AlgebraicTopology.Singular.standardFacetCarrierPrismResidual}, \href{https://github.com/Paul-Lez/HodgeConjecture/blob/main/HodgeConjecture/Lemmas/AlgebraicTopology/PuncturedEuclideanFundamentalClass.lean#L1630}{\texttt{HodgeConjecture/Lemmas/AlgebraicTopology/PuncturedEuclideanFundamentalClass.lean:1630}}.
\end{definition}

\begin{definition}[AlgebraicTopology.Singular.StandardFacetCarrierPrismEquation]\label{def:AlgebraicTopology.Singular.StandardFacetCarrierPrismEquation}
\lean{AlgebraicTopology.Singular.StandardFacetCarrierPrismEquation}\leanok
\uses{def:AlgebraicTopology.Singular.standardFacetCarrierIntersectionDiscrepancy, def:AlgebraicTopology.Singular.standardFacetCarrierPrismFaces}
The local acyclic-carrier prism equation in one degree.

\noindent\emph{Lean:} \texttt{def AlgebraicTopology.Singular.StandardFacetCarrierPrismEquation}, \href{https://github.com/Paul-Lez/HodgeConjecture/blob/main/HodgeConjecture/Lemmas/AlgebraicTopology/PuncturedEuclideanFundamentalClass.lean#L1787}{\texttt{HodgeConjecture/Lemmas/AlgebraicTopology/PuncturedEuclideanFundamentalClass.lean:1787}}.
\end{definition}

\begin{definition}[AlgebraicTopology.Singular.standardPuncturedFacetIntersectionIntegralTotalCycleFiller]\label{def:AlgebraicTopology.Singular.standardPuncturedFacetIntersectionIntegralTotalCycleFiller}
\lean{AlgebraicTopology.Singular.standardPuncturedFacetIntersectionIntegralTotalCycleFiller}\leanok
\uses{def:AlgebraicTopology.Singular.integralSingularHomologyEquivOfHomotopyEquiv, def:AlgebraicTopology.Singular.standardFacetIntersectionCenter, def:AlgebraicTopology.Singular.standardPuncturedFacetIntersectionHomeomorph}
A total selected filler for a positive-degree cycle in a proper facet intersection.  It
returns zero when its input is not a cycle, so it can be used in a structural recursion whose
cycle proof is established afterwards.

\noindent\emph{Lean:} \texttt{def AlgebraicTopology.Singular.standardPuncturedFacetIntersectionIntegralTotalCycleFiller}, \href{https://github.com/Paul-Lez/HodgeConjecture/blob/main/HodgeConjecture/Lemmas/AlgebraicTopology/PuncturedEuclideanFundamentalClass.lean#L1825}{\texttt{HodgeConjecture/Lemmas/AlgebraicTopology/PuncturedEuclideanFundamentalClass.lean:1825}}.
\end{definition}

\begin{definition}[AlgebraicTopology.Singular.standardFacetCarrierPrism]\label{def:AlgebraicTopology.Singular.standardFacetCarrierPrism}
\lean{AlgebraicTopology.Singular.standardFacetCarrierPrism}\leanok
\uses{def:AlgebraicTopology.Singular.standardFacetCarrierIntersectionDiscrepancy, def:AlgebraicTopology.Singular.standardPuncturedFacetIntersectionIntegralTotalCycleFiller, def:AlgebraicTopology.Singular.standardPuncturedFacetIntersectionMapOfSubset}
The canonical local acyclic-carrier prism, recursively obtained by filling the discrepancy
left after the prisms on all faces.

\noindent\emph{Lean:} \texttt{def AlgebraicTopology.Singular.standardFacetCarrierPrism}, \href{https://github.com/Paul-Lez/HodgeConjecture/blob/main/HodgeConjecture/Lemmas/AlgebraicTopology/PuncturedEuclideanFundamentalClass.lean#L1871}{\texttt{HodgeConjecture/Lemmas/AlgebraicTopology/PuncturedEuclideanFundamentalClass.lean:1871}}.
\end{definition}

\begin{definition}[AlgebraicTopology.Singular.standardFacetCarrierPrismHom]\label{def:AlgebraicTopology.Singular.standardFacetCarrierPrismHom}
\lean{AlgebraicTopology.Singular.standardFacetCarrierPrismHom}\leanok
\uses{def:AlgebraicTopology.Singular.standardFacetCarrierPrism, def:AlgebraicTopology.Singular.standardFacetCarrierPrismComponent}
The degree-raising carrier-prism family, extended by zero away from adjacent degrees.

\noindent\emph{Lean:} \texttt{def AlgebraicTopology.Singular.standardFacetCarrierPrismHom}, \href{https://github.com/Paul-Lez/HodgeConjecture/blob/main/HodgeConjecture/Lemmas/AlgebraicTopology/PuncturedEuclideanFundamentalClass.lean#L1951}{\texttt{HodgeConjecture/Lemmas/AlgebraicTopology/PuncturedEuclideanFundamentalClass.lean:1951}}.
\end{definition}

\begin{definition}[AlgebraicTopology.Singular.standardFacetCarrierDiscrepancyIntegralHomotopy]\label{def:AlgebraicTopology.Singular.standardFacetCarrierDiscrepancyIntegralHomotopy}
\lean{AlgebraicTopology.Singular.standardFacetCarrierDiscrepancyIntegralHomotopy}\leanok
\uses{def:AlgebraicTopology.Singular.StandardFacetCarrierPrismEquation, def:AlgebraicTopology.Singular.singularSubsetIntegralChainRetractionComponent, def:AlgebraicTopology.Singular.standardFacetCarrierDiscrepancyChainMap, def:AlgebraicTopology.Singular.standardFacetCarrierPrismHom, def:AlgebraicTopology.Singular.standardFacetCarrierPrismResidual}
The ambient carrier discrepancy is integrally chain-homotopic to zero.

\noindent\emph{Lean:} \texttt{def AlgebraicTopology.Singular.standardFacetCarrierDiscrepancyIntegralHomotopy}, \href{https://github.com/Paul-Lez/HodgeConjecture/blob/main/HodgeConjecture/Lemmas/AlgebraicTopology/PuncturedEuclideanFundamentalClass.lean#L1970}{\texttt{HodgeConjecture/Lemmas/AlgebraicTopology/PuncturedEuclideanFundamentalClass.lean:1970}}.
\end{definition}

\begin{definition}[AlgebraicTopology.Singular.standardFacetCarrierDiscrepancyRationalHomotopy]\label{def:AlgebraicTopology.Singular.standardFacetCarrierDiscrepancyRationalHomotopy}
\lean{AlgebraicTopology.Singular.standardFacetCarrierDiscrepancyRationalHomotopy}\leanok
\uses{def:AlgebraicTopology.Singular.rationalizeSimplicialChainHomotopy, def:AlgebraicTopology.Singular.standardFacetCarrierDiscrepancyIntegralHomotopy}
Rationalization of the acyclic-carrier homotopy.

\noindent\emph{Lean:} \texttt{def AlgebraicTopology.Singular.standardFacetCarrierDiscrepancyRationalHomotopy}, \href{https://github.com/Paul-Lez/HodgeConjecture/blob/main/HodgeConjecture/Lemmas/AlgebraicTopology/PuncturedEuclideanFundamentalClass.lean#L1998}{\texttt{HodgeConjecture/Lemmas/AlgebraicTopology/PuncturedEuclideanFundamentalClass.lean:1998}}.
\end{definition}

\begin{definition}[AlgebraicTopology.Singular.standardAffineBoundaryIntegralChainMap]\label{def:AlgebraicTopology.Singular.standardAffineBoundaryIntegralChainMap}
\lean{AlgebraicTopology.Singular.standardAffineBoundaryIntegralChainMap}\leanok
\uses{def:AlgebraicTopology.Singular.standardAffineBoundarySimplicialMap}
The integral affine-boundary chain map.

\noindent\emph{Lean:} \texttt{def AlgebraicTopology.Singular.standardAffineBoundaryIntegralChainMap}, \href{https://github.com/Paul-Lez/HodgeConjecture/blob/main/HodgeConjecture/Lemmas/AlgebraicTopology/PuncturedEuclideanFundamentalClass.lean#L2016}{\texttt{HodgeConjecture/Lemmas/AlgebraicTopology/PuncturedEuclideanFundamentalClass.lean:2016}}.
\end{definition}

\begin{definition}[AlgebraicTopology.Singular.standardFacetSmallBarycentricLastVertexIntegralChainMap]\label{def:AlgebraicTopology.Singular.standardFacetSmallBarycentricLastVertexIntegralChainMap}
\lean{AlgebraicTopology.Singular.standardFacetSmallBarycentricLastVertexIntegralChainMap}\leanok
\uses{def:AlgebraicTopology.Singular.CoverSmallIntegralSingularChainComplex, def:AlgebraicTopology.Singular.barycentricSubdivisionChainMapCanonical, def:AlgebraicTopology.Singular.standardPuncturedFacetCover, def:AlgebraicTopology.Singular.subdivisionLastVertexChainMap}
Barycentric subdivision followed by last vertex on the facet-small simplicial set.

\noindent\emph{Lean:} \texttt{def AlgebraicTopology.Singular.standardFacetSmallBarycentricLastVertexIntegralChainMap}, \href{https://github.com/Paul-Lez/HodgeConjecture/blob/main/HodgeConjecture/Lemmas/AlgebraicTopology/PuncturedEuclideanFundamentalClass.lean#L2024}{\texttt{HodgeConjecture/Lemmas/AlgebraicTopology/PuncturedEuclideanFundamentalClass.lean:2024}}.
\end{definition}

\begin{definition}[AlgebraicTopology.Singular.standardAffineBoundaryToFacetSmall]\label{def:AlgebraicTopology.Singular.standardAffineBoundaryToFacetSmall}
\lean{AlgebraicTopology.Singular.standardAffineBoundaryToFacetSmall}\leanok
\uses{def:AlgebraicTopology.Singular.coverSmallSingularSubcomplex, def:AlgebraicTopology.Singular.standardAffineBoundarySimplicialMap, def:AlgebraicTopology.Singular.standardPuncturedFacetCover}
The affine boundary simplicial map factored through cover-small singular simplices.

\noindent\emph{Lean:} \texttt{def AlgebraicTopology.Singular.standardAffineBoundaryToFacetSmall}, \href{https://github.com/Paul-Lez/HodgeConjecture/blob/main/HodgeConjecture/Lemmas/AlgebraicTopology/PuncturedEuclideanFundamentalClass.lean#L2201}{\texttt{HodgeConjecture/Lemmas/AlgebraicTopology/PuncturedEuclideanFundamentalClass.lean:2201}}.
\end{definition}

\begin{definition}[AlgebraicTopology.Singular.standardFacetCarrierRationalChainMap]\label{def:AlgebraicTopology.Singular.standardFacetCarrierRationalChainMap}
\lean{AlgebraicTopology.Singular.standardFacetCarrierRationalChainMap}\leanok
\uses{def:AlgebraicTopology.Singular.CoverSmallRationalSingularChainComplex, def:AlgebraicTopology.Singular.rationalizeSimplicialChainMap, def:AlgebraicTopology.Singular.standardFacetCarrierChainMap}
The rationalized carrier chain map.

\noindent\emph{Lean:} \texttt{def AlgebraicTopology.Singular.standardFacetCarrierRationalChainMap}, \href{https://github.com/Paul-Lez/HodgeConjecture/blob/main/HodgeConjecture/Lemmas/AlgebraicTopology/PuncturedEuclideanFundamentalClass.lean#L2377}{\texttt{HodgeConjecture/Lemmas/AlgebraicTopology/PuncturedEuclideanFundamentalClass.lean:2377}}.
\end{definition}

\begin{definition}[AlgebraicTopology.Singular.standardBarycentricLastVertexRationalChainMap]\label{def:AlgebraicTopology.Singular.standardBarycentricLastVertexRationalChainMap}
\lean{AlgebraicTopology.Singular.standardBarycentricLastVertexRationalChainMap}\leanok
\uses{def:AlgebraicTopology.Singular.barycentricSubdivisionChainMapCanonical, def:AlgebraicTopology.Singular.rationalizeSimplicialChainMap, def:AlgebraicTopology.Singular.subdivisionLastVertexChainMap}
The rational subdivision--last-vertex endomorphism.

\noindent\emph{Lean:} \texttt{def AlgebraicTopology.Singular.standardBarycentricLastVertexRationalChainMap}, \href{https://github.com/Paul-Lez/HodgeConjecture/blob/main/HodgeConjecture/Lemmas/AlgebraicTopology/PuncturedEuclideanFundamentalClass.lean#L2387}{\texttt{HodgeConjecture/Lemmas/AlgebraicTopology/PuncturedEuclideanFundamentalClass.lean:2387}}.
\end{definition}

\begin{definition}[AlgebraicTopology.Singular.standardBarycentricLastVertexRationalHomotopy]\label{def:AlgebraicTopology.Singular.standardBarycentricLastVertexRationalHomotopy}
\lean{AlgebraicTopology.Singular.standardBarycentricLastVertexRationalHomotopy}\leanok
\uses{def:AlgebraicTopology.Singular.barycentricSubdivisionLastVertexHomotopyCanonical, def:AlgebraicTopology.Singular.rationalizeSimplicialChainHomotopy, def:AlgebraicTopology.Singular.standardBarycentricLastVertexRationalChainMap}
The rational subdivision--last-vertex endomorphism is chain-homotopic to the identity.

\noindent\emph{Lean:} \texttt{def AlgebraicTopology.Singular.standardBarycentricLastVertexRationalHomotopy}, \href{https://github.com/Paul-Lez/HodgeConjecture/blob/main/HodgeConjecture/Lemmas/AlgebraicTopology/PuncturedEuclideanFundamentalClass.lean#L2395}{\texttt{HodgeConjecture/Lemmas/AlgebraicTopology/PuncturedEuclideanFundamentalClass.lean:2395}}.
\end{definition}

\begin{definition}[AlgebraicTopology.Singular.standardFacetSmallBarycentricLastVertexRationalChainMap]\label{def:AlgebraicTopology.Singular.standardFacetSmallBarycentricLastVertexRationalChainMap}
\lean{AlgebraicTopology.Singular.standardFacetSmallBarycentricLastVertexRationalChainMap}\leanok
\uses{def:AlgebraicTopology.Singular.CoverSmallRationalSingularChainComplex, def:AlgebraicTopology.Singular.rationalizeSimplicialChainMap, def:AlgebraicTopology.Singular.standardFacetSmallBarycentricLastVertexIntegralChainMap}
Rational subdivision followed by last vertex on the facet-small simplicial set.

\noindent\emph{Lean:} \texttt{def AlgebraicTopology.Singular.standardFacetSmallBarycentricLastVertexRationalChainMap}, \href{https://github.com/Paul-Lez/HodgeConjecture/blob/main/HodgeConjecture/Lemmas/AlgebraicTopology/PuncturedEuclideanFundamentalClass.lean#L2444}{\texttt{HodgeConjecture/Lemmas/AlgebraicTopology/PuncturedEuclideanFundamentalClass.lean:2444}}.
\end{definition}

\begin{definition}[AlgebraicTopology.Singular.standardFacetSmallBarycentricLastVertexRationalHomotopy]\label{def:AlgebraicTopology.Singular.standardFacetSmallBarycentricLastVertexRationalHomotopy}
\lean{AlgebraicTopology.Singular.standardFacetSmallBarycentricLastVertexRationalHomotopy}\leanok
\uses{def:AlgebraicTopology.Singular.barycentricSubdivisionLastVertexHomotopyCanonical, def:AlgebraicTopology.Singular.rationalizeSimplicialChainHomotopy, def:AlgebraicTopology.Singular.standardFacetSmallBarycentricLastVertexRationalChainMap}
On facet-small rational chains, subdivision followed by last vertex is homotopic to the
identity.

\noindent\emph{Lean:} \texttt{def AlgebraicTopology.Singular.standardFacetSmallBarycentricLastVertexRationalHomotopy}, \href{https://github.com/Paul-Lez/HodgeConjecture/blob/main/HodgeConjecture/Lemmas/AlgebraicTopology/PuncturedEuclideanFundamentalClass.lean#L2499}{\texttt{HodgeConjecture/Lemmas/AlgebraicTopology/PuncturedEuclideanFundamentalClass.lean:2499}}.
\end{definition}

\begin{definition}[AlgebraicTopology.Singular.standardFacetCarrierAffineRationalHomotopy]\label{def:AlgebraicTopology.Singular.standardFacetCarrierAffineRationalHomotopy}
\lean{AlgebraicTopology.Singular.standardFacetCarrierAffineRationalHomotopy}\leanok
\uses{def:AlgebraicTopology.Singular.coverSmallIntegralSingularChainInclusion, def:AlgebraicTopology.Singular.coverSmallRationalSingularChainInclusion, def:AlgebraicTopology.Singular.standardAffineBoundaryChainMap, def:AlgebraicTopology.Singular.standardAffineBoundaryIntegralChainMap, def:AlgebraicTopology.Singular.standardFacetCarrierDiscrepancyRationalHomotopy, def:AlgebraicTopology.Singular.standardFacetCarrierRationalChainMap, def:AlgebraicTopology.Singular.standardFacetSmallBarycentricLastVertexRationalHomotopy}
The rational carrier followed by affine realization is homotopic to the inclusion of
facet-small rational chains.

\noindent\emph{Lean:} \texttt{def AlgebraicTopology.Singular.standardFacetCarrierAffineRationalHomotopy}, \href{https://github.com/Paul-Lez/HodgeConjecture/blob/main/HodgeConjecture/Lemmas/AlgebraicTopology/PuncturedEuclideanFundamentalClass.lean#L2515}{\texttt{HodgeConjecture/Lemmas/AlgebraicTopology/PuncturedEuclideanFundamentalClass.lean:2515}}.
\end{definition}

\begin{definition}[AlgebraicTopology.Singular.standardAffineBoundaryChainRetraction]\label{def:AlgebraicTopology.Singular.standardAffineBoundaryChainRetraction}
\lean{AlgebraicTopology.Singular.standardAffineBoundaryChainRetraction}\leanok
\uses{def:AlgebraicTopology.Singular.coverSmallRationalChainHomotopyEquiv_of_openCover, def:AlgebraicTopology.Singular.standardFacetCarrierRationalChainMap}
A chain-level left inverse of affine realization, obtained by smallifying with respect to
the facet cover and then applying the explicit carrier map.

\noindent\emph{Lean:} \texttt{def AlgebraicTopology.Singular.standardAffineBoundaryChainRetraction}, \href{https://github.com/Paul-Lez/HodgeConjecture/blob/main/HodgeConjecture/Lemmas/AlgebraicTopology/PuncturedEuclideanFundamentalClass.lean#L2535}{\texttt{HodgeConjecture/Lemmas/AlgebraicTopology/PuncturedEuclideanFundamentalClass.lean:2535}}.
\end{definition}

\begin{definition}[AlgebraicTopology.Singular.standardAffineBoundaryChainRetractionHomotopy]\label{def:AlgebraicTopology.Singular.standardAffineBoundaryChainRetractionHomotopy}
\lean{AlgebraicTopology.Singular.standardAffineBoundaryChainRetractionHomotopy}\leanok
\uses{def:AlgebraicTopology.Singular.standardAffineBoundaryChainMap, def:AlgebraicTopology.Singular.standardAffineBoundaryChainRetraction, def:AlgebraicTopology.Singular.standardAffineBoundaryToFacetSmall, def:AlgebraicTopology.Singular.standardBarycentricLastVertexRationalHomotopy}
Affine realization followed by the explicit carrier retraction is chain-homotopic to the
identity.

\noindent\emph{Lean:} \texttt{def AlgebraicTopology.Singular.standardAffineBoundaryChainRetractionHomotopy}, \href{https://github.com/Paul-Lez/HodgeConjecture/blob/main/HodgeConjecture/Lemmas/AlgebraicTopology/PuncturedEuclideanFundamentalClass.lean#L2547}{\texttt{HodgeConjecture/Lemmas/AlgebraicTopology/PuncturedEuclideanFundamentalClass.lean:2547}}.
\end{definition}

\begin{definition}[AlgebraicTopology.Singular.standardAffineBoundaryChainCoretractionHomotopy]\label{def:AlgebraicTopology.Singular.standardAffineBoundaryChainCoretractionHomotopy}
\lean{AlgebraicTopology.Singular.standardAffineBoundaryChainCoretractionHomotopy}\leanok
\uses{def:AlgebraicTopology.Singular.standardAffineBoundaryChainRetraction, def:AlgebraicTopology.Singular.standardFacetCarrierAffineRationalHomotopy}
The explicit carrier retraction followed by affine realization is chain-homotopic to the
identity on punctured-space rational chains.

\noindent\emph{Lean:} \texttt{def AlgebraicTopology.Singular.standardAffineBoundaryChainCoretractionHomotopy}, \href{https://github.com/Paul-Lez/HodgeConjecture/blob/main/HodgeConjecture/Lemmas/AlgebraicTopology/PuncturedEuclideanFundamentalClass.lean#L2578}{\texttt{HodgeConjecture/Lemmas/AlgebraicTopology/PuncturedEuclideanFundamentalClass.lean:2578}}.
\end{definition}

\begin{definition}[AlgebraicTopology.Singular.standardAffineBoundaryChainHomotopyEquiv]\label{def:AlgebraicTopology.Singular.standardAffineBoundaryChainHomotopyEquiv}
\lean{AlgebraicTopology.Singular.standardAffineBoundaryChainHomotopyEquiv}\leanok
\uses{def:AlgebraicTopology.Singular.standardAffineBoundaryChainCoretractionHomotopy, def:AlgebraicTopology.Singular.standardAffineBoundaryChainRetractionHomotopy}
Affine realization of the standard simplicial boundary is a rational chain-homotopy
equivalence onto punctured Euclidean space.

\noindent\emph{Lean:} \texttt{def AlgebraicTopology.Singular.standardAffineBoundaryChainHomotopyEquiv}, \href{https://github.com/Paul-Lez/HodgeConjecture/blob/main/HodgeConjecture/Lemmas/AlgebraicTopology/PuncturedEuclideanFundamentalClass.lean#L2606}{\texttt{HodgeConjecture/Lemmas/AlgebraicTopology/PuncturedEuclideanFundamentalClass.lean:2606}}.
\end{definition}

\begin{definition}[AlgebraicTopology.Singular.standardAffineBoundaryHomologyMap]\label{def:AlgebraicTopology.Singular.standardAffineBoundaryHomologyMap}
\lean{AlgebraicTopology.Singular.standardAffineBoundaryHomologyMap}\leanok
\uses{def:AlgebraicTopology.Singular.standardAffineBoundaryChainMap}
The map on rational homology induced by the affine realization of a standard simplicial
sphere in punctured Euclidean space.

\noindent\emph{Lean:} \texttt{def AlgebraicTopology.Singular.standardAffineBoundaryHomologyMap}, \href{https://github.com/Paul-Lez/HodgeConjecture/blob/main/HodgeConjecture/Lemmas/AlgebraicTopology/PuncturedEuclideanFundamentalClass.lean#L2618}{\texttt{HodgeConjecture/Lemmas/AlgebraicTopology/PuncturedEuclideanFundamentalClass.lean:2618}}.
\end{definition}

\begin{definition}[AlgebraicTopology.Singular.standardAffineBoundaryHomologyRetraction]\label{def:AlgebraicTopology.Singular.standardAffineBoundaryHomologyRetraction}
\lean{AlgebraicTopology.Singular.standardAffineBoundaryHomologyRetraction}\leanok
\uses{def:AlgebraicTopology.Singular.standardAffineBoundaryChainRetraction}
The induced map on homology back to the simplicial boundary.

\noindent\emph{Lean:} \texttt{def AlgebraicTopology.Singular.standardAffineBoundaryHomologyRetraction}, \href{https://github.com/Paul-Lez/HodgeConjecture/blob/main/HodgeConjecture/Lemmas/AlgebraicTopology/PuncturedEuclideanFundamentalClass.lean#L2628}{\texttt{HodgeConjecture/Lemmas/AlgebraicTopology/PuncturedEuclideanFundamentalClass.lean:2628}}.
\end{definition}

\section{Injective sheaves are flasque}\label{sec:HodgeConjecture.Lemmas.AlgebraicTopology.InjectiveFlasque}
\noindent\emph{File:} \href{https://github.com/Paul-Lez/HodgeConjecture/blob/main/HodgeConjecture/Lemmas/AlgebraicTopology/InjectiveFlasque.lean}{\texttt{HodgeConjecture/Lemmas/AlgebraicTopology/InjectiveFlasque.lean}}.

For an open set \texttt{U}, sections of an additive sheaf \texttt{F} over \texttt{U} are represented by morphisms
from the sheafification of the free abelian presheaf on \texttt{yoneda.obj U}.  An inclusion of open
sets induces a monomorphism between these representing sheaves.  The extension property of an
injective sheaf therefore makes every restriction map surjective.

\begin{definition}[TopCat.Sheaf.freeAbelianYonedaSheaf]\label{def:TopCat.Sheaf.freeAbelianYonedaSheaf}
\lean{TopCat.Sheaf.freeAbelianYonedaSheaf}\leanok
The sheafified free abelian Yoneda functor on the poset of open subsets.

\noindent\emph{Lean:} \texttt{abbrev TopCat.Sheaf.freeAbelianYonedaSheaf}, \href{https://github.com/Paul-Lez/HodgeConjecture/blob/main/HodgeConjecture/Lemmas/AlgebraicTopology/InjectiveFlasque.lean#L40}{\texttt{HodgeConjecture/Lemmas/AlgebraicTopology/InjectiveFlasque.lean:40}}.
\end{definition}

\begin{definition}[TopCat.Sheaf.freeAbelianYonedaSheafHomEquiv]\label{def:TopCat.Sheaf.freeAbelianYonedaSheafHomEquiv}
\lean{TopCat.Sheaf.freeAbelianYonedaSheafHomEquiv}\leanok
\uses{def:TopCat.Sheaf.freeAbelianYonedaSheaf}
Morphisms from the free abelian Yoneda sheaf represent sections over the corresponding
open subset.

\noindent\emph{Lean:} \texttt{def TopCat.Sheaf.freeAbelianYonedaSheafHomEquiv}, \href{https://github.com/Paul-Lez/HodgeConjecture/blob/main/HodgeConjecture/Lemmas/AlgebraicTopology/InjectiveFlasque.lean#L48}{\texttt{HodgeConjecture/Lemmas/AlgebraicTopology/InjectiveFlasque.lean:48}}.
\end{definition}

\section{Descent for singular cochains}\label{sec:HodgeConjecture.Lemmas.AlgebraicTopology.SingularCochainSheafFlasque}
\noindent\emph{File:} \href{https://github.com/Paul-Lez/HodgeConjecture/blob/main/HodgeConjecture/Lemmas/AlgebraicTopology/SingularCochainSheafFlasque.lean}{\texttt{HodgeConjecture/Lemmas/AlgebraicTopology/SingularCochainSheafFlasque.lean}}.

Singular cochains in every degree are arbitrary functions on singular simplices. This gives
existence of gluings for compatible raw cochains, global raw representatives for the first plus
construction, and flasqueness of that first plus presheaf. The first plus construction is only the
separated reflection in general. Its flasqueness is therefore not a sheaf-level acyclicity or
derived-global-sections comparison theorem.

In degree zero, singular simplices are points and the raw presheaf is already the sheaf of all
functions. Consequently its sheafification is flasque in degree zero. No positive-degree
flasqueness claim is made for the second plus construction or abstract sheafification.

\begin{definition}[AlgebraicTopology.Singular.OpenSimplex]\label{def:AlgebraicTopology.Singular.OpenSimplex}
\lean{AlgebraicTopology.Singular.OpenSimplex}\leanok
Singular \texttt{n}-simplices in an open subset.

\noindent\emph{Lean:} \texttt{abbrev AlgebraicTopology.Singular.OpenSimplex}, \href{https://github.com/Paul-Lez/HodgeConjecture/blob/main/HodgeConjecture/Lemmas/AlgebraicTopology/SingularCochainSheafFlasque.lean#L50}{\texttt{HodgeConjecture/Lemmas/AlgebraicTopology/SingularCochainSheafFlasque.lean:50}}.
\end{definition}

\begin{definition}[AlgebraicTopology.Singular.singularChainOfSimplex]\label{def:AlgebraicTopology.Singular.singularChainOfSimplex}
\lean{AlgebraicTopology.Singular.singularChainOfSimplex}\leanok
\uses{def:AlgebraicTopology.Singular.OpenChains, def:AlgebraicTopology.Singular.OpenSimplex}
The singular chain consisting of one simplex with coefficient one.

\noindent\emph{Lean:} \texttt{def AlgebraicTopology.Singular.singularChainOfSimplex}, \href{https://github.com/Paul-Lez/HodgeConjecture/blob/main/HodgeConjecture/Lemmas/AlgebraicTopology/SingularCochainSheafFlasque.lean#L55}{\texttt{HodgeConjecture/Lemmas/AlgebraicTopology/SingularCochainSheafFlasque.lean:55}}.
\end{definition}

\begin{definition}[AlgebraicTopology.Singular.singularCochainToSimplexFunction]\label{def:AlgebraicTopology.Singular.singularCochainToSimplexFunction}
\lean{AlgebraicTopology.Singular.singularCochainToSimplexFunction}\leanok
\uses{def:AlgebraicTopology.Singular.OpenCochains, def:AlgebraicTopology.Singular.singularChainOfSimplex}
Evaluate a singular cochain on every basis simplex.

\noindent\emph{Lean:} \texttt{def AlgebraicTopology.Singular.singularCochainToSimplexFunction}, \href{https://github.com/Paul-Lez/HodgeConjecture/blob/main/HodgeConjecture/Lemmas/AlgebraicTopology/SingularCochainSheafFlasque.lean#L60}{\texttt{HodgeConjecture/Lemmas/AlgebraicTopology/SingularCochainSheafFlasque.lean:60}}.
\end{definition}

\begin{definition}[AlgebraicTopology.Singular.singularCochainOfSimplexFunction]\label{def:AlgebraicTopology.Singular.singularCochainOfSimplexFunction}
\lean{AlgebraicTopology.Singular.singularCochainOfSimplexFunction}\leanok
\uses{def:AlgebraicTopology.Singular.OpenCochains, def:AlgebraicTopology.Singular.OpenSimplex}
Construct a singular cochain from its values on basis simplices.

\noindent\emph{Lean:} \texttt{def AlgebraicTopology.Singular.singularCochainOfSimplexFunction}, \href{https://github.com/Paul-Lez/HodgeConjecture/blob/main/HodgeConjecture/Lemmas/AlgebraicTopology/SingularCochainSheafFlasque.lean#L67}{\texttt{HodgeConjecture/Lemmas/AlgebraicTopology/SingularCochainSheafFlasque.lean:67}}.
\end{definition}

\begin{definition}[AlgebraicTopology.Singular.singularCochainEquivSimplexFunction]\label{def:AlgebraicTopology.Singular.singularCochainEquivSimplexFunction}
\lean{AlgebraicTopology.Singular.singularCochainEquivSimplexFunction}\leanok
\uses{def:AlgebraicTopology.Singular.singularCochainOfSimplexFunction, def:AlgebraicTopology.Singular.singularCochainToSimplexFunction}
Singular cochains are precisely arbitrary functions on singular simplices.

\noindent\emph{Lean:} \texttt{def AlgebraicTopology.Singular.singularCochainEquivSimplexFunction}, \href{https://github.com/Paul-Lez/HodgeConjecture/blob/main/HodgeConjecture/Lemmas/AlgebraicTopology/SingularCochainSheafFlasque.lean#L139}{\texttt{HodgeConjecture/Lemmas/AlgebraicTopology/SingularCochainSheafFlasque.lean:139}}.
\end{definition}

\begin{definition}[AlgebraicTopology.Singular.openSimplexMap]\label{def:AlgebraicTopology.Singular.openSimplexMap}
\lean{AlgebraicTopology.Singular.openSimplexMap}\leanok
\uses{def:AlgebraicTopology.Singular.OpenSimplex}
An inclusion of open subsets sends a singular simplex to the same simplex in the larger
open subset.

\noindent\emph{Lean:} \texttt{def AlgebraicTopology.Singular.openSimplexMap}, \href{https://github.com/Paul-Lez/HodgeConjecture/blob/main/HodgeConjecture/Lemmas/AlgebraicTopology/SingularCochainSheafFlasque.lean#L150}{\texttt{HodgeConjecture/Lemmas/AlgebraicTopology/SingularCochainSheafFlasque.lean:150}}.
\end{definition}

\section{Naturality of the degree-zero homology augmentation}\label{sec:HodgeConjecture.Lemmas.AlgebraicTopology.HomologyZeroNaturality}
\noindent\emph{File:} \href{https://github.com/Paul-Lez/HodgeConjecture/blob/main/HodgeConjecture/Lemmas/AlgebraicTopology/HomologyZeroNaturality.lean}{\texttt{HodgeConjecture/Lemmas/AlgebraicTopology/HomologyZeroNaturality.lean}}.

The augmentation is natural for actual simplicial and continuous maps. Consequently a
continuous map between path-connected spaces induces an isomorphism on degree-zero homology;
for a nonempty source and path-connected target the induced map is an epimorphism.

The augmentation proof generalizes the additive-group-valued proof in
\texttt{StandardSimplexSingularComparison} to arbitrary coefficient objects.

\begin{definition}[SSet.chainComplexXZeroIsoCyclesZero]\label{def:SSet.chainComplexXZeroIsoCyclesZero}
\lean{SSet.chainComplexXZeroIsoCyclesZero}\leanok
In a nonnegative chain complex the degree-zero chains are the degree-zero cycles.

\noindent\emph{Lean:} \texttt{def SSet.chainComplexXZeroIsoCyclesZero}, \href{https://github.com/Paul-Lez/HodgeConjecture/blob/main/HodgeConjecture/Lemmas/AlgebraicTopology/HomologyZeroNaturality.lean#L32}{\texttt{HodgeConjecture/Lemmas/AlgebraicTopology/HomologyZeroNaturality.lean:32}}.
\end{definition}

\section{Homotopy invariance of relative singular homology}\label{sec:HodgeConjecture.Lemmas.AlgebraicTopology.RelativeHomotopyInvariance}
\noindent\emph{File:} \href{https://github.com/Paul-Lez/HodgeConjecture/blob/main/HodgeConjecture/Lemmas/AlgebraicTopology/RelativeHomotopyInvariance.lean}{\texttt{HodgeConjecture/Lemmas/AlgebraicTopology/RelativeHomotopyInvariance.lean}}.

The relative singular chain complex in \texttt{SingularCohomology} is the cokernel of the
inclusion of singular chains on the subspace. This file shows that a homotopy of maps of
topological pairs descends from the usual singular prism operator to this cokernel. Consequently,
homotopic maps of pairs induce the same map on relative singular homology.

\begin{definition}[TopPair.Homotopy.relativeChainProjectionComponentIsCokernel]\label{def:TopPair.Homotopy.relativeChainProjectionComponentIsCokernel}
\lean{TopPair.Homotopy.relativeChainProjectionComponentIsCokernel}\leanok
\uses{def:AlgebraicTopology.Singular.relativeChainProjection}
The component of the relative-chain projection in each degree is a cokernel.

\noindent\emph{Lean:} \texttt{def TopPair.Homotopy.relativeChainProjectionComponentIsCokernel}, \href{https://github.com/Paul-Lez/HodgeConjecture/blob/main/HodgeConjecture/Lemmas/AlgebraicTopology/RelativeHomotopyInvariance.lean#L135}{\texttt{HodgeConjecture/Lemmas/AlgebraicTopology/RelativeHomotopyInvariance.lean:135}}.
\end{definition}

\begin{definition}[TopPair.Homotopy.relativeChainHomotopyComponent]\label{def:TopPair.Homotopy.relativeChainHomotopyComponent}
\lean{TopPair.Homotopy.relativeChainHomotopyComponent}\leanok
\uses{def:TopPair.Homotopy.relativeChainProjectionComponentIsCokernel}
A component of the ordinary singular prism, descended to relative chains.

\noindent\emph{Lean:} \texttt{def TopPair.Homotopy.relativeChainHomotopyComponent}, \href{https://github.com/Paul-Lez/HodgeConjecture/blob/main/HodgeConjecture/Lemmas/AlgebraicTopology/RelativeHomotopyInvariance.lean#L149}{\texttt{HodgeConjecture/Lemmas/AlgebraicTopology/RelativeHomotopyInvariance.lean:149}}.
\end{definition}

\begin{definition}[TopPair.Homotopy.relativeChainHomotopy]\label{def:TopPair.Homotopy.relativeChainHomotopy}
\lean{TopPair.Homotopy.relativeChainHomotopy}\leanok
\uses{def:TopPair.Homotopy.relativeChainHomotopyComponent}
A homotopy of maps of topological pairs induces a chain homotopy on relative singular
chains.

\noindent\emph{Lean:} \texttt{def TopPair.Homotopy.relativeChainHomotopy}, \href{https://github.com/Paul-Lez/HodgeConjecture/blob/main/HodgeConjecture/Lemmas/AlgebraicTopology/RelativeHomotopyInvariance.lean#L253}{\texttt{HodgeConjecture/Lemmas/AlgebraicTopology/RelativeHomotopyInvariance.lean:253}}.
\end{definition}

\section{Normal-slice reduction for a product support}\label{sec:HodgeConjecture.Lemmas.AlgebraicTopology.NormalSlicePurity}
\noindent\emph{File:} \href{https://github.com/Paul-Lez/HodgeConjecture/blob/main/HodgeConjecture/Lemmas/AlgebraicTopology/NormalSlicePurity.lean}{\texttt{HodgeConjecture/Lemmas/AlgebraicTopology/NormalSlicePurity.lean}}.

Contracting the tangent coordinate gives a chain homotopy equivalence from
\texttt{(E × ℂ\^{}c, E × (ℂ\^{}c \textbackslash{} \{0\}))} to the normal point-complement pair, so the homology and
cohomology identifications come from the projection and the zero section. The distinguished
relative class is the image of the normalized \texttt{standardComplexLocalClass}. This is the
product-model computation for smooth-support purity.

\begin{definition}[AlgebraicTopology.Singular.normalSlicePair]\label{def:AlgebraicTopology.Singular.normalSlicePair}
\lean{AlgebraicTopology.Singular.normalSlicePair}\leanok
The complement of the zero normal slice in a product.

\noindent\emph{Lean:} \texttt{abbrev AlgebraicTopology.Singular.normalSlicePair}, \href{https://github.com/Paul-Lez/HodgeConjecture/blob/main/HodgeConjecture/Lemmas/AlgebraicTopology/NormalSlicePurity.lean#L28}{\texttt{HodgeConjecture/Lemmas/AlgebraicTopology/NormalSlicePurity.lean:28}}.
\end{definition}

\begin{definition}[AlgebraicTopology.Singular.normalSliceProjection]\label{def:AlgebraicTopology.Singular.normalSliceProjection}
\lean{AlgebraicTopology.Singular.normalSliceProjection}\leanok
\uses{def:AlgebraicTopology.Singular.normalSlicePair, def:AlgebraicTopology.Singular.standardComplexPuncturedPair}
Projection to the actual normal point-complement pair.

\noindent\emph{Lean:} \texttt{def AlgebraicTopology.Singular.normalSliceProjection}, \href{https://github.com/Paul-Lez/HodgeConjecture/blob/main/HodgeConjecture/Lemmas/AlgebraicTopology/NormalSlicePurity.lean#L32}{\texttt{HodgeConjecture/Lemmas/AlgebraicTopology/NormalSlicePurity.lean:32}}.
\end{definition}

\begin{definition}[AlgebraicTopology.Singular.normalSliceSection]\label{def:AlgebraicTopology.Singular.normalSliceSection}
\lean{AlgebraicTopology.Singular.normalSliceSection}\leanok
\uses{def:AlgebraicTopology.Singular.normalSlicePair, def:AlgebraicTopology.Singular.standardComplexPuncturedPair}
The zero tangent section preserves the punctured normal coordinate.

\noindent\emph{Lean:} \texttt{def AlgebraicTopology.Singular.normalSliceSection}, \href{https://github.com/Paul-Lez/HodgeConjecture/blob/main/HodgeConjecture/Lemmas/AlgebraicTopology/NormalSlicePurity.lean#L38}{\texttt{HodgeConjecture/Lemmas/AlgebraicTopology/NormalSlicePurity.lean:38}}.
\end{definition}

\begin{definition}[AlgebraicTopology.Singular.normalSliceContraction]\label{def:AlgebraicTopology.Singular.normalSliceContraction}
\lean{AlgebraicTopology.Singular.normalSliceContraction}\leanok
\uses{def:AlgebraicTopology.Singular.normalSliceProjection, def:AlgebraicTopology.Singular.normalSliceSection}
The explicit pair homotopy contracts only tangent coordinates. Its normal coordinate
is unchanged, so the complement condition holds throughout, including at the endpoints.

\noindent\emph{Lean:} \texttt{def AlgebraicTopology.Singular.normalSliceContraction}, \href{https://github.com/Paul-Lez/HodgeConjecture/blob/main/HodgeConjecture/Lemmas/AlgebraicTopology/NormalSlicePurity.lean#L48}{\texttt{HodgeConjecture/Lemmas/AlgebraicTopology/NormalSlicePurity.lean:48}}.
\end{definition}

\begin{definition}[AlgebraicTopology.Singular.normalSliceRelativeChainHomotopyEquiv]\label{def:AlgebraicTopology.Singular.normalSliceRelativeChainHomotopyEquiv}
\lean{AlgebraicTopology.Singular.normalSliceRelativeChainHomotopyEquiv}\leanok
\uses{def:AlgebraicTopology.Singular.normalSliceContraction, def:TopPair.Homotopy.relativeChainHomotopy}
Normal projection and zero section are inverse up to the actual relative prism homotopy.

\noindent\emph{Lean:} \texttt{def AlgebraicTopology.Singular.normalSliceRelativeChainHomotopyEquiv}, \href{https://github.com/Paul-Lez/HodgeConjecture/blob/main/HodgeConjecture/Lemmas/AlgebraicTopology/NormalSlicePurity.lean#L71}{\texttt{HodgeConjecture/Lemmas/AlgebraicTopology/NormalSlicePurity.lean:71}}.
\end{definition}

\begin{definition}[AlgebraicTopology.Singular.normalSliceRelativeHomologyIso]\label{def:AlgebraicTopology.Singular.normalSliceRelativeHomologyIso}
\lean{AlgebraicTopology.Singular.normalSliceRelativeHomologyIso}\leanok
\uses{def:AlgebraicTopology.Singular.RelativeHomology, def:AlgebraicTopology.Singular.normalSliceRelativeChainHomotopyEquiv}
The induced normal-slice isomorphism in every relative homology degree.

\noindent\emph{Lean:} \texttt{def AlgebraicTopology.Singular.normalSliceRelativeHomologyIso}, \href{https://github.com/Paul-Lez/HodgeConjecture/blob/main/HodgeConjecture/Lemmas/AlgebraicTopology/NormalSlicePurity.lean#L84}{\texttt{HodgeConjecture/Lemmas/AlgebraicTopology/NormalSlicePurity.lean:84}}.
\end{definition}

\section{Invariance of the complex local class under complex-linear coordinate changes}\label{sec:HodgeConjecture.Lemmas.AlgebraicTopology.ComplexLinearLocalClassInvariance}
\noindent\emph{File:} \href{https://github.com/Paul-Lez/HodgeConjecture/blob/main/HodgeConjecture/Lemmas/AlgebraicTopology/ComplexLinearLocalClassInvariance.lean}{\texttt{HodgeConjecture/Lemmas/AlgebraicTopology/ComplexLinearLocalClassInvariance.lean}}.

An invertible complex matrix is joined to the identity through invertible complex matrices.
The proof uses Mathlib's generation of invertible matrices by diagonal matrices and
transvections.  Diagonal entries are joined to one by \texttt{exp (t * log z)}, while a transvection is
joined to one by scaling its off-diagonal entry by \texttt{t}.

The resulting isotopy gives a homotopy of punctured pairs. Hence every complex-linear
automorphism acts trivially on the explicitly normalized standard complex local homology class.
This supplies the linear local-degree step needed to compare holomorphic coordinate charts; the
remaining nonlinear step is to replace a chart transition germ by its derivative.

\begin{definition}[Matrix.ComplexIsotopyToOne]\label{def:Matrix.ComplexIsotopyToOne}
\lean{Matrix.ComplexIsotopyToOne}\leanok
An isotopy from the identity matrix to \texttt{A}, all of whose members are invertible.

\noindent\emph{Lean:} \texttt{structure Matrix.ComplexIsotopyToOne}, \href{https://github.com/Paul-Lez/HodgeConjecture/blob/main/HodgeConjecture/Lemmas/AlgebraicTopology/ComplexLinearLocalClassInvariance.lean#L47}{\texttt{HodgeConjecture/Lemmas/AlgebraicTopology/ComplexLinearLocalClassInvariance.lean:47}}.
\end{definition}

\begin{definition}[Matrix.ComplexIsotopyToOne.mul]\label{def:Matrix.ComplexIsotopyToOne.mul}
\lean{Matrix.ComplexIsotopyToOne.mul}\leanok
\uses{def:Matrix.ComplexIsotopyToOne}
Pointwise multiplication concatenates isotopies algebraically.  This is the construction
used for products in the diagonal--transvection induction.

\noindent\emph{Lean:} \texttt{def Matrix.ComplexIsotopyToOne.mul}, \href{https://github.com/Paul-Lez/HodgeConjecture/blob/main/HodgeConjecture/Lemmas/AlgebraicTopology/ComplexLinearLocalClassInvariance.lean#L60}{\texttt{HodgeConjecture/Lemmas/AlgebraicTopology/ComplexLinearLocalClassInvariance.lean:60}}.
\end{definition}

\begin{definition}[Matrix.ComplexIsotopyToOne.diagonal]\label{def:Matrix.ComplexIsotopyToOne.diagonal}
\lean{Matrix.ComplexIsotopyToOne.diagonal}\leanok
\uses{def:Matrix.ComplexIsotopyToOne}
An invertible diagonal complex matrix is isotopic to the identity.

\noindent\emph{Lean:} \texttt{def Matrix.ComplexIsotopyToOne.diagonal}, \href{https://github.com/Paul-Lez/HodgeConjecture/blob/main/HodgeConjecture/Lemmas/AlgebraicTopology/ComplexLinearLocalClassInvariance.lean#L81}{\texttt{HodgeConjecture/Lemmas/AlgebraicTopology/ComplexLinearLocalClassInvariance.lean:81}}.
\end{definition}

\begin{definition}[Matrix.ComplexIsotopyToOne.transvection]\label{def:Matrix.ComplexIsotopyToOne.transvection}
\lean{Matrix.ComplexIsotopyToOne.transvection}\leanok
\uses{def:Matrix.ComplexIsotopyToOne}
A complex transvection is isotopic to the identity by scaling its off-diagonal entry.

\noindent\emph{Lean:} \texttt{def Matrix.ComplexIsotopyToOne.transvection}, \href{https://github.com/Paul-Lez/HodgeConjecture/blob/main/HodgeConjecture/Lemmas/AlgebraicTopology/ComplexLinearLocalClassInvariance.lean#L118}{\texttt{HodgeConjecture/Lemmas/AlgebraicTopology/ComplexLinearLocalClassInvariance.lean:118}}.
\end{definition}

\begin{definition}[AlgebraicTopology.Singular.complexMatrixMap]\label{def:AlgebraicTopology.Singular.complexMatrixMap}
\lean{AlgebraicTopology.Singular.complexMatrixMap}\leanok
The continuous self-map of complex affine space defined by a matrix.

\noindent\emph{Lean:} \texttt{def AlgebraicTopology.Singular.complexMatrixMap}, \href{https://github.com/Paul-Lez/HodgeConjecture/blob/main/HodgeConjecture/Lemmas/AlgebraicTopology/ComplexLinearLocalClassInvariance.lean#L189}{\texttt{HodgeConjecture/Lemmas/AlgebraicTopology/ComplexLinearLocalClassInvariance.lean:189}}.
\end{definition}

\begin{definition}[AlgebraicTopology.Singular.complexMatrixPuncturedMap]\label{def:AlgebraicTopology.Singular.complexMatrixPuncturedMap}
\lean{AlgebraicTopology.Singular.complexMatrixPuncturedMap}\leanok
The restriction of an invertible complex matrix to the complement of the origin.

\noindent\emph{Lean:} \texttt{def AlgebraicTopology.Singular.complexMatrixPuncturedMap}, \href{https://github.com/Paul-Lez/HodgeConjecture/blob/main/HodgeConjecture/Lemmas/AlgebraicTopology/ComplexLinearLocalClassInvariance.lean#L200}{\texttt{HodgeConjecture/Lemmas/AlgebraicTopology/ComplexLinearLocalClassInvariance.lean:200}}.
\end{definition}

\begin{definition}[AlgebraicTopology.Singular.complexMatrixPuncturedPairMap]\label{def:AlgebraicTopology.Singular.complexMatrixPuncturedPairMap}
\lean{AlgebraicTopology.Singular.complexMatrixPuncturedPairMap}\leanok
\uses{def:AlgebraicTopology.Singular.complexMatrixMap, def:AlgebraicTopology.Singular.complexMatrixPuncturedMap, def:AlgebraicTopology.Singular.standardComplexPuncturedPair}
An invertible complex matrix acts on complex affine space and its punctured subspace.

\noindent\emph{Lean:} \texttt{def AlgebraicTopology.Singular.complexMatrixPuncturedPairMap}, \href{https://github.com/Paul-Lez/HodgeConjecture/blob/main/HodgeConjecture/Lemmas/AlgebraicTopology/ComplexLinearLocalClassInvariance.lean#L219}{\texttt{HodgeConjecture/Lemmas/AlgebraicTopology/ComplexLinearLocalClassInvariance.lean:219}}.
\end{definition}

\begin{definition}[AlgebraicTopology.Singular.complexMatrixPuncturedPairHomotopy]\label{def:AlgebraicTopology.Singular.complexMatrixPuncturedPairHomotopy}
\lean{AlgebraicTopology.Singular.complexMatrixPuncturedPairHomotopy}\leanok
\uses{def:AlgebraicTopology.Singular.complexMatrixPuncturedPairMap, def:Matrix.ComplexIsotopyToOne}
An isotopy of invertible matrices induces a homotopy of their maps of punctured pairs.

\noindent\emph{Lean:} \texttt{def AlgebraicTopology.Singular.complexMatrixPuncturedPairHomotopy}, \href{https://github.com/Paul-Lez/HodgeConjecture/blob/main/HodgeConjecture/Lemmas/AlgebraicTopology/ComplexLinearLocalClassInvariance.lean#L231}{\texttt{HodgeConjecture/Lemmas/AlgebraicTopology/ComplexLinearLocalClassInvariance.lean:231}}.
\end{definition}

\section{Invariance of the complex local class under differentiable coordinate changes}\label{sec:HodgeConjecture.Lemmas.AlgebraicTopology.ComplexDifferentiableLocalClassInvariance}
\noindent\emph{File:} \href{https://github.com/Paul-Lez/HodgeConjecture/blob/main/HodgeConjecture/Lemmas/AlgebraicTopology/ComplexDifferentiableLocalClassInvariance.lean}{\texttt{HodgeConjecture/Lemmas/AlgebraicTopology/ComplexDifferentiableLocalClassInvariance.lean}}.

This file supplies the nonlinear local-degree step in the comparison of complex charts.  If a
continuous map fixes the origin, is complex differentiable there, and has invertible derivative,
then on a sufficiently small neighborhood its straight-line homotopy to the derivative avoids the
origin away from the origin.  Point excision lets us use this local homotopy to compare the induced
maps on the ambient local homology group.  Combining this with the complex-linear calculation shows
that the map preserves the standard complex local class.

The only global condition is the one required to even define a self-map of the point-complement
pair: the map has no zero away from the origin.  In chart applications this follows from
injectivity of the coordinate change.  No local homotopy or local-degree datum is assumed.

\begin{definition}[AlgebraicTopology.Singular.complexPuncturedPairMapOf]\label{def:AlgebraicTopology.Singular.complexPuncturedPairMapOf}
\lean{AlgebraicTopology.Singular.complexPuncturedPairMapOf}\leanok
\uses{def:AlgebraicTopology.Singular.standardComplexPuncturedPair}
A continuous map whose only zero is the origin defines a self-map of the punctured complex
affine-space pair.

\noindent\emph{Lean:} \texttt{def AlgebraicTopology.Singular.complexPuncturedPairMapOf}, \href{https://github.com/Paul-Lez/HodgeConjecture/blob/main/HodgeConjecture/Lemmas/AlgebraicTopology/ComplexDifferentiableLocalClassInvariance.lean#L45}{\texttt{HodgeConjecture/Lemmas/AlgebraicTopology/ComplexDifferentiableLocalClassInvariance.lean:45}}.
\end{definition}

\begin{definition}[AlgebraicTopology.Singular.complexNeighborhoodPuncturedPairMapOf]\label{def:AlgebraicTopology.Singular.complexNeighborhoodPuncturedPairMapOf}
\lean{AlgebraicTopology.Singular.complexNeighborhoodPuncturedPairMapOf}\leanok
\uses{def:AlgebraicTopology.Singular.neighborhoodPointComplementPair, def:AlgebraicTopology.Singular.standardComplexPuncturedPair}
A map defined continuously on a neighborhood and avoiding the origin off its distinguished
point gives a map from that neighborhood pair to the standard punctured pair.

\noindent\emph{Lean:} \texttt{def AlgebraicTopology.Singular.complexNeighborhoodPuncturedPairMapOf}, \href{https://github.com/Paul-Lez/HodgeConjecture/blob/main/HodgeConjecture/Lemmas/AlgebraicTopology/ComplexDifferentiableLocalClassInvariance.lean#L155}{\texttt{HodgeConjecture/Lemmas/AlgebraicTopology/ComplexDifferentiableLocalClassInvariance.lean:155}}.
\end{definition}

\begin{definition}[AlgebraicTopology.Singular.complexStraightLineLocalPairHomotopy]\label{def:AlgebraicTopology.Singular.complexStraightLineLocalPairHomotopy}
\lean{AlgebraicTopology.Singular.complexStraightLineLocalPairHomotopy}\leanok
\uses{def:AlgebraicTopology.Singular.complexMatrixPuncturedPairMap, def:AlgebraicTopology.Singular.complexNeighborhoodPuncturedPairMapOf, def:AlgebraicTopology.Singular.neighborhoodPointComplementPairMap}
The germ-local straight-line homotopy, requiring continuity and nonvanishing only on its
source neighborhood.

\noindent\emph{Lean:} \texttt{def AlgebraicTopology.Singular.complexStraightLineLocalPairHomotopy}, \href{https://github.com/Paul-Lez/HodgeConjecture/blob/main/HodgeConjecture/Lemmas/AlgebraicTopology/ComplexDifferentiableLocalClassInvariance.lean#L169}{\texttt{HodgeConjecture/Lemmas/AlgebraicTopology/ComplexDifferentiableLocalClassInvariance.lean:169}}.
\end{definition}

\begin{definition}[AlgebraicTopology.Singular.complexStraightLineNeighborhoodPairHomotopy]\label{def:AlgebraicTopology.Singular.complexStraightLineNeighborhoodPairHomotopy}
\lean{AlgebraicTopology.Singular.complexStraightLineNeighborhoodPairHomotopy}\leanok
\uses{def:AlgebraicTopology.Singular.complexMatrixPuncturedPairMap, def:AlgebraicTopology.Singular.complexPuncturedPairMapOf, def:AlgebraicTopology.Singular.neighborhoodPointComplementPairMap}
A puncture-preserving straight line from an invertible complex-linear map to a nonlinear
map gives a homotopy of pair maps after restricting the source to the neighborhood on which the
straight line avoids the origin.

\noindent\emph{Lean:} \texttt{def AlgebraicTopology.Singular.complexStraightLineNeighborhoodPairHomotopy}, \href{https://github.com/Paul-Lez/HodgeConjecture/blob/main/HodgeConjecture/Lemmas/AlgebraicTopology/ComplexDifferentiableLocalClassInvariance.lean#L290}{\texttt{HodgeConjecture/Lemmas/AlgebraicTopology/ComplexDifferentiableLocalClassInvariance.lean:290}}.
\end{definition}

\section{Chart-local class invariance under a differentiable transition germ}\label{sec:HodgeConjecture.Lemmas.AlgebraicTopology.ChartLocalFundamentalClassDifferentiableInvariance}
\noindent\emph{File:} \href{https://github.com/Paul-Lez/HodgeConjecture/blob/main/HodgeConjecture/Lemmas/AlgebraicTopology/ChartLocalFundamentalClassDifferentiableInvariance.lean}{\texttt{HodgeConjecture/Lemmas/AlgebraicTopology/ChartLocalFundamentalClassDifferentiableInvariance.lean}}.

This file turns the local nonlinear degree calculation into a chart-comparison theorem.  For two
complex charts through the same point, their compressed inverse-chart embeddings determine an
actual open partial homeomorphism of \texttt{ℂᵈ}.  If its derivative at the model origin is injective
and complex linear, the two explicitly normalized chart-local homology classes agree.

The proof constructs the overlap neighborhood, uses point excision to lift the standard class,
and proves the factorization through the transition as an equality of maps of topological pairs.
Thus no chart-compatibility or local-degree statement is assumed.

\begin{definition}[AlgebraicTopology.Singular.complexMatrixOfContinuousLinearMap]\label{def:AlgebraicTopology.Singular.complexMatrixOfContinuousLinearMap}
\lean{AlgebraicTopology.Singular.complexMatrixOfContinuousLinearMap}\leanok
The matrix of a complex continuous-linear endomorphism in the standard basis of a pi space.

\noindent\emph{Lean:} \texttt{def AlgebraicTopology.Singular.complexMatrixOfContinuousLinearMap}, \href{https://github.com/Paul-Lez/HodgeConjecture/blob/main/HodgeConjecture/Lemmas/AlgebraicTopology/ChartLocalFundamentalClassDifferentiableInvariance.lean#L168}{\texttt{HodgeConjecture/Lemmas/AlgebraicTopology/ChartLocalFundamentalClassDifferentiableInvariance.lean:168}}.
\end{definition}

\section{The normal derivative of a support-preserving holomorphic transition}\label{sec:HodgeConjecture.Lemmas.AlgebraicTopology.HolomorphicNormalTransition}
\noindent\emph{File:} \href{https://github.com/Paul-Lez/HodgeConjecture/blob/main/HodgeConjecture/Lemmas/AlgebraicTopology/HolomorphicNormalTransition.lean}{\texttt{HodgeConjecture/Lemmas/AlgebraicTopology/HolomorphicNormalTransition.lean}}.

Preservation of the zero-normal plane forces the tangent-to-normal derivative block to
vanish. Differentiating the actual local inverse identity then constructs inverse normal
blocks. No invertibility or orientation-preservation theorem is supplied for the normal
map; both are consequences of the given holomorphic support-preserving coordinate transition.

\begin{definition}[AlgebraicTopology.Singular.normalBlock]\label{def:AlgebraicTopology.Singular.normalBlock}
\lean{AlgebraicTopology.Singular.normalBlock}\leanok
The normal-to-normal block of an actual ambient derivative.

\noindent\emph{Lean:} \texttt{def AlgebraicTopology.Singular.normalBlock}, \href{https://github.com/Paul-Lez/HodgeConjecture/blob/main/HodgeConjecture/Lemmas/AlgebraicTopology/HolomorphicNormalTransition.lean#L31}{\texttt{HodgeConjecture/Lemmas/AlgebraicTopology/HolomorphicNormalTransition.lean:31}}.
\end{definition}

\begin{definition}[AlgebraicTopology.Singular.normalTransitionDerivativeEquiv]\label{def:AlgebraicTopology.Singular.normalTransitionDerivativeEquiv}
\lean{AlgebraicTopology.Singular.normalTransitionDerivativeEquiv}\leanok
\uses{def:AlgebraicTopology.Singular.normalBlock}
The normal derivative is complex-linearly invertible, with inverse the actual normal
block of the derivative of the inverse chart.

\noindent\emph{Lean:} \texttt{def AlgebraicTopology.Singular.normalTransitionDerivativeEquiv}, \href{https://github.com/Paul-Lez/HodgeConjecture/blob/main/HodgeConjecture/Lemmas/AlgebraicTopology/HolomorphicNormalTransition.lean#L115}{\texttt{HodgeConjecture/Lemmas/AlgebraicTopology/HolomorphicNormalTransition.lean:115}}.
\end{definition}

\begin{definition}[AlgebraicTopology.Singular.normalTransitionMap]\label{def:AlgebraicTopology.Singular.normalTransitionMap}
\lean{AlgebraicTopology.Singular.normalTransitionMap}\leanok
The actual transverse normal map of a coordinate transition.

\noindent\emph{Lean:} \texttt{def AlgebraicTopology.Singular.normalTransitionMap}, \href{https://github.com/Paul-Lez/HodgeConjecture/blob/main/HodgeConjecture/Lemmas/AlgebraicTopology/HolomorphicNormalTransition.lean#L154}{\texttt{HodgeConjecture/Lemmas/AlgebraicTopology/HolomorphicNormalTransition.lean:154}}.
\end{definition}

\begin{definition}[AlgebraicTopology.Singular.normalTransitionDomain]\label{def:AlgebraicTopology.Singular.normalTransitionDomain}
\lean{AlgebraicTopology.Singular.normalTransitionDomain}\leanok
Only normal parameters whose full points are in the transition domain are used.

\noindent\emph{Lean:} \texttt{def AlgebraicTopology.Singular.normalTransitionDomain}, \href{https://github.com/Paul-Lez/HodgeConjecture/blob/main/HodgeConjecture/Lemmas/AlgebraicTopology/HolomorphicNormalTransition.lean#L157}{\texttt{HodgeConjecture/Lemmas/AlgebraicTopology/HolomorphicNormalTransition.lean:157}}.
\end{definition}

\section{Moving the center of a complex-coordinate embedding}\label{sec:HodgeConjecture.Lemmas.AlgebraicTopology.CenteredComplexEmbeddingOrientation}
\noindent\emph{File:} \href{https://github.com/Paul-Lez/HodgeConjecture/blob/main/HodgeConjecture/Lemmas/AlgebraicTopology/CenteredComplexEmbeddingOrientation.lean}{\texttt{HodgeConjecture/Lemmas/AlgebraicTopology/CenteredComplexEmbeddingOrientation.lean}}.

For a continuous injective map \texttt{F : ℂ\^{}d → ℂ\^{}d}, the maps
\texttt{w ↦ F(w + t • v) - F(t • v)} form a homotopy of punctured pairs. Injectivity, not
differentiability at varying centers, prevents this homotopy from hitting zero.

In particular the standard radial chart compression preserves the exact complex local class
at every center. Complex differentiability is used only at the model origin, where the
derivative is a positive scalar times the identity. No assertion that radial compression is
holomorphic away from the origin is needed or made.

\begin{definition}[AlgebraicTopology.Singular.centeredComplexEmbeddingPair]\label{def:AlgebraicTopology.Singular.centeredComplexEmbeddingPair}
\lean{AlgebraicTopology.Singular.centeredComplexEmbeddingPair}\leanok
\uses{def:AlgebraicTopology.Singular.complexPuncturedPairMapOf}
Recenter a continuous injection at a chosen source point.

\noindent\emph{Lean:} \texttt{def AlgebraicTopology.Singular.centeredComplexEmbeddingPair}, \href{https://github.com/Paul-Lez/HodgeConjecture/blob/main/HodgeConjecture/Lemmas/AlgebraicTopology/CenteredComplexEmbeddingOrientation.lean#L42}{\texttt{HodgeConjecture/Lemmas/AlgebraicTopology/CenteredComplexEmbeddingOrientation.lean:42}}.
\end{definition}

\begin{definition}[AlgebraicTopology.Singular.centeredComplexEmbeddingPairHomotopy]\label{def:AlgebraicTopology.Singular.centeredComplexEmbeddingPairHomotopy}
\lean{AlgebraicTopology.Singular.centeredComplexEmbeddingPairHomotopy}\leanok
\uses{def:AlgebraicTopology.Singular.centeredComplexEmbeddingPair}
An explicit punctured-pair homotopy moving the center from zero to \texttt{v}.

\noindent\emph{Lean:} \texttt{def AlgebraicTopology.Singular.centeredComplexEmbeddingPairHomotopy}, \href{https://github.com/Paul-Lez/HodgeConjecture/blob/main/HodgeConjecture/Lemmas/AlgebraicTopology/CenteredComplexEmbeddingOrientation.lean#L51}{\texttt{HodgeConjecture/Lemmas/AlgebraicTopology/CenteredComplexEmbeddingOrientation.lean:51}}.
\end{definition}

\section{Normal fibers and ambient overlap compatibility}\label{sec:HodgeConjecture.Lemmas.AlgebraicTopology.NormalProjectionOverlap}
\noindent\emph{File:} \href{https://github.com/Paul-Lez/HodgeConjecture/blob/main/HodgeConjecture/Lemmas/AlgebraicTopology/NormalProjectionOverlap.lean}{\texttt{HodgeConjecture/Lemmas/AlgebraicTopology/NormalProjectionOverlap.lean}}.

Actual normal fibers factor both normal projections through maps of point-complement
pairs. Complex orientation invariance of the normal transition therefore gives agreement
of the actual ambient supported coclasses, not just an abstract transverse comparison.

\begin{definition}[AlgebraicTopology.Singular.chartNormalFiberPair]\label{def:AlgebraicTopology.Singular.chartNormalFiberPair}
\lean{AlgebraicTopology.Singular.chartNormalFiberPair}\leanok
\uses{def:AlgebraicTopology.Singular.neighborhoodPointComplementPair, def:AlgebraicTopology.Singular.neighborhoodSupportComplementPair}
The uncompressed normal fiber is a genuine map from a small normal point-complement
pair into the chosen ambient support-complement pair.

\noindent\emph{Lean:} \texttt{def AlgebraicTopology.Singular.chartNormalFiberPair}, \href{https://github.com/Paul-Lez/HodgeConjecture/blob/main/HodgeConjecture/Lemmas/AlgebraicTopology/NormalProjectionOverlap.lean#L71}{\texttt{HodgeConjecture/Lemmas/AlgebraicTopology/NormalProjectionOverlap.lean:71}}.
\end{definition}

\section{Sections of the homology sheaf from actual relative homology}\label{sec:HodgeConjecture.Lemmas.AlgebraicTopology.HomologySheafSection}
\noindent\emph{File:} \href{https://github.com/Paul-Lez/HodgeConjecture/blob/main/HodgeConjecture/Lemmas/AlgebraicTopology/HomologySheafSection.lean}{\texttt{HodgeConjecture/Lemmas/AlgebraicTopology/HomologySheafSection.lean}}.

Exact sheafification sends presheaf homology to the actual homology sheaf. The unit
therefore sends a relative singular homology class on an open support to a section of
the homology sheaf. The germ comparison records the actual restriction to local relative
homology, which is essential for preserving normalized fundamental classes.

\begin{definition}[AlgebraicTopology.Singular.additivePresheafStalkFunctor]\label{def:AlgebraicTopology.Singular.additivePresheafStalkFunctor}
\lean{AlgebraicTopology.Singular.additivePresheafStalkFunctor}\leanok
The stalk functor with its domain displayed as a functor category, to align the
canonical additive structures used in functor-category homology comparisons.

\noindent\emph{Lean:} \texttt{abbrev AlgebraicTopology.Singular.additivePresheafStalkFunctor}, \href{https://github.com/Paul-Lez/HodgeConjecture/blob/main/HodgeConjecture/Lemmas/AlgebraicTopology/HomologySheafSection.lean#L86}{\texttt{HodgeConjecture/Lemmas/AlgebraicTopology/HomologySheafSection.lean:86}}.
\end{definition}

\begin{definition}[AlgebraicTopology.Singular.additiveSheafStalkFunctor]\label{def:AlgebraicTopology.Singular.additiveSheafStalkFunctor}
\lean{AlgebraicTopology.Singular.additiveSheafStalkFunctor}\leanok
The exact sheaf stalk functor, with the categorical sheaf domain displayed.

\noindent\emph{Lean:} \texttt{abbrev AlgebraicTopology.Singular.additiveSheafStalkFunctor}, \href{https://github.com/Paul-Lez/HodgeConjecture/blob/main/HodgeConjecture/Lemmas/AlgebraicTopology/HomologySheafSection.lean#L96}{\texttt{HodgeConjecture/Lemmas/AlgebraicTopology/HomologySheafSection.lean:96}}.
\end{definition}

\begin{definition}[AlgebraicTopology.Singular.additivePresheafStalkSheafificationIso]\label{def:AlgebraicTopology.Singular.additivePresheafStalkSheafificationIso}
\lean{AlgebraicTopology.Singular.additivePresheafStalkSheafificationIso}\leanok
\uses{def:AlgebraicTopology.Singular.additivePresheafStalkFunctor, def:AlgebraicTopology.Singular.additiveSheafStalkFunctor}
Sheafification leaves stalks unchanged, naturally in the additive presheaf.

\noindent\emph{Lean:} \texttt{def AlgebraicTopology.Singular.additivePresheafStalkSheafificationIso}, \href{https://github.com/Paul-Lez/HodgeConjecture/blob/main/HodgeConjecture/Lemmas/AlgebraicTopology/HomologySheafSection.lean#L121}{\texttt{HodgeConjecture/Lemmas/AlgebraicTopology/HomologySheafSection.lean:121}}.
\end{definition}

\section{Local sections of flasque resolution comparisons}\label{sec:HodgeConjecture.Lemmas.AlgebraicTopology.FlasqueSheafSupportComparison}
\noindent\emph{File:} \href{https://github.com/Paul-Lez/HodgeConjecture/blob/main/HodgeConjecture/Lemmas/AlgebraicTopology/FlasqueSheafSupportComparison.lean}{\texttt{HodgeConjecture/Lemmas/AlgebraicTopology/FlasqueSheafSupportComparison.lean}}.

A quasi-isomorphism of bounded-below termwise-flasque complexes remains a
quasi-isomorphism on every open set. The proof restricts to that open set and
uses the proved flasque global-sections theorem. Consequently its direct image
under any continuous map is a sheaf quasi-isomorphism. This is a theorem about
flasque models, not an assertion that arbitrary direct image is exact.

The localization sequence is short exact on every open set for flasque
coefficients, so the actual supported-sections complex computes the fiber of
restriction even for the singular-cochain flasque model.

\begin{definition}[TopCat.Sheaf.openRestrictionGlobalSectionsIso]\label{def:TopCat.Sheaf.openRestrictionGlobalSectionsIso}
\lean{TopCat.Sheaf.openRestrictionGlobalSectionsIso}\leanok
\uses{def:TopCat.Sheaf.IsFlasque.BoundedBelowComplex.globalSectionsFunctor, def:TopCat.Sheaf.supportEvaluation}
Global sections after open restriction are actual sections on the ambient
open set, via the canonical equality of the image of the top open with \texttt{U}.

\noindent\emph{Lean:} \texttt{def TopCat.Sheaf.openRestrictionGlobalSectionsIso}, \href{https://github.com/Paul-Lez/HodgeConjecture/blob/main/HodgeConjecture/Lemmas/AlgebraicTopology/FlasqueSheafSupportComparison.lean#L35}{\texttt{HodgeConjecture/Lemmas/AlgebraicTopology/FlasqueSheafSupportComparison.lean:35}}.
\end{definition}

\begin{definition}[TopCat.Sheaf.supportRestrictionComplexShortComplexMap]\label{def:TopCat.Sheaf.supportRestrictionComplexShortComplexMap}
\lean{TopCat.Sheaf.supportRestrictionComplexShortComplexMap}\leanok
\uses{def:TopCat.Sheaf.supportRestrictionComplexShortComplex}
A coefficient-complex map induces the actual map of localization sequences.

\noindent\emph{Lean:} \texttt{def TopCat.Sheaf.supportRestrictionComplexShortComplexMap}, \href{https://github.com/Paul-Lez/HodgeConjecture/blob/main/HodgeConjecture/Lemmas/AlgebraicTopology/FlasqueSheafSupportComparison.lean#L131}{\texttt{HodgeConjecture/Lemmas/AlgebraicTopology/FlasqueSheafSupportComparison.lean:131}}.
\end{definition}

\section{Naturality of derived sections with closed support}\label{sec:HodgeConjecture.Lemmas.AlgebraicTopology.DerivedSheafSupportNaturality}
\noindent\emph{File:} \href{https://github.com/Paul-Lez/HodgeConjecture/blob/main/HodgeConjecture/Lemmas/AlgebraicTopology/DerivedSheafSupportNaturality.lean}{\texttt{HodgeConjecture/Lemmas/AlgebraicTopology/DerivedSheafSupportNaturality.lean}}.

Support enlargement comes from the actual kernel inclusion, and is then derived by
the universal property. In particular, forgetting support is the map to whole-space
support; no cohomology-level comparison morphism is supplied.

\begin{definition}[TopCat.Sheaf.openRestrictionImage]\label{def:TopCat.Sheaf.openRestrictionImage}
\lean{TopCat.Sheaf.openRestrictionImage}\leanok
The ambient open used when restricting to \texttt{U} and pushing forward again.

\noindent\emph{Lean:} \texttt{abbrev TopCat.Sheaf.openRestrictionImage}, \href{https://github.com/Paul-Lez/HodgeConjecture/blob/main/HodgeConjecture/Lemmas/AlgebraicTopology/DerivedSheafSupportNaturality.lean#L34}{\texttt{HodgeConjecture/Lemmas/AlgebraicTopology/DerivedSheafSupportNaturality.lean:34}}.
\end{definition}

\begin{definition}[TopCat.Sheaf.openRestrictionPushforwardMap]\label{def:TopCat.Sheaf.openRestrictionPushforwardMap}
\lean{TopCat.Sheaf.openRestrictionPushforwardMap}\leanok
\uses{def:TopCat.Sheaf.openRestrictionImage, def:TopCat.Sheaf.openRestrictionPushforward}
Further restriction from \texttt{U} to \texttt{V ⊆ U}, on the actual open pushforward functors.

\noindent\emph{Lean:} \texttt{def TopCat.Sheaf.openRestrictionPushforwardMap}, \href{https://github.com/Paul-Lez/HodgeConjecture/blob/main/HodgeConjecture/Lemmas/AlgebraicTopology/DerivedSheafSupportNaturality.lean#L43}{\texttt{HodgeConjecture/Lemmas/AlgebraicTopology/DerivedSheafSupportNaturality.lean:43}}.
\end{definition}

\begin{definition}[TopCat.Sheaf.sheafSectionsSupportedOutsideMap]\label{def:TopCat.Sheaf.sheafSectionsSupportedOutsideMap}
\lean{TopCat.Sheaf.sheafSectionsSupportedOutsideMap}\leanok
\uses{def:TopCat.Sheaf.liftSheafSectionsSupportedOutside, def:TopCat.Sheaf.openRestrictionPushforwardMap, def:TopCat.Sheaf.sheafSectionsSupportedOutsideInclusion}
The canonical inclusion of sections supported outside \texttt{U} into those supported
outside the smaller open \texttt{V}.

\noindent\emph{Lean:} \texttt{def TopCat.Sheaf.sheafSectionsSupportedOutsideMap}, \href{https://github.com/Paul-Lez/HodgeConjecture/blob/main/HodgeConjecture/Lemmas/AlgebraicTopology/DerivedSheafSupportNaturality.lean#L69}{\texttt{HodgeConjecture/Lemmas/AlgebraicTopology/DerivedSheafSupportNaturality.lean:69}}.
\end{definition}

\section{Sheafification commutes with actual open restriction}\label{sec:HodgeConjecture.Lemmas.AlgebraicTopology.OpenSheafification}
\noindent\emph{File:} \href{https://github.com/Paul-Lez/HodgeConjecture/blob/main/HodgeConjecture/Lemmas/AlgebraicTopology/OpenSheafification.lean}{\texttt{HodgeConjecture/Lemmas/AlgebraicTopology/OpenSheafification.lean}}.

Open-image functors lift covering sieves. The resulting sheafification
comparison is the actual sheafification lift of the restricted unit.

\begin{definition}[TopCat.Sheaf.openRestrictionSheafificationIso]\label{def:TopCat.Sheaf.openRestrictionSheafificationIso}
\lean{TopCat.Sheaf.openRestrictionSheafificationIso}\leanok
Sheafification followed by actual open restriction is canonically
sheafification of the restricted presheaf.

\noindent\emph{Lean:} \texttt{def TopCat.Sheaf.openRestrictionSheafificationIso}, \href{https://github.com/Paul-Lez/HodgeConjecture/blob/main/HodgeConjecture/Lemmas/AlgebraicTopology/OpenSheafification.lean#L48}{\texttt{HodgeConjecture/Lemmas/AlgebraicTopology/OpenSheafification.lean:48}}.
\end{definition}

\section{Singular cochains and actual restriction to open subspaces}\label{sec:HodgeConjecture.Lemmas.AlgebraicTopology.SingularCochainOpenRestriction}
\noindent\emph{File:} \href{https://github.com/Paul-Lez/HodgeConjecture/blob/main/HodgeConjecture/Lemmas/AlgebraicTopology/SingularCochainOpenRestriction.lean}{\texttt{HodgeConjecture/Lemmas/AlgebraicTopology/SingularCochainOpenRestriction.lean}}.

\begin{definition}[AlgebraicTopology.Singular.openSubspaceImageIso]\label{def:AlgebraicTopology.Singular.openSubspaceImageIso}
\lean{AlgebraicTopology.Singular.openSubspaceImageIso}\leanok
An open in an open subspace is identified with its actual ambient image.

\noindent\emph{Lean:} \texttt{def AlgebraicTopology.Singular.openSubspaceImageIso}, \href{https://github.com/Paul-Lez/HodgeConjecture/blob/main/HodgeConjecture/Lemmas/AlgebraicTopology/SingularCochainOpenRestriction.lean#L21}{\texttt{HodgeConjecture/Lemmas/AlgebraicTopology/SingularCochainOpenRestriction.lean:21}}.
\end{definition}

\begin{definition}[AlgebraicTopology.Singular.openSubspaceImageChainIso]\label{def:AlgebraicTopology.Singular.openSubspaceImageChainIso}
\lean{AlgebraicTopology.Singular.openSubspaceImageChainIso}\leanok
\uses{def:AlgebraicTopology.Singular.openSingularChainComplexFunctor, def:AlgebraicTopology.Singular.openSubspaceImageIso}
The actual ambient-image homeomorphisms induce chain isomorphisms.

\noindent\emph{Lean:} \texttt{def AlgebraicTopology.Singular.openSubspaceImageChainIso}, \href{https://github.com/Paul-Lez/HodgeConjecture/blob/main/HodgeConjecture/Lemmas/AlgebraicTopology/SingularCochainOpenRestriction.lean#L34}{\texttt{HodgeConjecture/Lemmas/AlgebraicTopology/SingularCochainOpenRestriction.lean:34}}.
\end{definition}

\begin{definition}[AlgebraicTopology.Singular.singularCochainPresheafOpenRestrictionIso]\label{def:AlgebraicTopology.Singular.singularCochainPresheafOpenRestrictionIso}
\lean{AlgebraicTopology.Singular.singularCochainPresheafOpenRestrictionIso}\leanok
\uses{def:AlgebraicTopology.Singular.openSubspaceImageChainIso, def:AlgebraicTopology.Singular.singularCochainPresheaf, def:HomologicalComplex.linearDualIso}
Restricting ambient raw singular cochains to the open subspace gives
intrinsic raw cochains there, by dualizing the actual image chain map.

\noindent\emph{Lean:} \texttt{def AlgebraicTopology.Singular.singularCochainPresheafOpenRestrictionIso}, \href{https://github.com/Paul-Lez/HodgeConjecture/blob/main/HodgeConjecture/Lemmas/AlgebraicTopology/SingularCochainOpenRestriction.lean#L51}{\texttt{HodgeConjecture/Lemmas/AlgebraicTopology/SingularCochainOpenRestriction.lean:51}}.
\end{definition}

\begin{definition}[AlgebraicTopology.Singular.singularCochainPresheafComplexOpenRestrictionIso]\label{def:AlgebraicTopology.Singular.singularCochainPresheafComplexOpenRestrictionIso}
\lean{AlgebraicTopology.Singular.singularCochainPresheafComplexOpenRestrictionIso}\leanok
\uses{def:AlgebraicTopology.Singular.singularCochainPresheafComplex, def:AlgebraicTopology.Singular.singularCochainPresheafOpenRestrictionIso}
The raw open-subspace comparison respects the actual singular
coboundaries, and hence is an isomorphism of complexes.

\noindent\emph{Lean:} \texttt{def AlgebraicTopology.Singular.singularCochainPresheafComplexOpenRestrictionIso}, \href{https://github.com/Paul-Lez/HodgeConjecture/blob/main/HodgeConjecture/Lemmas/AlgebraicTopology/SingularCochainOpenRestriction.lean#L87}{\texttt{HodgeConjecture/Lemmas/AlgebraicTopology/SingularCochainOpenRestriction.lean:87}}.
\end{definition}

\begin{definition}[AlgebraicTopology.Singular.singularCochainSheafOpenRestrictionIso]\label{def:AlgebraicTopology.Singular.singularCochainSheafOpenRestrictionIso}
\lean{AlgebraicTopology.Singular.singularCochainSheafOpenRestrictionIso}\leanok
\uses{def:AlgebraicTopology.Singular.singularCochainPresheafOpenRestrictionIso, def:AlgebraicTopology.Singular.singularCochainSheaf, def:TopCat.Sheaf.openRestrictionSheafificationIso}
Restriction of the ambient singular cochain sheaf is the intrinsic
singular cochain sheaf of the open subspace, with its unit normalization.

\noindent\emph{Lean:} \texttt{def AlgebraicTopology.Singular.singularCochainSheafOpenRestrictionIso}, \href{https://github.com/Paul-Lez/HodgeConjecture/blob/main/HodgeConjecture/Lemmas/AlgebraicTopology/SingularCochainOpenRestriction.lean#L102}{\texttt{HodgeConjecture/Lemmas/AlgebraicTopology/SingularCochainOpenRestriction.lean:102}}.
\end{definition}

\begin{definition}[AlgebraicTopology.Singular.singularCochainSheafComplexOpenRestrictionIso]\label{def:AlgebraicTopology.Singular.singularCochainSheafComplexOpenRestrictionIso}
\lean{AlgebraicTopology.Singular.singularCochainSheafComplexOpenRestrictionIso}\leanok
\uses{def:AlgebraicTopology.Singular.singularCochainSheafCoboundary, def:AlgebraicTopology.Singular.singularCochainSheafComplex, def:AlgebraicTopology.Singular.singularCochainSheafOpenRestrictionIso}
The actual open restriction of the ambient sheafified singular
cochains is the intrinsic complex on that open subspace.

\noindent\emph{Lean:} \texttt{def AlgebraicTopology.Singular.singularCochainSheafComplexOpenRestrictionIso}, \href{https://github.com/Paul-Lez/HodgeConjecture/blob/main/HodgeConjecture/Lemmas/AlgebraicTopology/SingularCochainOpenRestriction.lean#L130}{\texttt{HodgeConjecture/Lemmas/AlgebraicTopology/SingularCochainOpenRestriction.lean:130}}.
\end{definition}

\section{Cokernel preservation for a map with flasque source and kernel}\label{sec:HodgeConjecture.Lemmas.AlgebraicTopology.FlasqueCokernelPreservation}
\noindent\emph{File:} \href{https://github.com/Paul-Lez/HodgeConjecture/blob/main/HodgeConjecture/Lemmas/AlgebraicTopology/FlasqueCokernelPreservation.lean}{\texttt{HodgeConjecture/Lemmas/AlgebraicTopology/FlasqueCokernelPreservation.lean}}.

The sheaf-forgetful functor is not exact in general. For a map whose source and
kernel are flasque, its image is flasque and both associated short exact sequences
remain exact on every open. Consequently its actual cokernel is preserved by the
forgetful functor. This is a local hypothesis on one actual map, not a fabricated
exactness instance for the whole forgetful functor.

\begin{definition}[TopCat.Sheaf.IsFlasque.kernelImageShortComplex]\label{def:TopCat.Sheaf.IsFlasque.kernelImageShortComplex}
\lean{TopCat.Sheaf.IsFlasque.kernelImageShortComplex}\leanok
The actual kernel/source/image short complex.

\noindent\emph{Lean:} \texttt{def TopCat.Sheaf.IsFlasque.kernelImageShortComplex}, \href{https://github.com/Paul-Lez/HodgeConjecture/blob/main/HodgeConjecture/Lemmas/AlgebraicTopology/FlasqueCokernelPreservation.lean#L47}{\texttt{HodgeConjecture/Lemmas/AlgebraicTopology/FlasqueCokernelPreservation.lean:47}}.
\end{definition}

\begin{definition}[TopCat.Sheaf.IsFlasque.imageCokernelShortComplex]\label{def:TopCat.Sheaf.IsFlasque.imageCokernelShortComplex}
\lean{TopCat.Sheaf.IsFlasque.imageCokernelShortComplex}\leanok
The image/codomain/cokernel short exact sequence is the actual abelian image sequence.

\noindent\emph{Lean:} \texttt{def TopCat.Sheaf.IsFlasque.imageCokernelShortComplex}, \href{https://github.com/Paul-Lez/HodgeConjecture/blob/main/HodgeConjecture/Lemmas/AlgebraicTopology/FlasqueCokernelPreservation.lean#L70}{\texttt{HodgeConjecture/Lemmas/AlgebraicTopology/FlasqueCokernelPreservation.lean:70}}.
\end{definition}

\section{Localization between two nested closed supports}\label{sec:HodgeConjecture.Lemmas.AlgebraicTopology.NestedSheafSupportLocalization}
\noindent\emph{File:} \href{https://github.com/Paul-Lez/HodgeConjecture/blob/main/HodgeConjecture/Lemmas/AlgebraicTopology/NestedSheafSupportLocalization.lean}{\texttt{HodgeConjecture/Lemmas/AlgebraicTopology/NestedSheafSupportLocalization.lean}}.

For \texttt{V ⊆ U} open, the actual restrictions \texttt{F → j\_\{U*\}F\textbar{}U → j\_\{V*\}F\textbar{}V}
give a short exact sequence on injective coefficients:
\texttt{0 → Γ\_\{X∖U\} F → Γ\_\{X∖V\} F → ker(j\_\{U*\}F\textbar{}U → j\_\{V*\}F\textbar{}V) → 0}.
The last term is explicitly the sheaf of sections on \texttt{U} vanishing on \texttt{V},
viewed on \texttt{X}. Thus for closed supports \texttt{S ⊆ Z}, the sequence removes \texttt{S}
from \texttt{Z}, by taking \texttt{U = X∖S} and \texttt{V = X∖Z}.

Exactness persists on every open set, since the first restriction is onto for
injective (indeed flasque) coefficients. No local purity or dimension vanishing
is assumed here. The resulting canonical mapping-fiber comparison is the
algebraic localization step needed when extending over singular strata.

\begin{definition}[TopCat.Sheaf.sheafSectionsBetweenOpens]\label{def:TopCat.Sheaf.sheafSectionsBetweenOpens}
\lean{TopCat.Sheaf.sheafSectionsBetweenOpens}\leanok
\uses{def:TopCat.Sheaf.openRestrictionPushforwardMap}
Sections on \texttt{U} vanishing on \texttt{V}, pushed to the original ambient space.

\noindent\emph{Lean:} \texttt{def TopCat.Sheaf.sheafSectionsBetweenOpens}, \href{https://github.com/Paul-Lez/HodgeConjecture/blob/main/HodgeConjecture/Lemmas/AlgebraicTopology/NestedSheafSupportLocalization.lean#L38}{\texttt{HodgeConjecture/Lemmas/AlgebraicTopology/NestedSheafSupportLocalization.lean:38}}.
\end{definition}

\begin{definition}[TopCat.Sheaf.sheafSectionsBetweenOpensOnOpenIso]\label{def:TopCat.Sheaf.sheafSectionsBetweenOpensOnOpenIso}
\lean{TopCat.Sheaf.sheafSectionsBetweenOpensOnOpenIso}\leanok
\uses{def:TopCat.Sheaf.sheafSectionsBetweenOpens, def:TopCat.Sheaf.supportEvaluation}
On each ambient open set the last term is exactly the kernel of actual
restriction from its intersection with \texttt{U} to its intersection with \texttt{V}.

\noindent\emph{Lean:} \texttt{def TopCat.Sheaf.sheafSectionsBetweenOpensOnOpenIso}, \href{https://github.com/Paul-Lez/HodgeConjecture/blob/main/HodgeConjecture/Lemmas/AlgebraicTopology/NestedSheafSupportLocalization.lean#L58}{\texttt{HodgeConjecture/Lemmas/AlgebraicTopology/NestedSheafSupportLocalization.lean:58}}.
\end{definition}

\begin{definition}[TopCat.Sheaf.toSheafSectionsBetweenOpens]\label{def:TopCat.Sheaf.toSheafSectionsBetweenOpens}
\lean{TopCat.Sheaf.toSheafSectionsBetweenOpens}\leanok
\uses{def:TopCat.Sheaf.sheafSectionsBetweenOpens, def:TopCat.Sheaf.sheafSectionsSupportedOutside}
Restrict an ambient section vanishing on \texttt{V} to its section on \texttt{U}.

\noindent\emph{Lean:} \texttt{def TopCat.Sheaf.toSheafSectionsBetweenOpens}, \href{https://github.com/Paul-Lez/HodgeConjecture/blob/main/HodgeConjecture/Lemmas/AlgebraicTopology/NestedSheafSupportLocalization.lean#L73}{\texttt{HodgeConjecture/Lemmas/AlgebraicTopology/NestedSheafSupportLocalization.lean:73}}.
\end{definition}

\begin{definition}[TopCat.Sheaf.nestedSupportRestrictionShortComplex]\label{def:TopCat.Sheaf.nestedSupportRestrictionShortComplex}
\lean{TopCat.Sheaf.nestedSupportRestrictionShortComplex}\leanok
\uses{def:TopCat.Sheaf.sheafSectionsSupportedOutsideMap, def:TopCat.Sheaf.toSheafSectionsBetweenOpens}
The actual inclusion/restriction sequence for two nested complements.

\noindent\emph{Lean:} \texttt{def TopCat.Sheaf.nestedSupportRestrictionShortComplex}, \href{https://github.com/Paul-Lez/HodgeConjecture/blob/main/HodgeConjecture/Lemmas/AlgebraicTopology/NestedSheafSupportLocalization.lean#L100}{\texttt{HodgeConjecture/Lemmas/AlgebraicTopology/NestedSheafSupportLocalization.lean:100}}.
\end{definition}

\begin{definition}[TopCat.Sheaf.nestedSupportRestrictionComplexShortComplex]\label{def:TopCat.Sheaf.nestedSupportRestrictionComplexShortComplex}
\lean{TopCat.Sheaf.nestedSupportRestrictionComplexShortComplex}\leanok
\uses{def:TopCat.Sheaf.sheafSectionsSupportedOutsideMap, def:TopCat.Sheaf.toSheafSectionsBetweenOpens}
The coefficientwise nested-support sequence on a complex of sheaves.

\noindent\emph{Lean:} \texttt{def TopCat.Sheaf.nestedSupportRestrictionComplexShortComplex}, \href{https://github.com/Paul-Lez/HodgeConjecture/blob/main/HodgeConjecture/Lemmas/AlgebraicTopology/NestedSheafSupportLocalization.lean#L137}{\texttt{HodgeConjecture/Lemmas/AlgebraicTopology/NestedSheafSupportLocalization.lean:137}}.
\end{definition}

\begin{definition}[TopCat.Sheaf.nestedSupportRestrictionSectionsComplexShortComplex]\label{def:TopCat.Sheaf.nestedSupportRestrictionSectionsComplexShortComplex}
\lean{TopCat.Sheaf.nestedSupportRestrictionSectionsComplexShortComplex}\leanok
\uses{def:TopCat.Sheaf.nestedSupportRestrictionComplexShortComplex, def:TopCat.Sheaf.supportEvaluation}
The actual nested-support sequence of section complexes on an open set.

\noindent\emph{Lean:} \texttt{def TopCat.Sheaf.nestedSupportRestrictionSectionsComplexShortComplex}, \href{https://github.com/Paul-Lez/HodgeConjecture/blob/main/HodgeConjecture/Lemmas/AlgebraicTopology/NestedSheafSupportLocalization.lean#L152}{\texttt{HodgeConjecture/Lemmas/AlgebraicTopology/NestedSheafSupportLocalization.lean:152}}.
\end{definition}

\section{The open-set model of the last localization term}\label{sec:HodgeConjecture.Lemmas.AlgebraicTopology.NestedSheafSupportOnOpen}
\noindent\emph{File:} \href{https://github.com/Paul-Lez/HodgeConjecture/blob/main/HodgeConjecture/Lemmas/AlgebraicTopology/NestedSheafSupportOnOpen.lean}{\texttt{HodgeConjecture/Lemmas/AlgebraicTopology/NestedSheafSupportOnOpen.lean}}.

For \texttt{V ≤ U}, global sections of the sheaf of sections on \texttt{U} vanishing on \texttt{V} are
canonically sections on \texttt{U} of the sheaf of ambient sections vanishing on \texttt{V}, both sides
being the kernel of restriction \texttt{F(U) → F(V)}. The comparison is functorial in the
coefficient sheaf, so it identifies the last complex in nested-support localization with
the supported-section complex on the complement of the smaller closed support.

\begin{definition}[TopCat.Sheaf.openRestrictionPushforwardTopEvaluationIso]\label{def:TopCat.Sheaf.openRestrictionPushforwardTopEvaluationIso}
\lean{TopCat.Sheaf.openRestrictionPushforwardTopEvaluationIso}\leanok
\uses{def:TopCat.Sheaf.openRestrictionImage, def:TopCat.Sheaf.openRestrictionPushforward, def:TopCat.Sheaf.supportEvaluation}
The literal equality of opens identifies global sections of the open
pushforward with sections of the original sheaf on that open.

\noindent\emph{Lean:} \texttt{def TopCat.Sheaf.openRestrictionPushforwardTopEvaluationIso}, \href{https://github.com/Paul-Lez/HodgeConjecture/blob/main/HodgeConjecture/Lemmas/AlgebraicTopology/NestedSheafSupportOnOpen.lean#L38}{\texttt{HodgeConjecture/Lemmas/AlgebraicTopology/NestedSheafSupportOnOpen.lean:38}}.
\end{definition}

\begin{definition}[TopCat.Sheaf.sheafSectionsBetweenOpensInclusion]\label{def:TopCat.Sheaf.sheafSectionsBetweenOpensInclusion}
\lean{TopCat.Sheaf.sheafSectionsBetweenOpensInclusion}\leanok
\uses{def:TopCat.Sheaf.sheafSectionsBetweenOpens}
The last localization kernel includes into the actual sections on \texttt{U}.

\noindent\emph{Lean:} \texttt{def TopCat.Sheaf.sheafSectionsBetweenOpensInclusion}, \href{https://github.com/Paul-Lez/HodgeConjecture/blob/main/HodgeConjecture/Lemmas/AlgebraicTopology/NestedSheafSupportOnOpen.lean#L46}{\texttt{HodgeConjecture/Lemmas/AlgebraicTopology/NestedSheafSupportOnOpen.lean:46}}.
\end{definition}

\begin{definition}[TopCat.Sheaf.nestedSupportRestrictionTargetIso]\label{def:TopCat.Sheaf.nestedSupportRestrictionTargetIso}
\lean{TopCat.Sheaf.nestedSupportRestrictionTargetIso}\leanok
\uses{def:TopCat.Sheaf.openRestrictionImage, def:TopCat.Sheaf.openRestrictionPushforward}
The two target section groups are identified by their common actual open
\texttt{V}, using \texttt{V ≤ U}.

\noindent\emph{Lean:} \texttt{def TopCat.Sheaf.nestedSupportRestrictionTargetIso}, \href{https://github.com/Paul-Lez/HodgeConjecture/blob/main/HodgeConjecture/Lemmas/AlgebraicTopology/NestedSheafSupportOnOpen.lean#L54}{\texttt{HodgeConjecture/Lemmas/AlgebraicTopology/NestedSheafSupportOnOpen.lean:54}}.
\end{definition}

\begin{definition}[TopCat.Sheaf.sheafSectionsBetweenOpensGlobalIso]\label{def:TopCat.Sheaf.sheafSectionsBetweenOpensGlobalIso}
\lean{TopCat.Sheaf.sheafSectionsBetweenOpensGlobalIso}\leanok
\uses{def:TopCat.Sheaf.nestedSupportRestrictionTargetIso, def:TopCat.Sheaf.openRestrictionPushforwardTopEvaluationIso, def:TopCat.Sheaf.sheafSectionsBetweenOpensOnOpenIso, def:TopCat.Sheaf.sheafSectionsSupportedOutsideOnOpenIso}
Canonical kernel comparison between the last global localization term and
supported sections on the actual open complement.

\noindent\emph{Lean:} \texttt{def TopCat.Sheaf.sheafSectionsBetweenOpensGlobalIso}, \href{https://github.com/Paul-Lez/HodgeConjecture/blob/main/HodgeConjecture/Lemmas/AlgebraicTopology/NestedSheafSupportOnOpen.lean#L74}{\texttt{HodgeConjecture/Lemmas/AlgebraicTopology/NestedSheafSupportOnOpen.lean:74}}.
\end{definition}

\begin{definition}[TopCat.Sheaf.sheafSectionsBetweenOpensGlobalNatIso]\label{def:TopCat.Sheaf.sheafSectionsBetweenOpensGlobalNatIso}
\lean{TopCat.Sheaf.sheafSectionsBetweenOpensGlobalNatIso}\leanok
\uses{def:TopCat.Sheaf.sheafSectionsBetweenOpensGlobalIso, def:TopCat.Sheaf.sheafSectionsBetweenOpensInclusion, def:TopCat.Sheaf.sheafSectionsSupportedOutsideInclusion}
The comparison is a natural isomorphism of actual section functors.

\noindent\emph{Lean:} \texttt{def TopCat.Sheaf.sheafSectionsBetweenOpensGlobalNatIso}, \href{https://github.com/Paul-Lez/HodgeConjecture/blob/main/HodgeConjecture/Lemmas/AlgebraicTopology/NestedSheafSupportOnOpen.lean#L135}{\texttt{HodgeConjecture/Lemmas/AlgebraicTopology/NestedSheafSupportOnOpen.lean:135}}.
\end{definition}

\begin{definition}[TopCat.Sheaf.nestedSupportRestrictionLastComplexIso]\label{def:TopCat.Sheaf.nestedSupportRestrictionLastComplexIso}
\lean{TopCat.Sheaf.nestedSupportRestrictionLastComplexIso}\leanok
\uses{def:TopCat.Sheaf.nestedSupportRestrictionSectionsComplexShortComplex, def:TopCat.Sheaf.sheafSectionsBetweenOpensGlobalNatIso}
The last complex in global nested-support localization is the actual
supported-section complex on the complement \texttt{U} of the smaller support.

\noindent\emph{Lean:} \texttt{def TopCat.Sheaf.nestedSupportRestrictionLastComplexIso}, \href{https://github.com/Paul-Lez/HodgeConjecture/blob/main/HodgeConjecture/Lemmas/AlgebraicTopology/NestedSheafSupportOnOpen.lean#L142}{\texttt{HodgeConjecture/Lemmas/AlgebraicTopology/NestedSheafSupportOnOpen.lean:142}}.
\end{definition}

\section{Restriction naturality of the actual supported-section kernel comparison}\label{sec:HodgeConjecture.Lemmas.AlgebraicTopology.SupportedSectionRestrictionConeNaturality}
\noindent\emph{File:} \href{https://github.com/Paul-Lez/HodgeConjecture/blob/main/HodgeConjecture/Lemmas/AlgebraicTopology/SupportedSectionRestrictionConeNaturality.lean}{\texttt{HodgeConjecture/Lemmas/AlgebraicTopology/SupportedSectionRestrictionConeNaturality.lean}}.

Every map below is induced by the literal restriction of sheaf sections. In
particular the short-exact-sequence lift is the positive canonical lift; its
naturality is established before passing to homology or sheafification.

\begin{definition}[TopCat.Sheaf.sectionComplexRestrictionConeMap]\label{def:TopCat.Sheaf.sectionComplexRestrictionConeMap}
\lean{TopCat.Sheaf.sectionComplexRestrictionConeMap}\leanok
\uses{def:TopCat.Sheaf.sectionComplexRestriction}
The literal restriction on the cones of section restriction.

\noindent\emph{Lean:} \texttt{def TopCat.Sheaf.sectionComplexRestrictionConeMap}, \href{https://github.com/Paul-Lez/HodgeConjecture/blob/main/HodgeConjecture/Lemmas/AlgebraicTopology/SupportedSectionRestrictionConeNaturality.lean#L43}{\texttt{HodgeConjecture/Lemmas/AlgebraicTopology/SupportedSectionRestrictionConeNaturality.lean:43}}.
\end{definition}

\begin{definition}[TopCat.Sheaf.supportRestrictionSectionsComplexMap]\label{def:TopCat.Sheaf.supportRestrictionSectionsComplexMap}
\lean{TopCat.Sheaf.supportRestrictionSectionsComplexMap}\leanok
\uses{def:TopCat.Sheaf.sectionComplexRestriction, def:TopCat.Sheaf.supportRestrictionSectionsComplexShortComplex}
The actual map between the three terms of the supported-section sequence.

\noindent\emph{Lean:} \texttt{def TopCat.Sheaf.supportRestrictionSectionsComplexMap}, \href{https://github.com/Paul-Lez/HodgeConjecture/blob/main/HodgeConjecture/Lemmas/AlgebraicTopology/SupportedSectionRestrictionConeNaturality.lean#L58}{\texttt{HodgeConjecture/Lemmas/AlgebraicTopology/SupportedSectionRestrictionConeNaturality.lean:58}}.
\end{definition}

\section{Actual local restriction-cone naturality}\label{sec:HodgeConjecture.Lemmas.AlgebraicTopology.SingularCochainOpenConeNaturality}
\noindent\emph{File:} \href{https://github.com/Paul-Lez/HodgeConjecture/blob/main/HodgeConjecture/Lemmas/AlgebraicTopology/SingularCochainOpenConeNaturality.lean}{\texttt{HodgeConjecture/Lemmas/AlgebraicTopology/SingularCochainOpenConeNaturality.lean}}.

\begin{definition}[AlgebraicTopology.Singular.openInclusionPairMap]\label{def:AlgebraicTopology.Singular.openInclusionPairMap}
\lean{AlgebraicTopology.Singular.openInclusionPairMap}\leanok
\uses{def:AlgebraicTopology.Singular.openInclusionPair}
The actual map of pairs induced by a square of ambient open inclusions.

\noindent\emph{Lean:} \texttt{def AlgebraicTopology.Singular.openInclusionPairMap}, \href{https://github.com/Paul-Lez/HodgeConjecture/blob/main/HodgeConjecture/Lemmas/AlgebraicTopology/SingularCochainOpenConeNaturality.lean#L21}{\texttt{HodgeConjecture/Lemmas/AlgebraicTopology/SingularCochainOpenConeNaturality.lean:21}}.
\end{definition}

\begin{definition}[AlgebraicTopology.Singular.openRawSingularRestrictionConeMap]\label{def:AlgebraicTopology.Singular.openRawSingularRestrictionConeMap}
\lean{AlgebraicTopology.Singular.openRawSingularRestrictionConeMap}\leanok
\uses{def:AlgebraicTopology.Singular.openRawSingularRestrictionCone}
Actual restriction on raw local cones.

\noindent\emph{Lean:} \texttt{def AlgebraicTopology.Singular.openRawSingularRestrictionConeMap}, \href{https://github.com/Paul-Lez/HodgeConjecture/blob/main/HodgeConjecture/Lemmas/AlgebraicTopology/SingularCochainOpenConeNaturality.lean#L51}{\texttt{HodgeConjecture/Lemmas/AlgebraicTopology/SingularCochainOpenConeNaturality.lean:51}}.
\end{definition}

\begin{definition}[AlgebraicTopology.Singular.openSingularSheafRestrictionConeMap]\label{def:AlgebraicTopology.Singular.openSingularSheafRestrictionConeMap}
\lean{AlgebraicTopology.Singular.openSingularSheafRestrictionConeMap}\leanok
\uses{def:AlgebraicTopology.Singular.openSingularSheafRestrictionCone, def:AlgebraicTopology.Singular.singularCochainSheaf}
Actual restriction on sheaf-section local cones.

\noindent\emph{Lean:} \texttt{def AlgebraicTopology.Singular.openSingularSheafRestrictionConeMap}, \href{https://github.com/Paul-Lez/HodgeConjecture/blob/main/HodgeConjecture/Lemmas/AlgebraicTopology/SingularCochainOpenConeNaturality.lean#L60}{\texttt{HodgeConjecture/Lemmas/AlgebraicTopology/SingularCochainOpenConeNaturality.lean:60}}.
\end{definition}

\begin{definition}[AlgebraicTopology.Singular.forgottenRelativeCochainConeMap]\label{def:AlgebraicTopology.Singular.forgottenRelativeCochainConeMap}
\lean{AlgebraicTopology.Singular.forgottenRelativeCochainConeMap}\leanok
\uses{def:AlgebraicTopology.Singular.relativeCochainRestrictionInt, def:AlgebraicTopology.Singular.relativeDualCochainShortComplexIntMap}
The actual relative cone map after forgetting scalars termwise.

\noindent\emph{Lean:} \texttt{def AlgebraicTopology.Singular.forgottenRelativeCochainConeMap}, \href{https://github.com/Paul-Lez/HodgeConjecture/blob/main/HodgeConjecture/Lemmas/AlgebraicTopology/SingularCochainOpenConeNaturality.lean#L97}{\texttt{HodgeConjecture/Lemmas/AlgebraicTopology/SingularCochainOpenConeNaturality.lean:97}}.
\end{definition}

\begin{definition}[AlgebraicTopology.Singular.openRawSingularRestrictionConeIsoForgottenRelative]\label{def:AlgebraicTopology.Singular.openRawSingularRestrictionConeIsoForgottenRelative}
\lean{AlgebraicTopology.Singular.openRawSingularRestrictionConeIsoForgottenRelative}\leanok
\uses{def:AlgebraicTopology.Singular.openRawSingularRestrictionConeIsoRelative}
The raw local cone identified with the forgotten relative cone,
including the canonical comparison for forgetting scalars.

\noindent\emph{Lean:} \texttt{def AlgebraicTopology.Singular.openRawSingularRestrictionConeIsoForgottenRelative}, \href{https://github.com/Paul-Lez/HodgeConjecture/blob/main/HodgeConjecture/Lemmas/AlgebraicTopology/SingularCochainOpenConeNaturality.lean#L132}{\texttt{HodgeConjecture/Lemmas/AlgebraicTopology/SingularCochainOpenConeNaturality.lean:132}}.
\end{definition}

\begin{definition}[AlgebraicTopology.Singular.openRawSingularRestrictionConeHomologyIso]\label{def:AlgebraicTopology.Singular.openRawSingularRestrictionConeHomologyIso}
\lean{AlgebraicTopology.Singular.openRawSingularRestrictionConeHomologyIso}\leanok
\uses{def:AlgebraicTopology.Singular.openRawSingularRestrictionConeIsoForgottenRelative}
The actual homology comparison, retaining the scalar-forgetting
homology isomorphism as a separate canonical factor.

\noindent\emph{Lean:} \texttt{def AlgebraicTopology.Singular.openRawSingularRestrictionConeHomologyIso}, \href{https://github.com/Paul-Lez/HodgeConjecture/blob/main/HodgeConjecture/Lemmas/AlgebraicTopology/SingularCochainOpenConeNaturality.lean#L166}{\texttt{HodgeConjecture/Lemmas/AlgebraicTopology/SingularCochainOpenConeNaturality.lean:166}}.
\end{definition}

\begin{definition}[AlgebraicTopology.Singular.openRawToSingularSheafRestrictionConeHomologyIso]\label{def:AlgebraicTopology.Singular.openRawToSingularSheafRestrictionConeHomologyIso}
\lean{AlgebraicTopology.Singular.openRawToSingularSheafRestrictionConeHomologyIso}\leanok
\uses{def:AlgebraicTopology.Singular.ambientOpen, def:AlgebraicTopology.Singular.coverMemberToSmallRationalSingularChains, def:AlgebraicTopology.Singular.coveringSieveOpenFamily, def:AlgebraicTopology.Singular.coveringSieveOpenFamilyOn, def:AlgebraicTopology.Singular.globalRawSingularCochainComplexIso, def:AlgebraicTopology.Singular.openFamilySieve, def:AlgebraicTopology.Singular.openRawSingularCochainComplexIsoGlobal, def:AlgebraicTopology.Singular.openRawToSingularSheafRestrictionCone, def:AlgebraicTopology.Singular.openSimplexMap, def:AlgebraicTopology.Singular.openSingularCochainSheafComplexIsoGlobal, def:AlgebraicTopology.Singular.rationalCochainHomotopyEquivCoverSmall, def:AlgebraicTopology.Singular.rationalCochainRestrictionToCoverSmall, def:AlgebraicTopology.Singular.singularCochainComplexIsoTopOpen, def:AlgebraicTopology.Singular.singularCochainEquivSimplexFunction, def:AlgebraicTopology.Singular.singularCochainOpenConeDerivedCategory, def:AlgebraicTopology.Singular.singularCochainToPlusPlusPresheafComplex, def:AlgebraicTopology.Singular.topOpenToGlobalSingularCochainSheafComplex}
Homology of the actual raw-to-sheaf cone map.

\noindent\emph{Lean:} \texttt{def AlgebraicTopology.Singular.openRawToSingularSheafRestrictionConeHomologyIso}, \href{https://github.com/Paul-Lez/HodgeConjecture/blob/main/HodgeConjecture/Lemmas/AlgebraicTopology/SingularCochainOpenConeNaturality.lean#L251}{\texttt{HodgeConjecture/Lemmas/AlgebraicTopology/SingularCochainOpenConeNaturality.lean:251}}.
\end{definition}

\section{Naturality of actual supported singular-section relative cohomology}\label{sec:HodgeConjecture.Lemmas.AlgebraicTopology.SupportedSingularSectionNaturality}
\noindent\emph{File:} \href{https://github.com/Paul-Lez/HodgeConjecture/blob/main/HodgeConjecture/Lemmas/AlgebraicTopology/SupportedSingularSectionNaturality.lean}{\texttt{HodgeConjecture/Lemmas/AlgebraicTopology/SupportedSingularSectionNaturality.lean}}.

\begin{definition}[AlgebraicTopology.Singular.supportedSingularSectionConeHomologyIso]\label{def:AlgebraicTopology.Singular.supportedSingularSectionConeHomologyIso}
\lean{AlgebraicTopology.Singular.supportedSingularSectionConeHomologyIso}\leanok
\uses{def:AlgebraicTopology.Singular.ambientOpen, def:AlgebraicTopology.Singular.coveringSieveOpenFamilyOn, def:AlgebraicTopology.Singular.openFamilySieve, def:AlgebraicTopology.Singular.openSimplexMap, def:AlgebraicTopology.Singular.openSingularSheafRestrictionCone, def:AlgebraicTopology.Singular.singularCochainEquivSimplexFunction, def:AlgebraicTopology.Singular.singularCochainSheaf, def:AlgebraicTopology.Singular.supportedRationalSingularCochainComplex, def:TopCat.Sheaf.sectionComplexRestrictionExtendConeIso, def:TopCat.Sheaf.supportedSectionHomologyIsoRestrictionCone}
The actual two-step supported kernel and grading comparison, before the relative calculation.

\noindent\emph{Lean:} \texttt{def AlgebraicTopology.Singular.supportedSingularSectionConeHomologyIso}, \href{https://github.com/Paul-Lez/HodgeConjecture/blob/main/HodgeConjecture/Lemmas/AlgebraicTopology/SupportedSingularSectionNaturality.lean#L24}{\texttt{HodgeConjecture/Lemmas/AlgebraicTopology/SupportedSingularSectionNaturality.lean:24}}.
\end{definition}

\chapter{Lemmas: Analysis}

\section{Wedges of exact one-forms}\label{sec:HodgeConjecture.Lemmas.Analysis.Calculus.DifferentialForm.ExactWedge}
\noindent\emph{File:} \href{https://github.com/Paul-Lez/HodgeConjecture/blob/main/HodgeConjecture/Lemmas/Analysis/Calculus/DifferentialForm/ExactWedge.lean}{\texttt{HodgeConjecture/Lemmas/Analysis/Calculus/DifferentialForm/ExactWedge.lean}}.

Wedging the differentials of \texttt{p} scalar functions, all differentiated within one set, gives a
differential \texttt{p}-form on that set. This file proves it is closed, by identifying it with the
pullback of the constant volume form on \texttt{ℂ\^{}p} along the map assembling the functions and applying
naturality of the exterior derivative.

Multiplying such a wedge by an analytic scalar therefore has exterior derivative the wedge of the
scalar's differential with it — the only case of the Leibniz rule the holomorphic de Rham complex
needs.

\begin{definition}[DifferentialForm.exactWedgeWithin]\label{def:DifferentialForm.exactWedgeWithin}
\lean{DifferentialForm.exactWedgeWithin}\leanok
\uses{def:ContinuousAlternatingMap.wedgeCovectors}
A wedge of differentials of scalar functions, all differentiated within the same set.

\noindent\emph{Lean:} \texttt{def DifferentialForm.exactWedgeWithin}, \href{https://github.com/Paul-Lez/HodgeConjecture/blob/main/HodgeConjecture/Lemmas/Analysis/Calculus/DifferentialForm/ExactWedge.lean#L59}{\texttt{HodgeConjecture/Lemmas/Analysis/Calculus/DifferentialForm/ExactWedge.lean:59}}.
\end{definition}

\section{The holomorphic Poincaré lemma on a ball}\label{sec:HodgeConjecture.Lemmas.Analysis.Calculus.DifferentialForm.HolomorphicPoincare}
\noindent\emph{File:} \href{https://github.com/Paul-Lez/HodgeConjecture/blob/main/HodgeConjecture/Lemmas/Analysis/Calculus/DifferentialForm/HolomorphicPoincare.lean}{\texttt{HodgeConjecture/Lemmas/Analysis/Calculus/DifferentialForm/HolomorphicPoincare.lean}}.

Integration over the real radial parameter preserves complex analyticity. This is proved directly:
the radial homotopy operator of \texttt{Poincare} applied to a form given by a convergent formal
multilinear series is again given by such a series, obtained by integrating the homogeneous terms.

Consequently a closed analytic differential form on a ball has an analytic primitive there, and
translating the ball to the origin gives the same statement about a ball centred anywhere.

\begin{definition}[DifferentialForm.radialPrimitiveSeries]\label{def:DifferentialForm.radialPrimitiveSeries}
\lean{DifferentialForm.radialPrimitiveSeries}\leanok
The formal power series of the radial primitive.

\noindent\emph{Lean:} \texttt{def DifferentialForm.radialPrimitiveSeries}, \href{https://github.com/Paul-Lez/HodgeConjecture/blob/main/HodgeConjecture/Lemmas/Analysis/Calculus/DifferentialForm/HolomorphicPoincare.lean#L43}{\texttt{HodgeConjecture/Lemmas/Analysis/Calculus/DifferentialForm/HolomorphicPoincare.lean:43}}.
\end{definition}

\begin{definition}[DifferentialForm.translateForm]\label{def:DifferentialForm.translateForm}
\lean{DifferentialForm.translateForm}\leanok
Translate the base point of a differential form field to the origin.

\noindent\emph{Lean:} \texttt{def DifferentialForm.translateForm}, \href{https://github.com/Paul-Lez/HodgeConjecture/blob/main/HodgeConjecture/Lemmas/Analysis/Calculus/DifferentialForm/HolomorphicPoincare.lean#L342}{\texttt{HodgeConjecture/Lemmas/Analysis/Calculus/DifferentialForm/HolomorphicPoincare.lean:342}}.
\end{definition}

\begin{definition}[DifferentialForm.untranslateForm]\label{def:DifferentialForm.untranslateForm}
\lean{DifferentialForm.untranslateForm}\leanok
Translate a differential form field from origin-centered coordinates back to a center.

\noindent\emph{Lean:} \texttt{def DifferentialForm.untranslateForm}, \href{https://github.com/Paul-Lez/HodgeConjecture/blob/main/HodgeConjecture/Lemmas/Analysis/Calculus/DifferentialForm/HolomorphicPoincare.lean#L372}{\texttt{HodgeConjecture/Lemmas/Analysis/Calculus/DifferentialForm/HolomorphicPoincare.lean:372}}.
\end{definition}

\chapter{Lemmas: LinearAlgebra}

\section{Orientations of complex vector spaces}\label{sec:HodgeConjecture.Lemmas.LinearAlgebra.ComplexOrientation}
\noindent\emph{File:} \href{https://github.com/Paul-Lez/HodgeConjecture/blob/main/HodgeConjecture/Lemmas/LinearAlgebra/ComplexOrientation.lean}{\texttt{HodgeConjecture/Lemmas/LinearAlgebra/ComplexOrientation.lean}}.

Every complex-linear automorphism preserves either real orientation of its underlying real vector
space.  This file proves that fact from the determinant norm formula and defines the standard
orientation of \texttt{Fin n → ℂ}, with real and imaginary basis vectors interleaved.

The generic preservation theorem is phrased using mathlib's \texttt{Orientation}.  In particular, it is
the linear-algebra input needed to construct the constant-sign \texttt{Manifold.OrientationLift} proposed
in \href{https://github.com/leanprover-community/mathlib4/pull/35376}{mathlib4 PR \#35376}: apply
\texttt{Manifold.OrientationLift.compatible\_of\_det} to the positivity result after identifying a
holomorphic tangent coordinate change with the scalar restriction of a complex-linear
equivalence.

\begin{definition}[Complex.piBasisOneI]\label{def:Complex.piBasisOneI}
\lean{Complex.piBasisOneI}\leanok
The standard real basis of \texttt{Fin n → ℂ}, ordered
\texttt{(re z 0, im z 0, re z 1, im z 1, ...)}.

It is obtained by composing the standard complex basis of the function space with \texttt{basisOneI},
then reindexing the product basis by \texttt{finProdFinEquiv}.

\noindent\emph{Lean:} \texttt{def Complex.piBasisOneI}, \href{https://github.com/Paul-Lez/HodgeConjecture/blob/main/HodgeConjecture/Lemmas/LinearAlgebra/ComplexOrientation.lean#L74}{\texttt{HodgeConjecture/Lemmas/LinearAlgebra/ComplexOrientation.lean:74}}.
\end{definition}

\begin{definition}[Complex.piCoordCLE]\label{def:Complex.piCoordCLE}
\lean{Complex.piCoordCLE}\leanok
\uses{def:Complex.piBasisOneI}
The real-coordinate identification of \texttt{Fin n → ℂ} with \texttt{Fin (n * 2) → ℝ}, listing the real
and imaginary part of each complex coordinate consecutively.

This is the coordinate map of \texttt{piBasisOneI}. It is recorded as a continuous linear equivalence so
that continuity, linearity and \texttt{map\_eq\_zero\_iff} all come from Mathlib rather than being reproved
for the underlying function.

\noindent\emph{Lean:} \texttt{def Complex.piCoordCLE}, \href{https://github.com/Paul-Lez/HodgeConjecture/blob/main/HodgeConjecture/Lemmas/LinearAlgebra/ComplexOrientation.lean#L83}{\texttt{HodgeConjecture/Lemmas/LinearAlgebra/ComplexOrientation.lean:83}}.
\end{definition}

\chapter{Mathlib: Algebra}

\section{Dimension formulas for polynomial rings}\label{sec:HodgeConjecture.Mathlib.Algebra.PolynomialCatenary}
\noindent\emph{File:} \href{https://github.com/Paul-Lez/HodgeConjecture/blob/main/HodgeConjecture/Mathlib/Algebra/PolynomialCatenary.lean}{\texttt{HodgeConjecture/Mathlib/Algebra/PolynomialCatenary.lean}}.

This file develops dimension results for prime quotients of polynomial rings over a field.

\paragraph{Provenance}

The \texttt{Triangular} namespace below (from \texttt{namespace Triangular} to \texttt{end Triangular}) reproduces,
essentially character for character, the \texttt{equivT} section (lines 58-170) of Mathlib's
\texttt{Mathlib/RingTheory/NoetherNormalization.lean}, added to Mathlib in
\texttt{leanprover-community/mathlib4\#18247}, whose copyright header reads "Copyright (c) 2025 Sihan Su.
All rights reserved. Released under Apache 2.0 license as described in the file LICENSE. Authors:
Riccardo Brasca, Sihan Su, Wan Lin, Xiaoyang Su". Relative to that section the copy drops
\texttt{private}, rewrites or drops the docstrings, and renames two identifiers: \texttt{T\_leadingcoeff\_isUnit}
to \texttt{T\_leadingCoeff\_isUnit}, and the local hypothesis \texttt{coeff} to \texttt{coeff\_eq}. It exists only because
those declarations are \texttt{private} upstream and hence cannot be imported; of them only \texttt{Triangular.T}
and \texttt{Triangular.T\_leadingCoeff\_isUnit} are used below, in
\texttt{MvPolynomial.height\_add\_ringKrullDim\_quotient\_eq\_fin}. Mathlib pull request
\texttt{leanprover-community/mathlib4\#43700} makes those two declarations public; once it lands and the
Mathlib dependency is bumped, this block should be deleted and the upstream declarations used
instead.

\begin{definition}[PolynomialCatenary.Triangular.T1]\label{def:PolynomialCatenary.Triangular.T1}
\lean{PolynomialCatenary.Triangular.T1}\leanok
The triangular change of variables used in Noether normalization.

\noindent\emph{Lean:} \texttt{abbrev PolynomialCatenary.Triangular.T1}, \href{https://github.com/Paul-Lez/HodgeConjecture/blob/main/HodgeConjecture/Mathlib/Algebra/PolynomialCatenary.lean#L153}{\texttt{HodgeConjecture/Mathlib/Algebra/PolynomialCatenary.lean:153}}.
\end{definition}

\begin{definition}[PolynomialCatenary.Triangular.T]\label{def:PolynomialCatenary.Triangular.T}
\lean{PolynomialCatenary.Triangular.T}\leanok
\uses{def:PolynomialCatenary.Triangular.T1}
The triangular automorphism attached to a nonzero polynomial.

\noindent\emph{Lean:} \texttt{abbrev PolynomialCatenary.Triangular.T}, \href{https://github.com/Paul-Lez/HodgeConjecture/blob/main/HodgeConjecture/Mathlib/Algebra/PolynomialCatenary.lean#L163}{\texttt{HodgeConjecture/Mathlib/Algebra/PolynomialCatenary.lean:163}}.
\end{definition}

\section{Inclusion of a stupid truncation}\label{sec:HodgeConjecture.Mathlib.Algebra.Homology.StupidTruncation}
\noindent\emph{File:} \href{https://github.com/Paul-Lez/HodgeConjecture/blob/main/HodgeConjecture/Mathlib/Algebra/Homology/StupidTruncation.lean}{\texttt{HodgeConjecture/Mathlib/Algebra/Homology/StupidTruncation.lean}}.

For an embedding of complex shapes whose image is closed under the next differential, the stupid
truncation is canonically a subcomplex. This file constructs its inclusion into the original
complex and proves naturality. It supplies the inclusion noted as missing in Mathlib's
\texttt{Embedding.StupidTrunc}.

\begin{definition}[HomologicalComplex.stupidTruncInclusionApp]\label{def:HomologicalComplex.stupidTruncInclusionApp}
\lean{HomologicalComplex.stupidTruncInclusionApp}\leanok
The component of the inclusion from a stupid truncation. It is the canonical isomorphism in
the image of the embedding and the zero morphism outside that image.

\noindent\emph{Lean:} \texttt{def HomologicalComplex.stupidTruncInclusionApp}, \href{https://github.com/Paul-Lez/HodgeConjecture/blob/main/HodgeConjecture/Mathlib/Algebra/Homology/StupidTruncation.lean#L43}{\texttt{HodgeConjecture/Mathlib/Algebra/Homology/StupidTruncation.lean:43}}.
\end{definition}

\begin{definition}[HomologicalComplex.stupidTruncInclusion]\label{def:HomologicalComplex.stupidTruncInclusion}
\lean{HomologicalComplex.stupidTruncInclusion}\leanok
\uses{def:HomologicalComplex.stupidTruncInclusionApp}
The canonical inclusion of a stupid truncation whose retained degrees are closed under the
next differential.

\noindent\emph{Lean:} \texttt{def HomologicalComplex.stupidTruncInclusion}, \href{https://github.com/Paul-Lez/HodgeConjecture/blob/main/HodgeConjecture/Mathlib/Algebra/Homology/StupidTruncation.lean#L63}{\texttt{HodgeConjecture/Mathlib/Algebra/Homology/StupidTruncation.lean:63}}.
\end{definition}

\section{Applying an additive functor commutes with extension of a complex}\label{sec:HodgeConjecture.Mathlib.Algebra.Homology.MapExtend}
\noindent\emph{File:} \href{https://github.com/Paul-Lez/HodgeConjecture/blob/main/HodgeConjecture/Mathlib/Algebra/Homology/MapExtend.lean}{\texttt{HodgeConjecture/Mathlib/Algebra/Homology/MapExtend.lean}}.

The comparison is the identity in retained degrees and the canonical isomorphism between
zero objects in newly inserted degrees. In particular it fixes the normalization when
regrading a chain complex by \texttt{n ↦ -n} before or after a sheaf functor.

\begin{definition}[HomologicalComplex.mapExtendCanonicalXIso]\label{def:HomologicalComplex.mapExtendCanonicalXIso}
\lean{HomologicalComplex.mapExtendCanonicalXIso}\leanok
Degreewise comparison: identity on retained degrees, canonical zero-object comparison
on inserted degrees.

\noindent\emph{Lean:} \texttt{def HomologicalComplex.mapExtendCanonicalXIso}, \href{https://github.com/Paul-Lez/HodgeConjecture/blob/main/HodgeConjecture/Mathlib/Algebra/Homology/MapExtend.lean#L29}{\texttt{HodgeConjecture/Mathlib/Algebra/Homology/MapExtend.lean:29}}.
\end{definition}

\begin{definition}[HomologicalComplex.mapExtendCanonicalIso]\label{def:HomologicalComplex.mapExtendCanonicalIso}
\lean{HomologicalComplex.mapExtendCanonicalIso}\leanok
\uses{def:HomologicalComplex.mapExtendCanonicalXIso}
The canonical chain-complex comparison between functorial mapping and extension.

\noindent\emph{Lean:} \texttt{def HomologicalComplex.mapExtendCanonicalIso}, \href{https://github.com/Paul-Lez/HodgeConjecture/blob/main/HodgeConjecture/Mathlib/Algebra/Homology/MapExtend.lean#L46}{\texttt{HodgeConjecture/Mathlib/Algebra/Homology/MapExtend.lean:46}}.
\end{definition}

\chapter{Mathlib: AlgebraicGeometry}

\section{Algebra structure on section}\label{sec:HodgeConjecture.Mathlib.AlgebraicGeometry.Over.Basic}
\noindent\emph{File:} \href{https://github.com/Paul-Lez/HodgeConjecture/blob/main/HodgeConjecture/Mathlib/AlgebraicGeometry/Over/Basic.lean}{\texttt{HodgeConjecture/Mathlib/AlgebraicGeometry/Over/Basic.lean}}.

When working with schemes over a fixed ring, the section rings of the scheme have a natural
algebra structure that is not currently in Mathlib. We make this an \texttt{instance\_reducible} def
instead of an instance to avoid a possible diamond with the instance in
\texttt{Mathlib.AlgebraicGeometry.Group.Affine}.

\begin{definition}[AlgebraicGeometry.overSpecAlgebra]\label{def:AlgebraicGeometry.overSpecAlgebra}
\lean{AlgebraicGeometry.overSpecAlgebra}\leanok
The algebra structure on section rings of a scheme over an affine variety.

\noindent\emph{Lean:} \texttt{def AlgebraicGeometry.overSpecAlgebra}, \href{https://github.com/Paul-Lez/HodgeConjecture/blob/main/HodgeConjecture/Mathlib/AlgebraicGeometry/Over/Basic.lean#L37}{\texttt{HodgeConjecture/Mathlib/AlgebraicGeometry/Over/Basic.lean:37}}.
\end{definition}

\chapter{Mathlib: Analysis}

\section{Wedge products of continuous covectors}\label{sec:HodgeConjecture.Mathlib.Analysis.NormedSpace.WedgeCovectors}
\noindent\emph{File:} \href{https://github.com/Paul-Lez/HodgeConjecture/blob/main/HodgeConjecture/Mathlib/Analysis/NormedSpace/WedgeCovectors.lean}{\texttt{HodgeConjecture/Mathlib/Analysis/NormedSpace/WedgeCovectors.lean}}.

The wedge of \texttt{p} continuous linear functionals on a complex normed space is the alternating
\texttt{p}-form whose value on a family of vectors is the determinant of the matrix of pairings. This
file builds it by iterated alternatizing uncurry, proves the determinant formula, and derives the
multilinearity and alternation properties in each covector slot.

On \texttt{𝕜\^{}d} the coordinate projections wedge to the standard volume form, and every continuous
alternating \texttt{p}-form is a finite combination of wedges of coordinate projections. The coefficients
in that expansion are continuous linear functionals of the form, so a form depending analytically
on a parameter has analytic coefficients.

\begin{definition}[ContinuousAlternatingMap.wedgeCovectors]\label{def:ContinuousAlternatingMap.wedgeCovectors}
\lean{ContinuousAlternatingMap.wedgeCovectors}\leanok
The wedge of \texttt{p} continuous linear functionals, built by iterated alternatizing uncurry.

\noindent\emph{Lean:} \texttt{def ContinuousAlternatingMap.wedgeCovectors}, \href{https://github.com/Paul-Lez/HodgeConjecture/blob/main/HodgeConjecture/Mathlib/Analysis/NormedSpace/WedgeCovectors.lean#L46}{\texttt{HodgeConjecture/Mathlib/Analysis/NormedSpace/WedgeCovectors.lean:46}}.
\end{definition}

\begin{definition}[ContinuousAlternatingMap.standardVolumeForm]\label{def:ContinuousAlternatingMap.standardVolumeForm}
\lean{ContinuousAlternatingMap.standardVolumeForm}\leanok
\uses{def:ContinuousAlternatingMap.wedgeCovectors}
The standard constant volume form on \texttt{𝕜\^{}p}.

\noindent\emph{Lean:} \texttt{def ContinuousAlternatingMap.standardVolumeForm}, \href{https://github.com/Paul-Lez/HodgeConjecture/blob/main/HodgeConjecture/Mathlib/Analysis/NormedSpace/WedgeCovectors.lean#L134}{\texttt{HodgeConjecture/Mathlib/Analysis/NormedSpace/WedgeCovectors.lean:134}}.
\end{definition}

\begin{definition}[ContinuousAlternatingMap.covectorProduct]\label{def:ContinuousAlternatingMap.covectorProduct}
\lean{ContinuousAlternatingMap.covectorProduct}\leanok
The product of a tuple of coordinate covectors.

\noindent\emph{Lean:} \texttt{def ContinuousAlternatingMap.covectorProduct}, \href{https://github.com/Paul-Lez/HodgeConjecture/blob/main/HodgeConjecture/Mathlib/Analysis/NormedSpace/WedgeCovectors.lean#L168}{\texttt{HodgeConjecture/Mathlib/Analysis/NormedSpace/WedgeCovectors.lean:168}}.
\end{definition}

\begin{definition}[ContinuousAlternatingMap.alternatingFormEvaluation]\label{def:ContinuousAlternatingMap.alternatingFormEvaluation}
\lean{ContinuousAlternatingMap.alternatingFormEvaluation}\leanok
Evaluation at a fixed tuple is a continuous linear functional on continuous alternating
forms.

\noindent\emph{Lean:} \texttt{def ContinuousAlternatingMap.alternatingFormEvaluation}, \href{https://github.com/Paul-Lez/HodgeConjecture/blob/main/HodgeConjecture/Mathlib/Analysis/NormedSpace/WedgeCovectors.lean#L269}{\texttt{HodgeConjecture/Mathlib/Analysis/NormedSpace/WedgeCovectors.lean:269}}.
\end{definition}

\begin{definition}[ContinuousAlternatingMap.coordinateCoefficient]\label{def:ContinuousAlternatingMap.coordinateCoefficient}
\lean{ContinuousAlternatingMap.coordinateCoefficient}\leanok
The coefficient of an alternating form field in its finite coordinate-wedge expansion.

\noindent\emph{Lean:} \texttt{def ContinuousAlternatingMap.coordinateCoefficient}, \href{https://github.com/Paul-Lez/HodgeConjecture/blob/main/HodgeConjecture/Mathlib/Analysis/NormedSpace/WedgeCovectors.lean#L275}{\texttt{HodgeConjecture/Mathlib/Analysis/NormedSpace/WedgeCovectors.lean:275}}.
\end{definition}

\section{The radial homotopy operator on differential forms}\label{sec:HodgeConjecture.Mathlib.Analysis.Calculus.DifferentialForm.Poincare}
\noindent\emph{File:} \href{https://github.com/Paul-Lez/HodgeConjecture/blob/main/HodgeConjecture/Mathlib/Analysis/Calculus/DifferentialForm/Poincare.lean}{\texttt{HodgeConjecture/Mathlib/Analysis/Calculus/DifferentialForm/Poincare.lean}}.

This file develops the radial homotopy operator used in the Poincaré lemma.  For a differential
\texttt{(n + 1)}-form \texttt{ω} on a complex normed vector space, its radial contraction is

\texttt{Hω(x) = ∫ t in 0..1, t \^{} n • ω(t • x)(x, ·) dt}.

The construction is explicit.  In particular, exactness will not be introduced as an assumption.

\begin{definition}[DifferentialForm.radialIntegrand]\label{def:DifferentialForm.radialIntegrand}
\lean{DifferentialForm.radialIntegrand}\leanok
The integrand in the radial homotopy operator on an \texttt{(n + 1)}-form.

\noindent\emph{Lean:} \texttt{def DifferentialForm.radialIntegrand}, \href{https://github.com/Paul-Lez/HodgeConjecture/blob/main/HodgeConjecture/Mathlib/Analysis/Calculus/DifferentialForm/Poincare.lean#L179}{\texttt{HodgeConjecture/Mathlib/Analysis/Calculus/DifferentialForm/Poincare.lean:179}}.
\end{definition}

\begin{definition}[DifferentialForm.radialHomotopy]\label{def:DifferentialForm.radialHomotopy}
\lean{DifferentialForm.radialHomotopy}\leanok
\uses{def:DifferentialForm.radialIntegrand}
The radial homotopy operator, obtained by integrating contraction with the radial vector
field along scalar dilations.

\noindent\emph{Lean:} \texttt{def DifferentialForm.radialHomotopy}, \href{https://github.com/Paul-Lez/HodgeConjecture/blob/main/HodgeConjecture/Mathlib/Analysis/Calculus/DifferentialForm/Poincare.lean#L184}{\texttt{HodgeConjecture/Mathlib/Analysis/Calculus/DifferentialForm/Poincare.lean:184}}.
\end{definition}

\begin{definition}[DifferentialForm.radialIntegrandFDeriv]\label{def:DifferentialForm.radialIntegrandFDeriv}
\lean{DifferentialForm.radialIntegrandFDeriv}\leanok
The Fréchet derivative of the radial integrand with respect to its space variable.

\noindent\emph{Lean:} \texttt{def DifferentialForm.radialIntegrandFDeriv}, \href{https://github.com/Paul-Lez/HodgeConjecture/blob/main/HodgeConjecture/Mathlib/Analysis/Calculus/DifferentialForm/Poincare.lean#L190}{\texttt{HodgeConjecture/Mathlib/Analysis/Calculus/DifferentialForm/Poincare.lean:190}}.
\end{definition}

\begin{definition}[DifferentialForm.curryDerivativeAt]\label{def:DifferentialForm.curryDerivativeAt}
\lean{DifferentialForm.curryDerivativeAt}\leanok
The part of the derivative of a contraction that differentiates the alternating form.

\noindent\emph{Lean:} \texttt{def DifferentialForm.curryDerivativeAt}, \href{https://github.com/Paul-Lez/HodgeConjecture/blob/main/HodgeConjecture/Mathlib/Analysis/Calculus/DifferentialForm/Poincare.lean#L451}{\texttt{HodgeConjecture/Mathlib/Analysis/Calculus/DifferentialForm/Poincare.lean:451}}.
\end{definition}

\begin{definition}[DifferentialForm.radialPullback]\label{def:DifferentialForm.radialPullback}
\lean{DifferentialForm.radialPullback}\leanok
The pullback of an \texttt{(n + 1)}-form by scalar dilation.

\noindent\emph{Lean:} \texttt{def DifferentialForm.radialPullback}, \href{https://github.com/Paul-Lez/HodgeConjecture/blob/main/HodgeConjecture/Mathlib/Analysis/Calculus/DifferentialForm/Poincare.lean#L514}{\texttt{HodgeConjecture/Mathlib/Analysis/Calculus/DifferentialForm/Poincare.lean:514}}.
\end{definition}

\begin{definition}[DifferentialForm.radialPullbackDeriv]\label{def:DifferentialForm.radialPullbackDeriv}
\lean{DifferentialForm.radialPullbackDeriv}\leanok
The derivative of \texttt{radialPullback} in its real dilation parameter.

\noindent\emph{Lean:} \texttt{def DifferentialForm.radialPullbackDeriv}, \href{https://github.com/Paul-Lez/HodgeConjecture/blob/main/HodgeConjecture/Mathlib/Analysis/Calculus/DifferentialForm/Poincare.lean#L519}{\texttt{HodgeConjecture/Mathlib/Analysis/Calculus/DifferentialForm/Poincare.lean:519}}.
\end{definition}

\begin{definition}[DifferentialForm.zeroFormCoeff]\label{def:DifferentialForm.zeroFormCoeff}
\lean{DifferentialForm.zeroFormCoeff}\leanok
The scalar coefficient of a differential \texttt{0}-form.

\noindent\emph{Lean:} \texttt{def DifferentialForm.zeroFormCoeff}, \href{https://github.com/Paul-Lez/HodgeConjecture/blob/main/HodgeConjecture/Mathlib/Analysis/Calculus/DifferentialForm/Poincare.lean#L779}{\texttt{HodgeConjecture/Mathlib/Analysis/Calculus/DifferentialForm/Poincare.lean:779}}.
\end{definition}

\section{Normal coordinates from a split derivative}\label{sec:HodgeConjecture.Mathlib.Analysis.Calculus.SplitDerivativeNormalChart}
\noindent\emph{File:} \href{https://github.com/Paul-Lez/HodgeConjecture/blob/main/HodgeConjecture/Mathlib/Analysis/Calculus/SplitDerivativeNormalChart.lean}{\texttt{HodgeConjecture/Mathlib/Analysis/Calculus/SplitDerivativeNormalChart.lean}}.

A map with a split injective strict derivative extends to an actual local coordinate map
by adding vectors in the kernel of a left inverse. This is an application of the inverse
function theorem, not an assumed flattening equivalence. The zero-normal slice is exactly
the original parametrization on the constructed coordinate domain. Identifying that slice
with an entire geometric support additionally requires a local embedding assertion.

\begin{definition}[ContinuousLinearMap.splitKernelEquiv]\label{def:ContinuousLinearMap.splitKernelEquiv}
\lean{ContinuousLinearMap.splitKernelEquiv}\leanok
A specified left inverse splits the ambient vector space into the original tangent
space and the actual kernel of that left inverse.

\noindent\emph{Lean:} \texttt{def ContinuousLinearMap.splitKernelEquiv}, \href{https://github.com/Paul-Lez/HodgeConjecture/blob/main/HodgeConjecture/Mathlib/Analysis/Calculus/SplitDerivativeNormalChart.lean#L33}{\texttt{HodgeConjecture/Mathlib/Analysis/Calculus/SplitDerivativeNormalChart.lean:33}}.
\end{definition}

\begin{definition}[HasStrictFDerivAt.normalChart]\label{def:HasStrictFDerivAt.normalChart}
\lean{HasStrictFDerivAt.normalChart}\leanok
\uses{def:ContinuousLinearMap.splitKernelEquiv}
The actual inverse-function-theorem coordinate map \texttt{(v,n) ↦ f(v)+n}.

\noindent\emph{Lean:} \texttt{def HasStrictFDerivAt.normalChart}, \href{https://github.com/Paul-Lez/HodgeConjecture/blob/main/HodgeConjecture/Mathlib/Analysis/Calculus/SplitDerivativeNormalChart.lean#L86}{\texttt{HodgeConjecture/Mathlib/Analysis/Calculus/SplitDerivativeNormalChart.lean:86}}.
\end{definition}

\chapter{Mathlib: CategoryTheory}

\section{The short exact sequence of kernels of a composition with an epimorphism}\label{sec:HodgeConjecture.Mathlib.CategoryTheory.Abelian.KernelCompositionShortExact}
\noindent\emph{File:} \href{https://github.com/Paul-Lez/HodgeConjecture/blob/main/HodgeConjecture/Mathlib/CategoryTheory/Abelian/KernelCompositionShortExact.lean}{\texttt{HodgeConjecture/Mathlib/CategoryTheory/Abelian/KernelCompositionShortExact.lean}}.

For \texttt{A → B → C}, restriction to kernels gives
\texttt{0 → ker f → ker (f ≫ g) → ker g → 0} when \texttt{f} is an epimorphism.
The maps are the actual kernel maps, with their inclusion normalizations retained.

\begin{definition}[CategoryTheory.kernelCompositionShortComplex]\label{def:CategoryTheory.kernelCompositionShortComplex}
\lean{CategoryTheory.kernelCompositionShortComplex}\leanok
The canonical sequence of the three kernels in a composable pair.

\noindent\emph{Lean:} \texttt{abbrev CategoryTheory.kernelCompositionShortComplex}, \href{https://github.com/Paul-Lez/HodgeConjecture/blob/main/HodgeConjecture/Mathlib/CategoryTheory/Abelian/KernelCompositionShortExact.lean#L27}{\texttt{HodgeConjecture/Mathlib/CategoryTheory/Abelian/KernelCompositionShortExact.lean:27}}.
\end{definition}

\begin{definition}[CategoryTheory.kernelFactorizationShortComplex]\label{def:CategoryTheory.kernelFactorizationShortComplex}
\lean{CategoryTheory.kernelFactorizationShortComplex}\leanok
The same sequence with a specified composite, preserving the actual kernel objects.

\noindent\emph{Lean:} \texttt{def CategoryTheory.kernelFactorizationShortComplex}, \href{https://github.com/Paul-Lez/HodgeConjecture/blob/main/HodgeConjecture/Mathlib/CategoryTheory/Abelian/KernelCompositionShortExact.lean#L53}{\texttt{HodgeConjecture/Mathlib/CategoryTheory/Abelian/KernelCompositionShortExact.lean:53}}.
\end{definition}

\begin{definition}[CategoryTheory.kernelCompositionFactorizationIso]\label{def:CategoryTheory.kernelCompositionFactorizationIso}
\lean{CategoryTheory.kernelCompositionFactorizationIso}\leanok
\uses{def:CategoryTheory.kernelCompositionShortComplex, def:CategoryTheory.kernelFactorizationShortComplex}
Replacing the composite by its proved equality is a canonical kernel isomorphism.

\noindent\emph{Lean:} \texttt{def CategoryTheory.kernelCompositionFactorizationIso}, \href{https://github.com/Paul-Lez/HodgeConjecture/blob/main/HodgeConjecture/Mathlib/CategoryTheory/Abelian/KernelCompositionShortExact.lean#L64}{\texttt{HodgeConjecture/Mathlib/CategoryTheory/Abelian/KernelCompositionShortExact.lean:64}}.
\end{definition}

\begin{definition}[CategoryTheory.kernelFactorizationShortComplexMapIso]\label{def:CategoryTheory.kernelFactorizationShortComplexMapIso}
\lean{CategoryTheory.kernelFactorizationShortComplexMapIso}\leanok
\uses{def:CategoryTheory.kernelFactorizationShortComplex}
Kernel-preserving functors identify the actual sequence with that of the mapped
restriction maps. This does not require them to preserve epimorphisms in general.

\noindent\emph{Lean:} \texttt{def CategoryTheory.kernelFactorizationShortComplexMapIso}, \href{https://github.com/Paul-Lez/HodgeConjecture/blob/main/HodgeConjecture/Mathlib/CategoryTheory/Abelian/KernelCompositionShortExact.lean#L87}{\texttt{HodgeConjecture/Mathlib/CategoryTheory/Abelian/KernelCompositionShortExact.lean:87}}.
\end{definition}

